% !TeX program = pdflatex
% Samodzielny skrypt, wersja po audycie globalnym. Wersja do książki: grafy_i_wezly_tresc.tex
\documentclass[11pt,a4paper]{report}
\usepackage[T1]{fontenc}
\usepackage[utf8]{inputenc}
\usepackage[polish]{babel}
\usepackage{lmodern}
\usepackage{microtype}
\usepackage[margin=2.45cm]{geometry}
\usepackage{amsmath,amssymb,amsthm,mathtools}
\usepackage{enumitem}
\usepackage{xcolor}
\usepackage{tikz}
\usetikzlibrary{arrows.meta,calc,positioning}
\usepackage[unicode]{hyperref}
\hypersetup{colorlinks=true,linkcolor=blue!45!black,urlcolor=blue!45!black}
\setlength{\parindent}{0pt}
\setlength{\parskip}{4pt}
\setlength{\emergencystretch}{2em}
\setlist[enumerate]{topsep=3pt,itemsep=2pt}
\numberwithin{equation}{section}
\theoremstyle{plain}
\newtheorem{twierdzenie}{Twierdzenie}[section]
\newtheorem{lemat}[twierdzenie]{Lemat}
\newtheorem{stwierdzenie}[twierdzenie]{Stwierdzenie}
\newtheorem{wniosek}[twierdzenie]{Wniosek}
\theoremstyle{definition}
\newtheorem{definicja}[twierdzenie]{Definicja}
\newtheorem{przyklad}[twierdzenie]{Przykład}
\theoremstyle{remark}
\newtheorem{uwaga}[twierdzenie]{Uwaga}
\renewcommand{\proofname}{Dowód}
\begin{document}
\hypersetup{pageanchor=false}
\begin{titlepage}
\centering
\vspace*{3.6cm}
{\Huge\bfseries Teoria grafów i węzłów\par}
\vspace{0.45cm}
{\Large w topologii rozmaitości\par}
\vspace{0.9cm}
{\Large\color{blue!45!black} Od teorii grafów do chirurgii i rozmaitości czterowymiarowych\par}
\vfill
{\large Wersja zintegrowana, wersja po audycie globalnym — listopad 2026\par}
\vspace{1cm}
{\large Samodzielny skrypt matematyczny\par}
\end{titlepage}
\pagenumbering{roman}
\tableofcontents
\clearpage
\pagenumbering{arabic}
\hypersetup{pageanchor=true}
% Wstaw przez \input w książce posiadającej pakiety i środowiska zgodne z dotychczasowym skryptem.
\definecolor{tgNavy}{RGB}{30,52,88}
\definecolor{tgBlue}{RGB}{36,100,165}
\definecolor{tgTeal}{RGB}{18,140,132}
\definecolor{tgCoral}{RGB}{211,85,71}
\definecolor{tgGold}{RGB}{192,132,26}
\definecolor{tgPurple}{RGB}{123,83,160}
\tikzset{
  tgVertex/.style={circle,draw=tgNavy!80,fill=white,line width=.8pt,
      minimum size=6mm,inner sep=0pt,font=\small},
  tgBlueVertex/.style={tgVertex,draw=tgBlue,fill=tgBlue!13},
  tgRedVertex/.style={tgVertex,draw=tgCoral,fill=tgCoral!12},
  tgTealVertex/.style={tgVertex,draw=tgTeal,fill=tgTeal!12},
  tgEdge/.style={draw=tgNavy!82,line width=1.25pt},
  tgAccent/.style={draw=tgCoral,line width=2.5pt},
  tgTree/.style={draw=tgTeal,line width=2pt},
  tgGhost/.style={draw=tgNavy!35,densely dashed,line width=1.15pt}}

\chapter*{Jak czytać wspólny skrypt}
\addcontentsline{toc}{chapter}{Jak czytać wspólny skrypt}
Najpierw wprowadzamy grafy, związane z nimi macierze i realizacje topologiczne.
Następnie przechodzimy do geometrii rozmaitości, działań grup,
węzłów, chirurgii i topologii niskowymiarowej.
Wszystkie diagramy są edytowalne w TikZ. Dowody możliwe do przeprowadzenia
przy pomocy poznanego aparatu rozwijamy bez ukrywania istotnych kroków;
twierdzenia wymagające obszernej teorii dodatkowej oznaczamy jako
wyniki przyjęte z literatury. Przy rezultatach korzystających z teorii
indeksu, zwartości lub transwersalności rozróżniamy rachunek formalny,
skonstruowanie obiektu i zastosowanie zewnętrznego twierdzenia.
W całym tekście używamy polskiego terminu ``zanurzenie'' w znaczeniu
\emph{embedding}, odróżniając go od immersji.

\chapter{Podstawy teorii grafów}
\label{ch:tg-podstawy}

Teoria grafów bada układy obiektów i relacji między nimi. Wierzchołki będą
reprezentować obiekty, a krawędzie ich powiązania. Już na poziomie tego
prostego opisu pojawiają się trzy różne perspektywy: kombinatoryczna
(liczenie połączeń i dróg), algebraiczna (macierze) i topologiczna
(realizacja grafu jako przestrzeni jednowymiarowej). Pierwsze dwa rozdziały
poświęcamy części kombinatorycznej, nie zakładając znajomości żadnego
pojęcia z teorii grafów.

\section{Graf prosty, multigraf i realizacja geometryczna}
\label{sec:tg-obiekty}

\begin{definicja}[Graf prosty]
\emph{Skończonym grafem prostym nieskierowanym} nazywamy parę
$G=(V,E)$, gdzie $V$ jest niepustym, skończonym zbiorem
\emph{wierzchołków}, natomiast
$E\subseteq\binom V2=\{\{u,v\}:u,v\in V,\ u\ne v\}$
jest zbiorem \emph{krawędzi}. Zapis $uv\in E$ oznacza
$\{u,v\}\in E$. Wierzchołki $u,v$ są \emph{sąsiadujące},
gdy $uv\in E$; wierzchołek $u$ jest wówczas \emph{incydentny}
z krawędzią $uv$.
\end{definicja}

Z definicji żadna krawędź nie łączy wierzchołka z nim samym i żadna para
różnych wierzchołków nie jest połączona więcej niż jedną krawędzią.
Ważne jest jednak, by nie utożsamiać pojęcia grafu z jego rysunkiem:
przecięcia linii na rysunku nie są nowymi wierzchołkami, o ile ich
nie zaznaczono. Dwa różne rysunki mogą opisywać dokładnie ten sam graf.

\begin{definicja}[Multigraf z pętlami]
\emph{Skończonym multigrafem nieskierowanym z pętlami}
nazywamy trójkę $G=(V,E,\partial)$, w której $V$ i $E$ są
skończonymi zbiorami, a $\partial:E\to (V\times V)/\mathfrak S_2$
przyporządkowuje każdej krawędzi nieuporządkowaną parę
\emph{z zachowaniem krotności}. Zapis $\partial(e)=[(u,v)]$
oznacza orbitę pary $(u,v)$ przy zamianie współrzędnych;
w szczególności $[(u,u)]$ jest poprawną parą równych końców.
Jeśli $u=v$, krawędź $e$ nazywamy \emph{pętlą}.
Różne elementy $e,f\in E$ mogą mieć te same końce;
wówczas są \emph{krawędziami równoległymi}.
\end{definicja}

W całej pierwszej części, jeśli nie podano inaczej, \emph{graf} oznacza
graf prosty nieskierowany i skończony. Kiedy mówimy o multigrafie,
powiemy to jawnie. Później multigrafy będą potrzebne m.in.
przy diagramach węzłów i grafach dualnych triangulacji.

\begin{center}
\begin{tikzpicture}[scale=.91,line cap=round]
  \begin{scope}[shift={(-3.05,0)}]
    \node[tgBlueVertex] (a) at (0,0) {$a$};
    \node[tgBlueVertex] (b) at (2.0,0) {$b$};
    \node[tgBlueVertex] (c) at (1,1.55) {$c$};
    \draw[tgEdge] (a)--(b)--(c)--(a);
    \node[font=\small,text=tgNavy] at (1,-.70) {graf prosty};
  \end{scope}
  \begin{scope}[shift={(2.2,0)}]
    \node[tgRedVertex] (u) at (0,0) {$u$};
    \node[tgRedVertex] (v) at (2.0,0) {$v$};
    \draw[tgEdge] (u) to[bend left=26] (v);
    \draw[tgAccent] (u) to[bend right=26] (v);
    \draw[tgEdge] (u) .. controls (-1.10,1.15) and (-1.10,-.38) .. (u);
    \node[font=\small,text=tgNavy,align=center] at (1,-.85)
      {multigraf z pętlą};
  \end{scope}
\end{tikzpicture}

{\footnotesize Po lewej $E=\{ab,bc,ca\}$. Po prawej są dwie różne
krawędzie o końcach $u,v$ oraz pętla przy $u$.}
\end{center}

\begin{definicja}[Graf pełny i graf bez krawędzi]
Dla liczby całkowitej $n\ge 1$ przez $K_n$ oznaczamy graf
\emph{pełny} o wierzchołkach $\{1,\ldots,n\}$ i wszystkich
możliwych krawędziach. Graf \emph{bez krawędzi} na $n$
wierzchołkach ma zbiór krawędzi równy $\varnothing$;
można go oznaczyć przez $\overline{K_n}$.
Od grafu bez krawędzi odróżniamy \emph{graf zerowy}
$G_0=(\varnothing,\varnothing)$. W pierwszych rozdziałach,
zgodnie z konwencją definicji grafu prostego, rozważamy
grafy o niepustym zbiorze wierzchołków; $G_0$ pojawia się
jedynie przy operacjach usuwania i w konwencjach brzegowych.
\end{definicja}

\begin{przyklad}[Liczenie krawędzi grafu pełnego]
Każda z $n$ etykiet sąsiaduje w $K_n$ z $n-1$ pozostałymi.
Licząc uporządkowane pary końców, dostajemy $n(n-1)$,
ale każdą krawędź policzyliśmy dwukrotnie. Stąd
$|E(K_n)|=\binom n2=n(n-1)/2$.
Na przykład $K_4$ ma sześć krawędzi. Rysując go na płaszczyźnie
z dwoma przecinającymi się przekątnymi, nie tworzymy z ich
przecięcia piątego wierzchołka.
\end{przyklad}

\begin{definicja}[Izomorfizm]
Niech $G=(V,E)$ i $H=(W,F)$ będą grafami prostymi.
\emph{Izomorfizmem grafów} $G\to H$ nazywamy bijekcję
$f:V\to W$, dla której
$uv\in E\iff f(u)f(v)\in F$ dla wszystkich różnych $u,v\in V$.
Piszemy wtedy $G\cong H$.
\end{definicja}

Izomorfizm zmienia jedynie nazwy wierzchołków. Liczba wierzchołków,
liczba krawędzi, stopnie, cykle oraz liczba składowych spójnych
muszą zatem zgadzać się w grafach izomorficznych. Zgodność tych
wielkości nie wystarcza jednak zwykle do izomorfizmu.

\begin{przyklad}[Niezgodność wykryta przez trójkąty]
Graf graniastosłupa trójkątnego i graf pełny dwudzielny $K_{3,3}$
mają po sześć wierzchołków, dziewięć krawędzi i wszystkie stopnie
równe trzy. Pierwszy zawiera cykl długości trzy, natomiast drugi
nie zawiera żadnego cyklu nieparzystego. Nie są więc izomorficzne.
To przykład, dlaczego sam ciąg stopni nie klasyfikuje grafów.
\end{przyklad}

\begin{definicja}[Realizacja geometryczna]
Dla grafu prostego $G=(V,E)$ rozważmy przestrzeń
$V\sqcup\bigsqcup_{e\in E}[0,1]_e$ i utożsammy
końce odcinka $[0,1]_e$ z dwoma wierzchołkami
należącymi do krawędzi $e$. Otrzymany iloraz
$|G|$ nazywamy \emph{realizacją geometryczną} grafu.
\end{definicja}

Realizacja jest jednowymiarowym kompleksem CW. Jej definicja nie
zależy od tego, jak rysujemy graf na kartce. W szczególności
samo przecinanie się krzywych na rysunku nie prowadzi do
utożsamienia odpowiednich punktów w $|G|$.

\section{Stopnie wierzchołków i podstawowe rachunki}
\label{sec:tg-stopnie}

\begin{definicja}[Stopień]
Dla grafu prostego $G=(V,E)$ i $v\in V$ przez
$N_G(v)=\{u\in V:uv\in E\}$ oznaczamy
\emph{sąsiedztwo} $v$, a liczbę $d_G(v)=|N_G(v)|$
nazywamy \emph{stopniem} $v$. W multigrafie stopień
jest liczbą końców krawędzi incydentnych z danym
wierzchołkiem; każda pętla wnosi do stopnia dwa.
\end{definicja}

W grafie prostym zachodzi $0\le d_G(v)\le |V|-1$.
Stopień zero oznacza wierzchołek izolowany, a stopień jeden
--- liść. Dla skończonego niepustego grafu definiujemy
$\delta(G)=\min_{v\in V}d_G(v)$ i
$\Delta(G)=\max_{v\in V}d_G(v)$.

\begin{twierdzenie}[Lemat o uściskach dłoni]
Niech $G=(V,E)$ będzie skończonym multigrafem
nieskierowanym, z dopuszczeniem pętli, a $d_G(v)$
niech liczy każdą pętlę dwukrotnie. Wówczas
\[
\sum_{v\in V}d_G(v)=2|E|.
\]
W szczególności liczba wierzchołków nieparzystego
stopnia jest parzysta, także gdy multigraf ma pętle.
\end{twierdzenie}

\begin{proof}
Rozważmy zbiór incydencji, którego elementami są pary
$(v,h)$, gdzie $h$ jest końcem pewnej krawędzi przylegającym
do $v$. Dla każdej krawędzi są dokładnie dwa takie końce,
również wtedy, gdy jest ona pętlą. Liczba wszystkich
incydencji wynosi zatem $2|E|$. Jeśli natomiast
najpierw pogrupujemy incydencje według wierzchołków,
to przy $v$ otrzymamy dokładnie $d_G(v)$ końców.
Oba sposoby zliczania dają pierwszą równość.

Niech $V_{\mathrm{odd}}$ będzie zbiorem wierzchołków
nieparzystego stopnia. Redukując uzyskaną równość modulo
$2$, otrzymujemy
$0\equiv\sum_{v\in V}d_G(v)\equiv
|V_{\mathrm{odd}}|\pmod2$, co dowodzi drugiego twierdzenia.
\end{proof}

\begin{przyklad}[Obliczenie stopni]
Weźmy graf o wierzchołkach $a,b,c,d,e$ i krawędziach
$ab,ac,ad,bc,cd$. Wówczas
$(d(a),d(b),d(c),d(d),d(e))=(3,2,3,2,0)$.
Suma stopni wynosi $10=2\cdot5$. Dokładnie dwa
wierzchołki mają nieparzysty stopień. Wierzchołek $e$
jest izolowany.
\end{przyklad}

\begin{center}
\begin{tikzpicture}[scale=.96]
  \node[tgBlueVertex] (a) at (0,1.4) {$a$};
  \node[tgRedVertex] (b) at (-1.8,0) {$b$};
  \node[tgRedVertex] (c) at (1.8,0) {$c$};
  \node[tgBlueVertex] (d) at (0,-1.25) {$d$};
  \node[tgTealVertex] (e) at (3.5,1.05) {$e$};
  \draw[tgEdge] (a)--(b)--(c)--(a);
  \draw[tgEdge] (a)--(d)--(c);
  \node[font=\footnotesize,text=tgNavy] at (3.5,.45) {$d(e)=0$};
\end{tikzpicture}

{\footnotesize Przykład z pięcioma krawędziami; izolowany
wierzchołek nie wpływa na sumę stopni.}
\end{center}

\begin{wniosek}
Nie istnieje graf z dokładnie jednym wierzchołkiem
nieparzystego stopnia. Jeżeli graf prosty o $n$
wierzchołkach jest $r$-regularny (czyli każdy wierzchołek
ma stopień $r$), to $nr$ jest parzyste i
$|E|=nr/2$.
\end{wniosek}

\begin{proof}
Obie własności wynikają bezpośrednio z poprzedniego
twierdzenia: odpowiednio z parzystości liczby
wierzchołków nieparzystego stopnia i ze wzoru
$\sum_{v\in V}d(v)=nr=2|E|$.
\end{proof}

\section{Podgrafy, dopełnienie i macierz sąsiedztwa}
\label{sec:tg-podgrafy}

\begin{definicja}[Podgraf i podgraf indukowany]
Niech $G=(V,E)$ będzie grafem prostym. Graf
$H=(W,F)$ jest jego \emph{podgrafem}, jeśli
$W\subseteq V$ i $F\subseteq E\cap\binom W2$.
Dla $W\subseteq V$, $W\ne\varnothing$, graf
$G[W]=(W,E\cap\binom W2)$ nazywamy
\emph{podgrafem indukowanym przez $W$}.
\end{definicja}

W podgrafie można usuwać zarówno wierzchołki, jak i część
krawędzi pomiędzy pozostałymi wierzchołkami. W podgrafie
indukowanym zachowujemy \emph{wszystkie} krawędzie między
wybranymi wierzchołkami. Notacja $G-e$ oznacza usunięcie
samej krawędzi $e$; notacja $G-v$ usuwa wierzchołek $v$
wraz z wszystkimi jego krawędziami.

\begin{definicja}[Dopełnienie]
Dla grafu prostego $G=(V,E)$ jego
\emph{dopełnienie} $\overline G$ jest grafem na
tym samym zbiorze wierzchołków, w którym
$uv\in E(\overline G)$ wtedy i tylko wtedy,
gdy $u\ne v$ oraz $uv\notin E(G)$.
\end{definicja}

\begin{stwierdzenie}
Dla grafu prostego $G$ o $n$ wierzchołkach oraz
każdego $v\in V(G)$ zachodzi
$d_{\overline G}(v)=n-1-d_G(v)$, a ponadto
$|E(\overline G)|=\binom n2-|E(G)|$.
\end{stwierdzenie}

\begin{proof}
Spośród pozostałych $n-1$ wierzchołków dokładnie
$d_G(v)$ jest sąsiadami $v$ w $G$. Pozostałe
$n-1-d_G(v)$ są i tylko one sąsiadami $v$ w
$\overline G$. Każda nieuporządkowana para różnych
wierzchołków należy dokładnie do jednego ze zbiorów
$E(G)$, $E(\overline G)$, zatem sumy ich mocy dają
$\binom n2$.
\end{proof}

\begin{definicja}[Macierz sąsiedztwa]
Niech $G=(V,E)$ będzie grafem prostym o
$V=\{v_1,\ldots,v_n\}$. \emph{Macierzą sąsiedztwa}
nazywamy $A_G=(a_{ij})\in M_n(\mathbb Z)$,
gdzie $a_{ij}=1$, jeśli $v_iv_j\in E$, oraz
$a_{ij}=0$ w przeciwnym przypadku.
\end{definicja}

Macierz $A_G$ jest symetryczna i ma zerową przekątną.
Sumy wierszy to stopnie wierzchołków.

\begin{przyklad}[Ścieżka na czterech wierzchołkach]
Niech $V=\{1,2,3,4\}$, a $E=\{12,23,34\}$.
Wtedy
\[
A_G=\begin{pmatrix}
0&1&0&0\\1&0&1&0\\0&1&0&1\\0&0&1&0
\end{pmatrix}.
\]
Mnożenie $A_G(1,1,1,1)^T$ daje $(1,2,2,1)^T$,
a więc wektor stopni. Uogólnienie tej obserwacji
poprowadzi nas później do laplasjanu grafu i
jego widma.
\end{przyklad}

\begin{stwierdzenie}[Zliczanie spacerów macierzami]
Niech $G$ będzie grafem prostym o macierzy sąsiedztwa
$A_G$. Dla liczby całkowitej $k\ge0$ element
$(A_G^k)_{ij}$ jest równy liczbie ciągów
$(v_i=v_{i_0},v_{i_1},\ldots,v_{i_k}=v_j)$,
w których każda kolejna para jest krawędzią grafu.
Dla $k=0$ liczymy ciąg zerowej długości tylko wtedy,
gdy $i=j$.
\end{stwierdzenie}

\begin{proof}
Dla $k=0$ twierdzenie jest równością $A_G^0=I$.
Dla $k=1$ jest definicją macierzy sąsiedztwa.
Załóżmy, że jest prawdziwe dla $k$. Z definicji
mnożenia macierzy
$(A_G^{k+1})_{ij}=\sum_{r=1}^n(A_G^k)_{ir}a_{rj}$.
Każdy ciąg długości $k+1$ kończący się w $v_j$
ma dokładnie jeden przedostatni wierzchołek $v_r$.
Wybór poprzednich $k$ kroków daje $(A_G^k)_{ir}$
możliwości, a ostatni krok istnieje dokładnie przy
$a_{rj}=1$. Sumowanie po $r$ zlicza wszystkie
ciągi dokładnie raz. Indukcja jest zakończona.
\end{proof}

Ważna subtelność: takie ciągi mogą wielokrotnie odwiedzać
ten sam wierzchołek lub tę samą krawędź. Nie są więc z zasady
drogami prostymi ani cyklami.

\section{Spacery, szlaki, ścieżki i cykle}
\label{sec:tg-drogi}

\begin{definicja}[Spacer, szlak i ścieżka]
Niech $G=(V,E)$ będzie grafem prostym.
\emph{Spacerem} długości $k\ge0$ z $v_0$ do $v_k$
nazywamy ciąg $v_0,v_1,\ldots,v_k$ taki, że
$v_{i-1}v_i\in E$ dla $i=1,\ldots,k$.
Spacer nazywamy \emph{szlakiem}, jeśli żadna
krawędź nie jest przechodzona dwukrotnie.
Nazywamy go \emph{ścieżką} (\emph{drogą prostą}), jeśli
wierzchołki $v_0,\ldots,v_k$ są parami różne.
\emph{Spacer zamknięty} ma $v_0=v_k$.
\end{definicja}

Gdy w dalszym tekście pojawi się słowo ,,droga'' w znaczeniu
kombinatorycznym, będziemy mieli na myśli drogę prostą. W innych kontekstach matematycznych
,,droga'' może oznaczać dowolne ciągłe odwzorowanie odcinka;
rozróżnienie to jest istotne.

\begin{definicja}[Cykl]
\emph{Cyklem} długości $k\ge3$ w grafie prostym
nazywamy spacer zamknięty $v_0,v_1,\ldots,v_k=v_0$,
w którym $v_0,\ldots,v_{k-1}$ są parami różne.
Dwa zapisy powstałe przez zmianę punktu startowego
lub kierunku obiegu wyznaczają ten sam cykl.
Graf \emph{acykliczny} nie zawiera cyklu.
\end{definicja}

W multigrafach analogiczna definicja dopuszcza
cykle długości jeden (pętle) i dwa (dwie równoległe
krawędzie); te dwa przypadki zawsze podamy jawnie.

\begin{lemat}[Usuwanie zamkniętych fragmentów spaceru]
Niech $G$ będzie grafem skończonym i niech
$v_0,\ldots,v_k$ będzie spacerem z $u=v_0$
do $v=v_k$, gdzie $u\ne v$. Wtedy w $G$
istnieje droga prosta z $u$ do $v$ długości
nie większej niż $k$.
\end{lemat}

\begin{proof}
Jeśli w spacerze żaden wierzchołek się nie powtarza,
to jest on już drogą prostą. W przeciwnym razie istnieją
$i<j$ takie, że $v_i=v_j$. Usuwamy fragment od
$v_{i+1}$ do $v_j$: otrzymujemy spacer
$v_0,\ldots,v_i,v_{j+1},\ldots,v_k$
(jeśli $j=k$, spacer kończy się w $v_i=v$).
Zachowuje on oba końce i jest krótszy o $j-i$ kroków.
Powtarzamy tę operację. Długość jest nieujemną liczbą
całkowitą i za każdym razem maleje, więc proces
musi się zakończyć. Pozostały spacer nie powtarza
wierzchołków, czyli jest drogą prostą.
\end{proof}

\begin{uwaga}
Nie wolno w powyższym lemacie zastąpić ,,spaceru między
dwoma różnymi końcami'' przez dowolny niepusty spacer
zamknięty i wnioskować, że istnieje cykl. W drzewie
spacer $u,v,u$ jest zamknięty, ale nie zawiera cyklu.
Ten sam fragment krawędzi został przebyty dwa razy.
\end{uwaga}

\section{Spójność i składowe spójne}
\label{sec:tg-spojnosc}

\begin{definicja}[Spójność]
Wierzchołki $u,v$ grafu $G$ nazywamy
\emph{połączonymi}, jeśli $u=v$ albo istnieje
spacer od $u$ do $v$. Graf $G$ jest \emph{spójny},
jeśli każda para jego wierzchołków jest połączona.
Maksymalne względem inkluzji zbiory parami połączonych
wierzchołków, wraz z wszystkimi krawędziami między nimi,
nazywamy \emph{składowymi spójnymi} grafu.
\end{definicja}

\begin{stwierdzenie}
Dla każdego grafu $G=(V,E)$ relacja
$u\sim v\iff u,v$ są połączone jest relacją
równoważności na $V$. Jej klasami są dokładnie
zbiory wierzchołków składowych spójnych.
\end{stwierdzenie}

\begin{proof}
Relacja jest zwrotna, bo dopuszczamy spacer długości zero.
Jest symetryczna, bo każdy spacer można przejść w kierunku
odwrotnym. Jeśli spacer łączy $u$ z $v$, a drugi łączy
$v$ z $w$, ich sklejenie daje spacer z $u$ do $w$,
więc relacja jest przechodnia. Klasy relacji są
parami rozłączne i pokrywają $V$. Jeżeli wierzchołek
spoza danej klasy miałby krawędź do wierzchołka tej klasy,
byłby z nim połączony i sam należałby do klasy. Wobec
tego każda klasa z grafem indukowanym jest składową
spójną i jest maksymalna.
\end{proof}

\begin{przyklad}[Spójny podgraf nie musi być składową]
W cyklu $C_5$ dwa sąsiadujące wierzchołki wraz z
łączącą je krawędzią tworzą spójny podgraf. Nie jest
on jednak składową spójną całego $C_5$, ponieważ
pozostałe trzy wierzchołki można do niego dołączyć,
nie tracąc spójności. Składowa spójna oznacza
\emph{maksymalny} spójny fragment.
\end{przyklad}

\begin{center}
\begin{tikzpicture}[scale=.91]
  \begin{scope}[shift={(-3.0,0)}]
    \node[tgBlueVertex] (a) at (-.9,-.4) {$a$};
    \node[tgBlueVertex] (b) at (1.0,-.4) {$b$};
    \node[tgBlueVertex] (c) at (0,1.0) {$c$};
    \draw[tgEdge] (a)--(b)--(c)--(a);
    \draw[draw=tgBlue!50,rounded corners=10pt,line width=.9pt]
      (-1.48,-.92) rectangle (1.58,1.55);
  \end{scope}
  \begin{scope}[shift={(1.9,0)}]
    \node[tgRedVertex] (d) at (0,0) {$d$};
    \node[tgRedVertex] (e) at (1.65,0) {$e$};
    \draw[tgAccent] (d)--(e);
  \end{scope}
  \node[tgTealVertex] (f) at (4.65,1.2) {$f$};
  \node[font=\footnotesize,text=tgNavy] at (0,-1.5) {trzy składowe spójne};
\end{tikzpicture}
\end{center}

\begin{definicja}[Odległość]
Dla grafu $G$ i $u,v\in V(G)$ definiujemy
$d_G(u,v)$ jako najmniejszą długość drogi
prostej z $u$ do $v$, jeżeli taka droga istnieje,
oraz $d_G(u,v)=\infty$ w przeciwnym przypadku.
Przyjmujemy $d_G(v,v)=0$.
\end{definicja}

\begin{stwierdzenie}
Na zbiorze wierzchołków każdej składowej spójnej
grafu $G$ funkcja $d_G$ jest metryką.
\end{stwierdzenie}

\begin{proof}
Odległości są nieujemne, a odległość zero zachodzi
tylko dla identycznych wierzchołków. Odwrócenie
drogi daje symetrię. Dla $u,v,w$ w tej samej
składowej wybieramy najkrótsze drogi od $u$ do $v$
i od $v$ do $w$. Ich konkatenacja jest spacerem od
$u$ do $w$ długości $d_G(u,v)+d_G(v,w)$.
Na mocy lematu o usuwaniu powtórzeń istnieje droga
prosta nie dłuższa niż ten spacer. Zatem
$d_G(u,w)\le d_G(u,v)+d_G(v,w)$.
\end{proof}

W ten sposób nawet skończony graf otrzymuje naturalną
geometrię metryczną. W późniejszej części porównamy tę
metrykę z metryką długości na realizacji $|G|$.

\section{Grafy dwudzielne i parzystość cykli}
\label{sec:tg-dwudzielne}

\begin{definicja}[Graf dwudzielny]
Graf $G=(V,E)$ jest \emph{dwudzielny}, jeśli istnieją
rozłączne zbiory $V_0,V_1$ o sumie $V$ takie, że
każda krawędź ma jeden koniec w $V_0$, a drugi
w $V_1$. Parę $(V_0,V_1)$ nazywamy \emph{podziałem
dwudzielnym}. \emph{Pełny graf dwudzielny} $K_{r,s}$,
gdzie $r,s\ge1$, łączy wszystkimi możliwymi krawędziami
$r$ wierzchołków jednej klasy z $s$ wierzchołkami
drugiej klasy i nie zawiera innych krawędzi.
\end{definicja}

\begin{twierdzenie}[Kryterium cykli nieparzystych]
Dla skończonego grafu prostego $G$ następujące
warunki są równoważne:
\begin{enumerate}[label=(\roman*)]
\item $G$ jest dwudzielny;
\item $G$ nie zawiera cyklu nieparzystej długości.
\end{enumerate}
\end{twierdzenie}

\begin{proof}
Załóżmy (i). Idąc po kolejnych krawędziach, przy
każdym kroku zmieniamy klasę podziału z $V_0$ na
$V_1$ lub odwrotnie. Powrót do punktu początkowego
może więc nastąpić jedynie po parzystej liczbie
kroków. Żaden cykl nie ma długości nieparzystej.

Załóżmy teraz (ii). Wystarczy skonstruować podział
dla każdej składowej spójnej z osobna. Wybierzmy
wierzchołek $r$ i dla każdego $v$ oznaczmy przez
$\ell(v)=d_G(r,v)$ odległość od $r$. Ustalmy
$V_0=\{v:\ell(v)\text{ jest parzyste}\}$
oraz $V_1=\{v:\ell(v)\text{ jest nieparzyste}\}$.
Przypuśćmy, że istnieje krawędź $uv$ z końcami
o tej samej parzystości odległości. Wybierzmy
drzewo najkrótszych dróg od $r$: dla każdego
$v\ne r$ wybieramy poprzednika o odległości
$\ell(v)-1$. Wszystkie takie krawędzie tworzą
drzewo, gdyż odległość ściśle maleje w kierunku
korzenia, co wyklucza cykle, oraz każdy wierzchołek
dochodzi do $r$.

Niech $w$ będzie ostatnim wspólnym wierzchołkiem
dwóch dróg w tym drzewie od $r$ do $u$ i od $r$
do $v$. Jedyna droga drzewowa od $u$ do $v$
ma długość $\ell(u)+\ell(v)-2\ell(w)$, a więc
parzystą. Krawędź $uv$ nie należy do drzewa,
bo w drzewie łączyłaby poziomy o różnicy jeden.
Dodanie jej do drogi drzewowej daje cykl
nieparzystej długości, sprzeczny z (ii).
Zatem każda krawędź łączy wierzchołki o
różnych parzystościach, co dowodzi (i).
\end{proof}

\begin{center}
\begin{tikzpicture}[scale=.83]
 \begin{scope}[shift={(-2.8,0)}]
  \node[tgBlueVertex] (a) at (-.9,-.75) {$a$};
  \node[tgRedVertex] (b) at (.9,-.75) {$b$};
  \node[tgRedVertex] (c) at (-.9,.9) {$c$};
  \node[tgBlueVertex] (d) at (.9,.9) {$d$};
  \draw[tgEdge] (a)--(b)--(d)--(c)--(a);
  \node[font=\small,text=tgNavy] at (0,-1.35) {cykl $C_4$: dwudzielny};
 \end{scope}
 \begin{scope}[shift={(2.7,0)}]
  \node[tgBlueVertex] (u) at (0,1) {$u$};
  \node[tgRedVertex] (v) at (-1.12,-.8) {$v$};
  \node[tgRedVertex] (w) at (1.12,-.8) {$w$};
  \draw[tgEdge] (u)--(v)--(w)--(u);
  \node[font=\small,text=tgNavy] at (0,-1.35) {$C_3$: brak podziału};
 \end{scope}
\end{tikzpicture}
\end{center}

\begin{przyklad}[Szybkie rozpoznawanie dwudzielności]
Każde drzewo jest dwudzielne, bo nie zawiera żadnego
cyklu. Pełny graf $K_{r,s}$ jest dwudzielny bezpośrednio
z definicji. Dla $n\ge3$ cykl $C_n$ jest dwudzielny
dokładnie wtedy, gdy $n$ jest parzyste. W szczególności
$K_3$ i $K_5$ nie są dwudzielne.
\end{przyklad}

\section{Które ciągi liczb są ciągami stopni?}
\label{sec:tg-havel}

Nie każdy $n$-elementowy ciąg liczb całkowitych z przedziału
$\{0,\ldots,n-1\}$ jest ciągiem stopni grafu prostego.
Warunek parzystości sumy jest konieczny, lecz nie wystarcza.
Na przykład $(3,3,1,1)$ ma parzystą sumę $8$, ale nie jest
ciągiem stopni grafu prostego na czterech wierzchołkach:
dwa wierzchołki stopnia trzy muszą być połączone ze wszystkimi
innymi, więc pozostałe miałyby stopnie co najmniej dwa.

\begin{definicja}[Ciąg grafowy]
Skończony ciąg nieujemnych liczb całkowitych
$(d_1,\ldots,d_n)$ nazywamy \emph{grafowym},
jeśli istnieje graf prosty o wierzchołkach
$v_1,\ldots,v_n$, dla którego $d_G(v_i)=d_i$
dla każdego $i$.
\end{definicja}

\begin{twierdzenie}[Kryterium Havla--Hakimiego]
Niech $n\ge2$ i niech
$d_1\ge d_2\ge\cdots\ge d_n\ge0$ będą liczbami
całkowitymi. Jeśli $d_1>n-1$, ciąg nie jest grafowy.
Jeżeli $0\le d_1\le n-1$, to ciąg
$(d_1,\ldots,d_n)$ jest grafowy wtedy i tylko
wtedy, gdy ciąg
\[
(d_2-1,\ldots,d_{d_1+1}-1,
  d_{d_1+2},\ldots,d_n)
\]
jest grafowy (po ewentualnym przestawieniu elementów).
Dla $d_1=0$ otrzymujemy ciąg samych zer, który jest
grafowy.
\end{twierdzenie}

\begin{proof}
Jeśli $d_1>n-1$, wierzchołek prostego grafu
nie może mieć wymaganego stopnia. Załóżmy więc
$0\le d_1\le n-1$. Oznaczmy $d=d_1$.
W jedną stronę, jeśli istnieje graf na wierzchołkach
$v_2,\ldots,v_n$ o stopniach odpowiednio
$d_2-1,\ldots,d_{d+1}-1,d_{d+2},\ldots,d_n$,
to dołączamy nowy wierzchołek $v_1$ i krawędzie
$v_1v_2,\ldots,v_1v_{d+1}$. Stopnie wszystkich
wierzchołków stają się równe zadanym liczbom.

W trudniejszą stronę załóżmy, że istnieje graf $G$
o stopniach $d_1,\ldots,d_n$. Wykażemy za pomocą
\emph{zamiany dwóch krawędzi} (ang.\ $2$-switch), że
można go zmodyfikować bez zmiany stopni tak,
aby $v_1$ był sąsiadem dokładnie $v_2,\ldots,v_{d+1}$.
Wśród wszystkich realizacji z zadanymi stopniami
wybierzmy taką, w której $v_1$ ma najwięcej
sąsiadów w tym wyróżnionym zbiorze.

Gdyby nie był sąsiadem pewnego $v_j$ z indeksu
$2\le j\le d+1$, musiałby być sąsiadem pewnego
$v_k$ o $k>d+1$. Ponieważ $d_j\ge d_k$, pokażemy,
że da się zamienić te sąsiedztwa. Istnieje
wierzchołek $z\notin\{v_1,v_j,v_k\}$ taki,
że $v_jz\in E(G)$, lecz $v_kz\notin E(G)$.
Istotnie, rozważmy sąsiedztwa z pominięciem
wzajemnej krawędzi $v_jv_k$, jeżeli ona występuje.
Gdyby wszyscy pozostali sąsiedzi $v_j$ byli też
sąsiadami $v_k$, to $v_k$ miałby co najmniej
tylu takich sąsiadów co $v_j$ oraz dodatkowo
$v_1$ (który nie sąsiaduje z $v_j$).
Krawędź $v_jv_k$, jeśli jest obecna, zwiększa oba
stopnie po jednej jednostce. Dostajemy więc
$d_k\ge d_j+1$, sprzecznie z założeniem
$d_j\ge d_k$. Zatem szukany $z$ istnieje.

Usuwamy krawędzie $v_1v_k$ i $v_jz$ oraz
dodajemy $v_1v_j$ i $v_kz$. Nowych krawędzi
wcześniej nie było, więc nadal mamy graf prosty.
Każdy z czterech wierzchołków traci jedną i zyskuje
jedną krawędź, zatem wszystkie stopnie pozostają
bez zmian. Liczba sąsiadów $v_1$ w zbiorze
$\{v_2,\ldots,v_{d+1}\}$ wzrosła o jeden,
sprzecznie z wyborem realizacji. Stąd istnieje
realizacja, w której $v_1$ jest połączony
dokładnie z tymi $d$ wierzchołkami.
Usunięcie $v_1$ daje graf o stopniach wskazanych
w twierdzeniu. To kończy dowód.
\end{proof}

\begin{przyklad}[Obliczenie metodą Havla--Hakimiego]
Rozważmy ciąg $(3,3,2,2,2)$. Po usunięciu
pierwszej trójki zmniejszamy trzy kolejne wyrazy o jeden:
$(2,1,1,2)$. Sortując, otrzymujemy $(2,2,1,1)$.
Powtarzamy: usunięcie pierwszej dwójki daje
$(1,0,1)$, a po uporządkowaniu $(1,1,0)$.
Kolejny krok prowadzi do $(0,0)$, a zatem
pierwotny ciąg jest grafowy. Metoda stanowi
nie tylko test: odwrócenie kolejnych kroków
konstruuje graf o podanych stopniach.
\end{przyklad}

\begin{uwaga}[Zakres pierwszego rozdziału]
Dotychczas nie używaliśmy topologii algebraicznej ani
narzędzi różniczkowych. Za to każdy z poznanych obiektów
ma interpretację topologiczną: cykl staje się okręgiem,
składowe spójne mają sens dla realizacji $|G|$, a
właściwości drzew okażą się związane z trywialnością
pętli. W następnym rozdziale przygotujemy te związki
na poziomie czysto kombinatorycznym.
\end{uwaga}

\chapter{Drzewa, mosty i cykle fundamentalne}
\label{ch:tg-drzewa}

Drzewa są najprostszymi grafami spójnymi: można przemieszczać się
między dowolnymi dwoma wierzchołkami, ale nie ma możliwości
obchodzenia cykli. Dzięki tej sztywnej strukturze udowodnimy
wiele twierdzeń w kilku równoważnych postaciach. Później
zobaczymy, że wybór drzewa rozpinającego pozwala wyznaczać
bazę pierwszej homologii dowolnego grafu.

\section{Drzewa i ich równoważne charakteryzacje}
\label{sec:tg-def-drzewo}

\begin{definicja}[Drzewo i las]
\emph{Drzewem} nazywamy skończony, spójny, acykliczny
graf prosty. \emph{Lasem} nazywamy graf prosty
acykliczny, niekoniecznie spójny. W szczególności
każda składowa spójna lasu jest drzewem.
\end{definicja}

Dopuszczamy drzewo jednowierzchołkowe $K_1$; nie ma
ono krawędzi. Trzeba pamiętać, że acykliczność
nie oznacza braku zamkniętych spacerów (można
wracać tą samą krawędzią).

\begin{twierdzenie}[Jedyność drogi w drzewie]
Niech $G$ będzie skończonym, niepustym grafem prostym.
Wtedy $G$ jest drzewem wtedy i tylko wtedy, gdy
dla dowolnych różnych $u,v\in V(G)$ istnieje
dokładnie jedna droga prosta z $u$ do $v$.
\end{twierdzenie}

\begin{proof}
Załóżmy, że $G$ jest drzewem. Jego spójność zapewnia
istnienie spaceru między $u$ a $v$, a lemat o usuwaniu
powtórzeń --- istnienie drogi prostej. Gdyby istniały
dwie różne drogi proste $P$ i $Q$ z $u$ do $v$,
wybierzmy pierwszy wierzchołek $x$, w którym
rozchodzą się one przy przechodzeniu od $u$.
Idąc dalej po $P$, wybierzmy pierwszy wierzchołek
$y\ne x$ należący również do $Q$ po $x$;
taki wierzchołek istnieje, bo $v$ należy do obu
dróg. Fragmenty $P$ i $Q$ między $x$ i $y$
są wewnętrznie rozłączne i mają różne pierwsze
krawędzie. Ich suma tworzy cykl, co przeczy
acykliczności.

W drugą stronę, jeżeli pomiędzy każdymi dwoma
różnymi wierzchołkami jest dokładnie jedna droga
prosta, graf jest spójny. Gdyby zawierał cykl,
pomiędzy dowolnymi dwoma różnymi wierzchołkami
tego cyklu byłyby dwie drogi proste: dwa łuki
cyklu prowadzące w przeciwnych kierunkach.
Otrzymalibyśmy sprzeczność. Zatem graf jest
spójny i acykliczny, czyli jest drzewem.
\end{proof}

\begin{center}
\begin{tikzpicture}[scale=.92]
  \node[tgTealVertex] (a) at (-2.2,0) {$a$};
  \node[tgTealVertex] (b) at (-.7,1.15) {$b$};
  \node[tgTealVertex] (c) at (-.7,-1.15) {$c$};
  \node[tgTealVertex] (d) at (.8,1.5) {$d$};
  \node[tgTealVertex] (e) at (.8,.3) {$e$};
  \node[tgTealVertex] (f) at (.8,-1.35) {$f$};
  \draw[tgTree] (a)--(b)--(d) (b)--(e) (a)--(c)--(f);
  \draw[tgAccent] (d)--(b)--(a)--(c)--(f);
  \node[font=\small,text=tgNavy] at (2.3,.1) {droga $d\to f$};
\end{tikzpicture}

{\footnotesize W drzewie istnieje tylko jedna droga między
wierzchołkami $d$ i $f$ (zaznaczona na czerwono).}
\end{center}

\begin{lemat}[Liście w drzewie]
Niech $T$ będzie skończonym drzewem mającym
co najmniej dwa wierzchołki. Wtedy $T$ ma
co najmniej dwa różne wierzchołki stopnia jeden.
\end{lemat}

\begin{proof}
Wybierzmy drogę prostą $P=v_0v_1\ldots v_k$
o największej możliwej długości. Ponieważ drzewo
ma co najmniej dwa wierzchołki i jest spójne,
$k\ge1$. Pokażemy, że $d_T(v_0)=1$.
Gdyby $v_0$ miał sąsiada $w\ne v_1$, możliwe
byłyby dwa przypadki. Jeżeli $w$ leżałby na
$P$, krawędź $v_0w$ wraz z odpowiednim fragmentem
$P$ tworzyłaby cykl. Jeżeli $w$ nie leżałby
na $P$, droga $wv_0v_1\ldots v_k$ byłaby
dłuższa niż $P$. Oba przypadki są niemożliwe.
Zatem $v_0$ ma tylko jednego sąsiada $v_1$.
Tak samo $v_k$ ma stopień jeden, a ponieważ
$k\ge1$, oba końce są różne.
\end{proof}

\begin{twierdzenie}[Liczba krawędzi drzewa]
Dla każdego skończonego drzewa $T$ o $n\ge1$
wierzchołkach zachodzi $|E(T)|=n-1$.
\end{twierdzenie}

\begin{proof}
Dowodzimy przez indukcję po $n$. Dla $n=1$
drzewo nie ma krawędzi, więc wzór jest prawdziwy.
Niech $n\ge2$ i załóżmy tezę dla drzew o $n-1$
wierzchołkach. Z poprzedniego lematu $T$ ma
liść $v$ i dokładnie jedną krawędź $e$ incydentną
z $v$. Usuńmy $v$ oraz $e$. Pozostały graf $T-v$
nie ma cykli, bo jest podgrafem drzewa. Jest
też spójny: jedyna droga między dwoma dowolnymi
wierzchołkami różnymi od $v$ nie mogła przechodzić
przez liść $v$ jako wierzchołek wewnętrzny, bo
wymagałoby to dwóch różnych krawędzi przy $v$.
Stąd $T-v$ jest drzewem o $n-1$ wierzchołkach
i z założenia indukcyjnego ma $n-2$ krawędzie.
Po dołączeniu $e$ otrzymujemy $n-1$ krawędzi w $T$.
\end{proof}

\begin{wniosek}[Krawędzie lasu]
Niech $F$ będzie skończonym lasem o $n$ wierzchołkach
i $c$ składowych spójnych. Wtedy $|E(F)|=n-c$.
\end{wniosek}

\begin{proof}
Niech składowe mają odpowiednio $n_1,\ldots,n_c$
wierzchołków. Każda jest drzewem, zatem ma
$n_i-1$ krawędzi. Sumując, dostajemy
$|E(F)|=\sum_{i=1}^{c}(n_i-1)
=\sum_{i=1}^{c}n_i-c=n-c$.
\end{proof}

\begin{twierdzenie}[Charakteryzacja przez liczbę krawędzi]
Niech $G=(V,E)$ będzie skończonym grafem prostym
o $n\ge1$ wierzchołkach i $m$ krawędziach.
Następujące warunki są równoważne:
\begin{enumerate}[label=(\roman*)]
\item $G$ jest drzewem;
\item $G$ jest spójny oraz $m=n-1$;
\item $G$ jest acykliczny oraz $m=n-1$;
\item $G$ jest spójny i usunięcie dowolnej
krawędzi niszczy spójność;
\item $G$ jest acykliczny i dołączenie dowolnej
krawędzi łączącej dwa dotąd niesąsiadujące
wierzchołki tworzy cykl (jeśli takie pary istnieją).
\end{enumerate}
\end{twierdzenie}

\begin{proof}
Wykazaliśmy już (i)$\Rightarrow$(ii) oraz
(i)$\Rightarrow$(iii). Aby udowodnić
(ii)$\Rightarrow$(i), wybierzmy spośród wszystkich
spójnych podgrafów $G$ o pełnym zbiorze
wierzchołków taki, który ma najmniej krawędzi.
Nie zawiera cyklu: gdyby zawierał, usunięcie
jednej krawędzi cyklu nie zniszczyłoby spójności,
ponieważ można ją zastąpić pozostałym fragmentem
cyklu. Otrzymujemy drzewo rozpinające $T$.
Ma ono $n-1$ krawędzi; ponieważ sam $G$ także
ma $n-1$ krawędzi, musi zachodzić $E(G)=E(T)$.
Zatem $G=T$ jest drzewem.

Aby wykazać (iii)$\Rightarrow$(i), oznaczmy
przez $c$ liczbę składowych spójnych $G$.
Acykliczny $G$ jest lasem, zatem z poprzedniego
wniosku $m=n-c$. Z założenia $m=n-1$ wynika
$c=1$, a więc $G$ jest spójny i dlatego jest
drzewem.

Jeśli (i) zachodzi, to każde dwa wierzchołki
drzewa łączy jedyna droga prosta. Usunięcie
dowolnej krawędzi $uv$ niszczy tę drogę, czyli
wierzchołki $u$ i $v$ przestają być połączone;
mamy (iv). Odwrotnie, gdyby graf spełniający
(iv) zawierał cykl, usunięcie dowolnej krawędzi
tego cyklu zachowałoby spójność, co przeczy (iv).
Stąd (iv)$\Rightarrow$(i).

Na koniec załóżmy (i). Dla dwóch niesąsiadujących
wierzchołków $u,v$ istnieje w $G$ jedyna droga
prosta od $u$ do $v$. Dołączenie krawędzi $uv$
zamyka tę drogę w cykl, co daje (v).
Jeżeli natomiast (v) zachodzi, lecz $G$ nie
jest spójny, weźmy wierzchołki $u,v$ z różnych
składowych. Nie sąsiadują ze sobą, a dodanie
krawędzi $uv$ nie może utworzyć cyklu, bo
przed jej dodaniem nie było drogi między nimi.
Sprzeczność. Zatem $G$ jest spójny, a jako
acykliczny jest drzewem.
\end{proof}

\begin{uwaga}
W warunku (v) warto osobno sprawdzić przypadek grafu
pełnego. Jeżeli graf jest pełny i acykliczny,
może mieć co najwyżej dwa wierzchołki; wówczas
jest drzewem. Warunek o ,,każdej brakującej krawędzi''
pozostaje logicznie prawdziwy również wtedy, gdy
takich krawędzi nie ma.
\end{uwaga}

\section{Drzewa rozpinające}
\label{sec:tg-rozpinajace}

\begin{definicja}[Drzewo rozpinające]
Niech $G=(V,E)$ będzie grafem spójnym.
\emph{Drzewem rozpinającym} $G$ nazywamy
drugi graf $T=(V,E_T)$, gdzie $E_T\subseteq E$
oraz $T$ jest drzewem. Innymi słowy:
$T$ ma dokładnie te same wierzchołki co $G$
i pewien podzbiór jego krawędzi.
\end{definicja}

\begin{twierdzenie}[Istnienie drzewa rozpinającego]
Każdy skończony spójny graf prosty ma drzewo
rozpinające. Graf o $n$ wierzchołkach jest
spójny wtedy i tylko wtedy, gdy zawiera
dr łewo rozpinające. W szczególności każdy
spójny graf ma co najmniej $n-1$ krawędzi.
\end{twierdzenie}

\begin{proof}
Wśród wszystkich podgrafów $G$ o zbiorze
wierzchołków $V(G)$ wybieramy spójny podgraf
$T$ z minimalną liczbą krawędzi. Taki wybór
jest możliwy, bo sam $G$ jest jednym z
rozważanych podgrafów, a zbiór krawędzi jest
skończony. Gdyby $T$ zawierał cykl, można
by usunąć z niego jedną krawędź cyklu
bez utraty spójności, sprzecznie z minimalnością.
Zatem $T$ jest drzewem rozpinającym.

Jeżeli $G$ zawiera drzewo rozpinające, to
każde dwa wierzchołki łączy droga w $T$, a więc
również w $G$; jest on spójny. Ostatnia
nierówność wynika z $|E(T)|=n-1$ i z
$E(T)\subseteq E(G)$.
\end{proof}

\begin{przyklad}[Szkielet grafu z cyklami]
Dla grafu o wierzchołkach $a,b,c,d,e$ i krawędziach
$ab,bc,cd,de,ea,ac,ce$ jedno z drzew rozpinających
ma krawędzie $ab,bc,cd,de$; jest to droga
$a-b-c-d-e$. Pozostałe trzy krawędzie, po
ponownym dołączeniu, tworzą niezależne
fundamentalne cykle. Pojęcie ,,niezależności''
sformalizujemy w sekcji o przestrzeni cykli.
\end{przyklad}

\begin{center}
\begin{tikzpicture}[scale=.88]
 \node[tgBlueVertex] (a) at (-2.1,0) {$a$};
 \node[tgBlueVertex] (b) at (-.8,1.4) {$b$};
 \node[tgBlueVertex] (c) at (.9,1.4) {$c$};
 \node[tgBlueVertex] (d) at (2.1,0) {$d$};
 \node[tgBlueVertex] (e) at (0,-1.2) {$e$};
 \draw[tgGhost] (e)--(a) (a)--(c) (c)--(e);
 \draw[tgTree] (a)--(b)--(c)--(d)--(e);
 \node[font=\footnotesize,text=tgTeal] at (0,-1.78)
  {zielone: drzewo rozpinające; szare: krawędzie pozostałe};
\end{tikzpicture}
\end{center}

\begin{uwaga}[Konstrukcja algorytmiczna]
Drzewo rozpinające można otrzymać bez przeglądania
wszystkich podgrafów. Wybieramy wierzchołek
początkowy, następnie stopniowo dołączamy
krawędzie do jeszcze nieodwiedzonych wierzchołków.
Nigdy nie dołączamy krawędzi między dwiema
już odwiedzonymi pozycjami: taka krawędź
zamknęłaby cykl. Spójność gwarantuje, że
proces kończy się dopiero po odwiedzeniu
wszystkich wierzchołków. Wybór kolejności
może odpowiadać przeszukiwaniu wszerz (BFS)
lub w głąb (DFS).
\end{uwaga}

\section{Drzewo ukorzenione i odległości}
\label{sec:tg-ukorzenione}

\begin{definicja}[Korzeń, rodzic, głębokość]
\emph{Drzewem ukorzenionym} nazywamy parę
$(T,r)$, gdzie $T$ jest drzewem, a
$r\in V(T)$ jest wyróżnionym
\emph{korzeniem}. Dla $v\ne r$ jedyna droga
od $r$ do $v$ ma przedostatni wierzchołek
$p(v)$, zwany \emph{rodzicem} $v$. Każdy
wierzchołek $w$ o $p(w)=v$ jest \emph{dzieckiem}
$v$. \emph{Głębokością} $v$ nazywamy
$h(v)=d_T(r,v)$.
\end{definicja}

\begin{stwierdzenie}[Poziomy drzewa]
Niech $(T,r)$ będzie skończonym drzewem
ukorzenionym. Wtedy dla każdej krawędzi
$uv\in E(T)$ zachodzi $|h(u)-h(v)|=1$.
Każdy wierzchołek różny od $r$ ma dokładnie
jednego rodzica, a wzdłuż kolejnych rodziców
głębokość maleje o jeden aż do $r$.
\end{stwierdzenie}

\begin{proof}
Jeżeli $uv$ jest krawędzią, to
$|h(u)-h(v)|\le1$ z nierówności trójkąta
dla odległości grafowej. Gdyby głębokości
były równe, drogi od $r$ do $u$ i do $v$
musiałyby najpierw mieć wspólny odcinek,
a potem rozdzielić się (bo $u\ne v$).
Fragmenty obu dróg od ostatniego wspólnego
wierzchołka do $u,v$ oraz krawędź $uv$
utworzyłyby cykl, co jest niemożliwe.
Stąd różnica głębokości wynosi jeden.

Dla $v\ne r$ jedyna droga od $r$ do $v$
ma dokładnie jeden wierzchołek bezpośrednio
poprzedzający $v$ i jego głębokość wynosi
$h(v)-1$. Gdyby istniał inny sąsiad o
takiej głębokości, powstałyby dwie różne
drogi od $r$ do $v$, wbrew jedyności.
Kolejne zastosowanie rodzica zmniejsza
nieujemną liczbę całkowitą $h$ o jeden,
a więc kończy się dokładnie w korzeniu.
\end{proof}

\begin{center}
\begin{tikzpicture}[scale=.91]
  \node[tgRedVertex] (r) at (0,2.4) {$r$};
  \node[tgBlueVertex] (a) at (-2,1.15) {$a$};
  \node[tgBlueVertex] (b) at (1.75,1.15) {$b$};
  \node[tgTealVertex] (c) at (-3,0) {$c$};
  \node[tgTealVertex] (d) at (-1,0) {$d$};
  \node[tgTealVertex] (e) at (1,0) {$e$};
  \node[tgTealVertex] (f) at (2.55,0) {$f$};
  \draw[tgTree] (r)--(a) (r)--(b) (a)--(c) (a)--(d)
    (b)--(e) (b)--(f);
  \node[text=tgNavy,font=\footnotesize,anchor=west] at (3.25,2.4) {$h=0$};
  \node[text=tgNavy,font=\footnotesize,anchor=west] at (3.25,1.15) {$h=1$};
  \node[text=tgNavy,font=\footnotesize,anchor=west] at (3.25,0) {$h=2$};
\end{tikzpicture}

{\footnotesize Korzeń $r$ i trzy poziomy drzewa.
Każdy wierzchołek poza $r$ ma jednoznacznego rodzica.}
\end{center}

\begin{twierdzenie}[Konstrukcja BFS i odległości]
Niech $G=(V,E)$ będzie skończonym spójnym
grafem prostym, a $r\in V$. Rozważmy algorytm,
który nadaje $r$ poziom $0$, a pozostałym
wierzchołkom kolejne poziomy przez przeglądanie
sąsiadów w kolejności niemalejących poziomów.
Gdy wierzchołek $v$ zostaje pierwszy raz
odkryty jako sąsiad już odkrytego $u$,
nadajemy mu poziom $h(v)=h(u)+1$ i zapisujemy
krawędź $uv$ jako krawędź rodzicielską.
Wtedy $h(v)=d_G(r,v)$ dla każdego $v\in V$,
a krawędzie rodzicielskie tworzą drzewo
rozpinające $G$.
\end{twierdzenie}

\begin{proof}
Najpierw pokażmy, że każdemu wierzchołkowi
przypisany zostanie poziom. Gdyby po zakończeniu
procedury pozostał nieodkryty wierzchołek,
spójność grafu gwarantowałaby krawędź między
zbiorem odkrytych i nieodkrytych wierzchołków.
Algorytm musiałby ją rozważyć podczas
przeglądania wszystkich sąsiadów odkrytego
wierzchołka, więc jej drugi koniec również
zostałby odkryty. Sprzeczność.

Poziomy są rozpatrywane w kolejności $0,1,2,\ldots$.
W chwili odkrycia $v$ z poziomem $h(v)=k$
istnieje łańcuch kolejnych krawędzi
rodzicielskich długości $k$ prowadzący od $r$
do $v$. Wobec tego $d_G(r,v)\le k$.
Gdyby istniała krótsza droga długości $\ell<k$,
to indukcyjnie po kolejnych wierzchołkach tej
drogi algorytm odkryłby wierzchołek na pozycji
$j$ nie później niż przy przetwarzaniu poziomu
$j-1$. W szczególności $v$ otrzymałby poziom
co najwyżej $\ell$, sprzecznie z przypisaniem
poziomu $k$. Zatem $h(v)=d_G(r,v)$.

Każdy $v\ne r$ ma dokładnie jedną wybraną
krawędź rodzicielską. Wzdłuż tych krawędzi
poziom ściśle maleje w kierunku korzenia,
więc nie ma cykli. Możemy też dojść z
każdego wierzchołka do $r$, a więc podgraf
rodzicielski jest spójny i acykliczny.
Jest zatem drzewem rozpinającym.
\end{proof}

\begin{przyklad}[Rozbieżność między drzewem BFS a grafem]
W kwadracie $a-b-c-d-a$ wybierzmy korzeń $a$.
Mamy $d_G(a,b)=d_G(a,d)=1$ oraz $d_G(a,c)=2$.
Jedno z drzew BFS składa się z krawędzi
$ab,ad,bc$, drugie może zawierać $ab,ad,dc$.
Oba odtwarzają odległości od $a$, choć
nie są tym samym podgrafem. Nieużyta krawędź
zamyka cykl i jest celowo pomijana przez algorytm.
\end{przyklad}

\section{Mosty i przekroje krawędziowe}
\label{sec:tg-mosty}

\begin{definicja}[Most]
Niech $G$ będzie skończonym grafem prostym.
Krawędź $e\in E(G)$ nazywamy \emph{mostem},
jeżeli graf $G-e$ ma więcej składowych spójnych
niż $G$.
\end{definicja}

Usunięcie jednej krawędzi może zwiększyć liczbę składowych
co najwyżej o jeden, ponieważ jedynie składowa zawierająca
tę krawędź może ulec rozdzieleniu. W szczególności
w drzewie każda krawędź jest mostem.

\begin{twierdzenie}[Kryterium mostu]
Niech $G$ będzie skończonym grafem prostym,
a $e\in E(G)$. Krawędź $e$ jest mostem wtedy
i tylko wtedy, gdy nie należy do żadnego cyklu grafu.
\end{twierdzenie}

\begin{proof}
Niech $e=uv$ leży na cyklu $C$. Po jej
usunięciu pozostałe krawędzie cyklu wciąż
tworzą drogę z $u$ do $v$. Jeżeli spacer
w $G$ pomiędzy dowolnymi wierzchołkami
przechodził przez $e$, możemy każde takie
przejście zastąpić pozostałym fragmentem
$C$. Otrzymujemy spacer o tych samych
końcach w $G-e$. Zatem usunięcie $e$
nie zmienia składowych spójnych, więc $e$
nie jest mostem.

Załóżmy odwrotnie, że $e=uv$ nie jest
mostem. Wtedy $u$ i $v$ są połączone
w $G-e$, ponieważ w przeciwnym razie
usunięcie $e$ rozdzieliłoby ich wspólną
składową. Z lematu o redukcji spaceru
istnieje w $G-e$ droga prosta $P$ z $u$
do $v$. Ponieważ $e\notin P$, suma $P$
i krawędzi $e$ jest cyklem zawierającym $e$.
\end{proof}

\begin{definicja}[Przekrój krawędziowy]
Dla grafu $G=(V,E)$ i zbioru
$S\subseteq V$ definiujemy \emph{przekrój}
$\delta_G(S)=\{uv\in E:u\in S,\ v\in V\setminus S\}$.
Przekrój jest \emph{nietrywialny}, jeśli
$\varnothing\ne S\ne V$.
\end{definicja}

\begin{stwierdzenie}[Most jako przekrój jednoelementowy]
Niech $G$ będzie spójnym grafem prostym,
a $e\in E(G)$. Wówczas $e$ jest mostem
dokładnie wtedy, gdy istnieje niepusty
właściwy podzbiór $S\subset V(G)$ taki, że
$\delta_G(S)=\{e\}$.
\end{stwierdzenie}

\begin{proof}
Załóżmy, że $e=uv$ jest mostem. W grafie
$G-e$ wierzchołki $u,v$ należą do różnych
składowych. Niech $S$ będzie zbiorem
wierzchołków składowej zawierającej $u$.
Jest on niepusty i właściwy, a żadna krawędź
$G-e$ nie prowadzi od $S$ do jego dopełnienia.
Ponieważ $e$ łączy $u\in S$ z $v\notin S$,
mamy $\delta_G(S)=\{e\}$.

W drugą stronę, jeśli $\delta_G(S)=\{e\}$,
to po usunięciu $e$ nie istnieje żadna
krawędź łącząca $S$ z $V\setminus S$.
Oba zbiory są niepuste, więc w $G-e$
brak spaceru z dowolnego punktu jednego
zbioru do dowolnego punktu drugiego.
Graf nie jest spójny, a zatem $e$ jest mostem.
\end{proof}

\begin{center}
\begin{tikzpicture}[scale=.95]
  \node[tgBlueVertex] (a) at (-2.7,1.0) {$a$};
  \node[tgBlueVertex] (b) at (-1.1,1.0) {$b$};
  \node[tgBlueVertex] (c) at (-1.9,-.45) {$c$};
  \node[tgRedVertex] (d) at (.7,.35) {$d$};
  \node[tgRedVertex] (e) at (2.45,1.0) {$e$};
  \node[tgRedVertex] (f) at (2.45,-.55) {$f$};
  \draw[tgEdge] (a)--(b)--(c)--(a) (d)--(e)--(f)--(d);
  \draw[tgAccent] (b)--(d);
  \node[font=\small,text=tgCoral] at (-.1,1.12) {most $e$};
  \draw[tgPurple,densely dashed,line width=1pt]
     (-.18,-1.05)--(-.18,1.75);
  \node[font=\footnotesize,text=tgPurple] at (-.3,-1.35)
     {$\delta(S)=\{e\}$};
\end{tikzpicture}

{\footnotesize Jedyna krawędź między dwoma fragmentami
grafu jest mostem; krawędzie należące do trójkątów nim nie są.}
\end{center}

\begin{przyklad}[Mosty na schemacie układu połączeń]
Załóżmy, że dwa cykle $C_3$ połączono dokładnie jedną
krawędzią $e$. Krawędź $e$ jest mostem, ponieważ
nie leży na żadnym cyklu. Wszystkie sześć krawędzi
obu trójkątów leży na cyklach i nie jest mostami.
Po usunięciu $e$ graf rozpada się na dwie składowe.
\end{przyklad}

\section{Krawędzie spoza drzewa i cykle fundamentalne}
\label{sec:tg-fundamentalne}

\begin{definicja}[Cykl fundamentalny]
Niech $G$ będzie skończonym spójnym grafem
prostym i niech $T$ będzie jego drzewem
rozpinającym. Dla każdej krawędzi
$e=uv\in E(G)\setminus E(T)$ jedyna droga
$P_T(u,v)$ w $T$ wraz z krawędzią $e$
tworzy cykl $C_e$, zwany \emph{cyklem
fundamentalnym krawędzi $e$ względem $T$}.
\end{definicja}

Jeżeli graf ma $n$ wierzchołków i $m$ krawędzi,
każde drzewo rozpinające ma $n-1$ krawędzi,
a więc pozostaje dokładnie $m-n+1$ krawędzi
nie należących do drzewa. Tyle jest cykli
fundamentalnych, choć wszystkich cykli grafu
może być znacznie więcej.

\begin{center}
\begin{tikzpicture}[scale=1]
  \node[tgBlueVertex] (a) at (-2,0) {$a$};
  \node[tgBlueVertex] (b) at (-.65,1.1) {$b$};
  \node[tgBlueVertex] (c) at (.85,1.1) {$c$};
  \node[tgBlueVertex] (d) at (2.0,0) {$d$};
  \node[tgBlueVertex] (e) at (0,-1.2) {$e$};
  \draw[tgTree] (a)--(b)--(c)--(d)--(e);
  \draw[tgGhost] (a)--(e) (a)--(c);
  \draw[tgAccent] (c)--(e);
  \node[text=tgCoral,font=\footnotesize] at (1.25,-.35) {$e=ce$};
  \node[text=tgNavy,font=\footnotesize] at (0,-1.8)
    {dodanie $ce$ tworzy cykl $c-d-e-c$};
\end{tikzpicture}
\end{center}

\begin{definicja}[Przestrzeń krawędzi nad $\mathbb F_2$]
Dla skończonego grafu $G=(V,E)$ oznaczmy
przez $\mathbb F_2^E$ przestrzeń wszystkich
funkcji $x:E\to\mathbb F_2=\{0,1\}$, ze
zwykłym dodawaniem modulo $2$.
Zbiór krawędzi $F\subseteq E$ utożsamiamy
z funkcją charakterystyczną $\mathbf1_F$.
Dla dwóch zbiorów krawędzi suma odpowiada
\emph{różnicy symetrycznej}:
$F\triangle H=(F\setminus H)\cup(H\setminus F)$.
\end{definicja}

\begin{definicja}[Przestrzeń cykli]
\emph{Przestrzenią cykli} $\mathcal C(G)$ grafu
$G$ nazywamy podprzestrzeń $\mathbb F_2^E$
rozpiętą przez wektory $\mathbf1_C$,
gdzie $C$ przebiega wszystkie cykle proste grafu.
Jeśli graf nie ma cykli, przyjmujemy
$\mathcal C(G)=\{0\}$.
\end{definicja}

\begin{lemat}[Rozkład zbioru krawędzi o parzystych stopniach]
Niech $G=(V,E)$ będzie grafem prostym,
a $F\subseteq E$ takim zbiorem krawędzi,
że w podgrafie $(V,F)$ każdy wierzchołek ma
parzysty stopień. Wtedy $F$ jest rozłączną
sumą zbiorów krawędzi cykli prostych
(w szczególności $\mathbf1_F\in\mathcal C(G)$).
\end{lemat}

\begin{proof}
Dowodzimy przez indukcję po $|F|$. Gdy $F=\varnothing$,
teza jest oczywista. W przeciwnym razie
weźmy składową grafu $(V,F)$ zawierającą
co najmniej jedną krawędź. Każdy jej
nieizolowany wierzchołek ma stopień co najmniej
dwa, bo stopnie są parzyste. W skończonym
grafie o tej własności istnieje cykl:
rozpocznijmy spacer, w każdym kroku przechodząc
krawędzią inną niż ta, którą właśnie weszliśmy.
Nie możemy zakończyć ruchu w wierzchołku
stopnia co najmniej dwa; ponieważ graf jest
skończony, pewien wierzchołek zostanie
powtórzony. Fragment od pierwszego do
następnego wystąpienia powtarzającego się
wierzchołka, po ewentualnym odrzuceniu
wcześniejszego fragmentu, daje cykl prosty $C$.

Usuńmy wszystkie krawędzie $C$ z $F$.
Każdy wierzchołek na cyklu traci dokładnie
dwie krawędzie, a inne nie tracą żadnej.
Wobec tego w pozostałym podgrafie nadal
wszystkie stopnie są parzyste. Zbiór
krawędzi zmniejszył się, więc z założenia
indukcyjnego jest rozłączną sumą cykli.
Dołączając $C$, otrzymujemy rozkład całego $F$.
\end{proof}

\begin{twierdzenie}[Cykle fundamentalne tworzą bazę]
Niech $G=(V,E)$ będzie skończonym spójnym
grafem prostym, $|V|=n$, $|E|=m$,
a $T$ jego drzewem rozpinającym. Rodzina
$\{\mathbf1_{C_e}:e\in E\setminus E(T)\}$
jest bazą przestrzeni cykli $\mathcal C(G)$
nad $\mathbb F_2$. W szczególności
\[
\dim_{\mathbb F_2}\mathcal C(G)=m-n+1.
\]
\end{twierdzenie}

\begin{proof}
Najpierw pokażmy liniową niezależność.
Każdy cykl fundamentalny $C_e$ zawiera
jedną krawędź spoza drzewa, mianowicie
samą $e$; wszystkie inne jego krawędzie
należą do $T$. Dlatego w sumie
$\sum_{e\notin E(T)}a_e\mathbf1_{C_e}$
współrzędna przy krawędzi $f\notin E(T)$
jest dokładnie równa $a_f$. Jeśli cała
suma jest zerowa, każdy $a_f$ musi być
zerem. Rodzina jest więc liniowo niezależna.

Aby wykazać rozpinanie, weźmy dowolny
cykl prosty $C$ i rozważmy zbiór
$F=E(C)\triangle\bigtriangleup_{e\in E(C)\setminus E(T)} E(C_e)$.
Każda krawędź spoza drzewa należąca do $C$
pojawia się dokładnie dwa razy i znika w
różnicy symetrycznej; krawędzie spoza
drzewa nie należące do $C$ nie pojawiają się
wcale. Zatem $F\subseteq E(T)$.
Ponieważ cykl ma przy każdym wierzchołku
stopień parzysty, a własność parzystych
stopni zachowuje się przy różnicy symetrycznej,
$F$ ma wszędzie parzyste stopnie.
Ale niepusty las nie może mieć wszystkich
stopni parzystych: dowolna jego niepusta
składowa z krawędzią ma liść stopnia jeden.
Stąd $F=\varnothing$ i
$\mathbf1_C=\sum_{e\in E(C)\setminus E(T)}\mathbf1_{C_e}$.
W ten sposób każdy cykl leży w rozpiętości
cykli fundamentalnych, więc jest to baza.

Liczba jej elementów wynosi
$|E\setminus E(T)|=m-(n-1)=m-n+1$.
\end{proof}

\begin{przyklad}[Dwa cykle i ich różnica symetryczna]
Rozważmy kwadrat $a-b-c-d-a$ z przekątną $ac$.
Jako drzewo rozpinające weźmy krawędzie
$ab,bc,cd$. Krawędź $da$ tworzy fundamentalny
cykl $a-b-c-d-a$, a krawędź $ac$ tworzy
cykl $a-b-c-a$. Ich różnica symetryczna to
trójkąt $a-c-d-a$, gdyż krawędzie $ab$ i $bc$
występują w obu i znikają modulo $2$.
Mamy $n=4$, $m=5$, więc $\dim\mathcal C(G)=2$.
Wszystkie trzy cykle są zależne: suma ich
wektorów charakterystycznych wynosi zero.
\end{przyklad}

\section{Przekroje fundamentalne i parzystość przecięć}
\label{sec:tg-przekroje-fundamentalne}

\begin{definicja}[Przekrój fundamentalny]
Niech $G$ będzie skończonym spójnym grafem
prostym, a $T$ jego drzewem rozpinającym.
Dla krawędzi $f\in E(T)$ graf $T-f$
ma dokładnie dwie składowe o zbiorach
wierzchołków $S_f$ i $V\setminus S_f$.
Zbiór $D_f=\delta_G(S_f)$ nazywamy
\emph{przekrojem fundamentalnym $f$
względem $T$}.
\end{definicja}

\begin{stwierdzenie}[Parzysta liczba przecięć cyklu i przekroju]
Niech $G=(V,E)$ będzie skończonym grafem prostym,
$C$ jego cyklem prostym, a $S\subseteq V$.
Wtedy
\[
|E(C)\cap\delta_G(S)|\equiv0\pmod2.
\]
\end{stwierdzenie}

\begin{proof}
Przemieszczając się wzdłuż cyklu $C$, za
każdym razem, gdy przechodzimy krawędzią
z przekroju $\delta_G(S)$, zmieniamy
stronę podziału $S\sqcup(V\setminus S)$.
Przejście krawędzią spoza przekroju pozostawia
nas po tej samej stronie. Ponieważ spacer
jest zamknięty, po pełnym obiegu wracamy
do tej samej strony, od której zaczęliśmy.
Liczba zmian strony musi być więc parzysta.
Jest nią dokładnie liczba krawędzi
$E(C)\cap\delta_G(S)$.
\end{proof}

\begin{stwierdzenie}[Dualność fundamentalna]
Niech $G$ będzie spójnym grafem prostym,
$T$ jego drzewem rozpinającym, $e\notin E(T)$,
a $f\in E(T)$. Wtedy
\[
e\in D_f\quad\Longleftrightarrow\quad f\in C_e.
\]
Ponadto $D_f\cap E(T)=\{f\}$.
\end{stwierdzenie}

\begin{proof}
Usunięcie $f$ rozcina drzewo $T$ na
dwie składowe o zbiorach wierzchołków
$S_f$ i $V\setminus S_f$. Jedyna droga
drzewowa pomiędzy końcami $e$ przechodzi
przez $f$ dokładnie wtedy, gdy końce
$e$ leżą po przeciwnych stronach podziału.
Pierwszy warunek oznacza $f\in C_e$,
a drugi dokładnie $e\in D_f$.
W drzewie $T$ jedyną krawędzią
prowadzącą z $S_f$ do dopełnienia jest
$f$, bo gdyby była druga, po jej
dodaniu do $T-f$ powstałby cykl.
Wobec tego $D_f\cap E(T)=\{f\}$.
\end{proof}

\begin{uwaga}[Dlaczego wynik jest przygotowaniem do homologii]
W dalszej części wprowadzimy operator brzegowy
$\partial_1:\mathbb F_2^E\to\mathbb F_2^V$.
Warunek ,,każdy wierzchołek ma stopień
parzysty'' będzie równoważny $\partial_1x=0$.
Dzisiejsze twierdzenie o bazie cykli fundamentalnych
przyjmie więc postać obliczenia
$\dim\ker\partial_1=m-n+1$ dla grafu spójnego.
Po przejściu do realizacji $|G|$ otrzymamy
$H_1(|G|;\mathbb F_2)\cong\mathbb F_2^{m-n+1}$.
W szczególności $|T|$ dla drzewa $T$
nie ma nietrywialnej pierwszej homologii.
\end{uwaga}

\section{Liczba cyklomatyczna i geometryczna intuicja}
\label{sec:tg-cyklomatyczna}

\begin{definicja}[Liczba cyklomatyczna]
Niech $G$ będzie skończonym grafem prostym
o $n$ wierzchołkach, $m$ krawędziach
i $c$ składowych spójnych. Liczbę
$\beta(G)=m-n+c$ nazywamy
\emph{liczbą cyklomatyczną} grafu.
\end{definicja}

\begin{twierdzenie}
Dla każdego skończonego grafu prostego $G$
zachodzi
$\dim_{\mathbb F_2}\mathcal C(G)=m-n+c$.
W szczególności $\beta(G)\ge0$ i
$\beta(G)=0$ wtedy i tylko wtedy, gdy
$G$ jest lasem.
\end{twierdzenie}

\begin{proof}
Niech składowe spójne $G_1,\ldots,G_c$
mają odpowiednio $n_i$ wierzchołków
i $m_i$ krawędzi. Cykl nie może przebiegać
przez więcej niż jedną składową, więc
$\mathcal C(G)$ jest sumą prostą
przestrzeni $\mathcal C(G_i)$. Z twierdzenia
o bazie cykli fundamentalnych otrzymujemy
$\dim\mathcal C(G_i)=m_i-n_i+1$.
Sumując po $i$, dostajemy
$\dim\mathcal C(G)=\sum_i(m_i-n_i+1)=m-n+c$.
Wymiar przestrzeni jest nieujemny.
Jest równy zero dokładnie wtedy, gdy
żadna składowa nie zawiera cyklu, czyli
gdy cały graf jest lasem.
\end{proof}

\begin{przyklad}[Co mierzy liczba $\beta$?]
Dla drzewa na dziesięciu wierzchołkach mamy
$m=9$, $c=1$, stąd $\beta=0$.
Jeśli dołączymy do niego trzy nowe krawędzie
między dotąd niesąsiadującymi wierzchołkami,
to $n$ i $c$ pozostają bez zmian, a $m$
wzrasta o trzy. Dostajemy $\beta=3$:
nie oznacza to, że graf ma dokładnie trzy
cykle, tylko że przestrzeń cykli ma
wymiar trzy. Rzeczywistych cykli może
być znacznie więcej, ponieważ sumy
fundamentalnych cykli mogą tworzyć
nowe cykle.
\end{przyklad}

\begin{stwierdzenie}[Zmiana liczby cyklomatycznej]
Niech $G$ będzie skończonym grafem prostym,
a $e\in E(G)$. Wtedy:
\begin{enumerate}[label=(\roman*)]
\item jeżeli $e$ jest mostem, to $\beta(G-e)=\beta(G)$;
\item jeżeli $e$ nie jest mostem, to
$\beta(G-e)=\beta(G)-1$.
\end{enumerate}
\end{stwierdzenie}

\begin{proof}
Usunięcie krawędzi zmniejsza $m$ o jeden
i nie zmienia $n$. Jeśli $e$ jest mostem,
liczba składowych zwiększa się dokładnie
o jeden: $\beta(G-e)=(m-1)-n+(c+1)=m-n+c$.
Jeśli $e$ nie jest mostem, liczba
składowych pozostaje równa $c$, więc
$\beta(G-e)=(m-1)-n+c=\beta(G)-1$.
\end{proof}

\section*{Podsumowanie części podstawowej}
\addcontentsline{toc}{section}{Podsumowanie części podstawowej}

W dwóch pierwszych rozdziałach wprowadziliśmy grafy bez
odwoływania się do topologii algebraicznej. Znamy już ich
podstawowe modele, stopnie wierzchołków i macierze, wiemy,
jak budować drogi oraz rozpoznawać spójność i dwudzielność.
Rozumiemy również, dlaczego drzewa są minimalnymi grafami
spójnymi, a ich krawędzie stanowią szkielet dowolnego
spójnego grafu. Najważniejszy wynik przygotowujący dalszą
część ma postać $\dim\mathcal C(G)=m-n+c$: za chwilę
zinterpretujemy go jako obliczenie pierwszej grupy homologii.

Dalsze rozdziały rozbudują klasyczną teorię spójności,
cykle Eulera i skojarzenia. Następnie powrócimy do
macierzy incydencji, łańcuchów i operatora Laplace'a.

% ============================================================
% ITERACJA II: klasyczna teoria grafów — rozdziały 3--5
% Wszystkie rysunki w TikZ; brak osobnych zadań.
% ============================================================

\chapter{Przechodzenie po wszystkich krawędziach: teoria Eulera}
\label{ch:tg-euler}

Pierwsze dwa rozdziały dotyczyły przede wszystkim spójności i istnienia
ścieżek. Teraz stawiamy inne pytanie: czy można przejść przez \emph{każdą
krawędź dokładnie raz}? Problem ma historyczne źródło w zagadnieniu
mostów królewieckich. Prowadzi do kryterium całkowicie wyrażonego
przez stopnie wierzchołków, a także do konstruktywnego algorytmu.

\section{Szlaki i cykle Eulera}
\label{sec:tg-euler-def}

W tej sekcji dopuszczamy multigrafy nieskierowane, także z pętlami.
W zapisie spaceru wskazujemy kolejno zarówno wierzchołki, jak i
krawędzie: $v_0,e_1,v_1,\ldots,e_k,v_k$, gdzie końcami $e_i$ są
$v_{i-1}$ i $v_i$. Taki zapis odróżnia krawędzie równoległe.
\emph{Nośnikiem krawędziowym} grafu o niepustym zbiorze
krawędzi nazywamy podgraf o zbiorze wierzchołków złożonym
z końców wszystkich krawędzi i o pełnym zbiorze krawędzi.
Wierzchołki izolowane nie wpływają na możliwość przejścia
przez wszystkie krawędzie. Jeżeli $E=\varnothing$, nośnik
jest pustą przestrzenią i stosujemy osobną konwencję
ze wskazanej poniżej definicji.

\begin{definicja}[Szlak Eulera i cykl Eulera]
Niech $G=(V,E,\partial)$ będzie skończonym multigrafem nieskierowanym.
\emph{Szlakiem Eulera} nazywamy szlak, w którym każda krawędź
$e\in E$ występuje dokładnie raz. Jeżeli taki szlak jest zamknięty,
nazywamy go \emph{cyklem Eulera} (mimo że może wielokrotnie
odwiedzać ten sam wierzchołek). Dla $E=\varnothing$ i $V\ne\varnothing$ przyjmujemy,
że spacer długości zero, w dowolnym wierzchołku,
jest zamkniętym szlakiem Eulera. W tym szczególnym przypadku
nie należy wymagać spójności nośnika krawędziowego.
\end{definicja}

Termin ,,cykl Eulera'' jest tradycyjny: nie jest to cykl prosty
zdefiniowany wcześniej. Przykładowo obieg dwóch trójkątów
połączonych jednym wspólnym wierzchołkiem jest cyklem Eulera,
choć wspólny wierzchołek odwiedzamy więcej niż raz.

\begin{center}
\begin{tikzpicture}[scale=.95,line cap=round,>=Stealth]
  \node[tgBlueVertex] (a) at (-2,0) {$a$};
  \node[tgBlueVertex] (b) at (-1.1,1.2) {$b$};
  \node[tgBlueVertex] (c) at (0,0) {$c$};
  \node[tgBlueVertex] (d) at (1.1,1.2) {$d$};
  \node[tgBlueVertex] (e) at (2,0) {$e$};
  \draw[tgTree,-{Stealth[length=2mm]}] (a)--(b);
  \draw[tgTree,-{Stealth[length=2mm]}] (b)--(c);
  \draw[tgTree,-{Stealth[length=2mm]}] (c)--(a);
  \draw[tgAccent,-{Stealth[length=2mm]}] (c)--(d);
  \draw[tgAccent,-{Stealth[length=2mm]}] (d)--(e);
  \draw[tgAccent,-{Stealth[length=2mm]}] (e)--(c);
  \node[text=tgNavy,font=\footnotesize] at (0,-.65)
     {$a\to b\to c\to d\to e\to c\to a$};
\end{tikzpicture}

{\footnotesize Dwa trójkąty mają wspólny wierzchołek $c$.
Kolory pokazują dwa fragmenty jednego zamkniętego szlaku Eulera.}
\end{center}

\begin{lemat}[Parzystość stopni wewnętrznych]
\label{lem:tg-euler-parity}
Niech $v_0,e_1,v_1,\dots,e_k,v_k$ będzie szlakiem
w skończonym multigrafie nieskierowanym $G$. Jeżeli jest on
zamknięty, to w podgrafie złożonym z jego krawędzi każdy
wierzchołek ma parzysty stopień. Jeżeli $v_0\ne v_k$,
to stopnie $v_0,v_k$ w tym podgrafie są nieparzyste,
a wszystkie pozostałe są parzyste.
\end{lemat}

\begin{proof}
Każde pojawienie się wierzchołka jako punktu pośredniego
szlaku ma dwie incydentne krawędzie: tę, którą weszliśmy,
i tę, którą wyszliśmy. Krawędzie te są różne, ponieważ szlak
nie powtarza krawędzi. W ten sposób można pogrupować
wszystkie końce wykorzystanych krawędzi w pary.
Jeżeli początek i koniec są tym samym wierzchołkiem,
pierwszy wyjazd parujemy z ostatnim powrotem; również pętla
liczona jest jako dwa końce tej samej krawędzi. Przy różnych
końcach zostaje dokładnie jeden niesparowany koniec
przy $v_0$ i jeden przy $v_k$. To daje odpowiednie parzystości.
\end{proof}

\begin{twierdzenie}[Euler: charakterystyka zamknięta]
\label{thm:tg-euler-circuit}
Niech $G$ będzie skończonym multigrafem nieskierowanym
z co najmniej jedną krawędzią. Wówczas istnieje cykl Eulera
w $G$ wtedy i tylko wtedy, gdy
\begin{enumerate}[label=(\roman*)]
\item wszystkie wierzchołki dodatniego stopnia leżą w jednej
składowej spójnej;
\item dla każdego $v\in V(G)$ stopień $d_G(v)$ jest parzysty.
\end{enumerate}
\end{twierdzenie}

\begin{proof}
Jeżeli istnieje cykl Eulera, jego kolejne krawędzie łączą
wszystkie wierzchołki o dodatnim stopniu, a lemat
\ref{lem:tg-euler-parity} zapewnia parzystość ich stopni.

Załóżmy odwrotnie, że (i)--(ii) zachodzą. Wybieramy wierzchołek
$v_0$ dodatniego stopnia i rozpoczynamy szlak, wybierając
za każdym razem dowolną \emph{niewykorzystaną} krawędź.
Ponieważ graf jest skończony, proces musi się zatrzymać.
Nie może zakończyć się w wierzchołku $w\ne v_0$:
przed ewentualnym zakończeniem w $w$ liczba wykorzystanych
końców krawędzi incydentnych z $w$ jest nieparzysta,
a całkowity stopień $w$ jest parzysty. Pozostaje zatem
niewykorzystana krawędź wychodząca z $w$.
Otrzymaliśmy pewien zamknięty szlak $C$.

Jeżeli $C$ zawiera wszystkie krawędzie, dowód jest zakończony.
W przeciwnym razie istnieje niewykorzystana krawędź
przylegająca do wierzchołka $C$. Istotnie, gdyby takiej
krawędzi nie było, zbiór wierzchołków $C$ byłby odseparowany
od wszystkich pozostałych krawędzi, co przeczyłoby spójności
nośnika krawędziowego. Niech $w$ będzie takim wierzchołkiem.
W multigrafie złożonym z \emph{niewykorzystanych} krawędzi
wszystkie stopnie nadal są parzyste: od pierwotnego
parzystego stopnia odjęliśmy parzysty stopień pochodzący
z zamkniętego szlaku $C$. Możemy więc, zaczynając w $w$,
zbudować kolejny niepusty szlak zamknięty $C'$ złożony
wyłącznie z niewykorzystanych krawędzi.

W miejscu wystąpienia $w$ na $C$ wstawiamy cały obieg $C'$
i wracamy następnie do dalszej części $C$. Powstaje dłuższy
zamknięty szlak, nadal bez powtarzania krawędzi.
Przy każdym takim kroku zużywamy co najmniej jedną nową
krawędź. Po skończonej liczbie kroków uzyskamy cykl Eulera.
\end{proof}

\begin{uwaga}[Algorytm Hierholzera]
Dowód jest algorytmem: budujemy zamknięty szlak, wyszukujemy
na nim wierzchołek incydentny z niewykorzystaną krawędzią,
a następnie wplatamy nowy zamknięty szlak. Przy przechowywaniu
list krawędzi incydentnych można wykonać całą procedurę w czasie
$O(|V|+|E|)$, ponieważ każdą krawędź rozpatrujemy
stałą liczbę razy. Sama zasada ,,idź dowolną krawędzią''
bez etapu wplatania cykli nie gwarantuje natychmiastowego
wykorzystania wszystkich krawędzi.
\end{uwaga}

\begin{twierdzenie}[Euler: charakterystyka otwarta]
\label{thm:tg-euler-trail}
Niech $G$ będzie skończonym multigrafem nieskierowanym
z co najmniej jedną krawędzią, którego wszystkie wierzchołki
dodatniego stopnia należą do jednej składowej spójnej.
Wówczas $G$ ma szlak Eulera o różnych końcach wtedy
i tylko wtedy, gdy dokładnie dwa wierzchołki mają
nieparzysty stopień. Końcami dowolnego takiego szlaku
są właśnie te dwa wierzchołki.
\end{twierdzenie}

\begin{proof}
Konieczność daje lemat \ref{lem:tg-euler-parity}.
Jeśli dokładnie dwa wierzchołki $s,t$ mają nieparzysty
stopień, dołączamy \emph{nową} krawędź $e_*$ o końcach $s,t$,
nawet gdy były już połączone krawędzią. Otrzymany multigraf
ma wszystkie stopnie parzyste oraz spójny nośnik krawędziowy.
Na mocy twierdzenia \ref{thm:tg-euler-circuit} ma cykl Eulera.
Usuwając z tego cyklu wystąpienie $e_*$ i rozpoczynając
odpowiednio obieg, otrzymujemy szlak przechodzący po każdej
oryginalnej krawędzi dokładnie raz, od $s$ do $t$.
\end{proof}

\begin{przyklad}[Mosty królewieckie]
Model historycznego układu siedmiu mostów ma cztery
wierzchołki odpowiadające obszarom lądu. Ich stopnie
wynoszą $3,3,3,5$; suma $14$ odpowiada siedmiu krawędziom,
zgodnie z lematem o uściskach dłoni. Ponieważ wierzchołków
nieparzystych jest cztery, nie istnieje ani zamknięty,
ani otwarty szlak Eulera. Problem nie zależy od długości
mostów ani od ich rysunkowego położenia.
\end{przyklad}

\begin{center}
\begin{tikzpicture}[scale=.95,line cap=round]
 \node[tgBlueVertex] (a) at (-2,0) {$A$};
 \node[tgBlueVertex] (b) at (2,0) {$B$};
 \node[tgRedVertex] (c) at (0,1.45) {$C$};
 \node[tgRedVertex] (d) at (0,-1.45) {$D$};
 \draw[tgEdge] (a) to[bend left=19] (c);
 \draw[tgEdge] (a) to[bend right=19] (c);
 \draw[tgEdge] (b) to[bend left=19] (c);
 \draw[tgEdge] (b) to[bend right=19] (c);
 \draw[tgAccent] (a)--(d);
 \draw[tgAccent] (b)--(d);
 \draw[tgEdge] (c)--(d);
 \node[font=\footnotesize,text=tgNavy] at (0,-2.10)
    {$d(A)=3,\ d(B)=3,\ d(C)=5,\ d(D)=3$};
\end{tikzpicture}

{\footnotesize Jeden z abstrakcyjnych modeli problemu siedmiu mostów;
krawędzie równoległe odpowiadają różnym mostom.}
\end{center}

\section{Orientacje krawędzi i skierowane obiegi Eulera}
\label{sec:tg-orientacje}

\begin{definicja}[Graf skierowany i orientacja]
\emph{Graf skierowany} to para $(V,A)$, gdzie $V$ jest
skończonym zbiorem wierzchołków, a $A\subseteq V\times V$
jest zbiorem łuków; łuk $(u,v)$ zapisujemy jako $u\to v$.
\emph{Orientacja} grafu prostego nieskierowanego $G=(V,E)$
powstaje przez zastąpienie każdej krawędzi $\{u,v\}\in E$
dokładnie jednym z łuków $(u,v)$ lub $(v,u)$.
Stopień wejściowy $d^-(v)$ jest liczbą łuków kończących się
w $v$, a stopień wyjściowy $d^+(v)$ jest liczbą łuków
zaczynających się w $v$.
\end{definicja}

Dla grafów skierowanych rozważamy spacer zgodny z kierunkami
łuków. Graf skierowany nazywamy \emph{silnie spójnym}, gdy
każde dwa wierzchołki można połączyć skierowanymi ścieżkami
w obie strony. Przez spójność \emph{słabą} rozumiemy spójność
grafu otrzymanego po zapomnieniu kierunków łuków.

\begin{twierdzenie}[Skierowane obiegi Eulera]
Niech $D=(V,A)$ będzie skończonym grafem skierowanym
z co najmniej jednym łukiem. Wówczas ma on skierowany
zamknięty szlak Eulera wtedy i tylko wtedy, gdy
\begin{enumerate}[label=(\roman*)]
\item wszystkie wierzchołki incydentne z łukami należą do
jednej słabo spójnej składowej;
\item $d^+(v)=d^-(v)$ dla każdego $v\in V$.
\end{enumerate}
\end{twierdzenie}

\begin{proof}
W skierowanym zamkniętym szlaku każde wejście do wierzchołka
jest sparowane z wyjściem, więc otrzymujemy (ii).
Ponieważ szlak zawiera wszystkie łuki, zachodzi też (i).

Odwrotnie, wychodzimy z wierzchołka incydentnego z łukiem
i przedłużamy skierowany szlak bez powtarzania łuków.
Jeżeli zatrzymałby się w wierzchołku innym niż początek,
to miałby tam o jedno wykorzystane wejście więcej niż
wykorzystanych wyjść. Z równości $d^+=d^-$ wynikałoby
istnienie niewykorzystanego wyjścia, sprzeczność.
Otrzymujemy więc zamknięty szlak. Po jego usunięciu
równość stopni wejściowych i wyjściowych pozostaje prawdziwa.
Słaba spójność zapewnia, że jeżeli nie wykorzystaliśmy
wszystkich łuków, jeden z nich ma koniec w wierzchołku
już użytym w naszym szlaku. Z tego wierzchołka konstruujemy
następny zamknięty szlak i wplatamy go do poprzedniego.
Iteracja kończy się po wykorzystaniu wszystkich łuków.
\end{proof}

\begin{stwierdzenie}[Orientacja eulerowska]
Niech $G$ będzie skończonym multigrafem nieskierowanym.
Istnieje orientacja każdej jego krawędzi, dla której
$d^+(v)=d^-(v)$ przy wszystkich $v$, wtedy i tylko wtedy,
gdy każdy stopień $d_G(v)$ jest parzysty.
\end{stwierdzenie}

\begin{proof}
W każdej orientacji $d_G(v)=d^+(v)+d^-(v)$,
więc przy równości obu składników stopień jest parzysty.
Odwrotnie, w każdej nieizolowanej składowej grafu o
parzystych stopniach wybieramy cykl Eulera i orientujemy
krawędzie zgodnie z przebiegiem. Każde wejście na obiegu
odpowiada jednemu wyjściu. Wierzchołki izolowane mają
oba stopnie równe zero.
\end{proof}

\section{Silnie spójna orientacja i mosty}
\label{sec:tg-robbins}

\begin{twierdzenie}[Robbins]
\label{thm:tg-robbins}
Niech $G$ będzie skończonym spójnym grafem prostym
z co najmniej dwoma wierzchołkami. Istnieje silnie spójna
orientacja krawędzi $G$ wtedy i tylko wtedy,
gdy $G$ nie ma mostów.
\end{twierdzenie}

\begin{proof}
Jeżeli $e$ jest mostem, usunięcie go rozdziela graf na dwa
niepuste zbiory wierzchołków $S$ oraz $V\setminus S$.
Jedyny łuk łączący te części, niezależnie od wybranej
orientacji, umożliwia przejście tylko w jednym kierunku.
Silna spójność jest zatem niemożliwa.

Załóżmy, że $G$ nie ma mostów. Wykonujemy przeszukiwanie
w głąb (DFS) zaczynając od wierzchołka $r$ i otrzymujemy
drzewo rozpinające $T$. Orientujemy każdą krawędź drzewa
\emph{od korzenia ku potomkom}. Każdą pozostałą krawędź
orientujemy od potomka do przodka. Takie ustalenie
jest możliwe, ponieważ w nieskierowanym DFS każda
krawędź niedrzewowa łączy wierzchołek z jego przodkiem
(w razie potrzeby dowód: podczas pierwszego przeglądania
krawędzi jeden koniec jest jeszcze na stosie rekurencji;
gdyby oba końce należały do różnych, już zakończonych gałęzi,
krawędź zostałaby wcześniej użyta do odkrycia jednej z nich).

Z korzenia $r$ można dotrzeć do każdego wierzchołka
po skierowanych krawędziach drzewa. Pokażemy, że
z każdego wierzchołka można również dotrzeć do $r$.
Dla wierzchołka $v\ne r$ niech $p$ będzie jego rodzicem,
a $S_v$ zbiorem wszystkich potomków $v$ łącznie z $v$.
Krawędź $pv$ nie jest mostem, zatem istnieje inna krawędź
między $S_v$ i $V\setminus S_v$. Własność DFS zapewnia,
że jest ona niedrzewowa i łączy pewnego potomka $x\in S_v$
z przodkiem $y$ wierzchołka $v$ (ewentualnie $y=p$).
W orientacji możemy przejść od $v$ do $x$ w dół drzewa,
a potem łukiem $x\to y$. W ten sposób dotarliśmy do
wierzchołka o ściśle mniejszej głębokości niż $v$.
Powtarzając argument, osiągamy $r$, gdyż głębokość
jest nieujemną liczbą całkowitą i za każdym razem maleje.
Stąd dla dowolnych $u,w$ istnieje skierowany spacer
$u\to r\to w$, który można skrócić do skierowanej
ścieżki. Orientacja jest silnie spójna.
\end{proof}

\begin{przyklad}[Dlaczego obecność mostu przeszkadza]
Dwa trójkąty połączone pojedynczą krawędzią nie mają
silnie spójnej orientacji: każda podróż między trójkątami
musi użyć mostu, którego orientacja dopuszcza tylko jeden
kierunek. Z drugiej strony dwa trójkąty o wspólnym
wierzchołku nie mają mostów i można je zorientować
zgodnie z obiegiem każdego trójkąta, otrzymując
silnie spójny graf skierowany.
\end{przyklad}

\begin{center}
\begin{tikzpicture}[scale=.86,line cap=round]
\foreach \x/\y/\n in {-3/0/a,-2/1/b,-1/0/c,1/0/d,2/1/e,3/0/f}
  \node[tgBlueVertex] (\n) at (\x,\y) {$\n$};
\draw[tgTree] (a)--(b)--(c)--(a) (d)--(e)--(f)--(d);
\draw[tgAccent,line width=2.8pt] (c)--(d);
\node[font=\footnotesize,text=tgCoral] at (0,.43) {most};
\node[font=\footnotesize,text=tgNavy] at (0,-.73)
 {Jedna krawędź między blokami uniemożliwia silną spójność orientacji.};
\end{tikzpicture}
\end{center}

\chapter{Spójność, przekroje i twierdzenie Mengera}
\label{ch:tg-menger}

Spójność mówi jedynie, czy istnieje choć jedna ścieżka między
wierzchołkami. W zagadnieniach geometrycznych i kombinatorycznych
potrzebujemy często informacji silniejszej: \emph{ile wzajemnie
niezależnych ścieżek} łączy dwa punkty? Twierdzenie Mengera
wiąże odpowiedź z najmniejszą liczbą krawędzi albo wierzchołków,
które trzeba usunąć, aby przerwać połączenie.

\section{Przekroje krawędziowe i separatory wierzchołkowe}
\label{sec:tg-cuts}

\begin{definicja}[Przekrój krawędziowy]
Niech $G=(V,E)$ będzie grafem prostym, a
$\varnothing\ne S\subsetneq V$. \emph{Przekrojem krawędziowym}
wyznaczonym przez $S$ nazywamy zbiór
$\delta_G(S):=\{uv\in E:u\in S,\ v\in V\setminus S\}$.
Dla różnych wierzchołków $s,t$ zbiór krawędzi $F\subseteq E$
nazywamy \emph{separatorem krawędziowym $s,t$}, jeśli
w grafie $G-F$ nie ma ścieżki z $s$ do $t$.
\end{definicja}

\begin{definicja}[Spójność krawędziowa i wierzchołkowa]
Niech $G=(V,E)$ będzie spójnym grafem prostym z $|V|\ge2$.
\emph{Spójnością krawędziową} $\lambda(G)$ nazywamy
najmniejszą liczbę krawędzi, których usunięcie rozspójnia $G$.
\emph{Spójnością wierzchołkową} $\kappa(G)$ nazywamy
najmniejszą liczbę wierzchołków, których usunięcie
rozspójnia $G$ lub pozostawia graf z jednym wierzchołkiem.
Dla grafu niespójnego przyjmujemy $\lambda(G)=\kappa(G)=0$,
a dla $K_1$ obie liczby definiujemy jako zero.
\end{definicja}

Warunek dotyczący grafu jednoelementowego pozwala
otrzymać $\kappa(K_n)=n-1$. Inaczej dla grafu pełnego
nie byłoby żadnego zbioru rozspajającego w zwykłym sensie.

\begin{stwierdzenie}
\label{prop:tg-lambda-cut}
Niech $G=(V,E)$ będzie spójnym grafem prostym z $|V|\ge2$.
Wówczas
\[
\lambda(G)=\min_{\varnothing\ne S\subsetneq V}|\delta_G(S)|
\qquad\text{oraz}\qquad \lambda(G)\le \delta(G).
\]
\end{stwierdzenie}

\begin{proof}
Usunięcie wszystkich krawędzi $\delta_G(S)$ oddziela
$S$ od $V\setminus S$, więc
$\lambda(G)\le|\delta_G(S)|$.
Odwrotnie, jeśli $F$ rozspaja graf, wybieramy zbiór
wierzchołków jednej składowej $S$ grafu $G-F$.
Żadna krawędź $G$ między $S$ i jego dopełnieniem
nie mogła pozostać, więc $\delta_G(S)\subseteq F$.
Otrzymujemy nierówność przeciwną po minimalizacji.
Dla $S=\{v\}$ mamy $|\delta_G(S)|=d_G(v)$;
minimalizując stopień dostajemy $\lambda(G)\le\delta(G)$.
\end{proof}

\begin{definicja}[Separator wierzchołkowy]
Niech $s,t$ będą różnymi \emph{niesąsiadującymi} wierzchołkami
grafu prostego $G=(V,E)$. Zbiór $U\subseteq V\setminus\{s,t\}$
nazywamy \emph{separatorem wierzchołkowym $s,t$}, jeśli
w $G-U$ nie istnieje ścieżka z $s$ do $t$.
Liczbę $\kappa_G(s,t)$ definiujemy jako najmniejszą
moc takiego separatora.
\end{definicja}

Warunek niesąsiadowania jest istotny: gdy $st\in E$,
nie istnieje żaden separator wierzchołkowy nieobejmujący
$s$ ani $t$, ponieważ bezpośrednia krawędź nadal łączy końce.
Dlatego wersję wierzchołkową twierdzenia Mengera
najpierw sformułujemy dla wierzchołków niesąsiadujących.

\begin{center}
\begin{tikzpicture}[scale=.92,line cap=round]
 \node[tgBlueVertex] (s) at (-3,0) {$s$};
 \node[tgTealVertex] (u) at (-1,.95) {$u$};
 \node[tgTealVertex] (v) at (-1,-.95) {$v$};
 \node[tgRedVertex] (x) at (1,.95) {$x$};
 \node[tgRedVertex] (y) at (1,-.95) {$y$};
 \node[tgBlueVertex] (t) at (3,0) {$t$};
 \draw[tgEdge] (s)--(u) (s)--(v) (u)--(x) (u)--(y)
       (v)--(x) (v)--(y) (x)--(t) (y)--(t);
 \draw[tgAccent,densely dashed,line width=1.6pt]
       (0,-1.48)--(0,1.48);
 \node[font=\footnotesize,text=tgCoral] at (0,1.76)
       {przekrój o czterech krawędziach};
\end{tikzpicture}

{\footnotesize Przekrój krawędziowy zależy od podziału wierzchołków,
a nie od sposobu przecięcia linii na rysunku.}
\end{center}

\section{Sieci przepływowe: pomocnicze narzędzie dowodowe}
\label{sec:tg-flows}

\begin{definicja}[Sieć i przepływ]
\emph{Siecią przepływową} nazywamy skończony graf skierowany
$D=(V,A)$ bez pętli, z dwoma różnymi wyróżnionymi wierzchołkami
$s,t$ i funkcją \emph{przepustowości}
$c:A\to\mathbb R_{\ge0}$.
\emph{Przepływem} jest funkcja $f:A\to\mathbb R_{\ge0}$
spełniająca $0\le f(a)\le c(a)$ dla wszystkich łuków oraz
\emph{równowagę} $\sum_{a=(u,v)}f(a)=\sum_{a=(v,w)}f(a)$
dla każdego $v\notin\{s,t\}$. Jego \emph{wartość} to
$|f|:=\sum_{a=(s,w)}f(a)-\sum_{a=(u,s)}f(a)$.
Dla $S\subseteq V$ z $s\in S$ i $t\notin S$ definiujemy
pojemność przekroju
$c(S,V\setminus S):=\sum_{(u,v)\in A,\,u\in S,\,v\notin S}c(u,v)$.
\end{definicja}

W tej i następnych definicjach sumujemy po łukach
wskazanej orientacji. Jeśli w sieci są oba łuki $u\to v$
i $v\to u$, traktujemy je jako różne elementy.

\begin{lemat}[Wartość przepływu przez przekrój]
\label{lem:tg-cutbound}
Niech $f$ będzie przepływem w skończonej sieci z końcami
$s,t$. Dla dowolnego $S\subseteq V$ z $s\in S$ i $t\notin S$
zachodzi
\[
|f|=\sum_{\substack{(u,v)\in A\\u\in S,\,v\notin S}}f(u,v)
 -\sum_{\substack{(u,v)\in A\\u\notin S,\,v\in S}}f(u,v)
 \le c(S,V\setminus S).
\]
\end{lemat}

\begin{proof}
Sumujemy równania zachowania przepływu po wszystkich
wierzchołkach $S\setminus\{s\}$ i dodajemy do nich
definicję wartości przepływu w $s$. Każdy łuk o obu końcach
w $S$ jest liczony raz z plusem i raz z minusem, więc się
redukuje. Pozostają jedynie łuki wychodzące z $S$
i wchodzące do $S$, z przeciwnymi znakami. Daje to równość.
Nierówność wynika z nieujemności przepływu na łukach
wchodzących i ograniczeń $f\le c$ na łukach wychodzących.
\end{proof}

\begin{definicja}[Sieć resztowa]
Dla przepływu $f$ i oryginalnego łuku $a=(u,v)$
wprowadzamy dwie możliwości zmiany: przejście $u\to v$
o \emph{pojemności resztowej} $c(a)-f(a)$ oraz przejście
$v\to u$ o pojemności resztowej $f(a)$. Pomijamy
przejścia o zerowej pojemności. Różne łuki oryginalne
mogą dawać równoległe przejścia resztowe; pozostają one
rozróżnione. \emph{Ścieżką powiększającą} nazywamy
skierowaną ścieżkę z $s$ do $t$ w takiej sieci resztowej.
\end{definicja}

\begin{twierdzenie}[Maksymalny przepływ -- minimalny przekrój]
\label{thm:tg-maxflow}
Niech $D=(V,A)$ będzie skończoną siecią przepływową,
a $c:A\to\mathbb Z_{\ge0}$ funkcją całkowitoliczbowych
przepustowości. Wtedy istnieje całkowitoliczbowy przepływ
maksymalny $f$ i zachodzi
\[
\max_f |f|=\min_{\substack{s\in S\subseteq V\\t\notin S}}
 c(S,V\setminus S).
\]
\end{twierdzenie}

\begin{proof}[Dowód przez ścieżki powiększające]
Rozpoczynamy od przepływu zerowego. Dopóki istnieje
ścieżka powiększająca $P$ z $s$ do $t$, definiujemy
$\varepsilon$ jako najmniejszą dodatnią pojemność resztową
na tej ścieżce. Na oryginalnych łukach odpowiadających
krokom w przód zwiększamy przepływ o $\varepsilon$,
a na łukach odpowiadających krokom wstecz zmniejszamy
go o $\varepsilon$. Każdy przepływ pośredni spełnia
ograniczenia przepustowości. W wierzchołkach wewnętrznych
ścieżki dodane i odjęte wartości kompensują się,
więc równowaga zostaje zachowana. Wartość przepływu
zwiększa się dokładnie o $\varepsilon$.

Przy pojemnościach całkowitych również wszystkie przepływy
pośrednie i pojemności resztowe są całkowite, toteż
$\varepsilon\ge1$. Z lematu \ref{lem:tg-cutbound}
wartość przepływu jest ograniczona przez pojemność
dowolnego przekroju, a zatem algorytm zatrzyma się
po skończonej liczbie iteracji.

Niech $f$ będzie przepływem końcowym i niech $S$ będzie
zbiorem wierzchołków osiągalnych z $s$ po dodatnich
łukach resztowych. Skoro algorytm się zatrzymał, $t\notin S$.
Dla każdego oryginalnego łuku $u\to v$ z $u\in S$,
$v\notin S$, musi zachodzić $f(u,v)=c(u,v)$:
w przeciwnym razie istniałby dodatni resztowy krok
$u\to v$ i $v$ należałby do $S$.
Analogicznie dla oryginalnego łuku $v\to u$ z $v\notin S$,
$u\in S$ mamy $f(v,u)=0$, ponieważ jego dodatni przepływ
tworzyłby resztowy krok $u\to v$.
Z tożsamości przepływu przez przekrój otrzymujemy więc
$|f|=c(S,V\setminus S)$.

Lemat \ref{lem:tg-cutbound} zapewnia, że wartość żadnego
przepływu nie przekracza pojemności żadnego przekroju.
Znaleźliśmy przepływ i przekrój, dla których występuje
równość, zatem są one odpowiednio maksymalne i minimalne.
Każdy przepływ w konstrukcji jest całkowitoliczbowy.
\end{proof}

\begin{przyklad}[Rachunek przepływu i przekroju]
Rozważmy sieć z łukami $s\to a$, $s\to b$, $a\to t$,
$b\to t$ o pojemności $1$ oraz $a\to b$ o pojemności $2$.
Przepływ wysyłający jednostkę po drodze $s\to a\to t$
i jednostkę po drodze $s\to b\to t$ ma wartość $2$.
Przekrój $S=\{s\}$ ma pojemność $2$, więc na mocy
poprzedniego twierdzenia przepływ jest maksymalny.
Dodatkowy łuk $a\to b$ o większej pojemności
nie może zwiększyć przepływu: wąskim gardłem
są dwa łuki wychodzące z $s$.
\end{przyklad}

\begin{center}
\begin{tikzpicture}[scale=.95,>=Stealth,line cap=round]
 \node[tgBlueVertex] (s) at (-2.8,0) {$s$};
 \node[tgTealVertex] (a) at (0,1.15) {$a$};
 \node[tgTealVertex] (b) at (0,-1.15) {$b$};
 \node[tgBlueVertex] (t) at (2.8,0) {$t$};
 \draw[tgTree,-{Stealth[length=2.5mm]}] (s)--node[above,font=\footnotesize] {$1$}(a);
 \draw[tgTree,-{Stealth[length=2.5mm]}] (s)--node[below,font=\footnotesize] {$1$}(b);
 \draw[tgTree,-{Stealth[length=2.5mm]}] (a)--node[above,font=\footnotesize] {$1$}(t);
 \draw[tgTree,-{Stealth[length=2.5mm]}] (b)--node[below,font=\footnotesize] {$1$}(t);
 \draw[tgEdge,-{Stealth[length=2.5mm]}] (a)--node[right,font=\footnotesize] {$2$}(b);
 \draw[tgAccent,densely dashed] (-1.95,-1.60)--(-1.95,1.60);
 \node[text=tgCoral,font=\footnotesize] at (-1.95,1.95) {przekrój $S=\{s\}$};
\end{tikzpicture}

{\footnotesize Liczby przy łukach oznaczają przepustowości,
a zielone strzałki pokazują dwa niezależne jednostkowe przepływy.}
\end{center}

\section{Twierdzenie Mengera w wersji krawędziowej}
\label{sec:tg-menger-edges}

\begin{definicja}[Ścieżki krawędziowo rozłączne]
Niech $s,t$ będą różnymi wierzchołkami grafu prostego $G$.
Dwie ścieżki $s$--$t$ są \emph{krawędziowo rozłączne},
jeżeli nie mają żadnej wspólnej krawędzi; mogą natomiast
mieć wspólne wierzchołki pośrednie. Oznaczmy przez
$\rho_E(s,t)$ maksymalną liczność rodziny parami
krawędziowo rozłącznych ścieżek $s$--$t$.
Niech $\lambda_G(s,t)$ oznacza minimalną moc separatora
krawędziowego $s,t$.
\end{definicja}

\begin{twierdzenie}[Menger, wersja krawędziowa]
\label{thm:tg-menger-edge}
Dla dowolnego skończonego grafu prostego $G=(V,E)$
i różnych wierzchołków $s,t\in V$ zachodzi
$\rho_E(s,t)=\lambda_G(s,t)$.
\end{twierdzenie}

\begin{proof}
Każdy separator krawędziowy przecina każdą ścieżkę
$s$--$t$. Jeżeli ścieżki są parami krawędziowo rozłączne,
to różne ścieżki potrzebują różnych krawędzi separatora.
Stąd $\rho_E(s,t)\le\lambda_G(s,t)$.

Dla nierówności odwrotnej zastępujemy każdą krawędź
$\{u,v\}$ dwoma przeciwnymi łukami $u\to v$, $v\to u$
o pojemności $1$. Każdy przekrój $S$ tej sieci ma
pojemność równą $|\delta_G(S)|$: dla każdej krawędzi
przekraczającej granicę $S$ dokładnie jeden z dwóch
łuków wychodzi z $S$. Zatem najmniejsza pojemność
przekroju jest równa $\lambda_G(s,t)$.

Na mocy twierdzenia \ref{thm:tg-maxflow} istnieje
całkowitoliczbowy przepływ tej wartości. Jeśli na
obu przeciwnych łukach jakiejś krawędzi płynie
jednostka, usuwamy oba te przepływy. Równowaga
we wszystkich wierzchołkach i wartość całego
przepływu się nie zmieniają, a wzdłuż każdej
krawędzi nieskierowanej przepływ jest teraz
jednokierunkowy albo zerowy.

Rozłóżmy pozostały całkowity przepływ na ścieżki
i cykle, uzasadniając dokładnie ten krok. Jeśli w grafie
łuków o dodatnim przepływie istnieje skierowany cykl,
zmniejszamy przepływ na każdym jego łuku o najmniejszą
z występujących tam dodatnich wartości. Zachowuje to
równowagę w każdym wierzchołku i wartość całego przepływu,
a usuwa co najmniej jeden dodatni łuk. Po skończenie wielu
takich operacjach pozostaje przepływ dodatni wyłącznie
na skierowanym grafie acyklicznym.
Jeśli jego wartość jest dodatnia, z $s$ wychodzi łuk
o dodatnim przepływie. Przechodzimy po takich łukach.
Nie możemy się zapętlić z powodu acykliczności ani
zatrzymać w wierzchołku wewnętrznym: dodatni dopływ
i zachowanie przepływu wymuszają dodatni odpływ.
Dochodzimy zatem do $t$, otrzymując skierowaną ścieżkę.
Odejmujemy od niej jeden (całkowitoliczbowy) przepływ
jednostkowy i powtarzamy. Po dokładnie $|f|$ krokach
wartość przepływu jest zerowa, a pozostałe ewentualne
cykle można pominąć. Otrzymaliśmy $|f|$ ścieżek $s$--$t$.
Ponieważ każdy łuk ma przepływ co najwyżej $1$,
a przeciwne łuki tej samej krawędzi zostały skasowane,
otrzymane ścieżki nie mają wspólnych krawędzi.
Ich liczba jest równa wartości przepływu,
czyli $\lambda_G(s,t)$. To daje nierówność odwrotną.
\end{proof}

\begin{wniosek}
Niech $G$ będzie spójnym grafem prostym z co najmniej
dwoma wierzchołkami. Wtedy $\lambda(G)$ jest najmniejszą,
po wszystkich parach różnych wierzchołków $s,t$,
liczbą krawędziowo rozłącznych ścieżek łączących $s$ i $t$.
\end{wniosek}

\begin{proof}
Z twierdzenia \ref{thm:tg-menger-edge} liczba ścieżek
dla pary $s,t$ jest równa najmniejszej liczbie krawędzi,
których usunięcie rozdziela tę parę. Minimalizując po
wszystkich parach, otrzymujemy dokładnie najmniejszą
liczbę krawędzi rozspajających graf.
\end{proof}

\section{Twierdzenie Mengera w wersji wierzchołkowej}
\label{sec:tg-menger-vertices}

\begin{definicja}[Ścieżki wewnętrznie rozłączne]
Niech $s,t$ będą różnymi wierzchołkami grafu $G$.
Ścieżki $P_1,\dots,P_k$ z $s$ do $t$ nazywamy
\emph{wewnętrznie wierzchołkowo rozłącznymi}, jeżeli
żadne dwie z nich nie mają wspólnego wierzchołka
poza ewentualnie wspólnymi końcami $s,t$.
Niech $\rho_V(s,t)$ będzie największą możliwą
liczbą takich ścieżek.
\end{definicja}

\begin{twierdzenie}[Menger, wersja wierzchołkowa]
\label{thm:tg-menger-vertex}
Niech $G=(V,E)$ będzie skończonym grafem prostym,
a $s,t\in V$ różnymi \emph{niesąsiadującymi} wierzchołkami.
Wtedy $\rho_V(s,t)=\kappa_G(s,t)$.
\end{twierdzenie}

\begin{proof}
Każdy separator wierzchołkowy $s,t$ musi zawierać
przynajmniej jeden wierzchołek wewnętrzny każdej
ścieżki $s$--$t$. Jeśli takie ścieżki są wewnętrznie
rozłączne, separator musi mieć co najmniej tyle
wierzchołków, ile jest ścieżek. Stąd
$\rho_V(s,t)\le\kappa_G(s,t)$.

Zbudujmy skierowaną sieć, w której \emph{rozdzielimy
wierzchołki na wejścia i wyjścia}. Dla każdego
$v\in V\setminus\{s,t\}$ tworzymy wierzchołki
$v^-$ i $v^+$ połączone łukiem $v^-\to v^+$
o pojemności $1$. Końce $s,t$ pozostawiamy
nierozdzielone. Każdą krawędź $uv\in E$
zastępujemy dwoma przeciwnymi łukami,
łącząc \emph{wyjście} jednego końca z \emph{wejściem}
drugiego: odpowiednio $u^+\to v^-$ i $v^+\to u^-$,
gdzie dla $s,t$ symbol $z^\pm$ oznacza po prostu $z$.
Każdemu z tych łuków nadajemy pojemność
$M:=|V|+1$.

Ponieważ $s$ i $t$ nie sąsiadują, usunięcie wszystkich
pozostałych wierzchołków oddziela $s$ od $t$.
Istnieje zatem przekrój sieci o pojemności co najwyżej
$|V|-2<M$, utworzony przez łuki $v^-\to v^+$.
Każdy przekrój minimalny ma pojemność mniejszą niż $M$,
więc nie przecina żadnego z łuków o pojemności $M$.
W takim przekroju przecinamy wyłącznie łuki
$v^-\to v^+$ o pojemności $1$.

Odpowiadające im wierzchołki tworzą separator
$s,t$ w grafie wyjściowym: gdyby po ich usunięciu
istniała ścieżka $s$--$t$, odpowiadająca jej droga
w sieci od $s$ do $t$ nie przecinałaby minimalnego
przekroju. Odwrotnie, dla dowolnego separatora
$U\subseteq V\setminus\{s,t\}$ możemy zbudować
przekrój sieci o pojemności co najwyżej $|U|$:
usuwamy łuki $v^-\to v^+$ dla $v\in U$ i bierzemy
zbiór wierzchołków osiągalnych z $s$ w pozostałej
sieci; przez jego granicę mogą przechodzić
jedynie usunięte łuki o pojemności $1$.
Stąd minimalna pojemność przekroju sieci jest
równa $\kappa_G(s,t)$.

Twierdzenie \ref{thm:tg-maxflow} daje całkowitoliczbowy
przepływ tej wartości. Rozkład przepływu na ścieżki
jednostkowe i cykle przeprowadzamy tak samo jak
w dowodzie wersji krawędziowej. Każda jednostkowa
ścieżka przechodząca przez wierzchołek $v\ne s,t$
musi użyć łuku $v^-\to v^+$, a jego pojemność wynosi $1$.
Zatem różne ścieżki nie mogą mieć wspólnego wierzchołka
wewnętrznego. Po zapomnieniu orientacji łuków
i pominięciu sztucznych wierzchołków $v^\pm$
otrzymujemy $\kappa_G(s,t)$ wewnętrznie rozłącznych
ścieżek $s$--$t$ w $G$. Wraz z pierwszą nierównością
to kończy dowód.
\end{proof}

\begin{center}
\begin{tikzpicture}[scale=.94,>=Stealth,line cap=round]
 \node[tgBlueVertex] (u) at (-2.4,0) {$u$};
 \node[tgTealVertex] (mi) at (-.25,0) {$v^-$};
 \node[tgRedVertex] (po) at (1.2,0) {$v^+$};
 \node[tgBlueVertex] (w) at (3.2,0) {$w$};
 \draw[tgEdge,-{Stealth[length=2.6mm]}] (u)--(mi);
 \draw[tgAccent,-{Stealth[length=2.6mm]}] (mi)--node[above,font=\footnotesize] {pojemność $1$}(po);
 \draw[tgEdge,-{Stealth[length=2.6mm]}] (po)--(w);
 \node[text=tgNavy,font=\footnotesize] at (.42,-.69)
    {przejście przez $v$ zużywa jedną jednostkę pojemności};
\end{tikzpicture}

{\footnotesize Rozcięcie wierzchołka na wejście i wyjście;
kolorowy łuk koduje ograniczenie użycia wierzchołka.}
\end{center}

\begin{przyklad}[Ścieżki rozłączne w prostym grafie]
Niech $G$ ma wierzchołki $s,t,a,b,c$ i krawędzie
$sa,at,sb,bt,sc,ct$. Istnieją trzy wewnętrznie
wierzchołkowo rozłączne ścieżki $s\! -\!a\! -\! t$,
$s\! -\! b\! -\! t$ i $s\! -\! c\! -\! t$.
Usunięcie $a,b,c$ rozdziela $s,t$, więc
$\rho_V(s,t)=\kappa_G(s,t)=3$.
Jeśli dodamy krawędź $ab$, liczba ta nie wzrośnie:
separator $\{a,b,c\}$ nadal ma moc trzy.
Przykład pokazuje, że dodatkowe krawędzie nie zawsze
zwiększają liczbę niezależnych tras.
\end{przyklad}

\begin{center}
\begin{tikzpicture}[scale=.95,line cap=round]
 \node[tgBlueVertex] (s) at (-2.6,0) {$s$};
 \node[tgBlueVertex] (t) at (2.6,0) {$t$};
 \foreach \n/\y/\col in {a/1.35/tgBlue,b/0/tgTeal,c/-1.35/tgCoral}{
   \node[tgVertex,draw=\col,fill=\col!12] (\n) at (0,\y) {$\n$};
   \draw[draw=\col,line width=2pt] (s)--(\n)--(t);
 }
 \node[font=\footnotesize,text=tgNavy] at (0,-1.98)
 {Trzy ścieżki rozłączne wewnętrznie: $s\to a\to t$, $s\to b\to t$, $s\to c\to t$.};
\end{tikzpicture}
\end{center}

\section{Ile spójności zapewnia graf?}
\label{sec:tg-whitney}

\begin{twierdzenie}[Nierówności Whitneya]
\label{thm:tg-whitney}
Niech $G$ będzie skończonym spójnym grafem prostym
o co najmniej dwóch wierzchołkach. Wówczas
\[
\kappa(G)\le\lambda(G)\le\delta(G).
\]
\end{twierdzenie}

\begin{proof}
Najpierw udowodnimy nierówność $\kappa(G)\le\delta(G)$.
Wybierzmy wierzchołek $v$ o stopniu minimalnym
$d(v)=\delta(G)$ i usuńmy jego wszystkich sąsiadów.
W pozostałym grafie $v$ jest izolowany. Jeżeli pozostał
jeszcze inny wierzchołek, graf jest niespójny.
Jeżeli natomiast pozostał tylko $v$, to usunęliśmy
$|V|-1=\delta(G)$ wierzchołków i uzyskaliśmy
graf jednoelementowy. W obu przypadkach z definicji
spójności wierzchołkowej $\kappa(G)\le\delta(G)$.

Niech teraz $k=\lambda(G)$ i wybierzmy zbiór
$\varnothing\ne S\subsetneq V$ taki, że
$|\delta_G(S)|=k$; jego istnienie wynika ze
stwierdzenia \ref{prop:tg-lambda-cut}. Oznaczmy
$T=V\setminus S$ i $r=|S|$.

Jeśli $r>k$, usuńmy z $S$ wszystkie wierzchołki
będące końcami krawędzi przekroju $\delta_G(S)$.
Usunęliśmy co najwyżej $k$ wierzchołków,
w $S$ pozostał co najmniej jeden, a wszystkie
wierzchołki $T$ pozostały. Nie ma między nimi
żadnej krawędzi: każda taka krawędź wchodziłaby
do wybranego przekroju. Pozostały graf jest
niespójny, zatem $\kappa(G)\le k$.

Jeżeli $r\le k$, policzmy sumę stopni wierzchołków $S$.
Każdy z $r$ wierzchołków ma najwyżej $r-1$
sąsiadów w $S$, a łączna liczba krawędzi
z $S$ do $T$ wynosi $k$. Zatem
\[
r\delta(G)\le\sum_{v\in S}d_G(v)\le r(r-1)+k,
\qquad
\delta(G)\le r-1+\frac{k}{r}\le k.
\]
Ostatnia nierówność wynika z $1\le r\le k$:
po pomnożeniu przez $r$ jest równoważna
$(r-1)(r-k)\le0$. Z pierwszej części dowodu
$\kappa(G)\le\delta(G)\le k=\lambda(G)$.
Pozostała nierówność $\lambda(G)\le\delta(G)$
wynika ze stwierdzenia \ref{prop:tg-lambda-cut}.
\end{proof}



\begin{przyklad}
Dla cyklu $C_n$ przy $n\ge3$ mamy
$\kappa(C_n)=\lambda(C_n)=\delta(C_n)=2$.
Dla drzewa z co najmniej dwoma wierzchołkami mamy
$\lambda(G)=1$, a przy $|V|\ge3$ również $\kappa(G)=1$.
Dla grafu pełnego $K_n$ z $n\ge2$ wszystkie trzy
liczby są równe $n-1$.
\end{przyklad}

\chapter{Skojarzenia, systemy reprezentantów i cykle Hamiltona}
\label{ch:tg-skojarzenia}

Drugim podstawowym problemem kombinatorycznym jest wyszukiwanie
wielu krawędzi, które nie mogą ze sobą kolidować. Skojarzenia
opisują przydziały, równoległe parowania i układy niezależnych
wyborów. Na końcu rozdziału wrócimy do przechodzenia po grafach,
ale zamiast wszystkich \emph{krawędzi} spróbujemy odwiedzić
wszystkie \emph{wierzchołki} dokładnie raz. Takie jest
zasadnicze rozróżnienie między teorią Eulera i Hamiltona.

\section{Skojarzenia i ścieżki powiększające}
\label{sec:tg-matchings}

\begin{definicja}[Skojarzenie]
Niech $G=(V,E)$ będzie skończonym grafem prostym.
\emph{Skojarzeniem} nazywamy podzbiór $M\subseteq E$,
w którym żadne dwie krawędzie nie mają wspólnego końca.
Wierzchołek jest \emph{skojarzony}, jeśli jest końcem
pewnej krawędzi z $M$, a w przeciwnym razie jest
\emph{wolny}. Skojarzenie jest \emph{maksymalne}, gdy
nie można dodać do niego kolejnej krawędzi bez
naruszenia warunku; jest \emph{największe}, gdy ma
największą możliwą liczbę krawędzi. Jest \emph{doskonałe},
gdy każdy wierzchołek jest skojarzony.
\end{definicja}

Maksymalność nie oznacza największej liczności.
Przykładowo w ścieżce czterowierzchołkowej $a-b-c-d$
skojarzenie $\{bc\}$ jest maksymalne, lecz nie największe:
$\{ab,cd\}$ ma dwie krawędzie.

\begin{definicja}[Ścieżka naprzemienna i powiększająca]
Niech $M$ będzie skojarzeniem grafu $G$. \emph{Ścieżką
naprzemienną względem $M$} nazywamy ścieżkę, na której
krawędzie należące do $M$ i do $E\setminus M$
występują naprzemiennie. Jest ona \emph{powiększająca},
jeżeli oba jej końce są wolne, a pierwsza i ostatnia
krawędź nie należą do $M$. Dopuszczamy tylko ścieżki
z co najmniej jedną krawędzią.
\end{definicja}

\begin{lemat}[Odwrócenie ścieżki powiększającej]
Niech $P$ będzie ścieżką powiększającą względem
skojarzenia $M$ w grafie $G$. Wtedy
$M':=M\mathbin{\triangle}E(P)$ jest skojarzeniem,
a $|M'|=|M|+1$, gdzie $\triangle$ oznacza różnicę
symetryczną zbiorów.
\end{lemat}

\begin{proof}
Na wewnętrznych wierzchołkach $P$ jedna z dwóch
krawędzi ścieżki należy do $M$, a druga nie.
Po zamianie ról tych krawędzi każdy taki wierzchołek
nadal jest incydentny dokładnie z jedną krawędzią
nowego skojarzenia. Końce $P$ były wolne,
a teraz dostają po jednej krawędzi. Poza $P$
skojarzenie się nie zmienia. Ścieżka rozpoczyna się
i kończy krawędzią spoza $M$, zatem liczba usuniętych
krawędzi $M$ jest o jeden mniejsza od liczby dodanych.
Daje to $|M'|=|M|+1$.
\end{proof}

\begin{center}
\begin{tikzpicture}[scale=.85,line cap=round]
\foreach \i/\n in {0/a,1/b,2/c,3/d,4/e,5/f}
 {\node[tgVertex] (\n) at (1.3*\i,0) {$\n$};}
\draw[tgTree,line width=2.6pt] (b)--(c) (d)--(e);
\draw[tgAccent,line width=2.5pt] (a)--(b) (c)--(d) (e)--(f);
\node[font=\footnotesize,text=tgNavy] at (3.25,-.80)
 {Zielone krawędzie należą do $M$, czerwone — do $E\setminus M$.};
\end{tikzpicture}

{\footnotesize Po zamianie kolorów na tej ścieżce
otrzymujemy skojarzenie o jedną krawędź większe.}
\end{center}

\begin{twierdzenie}[Berge]
\label{thm:tg-berge}
Skojarzenie $M$ skończonego grafu prostego $G$
jest największe wtedy i tylko wtedy, gdy
nie istnieje ścieżka powiększająca względem $M$.
\end{twierdzenie}

\begin{proof}
Jeśli istnieje ścieżka powiększająca, poprzedni lemat
daje skojarzenie większe od $M$.

Załóżmy odwrotnie, że $M$ nie jest największe.
Wybieramy skojarzenie $N$ o $|N|>|M|$.
Rozważmy graf o zbiorze krawędzi $M\triangle N$.
Ponieważ każdemu wierzchołkowi przylega co najwyżej
jedna krawędź z $M$ i co najwyżej jedna z $N$,
w grafie tym stopień każdego wierzchołka jest co
najwyżej dwa. Jego nieizolowane składowe są zatem
ścieżkami albo cyklami. Ponieważ na każdej składowej
kolejne krawędzie muszą na przemian należeć do $M$
i $N$, każdy cykl ma parzystą liczbę krawędzi
z obu skojarzeń, a więc nie zmienia różnicy ich
liczności. Ścieżki mają równe liczby krawędzi
obu typów albo różnią się o jedną.

Ponieważ łącznie $|N|>|M|$, istnieje składowa
będąca ścieżką z jedną krawędzią $N$ więcej
niż $M$. Jej pierwsza i ostatnia krawędź
należą do $N\setminus M$. Jej końce są wolne
w $M$: gdyby któryś był skojarzony przez $M$,
odpowiednia krawędź musiałaby należeć do $M\triangle N$
i przedłużać ścieżkę (albo pokrywać się z krawędzią
$N$, co jest niemożliwe). Otrzymaliśmy ścieżkę
powiększającą względem $M$.
\end{proof}

\section{Grafy dwudzielne i twierdzenie Halla}
\label{sec:tg-hall}

\begin{definicja}[Sąsiedztwo zbioru]
Niech $G=(A\sqcup B,E)$ będzie skończonym grafem
dwudzielnym o ustalonych częściach $A$ i $B$.
Dla $S\subseteq A$ kładziemy
$N(S):=\{b\in B:\exists a\in S,\ ab\in E\}$.
Mówimy, że skojarzenie \emph{pokrywa $A$}, jeżeli
każdy wierzchołek $a\in A$ jest końcem pewnej
krawędzi skojarzenia.
\end{definicja}

\begin{twierdzenie}[Hall: twierdzenie o małżeństwach]
\label{thm:tg-hall}
Niech $G=(A\sqcup B,E)$ będzie skończonym grafem
dwudzielnym. Istnieje skojarzenie pokrywające
wszystkie wierzchołki zbioru $A$ wtedy i tylko wtedy,
gdy dla każdego $S\subseteq A$ zachodzi
\[
|N(S)|\ge |S|.
\]
\end{twierdzenie}

\begin{proof}[Dowód indukcyjny]
Konieczność jest natychmiastowa: różne wierzchołki
$S$ muszą być skojarzone z różnymi wierzchołkami
zbioru $N(S)$.

Dla dostateczności indukujemy względem $|A|$.
Jeśli $|A|=0$, skojarzenie puste działa.
Jeśli $|A|=1$, warunek Halla daje choć jednego
sąsiada, więc możemy wybrać jedną krawędź.
Załóżmy, że $|A|\ge2$ i twierdzenie zostało
udowodnione dla mniejszych części lewych.

\emph{Przypadek pierwszy:} istnieje niepusty, właściwy
zbiór $S\subsetneq A$ z $|N(S)|=|S|$.
W podgrafie indukowanym na $S\sqcup N(S)$
warunek Halla jest spełniony, więc z założenia
indukcyjnego istnieje skojarzenie pokrywające $S$.
Na grafie z częściami $A\setminus S$
i $B\setminus N(S)$ warunek również jest spełniony:
dla $T\subseteq A\setminus S$ mamy
$|N(S\cup T)|\ge|S|+|T|$, a zatem
$|N(T)\setminus N(S)|\ge|T|$.
Stąd drugie skojarzenie pokrywa $A\setminus S$.
Ponieważ zbiory wierzchołków użytych przez obydwa
skojarzenia są rozłączne, ich suma pokrywa $A$.

\emph{Przypadek drugi:} dla każdego niepustego
właściwego $S\subsetneq A$ zachodzi
$|N(S)|\ge|S|+1$. Wybierzmy dowolne $a\in A$.
Warunek Halla dla $\{a\}$ zapewnia istnienie
sąsiada $b\in B$. Usuwamy $a$ i $b$.
Dla dowolnego niepustego $T\subseteq A\setminus\{a\}$
zbiór $T$ jest właściwym podzbiorem $A$, więc
$|N(T)\setminus\{b\}|\ge|N(T)|-1\ge|T|$.
Dla $T=\varnothing$ warunek jest automatyczny.
Z założenia indukcyjnego otrzymujemy skojarzenie
pokrywające $A\setminus\{a\}$ w pozostałym grafie;
po dodaniu krawędzi $ab$ uzyskujemy skojarzenie
pokrywające $A$.
Oba przypadki wyczerpują możliwości.
\end{proof}

\begin{przyklad}[Przydział trzech osób do trzech stanowisk]
Niech $A=\{a_1,a_2,a_3\}$ oznacza osoby,
$B=\{b_1,b_2,b_3\}$ stanowiska, a krawędź
łączy osobę ze stanowiskiem, na którym może pracować.
Załóżmy, że $N(a_1)=\{b_1,b_2\}$,
$N(a_2)=\{b_2\}$, $N(a_3)=\{b_2,b_3\}$.
Warunek Halla jest spełniony: dla pojedynczych osób
każdy zbiór sąsiadów jest niepusty, dla par są
co najmniej dwa dostępne stanowiska, a dla całego $A$
są trzy. Jedynym sposobem przypisania $a_2$ jest
$b_2$, po czym można wybrać $a_1b_1$ i $a_3b_3$.
Dostajemy skojarzenie doskonałe.
\end{przyklad}

\begin{center}
\begin{tikzpicture}[scale=.95,line cap=round]
 \foreach \y/\x in {1.2/1,0/2,-1.2/3}{
  \node[tgBlueVertex] (a\x) at (-2.6,\y) {$a_{\x}$};
  \node[tgRedVertex] (b\x) at (2.6,\y) {$b_{\x}$};
 }
 \draw[tgTree,line width=2.7pt] (a1)--(b1) (a2)--(b2) (a3)--(b3);
 \draw[tgGhost] (a1)--(b2) (a3)--(b2);
 \node[font=\footnotesize,text=tgNavy] at (0,-1.75)
 {Zielone: skojarzenie pokrywające wszystkie trzy osoby.};
\end{tikzpicture}
\end{center}

\begin{stwierdzenie}[Wzór na deficyt Halla]
\label{prop:tg-hall-def}
Niech $G=(A\sqcup B,E)$ będzie skończonym grafem
dwudzielnym. Oznaczmy przez $\nu(G)$ maksymalną
liczność skojarzenia. Wtedy
\[
\nu(G)=|A|-\max_{S\subseteq A}\bigl(|S|-|N(S)|\bigr).
\]
\end{stwierdzenie}

\begin{proof}
Niech $M$ będzie dowolnym skojarzeniem, a $S\subseteq A$.
Co najwyżej $|N(S)|$ wierzchołków $S$ może być
skojarzonych z sąsiadami, więc co najmniej
$|S|-|N(S)|$ wierzchołków $A$ pozostaje wolnych.
Zatem $|M|\le|A|-|S|+|N(S)|$; minimalizując
prawą stronę, uzyskujemy górne oszacowanie.

Niech $d:=\max_{S\subseteq A}(|S|-|N(S)|)$.
Ponieważ $S=\varnothing$ daje zero, $d\ge0$.
Dodajmy do prawej części grafu $d$ nowych
wierzchołków, każdy połączony ze wszystkimi
wierzchołkami $A$. Dla każdego $S\subseteq A$
nowe sąsiedztwo ma moc $|N(S)|+d\ge|S|$.
Z twierdzenia Halla istnieje skojarzenie
pokrywające $A$ w powiększonym grafie.
Co najwyżej $d$ jego krawędzi prowadzi do nowych
wierzchołków. Po usunięciu tych krawędzi
pozostaje skojarzenie w oryginalnym grafie
o co najmniej $|A|-d$ krawędziach.
To daje nierówność odwrotną.
\end{proof}

\section{Pokrycia wierzchołkowe i twierdzenie Kőniga}
\label{sec:tg-konig}

\begin{definicja}[Pokrycie wierzchołkowe]
Dla grafu $G=(V,E)$ zbiór $C\subseteq V$ nazywamy
\emph{pokryciem wierzchołkowym}, jeśli każda krawędź
ma co najmniej jeden koniec w $C$.
Najmniejszą możliwą liczność takiego zbioru
oznaczamy przez $\tau(G)$.
\end{definicja}

Każde pokrycie wierzchołkowe musi zawierać przynajmniej
jeden koniec każdej krawędzi skojarzenia. Ponieważ
krawędzie skojarzenia są rozłączne w końcach,
$\nu(G)\le\tau(G)$ dla dowolnego grafu. Dla grafów
dwudzielnych zachodzi równość.

\begin{twierdzenie}[Kőnig]
\label{thm:tg-konig}
Niech $G=(A\sqcup B,E)$ będzie skończonym grafem
dwudzielnym. Wtedy minimalna moc pokrycia
wierzchołkowego jest równa maksymalnej liczności
skojarzenia: $\tau(G)=\nu(G)$.
\end{twierdzenie}

\begin{proof}
Wybieramy skojarzenie największe $M$ i oznaczamy
przez $A_0\subseteq A$ zbiór wierzchołków wolnych.
Rozważamy wszystkie ścieżki naprzemienne zaczynające
się w $A_0$ i przechodzące z $A$ do $B$ wyłącznie
krawędziami spoza $M$, a z $B$ do $A$ wyłącznie
krawędziami $M$. Niech $Z_A\subseteq A$ i $Z_B\subseteq B$
będą zbiorami wierzchołków osiągalnych w ten sposób.

Żaden wierzchołek $Z_B$ nie jest wolny: w przeciwnym
razie otrzymalibyśmy ścieżkę powiększającą,
sprzecznie z twierdzeniem \ref{thm:tg-berge}.
Zatem każdy $b\in Z_B$ jest skojarzony z pewnym
$a\in Z_A$. Ponadto każdy skojarzony $a\in Z_A$
ma partnera w $Z_B$ (do $a$ mogliśmy dojść
po krawędzi skojarzenia). Zbiory skojarzonych
wierzchołków w $Z_A$ i wszystkich wierzchołków
$Z_B$ są więc w bijekcji za pomocą krawędzi $M$.

Definiujemy $C:=(A\setminus Z_A)\cup Z_B$.
Pokażemy, że jest to pokrycie. Gdyby jakaś
krawędź $ab$ miała oba końce poza $C$,
to $a\in Z_A$ i $b\notin Z_B$.
Jeżeli $ab\notin M$, z definicji osiągalności
musielibyśmy mieć $b\in Z_B$. Jeśli zaś $ab\in M$,
to skoro $a\in Z_A$, również jego partner $b$
należy do $Z_B$. W obu przypadkach sprzeczność.

Każda krawędź $M$ ma dokładnie jeden koniec w $C$:
partnerzy z $Z_A$ i $Z_B$ przechodzą do $C$ tylko
po stronie $B$, a partnerzy poza nimi tylko po
stronie $A$. Żaden wierzchołek wolny nie należy
do $C$ (wolne w $A$ są w $Z_A$, a wolnych
w $Z_B$ nie ma). W konsekwencji $|C|=|M|$.
Ponieważ zawsze $\nu(G)\le\tau(G)$, a znaleźliśmy
pokrycie o mocy $|M|=\nu(G)$, otrzymujemy równość.
\end{proof}

\begin{uwaga}
Twierdzenia Halla, Bergego i Kőniga są odrębnymi opisami
tego samego zjawiska: przeszkoda w znalezieniu dużego
skojarzenia ma postać zbioru wierzchołków ze zbyt małą
liczbą dostępnych sąsiadów. W dalszej części skryptu
podobne idee będą wyrażane przez rząd macierzy,
przekroje i przepływy, a później przez grupy homologii.
\end{uwaga}

\section{Ścieżki i cykle Hamiltona}
\label{sec:tg-hamilton}

\begin{definicja}[Hamiltonowskość]
Niech $G$ będzie skończonym grafem prostym.
\emph{Ścieżką Hamiltona} nazywamy ścieżkę zawierającą
wszystkie wierzchołki $G$. \emph{Cyklem Hamiltona}
nazywamy cykl zawierający wszystkie wierzchołki $G$.
Graf z cyklem Hamiltona nazywamy \emph{hamiltonowskim}.
\end{definicja}

Obecność cyklu Eulera i Hamiltona to dwa różne warunki.
Cykl Eulera przechodzi przez wszystkie krawędzie,
a cykl Hamiltona przez wszystkie wierzchołki.
Cykl $C_n$ spełnia oba warunki. Dwa trójkąty o jednym
wspólnym wierzchołku mają cykl Eulera, ale nie mają
cyklu Hamiltona: usunięcie wspólnego wierzchołka
rozspaja graf, podczas gdy usunięcie dowolnego
wierzchołka z cyklu Hamiltona pozostawia ścieżkę
zawierającą wszystkie pozostałe wierzchołki.

\begin{center}
\begin{tikzpicture}[scale=.95,line cap=round]
 \foreach \x/\y/\n in {-2/0/a,-1/1.25/b,0/0/c,1/1.25/d,2/0/e}
   \node[tgBlueVertex] (\n) at (\x,\y) {$\n$};
 \draw[tgTree] (a)--(b)--(c)--(a);
 \draw[tgAccent] (c)--(d)--(e)--(c);
 \node[font=\footnotesize,text=tgNavy] at (0,-.65)
   {Euler: tak. Hamilton: nie (wierzchołek rozcinający $c$).};
\end{tikzpicture}
\end{center}

\begin{twierdzenie}[Dirac]
\label{thm:tg-dirac}
Niech $G$ będzie skończonym grafem prostym o
$n\ge3$ wierzchołkach i minimalnym stopniu
$\delta(G)\ge n/2$. Wtedy $G$ ma cykl Hamiltona.
\end{twierdzenie}

\begin{proof}
Najpierw zauważmy, że $G$ jest spójny.
Gdyby miał co najmniej dwie składowe, w każdej
z nich liczba wierzchołków byłaby co najmniej
$\delta(G)+1\ge n/2+1$. Suma liczebności dwóch
składowych przekroczyłaby $n$, sprzeczność.

Wybierzmy ścieżkę o największej liczbie wierzchołków:
$P=v_1,v_2,\ldots,v_k$. Wszystkich sąsiadów obu
końców $v_1,v_k$ musi zawierać $P$, bo w przeciwnym
razie przedłużylibyśmy ścieżkę.
Dla $i=1,\ldots,k-1$ rozważmy dwa warunki:
$v_1v_{i+1}\in E$ oraz $v_iv_k\in E$.
Pierwszy jest spełniony dla dokładnie $d(v_1)$ indeksów,
a drugi dla dokładnie $d(v_k)$ indeksów.
Ponieważ $d(v_1)+d(v_k)\ge n>k-1$, obydwa
zbiory indeksów muszą mieć punkt wspólny.
Wybierzmy takie $i$.

Wtedy ciąg
$v_1,v_2,\ldots,v_i,v_k,v_{k-1},\ldots,v_{i+1},v_1$
stanowi cykl prosty zawierający wszystkie wierzchołki
ścieżki $P$. W szczególności $k\ge3$.

Gdyby $k<n$, spójność $G$ zapewniałaby istnienie
krawędzi między tym cyklem a pewnym wierzchołkiem
$w\notin V(P)$ (bierzemy pierwszą krawędź ścieżki
wychodzącej poza cykl). Usuwając z cyklu jedną
krawędź incydentną z wierzchołkiem sąsiadującym
z $w$ i rozpoczynając obieg od przeciwnego końca
tej krawędzi, otrzymalibyśmy ścieżkę zawierającą
wszystkie $k$ wierzchołków cyklu oraz $w$.
Byłaby dłuższa od $P$, sprzeczność.
Zatem $k=n$, a znaleziony cykl jest Hamiltona.
\end{proof}

\begin{twierdzenie}[Ore]
\label{thm:tg-ore}
Niech $G$ będzie skończonym grafem prostym o
$n\ge3$ wierzchołkach. Jeżeli dla każdej pary
różnych niesąsiadujących wierzchołków $u,v$
zachodzi $d(u)+d(v)\ge n$, to $G$ jest hamiltonowski.
\end{twierdzenie}

\begin{proof}
Wykażemy najpierw pomocniczą zasadę: jeśli $u,v$
są niesąsiadujące i $d(u)+d(v)\ge n$, to
dodanie krawędzi $uv$ nie zmienia odpowiedzi
na pytanie o istnienie cyklu Hamiltona.
Jeden kierunek jest oczywisty: cykl Hamiltona
w $G$ pozostaje cyklem w $G+uv$.

Załóżmy, że w $G+uv$ istnieje cykl Hamiltona
używający nowej krawędzi $uv$ (inaczej już jest
cyklem w $G$). Po usunięciu tej krawędzi otrzymujemy
ścieżkę Hamiltona $P=x_1,x_2,\ldots,x_n$ w $G$
z $x_1=u$ i $x_n=v$. Dla $i=1,\ldots,n-1$
rozważmy warunki $ux_{i+1}\in E$ i $vx_i\in E$.
Liczby indeksów spełniających te warunki to
odpowiednio $d(u)$ i $d(v)$ (ponieważ $u,v$
nie sąsiadują). Ponieważ ich suma wynosi co
najmniej $n>n-1$, istnieje wspólny indeks $i$.
Wtedy
$u=x_1,x_2,\ldots,x_i,v=x_n,x_{n-1},\ldots,x_{i+1},u$
jest cyklem Hamiltona w $G$.

Dodawajmy teraz kolejno wszystkie brakujące
krawędzie grafu $G$. Przy każdym kroku suma stopni
końców dodawanej krawędzi jest co najmniej $n$,
bo stopnie mogą tylko wzrastać. Z udowodnionej
zasady dodawanie żadnej z krawędzi nie zmienia
hamiltonowskości. Ostatecznie otrzymujemy $K_n$,
który ma cykl Hamiltona. Zatem również $G$ go ma.
\end{proof}

\begin{przyklad}[Porównanie kryteriów Diraca i Ore'a]
Graf $K_{3,3}$ ma $n=6$ i wszystkie stopnie równe trzy,
więc spełnia warunek Diraca i jest hamiltonowski.
Z kolei cykl $C_5$ ma stopnie równe dwa; nie spełnia
warunku Diraca, choć sam jest cyklem Hamiltona.
Warunki Diraca i Ore'a są więc wystarczające, niekonieczne.

Istnieją także grafy, dla których działa kryterium Ore'a,
ale kryterium Diraca już nie. Rozważmy graf $K_5$
po usunięciu dwóch krawędzi o wspólnym końcu $a$.
Jego stopnie to $2,3,3,4,4$, zatem
$\delta(G)=2<5/2$. Jedyne pary niesąsiadujące
to końce dwóch usuniętych krawędzi; dla każdej
suma stopni wynosi $2+3=5=n$.
Z twierdzenia Ore'a wynika hamiltonowskość.
\end{przyklad}

\begin{center}
\begin{tikzpicture}[scale=.92,line cap=round]
  \foreach \ang/\name in {90/a,18/b,-54/c,-126/d,162/e}{
    \node[tgBlueVertex] (\name) at (\ang:1.7) {$\name$};
  }
  \foreach \u/\v in {a/d,a/e,b/c,b/d,b/e,c/d,c/e,d/e}
    \draw[tgEdge] (\u)--(\v);
  \node[font=\footnotesize,text=tgNavy] at (0,-2.2)
     {$K_5$ bez dwóch krawędzi przy $a$: wariant z warunkiem Ore'a.};
\end{tikzpicture}
\end{center}

\section*{Od kombinatoryki do algebry i topologii}
\addcontentsline{toc}{section}{Od kombinatoryki do algebry i topologii}

Dotychczas rozpoznawaliśmy cykle, drzewa, rozłączne ścieżki
i skojarzenia za pomocą metod kombinatorycznych. Teraz przechodzimy
do przestrzeni łańcuchów, macierzy incydencji i operatorów Laplace'a.
Pozwoli to zrozumieć te same obiekty jako jednowymiarowe kompleksy
CW, a ich cykle jako elementy grup homologii.


% ============================================================
% ITERACJA 3. Algebra, homologia, grafy jako CW, laplasjan.
% Tekst wymaga preambuły istniejącego skryptu.
% ============================================================

\chapter{Macierze incydencji, cykle i przekroje}
\label{ch:tg-algebra}

W poprzednich rozdziałach dowiedzieliśmy się, że liczba niezależnych
cykli grafu o $n$ wierzchołkach, $m$ krawędziach i $c$ składowych wynosi
$m-n+c$. Obecnie wyjaśnimy tę formułę za pomocą algebry liniowej.
Pozwoli to ponadto przejść od zbiorów krawędzi do łańcuchów całkowitych,
co będzie pierwszym rzeczywistym zastosowaniem teorii homologii.

\section{Orientacja pomocnicza i operator brzegowy}
\label{sec:tg-boundary}

W całym rozdziale $G=(V,E)$ oznacza skończony graf prosty, z $V\ne\varnothing$.
\emph{Pomocnicza orientacja} polega na wybraniu dla każdej krawędzi
$e=\{u,v\}$ jednego z dwóch uporządkowań końców: $e:u\to v$ albo
$e:v\to u$. Nie zmieniamy w ten sposób grafu nieskierowanego;
wybór jest jedynie potrzebny do ustalenia znaków.

\begin{definicja}[Łańcuchy wymiaru zero i jeden]
Niech $R$ będzie jednym z pierścieni $\mathbb Z$, $\mathbb Q$ lub
$\mathbb R$. Definiujemy wolne moduły
$C_0(G;R):=R^V$ i $C_1(G;R):=R^E$, których bazy standardowe
oznaczamy odpowiednio przez $[v]$, $v\in V$, oraz $[e]$, $e\in E$.
Dla pomocniczej orientacji $e:u\to v$ określamy $R$-liniowy
\emph{operator brzegowy} $\partial_1:C_1(G;R)\to C_0(G;R)$ przez
$\partial_1[e]=[v]-[u]$.
\end{definicja}

Można myśleć o $[e]$ jako o odcinku posiadającym początek i koniec.
W operatorze brzegowym koniec występuje z plusem, a początek z minusem.
Jeśli odwrócimy pomocniczą orientację, odpowiadająca jej kolumna
zmieni znak; rząd macierzy ani opis jądra nie ulegną jednak zmianie.

\begin{definicja}[Macierz incydencji]
Ustalmy kolejność $v_1,\ldots,v_n$ wierzchołków i $e_1,\ldots,e_m$
krawędzi oraz orientację każdej krawędzi. Macierzą incydencji
$B_G=(b_{ij})\in\mathbb Z^{n\times m}$ nazywamy macierz operatora
$\partial_1$, to znaczy $b_{ij}=1$, gdy $v_i$ jest końcem $e_j$,
$b_{ij}=-1$, gdy jest początkiem $e_j$, i $b_{ij}=0$ w pozostałych
przypadkach.
\end{definicja}

\begin{center}
\begin{tikzpicture}[scale=.97,>=Stealth,line cap=round]
 \node[tgBlueVertex] (a) at (-2,0) {$a$};
 \node[tgBlueVertex] (b) at (0,1.4) {$b$};
 \node[tgBlueVertex] (c) at (2,0) {$c$};
 \draw[tgTree,-{Stealth[length=2.8mm]}] (a)--node[above left,font=\small] {$e_1$}(b);
 \draw[tgTree,-{Stealth[length=2.8mm]}] (b)--node[above right,font=\small] {$e_2$}(c);
 \draw[tgAccent,-{Stealth[length=2.8mm]}] (a)--node[below,font=\small] {$e_3$}(c);
 \node[text=tgNavy,font=\footnotesize] at (0,-.92)
   {Wybór strzałek ustala znaki, a nie typ grafu.};
\end{tikzpicture}
\end{center}

\begin{przyklad}[Trójkąt: macierz i brzeg cyklu]
Dla orientacji $e_1:a\to b$, $e_2:b\to c$, $e_3:a\to c$ mamy
\[
B_G=\begin{pmatrix}-1&0&-1\\1&-1&0\\0&1&1\end{pmatrix}.
\]
Kolumny sumują się do zera. Wektor $z=(1,1,-1)^T$ spełnia
$B_Gz=0$: jest to zamknięty obieg $a\to b\to c\to a$, w którym
trzecią krawędź przechodzimy przeciwnie do jej orientacji pomocniczej.
Natomiast $B_G(1,0,0)^T=(-1,1,0)^T$ opisuje łańcuch mający początek
w $a$ i koniec w $b$.
\end{przyklad}

\section{Rząd macierzy i homologia zerowymiarowa}
\label{sec:tg-rank-incidence}

\begin{twierdzenie}[Rząd macierzy incydencji]
\label{thm:tg-rank}
Niech $G$ będzie skończonym grafem prostym o $n\ge1$ wierzchołkach,
$m$ krawędziach i $c$ składowych spójnych. Niech $B_G$ będzie
macierzą incydencji dla dowolnej pomocniczej orientacji. Wtedy
$\operatorname{rank}_{\mathbb R}B_G=n-c$ i
$\dim_{\mathbb R}\ker B_G=m-n+c$.
\end{twierdzenie}

\begin{proof}
Oznaczmy zbiory wierzchołków składowych przez $V_1,\ldots,V_c$.
Każda kolumna $B_G$ ma sumę współrzędnych równą zero na każdej
składowej. Zatem obraz $B_G$ zawiera się w podprzestrzeni
\[
W=\left\{x\in\mathbb R^V:\sum_{v\in V_i}x_v=0
\text{ dla }i=1,\ldots,c\right\}.
\]
Warunki definiujące $W$ są niezależne, ponieważ zbiory $V_i$ są
niepuste i rozłączne; zatem $\dim W=n-c$.

Pokażemy teraz równość $\operatorname{im}B_G=W$.
W każdej składowej wybierzmy korzeń $r_i$ i drzewo rozpinające.
Dla $v\in V_i\setminus\{r_i\}$ istnieje jednoznaczna ścieżka
w drzewie od $r_i$ do $v$. Sumując jej krawędzie ze współczynnikami
$+1$ albo $-1$ według zgodności z orientacją pomocniczą, otrzymujemy
łańcuch $z_v\in C_1(G;\mathbb R)$ o brzegu
$\partial_1z_v=[v]-[r_i]$: wkłady pośrednich wierzchołków znoszą się
parami. Wektory $[v]-[r_i]$, dla wszystkich $v\ne r_i$, tworzą bazę
$W$. W szczególności $W\subseteq\operatorname{im}B_G$, co daje
pierwszą tezę. Z twierdzenia o rzędzie i jądrze otrzymujemy
$\dim\ker B_G=m-(n-c)=m-n+c$.
\end{proof}

\begin{uwaga}[Przypadki brzegowe]
Dla pojedynczego wierzchołka bez krawędzi macierz ma wymiar
$1\times0$, rząd zero i jądro zerowe. Dla grafu mającego tylko
wierzchołki izolowane zachodzi $c=n$ i $B_G=0$.
Nie trzeba więc wyłączać tych przypadków z twierdzenia.
\end{uwaga}

\begin{definicja}[Homologie grafu]
Dla skończonego grafu $G$ i $R\in\{\mathbb Z,\mathbb Q,\mathbb R\}$
przyjmujemy $C_k(G;R)=0$ dla $k\notin\{0,1\}$.
Definiujemy $H_1(G;R):=\ker\partial_1$ oraz
$H_0(G;R):=C_0(G;R)/\operatorname{im}\partial_1$.
Są to grupy (odpowiednio moduły) homologii komórkowej
jednowymiarowej realizacji geometrycznej $|G|$.
\end{definicja}

Powyższa definicja jest szczególnym przypadkiem ogólnego kompleksu
łańcuchowego. Ponieważ $C_2=0$, grupa pierwszych brzegów
$\operatorname{im}\partial_2$ jest zerowa, stąd $H_1=\ker\partial_1$.

\begin{twierdzenie}[Homologia grafu ze współczynnikami całkowitymi]
\label{thm:tg-homology}
Niech $G$ będzie skończonym grafem prostym o $n$ wierzchołkach,
$m$ krawędziach i $c$ składowych spójnych. Wówczas
\[
H_0(G;\mathbb Z)\cong\mathbb Z^c,\qquad
H_1(G;\mathbb Z)\cong\mathbb Z^{m-n+c},\qquad
H_k(G;\mathbb Z)=0\quad(k\ge2).
\]
Izomorfizmy nie wymagają wyboru orientacji krawędzi, choć
konkretne bazy pierwszej homologii zależą od wyboru drzewa.
\end{twierdzenie}

\begin{proof}
Rozważmy homologię zerową. Dla każdej składowej $V_i$ określamy
$\varepsilon_i:\mathbb Z^V\to\mathbb Z$, sumując współczynniki
wierzchołków należących do $V_i$. Łączny homomorfizm
$\varepsilon=(\varepsilon_1,\ldots,\varepsilon_c)$ jest surjektywny
i znika na $\operatorname{im}\partial_1$. Ponadto jego jądro jest
równe temu obrazowi: jeśli $x\in\mathbb Z^V$ ma sumę współczynników
zero na każdej składowej, to można zapisać
$x=\sum_i\sum_{v\in V_i\setminus\{r_i\}}x_v([v]-[r_i])$.
Każdy składnik jest brzegiem całkowitego łańcucha odpowiadającego
ścieżce drzewowej od $r_i$ do $v$. Zatem
$\mathbb Z^V/\operatorname{im}\partial_1\cong\mathbb Z^c$.

Dla homologii pierwszej wybieramy las rozpinający $T$.
Dla każdej krawędzi $e\in E\setminus E(T)$ niech $z_e$
będzie całkowitym łańcuchem złożonym z $e$ ze współczynnikiem
$1$ i z odpowiednio zorientowanej ścieżki w $T$ zamykającej jego
końce. Wtedy $\partial_1z_e=0$. Wektory $z_e$ są liniowo
niezależne nad $\mathbb Z$, gdyż współrzędna przy krawędzi
$f\notin E(T)$ wektorów $z_e$ wynosi $\delta_{ef}$.

Aby udowodnić, że tworzą bazę całego jądra, weźmy
$z\in\ker\partial_1$ i odejmijmy od niego łańcuch
$\sum_{e\notin E(T)}z_e\,z(e)$, gdzie $z(e)\in\mathbb Z$
jest współczynnikiem $z$ przy $e$. Różnica jest cyklem,
którego nośnik leży w lesie $T$. Taki cykl musi być zerowy:
jeśli jego nośnik jest niepusty, wybieramy liść tego nośnika;
współczynnik przy jego jedynej krawędzi pojawia się w brzegu
z niezerowym znakiem, sprzecznie z $\partial_1z=0$.
Zatem $\{z_e:e\notin E(T)\}$ jest bazą całkowitego
jądra, liczącą $m-(n-c)$ elementów. W wyższych wymiarach
brak komórek, więc pozostałe grupy znikają.
\end{proof}

\section{Przestrzeń cykli modulo dwa i przestrzeń cięć}
\label{sec:tg-cut-cycle}

Do opisu cyklu jako zbioru krawędzi nie potrzebujemy znaków.
Wystarczy przejść do $\mathbb F_2$, w którym $-1=1$.
Niech $B_2$ oznacza macierz incydencji zredukowaną modulo $2$.

\begin{stwierdzenie}[Porównanie z przestrzenią cykli]
Niech $G$ będzie skończonym grafem prostym i niech
$\mathcal C(G)$ będzie przestrzenią cykli zdefiniowaną
w rozdziale \ref{ch:tg-drzewa}. Wtedy
$\mathcal C(G)=\ker(B_2:\mathbb F_2^E\to\mathbb F_2^V)$.
W szczególności ma ona wymiar $m-n+c$.
\end{stwierdzenie}

\begin{proof}
Dla zbioru krawędzi $F\subseteq E$ współrzędna
$B_2\mathbf1_F$ przy wierzchołku $v$ jest równa
stopniowi $v$ w $(V,F)$ modulo $2$.
Każdy cykl prosty ma wszystkie stopnie parzyste, a więc
$\mathcal C(G)\subseteq\ker B_2$.
Odwrotnie, wektor należący do jądra odpowiada zbiorowi
krawędzi o parzystych stopniach. Lemat o rozkładzie takiego
zbioru na cykle, udowodniony w rozdziale
\ref{ch:tg-drzewa}, daje jego przynależność do
$\mathcal C(G)$. Wymiar otrzymujemy z twierdzenia
o bazie cykli fundamentalnych.
\end{proof}

\begin{definicja}[Przestrzeń przekrojów]
Dla grafu $G=(V,E)$ nad $\mathbb F_2$ określamy
$\mathcal D(G)\subseteq\mathbb F_2^E$ jako przestrzeń
liniową rozpiętą przez funkcje charakterystyczne
przekrojów $\delta_G(S)$, gdzie $S\subseteq V$.
\end{definicja}

\begin{twierdzenie}[Ortogonalność cykli i przekrojów]
\label{thm:tg-orthogonal}
Dla dowolnego skończonego grafu prostego $G$ zachodzi
\[
\mathcal D(G)=\operatorname{im}B_2^T
  =\mathcal C(G)^\perp,\qquad
\mathcal C(G)=\mathcal D(G)^\perp,
\]
gdzie ortogonalność jest rozumiana względem
standardowego iloczynu współrzędnych nad $\mathbb F_2$.
Ponadto $\dim\mathcal D(G)=n-c$.
\end{twierdzenie}

\begin{proof}
Dla $S\subseteq V$ wektor $B_2^T\mathbf1_S$
ma przy krawędzi $uv$ współrzędną równą
$\mathbf1_S(u)+\mathbf1_S(v)$ modulo dwa. Jest ona równa
jeden dokładnie wtedy, gdy krawędź przecina podział
$S\sqcup(V\setminus S)$. To pokazuje
$\mathcal D(G)=\operatorname{im}B_2^T$.

Jeśli $z\in\ker B_2$ i $y=B_2^Tf$, to
$\langle z,y\rangle=\langle B_2z,f\rangle=0$.
Stąd $\mathcal D(G)\subseteq\mathcal C(G)^\perp$.
Rząd macierzy incydencji nad $\mathbb F_2$ jest nadal
równy $n-c$: dowód z drzewami działa również modulo dwa.
Tak więc $\dim\mathcal D(G)=n-c$, podczas gdy
$\dim\mathcal C(G)^\perp=m-(m-n+c)=n-c$.
Daje to pierwszą równość, a drugą otrzymujemy przez
ponowne przejście do dopełnienia ortogonalnego,
ponieważ standardowe parowanie na $\mathbb F_2^E$
jest niezdegenerowane.
\end{proof}

\begin{przyklad}[Kwadrat z przekątną: dwa niezależne cykle]
Niech $V=\{1,2,3,4\}$, a krawędzie będą skierowane
$e_1:1\to2$, $e_2:2\to3$, $e_3:3\to4$,
$e_4:4\to1$, $e_5:1\to3$. Wówczas
\[
B=\begin{pmatrix}
-1&0&0&1&-1\\1&-1&0&0&0\\
0&1&-1&0&1\\0&0&1&-1&0
\end{pmatrix}.
\]
Dla drzewa rozpinającego $T=\{e_1,e_2,e_3\}$
dwa fundamentalne cykle mają wektory
$z_4=(1,1,1,1,0)^T$ i $z_5=(-1,-1,0,0,1)^T$.
Bezpośrednie mnożenie daje $Bz_4=Bz_5=0$.
Ponieważ $m-n+1=5-4+1=2$, tworzą one bazę
$H_1(G;\mathbb Z)\cong\mathbb Z^2$.

W wersji modulo dwa mamy wektory
$\bar z_4=(1,1,1,1,0)$ oraz $\bar z_5=(1,1,0,0,1)$.
Ich suma odpowiada trójkątowi na wierzchołkach $1,3,4$.
Wyróżnione dwa cykle są bazą, ale nie są wszystkimi
cyklami grafu.
\end{przyklad}

\begin{center}
\begin{tikzpicture}[scale=1.0,line cap=round]
 \node[tgBlueVertex] (a) at (-1.6,1.0) {$1$};
 \node[tgBlueVertex] (b) at (1.6,1.0) {$2$};
 \node[tgBlueVertex] (c) at (1.6,-1.0) {$3$};
 \node[tgBlueVertex] (d) at (-1.6,-1.0) {$4$};
 \draw[tgTree] (a)--node[above,font=\scriptsize] {$e_1$}(b);
 \draw[tgTree] (b)--node[right,font=\scriptsize] {$e_2$}(c);
 \draw[tgTree] (c)--node[below,font=\scriptsize] {$e_3$}(d);
 \draw[tgAccent] (d)--node[left,font=\scriptsize] {$e_4$}(a);
 \draw[tgAccent] (a)--node[pos=.68,above,font=\scriptsize] {$e_5$}(c);
 \node[text=tgNavy,font=\footnotesize] at (0,-1.53)
   {Zielone: drzewo $T$. Czerwone: krawędzie generujące cykle.};
\end{tikzpicture}
\end{center}

\begin{uwaga}[Ortogonalność nad $\mathbb F_2$ nie oznacza sumy prostej]
Równości $\mathcal C(G)=\mathcal D(G)^\perp$ oraz
$\mathcal D(G)=\mathcal C(G)^\perp$ nie implikują,
że ich przecięcie jest zerowe. Przykładowo w cyklu
$C_4$ zbiór wszystkich czterech krawędzi jest cyklem,
a jednocześnie jest przekrojem $\delta(\{v_1,v_3\})$.
Zatem należy odróżnić te równości nad $\mathbb F_2$
od rzeczywistego ortogonalnego rozkładu na gradienty i cykle.
\end{uwaga}

\section{Rozkład ortogonalny w przestrzeni krawędzi}
\label{sec:tg-linear-hodge}

Zastąpmy teraz $\mathbb F_2$ przez $\mathbb R$ i nadajmy
przestrzeniom $C_0(G;\mathbb R)$, $C_1(G;\mathbb R)$
standardowe iloczyny skalarne, w których bazy wierzchołków
i krawędzi są ortonormalne. Wtedy adjungowanym operatorem
do $B$ jest $B^T$.

\begin{twierdzenie}[Ortogonalny rozkład Hodge'a na grafie]
\label{thm:tg-edge-hodge}
Dla skończonego grafu prostego $G$ zachodzi suma prosta ortogonalna
\[
\mathbb R^E=\operatorname{im}B^T\ \oplus^\perp\ \ker B.
\]
W szczególności każdy łańcuch krawędziowy $x$ można
jednoznacznie zapisać jako $x=B^Tf+h$ z $Bh=0$.
Wektor $f$ jest wyznaczony z dokładnością do dodania funkcji
stałej na każdej składowej grafu.
\end{twierdzenie}

\begin{proof}
Dla $h\in\ker B$ i dowolnego $f\in\mathbb R^V$ mamy
$\langle h,B^Tf\rangle=\langle Bh,f\rangle=0$.
Zatem obie podprzestrzenie są ortogonalne.
Ich wymiary są równe odpowiednio
$m-n+c$ i $n-c$, co wynika z twierdzenia
\ref{thm:tg-rank}. Suma wymiarów jest równa $m$,
więc otrzymujemy rozkład całej przestrzeni.

Jeśli $x=B^Tf+h$, zastosowanie $B$ daje
$BB^Tf=Bx$; istnienie $f$ zapewnia
$Bx\in\operatorname{im}B=\operatorname{im}BB^T$.
Ostatnia równość zachodzi, gdyż
$\ker(BB^T)=\ker B^T$ (wynika to z tożsamości
$\langle BB^Tf,f\rangle=\|B^Tf\|^2$), a więc
$BB^T$ ma ten sam rząd co $B$ i obraz zawarty w
$\operatorname{im}B$.
Jednoznaczność składnika $B^Tf$ wynika z ortogonalności.
Jeśli $B^Tf=B^Tg$, to $B^T(f-g)=0$, co oznacza,
że funkcja $f-g$ przyjmuje tę samą wartość na obu
końcach każdej krawędzi, a więc jest stała na składowych.
\end{proof}

Zauważmy istotne rozróżnienie: przestrzenie cykli nad
$\mathbb F_2$ i $\mathbb R$ mają ten sam wymiar,
ale nie są tym samym obiektem. Nad $\mathbb R$ kierunek
przebiegu krawędzi i znaki współczynników są istotne;
nad $\mathbb F_2$ znak nie występuje. W następnym
rozdziale utożsamimy część $\ker B$ z pierwszą homologią
przestrzeni topologicznej.

\chapter{Graf jako przestrzeń topologiczna}
\label{ch:tg-topology}

Dotychczas traktowaliśmy graf jako zbiór wierzchołków i krawędzi.
Od teraz każdą krawędź zastępujemy rzeczywistym odcinkiem.
Tak powstała przestrzeń jest jednowymiarowym kompleksem CW.
Dla topologii ważne jest, że nawet jeśli narysujemy graf
na płaszczyźnie z przecięciami, nie wprowadzamy tam
nowych wierzchołków bez wyraźnej identyfikacji.

\section{Kompleks CW, charakterystyka Eulera i zwijanie drzewa}
\label{sec:tg-collapse-tree}

Przypomnijmy konstrukcję realizacji geometrycznej $|G|$
z pierwszego rozdziału: dla każdego $v\in V$ bierzemy
$0$-komórkę, a dla każdej krawędzi $e$ odcinek z końcami
doklejonymi do odpowiednich wierzchołków. Krawędzie
równoległe i pętle również dają poprawne CW-kompleksy,
gdy zachowamy osobne odcinki i odrębnie określimy mapy
przyklejania ich końców.

\begin{definicja}[Charakterystyka Eulera]
Dla skończonego kompleksu CW $X$ oznaczamy przez $c_k(X)$
liczbę jego otwartych $k$-komórek i definiujemy
$\chi(X)=\sum_{k\ge0}(-1)^k c_k(X)$.
Dla realizacji skończonego grafu $G$ otrzymujemy
$\chi(|G|)=|V(G)|-|E(G)|$.
\end{definicja}

\begin{lemat}[Zwinięcie kontraktowalnego podkompleksu]
\label{lem:tg-collapse-CW}
Niech $X$ będzie kompleksem CW, a $A\subseteq X$
niepustym, kontraktowalnym podkompleksem CW.
Wtedy naturalne odwzorowanie ilorazowe $q:X\to X/A$
jest równoważnością homotopijną.
\end{lemat}

\begin{proof}
Ponieważ $(X,A)$ jest parą CW, ma własność przedłużania
homotopii: każdą homotopię $A\times[0,1]\to X$,
rozpoczynającą się od włożenia $A\hookrightarrow X$,
można przedłużyć do homotopii $H:X\times[0,1]\to X$
z $H_0=\operatorname{id}_X$. Wybieramy kontrakcję
$A$ do punktu $a_0\in A$, tak aby $H_t(A)\subseteq A$
i $H_1(A)=\{a_0\}$. Wtedy $H_1$ jest stałe na $A$,
więc faktoryzuje się przez iloraz jako $H_1=g\circ q$
dla pewnego ciągłego $g:X/A\to X$.

Z jednej strony $gq=H_1\simeq\operatorname{id}_X$.
Z drugiej strony $qH_t$ jest stałe na $A$ dla każdego $t$,
więc daje homotopię na ilorazie od
$\operatorname{id}_{X/A}$ do $qg$.
Stąd również $qg\simeq\operatorname{id}_{X/A}$.
Obie równości homotopijne dowodzą tezy.
\end{proof}

\begin{uwaga}
W dowodzie użyto standardowej własności par CW.
Zwijanie dowolnego kontraktowalnego podzbioru przestrzeni
bez założenia o podkompleksie lub odpowiednim warunku
kofibracji nie musi dawać równoważności homotopijnej. Warunek ten jest istotny.
\end{uwaga}

\begin{twierdzenie}[Graf jest homotopijnie równoważny bukietowi okręgów]
\label{thm:tg-wedge}
Niech $G$ będzie skończonym spójnym grafem prostym
z $n\ge1$ wierzchołkami i $m$ krawędziami.
Dla dowolnego drzewa rozpinającego $T\subseteq G$
przestrzeń $|G|$ jest homotopijnie równoważna bukietowi
$\bigvee_{j=1}^{m-n+1}S^1$.
Przy zerowej liczbie okręgów bukiet oznacza pojedynczy punkt.
\end{twierdzenie}

\begin{proof}
Drzewo $T$ jest kontraktowalne. Aby to zobaczyć,
można kolejno usuwać otwarte krawędzie przylegające
do liścia, przesuwając każdy liść wzdłuż jego jedynej
krawędzi do sąsiedniego wierzchołka. Ponieważ drzewo
jest skończone, po skończonej liczbie deformacji
pozostaje pojedynczy wierzchołek.

Z lematu \ref{lem:tg-collapse-CW} iloraz
$q:|G|\to|G|/|T|$ jest równoważnością homotopijną.
W ilorazie wszystkie wierzchołki drzewa, czyli wszystkie
wierzchołki $G$, stają się jednym punktem.
Każda krawędź spoza $T$ pozostaje odcinkiem, którego
oba końce zostały ze sobą utożsamione, a więc tworzy
okrąg. Różne okręgi mają wspólny tylko ten punkt.
Ich liczba wynosi $m-(n-1)=m-n+1$.
\end{proof}

\begin{center}
\begin{tikzpicture}[scale=.96,line cap=round,>=Stealth]
 \node[tgBlueVertex] (a) at (-3.9,.8) {$a$};
 \node[tgBlueVertex] (b) at (-2.1,1.25) {$b$};
 \node[tgBlueVertex] (c) at (-1.8,-.65) {$c$};
 \node[tgBlueVertex] (d) at (-3.7,-.8) {$d$};
 \draw[tgTree] (a)--(b)--(c)--(d);
 \draw[tgAccent] (d)--(a) (a)--(c);
 \draw[tgTeal,-{Stealth[length=2.6mm]},line width=1.6pt] (-.85,.15)--(.25,.15);
 \node[font=\footnotesize,text=tgNavy] at (-.25,.51) {zwiń $T$};
 \node[tgBlueVertex] (o) at (2.0,.15) {$*$};
 \draw[tgAccent,line width=2pt] (o) .. controls (1.2,1.9) and (3.6,1.9) .. (o);
 \draw[tgAccent,line width=2pt] (o) .. controls (1.2,-1.4) and (3.6,-1.4) .. (o);
 \node[font=\footnotesize,text=tgNavy] at (1.95,-1.78) {bukiet dwóch okręgów};
\end{tikzpicture}

{\footnotesize Zwijamy zielone drzewo. Dwie czerwone krawędzie
pozostają generatorami dwóch niezależnych okręgów.}
\end{center}

\section{Grupa fundamentalna grafu}
\label{sec:tg-pi1}

\begin{definicja}[Grupa wolna]
Dla zbioru $S$ przez $F(S)$ oznaczamy grupę, której elementami
są zredukowane słowa z liter $s$ i $s^{-1}$, $s\in S$,
z mnożeniem przez konkatenację i redukcję sąsiadujących par
$s s^{-1}$ lub $s^{-1}s$. Dla $|S|=r<\infty$ piszemy $F_r$.
Grupą $F_0$ jest grupa trywialna.
\end{definicja}

\begin{twierdzenie}[Grupa fundamentalna skończonego grafu]
\label{thm:tg-free-pi1}
Niech $G$ będzie skończonym spójnym grafem prostym,
a $x_0\in|G|$ dowolnym punktem bazowym. Jeżeli
$\beta=m-n+1$, to $\pi_1(|G|,x_0)\cong F_\beta$.
W szczególności $|G|$ jest po prostu spójna wtedy i tylko
wtedy, gdy $G$ jest drzewem.
\end{twierdzenie}

\begin{proof}
Wybierzmy drzewo rozpinające $T$ i wierzchołek bazowy
$r\in T$. Poprzednie twierdzenie daje równoważność
homotopijną $|G|\simeq\bigvee_\beta S^1$.
Grupa fundamentalna bukietu $\beta$ okręgów jest wolna
na $\beta$ generatorach. Przypomnijmy uzasadnienie
przez twierdzenie van Kampena: dla każdego okręgu
wybieramy otoczenie jego wnętrza wraz z niewielkim wspólnym
otoczeniem punktu sklejenia. Każde takie otoczenie wnosi
jeden generator $\mathbb Z$, a przecięcia, po odpowiednim
zmniejszeniu otoczeń, są ściągalne, więc nie nakładają
żadnych dodatkowych relacji. Iteracja van Kampena
daje wolny iloczyn $\mathbb Z*\cdots*\mathbb Z=F_\beta$.

Generatorowi odpowiada pętla: od $r$ idziemy po jedynej
ścieżce drzewowej do początku wybranej krawędzi
$e\notin T$, przechodzimy $e$ i wracamy po drzewie do $r$.
Zmiana punktu bazowego wewnątrz spójnego $|G|$ nie zmienia
typu izomorfizmu grupy fundamentalnej.

Jeśli $\beta=0$, graf nie ma cykli i jest drzewem.
Gdy $\beta\ge1$, grupa $F_\beta$ zawiera nietrywialną
pętlę, a więc $|G|$ nie jest po prostu spójna.
\end{proof}

\begin{przyklad}[Trójkąt i graf theta]
Dla trójkąta $K_3$ mamy $n=m=3$ i $\beta=1$.
Zatem $|K_3|\simeq S^1$ oraz
$\pi_1(|K_3|)\cong\mathbb Z$.

Graf theta składa się z dwóch wierzchołków połączonych
trzema różnymi krawędziami (jest to multigraf bez pętli).
Ma $n=2$, $m=3$, więc $\beta=2$.
Jego realizacja jest homotopijnie równoważna bukietowi
dwóch okręgów, a grupa fundamentalna jest izomorficzna
z wolną grupą $F_2$, która jest nieprzemienna.
To pokazuje, że przestrzenie jednowymiarowe mogą mieć
nietrywialne, nawet nieprzemienne grupy fundamentalne.
\end{przyklad}

\begin{center}
\begin{tikzpicture}[scale=.98,line cap=round]
 \node[tgBlueVertex] (u) at (-1.75,0) {$u$};
 \node[tgBlueVertex] (v) at (1.75,0) {$v$};
 \draw[tgBlue,line width=2.1pt] (u) to[bend left=48] (v);
 \draw[tgTeal,line width=2.1pt] (u)--(v);
 \draw[tgCoral,line width=2.1pt] (u) to[bend right=48] (v);
 \node[text=tgNavy,font=\footnotesize] at (0,-1.22) {graf theta, $\pi_1\cong F_2$};
\end{tikzpicture}
\end{center}

\section{Nakrycie uniwersalne i drzewa}
\label{sec:tg-universal-cover}

Sytuacja staje się jeszcze ciekawsza, gdy zastąpimy skończony graf
jego nakryciem. Nakrycie jest lokalnie takie samo, ale może mieć
nieskończenie wiele wierzchołków i krawędzi.

\begin{definicja}[Nakrycie grafów]
Dla grafów $\widetilde G,G$ o realizacjach geometrycznych
odwzorowanie ciągłe i surjektywne
$p:|\widetilde G|\to|G|$ nazywamy
\emph{nakryciem grafowym}, jeżeli przenosi wierzchołki na
wierzchołki, każdą otwartą krawędź homeomorficznie na
otwartą krawędź, jest zgodne z odwzorowaniem końców
krawędzi oraz dla każdego wierzchołka
$\widetilde v$ indukuje bijekcję między zbiorami
\emph{incydencji} (końców krawędzi wychodzących z
$\widetilde v$) i $p(\widetilde v)$.
Dopuszczamy także grafy nieskończone, lokalnie skończone;
każda pętla wnosi do zbioru incydencji dwa końce.
\end{definicja}

Dla realizacji jednowymiarowych podany warunek jest dokładnie
lokalnym warunkiem nakryciowym: małe gwiazdowe otoczenie
każdego wierzchołka ma rozłączne kopie w przestrzeni nakrywającej.

\begin{twierdzenie}[Nakrycie uniwersalne realizacji grafu]
\label{thm:tg-universal}
Każdy spójny, skończony graf $G$ posiada spójne nakrycie
$p:|\widetilde G|\to|G|$, w którym $\widetilde G$ jest drzewem.
Nakrycie to jest uniwersalne.
\end{twierdzenie}

\begin{proof}[Konstrukcja ze zredukowanych ścieżek]
Wybierzmy wierzchołek $r\in G$ i orientację pomocniczą
krawędzi. Wierzchołkami $\widetilde G$ uczynimy wszystkie
\emph{zredukowane spacery krawędziowe} rozpoczynające się w $r$:
nie dopuszczamy dwóch kolejnych kroków po tej samej krawędzi
w przeciwnych kierunkach. Spacer pusty jest korzeniem $\widetilde r$.
Każdy taki spacer kończy się w pewnym wierzchołku $v\in G$.
Każdy niepusty zredukowany spacer $w$ określa dokładnie
jedną krawędź w $\widetilde G$, łączącą $w$ ze spacerem
powstałym przez skreślenie jego ostatniego kroku.
Tym samym nie liczymy tej samej krawędzi dwukrotnie.
Z punktu $w$ można wykonać każdy krok po incydentnej
krawędzi $e$ w kierunku od wierzchołka końcowego $w$:
jeśli jest to odwrotność ostatniego kroku,
przechodzimy do poprzednika, a w przeciwnym razie
do spaceru uzyskanego przez dopisanie tego kroku.
W przypadku pętli rozróżniamy oba kierunki jej przejścia.
Różne krawędzie równoległe w $G$ pozostają rozróżnione.

Rzutowanie $p$ przyporządkowuje spacerowi jego koniec w $G$,
a każdej nowej krawędzi odpowiadającą jej krawędź bazową.
Przy każdym wierzchołku $w$ i każdym możliwym kroku w $G$
istnieje dokładnie jedno odpowiednie przejście w $\widetilde G$.
Stąd $p$ spełnia warunek lokalnej bijekcji i jest nakryciem.

Każdy niepusty zredukowany spacer ma jednego poprzednika:
otrzymujemy go przez usunięcie ostatniego kroku.
Każda inna krawędź wychodząca z tego wierzchołka wydłuża
spacer o jeden krok. Wobec tego każdy wierzchołek można
połączyć z $\widetilde r$ przez jednoznaczny ciąg poprzedników.
Cykl w $\widetilde G$ musiałby mieć wierzchołek największej
długości spaceru; miałby on na cyklu dwóch różnych sąsiadów
o mniejszej długości, mimo że istnieje tylko jeden poprzednik.
To niemożliwe. Zatem $\widetilde G$ jest spójnym drzewem.

Realizacja drzewa jest kontraktowalna, więc po prostu spójna.
Z definicji nakrycia uniwersalnego otrzymujemy wniosek.
\end{proof}

\begin{uwaga}[Skończoność nakrycia]
Jeśli bazowy graf jest drzewem, jego nakrycie uniwersalne
jest z nim izomorficzne. Jeśli graf ma cykl, nakrycie
uniwersalne jest nieskończone: zamknięty zredukowany obieg
można przechodzić dowolną liczbę razy, otrzymując
coraz dłuższe różne zredukowane spacery.
\end{uwaga}

\section{Nakrycia skończone i wzór Nielsena--Schreiera}
\label{sec:tg-finite-cover}

\begin{twierdzenie}[Charakterystyka Eulera nakrycia]
Niech $G$ będzie skończonym spójnym grafem, a
$p:|\widehat G|\to|G|$ spójnym, $d$-kartkowym nakryciem,
gdzie $d\ge1$ i $\widehat G$ jest grafem skończonym.
Wtedy $|V(\widehat G)|=d|V(G)|$ oraz
$|E(\widehat G)|=d|E(G)|$, a zatem
$\chi(|\widehat G|)=d\chi(|G|)$.
Jeśli $\pi_1(|G|)\cong F_r$, to
$\pi_1(|\widehat G|)\cong F_{1+d(r-1)}$.
\end{twierdzenie}

\begin{proof}
Każdy wierzchołek $v\in G$ ma dokładnie $d$ przeciwobrazów,
bo liczba kartek nakrycia jest stała na spójnej bazie.
Nad wnętrzem każdej krawędzi otrzymujemy dokładnie $d$
rozłącznych otwartych odcinków, stanowiących wnętrza
krawędzi $\widehat G$. To daje pierwsze dwie równości.
Odejmując liczby komórek wymiaru jeden od liczby
komórek wymiaru zero, dostajemy równość dla $\chi$.

Graf $\widehat G$ jest spójny, zatem jego liczba cyklomatyczna
wynosi
$|E(\widehat G)|-|V(\widehat G)|+1
=d(m-n)+1=1+d(r-1)$.
Z twierdzenia \ref{thm:tg-free-pi1} grupa fundamentalna
$\widehat G$ jest wolna o tej liczbie generatorów.
\end{proof}

\begin{przyklad}[Nakrycie trzykartkowe bukietu]
Niech $G$ będzie bukietem dwóch okręgów, przedstawionym jako
graf z jednym wierzchołkiem i dwiema pętlami. Dla dowolnego
spójnego nakrycia trzykartkowego $\widehat G$ mamy
$|V(\widehat G)|=3$, $|E(\widehat G)|=6$, a więc
$\pi_1(|\widehat G|)\cong F_{6-3+1}=F_4$.
W języku grup oznacza to, że podgrupa indeksu trzy
grupy wolnej rzędu dwa jest wolna rzędu cztery.
\end{przyklad}

\begin{uwaga}[Subdywizje a nakrycia multigrafów]
Twierdzenie o grupie fundamentalnej sformułowaliśmy wcześniej
dla grafów prostych. Jego dowód działa również dla multigrafów
i pętli: wystarczy podzielić każdą krawędź w dwóch różnych
punktach, otrzymując graf prosty o tej samej realizacji
topologicznej. Liczba $m-n+c$ nie zmienia się przy pojedynczej
subdywizji. Dlatego obliczenie dla bukietu z pętlami
jest szczególnym przypadkiem tej samej teorii.
\end{uwaga}

\section{Metryka długości na grafie}
\label{sec:tg-metric}

\begin{definicja}[Metryka długości]
Niech $G$ będzie spójnym grafem skończonym. Nadajmy każdej
krawędzi długość jeden i utożsammy ją z odcinkiem $[0,1]$.
Długość krzywej w $|G|$ liczymy odcinkami euklidesowymi;
przez $d_{|G|}(x,y)$ oznaczamy infimum długości krzywych
łączących $x,y\in|G|$.
\end{definicja}

\begin{stwierdzenie}[Zgodność metryki grafowej z metryką długości]
Dla spójnego skończonego grafu $G$ i dowolnych wierzchołków
$u,v\in V(G)$ zachodzi $d_{|G|}(u,v)=d_G(u,v)$.
\end{stwierdzenie}

\begin{proof}
Kombinatoryczna najkrótsza ścieżka długości $k=d_G(u,v)$
po sklejeniu odcinków jest krzywą w $|G|$ długości $k$,
a zatem $d_{|G|}(u,v)\le k$.

Dla nierówności odwrotnej określmy na wierzchołkach funkcję
$f(w)=d_G(u,w)$. Na końcach każdej krawędzi $ab$ mamy
$|f(a)-f(b)|\le1$, na mocy nierówności trójkąta dla
metryki grafowej. Przedłużmy $f$ liniowo na każdej krawędzi.
Tak uzyskana funkcja jest $1$-Lipschitzowska na każdym
odcinku o długości jeden. Wzdłuż dowolnej prostowalnej
krzywej $\gamma$ w $|G|$ zmiana jej wartości nie może
przekroczyć długości $\gamma$: wynika to przez podział
krzywej na fragmenty zawarte w krawędziach, a następnie
przez aproksymację dowolnej krzywej skończonymi takimi
fragmentami. W konsekwencji
$k=|f(v)-f(u)|\le\operatorname{len}(\gamma)$.
Biorąc infimum po wszystkich $\gamma$, otrzymujemy
$d_{|G|}(u,v)\ge k$.
\end{proof}

\chapter{Laplasjan grafu, dyskretna teoria Hodge'a i twierdzenie Kirchhoffa}
\label{ch:tg-laplacian}

W pierwszym rozdziale pojawiła się macierz sąsiedztwa, a w szóstym
macierz incydencji. Z ich połączenia otrzymamy operator, który
jest dyskretnym odpowiednikiem laplasjanu na rozmaitości.
W niniejszym rozdziale wykazujemy związek tego operatora
ze spójnością, potencjałem harmonicznym, rozkładem Hodge'a
oraz liczbą drzew rozpinających.

\section{Definicja i energia Dirichleta}
\label{sec:tg-lap-definition}

Niech $G=(V,E)$ będzie skończonym grafem prostym.
Każdej krawędzi nadajmy dodatnią wagę $w_e>0$.
Macierz diagonalną tych wag oznaczamy przez $W$.

\begin{definicja}[Laplasjan ważony]
Niech $B\in\mathbb R^{n\times m}$ będzie macierzą incydencji
grafu $G$ dla dowolnej orientacji pomocniczej. Definiujemy
\emph{laplasjan ważony} $L:=BWB^T:\mathbb R^V\to\mathbb R^V$.
Gdy wszystkie $w_e=1$, piszemy $L_G:=BB^T$ i mówimy
po prostu o laplasjanie grafu.
\end{definicja}

Dla grafu prostego nieważonego macierz ma postać
$L_G=D-A_G$, gdzie $D$ jest diagonalna i $D_{ii}=d_G(v_i)$.
Ważone elementy diagonalne są sumami wag krawędzi
przylegających do wierzchołka, a pozadiagonalny element
$L_{uv}$ wynosi $-w_{uv}$ dla krawędzi $uv$ i zero w pozostałych
przypadkach. Odwrócenie pomocniczej orientacji zmienia znak
kolumny $B$, ale nie zmienia iloczynu $BWB^T$.

\begin{twierdzenie}[Tożsamość energii]
\label{thm:tg-energy}
Dla dowolnego $f\in\mathbb R^V$ zachodzi
\[
f^TLf=\sum_{\{u,v\}\in E}w_{uv}(f(u)-f(v))^2\ge0.
\]
Stąd $L$ jest macierzą symetryczną dodatnio półokreśloną, a
$\ker L$ składa się dokładnie z funkcji stałych na każdej
składowej spójnej. W szczególności $\dim\ker L=c$ i
$\operatorname{rank}L=n-c$.
\end{twierdzenie}

\begin{proof}
Ponieważ $L=BWB^T$, dostajemy
$f^TLf=(B^Tf)^TW(B^Tf)$. Dla krawędzi skierowanej
pomocniczo $u\to v$ jej współrzędna wektora $B^Tf$
wynosi $f(v)-f(u)$. Po rozwinięciu iloczynu otrzymujemy
wskazaną sumę kwadratów. Jest ona nieujemna, gdyż wagi są
dodatnie. Ponieważ dla macierzy symetrycznej dodatnio
półokreślonej równość $f^TLf=0$ jest równoważna
$Lf=0$, mamy $Lf=0$ wtedy i tylko wtedy, gdy na
końcach każdej krawędzi $f$ przyjmuje równe wartości.
To oznacza dokładnie stałość na składowych spójnych.
Ich wskaźniki stanowią bazę jądra, więc jego wymiar wynosi
$c$, a rząd wynosi $n-c$.
\end{proof}

\begin{przyklad}[Laplasjan ścieżki]
Dla ścieżki $1-2-3$ bez dodatkowych wag otrzymujemy
\[
L=\begin{pmatrix}1&-1&0\\-1&2&-1\\0&-1&1\end{pmatrix},
\quad
f^TLf=(f_1-f_2)^2+(f_2-f_3)^2.
\]
Wektor stały $(1,1,1)^T$ należy do jądra.
Wektory $(1,0,-1)^T$ i $(1,-2,1)^T$ są wektorami własnymi
odpowiadającymi wartościom odpowiednio $1$ i $3$.
Wartości własne laplasjanu wynoszą więc $0,1,3$.
\end{przyklad}

\section{Harmoniczność, zasada maksimum i zagadnienie Dirichleta}
\label{sec:tg-harmonic}

\begin{definicja}[Funkcja harmoniczna we wierzchołku]
Dla ważonego grafu $G$ i funkcji $f:V\to\mathbb R$
mówimy, że $f$ jest \emph{harmoniczna w wierzchołku $v$},
jeżeli $(Lf)(v)=0$. Oznaczmy przez $d_w(v):=\sum_{u:\,uv\in E}w_{uv}$
jego ważony stopień.
\end{definicja}

Jeżeli $v$ jest izolowany, to $(Lf)(v)=0$ dla każdej funkcji,
ale wzór średniej ważonej nie jest wtedy zdefiniowany.
Dalej stosujemy go wyłącznie dla $d_w(v)>0$.

Warunek harmoniczności w nieizolowanym $v$ jest równoważny równości
\[
f(v)=\frac{1}{d_w(v)}\sum_{u:\,uv\in E}w_{uv}f(u),
\]
to znaczy wartość w wierzchołku jest średnią ważoną wartości
na jego sąsiadach. Jest to dyskretny odpowiednik
własności wartości średniej funkcji harmonicznych.

\begin{twierdzenie}[Zasada maksimum na grafie]
Niech $G$ będzie skończonym spójnym grafem ważonym, którego
wszystkie wagi są dodatnie. Niech $\varnothing\ne S\subseteq V$
będzie zbiorem wierzchołków brzegowych. Jeżeli funkcja
$f:V\to\mathbb R$ jest harmoniczna w każdym wierzchołku
z $V\setminus S$, to osiąga maksimum i minimum na $S$.
\end{twierdzenie}

\begin{proof}
Niech $M=\max_{v\in V}f(v)$ i niech $U$ będzie zbiorem
wierzchołków, w których funkcja przyjmuje wartość $M$.
Jeżeli $U\cap S\ne\varnothing$, maksimum jest osiągnięte
na brzegu. W przeciwnym razie każdy $v\in U$ jest
wewnętrzny i harmoniczny. Równość średniej ważonej
$f(v)=\sum_{u\sim v}w_{uv}f(u)/d_w(v)=M$ przy dodatnich wagach
implikuje, że każdy sąsiad $u$ wierzchołka $v$ też ma
wartość $M$. Zatem $U$ jest zamknięty na sąsiadów.
Spójność grafu prowadzi wtedy do $U=V$, co jest
sprzeczne z $U\cap S=\varnothing$.
Tak więc maksimum jest osiągnięte na $S$.
Stosując ten sam argument do $-f$, otrzymujemy
analogiczne stwierdzenie dla minimum.
\end{proof}

\begin{twierdzenie}[Dyskretne zagadnienie Dirichleta]
\label{thm:tg-dirichlet}
Niech $G=(V,E)$ będzie skończonym grafem z dodatnimi
wagami, a $S\subseteq V$ niepustym zbiorem przecinającym
każdą składową spójną. Dla dowolnej funkcji brzegowej
$g:S\to\mathbb R$ istnieje dokładnie jedna funkcja
$f:V\to\mathbb R$ taka, że $f|_S=g$ oraz
$(Lf)(v)=0$ dla każdego $v\in V\setminus S$.
Ponadto $f$ jest jednoznacznym minimizatorem energii
$E(h):=\tfrac12 h^TLh$ wśród wszystkich $h$ o $h|_S=g$.
\end{twierdzenie}

\begin{proof}
Niech $I=V\setminus S$. Jeśli $I=\varnothing$, funkcja $f=g$
jest jedynym rozwiązaniem i twierdzenie jest oczywiste.
W przeciwnym razie rozpiszmy macierz $L$ według podziału
$V=I\sqcup S$ na bloki $L_{II},L_{IS},L_{SI},L_{SS}$.
Równanie $Lf=0$ na $I$ przy zadanych wartościach
$f_S=g$ ma postać $L_{II}f_I=-L_{IS}g$.

Pokażemy, że $L_{II}$ jest dodatnio określona. Niech
$x\in\mathbb R^I$, $x\ne0$, a $\widetilde x\in\mathbb R^V$
będzie jej przedłużeniem przez zero na $S$.
Wówczas $x^TL_{II}x=\widetilde x^TL\widetilde x\ge0$.
Równość oznaczałaby, na mocy twierdzenia
\ref{thm:tg-energy}, że $\widetilde x$ jest stała
na każdej składowej. Ponieważ każda składowa spotyka $S$,
a tam $\widetilde x=0$, wówczas $\widetilde x=0$,
sprzecznie z $x\ne0$. A więc $L_{II}$ jest dodatnio
określona i odwracalna; równanie ma jednoznaczne rozwiązanie.

Każda funkcja $h$ o tych samych wartościach na $S$
ma postać $h=f+\eta$ z $\eta|_S=0$.
Rozwinięcie energii daje
$E(f+\eta)=E(f)+\eta^TLf+\tfrac12\eta^TL\eta$.
Wyraz mieszany jest zerowy, ponieważ $Lf$ znika na
$I$, a $\eta$ na $S$. Pozostały składnik jest dodatni
dla $\eta\ne0$ na mocy dodatniej określoności $L_{II}$.
Otrzymujemy jednoznaczność minimum.
\end{proof}

\begin{przyklad}[Napięcie na łańcuchu trzech wierzchołków]
Na ścieżce $a-b-c$ ustalamy $f(a)=0$, $f(c)=1$
i przyjmujemy wagi jednostkowe. Harmoniczność w $b$ daje
$2f(b)-f(a)-f(c)=0$, więc $f(b)=1/2$.
Energia wynosi $E(f)=\tfrac12((1/2)^2+(1/2)^2)=1/4$.
Każda inna wartość $f(b)=t$ daje
$E(t)=\tfrac12(t^2+(1-t)^2)=t^2-t+1/2$,
które przyjmuje minimum dokładnie dla $t=1/2$.
\end{przyklad}

\begin{center}
\begin{tikzpicture}[scale=1.05,line cap=round]
 \node[tgBlueVertex] (a) at (-2.1,0) {$0$};
 \node[tgTealVertex] (b) at (0,0) {$\frac12$};
 \node[tgRedVertex] (c) at (2.1,0) {$1$};
 \draw[tgTree,line width=3pt] (a)--(b);
 \draw[tgAccent,line width=3pt] (b)--(c);
 \node[text=tgNavy,font=\footnotesize] at (0,-.67)
   {Wartość harmoniczna w środku to średnia wartości na końcach.};
\end{tikzpicture}
\end{center}

\section{Dyskretne operatory Hodge'a}
\label{sec:tg-hodge-operators}

Dla przestrzeni euklidesowych $C^0(G):=\mathbb R^V$
i $C^1(G):=\mathbb R^E$ wprowadźmy operator
różniczki dyskretnej $d_0:=B^T:C^0\to C^1$.
Dla krawędzi $e:u\to v$ otrzymujemy
$(d_0f)(e)=f(v)-f(u)$.
Operator sprzężony wynosi $d_0^*=B:C^1\to C^0$.

\begin{definicja}[Laplasjany Hodge'a na grafie]
Dla skończonego grafu prostego z jednostkowymi wagami
definiujemy operator na $0$-kołańcuchach
$\Delta_0:=d_0^*d_0=BB^T$ i operator na
$1$-kołańcuchach $\Delta_1:=d_0d_0^*=B^TB$.
$k$-kołańcuch $\omega$ nazywamy harmonicznym,
gdy $\Delta_k\omega=0$.
\end{definicja}

\begin{twierdzenie}[Harmoniczne reprezentanty kohomologii]
\label{thm:tg-hodge-harmonic}
Niech $G$ będzie skończonym grafem prostym o $n$ wierzchołkach,
$m$ krawędziach i $c$ składowych. Wtedy
\[
\ker\Delta_0=\ker B^T\cong\mathbb R^c,
\qquad
\ker\Delta_1=\ker B\cong\mathbb R^{m-n+c}.
\]
Ponadto każda klasa w
$H^1(G;\mathbb R):=C^1(G)/\operatorname{im}d_0$
ma dokładnie jednego przedstawiciela należącego do
$\ker\Delta_1$. Ponieważ graf jest kompleksem
jednowymiarowym, operator $d_1:C^1\to C^2=0$ jest zerowy,
a więc $H^1=\ker d_1/\operatorname{im}d_0$.
\end{twierdzenie}

\begin{proof}
Ponieważ $\langle\Delta_0f,f\rangle=\|B^Tf\|^2$,
warunek $\Delta_0f=0$ jest równoważny $B^Tf=0$.
Oznacza to stałość $f$ na każdej składowej,
a więc wymiar jądra wynosi $c$.
Analogicznie $\langle\Delta_1h,h\rangle=\|Bh\|^2$,
zatem $\ker\Delta_1=\ker B$; jego wymiar
wynosi $m-n+c$ z twierdzenia \ref{thm:tg-rank}.

Z rozkładu ortogonalnego
$C^1=\operatorname{im}B^T\oplus^\perp\ker B$
każdy $x\in C^1$ rozkłada się jednoznacznie na
$x=d_0f+h$, gdzie $h\in\ker B$. W klasie $[x]$
modulo obrazu $d_0$ przedstawiciel $h$ istnieje,
a gdyby były dwa, ich różnica należałaby zarówno
do $\operatorname{im}B^T$, jak i do jego dopełnienia
ortogonalnego. Musiałaby więc być zerowa.
\end{proof}

\begin{uwaga}[Porównanie z rozmaitościami]
W teorii Hodge'a na rozmaitości zewnętrzna różniczka
ma postać $d:\Omega^k\to\Omega^{k+1}$, a laplasjan
Hodge'a to $\Delta=dd^*+d^*d$. Na grafie mamy tylko
wymiary zero i jeden, więc $d_1=0$: stąd na
$C^0$ pozostaje $d_0^*d_0$, a na $C^1$ -- $d_0d_0^*$.
To szczególny, skończeniewymiarowy model twierdzenia,
że klasy kohomologii mają jednoznacznych harmonicznych
reprezentantów.
\end{uwaga}

\begin{stwierdzenie}[Dodatnie widma dwóch laplasjanów]
Niech $B$ będzie rzeczywistą macierzą incydencji grafu.
Dodatnie wartości własne macierzy $BB^T$ i $B^TB$
są takie same, z tymi samymi krotnościami.
\end{stwierdzenie}

\begin{proof}
Jeśli $BB^Tf=\lambda f$ z $\lambda>0$, to
$B^Tf\ne0$ (gdyby $B^Tf=0$, lewa strona byłaby zerowa),
a ponadto
$B^TB(B^Tf)=B^T(BB^Tf)=\lambda B^Tf$.
W drugą stronę, dla $B^TBh=\lambda h$ wektor
$Bh$ jest niezerowy i spełnia
$BB^T(Bh)=\lambda Bh$.
Przekształcenia $B^T$ i $B$ tworzą wzajemnie odwrotne
izomorfizmy przestrzeni własnych po przeskalowaniu
przez $\lambda^{-1}$. Stąd zgadzają się krotności.
\end{proof}

\section{Twierdzenie macierzowo-drzewowe Kirchhoffa}
\label{sec:tg-matrix-tree}

\begin{definicja}[Liczba drzew rozpinających]
Dla skończonego grafu prostego $G$ przez $\tau(G)$
oznaczamy liczbę jego drzew rozpinających.
Jeśli $G$ jest niespójny, przyjmujemy $\tau(G)=0$.
Dla grafu jednoelementowego $K_1$ przyjmujemy
$\tau(K_1)=1$, ponieważ jego jedynym drzewem
rozpinającym jest on sam.
\end{definicja}

\begin{twierdzenie}[Kirchhoff: liczba drzew z minora laplasjanu]
\label{thm:tg-kirchhoff}
Niech $G=(V,E)$ będzie skończonym spójnym grafem prostym,
$|V|=n\ge1$, z dodatnimi wagami $w_e$ na krawędziach.
Niech $L=BWB^T$, a $L^{(r)}$ macierzą otrzymaną
przez skreślenie wiersza i kolumny odpowiadających
wierzchołkowi $r\in V$. Wówczas
\[
\det L^{(r)}=\sum_{T}\prod_{e\in E(T)}w_e,
\]
gdzie suma biegnie po wszystkich drzewach rozpinających $T$.
Dla wag jednostkowych otrzymujemy
$\det L^{(r)}=\tau(G)$, niezależnie od wyboru $r$.
Dla $n=1$ wyznacznik pustej macierzy przyjmujemy równy $1$.
\end{twierdzenie}

\begin{proof}
Przypadek $n=1$ jest natychmiastowy. Niech $n\ge2$.
Usunięcie $r$-tego wiersza z macierzy $B$ daje
$(n-1)\times m$ macierz $B^{(r)}$. Mamy
$L^{(r)}=B^{(r)}W(B^{(r)})^T$.
Z wzoru Cauchy'ego--Bineta na wyznacznik iloczynu
macierzy prostokątnych otrzymujemy
\[
\det L^{(r)}=\sum_{\substack{S\subseteq E\\|S|=n-1}}
  \det(B^{(r)}_S)^2\prod_{e\in S}w_e,
\]
gdzie $B^{(r)}_S$ oznacza kwadratową podmacierz
złożoną z kolumn oznaczonych przez $S$.

Udowodnimy, że $\det(B^{(r)}_S)=0$, jeżeli $(V,S)$
nie jest drzewem, a $\det(B^{(r)}_S)=\pm1$, gdy jest
drzewem. Jeśli $(V,S)$ zawiera cykl, kolumny
incydencji tego cyklu są liniowo zależne: po wybraniu
znaków zgodnych z jego przebiegiem ich suma wynosi zero.
Zatem wyznacznik znika. Jeśli zaś $(V,S)$ nie zawiera
cyklu i ma $n-1$ krawędzi, jest lasem z
$n-(n-1)=1$ składową, czyli drzewem.

Niech więc $(V,S)$ będzie drzewem. W drzewie z co najmniej
dwoma wierzchołkami istnieje liść $v\ne r$.
W macierzy $B^{(r)}_S$ wiersz odpowiadający $v$
ma dokładnie jeden niezerowy element, równy $1$ albo $-1$.
Rozwinięcie wyznacznika względem tego wiersza sprowadza
go, z dokładnością do znaku, do analogicznego wyznacznika
dla drzewa z usuniętym liściem $v$ i przyległą krawędzią.
Powtarzając procedurę, dochodzimy do pustej macierzy
o wyznaczniku $1$. Stąd $\det(B^{(r)}_S)=\pm1$.

Do sumy Cauchy'ego--Bineta wkład dają więc dokładnie
zbiory krawędzi tworzące drzewa rozpinające, a ich
wkład wynosi $\prod_{e\in E(T)}w_e$.
Po zsumowaniu otrzymujemy tezę.
Warto zauważyć, że argument nie wymaga w istocie dodatniości
wag: jest to tożsamość wielomianowa w zmiennych $(w_e)_{e\in E}$.
Założenie $w_e>0$ zapewnia dodatkowo dodatniość każdego
składnika sumy i nieosobliwość zredukowanego laplasjanu
spójnego grafu. Sam wzór wyznacznikowy obowiązuje również
dla wag zerowych i ujemnych.
\end{proof}

\begin{przyklad}[Cztery wierzchołki na cyklu]
Dla cyklu $C_4$ z krawędziami jednostkowymi mamy
\[
L=\begin{pmatrix}
2&-1&0&-1\\-1&2&-1&0\\0&-1&2&-1\\-1&0&-1&2
\end{pmatrix}.
\]
Skreślając czwarty wiersz i kolumnę, dostajemy
\[
\det\begin{pmatrix}2&-1&0\\-1&2&-1\\0&-1&2\end{pmatrix}=4.
\]
Z drugiej strony każde drzewo rozpinające $C_4$
powstaje przez usunięcie dokładnie jednej z czterech
krawędzi. Takich drzew jest cztery, więc oba rachunki
się zgadzają.
\end{przyklad}

\begin{przyklad}[Ważony trójkąt]
Niech wagi trzech krawędzi trójkąta o wierzchołkach
$1,2,3$ wynoszą $w_{12}=2$, $w_{23}=w_{31}=1$.
Wtedy
\[
L=\begin{pmatrix}3&-2&-1\\-2&3&-1\\-1&-1&2\end{pmatrix},
\qquad
\det L^{(3)}=\det\begin{pmatrix}3&-2\\-2&3\end{pmatrix}=5.
\]
Trzy drzewa rozpinające składają się z dwóch spośród
trzech krawędzi, a ich iloczyny wag wynoszą $2$, $2$
i $1$. Suma $2+2+1=5$ zgadza się z wynikiem macierzowym.
\end{przyklad}

\begin{center}
\begin{tikzpicture}[scale=.95,line cap=round]
  \node[tgBlueVertex] (a) at (-1.7,-.95) {$1$};
  \node[tgBlueVertex] (b) at (1.7,-.95) {$2$};
  \node[tgBlueVertex] (c) at (0,1.40) {$3$};
  \draw[tgAccent,line width=3.0pt] (a)--node[below,font=\small] {$2$}(b);
  \draw[tgTree,line width=2.6pt] (b)--node[above right,font=\small] {$1$}(c);
  \draw[tgBlue,line width=2.6pt] (c)--node[above left,font=\small] {$1$}(a);
  \node[text=tgNavy,font=\footnotesize] at (0,-1.65)
    {Waga drzewa = iloczyn wag jego krawędzi.};
\end{tikzpicture}
\end{center}

\section{Wartości własne laplasjanu i liczba drzew}
\label{sec:tg-spectrum-tree}

\begin{twierdzenie}[Spektralny wzór Kirchhoffa]
\label{thm:tg-spectral-tree}
Niech $G$ będzie spójnym grafem prostym o $n\ge2$
wierzchołkach, z jednostkowymi wagami, a
$0=\lambda_1<\lambda_2\le\cdots\le\lambda_n$
niech będą wartościami własnymi jego laplasjanu
liczonymi z krotnościami. Wtedy
\[
\tau(G)=\frac1n\prod_{j=2}^n\lambda_j.
\]
\end{twierdzenie}

\begin{proof}
Z twierdzenia o energii i spójności jądro $L$ jest
jednowymiarowe, rozpięte przez $\mathbf1=(1,\ldots,1)^T$.
Z twierdzenia Kirchhoffa każdy główny kofaktor $L$
wynosi $\tau(G)$. Co więcej, macierz dopełnień
algebraicznych $\operatorname{adj}(L)$ ma wszystkie
kolumny w $\ker L$, ponieważ
$L\operatorname{adj}(L)=(\det L)I=0$.
Symetria $L$ i zgodność diagonalnych elementów
$\operatorname{adj}(L)$ z $\tau(G)$ dają
$\operatorname{adj}(L)=\tau(G)\mathbf1\mathbf1^T$.

Rozważmy macierz
$M=L+\tfrac1n\mathbf1\mathbf1^T$.
Na prostej rozpiętej przez $\mathbf1$ działa ona jako
identyczność, a na dopełnieniu ortogonalnym jak $L$.
Zatem jej wartości własne to $1,\lambda_2,\ldots,\lambda_n$
i $\det M=\prod_{j=2}^n\lambda_j$.
Z drugiej strony wzór na wyznacznik po dodaniu
macierzy rzędu jeden daje
\[
\det M=\det L+\frac1n\mathbf1^T\operatorname{adj}(L)\mathbf1
=\frac1n\tau(G)n^2=n\tau(G).
\]
Porównanie obydwu wyrażeń kończy dowód.
\end{proof}

\begin{wniosek}[Wzór Cayleya]
Dla grafu pełnego $K_n$, $n\ge2$, liczba drzew
rozpinających wynosi $\tau(K_n)=n^{n-2}$.
\end{wniosek}

\begin{proof}
Macierz sąsiedztwa $K_n$ ma wszystkie elementy
pozadiagonalne równe $1$, natomiast stopnie wynoszą
$n-1$. Stąd $L=nI-J$, gdzie $J=\mathbf1\mathbf1^T$.
Na wektorze $\mathbf1$ operator znika, a na
$\mathbf1^\perp$ działa jako mnożenie przez $n$.
Zatem wartości własne to $0$ oraz $n$ o krotności $n-1$.
Ze wzoru spektralnego otrzymujemy
$\tau(K_n)=\frac1n n^{n-1}=n^{n-2}$.
\end{proof}

\begin{przyklad}[Widmo cyklu i grafu z przekątną]
Dla $C_4$ wartości własne laplasjanu wynoszą
$0,2,2,4$. Otrzymujemy $\tau(C_4)=\frac14(2\cdot2\cdot4)=4$.
Dodajmy jedną przekątną do kwadratu. Otrzymamy graf
$K_4$ bez jednej krawędzi, którego widmo to $0,2,4,4$.
Stąd $\tau=\frac14(2\cdot4\cdot4)=8$.
Dodanie pojedynczej krawędzi podwoiło więc liczbę
rozpinających drzew w tym przykładzie.
\end{przyklad}

\section{Efektywna rezystancja i zasada minimalnej energii}
\label{sec:tg-resistance}

Laplasjan może również opisywać sieć elektryczną: wagę $w_e>0$
interpretujemy jako przewodność krawędzi, a jej odwrotność
$r_e=1/w_e$ jako opór. Ustalenie kierunków krawędzi jest wyłącznie
pomocnicze; to równania przepływu, a nie te kierunki, wyznaczają
prąd w sieci.

\begin{definicja}[Efektywna rezystancja między wierzchołkami]
Niech $G=(V,E)$ będzie skończonym spójnym grafem prostym,
$|V|\ge2$, z dodatnimi wagami $w_e$. Niech $s,t\in V$
będą różnymi wierzchołkami, a $b=[s]-[t]\in\mathbb R^V$.
Wybieramy dowolne rozwiązanie równania $Lf=b$,
gdzie $L=BWB^T$ jest ważonym laplasjanem.
Liczbę $R_{\mathrm{eff}}(s,t):=f(s)-f(t)$ nazywamy
\emph{efektywną rezystancją} między $s$ i $t$.
\end{definicja}

\begin{stwierdzenie}[Istnienie rozwiązania i dodatniość oporu]
W sytuacji z poprzedniej definicji równanie $Lf=b$ ma
rozwiązanie, a każde dwa rozwiązania różnią się funkcją stałą.
Liczba $R_{\mathrm{eff}}(s,t)$ nie zależy od wyboru rozwiązania
i jest ściśle dodatnia. Ponadto
\[
R_{\mathrm{eff}}(s,t)=f^TLf
   =\sum_{uv\in E}w_{uv}(f(u)-f(v))^2.
\]
\end{stwierdzenie}

\begin{proof}
Ponieważ graf jest spójny, $\ker L$ jest prostą
rozpiętą przez wektor stały $\mathbf1$. Jako że $L$
jest symetryczna, mamy
$\operatorname{im}L=(\ker L)^\perp
=\{x\in\mathbb R^V:\sum_vx_v=0\}$.
Wektor $b=[s]-[t]$ ma sumę współrzędnych równą zero,
więc należy do obrazu $L$; rozwiązanie $f$ istnieje.
Dwa rozwiązania różnią się elementem $\ker L$,
czyli stałą, której dodanie nie zmienia $f(s)-f(t)$.

Mnożąc równanie $Lf=b$ skalarnie przez $f$,
otrzymujemy $f^TLf=f^Tb=f(s)-f(t)$.
Z twierdzenia \ref{thm:tg-energy} jest to suma
nieujemnych składników wskazana w tezie.
Nie może wynosić zero: wtedy $f$ byłaby stała,
a więc $Lf=0$, wbrew $b\ne0$. W konsekwencji
$R_{\mathrm{eff}}(s,t)>0$.
\end{proof}

\begin{twierdzenie}[Zasada Thomsona: minimalna energia przepływu]
Niech $G,B,W,s,t,b$ będą jak wyżej. Przez
$\mathcal J:=\{j\in\mathbb R^E:Bj=b\}$ oznaczmy
zbiór wszystkich prądów krawędziowych realizujących
jednostkowy przepływ między $s$ i $t$.
Wtedy
\[
R_{\mathrm{eff}}(s,t)
=\min_{j\in\mathcal J}\sum_{e\in E}\frac{j_e^2}{w_e}.
\]
Minimum jest osiągane dokładnie przez prąd
$j_0=WB^Tf$, gdzie $f$ jest dowolnym rozwiązaniem $Lf=b$.
\end{twierdzenie}

\begin{proof}
Dla $j_0=WB^Tf$ mamy
$Bj_0=BWB^Tf=Lf=b$, więc $j_0\in\mathcal J$.
Każdy inny prąd $j\in\mathcal J$ możemy jednoznacznie
zapisać jako $j=j_0+h$ z $Bh=0$.
Rozwijając formę kwadratową energii,
otrzymujemy
\[
\begin{aligned}
 j^TW^{-1}j
 &=j_0^TW^{-1}j_0+2h^TW^{-1}j_0+h^TW^{-1}h\\
 &=f^TBWB^Tf+2h^TB^Tf+h^TW^{-1}h\\
 &=R_{\mathrm{eff}}(s,t)+h^TW^{-1}h.
\end{aligned}
\]
W drugim przejściu użyliśmy $j_0=WB^Tf$,
a w trzecim $Bh=0$ i poprzedniego stwierdzenia.
Ponieważ $W^{-1}$ jest dodatnio określona,
ostatni składnik jest nieujemny i znika wyłącznie
dla $h=0$. Zatem $j_0$ jest jedynym minimizatorem,
a wartość minimum wynosi $R_{\mathrm{eff}}(s,t)$.
\end{proof}

\begin{przyklad}[Dwa równoległe sposoby przejścia prądu]
Dla ścieżki $s-a-t$ o przewodnościach równych jeden
przyjmujemy $b=[s]-[t]$. Rozwiązaniem $Lf=b$
jest funkcja $f(s)=1$, $f(a)=0$, $f(t)=-1$.
Stąd $R_{\mathrm{eff}}(s,t)=1-(-1)=2$:
dwa opory jednostkowe połączone szeregowo sumują się.

Dla trójkąta $K_3$ z wierzchołkami $s,t,a$
i wagami jednostkowymi mamy $L=3I-J$.
Wektor $f=(1/3,-1/3,0)$ w kolejności $(s,t,a)$
spełnia $Lf=(1,-1,0)=b$.
Wynika stąd
$R_{\mathrm{eff}}(s,t)=2/3$.
Bezpośrednia krawędź $st$ stanowi drogę o oporze $1$,
a druga możliwość $s-a-t$ ma opór $2$;
połączenie równoległe daje $(1^{-1}+2^{-1})^{-1}=2/3$.
Dodanie krawędzi zmniejszyło opór efektywny.
\end{przyklad}

\begin{center}
\begin{tikzpicture}[scale=.97,line cap=round]
  \begin{scope}[shift={(-3.2,0)}]
    \node[tgBlueVertex] (s) at (-1.1,-.2) {$s$};
    \node[tgTealVertex] (a) at (0,1.05) {$a$};
    \node[tgRedVertex] (t) at (1.1,-.2) {$t$};
    \draw[tgTree,line width=2.6pt] (s)--(a)--(t);
    \node[font=\footnotesize,text=tgNavy] at (0,-.82) {$R_{st}=2$};
  \end{scope}
  \begin{scope}[shift={(3.2,0)}]
    \node[tgBlueVertex] (s) at (-1.1,-.2) {$s$};
    \node[tgTealVertex] (a) at (0,1.05) {$a$};
    \node[tgRedVertex] (t) at (1.1,-.2) {$t$};
    \draw[tgTree,line width=2.6pt] (s)--(a)--(t);
    \draw[tgAccent,line width=2.6pt] (s)--(t);
    \node[font=\footnotesize,text=tgNavy] at (0,-.82) {$R_{st}=2/3$};
  \end{scope}
\end{tikzpicture}

{\footnotesize Dodanie czerwonej krawędzi jednostkowej
zwiększa przewodność pomiędzy $s$ i $t$.}
\end{center}


\section{Luka spektralna i geometryczne znaczenie laplasjanu}
\label{sec:tg-spectral-gap}

\begin{definicja}[Druga wartość własna]
Niech $G$ będzie spójnym grafem prostym o $n\ge2$
wierzchołkach i laplasjanie nieważonym $L$.
Liczbę $\lambda_2(G)$, czyli drugą najmniejszą
wartość własną $L$, nazywamy \emph{spójnością algebraiczną}
(lub \emph{luką spektralną} laplasjanu).
\end{definicja}

\begin{twierdzenie}[Zasada Rayleigha]
Niech $G$ będzie jak wyżej. Wówczas
\[
\lambda_2(G)=\min_{\substack{x\in\mathbb R^V\setminus\{0\}\\
               \sum_{v\in V}x_v=0}}
\frac{\sum_{uv\in E}(x_u-x_v)^2}{\sum_{v\in V}x_v^2}>0.
\]
\end{twierdzenie}

\begin{proof}
Macierz $L$ jest rzeczywista i symetryczna, więc posiada
ortogonalną bazę wektorów własnych. Wektor
$\mathbf1/\sqrt n$ odpowiada wartości własnej zero.
Każdy $x\perp\mathbf1$ możemy zapisać w pozostałej
części tej bazy jako $x=\sum_{j=2}^na_ju_j$.
Wtedy $x^TLx=\sum_{j=2}^n\lambda_ja_j^2$,
a $\|x\|^2=\sum_{j=2}^na_j^2$.
Iloraz Rayleigha jest więc średnią ważoną liczb
$\lambda_2,\ldots,\lambda_n$ o nieujemnych wagach,
a minimum przyjmuje na wektorze własnym dla
$\lambda_2$. Związek z sumą po krawędziach wynika
z twierdzenia \ref{thm:tg-energy}.
W spójnym grafie jedynymi wektorami zerowej energii
są stałe. W przestrzeni prostopadłej do $\mathbf1$
nie ma niezerowego takiego wektora, stąd $\lambda_2>0$.
\end{proof}

Liczba $\lambda_2$ mierzy opór wobec zmian funkcji między
połączonymi wierzchołkami. W następnym etapie, po wprowadzeniu
grafów planarnych i osadzeń grafów w powierzchniach,
połączymy laplasjan z topologią komórek wyższych wymiarów,
kompleksami dualnymi i dalszymi przykładami dyskretnej geometrii.

% ============================================================
% Iteracja 4, rozdział 9
% ============================================================
\chapter{Grafy planarne, ściany i wzór Eulera}
\label{ch:tg-planarne}

Dotychczas krawędzie grafu były abstrakcyjnymi odcinkami. Od tej chwili
istotne będzie również to, \emph{jak} graf umieszczamy w powierzchni.
Na przykład narysowanie dwóch krawędzi przecinających się na kartce
nie oznacza, że ich punkt przecięcia jest wierzchołkiem grafu.
Właśnie to rozróżnienie prowadzi do planarności, ścian, dualności
i do związku z charakterystyką Eulera powierzchni.

\section{Rysunek grafu a zanurzenie topologiczne}
\label{sec:tg-osadzanie-planarne}

\begin{definicja}[Zanurzenie grafu w powierzchnię]
Niech $G$ będzie skończonym grafem, dopuszczającym również pętle
i krawędzie równoległe, a $S$ powierzchnią topologiczną.
\emph{Zanurzeniem topologicznym grafu w $S$} (ang. \emph{embedding})
nazywamy ciągłe odwzorowanie różnowartościowe $\iota:|G|\to S$.
Z powodu zwartości $|G|$ i własności Hausdorffa przestrzeni $S$
jest ono homeomorfizmem na swój obraz. Graf nazywamy \emph{planarnym},
jeżeli istnieje takie $\iota$ dla $S=\mathbb R^2$.
W przypadku $S=\mathbb R^2$ konkretny graf wraz z wybranym
zanurzeniem nazywamy \emph{grafem płaskim}. Wnętrza obrazów różnych
krawędzi są rozłączne i nie zawierają obrazów wierzchołków;
dwie krawędzie mogą mieć wspólny punkt wyłącznie wtedy, gdy
mają wspólny koniec.
\end{definicja}

\begin{uwaga}[Płaszczyzna i sfera]
Graf skończony jest planarny wtedy i tylko wtedy, gdy można go
zanurzyć w sferze $S^2$. W jedną stronę stosujemy dołączenie
punktu w nieskończoności $\mathbb R^2\subset S^2$.
W drugą wybieramy punkt $p\in S^2\setminus\iota(|G|)$
(taki punkt istnieje, ponieważ skończona suma zanurzonych łuków
i punktów ma puste wnętrze w powierzchni)
i identyfikujemy $S^2\setminus\{p\}$ z $\mathbb R^2$.
Na sferze żadna ściana nie jest wyróżniona jako zewnętrzna.
\end{uwaga}

\begin{definicja}[Ściana]
Niech $\iota:|G|\hookrightarrow S^2$ będzie zanurzeniem grafu
skończonego. \emph{Ścianami} są składowe spójne dopełnienia
$S^2\setminus\iota(|G|)$. Ich liczbę oznaczamy przez $f(G,\iota)$.
Dla grafu płaskiego jedna ze ścian jest nieograniczona.
\end{definicja}

\begin{przyklad}[Ten sam graf, różne rysunki]
Graf $K_4$ można narysować z przecięciami na rysunku czworokąta
z obiema przekątnymi. Nie dowodzi to nieplanarności: wystarczy
umieścić trzy wierzchołki na brzegu trójkąta, a czwarty w jego wnętrzu.
Otrzymujemy wtedy cztery ściany na sferze, wszystkie ograniczone
trójkątami.
\end{przyklad}

\begin{center}
\begin{tikzpicture}[line cap=round,line join=round,scale=.95]
  \begin{scope}[shift={(-3.3,0)}]
    \coordinate (a) at (-1.15,-.9); \coordinate (b) at (1.15,-.9);
    \coordinate (c) at (0,1.30); \coordinate (d) at (0,-.15);
    \fill[tgBlue!8] (a)--(b)--(d)--cycle;
    \fill[tgTeal!12] (a)--(c)--(d)--cycle;
    \fill[tgGold!17] (b)--(c)--(d)--cycle;
    \draw[tgEdge] (a)--(b)--(c)--cycle (a)--(d)--(b) (c)--(d);
    \foreach \pt/\name in {a/a,b/b,c/c,d/d}
      \node[tgBlueVertex] at (\pt) {$\name$};
    \node[text=tgNavy,font=\small] at (0,-1.58) {$K_4$ bez przecięć};
  \end{scope}
  \begin{scope}[shift={(3.3,0)}]
    \coordinate (a) at (-1.10,-.95);\coordinate (b) at (1.10,-.95);
    \coordinate (c) at (1.10,.95);\coordinate (d) at (-1.10,.95);
    \draw[tgEdge] (a)--(b)--(c)--(d)--cycle;
    \draw[tgAccent] (a)--(c) (b)--(d);
    \foreach \pt/\name in {a/a,b/b,c/c,d/d}
      \node[tgRedVertex] at (\pt) {$\name$};
    \node[text=tgNavy,font=\small] at (0,-1.58) {inny rysunek $K_4$};
  \end{scope}
\end{tikzpicture}

{\footnotesize Czerwone przekątne po prawej przecinają się, lecz ten punkt
nie jest dodatkowym wierzchołkiem. Rysunek po lewej dowodzi planarności $K_4$.}
\end{center}

\section{Wzór Eulera: pełny dowód kombinatoryczny}
\label{sec:tg-planarny-euler}

\begin{twierdzenie}[Euler dla grafów zanurzonych w sferze]
\label{thm:tg-euler-plane}
Niech $G$ będzie skończonym grafem o $n\ge1$ wierzchołkach,
$m$ krawędziach i $c$ składowych spójnych, zanurzonym w $S^2$.
Jeżeli $f$ oznacza liczbę ścian tego zanurzenia, to
\[
 n-m+f=1+c.
\]
W szczególności dla grafu spójnego zachodzi $n-m+f=2$.
Wynik jest prawdziwy także dla pętli i krawędzi równoległych.
\end{twierdzenie}

\begin{proof}
Ustalamy jedno zanurzenie i dodajemy jego krawędzie po kolei,
nie zmieniając ich położenia. Zaczynamy od $n$ izolowanych
wierzchołków. Mamy wtedy $m_0=0$, $c_0=n$, a dopełnienie
skończonego zbioru punktów w $S^2$ jest spójne, więc $f_0=1$.
Stąd $n-m_0+f_0=n+1=1+c_0$.

Przypuśćmy, że do rozważonego podgrafu dodajemy krawędź $e$.
Jej wnętrze leży w jednej ze ścian dotychczasowego rysunku.
Rozpatrzmy dwa przypadki.

Jeżeli końce $e$ należą do różnych składowych aktualnego grafu,
to $e$ scala te składowe: liczba $c$ maleje o jeden,
zaś liczba ścian się nie zmienia. Ostatni fakt można
uzasadnić, przesuwając mały łuk łączący dwa rozłączne
podzbiory grafu: taki łuk nie zamyka krzywej Jordana
ani nie rozcina na dwie składowe żadnej ściany.

Jeżeli końce $e$ są już połączone ścieżką w aktualnym
grafie (dopuszczamy też pętlę), to suma tej ścieżki
i nowej krawędzi zawiera krzywą prostą zamkniętą
w sensie topologicznym. Wybieramy prostą ścieżkę $P$ w dotychczasowym grafie
łączącą końce $e$; jeśli $e$ jest pętlą, bierzemy ścieżkę
stałą w jej wierzchołku. Obraz $P$ połączony z obrazem
$e$ tworzy krzywą Jordana. Jej dopełnienie ma dwie składowe,
a nowa krawędź leży w jednej dotychczasowej ścianie.
Obie strony $e$ należą po dołączeniu do różnych ścian,
a żadna inna ściana nie zostaje rozcięta. Zatem $f$ wzrasta
o jeden, podczas gdy $c$ pozostaje bez zmiany.

W pierwszym przypadku $m$ wzrasta o jeden, $c$ maleje
o jeden, $f$ się nie zmienia; w drugim $m$ i $f$
wzrastają jednocześnie o jeden. W obu sytuacjach
równość $n-m+f=1+c$ pozostaje prawdziwa. Indukcja
po liczbie krawędzi kończy dowód.
\end{proof}

\begin{uwaga}[Co oznacza brzeg ściany]
Brzeg ściany nie musi być cyklem prostym. W drzewie
cała sfera poza grafem jest jedną ścianą, a przejście
wokół jej brzegu używa każdej krawędzi dwa razy.
W ogólności śledzimy \emph{obchody ścianowe}: zamknięte
spacery krawędziowe, w których liczymy każdą stronę
krawędzi osobno. Suma ich długości, liczonych z krotnościami,
wynosi $2m$, także gdy występują mosty i pętle.
Dla grafów niespójnych pojedyncza ściana może mieć kilka
składowych brzegowych.
\end{uwaga}

\begin{przyklad}[Drzewo, cykl i kwadrat z przekątną]
Każde drzewo o $n$ wierzchołkach ma $m=n-1$ krawędzi.
Euler daje $f=1$. Cykl $C_n$ ma $n$ krawędzi i dwie ściany:
wnętrze i zewnętrze. Kwadrat z jedną przekątną ma
$n=4$, $m=5$, więc $f=2-4+5=3$.
Dwie wewnętrzne ściany są trójkątne, trzecia jest zewnętrzna.
\end{przyklad}

\begin{center}
\begin{tikzpicture}[scale=.95,line cap=round,line join=round]
  \coordinate (a) at (-2,0);\coordinate (b) at (2,0);
  \coordinate (c) at (2,2.0);\coordinate (d) at (-2,2.0);
  \fill[tgBlue!11] (a)--(b)--(c)--cycle;
  \fill[tgTeal!11] (a)--(c)--(d)--cycle;
  \draw[tgEdge] (a)--(b)--(c)--(d)--cycle;
  \draw[tgAccent] (a)--(c);
  \foreach \pt in {a,b,c,d} \node[tgBlueVertex] at (\pt) {};
  \node[text=tgBlue] at (.9,.55) {$F_1$};
  \node[text=tgTeal] at (-.9,1.45) {$F_2$};
  \node[text=tgNavy] at (0,-.65) {$F_3$ -- ściana zewnętrzna};
\end{tikzpicture}
\end{center}

\section{Oszacowania liczby krawędzi i nieplanarność}
\label{sec:tg-planarne-granice}

\begin{lemat}[Stopnie ścianowe]
Niech $G$ będzie spójnym, prostym grafem płaskim,
$n\ge3$, o $m$ krawędziach i $f$ ścianach.
Długość obchodu brzegowego każdej ściany, liczona
z krotnościami skierowanych wystąpień krawędzi, jest
co najmniej trzy, a suma długości wszystkich obchodów
to $2m$. Jeśli $G$ jest ponadto dwudzielny, każdy
obchód jest parzystej długości, a więc ma długość co najmniej cztery.
\end{lemat}

\begin{proof}
Każda krawędź ma dwie strony, które należą do ścian
(niekoniecznie różnych). Idąc po obwodzie każdej ściany,
przechodzimy każdą stronę dokładnie raz, skąd suma długości
wynosi $2m$. Obchód długości jeden wymagałby pętli, co jest
sprzeczne z prostotą. Obchód długości dwa mógłby powstać
z dwóch krawędzi równoległych albo z pojedynczej krawędzi
przechodzonej dwukrotnie. Jeśli zaś obchód składa się z dwukrotnego przejścia
jednej krawędzi, krawędź ta jest mostem, a obydwa jej końce
mają stopień jeden w odpowiednim brzegu ściany.
W przypadku grafu spójnego oznaczałoby to, że graf jest
pojedynczą krawędzią $K_2$; wyklucza to $n\ge3$.
Zatem każda długość jest co najmniej trzy.

W grafie dwudzielnym dowolny zamknięty spacer przechodzi
naprzemiennie pomiędzy obiema klasami wierzchołków,
ma więc długość parzystą. Wobec poprzedniego oszacowania
długość każdej ściany jest nie mniejsza niż cztery.
\end{proof}

\begin{twierdzenie}[Ograniczenia planarne]
\label{thm:tg-planarne-bound}
Jeżeli $G$ jest spójnym grafem prostym, planarnym,
o $n\ge3$ wierzchołkach i $m$ krawędziach, to
$m\le3n-6$. Jeżeli ponadto $G$ jest dwudzielny,
to $m\le2n-4$.
\end{twierdzenie}

\begin{proof}
Ustalmy zanurzenie płaskie. Poprzedni lemat daje
$3f\le2m$. Ze wzoru Eulera $f=2-n+m$,
więc $3(2-n+m)\le2m$, co po uporządkowaniu
daje $m\le3n-6$. W przypadku grafu dwudzielnego
mamy mocniejsze $4f\le2m$, a stąd
$4(2-n+m)\le2m$ oraz $m\le2n-4$.
\end{proof}

\begin{przyklad}[Dwie klasyczne przeszkody]
Graf $K_5$ ma pięć wierzchołków i dziesięć krawędzi,
lecz $10>3\cdot5-6=9$, więc nie jest planarny.
Graf $K_{3,3}$ ma sześć wierzchołków i dziewięć krawędzi,
jest dwudzielny, a $9>2\cdot6-4=8$, więc również
nie jest planarny. Samo ograniczenie $m\le3n-6$
nie wystarczyłoby do wykrycia nieplanarności $K_{3,3}$.
\end{przyklad}

\begin{center}
\begin{tikzpicture}[scale=.82,line cap=round,>=Stealth]
 \begin{scope}[shift={(-3.3,0)}]
  \foreach \i in {0,...,4} {
    \coordinate (p\i) at ({90-72*\i}:1.2);
  }
  \foreach \i in {0,...,4} {
    \foreach \j in {0,...,4} {
      \ifnum\i<\j \draw[tgEdge] (p\i)--(p\j);\fi
    }
  }
  \foreach \i in {0,...,4} \node[tgBlueVertex] at (p\i) {};
  \node[text=tgNavy,font=\small] at (0,-1.80) {$K_5$: $10>9$};
 \end{scope}
 \begin{scope}[shift={(3.3,0)}]
  \foreach \i in {1,2,3} {
    \coordinate (u\i) at (-1.25,{2-\i});
    \coordinate (v\i) at (1.25,{2-\i});
  }
  \foreach \i in {1,2,3} \foreach \j in {1,2,3}
      \draw[tgEdge] (u\i)--(v\j);
  \foreach \i in {1,2,3} {
    \node[tgBlueVertex] at (u\i) {};
    \node[tgRedVertex] at (v\i) {};
  }
  \node[text=tgNavy,font=\small] at (0,-1.80) {$K_{3,3}$: $9>8$};
 \end{scope}
\end{tikzpicture}

{\footnotesize Rysunki mają przecięcia niebędące wierzchołkami.
Nieplanarność tych grafów wynika z dowodów liczbowych, nie z wyglądu rysunków.}
\end{center}

\section{Kolorowanie ścian i graf dwudzielny w planie}
\label{sec:tg-planarne-color}

\begin{stwierdzenie}[Dwukolorowanie ścian]
Niech $G\subset S^2$ będzie spójnym grafem płaskim,
przy czym każdy wierzchołek ma parzysty stopień.
Wtedy ściany można pokolorować dwoma kolorami w taki sposób,
że ściany po przeciwnych stronach każdej krawędzi mają różne kolory.
W szczególności żadna krawędź takiego grafu nie jest mostem.
\end{stwierdzenie}

\begin{proof}
Każdy spójny graf, w którym wszystkie stopnie są parzyste,
ma zamknięty szlak Eulera. Żadna krawędź nie jest mostem,
ponieważ usunięcie mostu rozdzielałoby graf na dwie części,
każda mająca dokładnie jeden wierzchołek o nieparzystym stopniu
po usunięciu tej krawędzi, wbrew lematowi o uściskach dłoni.

Przypiszmy ścianie zewnętrznej kolor zero.
Dla dowolnej innej ściany wybieramy drogę w sferze
od ustalonego punktu ściany zewnętrznej do jej wnętrza,
przecinającą graf poprzecznie w skończenie wielu punktach
wnętrz krawędzi. Kolor określamy jako parzystość liczby
przecięć. Jest ona niezależna od drogi: różnica dwóch
takich dróg jest zamkniętą krzywą w sferze, a liczba
jej przecięć z cyklem eulerowskim jest parzysta.
Istotnie, mod $2$ cykl Eulera jest sumą cykli prostych,
a dowolna zamknięta krzywa w $S^2$ przecina każdy taki
cykl parzyście (twierdzenie Jordana i ogólne położenie).
Przekroczenie pojedynczej krawędzi zmienia parzystość,
a więc także kolor. To daje żądaną kolorystykę.
\end{proof}

\begin{uwaga}
Warunek parzystych stopni wierzchołków nie jest równoważny
dwudzielności grafu: cykl trójkątny ma wszystkie stopnie
parzyste, lecz nie jest dwudzielny. Stwierdzenie dotyczy
kolorowania \emph{ścian}, a nie wierzchołków.
\end{uwaga}


% ============================================================
% Iteracja 4, rozdział 10
% ============================================================
\chapter{Kryteria planarności, podziały krawędzi i kolorowanie}
\label{ch:tg-kuratowski}

\section{Podgrafy, podziały krawędzi i minory}
\label{sec:tg-rozne-redukcje}

\begin{definicja}[Podział krawędzi i podgraf topologiczny]
Niech $H$ będzie skończonym grafem. \emph{Podziałem krawędzi}
$uv\in E(H)$ nazywamy zastąpienie tej krawędzi ścieżką
$u-w-v$, gdzie $w$ jest nowym wierzchołkiem stopnia dwa.
Graf $H'$ jest \emph{subdywizją} $H$, jeżeli powstaje ze skończonego
ciągu takich podziałów. Mówimy, że $H$ jest \emph{podgrafem
topologicznym} grafu $G$, jeżeli pewna subdywizja $H$
jest izomorficzna z podgrafem $G$.
\end{definicja}

\begin{stwierdzenie}[Niezmienniczość planarności przy subdywizji]
Niech $H'$ będzie grafem otrzymanym z grafu $H$ przez podział
jednej krawędzi. Wtedy $H$ jest planarny wtedy i tylko wtedy,
gdy $H'$ jest planarny.
\end{stwierdzenie}

\begin{proof}
Jeśli znamy zanurzenie $H$, wybieramy dowolny punkt we wnętrzu
obrazu dzielonej krawędzi i traktujemy go jako nowy wierzchołek.
Tak powstaje zanurzenie $H'$.
W drugą stronę, w zanurzeniu $H'$ nowy wierzchołek ma stopień
dwa. Usuwamy jego oznaczenie i traktujemy dwie incydentne
krawędzie jako jeden łuk prosty topologicznie.
Nie pojawiają się nowe przecięcia, więc otrzymujemy zanurzenie $H$.
\end{proof}

\begin{definicja}[Minor grafu]
Niech $G$ i $H$ będą skończonymi grafami. Mówimy, że
$H$ jest \emph{minorem} $G$, jeśli można otrzymać $H$
z $G$ przez skończony ciąg operacji usuwania wierzchołków,
usuwania krawędzi i \emph{kontrakcji} krawędzi (utożsamienia
obu jej końców, po usunięciu tej krawędzi).
Po kontrakcji dopuszczamy pętle i krawędzie równoległe,
które można następnie usunąć.
\end{definicja}

\begin{stwierdzenie}[Planarność jest zamknięta na minory]
Jeżeli skończony graf $G$ jest planarny i $H$ jest jego minorem,
to $H$ jest planarny.
\end{stwierdzenie}

\begin{proof}
Usuwanie krawędzi i wierzchołków nie wprowadza przecięć.
Kontrakcję krawędzi realizujemy w jej małym regularnym
otoczeniu, przesuwając jeden koniec ku drugiemu wzdłuż
zanurzonego łuku. Wszystkie inne krawędzie incydentne
przesuwamy razem z ich końcami, zachowując kolejność
przy wierzchołkach. Po operacji otrzymujemy zanurzenie
nowego multigrafu. Stąd każda operacja minorowa zachowuje planarność.
\end{proof}

\begin{twierdzenie}[Kuratowski]
\label{thm:tg-kuratowski}
Niech $G$ będzie skończonym grafem prostym. Następujące
warunki są równoważne:
\begin{enumerate}[label=(\roman*)]
\item $G$ jest planarny;
\item $G$ nie zawiera podgrafu izomorficznego z żadną
subdywizją $K_5$ ani $K_{3,3}$.
\end{enumerate}
\end{twierdzenie}

\begin{proof}[Dowód implikacji \texorpdfstring{$(i)\Rightarrow(ii)$}{(i) => (ii)}]
Gdyby planarny $G$ zawierał podgraf będący subdywizją
$K_5$ lub $K_{3,3}$, sam ten podgraf byłby planarny
(jako podgraf grafu planarnego). Cofając podziały krawędzi,
uzyskalibyśmy zanurzenie $K_5$ albo $K_{3,3}$.
Sprzeciwia się to obliczeniu z poprzedniego rozdziału,
w którym wykazaliśmy ich nieplanarność.
\end{proof}

\begin{uwaga}[Status trudnej implikacji]
Odwrotna implikacja $(ii)\Rightarrow(i)$ jest zasadniczą,
znacznie trudniejszą częścią twierdzenia Kuratowskiego.
Wymaga analizy dwuspójnych i trójspójnych bloków grafu
oraz możliwości rozbudowania zanurzenia o kolejne ścieżki.
Nie wynika z samego wzoru Eulera ani z przedstawionego
powyżej dowodu. Przyjmujemy ją tutaj jako klasyczne
twierdzenie z teorii grafów planarnych; jej pełny dowód
można znaleźć w \emph{Graph Theory} R. Diestela oraz
w podręczniku J. A. Bondy'ego i U. S. R. Murty'ego.
\end{uwaga}

\begin{twierdzenie}[Wagner -- sformułowanie przez minory]
\label{thm:tg-wagner}
Skończony graf prosty $G$ jest planarny wtedy i tylko wtedy,
gdy nie zawiera $K_5$ ani $K_{3,3}$ jako minora.
\end{twierdzenie}

\begin{uwaga}
Jeśli $G$ nie jest planarny, to z twierdzenia Kuratowskiego
zawiera subdywizję $K_5$ lub $K_{3,3}$; kontrakcja
podzielonych krawędzi daje stosowny minor.
Jeśli natomiast $G$ zawiera taki minor, nie może być
planarny na mocy udowodnionej wyżej niezmienniczości
planarności przy operacjach minorowych.
Zatem po przyjęciu pełnego twierdzenia Kuratowskiego
kryterium Wagnera jest już uzasadnione. Niezależny,
klasyczny dowód Wagnera jest osobnym głębokim wynikiem. Nie należy utożsamiać relacji
,,jest minorem'' z relacją ,,jest podgrafem topologicznym'':
są to różne pojęcia, mimo że w kryterium planarności
prowadzą do tego samego zbioru dwóch przeszkód.
\end{uwaga}

\begin{przyklad}[Dlaczego sama nierówność Eulera nie wystarcza]
Dla grafu $K_{3,3}$ mamy $n=6$ i $m=9$.
Ograniczenie $m\le3n-6$ daje tylko $9\le12$ i nie wykrywa
nieplanarności. Dopiero wykorzystanie dwudzielności,
a więc braku cykli nieparzystych, daje sprzeczność
$9\le2\cdot6-4=8$. Dla innych grafów mogą być potrzebne
jeszcze silniejsze narzędzia, takie jak twierdzenie Kuratowskiego.
\end{przyklad}

\section{Grafy zewnętrznie planarne}
\label{sec:tg-outerplanar}

\begin{definicja}[Zewnętrzna planarność]
Graf skończony nazywamy \emph{zewnętrznie planarnym}, jeśli
ma zanurzenie w płaszczyźnie, w którym każdy wierzchołek
leży na brzegu ściany nieograniczonej.
\end{definicja}

\begin{twierdzenie}[Liczba krawędzi grafu zewnętrznie planarnego]
Niech $G$ będzie prostym, spójnym grafem zewnętrznie planarnym,
o $n\ge2$ wierzchołkach i $m$ krawędziach. Wtedy
$m\le2n-3$.
\end{twierdzenie}

\begin{proof}
Ustalmy zanurzenie z wszystkimi wierzchołkami na brzegu
ściany zewnętrznej. Jeśli $n=2$, mamy co najwyżej jedną
krawędź i wynik jest oczywisty. Przyjmijmy $n\ge3$.
Najpierw uzupełnimy graf do grafu dwuspójnego bez zmiany
zbioru wierzchołków. Jeśli istnieje wierzchołek rozcinający
$v$, obchód brzegu zewnętrznego odwiedza $v$ co najmniej
dwukrotnie. Między dwoma kolejnymi wystąpieniami $v$
wybieramy sąsiadów $a,b$ należących do różnych bloków
przy $v$. Mały łuk w ścianie zewnętrznej, omijający $v$,
pozwala poprowadzić krawędź $ab$ bez przecięcia innych
krawędzi. Krawędź ta zmniejsza liczbę bloków w rozkładzie
blokowym. Powtarzając operację, otrzymujemy graf dwuspójny
na tym samym zbiorze wierzchołków, nadal zewnętrznie planarny.
W grafie dwuspójnym obchód ściany zewnętrznej jest cyklem
prostym: powtórzenie wierzchołka na obwodzie dawałoby punkt
rozdzielający. Ponieważ każdy wierzchołek leży na brzegu
zewnętrznym, jest to cykl długości $n$.

Każda pozostała ściana jest ograniczona cyklem prostym.
Jeżeli ma co najmniej cztery boki, można dorysować przekątną
między dwoma niekolejnymi wierzchołkami tego cyklu,
wewnątrz otwartego dysku będącego ścianą. Nie powstaje
przecięcie ani zmiana zewnętrznego obwodu. Operację
powtarzamy, aż każda wewnętrzna ściana będzie trójkątem.
Ponieważ za każdym razem zwiększamy liczbę krawędzi,
a graf pozostaje prosty na $n$ wierzchołkach, proces
kończy się po skończonej liczbie kroków.

Niech $f$ oznacza liczbę jego ścian na sferze.
Suma długości ich obwodów wynosi $2m$, przy czym
ściana zewnętrzna ma obwód długości $n$,
a każda z $f-1$ ścian wewnętrznych długości trzy.
Stąd $2m=n+3(f-1)$.
Z równania Eulera $f=2-n+m$ dostajemy
$2m=n+3(1-n+m)$, czyli $m=2n-3$.
Usunięcie dołożonych krawędzi daje nierówność dla grafu $G$.
Dla $n=2$ zachodzi po prostu $m\le1$, co zgadza się ze wzorem.
\end{proof}

\begin{uwaga}
W poprzednim dowodzie korzystamy z elementarnego faktu,
że każde zewnętrznie planarne zanurzenie można uzupełnić
do maksymalnego zewnętrznie planarnego grafu na tym samym
zbiorze wierzchołków. Można go również udowodnić
indukcyjnie przez rozkład na bloki i dołączanie kolejnych
krawędzi na brzegu ściany zewnętrznej. Samo liczenie
$2m=n+3(f-1)$ dotyczy dopiero grafu \emph{uzupełnionego},
a nie dowolnego grafu zewnętrznie planarnego.
\end{uwaga}

\section{Kolorowanie wierzchołków: twierdzenie o pięciu kolorach}
\label{sec:tg-five-color}

\begin{definicja}[Poprawne kolorowanie]
Niech $G=(V,E)$ będzie grafem prostym, a $q\ge1$ liczbą całkowitą.
\emph{Poprawnym $q$-kolorowaniem wierzchołków} nazywamy
funkcję $c:V\to\{1,\ldots,q\}$, dla której
$c(u)\ne c(v)$ przy każdej krawędzi $uv\in E$.
\emph{Liczbą chromatyczną} $\chi(G)$ nazywamy najmniejsze
takie $q$. Dla grafu o co najmniej jednym wierzchołku
bez krawędzi mamy $\chi(G)=1$, ponieważ wszystkie
wierzchołki mogą otrzymać ten sam kolor.
Jedynie dla grafu zerowego $G_0$ przyjmujemy $\chi(G_0)=0$.
Przykładowo $\chi(\overline{K_4})=1$, natomiast
$\chi(K_4)=4$; ponadto $\chi(G_0)=0$.
Nie należy mylić liczby chromatycznej z liczbą wszystkich
poprawnych kolorowań: dla $\overline{K_4}$ jest ich $q^4$
przy $q$ oznaczonych kolorach.
\end{definicja}

\begin{lemat}[Wierzchołek małego stopnia]
Każdy skończony prosty graf planarny o co najmniej jednym
wierzchołku ma wierzchołek stopnia co najwyżej pięć.
\end{lemat}

\begin{proof}
Dla $n\le2$ teza jest oczywista. Jeśli $n\ge3$, możemy
najpierw dołożyć krawędzie łączące składowe w planie,
uzyskując graf spójny i prosty. Twierdzenie
\ref{thm:tg-planarne-bound} daje $2m\le6n-12<6n$.
Z lematu o uściskach dłoni średni stopień wynosi $2m/n<6$.
Zatem co najmniej jeden stopień jest nie większy niż pięć.
\end{proof}

\begin{twierdzenie}[Twierdzenie o pięciu kolorach]
\label{thm:tg-5color}
Każdy skończony graf prosty planarny można poprawnie
pokolorować co najwyżej pięcioma kolorami.
\end{twierdzenie}

\begin{proof}
Stosujemy indukcję po liczbie wierzchołków. Dla
$n\le5$ twierdzenie jest oczywiste. Niech $G$
będzie planarny o $n>5$ wierzchołkach. Z lematu wybieramy
wierzchołek $v$ stopnia co najwyżej pięć. Z założenia
indukcyjnego graf $G-v$ ma poprawne pięciokolorowanie.
Jeżeli sąsiedzi $v$ używają co najwyżej czterech kolorów,
możemy nadać $v$ piąty wolny kolor. Rozważmy więc
jedyny istotny przypadek: $v$ ma pięciu różnych sąsiadów
$v_1,\ldots,v_5$, ustawionych w tej kolejności wokół $v$,
a ich kolory to odpowiednio $1,2,3,4,5$.

Rozważmy podgraf $H_{13}$ grafu $G-v$ indukowany przez
wierzchołki o kolorach $1$ i $3$. Jeżeli $v_1$ i $v_3$
należą do różnych składowych $H_{13}$, zamieniamy
kolory $1$ i $3$ na składowej zawierającej $v_1$.
Żadna krawędź nie naruszy poprawności kolorowania,
a kolor $1$ zniknie z sąsiedztwa $v$; kolorujemy wtedy
$v$ kolorem $1$.

W przeciwnym przypadku istnieje ścieżka w $H_{13}$
łącząca $v_1$ z $v_3$. Razem z krawędziami
$vv_1$ i $vv_3$ tworzy ona krzywą Jordana.
W ustalonym zanurzeniu planarnym sąsiedzi $v_2$ i $v_4$
leżą po przeciwnych stronach tej krzywej, ponieważ
ich krawędzie do $v$ wychodzą pomiędzy dwiema wskazanymi
krawędziami w przeciwległych sektorach.
Zatem nie może istnieć ścieżka złożona z wierzchołków
kolorów $2$ i $4$, łącząca $v_2$ z $v_4$:
przecinałaby poprzednią krzywą, lecz jej kolory są
rozłączne z $\{1,3\}$, a wierzchołka $v$ nie zawiera.
Możemy więc zamienić kolory $2$ i $4$ w składowej
$H_{24}$ zawierającej $v_2$. Kolor $2$ znika
z sąsiedztwa $v$ i można go przypisać wierzchołkowi $v$.
W obu przypadkach przedłużamy kolorowanie, co kończy indukcję.
\end{proof}

\begin{center}
\begin{tikzpicture}[scale=.87,line cap=round]
  \node[tgBlueVertex] (v) at (0,0) {$v$};
  \foreach \i/\ang/\col in {1/90/tgBlue,2/18/tgTeal,3/-54/tgCoral,4/-126/tgGold,5/162/tgPurple} {
    \node[tgVertex,draw=\col,fill=\col!15] (n\i) at (\ang:2.0) {$\i$};
    \draw[tgEdge] (v)--(n\i);
  }
  \draw[tgAccent,very thick,rounded corners] (n1) .. controls (1.2,2.8) and (2.7,-1.1) .. (n3);
  \node[text=tgNavy,font=\footnotesize,align=center] at (0,-2.75)
    {Łańcuch kolorów 1--3 rozdziela sąsiadów 2 i 4;\\ nie może go przeciąć łańcuch kolorów 2--4.};
\end{tikzpicture}
\end{center}

\begin{uwaga}[Twierdzenie o czterech kolorach]
Wynik silniejszy mówi, że każdemu skończonemu grafowi planarnemu
wystarczają cztery kolory. Jego znany dowód korzysta
z obszernej analizy przypadków i weryfikacji komputerowej;
nie przedstawiamy go jako prostego następstwa powyższego
argumentu. Dowód pięciokolorowalności jest natomiast
w pełni elementarny i ilustruje znaczenie twierdzenia Jordana.
\end{uwaga}


% ============================================================
% Iteracja 4, rozdział 11
% ============================================================
\chapter{Grafy na powierzchniach, grafy dualne i homologia komórkowa}
\label{ch:tg-powierzchnie}

\section{Zanurzenia komórkowe i rodzaj powierzchni}
\label{sec:tg-cellular-surface}

\begin{definicja}[Zanurzenie komórkowe]
Niech $\Sigma$ będzie zwartą, spójną powierzchnią bez brzegu,
a $G$ skończonym spójnym grafem. Zanurzenie
$\iota:|G|\hookrightarrow\Sigma$ nazywamy \emph{komórkowym},
jeżeli każda składowa dopełnienia $\Sigma\setminus\iota(|G|)$
jest homeomorficzna z otwartym dyskiem. Te składowe nazywamy
\emph{ścianami} zanurzenia; po dołączeniu ich jako dwukomórek
otrzymujemy strukturę CW na $\Sigma$.
\end{definicja}

\begin{uwaga}[Dlaczego wymagana jest komórkowość]
Graf jednowierzchołkowy zanurzony w torusie ma dopełnienie
homeomorficzne z torusem bez punktu, a nie z otwartym dyskiem.
Nie wolno więc liczyć tego dopełnienia jako pojedynczej
dwukomórki. Równość $n-m+f=\chi(\Sigma)$ wymaga
zanurzenia komórkowego. W szczególności nie stosujemy jej
automatycznie do dowolnego rysunku grafu na torusie.
\end{uwaga}

\begin{twierdzenie}[Euler--Poincaré dla powierzchni orientowalnej]
\label{thm:tg-surface-euler}
Niech $\Sigma_g$ będzie zwartą, spójną, orientowalną powierzchnią
bez brzegu rodzaju $g\ge0$. Niech spójny graf $G$ będzie
zanurzony komórkowo w $\Sigma_g$, a $n,m,f$ oznaczają
odpowiednio liczby wierzchołków, krawędzi i ścian tego zanurzenia.
Wtedy
\[
 n-m+f=2-2g.
\]
\end{twierdzenie}

\begin{proof}
Z założenia wierzchołki, otwarte krawędzie i ściany
tworzą rozkład CW powierzchni $\Sigma_g$ na $n$
komórek wymiaru zero, $m$ komórek wymiaru jeden i $f$
komórek wymiaru dwa. Z definicji charakterystyki Eulera
$\chi(\Sigma_g)=n-m+f$. Klasyfikacja zwartych, spójnych,
orientowalnych powierzchni mówi, że
$\Sigma_g\cong\#^g(S^1\times S^1)$ i
$\chi(\Sigma_g)=2-2g$. Łącząc te równości,
otrzymujemy tezę. W tym dowodzie przyjmujemy znany,
wcześniej udowodniony w topologii powierzchni wzór
na charakterystykę Eulera \(\Sigma_g\); nie jest on
ukrytą konsekwencją samej definicji grafu.
\end{proof}

\begin{stwierdzenie}[Oszacowanie rodzaju z liczby krawędzi]
Niech prosty spójny graf $G$ o $n\ge3$ wierzchołkach
i $m$ krawędziach ma komórkowe zanurzenie w orientowalnej
powierzchni $\Sigma_g$. Wtedy
\[
 m\le3n-6+6g.
\]
Jeżeli $G$ jest dwudzielny, to
$m\le2n-4+4g$.
\end{stwierdzenie}

\begin{proof}
Każdy obchód ścianowy ma długość co najmniej trzy,
a w przypadku grafu dwudzielnego co najmniej cztery.
Jest to lokalny argument identyczny jak dla sfery;
założenie spójności i $n\ge3$ wyklucza ściany długości
dwa w grafie prostym. Zatem odpowiednio $3f\le2m$
lub $4f\le2m$. Podstawiając
$f=2-2g-n+m$ z poprzedniego twierdzenia,
otrzymujemy obie nierówności po uporządkowaniu.
\end{proof}

\begin{przyklad}[Dlaczego $K_5$ potrzebuje uchwytu]
Dla $K_5$ mamy $n=5$ i $m=10$. Gdyby istniało
komórkowe zanurzenie w $\Sigma_0=S^2$, musiałaby zachodzić
nierówność $10\le9$, co jest niemożliwe.
Oszacowanie dla ogólnego rodzaju daje
$10\le9+6g$, a więc $g\ge1$.
Istnieje zanurzenie $K_5$ w torusie, więc minimalny
rodzaj orientowalny tego grafu jest równy jeden.
Podobnie $K_{3,3}$ ma minimalny rodzaj orientowalny jeden.
Samo oszacowanie daje dolne ograniczenia; stwierdzenie
istnienia zanurzeń w torusie jest dodatkowym faktem
geometrycznym, a nie wnioskiem z nierówności.
\end{przyklad}

\section{Cykliczne porządki incydencji i grafy wstęgowe}
\label{sec:tg-rotation}

Samo abstrakcyjne $G=(V,E)$ nie zapamiętuje, jak krawędzie
są ułożone wokół wierzchołków na zorientowanej powierzchni.
Dane te można zakodować skończoną rodziną permutacji.

\begin{definicja}[Półkrawędź i układ rotacji]
Niech $G$ będzie skończonym spójnym multigrafem. Każda
krawędź daje dwie \emph{półkrawędzie} (ang. \emph{darts});
ich zbiór oznaczamy przez $D$. Także pętla daje dwie
różne półkrawędzie incydentne z tym samym wierzchołkiem.
Przez $\alpha:D\to D$ oznaczamy inwolucję bez punktów stałych,
która zamienia półkrawędzie tej samej krawędzi.
\emph{Układem rotacji} nazywamy permutację $\sigma$ zbioru $D$,
której cykle odpowiadają cyklicznej kolejności półkrawędzi
wokół poszczególnych wierzchołków. Dla wierzchołka
izolowanego kolejność jest pusta; w dalszych rachunkach
zakładamy graf spójny z co najmniej jedną krawędzią.
\end{definicja}

\begin{twierdzenie}[Powierzchnia wyznaczona przez układ rotacji]
\label{thm:tg-rotation-genus}
Niech $G$ będzie skończonym spójnym multigrafem o $n$
wierzchołkach i $m\ge1$ krawędziach, a $\alpha$ inwolucją parującą półkrawędzie i $\sigma$
permutacją, której cyklami są dokładnie zbiory półkrawędzi
przy kolejnych wierzchołkach, z zadanym porządkiem cyklicznym. Istnieje zorientowana, zamknięta,
spójna powierzchnia $\Sigma_g$ oraz komórkowe zanurzenie
$G\hookrightarrow\Sigma_g$ realizujące te porządki.
Jeżeli $f$ oznacza liczbę cykli permutacji
$\varphi:=\sigma\circ\alpha$ na $D$, to
\[
 2-2g=n-m+f,
 \qquad g=\frac{2-n+m-f}{2}.
\]
\end{twierdzenie}

\begin{proof}[Konstrukcja przez otoczenie wstęgowe]
Każdy wierzchołek zastępujemy małym zorientowanym
dyskiem. Na jego brzegu wyznaczamy rozłączne odcinki
odpowiadające incydentnym półkrawędziom, zachowując
kolejność cykliczną zapisaną w $\sigma$.
Dla każdej krawędzi doklejamy prostokąt
$[0,1]\times[0,1]$ do dwóch odpowiednich odcinków
brzegowych dysków. Doklejenia dobieramy w sposób
zachowujący orientację powstającej powierzchni z brzegiem.
Dopuszczamy, że oba końce prostokąta spotykają ten sam dysk,
co modeluje pętlę.

Uzyskujemy spójną, orientowalną powierzchnię wstęgową $N$
(z brzegiem), która deformacyjnie retraktuje się na $|G|$.
Śledzenie jej brzegu polega na przejściu wzdłuż danej
półkrawędzi na drugi koniec tej samej krawędzi (działanie
$\alpha$), a następnie do kolejnej półkrawędzi przy
wierzchołku (działanie $\sigma$). A więc składowe brzegu
$N$ odpowiadają dokładnie cyklom $\varphi=\sigma\alpha$.
Jest ich $f$.

Do każdej składowej brzegu $N$ doklejamy dysk.
Otrzymujemy zwartą, spójną powierzchnię zorientowaną
bez brzegu, w której $G$ jest szkieletem jednowymiarowym,
a wnętrza dołączonych dysków są komórkowymi ścianami.
Ponieważ $N$ ma charakterystykę $n-m$, po dołączeniu
$f$ dysków mamy $\chi(\Sigma)=n-m+f$.
Klasyfikacja powierzchni daje $\chi(\Sigma)=2-2g$,
co kończy dowód.
\end{proof}

\begin{przyklad}[Dwa układy rotacji tego samego grafu theta]
Rozważmy multigraf z wierzchołkami $u,v$ połączonymi trzema
równoległymi krawędziami. Niech $a_i$ i $b_i$ będą ich
półkrawędziami przy odpowiednio $u$ i $v$.
Wtedy $\alpha=(a_1b_1)(a_2b_2)(a_3b_3)$.
W obu modelach przy $u$ ustalamy porządek
$(a_1a_2a_3)$.

Jeśli przy $v$ przyjmiemy porządek $(b_1b_3b_2)$,
permutacja $\varphi=\sigma\alpha$ ma trzy cykle
$(a_1b_3)$, $(a_2b_1)$, $(a_3b_2)$.
Zatem $f=3$ oraz $2-2g=2-3+3=2$; otrzymujemy sferę.
Jeśli natomiast przy $v$ przyjmiemy ten sam porządek
$(b_1b_2b_3)$, $\varphi$ ma jeden cykl długości sześć:
$(a_1b_2a_3b_1a_2b_3)$. Teraz $f=1$ i
$2-2g=2-3+1=0$; otrzymujemy torus.
To samo abstrakcyjne $G$ ma więc dwa komórkowe zanurzenia
o różnych rodzajach powierzchni.
\end{przyklad}

\begin{center}
\begin{tikzpicture}[scale=.95,line cap=round]
  \begin{scope}[shift={(-2.9,0)}]
    \node[tgBlueVertex] (u) at (-1.05,0) {$u$};
    \node[tgRedVertex] (v) at (1.05,0) {$v$};
    \draw[tgBlue,line width=1.8pt] (u) to[bend left=43] (v);
    \draw[tgTeal,line width=1.8pt] (u)--(v);
    \draw[tgCoral,line width=1.8pt] (u) to[bend right=43] (v);
    \node[text=tgNavy,font=\small] at (0,-1.38) {$f=3$, $g=0$};
    \node[text=tgPurple,font=\footnotesize] at (0,-1.72) {$(b_1b_3b_2)$};
  \end{scope}
  \begin{scope}[shift={(2.9,0)}]
    \node[tgBlueVertex] (u) at (-1.05,0) {$u$};
    \node[tgRedVertex] (v) at (1.05,0) {$v$};
    \draw[tgBlue,line width=1.8pt] (u) to[bend left=43] (v);
    \draw[tgTeal,line width=1.8pt] (u)--(v);
    \draw[tgCoral,line width=1.8pt] (u) to[bend right=43] (v);
    \node[text=tgNavy,font=\small] at (0,-1.38) {$f=1$, $g=1$};
    \node[text=tgPurple,font=\footnotesize] at (0,-1.72) {$(b_1b_2b_3)$};
  \end{scope}
  \node[text=tgNavy,font=\footnotesize,align=center] at (0,-2.48)
    {Różnica leży w cyklicznym porządku \emph{półkrawędzi};\\
     dwuwymiarowy rysunek nie zastępuje danych rotacji.};
\end{tikzpicture}
\end{center}

\section{Graf dualny na sferze}
\label{sec:tg-dual}

\begin{definicja}[Graf dualny]
Niech $G\subset S^2$ będzie spójnym, skończonym grafem
płaskim o zbiorze ścian $F$. W każdej ścianie wybieramy
punkt $v_F$. Dla każdej krawędzi $e$ grafu $G$ rysujemy
łuk $e^*$ łączący punkty dwóch ścian leżących po
przeciwnych stronach $e$, przecinający $e$ dokładnie
raz i nieprzecinający innych krawędzi pierwotnych.
Graf z wierzchołkami $v_F$ i krawędziami $e^*$
nazywamy \emph{grafem dualnym} i oznaczamy $G^*$.
Jeśli obie strony $e$ należą do tej samej ściany,
$e^*$ jest pętlą; dualny graf może mieć krawędzie równoległe.
\end{definicja}

\begin{stwierdzenie}[Podstawowe własności dualności]
\label{thm:tg-dual-basic}
Dla skończonego spójnego grafu płaskiego $G\subset S^2$
graf dualny $G^*$ jest spójny i ma
$|V(G^*)|=f(G)$ oraz $|E(G^*)|=|E(G)|$.
Ponadto $(G^*)^*$ jest równoważny $G$ jako graf zanurzony
w sferze. Krawędź $e\in E(G)$ jest mostem wtedy i tylko
wtedy, gdy $e^*$ jest pętlą w $G^*$.
\end{stwierdzenie}

\begin{proof}
Pierwsze dwie równości wynikają z konstrukcji.
Aby wykazać spójność, bierzemy dwa punkty w różnych
ścianach i łączymy je krzywą na sferze, omijającą
wierzchołki pierwotnego grafu i przecinającą jego
krawędzie poprzecznie. Zapisując kolejne napotkane
ściany, otrzymujemy spacer w $G^*$.

Przy skonstruowaniu drugiego grafu dualnego można
wybrać jego wierzchołki w pobliżu pierwotnych
wierzchołków $G$ i prowadzić krawędzie wzdłuż
pierwotnych krawędzi. Daje to izomorfizm zanurzeń.

Jeżeli $e$ jest mostem, usunięcie jej wnętrza nie
rozdziela żadnej ściany na dwie; oba jej boki należą
do tej samej ściany, więc $e^*$ jest pętlą.
Odwrotnie, jeśli oba boki $e$ są po tej samej stronie,
to $e$ nie może należeć do żadnego cyklu prostego:
cykl tworzyłby krzywą Jordana, której dwa boki
przy $e$ należą do różnych ścian. Z kryterium mostu
wynika wtedy, że $e$ jest mostem.
\end{proof}

\begin{center}
\begin{tikzpicture}[scale=.88,line cap=round,line join=round]
 \begin{scope}[shift={(-3.1,0)}]
  \coordinate (a) at (-1.3,-1);\coordinate (b) at (1.3,-1);
  \coordinate (c) at (1.3,1);\coordinate (d) at (-1.3,1);
  \fill[tgBlue!9] (a)--(b)--(c)--cycle;
  \fill[tgTeal!9] (a)--(c)--(d)--cycle;
  \draw[tgEdge] (a)--(b)--(c)--(d)--cycle;
  \draw[tgAccent,line width=2pt] (a)--(c);
  \foreach \pt in {a,b,c,d} \node[tgBlueVertex] at (\pt) {};
  \node[text=tgBlue] at (.52,-.44) {$F_1$};
  \node[text=tgTeal] at (-.5,.45) {$F_2$};
  \node[text=tgNavy,font=\footnotesize] at (0,-1.55) {$G$};
 \end{scope}
 \begin{scope}[shift={(3.1,0)}]
  \node[tgBlueVertex] (f1) at (0,-.6) {$F_1$};
  \node[tgTealVertex] (f2) at (0,.8) {$F_2$};
  \node[tgRedVertex] (f3) at (1.8,.1) {$F_3$};
  \draw[tgAccent] (f1)--(f2);
  \draw[tgBlue,line width=1.7pt] (f1) to[bend right=27] (f3);
  \draw[tgBlue,line width=1.7pt] (f1) to[bend left=27] (f3);
  \draw[tgTeal,line width=1.7pt] (f2) to[bend left=27] (f3);
  \draw[tgTeal,line width=1.7pt] (f2) to[bend right=27] (f3);
  \node[text=tgNavy,font=\footnotesize] at (.7,-1.55) {$G^*$};
 \end{scope}
\end{tikzpicture}

{\footnotesize Dualność zamienia ściany w wierzchołki i zachowuje
liczbę krawędzi. Dwie krawędzie między tymi samymi ścianami
stają się równoległymi krawędziami dualnymi.}
\end{center}

\begin{twierdzenie}[Drzewo dualne dopełniające]
\label{thm:tg-dual-spanning}
Niech $G\subset S^2$ będzie spójnym grafem płaskim,
a $T$ jego drzewem rozpinającym. Wtedy zbiór
$\{e^*:e\in E(G)\setminus E(T)\}$ tworzy drzewo
rozpinające w $G^*$.
\end{twierdzenie}

\begin{proof}
Z twierdzenia Eulera $f=m-n+2$, a więc liczba krawędzi
dualnych, które pozostawiamy, wynosi
$m-(n-1)=f-1=|V(G^*)|-1$.
Wystarczy zatem wykazać spójność wybranego podgrafu dualnego.
Gdyby nie był spójny, istniałby niepusty przekrój
$\delta_{G^*}(S)$, którego żadna krawędź nie została
wybrana. Oznaczałoby to, że wszystkie jego krawędzie
pierwotne należą do $T$. Jednak przekrój w grafie
dualnym odpowiada zbiorowi cykli w grafie pierwotnym
(najprostszy przypadek to obwód brzegu sumy ścian
należących do $S$). Niech $A$ będzie zbiorem pierwotnych krawędzi odpowiadających
temu przekrojowi. Każdy wierzchołek pierwotny ma parzystą
liczbę incydencji z $A$: kiedy przechodzimy wokoło niego
po kolejnych sektorach, liczba zmian między ścianami z $S$
i spoza $S$ jest parzysta. Dotyczy to również pętli,
których dwie półkrawędzie liczymy osobno. Zatem $A$
wyznacza cykl nad $\mathbb F_2$. Ponieważ przekrój jest
niepusty, $A\ne\varnothing$ i graf $(V,A)$ zawiera cykl
(w razie pętli jest to cykl długości jeden). Wszystkie
krawędzie $A$ należałyby do drzewa $T$, sprzeczność.
Wybrany podgraf dualny jest więc spójny i ma $f-1$
krawędzi, czyli jest drzewem rozpinającym.
\end{proof}

\section{Od ścian do kompleksu łańcuchowego powierzchni}
\label{sec:tg-surface-chains}

\begin{definicja}[Operator brzegu dwukomórek]
Niech skończony spójny graf $G$ będzie komórkowo zanurzony
w zorientowanej zamkniętej powierzchni $\Sigma$.
Oznaczmy przez $F$ zbiór ścian. Definiujemy grupy
$C_2:=\mathbb Z^F$, $C_1:=\mathbb Z^E$ i
$C_0:=\mathbb Z^V$. Operator $\partial_1:C_1\to C_0$
jest macierzą incydencji z poprzednich rozdziałów,
a $\partial_2:C_2\to C_1$ wysyła zorientowaną ścianę
na sumę zorientowanych krawędzi jej obchodu brzegowego,
ze znakami zależnymi od kierunku przejścia.
\end{definicja}

\begin{stwierdzenie}[Warunek łańcuchowy]
Dla operatorów $\partial_2$ i $\partial_1$ określonych
powyżej zachodzi $\partial_1\partial_2=0$.
\end{stwierdzenie}

\begin{proof}
Każda ściana ma obchód będący zamkniętym spacerem.
Pod zastosowaniem $\partial_1$ wkład przy każdym kroku
jest różnicą końca i początku krawędzi. W sumie
wierzchołek będący końcem jednego kroku jest zarazem
początkiem następnego; wkłady teleskopowo się znoszą.
Zatem brzeg brzegu każdej komórki wymiaru dwa jest zerowy,
a przez liniowość zachodzi $\partial_1\partial_2=0$.
\end{proof}

\begin{twierdzenie}[Homologia komórkowa powierzchni orientowalnej]
\label{thm:tg-surface-homology}
Niech $G$ będzie skończonym spójnym grafem z komórkowym
zanurzeniem w zorientowanej, spójnej zamkniętej powierzchni
$\Sigma_g$. Dla powyższego kompleksu łańcuchowego
otrzymujemy
\[
H_0(\Sigma_g;\mathbb R)\cong\mathbb R,
\qquad H_1(\Sigma_g;\mathbb R)\cong\mathbb R^{2g},
\qquad H_2(\Sigma_g;\mathbb R)\cong\mathbb R.
\]
\end{twierdzenie}

\begin{proof}
Graf jest spójny, zatem $\operatorname{rank}\partial_1=n-1$
i $\dim\ker\partial_1=m-n+1$. Udowodnimy, że
$\ker\partial_2$ jest jednowymiarowe. W orientowanej
powierzchni każda krawędź sąsiaduje z dwiema stronami
ścian, z przeciwnymi znakami w ich obwodach. Jeśli
$z=\sum_{F\in\mathcal F}a_F[F]$ spełnia
$\partial_2z=0$, to po porównaniu współczynnika
przy dowolnej krawędzi między ścianami $F,F'$ dostajemy
$a_F=a_{F'}$. Dualny graf ścian jest spójny:
dowolne dwie ściany można połączyć łukiem przecinającym
graf transwersalnie i omijającym wierzchołki.
Zatem wszystkie $a_F$ są równe. Z drugiej strony
suma wszystkich zorientowanych ścian ma brzeg zero,
bo każda krawędź pojawia się dwukrotnie z przeciwnymi
znakami. Stąd
$\ker\partial_2=\mathbb R\sum_F[F]$, a więc
$\operatorname{rank}\partial_2=f-1$.

Wynika z tego
\[
\dim H_1=\dim\ker\partial_1-\operatorname{rank}\partial_2
=(m-n+1)-(f-1)=m-n-f+2=2g,
\]
gdzie ostatnia równość wynika z twierdzenia Eulera--Poincarégo.
Ponadto $H_2=\ker\partial_2\cong\mathbb R$,
a $H_0=C_0/\operatorname{im}\partial_1\cong\mathbb R$.
To wyprowadza homologiczne znaczenie rodzaju powierzchni
bez pomijania roli operatora $\partial_2$.
\end{proof}

\begin{przyklad}[Sfera i torus -- co zmienia doklejenie ścian]
Rozważmy trójkąt jako graf na sferze. Ma jedną niezależną
klasę w $H_1(|G|;\mathbb Z)\cong\mathbb Z$, reprezentowaną
przez obieg trójkąta. Po dołączeniu dwóch ścian
(wewnętrznej i zewnętrznej) otrzymujemy sferę.
Brzeg pierwszej ściany jest dodatnim obiegiem trójkąta,
a drugiej jego przeciwnym obiegiem, więc
$H_1(S^2;\mathbb Z)=0$.

Standardowy model torusa ma z kolei jedną zerokomórkę,
dwie jednokomórki $a,b$ i jedną dwukomórkę dołączoną
wzdłuż słowa $aba^{-1}b^{-1}$. Operator $\partial_1$ znika,
a $\partial_2$ również znika, ponieważ suma wykładników
każdego generatora w słowie przyklejającym jest zerowa.
Stąd $H_1(T^2;\mathbb Z)\cong\mathbb Z^2$, choć grupa
fundamentalna samego szkieletu jednowymiarowego jest
wolna: $\pi_1(S^1\vee S^1)\cong F_2$.
Dołączenie dwukomórki dokłada relację
$aba^{-1}b^{-1}=1$, która zmienia grupę fundamentalną
na $\pi_1(T^2)\cong\mathbb Z^2$.
\end{przyklad}

\begin{center}
\begin{tikzpicture}[scale=.90,>=Stealth,line cap=round]
  \fill[tgBlue!8] (0,0) rectangle (4,2.7);
  \draw[tgEdge] (0,0) rectangle (4,2.7);
  \draw[tgAccent,-{Stealth[length=2.8mm]},line width=2.2pt] (.25,0)--(3.65,0);
  \draw[tgAccent,-{Stealth[length=2.8mm]},line width=2.2pt] (.25,2.7)--(3.65,2.7);
  \draw[tgTeal,-{Stealth[length=2.8mm]},line width=2.2pt] (0,.25)--(0,2.35);
  \draw[tgTeal,-{Stealth[length=2.8mm]},line width=2.2pt] (4,.25)--(4,2.35);
  \node[text=tgCoral,font=\small] at (2,-.35) {$a$};
  \node[text=tgCoral,font=\small] at (2,3.08) {$a$};
  \node[text=tgTeal,font=\small] at (-.36,1.35) {$b$};
  \node[text=tgTeal,font=\small] at (4.36,1.35) {$b$};
  \node[text=tgNavy,align=center] at (2,1.35) {$aba^{-1}b^{-1}$\\[2pt] jedna ściana};
\end{tikzpicture}

{\footnotesize Kwadrat fundamentalny torusa: przeciwległe boki
utożsamiamy zgodnie ze strzałkami, otrzymując $n=1$, $m=2$, $f=1$.}
\end{center}


% ============================================================
% Iteracja 4, rozdział 12
% ============================================================
\chapter{Wielomian Tutte'a, kolorowania i dualność algebraiczna}
\label{ch:tg-tutte}

Graf planarny daje jednocześnie obraz geometryczny i zbiór danych
kombinatorycznych. Wielomian Tutte'a jest niezmiennikiem, który
łączy liczenie drzew rozpinających, kolorowania, przepływy
i dualność grafów płaskich. Zaczniemy od definicji przez
sumowanie po wszystkich podzbiorach krawędzi, a następnie
wyprowadzimy efektywny algorytm obliczania tego wielomianu.

\section{Usuwanie i kontrakcja krawędzi}
\label{sec:tg-deletion-contraction}

\begin{definicja}[Dwie operacje na krawędzi]
Niech $G=(V,E)$ będzie skończonym multigrafem, a $e\in E$.
Przez $G\setminus e$ oznaczamy graf po usunięciu krawędzi $e$
bez usuwania jej końców. Jeśli $e$ nie jest pętlą,
przez $G/e$ oznaczamy graf powstały przez kontrakcję $e$:
utożsamiamy jej dwa końce i usuwamy $e$; pozostałe krawędzie
zachowują swoje wystąpienia, także jeśli po utożsamieniu
zamieniają się w pętle. Przy liczeniu wierzchołków grafu $G/e$
liczba wierzchołków maleje dokładnie o jeden.
\end{definicja}

\begin{definicja}[Rząd i zerowość cykliczna zbioru krawędzi]
Niech $G=(V,E)$ będzie skończonym multigrafem, a $A\subseteq E$.
Przez $k_G(A)$ oznaczamy liczbę składowych grafu $(V,A)$,
łącznie z izolowanymi wierzchołkami. Definiujemy
$r_G(A):=|V|-k_G(A)$ oraz
$\nu_G(A):=|A|-r_G(A)=|A|-|V|+k_G(A)$.
Liczba $\nu_G(A)$ jest wymiarem przestrzeni cykli grafu $(V,A)$.
\end{definicja}

\begin{definicja}[Wielomian Tutte'a]
\label{def:tg-tutte}
Dla dowolnego skończonego multigrafu $G=(V,E)$ definiujemy
\[
T_G(x,y):=\sum_{A\subseteq E}
 (x-1)^{k_G(A)-k_G(E)}
 (y-1)^{|A|-|V|+k_G(A)}\in\mathbb Z[x,y].
\]
Wykładniki są nieujemne: dodawanie krawędzi nie zwiększa liczby
składowych, a $|A|-|V|+k_G(A)$ jest wymiarem przestrzeni cykli.
Dla grafu bez krawędzi mamy $T_G=1$, niezależnie od liczby
wierzchołków izolowanych.
\end{definicja}

\begin{twierdzenie}[Rekurencja usuwania i kontrakcji]
\label{thm:tg-tutte-recursion}
Niech $G$ będzie skończonym multigrafem, a $e\in E(G)$.
Wtedy zachodzą następujące tożsamości:
\begin{enumerate}[label=(\roman*)]
\item jeśli $e$ nie jest ani mostem, ani pętlą, to
$T_G=T_{G\setminus e}+T_{G/e}$;
\item jeśli $e$ jest mostem, to $T_G=x\,T_{G/e}$;
\item jeśli $e$ jest pętlą, to $T_G=y\,T_{G\setminus e}$.
\end{enumerate}
Wraz z wartością $T=1$ dla grafu bez krawędzi reguły te
jednoznacznie określają wielomian Tutte'a.
\end{twierdzenie}

\begin{proof}
Podzielmy sumę definiującą $T_G$ na dwa rodzaje składników:
podzbiorów $A\subseteq E\setminus\{e\}$ i podzbiorów
postaci $A=B\cup\{e\}$, gdzie $B\subseteq E\setminus\{e\}$.

Gdy $e$ nie jest mostem ani pętlą, liczba składowych całego
grafu nie zmienia się przy usunięciu ani przy kontrakcji $e$.
Składniki nie zawierające $e$ są zatem dokładnie składnikami
$T_{G\setminus e}$. Dla składników zawierających $e$ mamy
$k_G(B\cup\{e\})=k_{G/e}(B)$, ponieważ kontrakcja
nie zmienia liczby składowych. Zarazem
$|V(G/e)|=|V(G)|-1$, więc
\[
|B\cup\{e\}|-|V(G)|+k_G(B\cup\{e\})
=|B|-|V(G/e)|+k_{G/e}(B).
\]
Składniki z $e$ dają zatem dokładnie $T_{G/e}$.

Jeśli $e$ jest mostem, każdy podzbiór $B\subseteq E\setminus\{e\}$
ma o jedną składową więcej niż $B\cup\{e\}$.
Kontrakcja $e$ daje bijekcję między podzbiorami krawędzi
$G\setminus e$ a podzbiorami krawędzi $G/e$, a porównanie
obu składników związanych z tym samym $B$ daje czynnik
$(x-1)+1=x$. Stąd $T_G=xT_{G/e}$.
Jeżeli $e$ jest pętlą, jej dołożenie nie zmienia liczby
składowych, lecz zwiększa $\nu_G(A)$ o jeden.
Dwa odpowiadające sobie składniki dają czynnik
$1+(y-1)=y$, co daje $T_G=yT_{G\setminus e}$.

Każda z powyższych rekurencji zmniejsza liczbę krawędzi,
a więc pozwala obliczyć wielomian po skończonej liczbie
kroków; to również dowodzi jednoznaczności.
\end{proof}

\begin{przyklad}[Drzewo, cykl i trzy krawędzie równoległe]
Jeżeli $G$ jest drzewem o $n\ge1$ wierzchołkach,
każda krawędź jest mostem. Kontraktując je kolejno,
otrzymujemy $T_G(x,y)=x^{n-1}$.

Dla trójkąta wybierzmy krawędź $e$.
Graf $C_3\setminus e$ jest ścieżką o dwóch krawędziach,
z wielomianem $x^2$. Graf $C_3/e$ składa się z dwóch
równoległych krawędzi; ich wielomian to $x+y$.
Stąd
$T_{C_3}=x^2+x+y$.
Analogiczna rekurencja daje dla $n\ge2$
\[
 T_{C_n}(x,y)=x^{n-1}+x^{n-2}+\cdots+x+y.
\]
Dla grafu theta złożonego z dwóch wierzchołków
i trzech równoległych krawędzi otrzymujemy
$T_{\Theta}(x,y)=x+y+y^2$.
\end{przyklad}

\begin{center}
\begin{tikzpicture}[scale=.90,>=Stealth,line cap=round]
 \begin{scope}[shift={(-3.3,0)}]
  \node[tgBlueVertex] (a) at (-1,-.7) {$a$};
  \node[tgBlueVertex] (b) at (1,-.7) {$b$};
  \node[tgBlueVertex] (c) at (0,1.05) {$c$};
  \draw[tgEdge] (a)--(b)--(c);
  \draw[tgAccent] (c)--(a);
  \node[text=tgNavy,font=\small] at (0,-1.45) {$C_3$};
 \end{scope}
 \draw[tgTeal,-{Stealth[length=3mm]},line width=1.5pt] (-1.9,0)--(-1.1,0);
 \node[text=tgTeal,font=\footnotesize] at (-1.5,.45) {$\setminus e$};
 \begin{scope}[shift={(.1,0)}]
  \node[tgBlueVertex] (a) at (-.7,-.7) {};
  \node[tgBlueVertex] (b) at (.7,-.7) {};
  \node[tgBlueVertex] (c) at (0,1.05) {};
  \draw[tgTree] (a)--(b)--(c);
  \node[text=tgNavy,font=\small] at (0,-1.45) {$x^2$};
 \end{scope}
 \node[text=tgTeal,font=\small] at (2.0,.5) {$+$};
 \begin{scope}[shift={(4.0,0)}]
  \node[tgBlueVertex] (a) at (-.7,0) {};
  \node[tgBlueVertex] (b) at (.7,0) {};
  \draw[tgBlue,line width=2pt] (a) to[bend left=38] (b);
  \draw[tgAccent,line width=2pt] (a) to[bend right=38] (b);
  \node[text=tgNavy,font=\small] at (0,-1.45) {$x+y$};
 \end{scope}
\end{tikzpicture}

{\footnotesize Rekurencja Tutte'a dla trójkąta: usunięcie krawędzi
daje drzewo, kontrakcja daje dwie krawędzie równoległe.}
\end{center}

\section{Drzewa rozpinające jako specjalizacja wielomianu}
\label{sec:tg-tutte-trees}

\begin{twierdzenie}[Wartość $T_G(1,1)$]
Niech $G$ będzie skończonym spójnym multigrafem.
Wtedy $T_G(1,1)=\tau(G)$, gdzie $\tau(G)$ oznacza
liczbę drzew rozpinających, z uwzględnieniem krotności
krawędzi równoległych.
\end{twierdzenie}

\begin{proof}
Dla grafu spójnego $k_G(E)=1$. W sumie definiującej
$T_G(1,1)$ składnik odpowiadający $A\subseteq E$
nie znika wtedy i tylko wtedy, gdy obie potęgi zerowej
podstawy mają wykładnik zero, czyli
$k_G(A)=1$ oraz $|A|-|V|+k_G(A)=0$.
Te warunki oznaczają, że $(V,A)$ jest spójnym grafem
bez cykli, czyli drzewem rozpinającym. Każdy taki
składnik jest równy jeden ($0^0$ rozumiemy tu
jako wartość jednomianu z wykładnikiem zero).
Suma liczy dokładnie drzewa rozpinające.
\end{proof}

\begin{uwaga}
Twierdzenie Kirchhoffa daje tę samą liczbę z wyznacznika
minora laplasjanu. Mamy zatem dwie niezależne metody
obliczania $\tau(G)$: algebraiczną (wyznacznik) i
kombinatoryczną (wielomian Tutte'a). Dla grafu niespójnego
$T_G(1,1)$ liczy maksymalne lasy rozpinające, po jednym
drzewie rozpinającym na każdą składową.
\end{uwaga}

\begin{stwierdzenie}[Dwie dalsze specjalizacje]
Dla skończonego multigrafu $G=(V,E)$ zachodzą równości:
$T_G(2,2)=2^{|E|}$,
$T_G(2,1)=N_{\mathrm{las}}(G)$,
gdzie $N_{\mathrm{las}}(G)$ oznacza liczbę podzbiorów
krawędzi nie zawierających cyklu (w tym podzbiór pusty).
\end{stwierdzenie}

\begin{proof}
Dla $x=y=2$ w definicji każdy składnik ma wartość jeden,
a podzbiorów $E$ jest $2^{|E|}$. Przy $x=2$, $y=1$
pozostają dokładnie składniki z $\nu_G(A)=0$, czyli
podgrafy $(V,A)$ będące lasami. Każdy taki składnik
wnosi jeden. Pętle nigdy nie należą do lasu i słusznie
nie pojawiają się w tej sumie.
\end{proof}

\section{Wielomian chromatyczny i jego wzór przez Tutte'a}
\label{sec:tg-chromatic}

\begin{uwaga}[Konwencje dla grafu zerowego i pętli]
W tej części dopuszczamy również graf zerowy: ma on dokładnie
jedno kolorowanie pustą funkcją i $P_{G_0}(q)=1$, ale jego
liczbę chromatyczną definiujemy osobno jako zero. Jeśli graf
posiada pętlę, nie istnieje poprawne kolorowanie żadną liczbą
kolorów, ponieważ wierzchołek nie może mieć koloru różnego
od samego siebie. Zatem jego wielomian chromatyczny jest
wielomianem zerowym. Krawędzie równoległe nie zmieniają
liczby poprawnych kolorowań.
\end{uwaga}

\begin{definicja}[Wielomian chromatyczny]
Niech $G=(V,E)$ będzie skończonym multigrafem.
Dla $q\in\mathbb N$, $q\ge1$, przez $P_G(q)$ oznaczamy
liczbę poprawnych kolorowań wierzchołków kolorami
$1,\ldots,q$, tzn. funkcji $c:V\to\{1,\ldots,q\}$
spełniających $c(u)\ne c(v)$ dla każdej krawędzi $uv$.
Gdy $G$ ma pętlę, przyjmujemy $P_G(q)=0$.
\end{definicja}

\begin{twierdzenie}[Wzór inkluzji i ekskluzji]
\label{thm:tg-chromatic-inclusion}
Dla dowolnego skończonego multigrafu $G=(V,E)$ i każdej
liczby całkowitej $q\ge1$ zachodzi
\[
 P_G(q)=\sum_{A\subseteq E}(-1)^{|A|}q^{k_G(A)}.
\]
Prawa strona jest wielomianem w $q$ o współczynnikach
całkowitych, wyznaczającym liczbę kolorowań dla każdego $q\ge1$.
\end{twierdzenie}

\begin{proof}
Wszystkich funkcji kolorujących wierzchołki jest $q^{|V|}$.
Dla każdej krawędzi $e=uv$ oznaczmy przez $B_e$
zbiór funkcji, które mają $c(u)=c(v)$.
Poprawne kolorowania to dokładnie funkcje nienależące
do żadnego z $B_e$. Z zasady włączeń i wyłączeń ich liczba
to $\sum_{A\subseteq E}(-1)^{|A|}\bigl|\bigcap_{e\in A}B_e\bigr|$.
Dla ustalonego $A$ warunki $c(u)=c(v)$ na krawędziach $A$
wymuszają stałość kolorowania na każdej składowej $(V,A)$.
Składowych jest $k_G(A)$ i każda może niezależnie otrzymać
jeden z $q$ kolorów. Dlatego
$\bigl|\bigcap_{e\in A}B_e\bigr|=q^{k_G(A)}$.
Podstawienie kończy dowód, również w przypadku pętli.
\end{proof}

\begin{twierdzenie}[Wielomian chromatyczny jako specjalizacja Tutte'a]
\label{thm:tg-chrom-tutte}
Niech $G=(V,E)$ będzie skończonym multigrafem o $n=|V|$
wierzchołkach i $c=k_G(E)$ składowych. Wtedy
\[
P_G(q)=(-1)^{n-c}q^c\,T_G(1-q,0)
\]
jako tożsamość wielomianów w $\mathbb Z[q]$.
\end{twierdzenie}

\begin{proof}
Podstawiając $x=1-q$, $y=0$ w definicji Tutte'a,
otrzymujemy
\[
(-1)^{n-c}q^c T_G(1-q,0)
=\sum_{A\subseteq E}
 (-1)^{n-c+k_G(A)-c+|A|-n+k_G(A)}q^{k_G(A)}.
\]
Wykładnik znaku $-1$ różni się od $|A|$ o liczbę parzystą
$2(k_G(A)-c)$, dlatego każdy składnik ma postać
$(-1)^{|A|}q^{k_G(A)}$. Poprzednie twierdzenie
identyfikuje sumę z $P_G(q)$.
\end{proof}

\begin{przyklad}[Kolorowania trójkąta]
Ponieważ $T_{C_3}(x,y)=x^2+x+y$, mamy
\[
P_{C_3}(q)=q\bigl((1-q)^2+(1-q)\bigr)
=q(q-1)(q-2).
\]
Odpowiada to bezpośredniemu rachunkowi: pierwszy
wierzchołek kolorujemy na $q$ sposobów, drugi
na $q-1$, a trzeci na $q-2$ sposoby.
W szczególności $\chi(C_3)=3$.
\end{przyklad}

\section{Dualność wielomianu Tutte'a}
\label{sec:tg-tutte-duality}

\begin{twierdzenie}[Dualność Tutte'a]
\label{thm:tg-tutte-dual}
Niech $G\subset S^2$ będzie skończonym spójnym grafem
płaskim, a $G^*$ jego dualnym multigrafem. Wtedy
\[
 T_{G^*}(x,y)=T_G(y,x).
\]
\end{twierdzenie}

\begin{proof}
Dla grafu płaskiego dualność zamienia usuwanie i kontrakcję
krawędzi, z zachowaniem naturalnego przyporządkowania
pozostałych krawędzi: jeśli $e$ nie jest pętlą, to
$(G/e)^*$ odpowiada $G^*\setminus e^*$, a jeśli $e$
nie jest mostem, to $(G\setminus e)^*$ odpowiada
$G^*/e^*$. Dla krawędzi będącej mostem $e^*$ jest
pętlą, natomiast dla pętli $e^*$ jest mostem.
Te tożsamości widać, śledząc dwie ściany po obu
stronach krawędzi: ich scalenie w dualu odpowiada
usunięciu krawędzi w grafie pierwotnym, zaś
scalenie końców krawędzi pierwotnej odpowiada
usunięciu jej krawędzi dualnej.

Stosujemy indukcję po liczbie krawędzi. Dla grafu
bez krawędzi $G=K_1$, a obie strony są równe jeden.
Jeśli $e$ jest zwykłą krawędzią (ani pętlą,
ani mostem), zastosowanie rekurencji do $G$ i $G^*$,
a następnie założenia indukcyjnego, daje
\[
\begin{aligned}
T_{G^*}(x,y)
 &=T_{G^*\setminus e^*}(x,y)+T_{G^*/e^*}(x,y)\\
 &=T_{(G/e)^*}(x,y)+T_{(G\setminus e)^*}(x,y)\\
 &=T_{G/e}(y,x)+T_{G\setminus e}(y,x)=T_G(y,x).
\end{aligned}
\]
Jeśli $e$ jest mostem, to $e^*$ jest pętlą i
$T_G(y,x)=yT_{G/e}(y,x)$, a
$T_{G^*}(x,y)=yT_{G^*\setminus e^*}(x,y)$.
Indukcja oraz $(G/e)^*=G^*\setminus e^*$ dają równość.
Przypadek pętli jest symetryczny. To kończy dowód.
\end{proof}

\begin{przyklad}[Trójkąt i trzy krawędzie równoległe]
Na sferze dualem trójkąta $C_3$ jest multigraf złożony
z dwóch wierzchołków połączonych trzema równoległymi
krawędziami. Zatem
$T_{C_3}=x^2+x+y$ oraz
$T_{C_3^*}=x+y+y^2=T_{C_3}(y,x)$.
Dualność zamieniła stopień cykliczny na strukturę
przekrojów krawędziowych.
\end{przyklad}

\section{Kolorowanie grafu dualnego a przepływy}
\label{sec:tg-flows-duality}

\begin{definicja}[Przepływ nigdzie niezerowy modulo $q$]
Niech $G=(V,E)$ będzie skończonym grafem z wybraną
orientacją krawędzi, a $q\ge2$ liczbą całkowitą.
\emph{Przepływem o wartościach w $\mathbb Z/q\mathbb Z$}
nazywamy wektor $j\in(\mathbb Z/q\mathbb Z)^E$
spełniający $Bj=0$. Jest on \emph{nigdzie niezerowy}, jeżeli
$j_e\ne0$ dla każdej krawędzi. Liczbę takich przepływów
oznaczamy przez $F_G(q)$.
\end{definicja}

\begin{twierdzenie}[Dualność kolorowań i przepływów]
\label{thm:tg-flow-color-dual}
Niech $G\subset S^2$ będzie skończonym spójnym grafem płaskim,
a $G^*$ jego grafem dualnym. Dla każdego $q\ge2$ zachodzi
\[
 P_{G^*}(q)=q\,F_G(q).
\]
\end{twierdzenie}

\begin{proof}
Ustalmy orientację sfery i pomocniczą orientację każdej
krawędzi $e$ grafu $G$. Niech $F_L(e)$ i $F_R(e)$ oznaczają
ściany po lewej i prawej stronie zorientowanej krawędzi.
Z kolorowania $c:F(G)\to\mathbb Z/q\mathbb Z$ tworzymy
wektor $j_e=c(F_R(e))-c(F_L(e))$.
Wokół każdego wierzchołka sumy przychodzących i wychodzących
wartości $j_e$ teleskopowo się znoszą, a więc $Bj=0$.
Ponadto $j_e\ne0$ dokładnie wtedy, gdy ściany po obu
stronach $e$ mają różne kolory, czyli gdy $c$ jest
poprawnym kolorowaniem grafu $G^*$.

Wykażemy, że każdy przepływ $j$ pochodzi z takiego
kolorowania. Wybieramy jedną ścianę i nadajemy jej
arbitralny kolor. Przechodząc w grafie dualnym przez
krawędź $e^*$, wyznaczamy kolor nowej ściany
z różnicy równej $j_e$, z odpowiednim znakiem.
Nie zależy to od wyboru drogi w grafie dualnym:
dla każdej zamkniętej ścieżki w $G^*$ suma odpowiednich
różnic jest zerowa. Rzeczywiście, prosty cykl dualny
rozdziela sferę na dwie części, a jego krawędzie
pierwotne tworzą przekrój w $G$. Sumując równania
Kirchhoffa $Bj=0$ po wierzchołkach jednej strony
tego przekroju, otrzymujemy dokładnie zerowanie
algebraicznej sumy przepływu przez cykl dualny.
Dowolny zamknięty spacer dualny rozkłada się na cykle
proste, a przejście przez dwie przeciwne orientacje tej
samej krawędzi wnosi sumę zero. Dla pętli dualnej,
która odpowiada mostowi pierwotnemu, każde rozwiązanie
$Bj=0$ ma na tym moście wartość zero (sumujemy równania
równowagi po jednej składowej po jego usunięciu).
Tym samym także w obecności mostów i krawędzi
równoległych kolor ściany nie zależy od wybranej trasy.

Dla ustalonego $j$ wszystkie kolorowania różnią się
o wspólną stałą dodaną do kolorów wszystkich ścian,
ponieważ $G^*$ jest spójny. Takich stałych jest dokładnie
$q$. Stąd $P_{G^*}(q)=qF_G(q)$.
\end{proof}

\begin{przyklad}[Przepływy na grafie theta]
Graf theta z trzema równoległymi krawędziami ma dualny
trójkąt. Dlatego
\[
 F_{\Theta}(q)=\frac{P_{C_3}(q)}{q}=(q-1)(q-2).
\]
Można to policzyć również bezpośrednio. Nadajmy wszystkim
trzem krawędziom orientację od $u$ do $v$.
Warunek przepływu to $j_1+j_2+j_3=0$ modulo $q$.
Wybieramy $j_1\ne0$ na $q-1$ sposobów, następnie
$j_2\ne0,-j_1$ na $q-2$ sposoby; wtedy
$j_3=-(j_1+j_2)$ jest jednoznacznie wyznaczone i niezerowe.
Zgodność z rachunkiem dualnym jest całkowita.
\end{przyklad}

\section*{Przejście do dalszej geometrii}
Wprowadziliśmy dwie różne operacje tworzenia obiektów
z grafu: samo nadanie długości krawędzi prowadziło
do przestrzeni metrycznej i laplasjanu, natomiast
zanurzenie komórkowe w powierzchni dodaje nowe dane
wymiaru dwa, opisane przez $\partial_2$ i przez porządki
rotacji. W następnych rozdziałach wykorzystamy ten punkt widzenia
w triangulacjach rozmaitości, grafach Reeba oraz dyskretnej
teorii Morse'a. Grafy plumbingowe pojawią się w dalszej części skryptu.

\section*{Wskazówki bibliograficzne do części planarnej}
Klasyczne twierdzenia o planarności, w tym pełne dowody
twierdzeń Kuratowskiego i Wagnera, można znaleźć w książce
R. Diestela, \emph{Graph Theory}. Szczegółowa analiza
grafów na powierzchniach, układów rotacji i grafów wstęgowych
znajduje się w monografii B. Mohara i C. Thomassena,
\emph{Graphs on Surfaces}. Wielomianowi Tutte'a oraz jego
związkom z kolorowaniami i przepływami poświęcone są m.in.
rozdziały \emph{Handbook of the Tutte Polynomial and Related Topics}.
Wszystkie te źródła traktujemy jako materiał uzupełniający;
definicje i rachunki używane w skrypcie pozostają podane
bezpośrednio w tekście.

% ================================================================
% ITERACJA 5. CZESC TOPOLOGICZNA
% ================================================================

\chapter{Triangulacje i grafy dualne kompleksów symplicjalnych}
\label{ch:tg-triangulacje}

Graf jest jednowymiarowym kompleksem symplicjalnym. Aby opisywać
powierzchnie i rozmaitości, potrzebujemy także dwuwymiarowych
trójkątów i sympleksów wyższych wymiarów. Zaczniemy od definicji,
a następnie sprawdzimy, jakie dodatkowe dane musi zawierać
triangulacja, żeby wyznaczać powierzchnię, a nie tylko jej graf.

\section{Kompleks symplicjalny i realizacja geometryczna}
\label{sec:tg-kompleks-symplicjalny}

\begin{definicja}[Abstrakcyjny kompleks symplicjalny]
Niech $V$ będzie skończonym zbiorem. \emph{Abstrakcyjnym
kompleksem symplicjalnym} na $V$ nazywamy rodzinę $K$
niepustych, skończonych podzbiorów $V$ taką, że jeżeli
$\sigma\in K$ i $\varnothing\ne\tau\subseteq\sigma$,
to $\tau\in K$. Zakładamy, że każdy $v\in V$ należy do $K$
jako jednoelementowy sympleks. Element $\sigma\in K$
o $p+1$ wierzchołkach nazywamy \emph{$p$-sympleksem}.
\emph{Wymiarem} $K$ jest
$\dim K:=\max\{|\sigma|-1:\sigma\in K\}$.
\end{definicja}

Wierzchołki są zatem sympleksami wymiaru zero, krawędzie
sympleksami wymiaru jeden, a trójkąty sympleksami wymiaru dwa.
Warunek domknięcia na niepuste podzbiory mówi, że jeśli
trójkąt należy do kompleksu, to należą do niego także jego
trzy krawędzie i trzy wierzchołki. Nie utożsamiamy
trójkąta z samym jego brzegiem.

\begin{definicja}[Realizacja geometryczna]
Niech $K$ będzie skończonym kompleksem na zbiorze $V$.
Jego \emph{realizacją geometryczną} nazywamy
\[
|K|=\left\{(t_v)_{v\in V}\in\mathbb R^V:\
t_v\ge0,\quad\sum_{v\in V}t_v=1,\quad
\{v:t_v>0\}\in K\right\}.
\]
Podzbiór, w którym $t_v=0$ dla $v\notin\sigma$, jest
geometrycznym sympleksem odpowiadającym $\sigma\in K$.
\end{definicja}

\begin{przyklad}[Trójkąt i jego brzeg]
Niech $V=\{a,b,c\}$. Kompleks zawierający wszystkie niepuste
podzbiory $V$ realizuje pełny trójkąt, homeomorficzny z dyskiem.
Jeżeli usuniemy jedyny dwusympleks $\{a,b,c\}$, pozostawiając
wierzchołki i trzy krawędzie, otrzymamy brzeg trójkąta,
homeomorficzny z $S^1$. Obie realizacje mają identyczny
jednowymiarowy szkielet, lecz różne grupy fundamentalne
i homologie. Sam graf nie odróżnia zatem tych obiektów.
\end{przyklad}

\begin{samepage}
\begin{center}
\begin{minipage}{.96\linewidth}\centering
\begin{tikzpicture}[scale=1.05,line cap=round]
\begin{scope}[shift={(-2.8,0)}]
 \coordinate (a) at (-1,-.7);\coordinate (b) at (1,-.7);
 \coordinate (c) at (0,1.05);
 \fill[tgBlue!15] (a)--(b)--(c)--cycle;
 \draw[tgEdge] (a)--(b)--(c)--cycle;
 \foreach \p in {a,b,c} \node[tgBlueVertex] at (\p) {};
 \node[text=tgNavy,font=\small] at (0,-1.25) {$|K|\cong D^2$};
\end{scope}
\begin{scope}[shift={(2.8,0)}]
 \coordinate (a) at (-1,-.7);\coordinate (b) at (1,-.7);
 \coordinate (c) at (0,1.05);
 \draw[tgAccent] (a)--(b)--(c)--cycle;
 \foreach \p in {a,b,c} \node[tgRedVertex] at (\p) {};
 \node[text=tgNavy,font=\small] at (0,-1.25) {$|\partial K|\cong S^1$};
\end{scope}
\end{tikzpicture}\par\smallskip

{\footnotesize Ten sam graf krawędziowy: po lewej z dołączoną
2-komórką, po prawej bez niej.}
\end{minipage}
\end{center}
\end{samepage}

\begin{definicja}[Szkielet i ogniwo]
Dla $p\ge0$ przez $K^{(p)}$ oznaczamy podkompleks
złożony ze wszystkich sympleksów wymiaru co najwyżej $p$.
Dla sympleksu $\sigma\in K$ definiujemy jego
\emph{ogniwo} (ang. \emph{link}) przez
$\operatorname{Lk}_K(\sigma)
=\{\tau\subseteq V\setminus\sigma:
\ \tau\cup\sigma\in K\}$.
W tej formule dopuszczamy $\tau=\varnothing$, choć
pusty sympleks nie był elementem rodziny $K$ w przyjętej
na początku konwencji. Nie dodaje on punktów
realizacji geometrycznej ogniwa. \emph{Gwiazdą domkniętą}
$\operatorname{St}_K(\sigma)$ nazywamy podkompleks
wszystkich ścian sympleksów zawierających $\sigma$.
\end{definicja}

Użycie sympleksu pustego w ogniwie jest wygodną konwencją
kombinatoryczną. W realizacji geometrycznej nie daje on
nowych punktów. Gwiazda wierzchołka jest stożkiem nad
jego ogniwem. Ten prosty fakt pozwoli rozpoznawać lokalną
topologię triangulowanych powierzchni.

\section{Triangulowana powierzchnia i warunek ogniwa}
\label{sec:tg-link-manifold}

\begin{definicja}[Triangulacja]
Niech $M$ będzie przestrzenią topologiczną. \emph{Triangulacją}
$M$ nazywamy parę $(K,h)$, gdzie $K$ jest skończonym
kompleksem symplicjalnym, a $h:|K|\to M$ homeomorfizmem.
Ponieważ $K$ jest skończony, taka triangulacja może istnieć
wyłącznie dla zwartej $M$; w dalszych twierdzeniach rozważamy
powierzchnie zwarte i triangulacje skończone. Gdy $M$ jest
powierzchnią, wierzchołki, krawędzie i trójkąty $K$ nazywamy
elementami jej triangulacji. Triangulację
nazywamy \emph{czystą dwuwymiarową}, gdy każda krawędź
i każdy wierzchołek są ścianami pewnego dwusympleksu.
\end{definicja}

Nie każdy skończony kompleks złożony z trójkątów jest
powierzchnią. Dwa trójkąty sklejone tylko w jednym
wierzchołku mają w tym punkcie otoczenie niehomeomorficzne
z dyskiem. Potrzebujemy zatem lokalnego kryterium.

\begin{twierdzenie}[Kryterium ogniwa dla powierzchni zamkniętej]
\label{thm:tg-link-surface}
Niech $K$ będzie skończonym, spójnym, czystym
kompleksem symplicjalnym wymiaru dwa. Następujące warunki
są równoważne:
\begin{enumerate}[label=(\roman*)]
\item $|K|$ jest zamkniętą powierzchnią topologiczną;
\item dla każdego wierzchołka $v\in K$ realizacja
$|\operatorname{Lk}_K(v)|$ jest okręgiem.
\end{enumerate}
W szczególności, jeżeli zachodzą te warunki, każda
krawędź należy do dokładnie dwóch trójkątów.
\end{twierdzenie}

\begin{proof}
Załóżmy najpierw (ii). Gwiazda każdego wierzchołka jest
stożkiem nad okręgiem, a więc jest dyskiem. Wnętrza trójkątów
mają otoczenia dyskowe wprost z geometrii sympleksu.
Każda krawędź należy do dwóch trójkątów: w ogniwie jednego
jej końca odpowiada jej wierzchołek, który w cyklu ma
dokładnie dwóch sąsiadów. Zatem również punkt we wnętrzu
krawędzi ma otoczenie powstałe z dwóch półdysków i jest
punktem wewnętrznym powierzchni. Wszystkie punkty mają
więc otoczenia homeomorficzne z $\mathbb R^2$. Ponieważ
$|K|$ jest zwarty, spójny i Hausdorffa, jest zamkniętą
powierzchnią.

W drugą stronę załóżmy (i). Dla punktu $x$ we wnętrzu
krawędzi $e$ małe otoczenie w $|K|$ jest sklejeniem
$k$ półdysków wzdłuż ich średnic, gdzie $k$ jest liczbą
trójkątów zawierających $e$. Jego ogniwo ma $k$ punktów,
a ogniwo punktu w powierzchni bez brzegu ma dokładnie
dwa punkty. Zatem $k=2$: każda krawędź przylega do
dokładnie dwóch trójkątów. Link dowolnego
wierzchołka jest skończonym grafem, którego każdy
wierzchołek ma stopień dwa, a zatem jest sumą rozłączną
cykli. Gdyby składał się z kilku cykli, małe otoczenie
wierzchołka po usunięciu tego wierzchołka miałoby kilka
składowych, podczas gdy przebity dysk jest spójny.
Zatem ogniwo jest jednym cyklem. W dowodzie korzystamy
z faktu, że stożek nad ogniwem opisuje małe otoczenie
wierzchołka w realizacji kompleksu.
\end{proof}

\begin{przyklad}[Brzeg czworościanu]
Niech $K$ składa się z czterech trójkątów, sześciu krawędzi
i czterech wierzchołków tworzących brzeg czworościanu.
Każde ogniwo wierzchołka jest trójkątnym cyklem, więc
$|K|$ jest zamkniętą powierzchnią. Mamy
$\chi(|K|)=4-6+4=2$, a ponieważ $|K|$ jest brzegiem
wypukłego wielościanu, jest homeomorficzny z $S^2$.
Nie należy mylić liczby trójkątów $4$ z liczbą ścian
grafowego zanurzenia: tu są to komórki wymiaru dwa
konkretnego kompleksu.
\end{przyklad}

\begin{samepage}
\begin{center}
\begin{minipage}{.96\linewidth}\centering
\begin{tikzpicture}[scale=1.05,line join=round,line cap=round]
  \coordinate (a) at (-1.3,-.8);\coordinate (b) at (1.3,-.8);
  \coordinate (c) at (.35,1.4);\coordinate (d) at (-.22,.13);
  \fill[tgBlue!11] (a)--(b)--(c)--cycle;
  \fill[tgCoral!18] (a)--(b)--(d)--cycle;
  \fill[tgTeal!18] (a)--(c)--(d)--cycle;
  \fill[tgGold!20] (b)--(c)--(d)--cycle;
  \draw[tgEdge] (a)--(b)--(c)--cycle;
  \draw[tgEdge] (a)--(d)--(b) (c)--(d);
  \foreach \p in {a,b,c,d} \node[tgBlueVertex] at (\p) {};
  \node[font=\small,text=tgNavy] at (0,-1.38) {$V=4,\ E=6,\ F=4,\ \chi=2$};
\end{tikzpicture}\par\smallskip

{\footnotesize Schemat kombinatoryczny powierzchni czworościanu;
kolory oznaczają cztery trójkąty.}
\end{minipage}
\end{center}
\end{samepage}

\section{Łańcuchy symplicjalne i związek z homologią}
\label{sec:tg-simplicial-chain}

\begin{definicja}[Operator brzegu]
Dla skończonego kompleksu symplicjalnego $K$ niech
$C_p(K;\mathbb Z)$ będzie wolną grupą abelową generowaną
przez zorientowane $p$-sympleksy, z relacją
$[v_{\pi(0)},\ldots,v_{\pi(p)}]
=\operatorname{sgn}(\pi)[v_0,\ldots,v_p]$.
Dla $p\ge1$ definiujemy
\[
\partial_p[v_0,\ldots,v_p]
=\sum_{i=0}^p(-1)^i
[v_0,\ldots,\widehat v_i,\ldots,v_p],
\]
a $\partial_0=0$. Homologie symplicjalne są
$H_p(K;\mathbb Z):=\ker\partial_p/\operatorname{im}\partial_{p+1}$.
\end{definicja}

\begin{stwierdzenie}[Brzeg brzegu]
Dla każdego skończonego kompleksu symplicjalnego $K$
i każdego $p\ge2$ zachodzi $\partial_{p-1}\partial_p=0$.
\end{stwierdzenie}

\begin{proof}
W rozwinięciu $\partial_{p-1}\partial_p$ każda ściana
powstała przez usunięcie dwóch różnych wierzchołków
pojawia się dwukrotnie. Jeśli najpierw usuniemy wierzchołek
o indeksie $i$, a następnie o indeksie $j>i$, znak jest
$(-1)^i(-1)^{j-1}$. Przy odwrotnej kolejności znak jest
$(-1)^j(-1)^i$. Ich suma wynosi zero. Sumując po wszystkich
parach $i<j$, otrzymujemy równość operatorową.
\end{proof}

\begin{przyklad}[Co robi dołączenie trójkąta]
Dla brzegu trójkąta o wierzchołkach $0,1,2$ jedynym
niezależnym cyklem jednowymiarowym jest
$z=[0,1]+[1,2]-[0,2]$, zatem
$H_1(\partial\Delta^2;\mathbb Z)\cong\mathbb Z$.
Po dołączeniu trójkąta $[0,1,2]$ otrzymujemy
$z=\partial_2[0,1,2]$, więc klasa $z$ zanika i
$H_1(\Delta^2;\mathbb Z)=0$. Jednocześnie $H_0\cong\mathbb Z$
w obu przypadkach. To bezpośredni rachunek ilustrujący
różnicę między grafem i wypełnionym kompleksem.
\end{przyklad}

\begin{stwierdzenie}[Charakterystyka Eulera z homologii]
Niech $K$ będzie skończonym kompleksem symplicjalnym,
a $f_p$ liczbą jego $p$-sympleksów. Jeżeli
$b_p=\dim_{\mathbb R}H_p(K;\mathbb R)$, to
\[
\sum_{p\ge0}(-1)^pf_p=\sum_{p\ge0}(-1)^pb_p.
\]
\end{stwierdzenie}

\begin{proof}
Niech $r_p$ oznacza rząd operatora
$\partial_p:C_p(K;\mathbb R)\to C_{p-1}(K;\mathbb R)$;
przyjmujemy $r_0=0$. Z twierdzenia o rzędzie i jądrze
oraz definicji homologii mamy
$b_p=(f_p-r_p)-r_{p+1}$, a więc
$f_p=b_p+r_p+r_{p+1}$. Po wymnożeniu przez $(-1)^p$
i zsumowaniu wkłady $r_p$ znoszą się parami, pozostawiając
sumę liczb Bettiego. Wszystkie sumy są skończone.
\end{proof}

\section{Graf dualny triangulacji i dane sklejania}
\label{sec:tg-dual-triangulation}

\begin{definicja}[Graf dualny triangulacji powierzchni]
Niech $K$ będzie skończoną triangulacją spójnej,
zamkniętej powierzchni. \emph{Graf dualny triangulacji}
ma po jednym wierzchołku dla każdego trójkąta $K$
i po jednej krawędzi dla każdej krawędzi $K$, łącząc
wierzchołki dwóch trójkątów incydentnych z daną krawędzią.
\end{definicja}

\begin{stwierdzenie}
Dla triangulacji powierzchni zamkniętej graf dualny jest
spójny i trójregularny, czyli każdy jego wierzchołek ma
stopień trzy (z krotnościami). Jeżeli triangulacja ma
$f_2$ trójkątów i $f_1$ krawędzi, to
$3f_2=2f_1$.
\end{stwierdzenie}

\begin{proof}
Każdy trójkąt ma trzy boki, a każda krawędź zamkniętej
powierzchni jest wspólna dokładnie dwóm trójkątom.
Stąd stopień dualnego wierzchołka jest równy trzy,
a podwójne zliczenie par (trójkąt, jego bok) daje
$3f_2=2f_1$. Spójność wynika stąd, że dwa punkty we
wnętrzach dowolnych trójkątów można połączyć ścieżką
na spójnej powierzchni, omijającą wierzchołki triangulacji
i przechodzącą poprzecznie przez krawędzie. Kolejne
napotkane trójkąty wyznaczają ścieżkę w grafie dualnym.
\end{proof}

\begin{uwaga}[Dlaczego sam graf dualny nie wystarczy]
Graf dualny mówi, które trójkąty sąsiadują, ale nie pamięta,
\emph{które ich boki} są sklejane, w jakiej kolejności i czy
sklejenie odwraca orientację. Rozkłady powierzchni za pomocą
trójkątów z identyfikowanymi bokami bywają wygodniej opisane
przez \emph{kompleksy typu $\Delta$}, które nie muszą być
abstrakcyjnymi kompleksami symplicjalnymi. Dla przykładu
sferę można uzyskać przez sklejenie dwóch trójkątów wzdłuż
odpowiadających boków, a torus też ma model z dwóch
trójkątów z odpowiednio identyfikowanymi bokami. W obu
modelach graf dualny jest grafem dwóch wierzchołków
i trzech krawędzi równoległych, mimo że
$\chi(S^2)=2$, a $\chi(T^2)=0$. Rozróżniają je dopiero
dane sklejeń; są to modele typu $\Delta$, nie simplicjalne
triangulacje z dwiema dwukomórkami.
\end{uwaga}

\section{Ruchy Pachnera na triangulacjach powierzchni}
\label{sec:tg-pachner}

\begin{definicja}[Lokalne ruchy bistellarne w wymiarze dwa]
Niech $K$ będzie triangulacją zamkniętej powierzchni.
Ruch $1\leftrightarrow3$ zastępuje pojedynczy trójkąt
trzema trójkątami przez wprowadzenie nowego wierzchołka
w jego wnętrzu (lub wykonuje operację odwrotną).
Ruch $2\leftrightarrow2$ usuwa wspólną przekątną dwóch
różnych trójkątów, których suma jest podkompleksem
homeomorficznym z domkniętym dyskiem o brzegu złożonym
z czterech różnych krawędzi, po czym wprowadza drugą
przekątną tego czworokąta. Zakładamy, że ruch nie tworzy
niedozwolonego wielokrotnego sympleksu, a wynik jest
ponownie kompleksem symplicjalnym triangulującym tę samą
powierzchnię. W szczególności w ruchu $2\to2$ nowa
przekątna nie może już należeć do kompleksu poza
rozważanym dyskiem.
\end{definicja}

\begin{samepage}
\begin{center}
\begin{minipage}{.96\linewidth}\centering
\begin{tikzpicture}[scale=.94,line cap=round,line join=round,>=Stealth]
 \begin{scope}[shift={(-3,0)}]
  \coordinate (a) at (-1,-.7);\coordinate (b) at (1,-.7);
  \coordinate (c) at (0,1.0);\coordinate (o) at (0,-.1);
  \fill[tgBlue!12] (a)--(b)--(c)--cycle;
  \draw[tgEdge] (a)--(b)--(c)--cycle;
  \node[tgBlueVertex] at (a) {};\node[tgBlueVertex] at (b) {};\node[tgBlueVertex] at (c) {};
  \node[text=tgNavy,font=\small] at (0,-1.2) {$1$ trójkąt};
 \end{scope}
 \draw[tgTeal,-{Stealth[length=3mm]},line width=1.5pt] (-1.45,0)--(-.75,0);
 \begin{scope}[shift={(1,0)}]
  \coordinate (a) at (-1,-.7);\coordinate (b) at (1,-.7);
  \coordinate (c) at (0,1.0);\coordinate (o) at (0,-.1);
  \fill[tgTeal!12] (a)--(b)--(o)--cycle;
  \fill[tgBlue!12] (a)--(c)--(o)--cycle;
  \fill[tgGold!16] (b)--(c)--(o)--cycle;
  \draw[tgEdge] (a)--(b)--(c)--cycle;
  \draw[tgAccent] (a)--(o)--(b) (c)--(o);
  \foreach \p in {a,b,c,o} \node[tgBlueVertex] at (\p) {};
  \node[text=tgNavy,font=\small] at (0,-1.2) {$3$ trójkąty};
 \end{scope}
\end{tikzpicture}\par\smallskip

{\footnotesize Ruch $1\to3$: powierzchnia i jej brzeg pozostają te same;
zmienia się tylko podział symplicjalny dysku.}
\end{minipage}
\end{center}
\end{samepage}

\begin{samepage}
\begin{center}
\begin{minipage}{.96\linewidth}\centering
\begin{tikzpicture}[scale=.84,line cap=round,line join=round,>=Stealth]
 \begin{scope}[shift={(-2.5,0)}]
  \coordinate (a) at (-1.25,-.85);\coordinate (b) at (1.25,-.85);
  \coordinate (c) at (1.05,.85);\coordinate (d) at (-1.05,.85);
  \fill[tgTeal!14] (a)--(b)--(c)--cycle;
  \fill[tgBlue!14] (a)--(c)--(d)--cycle;
  \draw[tgEdge] (a)--(b)--(c)--(d)--cycle;
  \draw[tgAccent] (a)--(c);
  \foreach \p in {a,b,c,d} \node[tgBlueVertex] at (\p) {};
  \node[font=\small,text=tgNavy] at (0,-1.30) {przekątna $ac$};
 \end{scope}
 \draw[tgTeal,-{Stealth[length=3mm]},line width=1.5pt] (-.64,0)--(.64,0);
 \begin{scope}[shift={(2.5,0)}]
  \coordinate (a) at (-1.25,-.85);\coordinate (b) at (1.25,-.85);
  \coordinate (c) at (1.05,.85);\coordinate (d) at (-1.05,.85);
  \fill[tgGold!15] (a)--(b)--(d)--cycle;
  \fill[tgCoral!12] (b)--(c)--(d)--cycle;
  \draw[tgEdge] (a)--(b)--(c)--(d)--cycle;
  \draw[tgAccent] (b)--(d);
  \foreach \p in {a,b,c,d} \node[tgRedVertex] at (\p) {};
  \node[font=\small,text=tgNavy] at (0,-1.30) {przekątna $bd$};
 \end{scope}
\end{tikzpicture}\par\smallskip

{\footnotesize Ruch $2\leftrightarrow2$: dwie triangulacje tego samego
czworokąta, z niezmienionym brzegiem.}
\end{minipage}
\end{center}
\end{samepage}

\begin{stwierdzenie}[Niezmienniczość topologii i charakterystyki Eulera]
Każdy dopuszczalny ruch $1\leftrightarrow3$ i
$2\leftrightarrow2$ zachowuje typ homeomorficzny
triangulowanej powierzchni oraz jej charakterystykę Eulera.
\end{stwierdzenie}

\begin{proof}
Obie operacje odbywają się wewnątrz domkniętego dysku
triangulacji i nie zmieniają sposobu sklejenia jego brzegu
z resztą powierzchni. Ruchy zastępują dysk inną
triangulacją tego samego dysku, więc tożsamość na brzegu
przedłuża się do homeomorfizmu całej powierzchni.
W ruchu $1\to3$ przybywa jeden wierzchołek, trzy krawędzie
i dwie ściany dwuwymiarowe netto. Zmiana $f_0-f_1+f_2$
wynosi $1-3+2=0$. Ruch $2\to2$ nie zmienia żadnej z tych
liczb; odwrotne ruchy zachowują je również.
\end{proof}

\begin{stwierdzenie}[Subdywizja krawędzi przez dwa ruchy]
Niech $e=ab$ będzie krawędzią triangulacji zamkniętej
powierzchni, przylegającą do dwóch trójkątów $abc$
i $abd$ o różnych wierzchołkach $c,d$. Wtedy podział krawędzi $ab$ przez nowy wierzchołek
jest kombinatorycznie równoważny wykonaniu jednego
ruchu $1\to3$ oraz jednego ruchu $2\to2$.
\end{stwierdzenie}

\begin{proof}
Wstawiamy nowy wierzchołek $x$ do wnętrza trójkąta $abc$.
Po ruchu $1\to3$ trójkąty tej części to $abx$, $bcx$, $cax$;
trójkąt $abd$ pozostaje bez zmian. Wspólna krawędź $ab$
trójkątów $abx$ i $abd$ ma teraz przeciwległe wierzchołki
$x,d$. Wykonujemy ruch $2\to2$, zamieniając przekątną
$ab$ na $xd$. Powstają trójkąty $adx$, $bdx$, $bcx$, $cax$.
Są one dokładnie czterema trójkątami otrzymanymi przez
podział dwóch starych trójkątów wzdłuż krawędzi $ab$.
Zależność od położenia geometrycznego $x$ jest pozorna:
liczy się kombinatoryka kompleksu, a oba podziały
dysku są homeomorficzne względnie do jego brzegu.
\end{proof}

\begin{twierdzenie}[Pachner; wynik przyjmowany]
Niech $K,L$ będą dwiema skończonymi triangulacjami
tej samej zwartej, zamkniętej rozmaitości PL
wymiaru dwa (lub, równoważnie w tym wymiarze,
dwiema PL-równoważnymi triangulacjami powierzchni).
Wtedy istnieje skończony ciąg dopuszczalnych ruchów
$1\leftrightarrow3$ i $2\leftrightarrow2$ prowadzący
od $K$ do kompleksu izomorficznego z $L$.
\end{twierdzenie}

\begin{uwaga}[Status dowodu]
Udowodniliśmy lokalną poprawność ruchów i ich wpływ na
$\chi$, ale twierdzenie o istnieniu ciągu łączącego
\emph{dowolne} dwie triangulacje jest znacznie silniejsze.
Wymaga teorii podziałów gwiaździstych i wspólnych
podziałów PL (lub pełnego argumentu Pachnera).
Przyjmujemy je jako klasyczny wynik i nie podszywamy
się pod jego kompletny dowód.
\end{uwaga}

\chapter{Grafy Reeba i funkcje Morse'a na powierzchniach}
\label{ch:tg-reeb}

\section{Od poziomic funkcji do grafu}
\label{sec:tg-reeb-def}

\begin{definicja}[Graf Reeba]
Niech $M$ będzie zwartą, spójną rozmaitością gładką,
a $f:M\to\mathbb R$ funkcją gładką. Na $M$ definiujemy
$x\sim_f y$ wtedy i tylko wtedy, gdy $f(x)=f(y)$
oraz $x,y$ należą do tej samej spójnej składowej
zbioru poziomicowego $f^{-1}(f(x))$.
Przestrzeń ilorazową $R_f:=M/{\sim_f}$ nazywamy
\emph{przestrzenią Reeba} funkcji $f$. Odwzorowanie
ilorazowe oznaczamy $q_f:M\to R_f$.
Funkcja $f$ indukuje ciągłą funkcję
$\bar f:R_f\to\mathbb R$ spełniającą
$f=\bar f\circ q_f$.
\end{definicja}

Nie każda przestrzeń Reeba funkcji gładkiej jest
skończonym grafem. Najpierw przyjmiemy warunki,
które pozwalają udowodnić tę własność i opisać
stopnie jej wierzchołków.

\begin{definicja}[Prosta funkcja Morse'a na powierzchni]
Niech $M$ będzie zwartą powierzchnią bez brzegu.
Funkcję $f:M\to\mathbb R$ nazywamy \emph{funkcją Morse'a},
gdy jej wszystkie punkty krytyczne są niezdegenerowane.
Nazywamy ją \emph{prostą}, jeżeli dodatkowo różne punkty
krytyczne mają różne wartości. Indeks punktu krytycznego
jest liczbą ujemnych wartości własnych hesjanu w tym punkcie.
Na powierzchni wynosi on $0$, $1$ lub $2$.
\end{definicja}

\begin{stwierdzenie}[Regularne przedziały poziomów]
Niech $M$ będzie zwartą rozmaitością bez brzegu,
a $f:M\to\mathbb R$ gładką funkcją. Jeśli
$[a,b]$ nie zawiera wartości krytycznych i
$f^{-1}([a,b])\ne\varnothing$, to istnieje
homeomorfizm, a nawet dyfeomorfizm
$f^{-1}([a,b])\cong f^{-1}(a)\times[a,b]$
zachowujący wartości funkcji $f$ jako drugą współrzędną.
\end{stwierdzenie}

\begin{proof}
Wybieramy metrykę riemannowską na $M$ i rozważamy
pole $X:=\nabla f/\|\nabla f\|^2$ na otoczeniu
$f^{-1}([a,b])$. Mianownik nigdzie tam nie znika.
Dla krzywej całkowej $\gamma$ pola $X$ mamy
$\frac{\mathrm d}{\mathrm dt}(f\circ\gamma)
=\mathrm df(X)=1$.
A zatem przepływ $\Phi_t$ przesuwa wartość funkcji
o dokładnie $t$. Zwarta warstwa $f^{-1}([a,b])$
zapewnia istnienie wszystkich odcinków przepływu
potrzebnych do przejścia między poziomami $a$ i $b$.
Odwzorowanie $(x,t)\mapsto\Phi_{t-a}(x)$,
$x\in f^{-1}(a)$, ma odwrotność
$y\mapsto(\Phi_{a-f(y)}(y),f(y))$.
Gładkość obu odwzorowań wynika z gładkiej zależności
przepływu od warunków początkowych.
\end{proof}

\begin{twierdzenie}[Skończoność i lokalna struktura grafu Reeba]
\label{thm:tg-reeb-finite}
Niech $M$ będzie zwartą, spójną, orientowalną
powierzchnią bez brzegu, a $f:M\to\mathbb R$
prostą funkcją Morse'a. Wtedy $R_f$ jest
homeomorficzna ze skończonym, spójnym grafem.
Przy kanonicznym wyborze wierzchołków nad punktami
krytycznymi wierzchołek odpowiadający minimum lub
maksimum ma stopień jeden, a odpowiadający punktowi
siodłowemu ma stopień trzy. Wnętrza krawędzi
odpowiadają rodzinom spójnych składowych regularnych
poziomic, które są okręgami.
\end{twierdzenie}

\begin{proof}
Punkty krytyczne funkcji Morse'a są izolowane;
ponieważ $M$ jest zwarta, jest ich skończenie wiele.
Wybieramy regularne liczby
$a_0<a_1<\cdots<a_N$ rozdzielające kolejne wartości
krytyczne. Na każdym zamkniętym przedziale regularnym
poprzednie stwierdzenie identyfikuje podzbiór powierzchni
z produktem poziomicy i przedziału. Regularna poziomica
na zwartej powierzchni bez brzegu jest skończoną sumą
okręgów. Po zidentyfikowaniu każdej składowej okręgu
z punktem produkt przechodzi w skończoną sumę odcinków.

W otoczeniu punktu krytycznego lemat Morse'a daje lokalne
modele $f(p)+x^2+y^2$, $f(p)-x^2+y^2$ i
$f(p)-x^2-y^2$. Aby przejść od lokalnego modelu do
informacji o spójnych składowych całej poziomicy,
wybieramy $\varepsilon>0$ tak małe, by przedział
$[f(p)-\varepsilon,f(p)+\varepsilon]$ nie zawierał
innej wartości krytycznej. Składowa tego pasa zawierająca
$p$ jest przy ekstremum dyskiem, a przy siodle
powierzchnią typu pary spodni: ma trzy brzegowe okręgi,
z których jeden leży po jednej stronie poziomu
krytycznego, dwa po drugiej. Pozostałe składowe pasa
są cylindrami, bo nie zawierają punktów krytycznych.
Po zredukowaniu spójnych składowych poziomic do punktów
dysk daje wierzchołek stopnia jeden, para spodni
wierzchołek stopnia trzy, a cylindry odcinki
przechodzące przez rozważany przedział bez nowych
rozgałęzień. Użyliśmy tu orientowalności powierzchni
oraz tego, że poziom krytyczny zawiera dokładnie jeden
punkt krytyczny.
Ponieważ mamy skończenie wiele krytycznych poziomów,
po sklejeniu skończenie wielu odcinków i wierzchołków
uzyskujemy skończony graf. Jego spójność wynika z
ciągłości i surjektywności odwzorowania ilorazowego $q_f$.
\end{proof}

\begin{samepage}
\begin{center}
\begin{minipage}{.96\linewidth}\centering
\begin{tikzpicture}[scale=1,>=Stealth,line cap=round]
 \begin{scope}[shift={(-2.6,0)}]
  \draw[tgEdge] (0,-1.25)--(0,1.25);
  \node[tgBlueVertex] at (0,-1.25) {};
  \node[tgRedVertex] at (0,1.25) {};
  \node[text=tgNavy,font=\small] at (0,-1.65) {sfera: odcinek};
 \end{scope}
 \begin{scope}[shift={(2.5,0)}]
  \node[tgBlueVertex] (b) at (0,-1.25) {};
  \node[tgTealVertex] (s1) at (0,-.35) {};
  \node[tgTealVertex] (s2) at (0,.65) {};
  \node[tgRedVertex] (t) at (0,1.4) {};
  \draw[tgEdge] (b)--(s1) (s2)--(t);
  \draw[tgBlue,line width=2pt] (s1) .. controls (-1.10,-.35) and (-1.10,.65) .. (s2);
  \draw[tgCoral,line width=2pt] (s1) .. controls (1.10,-.35) and (1.10,.65) .. (s2);
  \node[text=tgNavy,font=\small] at (0,-1.65) {torus: jeden cykl};
 \end{scope}
\end{tikzpicture}\par\smallskip

{\footnotesize Schemat grafów Reeba standardowych funkcji wysokości;
nie są to rysunki poziomic w samej powierzchni.}
\end{minipage}
\end{center}
\end{samepage}

\section{Rodzaj powierzchni zapisany w grafie Reeba}
\label{sec:tg-reeb-genus}

\begin{uwaga}[Znaczenie prostoty funkcji Morse'a]
W wynikach o grafach Reeba zakładamy, że każda wartość
krytyczna odpowiada dokładnie jednemu punktowi krytycznemu.
Bez tego założenia kilka lokalnych zmian zbioru poziomicowego
może zachodzić jednocześnie; graf Reeba nadal bywa skończony,
ale opis stopni jego wierzchołków wymaga wtedy dodatkowej
analizy. Wzoru na liczbę cykli poniżej nie należy przenosić
bez uzasadnienia na funkcje zdegenerowane ani na powierzchnie
z brzegiem.
\end{uwaga}

\begin{twierdzenie}[Pierwsza liczba Bettiego grafu Reeba]
\label{thm:tg-reeb-betti}
Niech $M=\Sigma_g$ będzie zwartą, spójną, orientowalną
powierzchnią bez brzegu rodzaju $g\ge0$, a
$f:M\to\mathbb R$ prostą funkcją Morse'a. Wtedy
\[
\dim H_1(R_f;\mathbb R)=g.
\]
W szczególności graf Reeba na sferze jest drzewem,
a na torusie ma dokładnie jeden niezależny cykl.
\end{twierdzenie}

\begin{proof}
Niech $c_i$ oznacza liczbę punktów krytycznych indeksu
$i\in\{0,1,2\}$. Z poprzedniego twierdzenia $R_f$ ma
$c_0+c_2$ wierzchołków stopnia jeden oraz $c_1$
wierzchołków stopnia trzy. Jeśli $v$ i $e$ oznaczają
liczbę jego wierzchołków i krawędzi, to
$v=c_0+c_1+c_2$ i
$2e=(c_0+c_2)+3c_1$.
Dla grafu spójnego otrzymujemy więc
\[
\dim H_1(R_f;\mathbb R)=e-v+1
=1+\frac{c_1-c_0-c_2}{2}.
\]
Z klasycznej teorii Morse'a dla powierzchni (albo
z rozkładu na uchwyty podpoziomic) wiemy, że
$\chi(\Sigma_g)=c_0-c_1+c_2=2-2g$:
minimum dodaje dysk, siodło uchwyt jednowymiarowy,
a maksimum dysk zamykający składową brzegu, co zmienia
charakterystykę Eulera odpowiednio o $+1,-1,+1$.
Zatem $c_1-c_0-c_2=2g-2$. Po podstawieniu uzyskujemy
$\dim H_1(R_f;\mathbb R)=g$, jak należało dowieść.
\end{proof}

\begin{przyklad}[Sfera i torus]
Standardowa funkcja wysokości na sferze ma jedno minimum
i jedno maksimum, a więc $(c_0,c_1,c_2)=(1,0,1)$.
Graf Reeba ma dwa wierzchołki i jedną krawędź,
zatem $\beta_1=0$. Dla torusa ustawionego pionowo
standardowa funkcja wysokości ma jedno minimum,
dwa siodła i jedno maksimum, czyli $(1,2,1)$.
Wtedy $v=4$, $e=4$ i $\beta_1=4-4+1=1$.
Dodawanie par krytycznych punktów podczas deformacji
funkcji może zwiększyć liczby $c_i$, ale równość
$c_0-c_1+c_2=2-2g$ i wartość $\beta_1=g$ pozostają.
\end{przyklad}

\begin{uwaga}[Zakres twierdzenia]
Ostatnia równość jest szczególną własnością prostych
funkcji Morse'a na \emph{orientowalnych powierzchniach
zamkniętych}. Dla rozmaitości wyższego wymiaru
lub funkcji o zdegenerowanych krytycznych poziomach
graf Reeba może mieć inną kombinatorykę, a powyższe
liczenie stopni wierzchołków nie ma zastosowania.
\end{uwaga}

\section{Graf Reeba a rozkład na uchwyty}
\label{sec:tg-reeb-handles}

\begin{przyklad}[Siodło łączy lub rozdziela]
Niech $a<b$ będą regularnymi wartościami funkcji
Morse'a $f:\Sigma_g\to\mathbb R$ i niech przedział
$(a,b)$ zawiera dokładnie jeden punkt krytyczny indeksu
jeden. Podpoziomica $f^{-1}(( -\infty,a])$ po przekroczeniu
siodła otrzymuje jeden uchwyt jednowymiarowy
$D^1\times D^1$, przyczepiony wzdłuż
$\partial D^1\times D^1$. Jeżeli jego końce są przyczepione
do dwóch różnych składowych brzegowych podpoziomicy,
dwie składowe regularnego poziomu łączą się w jedną.
Jeżeli końce leżą na tej samej składowej brzegu,
powstają dwie składowe poziomu. Obie operacje mają
w grafie Reeba lokalny obraz wierzchołka stopnia trzy.
\end{przyklad}

\begin{samepage}
\begin{center}
\begin{minipage}{.96\linewidth}\centering
\begin{tikzpicture}[scale=.98,>=Stealth,line cap=round,line join=round]
 \begin{scope}[shift={(-2.6,0)}]
  \draw[tgBlue,line width=2pt] (-1,-.8) .. controls (-.8,-.1) and (-.6,.25) .. (0,1);
  \draw[tgTeal,line width=2pt] (1,-.8) .. controls (.8,-.1) and (.6,.25) .. (0,1);
  \node[tgRedVertex] at (0,1) {};
  \node[text=tgNavy,font=\small] at (0,-1.25) {połączenie: $2\to1$};
 \end{scope}
 \begin{scope}[shift={(2.6,0)}]
  \draw[tgBlue,line width=2pt] (0,-.8) .. controls (-.6,-.05) and (-.8,.45) .. (-1,1);
  \draw[tgTeal,line width=2pt] (0,-.8) .. controls (.6,-.05) and (.8,.45) .. (1,1);
  \node[tgRedVertex] at (0,-.8) {};
  \node[text=tgNavy,font=\small] at (0,-1.25) {rozdzielenie: $1\to2$};
 \end{scope}
\end{tikzpicture}\par\smallskip

{\footnotesize Lokalny obraz wierzchołka siodłowego
w grafie Reeba, zorientowany wzrostem wartości funkcji.}
\end{minipage}
\end{center}
\end{samepage}

\begin{uwaga}[Czego graf Reeba nie zachowuje]
Iloraz $R_f$ zapamiętuje liczbę i sposób łączenia
spójnych składowych poziomic, ale nie pełną geometrię
metryczną powierzchni. Dwie różne funkcje Morse'a
mogą mieć izomorficzne grafy Reeba; dodatkowe dane,
na przykład wartości krytyczne i rodzaje składowych
podpoziomic, pozwalają na dokładniejszą klasyfikację.
Na powierzchni orientowalnej sam pierwszy wymiar homologii
grafu wyznacza rodzaj powierzchni na mocy
Twierdzenia~\ref{thm:tg-reeb-betti}.
\end{uwaga}

\chapter{Dyskretna teoria Morse'a Formana}
\label{ch:tg-forman}

W klasycznej teorii Morse'a topologia rozmaitości jest
odczytywana z punktów krytycznych funkcji gładkiej.
W kompleksie symplicjalnym nie mamy pochodnych,
ale dysponujemy relacją bycia ścianą innego sympleksu.
Dyskretna teoria Morse'a wykorzystuje \emph{acykliczne
parowania komórek} do usuwania zbędnych par sympleksów.
Przedstawimy najpierw tę konstrukcję, a następnie
wyjaśnimy dokładnie, które wnioski da się udowodnić
samą algebrą kompleksów łańcuchowych.

\section{Relacja incydencji i parowanie komórek}
\label{sec:tg-forman-matching}

\begin{definicja}[Pokrycie w posetcie ścian]
Niech $K$ będzie skończonym kompleksem symplicjalnym.
Mówimy, że sympleks $\beta$ \emph{pokrywa} sympleks
$\alpha$, i piszemy $\alpha\prec\beta$, gdy
$\alpha\subset\beta$ oraz
$\dim\beta=\dim\alpha+1$. \emph{Graf Hassego}
relacji ścian ma wierzchołki odpowiadające niepustym
sympleksom i krawędź między $\alpha$ i $\beta$
dokładnie wtedy, gdy $\alpha\prec\beta$.
\end{definicja}

\begin{definicja}[Dyskretne pole gradientowe]
\label{def:tg-discrete-morse-field}
Niech $K$ będzie skończonym kompleksem symplicjalnym.
\emph{Parowaniem} $\mathcal V$ nazywamy rodzinę par
$(\alpha^p,\beta^{p+1})$, w których
$\alpha\prec\beta$ i żaden sympleks nie występuje
w więcej niż jednej parze. Sympleks niewystępujący
w żadnej parze nazywamy \emph{krytycznym}.
\emph{$\mathcal V$-ścieżką wymiaru $p$} nazywamy ciąg
\[
\alpha_0^p\prec\beta_0^{p+1}\succ\alpha_1^p
\prec\beta_1^{p+1}\succ\cdots\succ\alpha_r^p,
\]
w którym każda $\alpha_i\prec\beta_i$ należy do
$\mathcal V$, zaś $\alpha_{i+1}\prec\beta_i$ i
$\alpha_{i+1}\ne\alpha_i$. Parowanie jest
\emph{acykliczne}, gdy nie istnieje niebanalna
zamknięta $\mathcal V$-ścieżka z
$\alpha_r=\alpha_0$.
\end{definicja}

To ostatnie ograniczenie jest kluczowe. Samo parowanie
może krążyć w zamkniętym cyklu i nie pozwalać na
monotoniczne upraszczanie kompleksu. Acykliczność
stanowi kombinatoryczny odpowiednik braku zamkniętych
trajektorii gradientowych.

\begin{stwierdzenie}[Charakteryzacja przez graf skierowany]
\label{thm:tg-morse-acyclic}
Niech $\mathcal V$ będzie parowaniem sympleksów
skończonego kompleksu $K$. W grafie Hassego skierujmy
każdą krawędź nieparowaną od sympleksu mniejszego do
większego, a parowaną od większego do mniejszego.
Wtedy $\mathcal V$ jest acykliczne wtedy i tylko wtedy,
gdy tak otrzymany graf skierowany nie ma cyklu skierowanego.
\end{stwierdzenie}

\begin{proof}
Zamknięta $\mathcal V$-ścieżka daje cykl skierowany:
od $\beta_i$ przechodzimy strzałką parowaną do
$\alpha_i$, następnie strzałką nieparowaną do
$\beta_{i-1}$ (indeksy czytamy cyklicznie).
W drugą stronę niech $C$ będzie cyklem skierowanym
w grafie Hassego i wybierzmy na nim sympleks $\beta$
o maksymalnym wymiarze $p+1$. Strzałka wychodząca
z $\beta$ nie może prowadzić do nadsympleksu,
bo przeczyłoby to maksymalności wymiaru. Prowadzi
więc do pewnej ściany $\alpha$ wymiaru $p$ po
krawędzi parowanej. Ponieważ $\alpha$ jest już
sparowana z $\beta$, jej następna strzałka nie może
prowadzić w dół: oznaczałoby to drugą parę zawierającą
$\alpha$. Musi zatem prowadzić po krawędzi
nieparowanej do innego sympleksu wymiaru $p+1$.
Powtarzając rozumowanie, przechodzimy naprzemiennie
między wymiarami $p+1$ i $p$ (bez możliwości przejścia
na większy wymiar), aż cykl się zamknie. Otrzymaliśmy
niebanalną zamkniętą $\mathcal V$-ścieżkę.
\end{proof}

\begin{uwaga}[Dyskretna funkcja Morse'a]
Równoważny opis można podać przy pomocy funkcji
$F$ na zbiorze sympleksów, dla której każdy sympleks
ma co najwyżej jedną bezpośrednią ścianę o wartości
$\ge F(\sigma)$ i co najwyżej jeden bezpośredni
nadsympleks o wartości $\le F(\sigma)$.
Po zadaniu acyklicznego parowania taki $F$ można
skonstruować, przypisując liczbom wierzchołków
skierowanego grafu Hassego porządek topologiczny.
Dla krawędzi nieparowanej wartości rosną od ściany
do nadsympleksu, a dla parowanej kierunek nierówności
jest odwrotny. W dalszym wykładzie wygodniej
pracować bezpośrednio z parowaniem.
\end{uwaga}

\begin{samepage}
\begin{center}
\begin{minipage}{.96\linewidth}\centering
\begin{tikzpicture}[scale=.91,>=Stealth,line cap=round]
 \node[tgTealVertex] (v0) at (0,0) {$a$};
 \node[tgBlueVertex] (v1) at (2,0) {$b$};
 \node[tgBlueVertex] (v2) at (4,0) {$c$};
 \draw[tgTree,->] (v1)--(v0);
 \draw[tgTree,->] (v2)--(v1);
 \node[text=tgNavy,font=\small] at (2,-.63)
   {pary: $(b,ab)$ i $(c,bc)$};
 \node[text=tgTeal,font=\small] at (0,.52) {krytyczny};
\end{tikzpicture}\par\smallskip

{\footnotesize Dla drzewa wybór rodzica każdego niekorzeniowego
wierzchołka daje acykliczne parowanie.}
\end{minipage}
\end{center}
\end{samepage}

\section{Dyskretny gradient na grafie i drzewa rozpinające}
\label{sec:tg-forman-on-graph}

\begin{twierdzenie}[Parowanie wyznaczone przez drzewo]
\label{thm:tg-forman-tree}
Niech $G$ będzie skończonym spójnym grafem z $n\ge1$
wierzchołkami i $m$ krawędziami, a $T$ drzewem
rozpinającym $G$. Traktujemy $G$ jako jednowymiarowy
kompleks CW; ewentualne pętle i krawędzie równoległe
można uprzednio podzielić na prosty kompleks symplicjalny.
Wybierzmy korzeń $r\in V(T)$ i dla każdego $v\ne r$
sparujmy $v$ z krawędzią łączącą $v$ z jego rodzicem
w drzewie $T$. Otrzymane parowanie jest acykliczne.
Ma dokładnie jeden krytyczny wierzchołek oraz
$m-n+1$ krytycznych krawędzi.
\end{twierdzenie}

\begin{proof}
Każdy wierzchołek $v\ne r$ występuje dokładnie raz,
a każda krawędź drzewa również dokładnie raz jako
para z jej dalszym od korzenia końcem. Krawędzie
spoza $T$ nie są sparowane i są krytyczne; ich liczba
to $m-(n-1)=m-n+1$. Korzeń jest jedynym krytycznym
wierzchołkiem.

Każda $\mathcal V$-ścieżka wymiaru zero biegnie
od wierzchołka do jego rodzica: jeśli
$(v,e_v)$ jest parą, drugim końcem $e_v$ jest
właśnie rodzic $v$. Odległość w drzewie od korzenia
maleje więc o jeden przy każdym kroku takiej ścieżki.
Żadna niebanalna $\mathcal V$-ścieżka nie może być
zamknięta. W wymiarze jeden nie ma par krawędź--trójkąt,
więc innych zamkniętych ścieżek również nie ma.
\end{proof}

\begin{przyklad}[Trójkątny cykl]
Dla grafu $C_3$ o wierzchołkach $a,b,c$ wybieramy
drzewo $T=\{ab,bc\}$ i korzeń $a$.
Pary to $(b,ab)$ oraz $(c,bc)$. Krytyczny pozostaje
wierzchołek $a$ i krawędź $ac$. Otrzymujemy więc
jedną krytyczną komórkę wymiaru zero i jedną
wymiaru jeden, co odpowiada homotopijnemu typowi $S^1$.
Jeżeli teraz dołączymy wypełniający trójkąt $abc$,
możemy dodatkowo sparować krawędź $ac$ z trójkątem.
Pozostaje tylko krytyczny wierzchołek $a$, zgodnie
z kontraktowalnością dysku.
\end{przyklad}

\section{Elementarny kolaps a zachowanie typu homotopijnego}
\label{sec:tg-forman-collapse}

\begin{definicja}[Para wolna i elementarny kolaps]
Niech $K$ będzie skończonym kompleksem symplicjalnym.
Para $\alpha\prec\beta$ nazywa się \emph{wolna},
gdy $\beta$ jest jedynym maksymalnym sympleksem $K$
zawierającym $\alpha$. Usunięcie obu sympleksów
$\alpha$ i $\beta$ (bez usuwania ich pozostałych
ścian) nazywamy \emph{elementarnym kolapsem}.
Zakładamy, że usuwane sympleksy tworzą rzeczywiście
parę kowymiaru jeden; wtedy pozostały zbiór sympleksów
nadal jest kompleksem.
\end{definicja}

\begin{stwierdzenie}[Kolaps nie zmienia typu homotopijnego]
Jeżeli $K'$ powstaje ze skończonego kompleksu
symplicjalnego $K$ przez elementarny kolaps wolnej
pary $\alpha\prec\beta$, to włączenie
$|K'|\hookrightarrow|K|$ jest równoważnością homotopijną,
a nawet istnieje silna retrakcja deformacyjna
$|K|$ na $|K'|$.
\end{stwierdzenie}

\begin{proof}
Realizację sympleksu $\beta$ traktujemy jako kulę
geometryczną $D^{p+1}$. Wnętrze ściany $\alpha$
stanowi otwartą część jej brzegu, do której nie
przylega inny sympleks maksymalny. Kula $\beta$
bez wnętrza $\alpha$ i bez własnego wnętrza
pozostawia pozostałe ściany, które tworzą podkompleks
będący dyskiem wymiaru $p$. Można deformować całą
kulę $\beta$ na sumę tych pozostałych ścian przez
przesuwanie punktów wzdłuż odcinków prowadzących
od punktu wybranego na wnętrzu wolnej ściany ku
przeciwległej części brzegu. Jest to standardowa
retrakcja PL dysku na brzeg z usuniętą otwartą ścianą,
wykonywana relatywnie do pozostałych ścian.
Wobec wolności $\alpha$ deformacja jest identycznością
na wszystkich pozostałych sympleksach i skleja się
z identycznością poza $\beta$. Dostajemy silną
retrakcję deformacyjną $|K|$ na $|K'|$.
\end{proof}

\begin{przyklad}[Dysk z jednym krytycznym wierzchołkiem]
Niech $K$ będzie pełnym trójkątem $abc$. Najpierw
usuwamy wolną parę $(ac,abc)$, pozostawiając dwa
odcinki $ab$ i $bc$. Następnie usuwamy
$(c,bc)$, potem $(b,ab)$. Otrzymujemy jednoelementowy
kompleks $\{a\}$. Każdy krok jest elementarnym
kolapsem. W ten sposób nie tylko obliczamy homologie,
ale wykazujemy kontraktowalność geometryczną całej
realizacji $|K|$.
\end{przyklad}

\section{Komórki krytyczne, kompleks Morse'a i nierówności}
\label{sec:tg-forman-inequalities}

\begin{uwaga}[Acykliczne parowanie a elementarny kolaps]
Nie każda para z acyklicznego parowania jest wolną parą
w wyjściowym kompleksie. Nie wolno więc usuwać wszystkich
sparowanych komórek w dowolnej kolejności jako kolapsów
bez wcześniejszego sprawdzenia warunku wolnej ściany.
Twierdzenie Formana konstruuje odpowiedni model CW,
a jego dowód wymaga uporządkowania par zgodnie z dyskretnym
przepływem gradientowym. Jest to istotnie mocniejszy wynik
niż elementarne twierdzenie o pojedynczym kolapsie.
\end{uwaga}

\begin{twierdzenie}[Forman; podstawowe twierdzenie]
\label{thm:tg-forman-main}
Niech $K$ będzie skończonym regularnym kompleksem CW,
a $\mathcal V$ acyklicznym parowaniem jego komórek.
Niech $c_p$ oznacza liczbę niesparowanych komórek
wymiaru $p$. Wtedy $|K|$ jest homotopijnie równoważny
kompleksowi CW mającemu dokładnie $c_p$ komórek
wymiaru $p$ dla każdego $p$.
W szczególności istnieje skończony kompleks łańcuchowy
Morse'a nad dowolnym ciałem $\mathbb F$, którego grupy łańcuchów
w stopniu $p$ mają wymiar $c_p$, a homologie są
izomorficzne z homologiami $|K|$.
\end{twierdzenie}

\begin{uwaga}[Status twierdzenia Formana]
Jest to głęboki rezultat; nie wynika on wyłącznie z
istnienia parowania. Dowód wykorzystuje porządkowanie
acyklicznych par i odpowiednie przekształcenia komórkowe
(algebraiczną wersję można skonstruować przez eliminację
par w kompleksie łańcuchowym). W szczególności pary
występujące w polu gradientowym \emph{nie muszą być
od razu wolnymi parami} w pierwotnym kompleksie.
Powyżej udowodniono bezpośrednio tylko zachowanie
homotopii dla rzeczywistych elementarnych kolapsów.
Pełne twierdzenie Formana przyjmujemy tutaj z literatury.
\end{uwaga}

\begin{twierdzenie}[Słabe i silne nierówności Morse'a]
\label{thm:tg-forman-ineq}
Niech $K$ będzie skończonym kompleksem symplicjalnym,
$\mathcal V$ acyklicznym parowaniem, a $c_p$ liczbą
krytycznych $p$-sympleksów. Niech
$b_p=\dim_{\mathbb F}H_p(|K|;\mathbb F)$
dla ustalonego ciała $\mathbb F$. Wtedy dla każdego
$p\ge0$ zachodzi $c_p\ge b_p$, a ponadto
\[
 \sum_{j=0}^p(-1)^{p-j}c_j
 \ge \sum_{j=0}^p(-1)^{p-j}b_j,
 \qquad
 \sum_{j\ge0}(-1)^jc_j=\chi(|K|).
\]
\end{twierdzenie}

\begin{proof}
Stosujemy twierdzenie Formana, które daje skończony
kompleks łańcuchowy Morse'a $(M_*,d_*)$ z
$\dim M_p=c_p$ i homologiami o wymiarach $b_p$.
Niech $r_p=\operatorname{rank}d_p$; przyjmujemy
$r_0=0$ i $r_p=0$ ponad wymiarem kompleksu.
Z twierdzenia o rzędzie oraz definicji homologii mamy
$b_p=(c_p-r_p)-r_{p+1}$, zatem
\[
 c_p=b_p+r_p+r_{p+1}.
\]
Od razu dostajemy $c_p\ge b_p$.
Po zsumowaniu naprzemiennym od zera do $p$
wszystkie składniki $r_j$ w środku teleskopowo
się znoszą i pozostaje $r_{p+1}\ge0$. Stąd silne
nierówności. Sumując po wszystkich wymiarach,
ponieważ najwyższy rząd jest zerowy, uzyskujemy
$\sum(-1)^pc_p=\sum(-1)^pb_p=\chi(|K|)$,
gdzie ostatnia równość jest wzorem Eulera--Poincarégo
udowodnionym w rozdziale o triangulacjach.
\end{proof}

\begin{przyklad}[Sfera z czterech trójkątów]
Rozważmy triangulację brzegu czworościanu o wierzchołkach
$0,1,2,3$. Parujemy wierzchołki $1,2,3$
odpowiednio z krawędziami $01,02,03$ oraz krawędzie
$12,13,23$ z trójkątami $012,013,023$.
Niesparowane są dokładnie wierzchołek $0$ i trójkąt
$123$. Parowanie jest acykliczne: w wymiarze zero
wszystkie ścieżki kierują się w stronę korzenia $0$,
a w wymiarze jeden po wejściu w sparowany trójkąt
przechodzimy tylko do krawędzi incydentnych z $0$,
które nie są sparowane z trójkątami, więc ścieżka
nie może się zamknąć. Otrzymujemy
$(c_0,c_1,c_2)=(1,0,1)$, zgodne z
$H_0(S^2)=H_2(S^2)\cong\mathbb Z$ i $H_1(S^2)=0$.
\end{przyklad}

\begin{samepage}
\begin{center}
\begin{minipage}{.96\linewidth}\centering
\begin{tikzpicture}[scale=1.0,line join=round,line cap=round]
 \coordinate (a) at (-1.15,-.7);\coordinate (b) at (1.2,-.7);
 \coordinate (c) at (.3,1.35);\coordinate (d) at (-.2,.18);
 \fill[tgGold!20] (b)--(c)--(d)--cycle;
 \draw[tgEdge] (a)--(b)--(c)--cycle (a)--(d)--(b) (c)--(d);
 \node[tgTealVertex] at (a) {$0$};
 \node[tgBlueVertex] at (b) {$1$};
 \node[tgBlueVertex] at (c) {$2$};
 \node[tgBlueVertex] at (d) {$3$};
 \node[font=\small,text=tgTeal] at (-1.55,-1.25) {krytyczne $0$};
 \node[font=\small,text=tgGold] at (1.6,.6) {krytyczne $123$};
\end{tikzpicture}\par\smallskip

{\footnotesize Na brzegu czworościanu pozostają dwie komórki
krytyczne: wierzchołek i ściana.}
\end{minipage}
\end{center}
\end{samepage}

\section{Zestawienie teorii ciągłej i dyskretnej}
\label{sec:tg-morse-bridge}

Oba rodzaje teorii Morse'a prowadzą do opisu topologii
za pomocą niewielkiej liczby danych krytycznych,
ale nie są definicyjnie tożsame. Klasyczna funkcja
Morse'a na rozmaitości jest różniczkowalna, a jej
punkty krytyczne są badane przy użyciu hesjanu.
Dyskretne pole gradientowe działa na zbiorze komórek,
a warunkiem podstawowym jest brak zamkniętych
$\mathcal V$-ścieżek. W obu przypadkach otrzymujemy
nierówności wiążące liczby krytycznych obiektów
z liczbami Bettiego i charakterystyką Eulera.

\begin{uwaga}[Minimalność a doskonałość]
Parowanie nazywamy \emph{doskonałym nad ciałem}
$\mathbb F$, gdy $c_p=b_p$ dla wszystkich $p$.
Jest to warunek silniejszy niż samo spełnianie równości
$\sum(-1)^pc_p=\chi$. Dla grafu spójnego konstrukcja
z drzewa rozpinającego daje doskonałe parowanie
($c_0=1$, $c_1=m-n+1$). Dla brzegu czworościanu
otrzymujemy również doskonałe parowanie z dwiema
komórkami krytycznymi. W bardziej złożonych
kompleksach istnienie doskonałego parowania nie jest
zagwarantowane.
\end{uwaga}

\chapter{Od triangulacji do topologii i geometrii rozmaitości}
\label{ch:tg-most-geometria}

Ostatni rozdział tej iteracji scala trzy perspektywy:
triangulację jako komórkowy model przestrzeni,
graf dualny jako zapis sąsiedztwa komórek oraz
operatory łańcuchowe jako algebraiczne narzędzia
wykrywania cykli. Nie zastępujemy nimi pełnej
geometrii różniczkowej, lecz budujemy formalne
przejście do następnych części książki.

\section{Graf dualny nie koduje całej triangulacji}
\label{sec:tg-limit-dual}

Graf dualny triangulacji powierzchni zamkniętej jest
trójregularny. Odwrotna konstrukcja nie jest jednak
jednoznaczna: do odtworzenia triangulacji trzeba wiedzieć,
które konkretne boki trójkątów są utożsamiane oraz
w jakich kierunkach. Jest to analogiczne do rozróżnienia
samego abstrakcyjnego grafu i jego układu rotacji.

\begin{przyklad}[Dwie ściany, dwie powierzchnie]
Bierzemy dwa abstrakcyjne trójkąty i sklejamy ich
boki parami. Przy sklejeniu odpowiadających boków
z przeciwną orientacją na brzegach powstaje sfera:
ma trzy klasy wierzchołków, trzy klasy krawędzi
i dwie ściany, zatem $\chi=3-3+2=2$.
Istnieje również sklejenie tych samych dwóch trójkątów
w kwadracie fundamentalnym torusa z identyfikacją
przeciwległych boków. Wtedy jest jedna klasa wierzchołków,
trzy klasy krawędzi i dwie ściany, więc
$\chi=1-3+2=0$. Otrzymujemy torus.
W obu przypadkach abstrakcyjny multigraf dualny
ma dwa wierzchołki i trzy równoległe krawędzie.
Różnica topologiczna nie jest widoczna w grafie
dualnym bez dodatkowych informacji o sklejaniu.
Trzeba podkreślić, że otrzymane identyfikacjami
kompleksy są obiektami typu $\Delta$, nie zaś
abstrakcyjnymi kompleksami symplicjalnymi z dwoma
różnymi trójkątami.
\end{przyklad}

\begin{stwierdzenie}[Rekonstrukcja z danych incydencji]
Jeżeli dla skończonej rodziny zorientowanych trójkątów
zadano pary sklejanych boków oraz homeomorfizmy
tych boków, to ich iloraz jest jednoznacznie określoną
przestrzenią topologiczną złożoną z dwukomórek.
Jest ona powierzchnią zamkniętą wtedy i tylko wtedy,
gdy każda krawędź ma dwa incydentne boki, a ogniwo
każdego wierzchołka w otrzymanym kompleksie typu
$\Delta$ jest okręgiem.
\end{stwierdzenie}

\begin{proof}
Iloraz jest określony jako przestrzeń otrzymana przez
identyfikację punktów brzegowych zgodnie z zadanymi
homeomorfizmami. Wnętrza trójkątów mają otoczenia
dyskowe. Punkt we wnętrzu sklejonej krawędzi ma
otoczenie dyskowe dokładnie przy sklejeniu dwóch
półdysków bez rozgałęzień. Wierzchołek ma otoczenie
będące stożkiem nad ogniwem, zatem jest punktem
powierzchni dokładnie wtedy, gdy ogniwo jest
okręgiem. Zwartość i własność Hausdorffa wynikają
ze skończoności układu i sklejeń wzdłuż zwartych
odcinków. Kryterium jest zatem lokalne.
\end{proof}

\section{Laplacjan komórkowy i dyskretna teoria Hodge'a}
\label{sec:tg-cellular-hodge}

\begin{uwaga}[Dane dualne nie wyznaczają topologii]
Sam graf dualny triangulacji nie wyznacza na ogół typu
homeomorficznego rozmaitości. Aby odtworzyć kompleks
w wymiarze co najmniej dwa, potrzebujemy informacji,
które ściany są utożsamiane, jak działają odwzorowania
sklejające oraz jakie są ogniwa niżej wymiarowych sympleksów.
Warunki lokalne, takie jak sferyczność ogniw, nie wynikają
z samej listy par przylegających komórek.
\end{uwaga}

\begin{definicja}[Kołańcuchy i różniczka komórkowa]
Niech $K$ będzie skończonym zorientowanym
kompleksem symplicjalnym wymiaru $d$. Niech
$C^p(K;\mathbb R):=\operatorname{Hom}(C_p(K;\mathbb R),\mathbb R)$.
Operatory brzegowe $\partial_p$ wyznaczają
\emph{różniczki kobrzegowe}
$d_p:=\partial_{p+1}^{\,T}:C^p\to C^{p+1}$
po wyborze ortonormalnych baz zorientowanych sympleksów.
Przy ustalonych dodatnich iloczynach skalarnych
w przestrzeniach $C^p$ przez $d_p^*$ oznaczamy
operator sprzężony. \emph{Laplasjanem Hodge'a}
na $p$-kołańcuchach nazywamy
\[
\Delta_p:=d_{p-1}d_{p-1}^*+d_p^*d_p,
\]
przyjmując brakujące składniki na końcach
kompleksu za zero.
\end{definicja}

\begin{stwierdzenie}[Dodatniość i harmoniczność]
Dla każdego $p$ laplasjan $\Delta_p$ jest
sam sprzężony i nieujemny oraz
\[
\ker\Delta_p=\ker d_p\cap\ker d_{p-1}^*.
\]
\end{stwierdzenie}

\begin{proof}
Dla dowolnego $\omega\in C^p$ mamy
\[
\langle\Delta_p\omega,\omega\rangle
=\|d_p\omega\|^2+\|d_{p-1}^*\omega\|^2.
\]
Każdy składnik jest nieujemny, więc operator jest
nieujemny; jest samosprzężony jako suma operatorów
postaci $AA^*$ oraz $B^*B$. Jeżeli $\Delta_p\omega=0$,
oba kwadraty norm muszą znikać, zatem
$d_p\omega=0$ i $d_{p-1}^*\omega=0$.
Odwrotnie oba równania implikują
$\Delta_p\omega=0$.
\end{proof}

\begin{twierdzenie}[Skończeniewymiarowe twierdzenie Hodge'a]
\label{thm:tg-finite-hodge}
Niech $K$ będzie skończonym kompleksem symplicjalnym,
a w każdej grupie $C^p(K;\mathbb R)$ ustalmy dodatni
iloczyn skalarny. Wtedy zachodzi ortogonalny rozkład
\[
C^p=\operatorname{im}d_{p-1}
\ \oplus^{\perp}\ \ker\Delta_p
\ \oplus^{\perp}\ \operatorname{im}d_p^*,
\]
a naturalne odwzorowanie
$\ker\Delta_p\to H^p(K;\mathbb R)$,
$\omega\mapsto[\omega]$, jest izomorfizmem.
\end{twierdzenie}

\begin{proof}
Z tożsamości $d_pd_{p-1}=0$ dostajemy
$\langle d_{p-1}u,d_p^*v\rangle
=\langle d_pd_{p-1}u,v\rangle=0$.
Zatem $U:=\operatorname{im}d_{p-1}$ oraz
$W:=\operatorname{im}d_p^*$ są ortogonalne.
Ich wspólne dopełnienie ortogonalne to
$U^\perp\cap W^\perp
=\ker d_{p-1}^*\cap\ker d_p=\ker\Delta_p$.
Skończony wymiar pozwala więc rozłożyć $C^p$
wskazaną sumą ortogonalną.

Jeżeli $\omega\in\ker\Delta_p$ reprezentuje zerową
klasę kohomologii, to $\omega=d_{p-1}u\in U$,
a równocześnie $\omega\in U^\perp$, zatem $\omega=0$.
To daje injektywność. Dla surjektywności wybierzmy
zamknięty kołańcuch $\eta\in\ker d_p$. Rozłóżmy go
jako $\eta=d_{p-1}u+h+d_p^*v$.
Po zastosowaniu $d_p$ dostajemy $d_pd_p^*v=0$,
a zatem $\|d_p^*v\|^2
=\langle d_pd_p^*v,v\rangle=0$.
Stąd $\eta-d_{p-1}u=h\in\ker\Delta_p$,
co oznacza, że klasa $[\eta]$ ma reprezentanta
harmonicznego. Dowód jest zakończony.
\end{proof}

\begin{przyklad}[Trójkąt, jego wypełnienie i wymiar harmoniki]
Dla grafu cyklicznego $C_3$ mamy
$\dim H^1(C_3;\mathbb R)=1$, zatem
$\dim\ker\Delta_1=1$. Po dołączeniu
pełnego trójkąta pojawia się niezerowy operator
$d_1:C^1\to C^2$; klasa cyklu brzegowego staje
się brzegiem dwukomórki i
$H^1(\Delta^2;\mathbb R)=0$.
W konsekwencji $\ker\Delta_1=0$.
Rozróżnienie następuje na poziomie $d_1$,
nie na poziomie samego grafu wierzchołków i krawędzi.
\end{przyklad}

\section{Co będzie potrzebne w topologii trój- i czterowymiarowej}
\label{sec:tg-roadmap-topology}

\begin{uwaga}[Grafy diagramów, rozkładów i plumbing]
W kolejnym etapie przejdziemy od kombinatoryki
kompleksów do struktur bardziej geometrycznych.
Grafy Cayleya i grafy grup pozwolą kodować działania
i rozkłady grup fundamentalnych. Grafy Taita
przypisane diagramom węzłów pozwolą przenieść
niektóre rachunki z teorii węzłów na wielomiany grafowe.
Grafy rozkładów JSJ opisują sposób sklejania części
trójrozmaitości. Wreszcie ważone grafy plumbingowe
prowadzą do czterorozmaitości, których formy przecięcia
można odczytać z macierzy związanej z grafem.
Na razie istotne jest rozróżnienie trzech rodzajów danych:
(1) abstrakcyjnego grafu, (2) zanurzenia i danych sklejania,
(3) dodatkowych wag, orientacji i ramek geometrycznych.
Bez tego rozróżnienia późniejsze wzory na formy przecięcia
byłyby matematycznie nieprecyzyjne.
\end{uwaga}

\section*{Wskazówki bibliograficzne do iteracji piątej}
Standardowe wprowadzenie do kompleksów symplicjalnych i
homologii zawiera A. Hatcher, \emph{Algebraic Topology}.
Ruchy bistellarne i triangulacje rozmaitości omawia
U. Pachner, \emph{PL Homeomorphic Manifolds are Equivalent
by Elementary Shellings}.
Dyskretna teoria Morse'a jest przedstawiona w pracy
R. Formana, \emph{Morse Theory for Cell Complexes}.
Grafy Reeba i ich zastosowania do funkcji Morse'a
opisuje literatura topologii obliczeniowej, m.in.
H. Edelsbrunner i J. Harer, \emph{Computational Topology}.
Znane w skrypcie wyniki o charakterystyce Eulera,
homologii i rozkładach na uchwyty stanowią
przygotowanie do tych konstrukcji.


\chapter{Grafy Cayleya, działania na drzewach i teoria Bassa--Serre'a}
\label{ch:tg-bass-serre}

Dotychczas traktowaliśmy graf jako skończoną przestrzeń kombinatoryczną lub
jednowymiarowy kompleks CW. Teraz dopuszczamy także grafy nieskończone. Niektóre z nich,
jak graf Cayleya grupy względem skończonego zbioru generatorów,
są lokalnie skończone: każdy wierzchołek ma skończony stopień.
Natomiast drzewa pojawiające się w teorii Bassa--Serre'a mogą
mieć wierzchołki nieskończonego stopnia. Grafy nieskończone
występują naturalnie jako nakrycia kompleksów oraz przestrzenie,
na których działają grupy. Działania na drzewach pozwalają
odczytywać rozkłady grup fundamentalnych, a następnie rozkłady rozmaitości.

\section{Grafy Cayleya i prezentacje grup}
\label{sec:tg-cayley}

\begin{definicja}[Graf Cayleya]
Niech $G$ będzie grupą, a $S\subseteq G\setminus\{1\}$
skończonym, symetrycznym zbiorem, czyli $S=S^{-1}$.
\emph{Grafem Cayleya} $\operatorname{Cay}(G,S)$ nazywamy graf,
którego zbiorem wierzchołków jest $G$, a krawędzie łączą
$g$ z $gs$ dla wszystkich $g\in G$ i $s\in S$.
Jeżeli $s=s^{-1}$, odpowiednią nieorientowaną krawędź liczymy
jednokrotnie. Gdy potrzebujemy zachować pełną informację o generatorach,
rozpatrujemy graf skierowany z łukami $g\to gs$
etykietowanymi przez $s$. W niezorientowanym grafie
utożsamiamy łuk $(g,s)$ z łukiem $(gs,s^{-1})$.
\end{definicja}

\begin{stwierdzenie}
Niech $G$ będzie grupą, a $S\subseteq G\setminus\{1\}$
skończonym zbiorem symetrycznym. Graf $\operatorname{Cay}(G,S)$
jest spójny wtedy i tylko wtedy, gdy $S$ generuje $G$.
\end{stwierdzenie}
\begin{proof}
Jeżeli $S$ generuje $G$, to dla dowolnego $g\in G$ istnieje
rozkład $g=s_1\cdots s_k$, gdzie $s_i\in S$. Kolejne
wierzchołki $1,s_1,s_1s_2,\ldots,g$ tworzą spacer od
jedności do $g$. Ponieważ z każdego $g$ otrzymujemy
spacer do $h$ przez translację spaceru od $1$ do
$g^{-1}h$, graf jest spójny. Odwrotnie, jeżeli istnieje
spacer $1=g_0,g_1,\ldots,g_k=g$, to dla każdego $i$
z definicji krawędzi $g_{i-1}^{-1}g_i\in S$.
Po wymnożeniu równości dostajemy
$g=(g_0^{-1}g_1)\cdots(g_{k-1}^{-1}g_k)$, a więc
każdy element należy do podgrupy generowanej przez $S$.
\end{proof}

\begin{przyklad}[Trzy rodziny grafów Cayleya]
Dla $G=\mathbb Z$ i $S=\{-1,1\}$ otrzymujemy
nieskończoną prostą, z wierzchołkami oznaczonymi liczbami całkowitymi.
Dla $G=\mathbb Z/n\mathbb Z$, $n\ge3$, oraz $S=\{-1,1\}$
otrzymujemy cykl $C_n$. Dla $G=\mathbb Z^2$ i
$S=\{(\pm1,0),(0,\pm1)\}$ powstaje kwadratowa sieć
płaszczyzny. W dwóch ostatnich przykładach odczytujemy
z grafu odpowiednio relację $a^n=1$ i przemienność
generatorów odpowiadających przesunięciom poziomym
oraz pionowym.
\end{przyklad}

\begin{center}
\begin{tikzpicture}[line cap=round,>=Stealth,scale=.94]
 \foreach \x in {-2,-1,0,1,2} {
   \node[tgBlueVertex] (z\x) at (\x*.58,0) {};
   \node[font=\scriptsize,text=tgNavy] at (\x*.58,-.38) {$\x$};
 }
 \draw[tgEdge] (z-2)--(z-1)--(z0)--(z1)--(z2);
 \draw[tgEdge,dashed] (-1.70,0)--(-1.38,0) (1.38,0)--(1.70,0);
 \node[font=\small,text=tgNavy] at (0,.62) {$\operatorname{Cay}(\mathbb Z,\{\pm1\})$};
 \begin{scope}[shift={(4.45,0)}]
  \foreach \a in {0,60,...,300} {
    \node[tgTealVertex] (c\a) at (\a:1.02) {};
    \node[font=\scriptsize,text=tgNavy] at (\a:1.37) {$\pgfmathparse{int(\a/60)}\pgfmathresult$};
  }
  \foreach \a [evaluate=\a as \b using {int(mod(\a+60,360))}] in {0,60,...,300} {
    \draw[tgTree] (c\a)--(c\b);
  }
  \node[font=\small,text=tgNavy] at (0,-1.72) {$\operatorname{Cay}(\mathbb Z/6\mathbb Z,\{\pm1\})$};
 \end{scope}
\end{tikzpicture}
\end{center}

\begin{twierdzenie}[Graf Cayleya grupy wolnej]
Niech $F_r=\langle a_1,\ldots,a_r\rangle$, $r\ge1$,
będzie grupą wolną, a $S=\{a_1^{\pm1},\ldots,a_r^{\pm1}\}$.
Wtedy $\operatorname{Cay}(F_r,S)$ jest drzewem,
a dla $r\ge2$ każdy wierzchołek ma stopień $2r$.
\end{twierdzenie}
\begin{proof}
Z poprzedniego stwierdzenia graf jest spójny. Gdyby istniał
cykl prosty $g_0,g_1,\ldots,g_k=g_0$, $k\ge1$, to
kolejne etykiety krawędzi dawałyby relację
$s_1\cdots s_k=1$ w $F_r$. Ponieważ cykl nie zawraca
natychmiast tą samą krawędzią, słowo nie zawiera sąsiedniej
pary $a_i a_i^{-1}$ ani $a_i^{-1}a_i$; jest więc słowem
zredukowanym i niepustym. W grupie wolnej słowo takie
nie reprezentuje jedności, co daje sprzeczność.
Spójny graf bez cykli jest drzewem. Dla $r\ge2$
różne generatory i ich odwrotności dają $2r$
różnych sąsiadów każdego wierzchołka.
\end{proof}

\section{Działania grup na drzewach i ilorazy}
\label{sec:tg-tree-actions}

\begin{definicja}[Działanie bez inwersji]
Niech grupa $G$ działa przez automorfizmy na drzewie $T$.
Mówimy, że działanie jest \emph{bez inwersji}, jeżeli
żaden element $g\in G$ nie zamienia końców krawędzi,
pozostawiając tę krawędź jako zbiór niezmienniczy.
Stabilizatorem wierzchołka $v$ jest
$G_v=\{g\in G:gv=v\}$, a krawędzi $e$ stabilizator
$G_e=\{g\in G:ge=e\}$.
\end{definicja}

\begin{stwierdzenie}[Wolne działania na drzewach]
Niech grupa $G$ działa swobodnie i bez inwersji na drzewie
$T$ przez automorfizmy grafu. Wtedy odwzorowanie
$T\to G\backslash T$ jest nakryciem grafów, $T$ jest
nakryciem uniwersalnym ilorazu, a grupa $G$ jest wolna.
Jeżeli iloraz jest skończonym grafem spójnym z $v$
wierzchołkami i $e$ krawędziami, to $G\cong F_{e-v+1}$.
\end{stwierdzenie}
\begin{proof}
Swoboda i brak inwersji oznaczają, że żadne nietrywialne
przekształcenie nie identyfikuje dwóch półkrawędzi
wychodzących z tego samego wierzchołka. Faktycznie,
jeśli $g$ przenosi $e_1$ na $e_2$, a ich wspólny początek
$v$ pozostaje w tej samej orbicie, to dla lokalnej
identyfikacji wybieramy element stabilizujący $v$;
ze swobody musi być on jednością. Gwiazdy wierzchołków
odwzorowują się więc bijektywnie na gwiazdy obrazów:
to definicja nakrycia grafów. Drzewo jest po prostu
spójne, zatem nakrycie jest uniwersalne. Z teorii
nakryć mamy $G\cong\pi_1(G\backslash T)$, a grupa
fundamentalna spójnego grafu jest wolna. Jeżeli
iloraz jest skończony, wzór na rząd otrzymujemy z
liczby krawędzi poza drzewem rozpinającym: $e-v+1$.
\end{proof}

\section{Grafy grup i twierdzenie van Kampena}
\label{sec:tg-graphs-of-groups}

\begin{definicja}[Graf grup]
Niech $\Gamma$ będzie spójnym grafem, którego każdą
krawędź rozpatrujemy w dwóch przeciwnych orientacjach
$e$ oraz $\bar e$. \emph{Graf grup} składa się z grup
$G_v$ przypisanych wierzchołkom, grup $G_e=G_{\bar e}$
przypisanych krawędziom oraz monomorfizmów
$\iota_{e,-}:G_e\hookrightarrow G_{o(e)}$ i
$\iota_{e,+}:G_e\hookrightarrow G_{t(e)}$.
Zmiana orientacji krawędzi wymienia oba monomorfizmy.
\end{definicja}

\begin{twierdzenie}[Dwa podstawowe przypadki van Kampena]
\label{thm:tg-vk-amalgam-hnn}
\begin{enumerate}[label=(\roman*)]
\item Niech $X=A\cup B$ będzie łukowo spójną
przestrzenią, gdzie $A,B,A\cap B$ są otwarte,
łukowo spójne i zawierają wspólny punkt bazowy.
Załóżmy, że homomorfizmy
$\pi_1(A\cap B)\to\pi_1(A)$ i
$\pi_1(A\cap B)\to\pi_1(B)$ są injekcjami.
Wtedy
$\pi_1(X)\cong\pi_1(A)*_{\pi_1(A\cap B)}\pi_1(B)$.
\item Niech $Y,Z$ będą łukowo spójnymi przestrzeniami,
a $\iota_-,\iota_+\colon Z\to Y$ ciągłymi odwzorowaniami
indukującymi monomorfizmy na grupach fundamentalnych
(po wyborze zgodnych punktów bazowych). Jeśli przestrzeń
powstaje przez dołączenie cylindra $Z\times[0,1]$ do $Y$
za pomocą $\iota_-$ i $\iota_+$ na jego końcach,
to jej grupa fundamentalna jest
rozszerzeniem HNN grupy $\pi_1(Y)$ po odpowiednio
utożsamionych obrazach $\pi_1(Z)$.
\end{enumerate}
\end{twierdzenie}
\begin{proof}
W przypadku (i) klasyczne twierdzenie van Kampena stwierdza, że
$\pi_1(X)$ jest ilorazem iloczynu wolnego
$\pi_1(A)*\pi_1(B)$ przez relacje utożsamiające
oba obrazy każdego elementu $\pi_1(A\cap B)$.
Ponieważ oba odwzorowania są injektywne, jest to
właśnie iloczyn wolny z amalgamacją. W przypadku (ii) przestrzenie $Y$ i $Z$ są łukowo spójne,
a iloraz powstaje ze $Y\sqcup(Z\times[0,1])$ przez
utożsamienie $(z,0)$ z $\iota_-(z)$ i $(z,1)$ z
$\iota_+(z)$; oba homomorfizmy $\pi_1(Z)\to\pi_1(Y)$
są injekcjami. Ustalmy $z_0\in Z$, $y_0\in Y$
oraz ścieżki w $Y$ od $y_0$ do obrazów $z_0$.
Odcinek $\{z_0\}\times[0,1]$, wraz z tymi ścieżkami,
tworzy nową pętlę, której klasę oznaczamy przez $t$.
Po rozcięciu cylindra w jego środku otrzymujemy pokrycie
z niespójnym przecięciem. Grupoidowa wersja twierdzenia
van Kampena (lub zwykła wersja po dołączeniu pomocniczego
łuku bazowego) daje jedną nową literę $t$ i relacje
utożsamiające obrazy pętli w dwóch końcach cylindra.
Dostajemy prezentację
\[
\langle \pi_1(Y),t\mid t\iota_-(c)t^{-1}=\iota_+(c)
\ \text{dla wszystkich }c\in\pi_1(Z)\rangle.
\]
Ta prezentacja definiuje rozszerzenie HNN. Przy zmianie ścieżek bazowych odpowiednie monomorfizmy
zmieniają się przez sprzężenie, zatem otrzymane prezentacje
nadal opisują izomorficzne grupy. W szczególności nie można
pominąć nowego generatora $t$: inaczej utracilibyśmy
informację o cyklu w grafie sklejania.
\end{proof}

\begin{przyklad}[Iloczyn wolny i torus jako HNN]
Gdy graf grup jest odcinkiem, grupy na jego końcach to
$\mathbb Z=\langle a\rangle$ i
$\mathbb Z=\langle b\rangle$, a grupa krawędzi jest
trywialna, otrzymujemy $\mathbb Z*\mathbb Z\cong F_2$.
Gdy graf jest jedną pętlą, grupa wierzchołka i krawędzi
są równe $\mathbb Z$, a oba monomorfizmy grupy krawędzi
w grupę wierzchołka są identycznościami, otrzymujemy
\[
\langle a,t\mid tat^{-1}=a\rangle\cong\mathbb Z^2.
\]
Pierwszy przykład odpowiada sklejaniu dwóch okręgów
w punkcie, drugi jest przykładem cylindrycznego
sklejania prowadzącego do torusa.
\end{przyklad}

\begin{center}
\begin{tikzpicture}[>=Stealth,scale=1]
 \node[tgBlueVertex] (a) at (-2.5,0) {};
 \node[tgRedVertex] (b) at (-.5,0) {};
 \draw[tgEdge] (a)--node[above,font=\small] {$G_e$}(b);
 \node[below=5pt of a,font=\scriptsize] {$G_{v_1}$};
 \node[below=5pt of b,font=\scriptsize] {$G_{v_2}$};
 \node[tgTealVertex] (c) at (2.5,0) {};
 \draw[tgAccent] (c) .. controls (1.65,.95) and (3.35,.95) .. (c);
 \node[font=\scriptsize,text=tgNavy] at (2.5,1.16) {jedna krawędź pętlowa};
 \node[below=5pt of c,font=\scriptsize] {$G_v$};
 \node[text=tgNavy,font=\scriptsize] at (-1.5,-.80) {amalgamacja};
 \node[text=tgNavy,font=\scriptsize] at (2.5,-.80) {rozszerzenie HNN};
\end{tikzpicture}
\end{center}

\begin{uwaga}[Wymagane dane grafu grup]
Nie wystarcza znajomość abstrakcyjnego grafu ilorazowego
$G\backslash T$. Konstrukcja potrzebuje stabilizatorów
wierzchołków i krawędzi, homomorfizmów włączających oraz
wybranych elementów realizujących zmianę przedstawicieli orbit.
Dwa działania o jednakowym grafie ilorazowym mogą dawać
grupy podstawowe o różnych prezentacjach.
Ponadto drzewo Bassa--Serre'a nie musi być lokalnie skończone:
stopnie wierzchołków zależą od indeksów podgrup krawędziowych.
\end{uwaga}

\begin{twierdzenie}[Fundamentalne twierdzenie Bassa--Serre'a]
Niech $G$ działa bez inwersji na drzewie $T$.
Po wybraniu przedstawicieli orbit wierzchołków i krawędzi,
dróg do tych przedstawicieli i odpowiednich elementów grupy,
działanie określa graf grup na grafie ilorazowym
$G\backslash T$. Jego grupami wierzchołków i krawędzi
są stabilizatory, a monomorfizmy krawędziowe pochodzą
z inkluzji stabilizatorów po potrzebnych sprzężeniach. Grupa fundamentalna tego
grafu grup jest izomorficzna z $G$. Odwrotnie, każdy graf
grup posiada drzewo Bassa--Serre'a, na którym jego grupa
fundamentalna działa bez inwersji, a otrzymany graf grup
jest izomorficzny z wyjściowym (po odpowiednich wyborach).
\end{twierdzenie}

\begin{uwaga}[Zakres dowodu]
Powyższy wynik jest bardziej ogólny niż twierdzenie van
Kampena o amalgamacji i HNN. Do pełnego dowodu trzeba
skonstruować normalną postać słów i drzewo Bassa--Serre'a
jako graf klas bocznych oraz sprawdzić, że nie ma w nim
cykli. Są to dodatkowe wyniki, które tutaj przyjmujemy.
Dowiedzione wcześniej przypadki pokazują natomiast,
dlaczego w rozkładach topologicznych pojawiają się właśnie
iloczyny wolne z amalgamacją i rozszerzenia HNN.
\end{uwaga}

\begin{stwierdzenie}[Stopnie w drzewie Bassa--Serre'a]
Niech $\mathcal G$ będzie spójnym grafem grup o grafie bazowym
$\Gamma$, a $T$ jego drzewem Bassa--Serre'a. Dla
wierzchołka $\widetilde v\in T$ nad $v\in V(\Gamma)$
liczba półkrawędzi wychodzących z $\widetilde v$ wynosi
\[
\deg_T(\widetilde v)=
\sum_{e\colon o(e)=v}[G_v:\iota_{e,-}(G_e)],
\]
przy czym sumujemy po wszystkich zorientowanych
półkrawędziach grafu $\Gamma$ rozpoczynających się w $v$;
w pętli obie orientacje są osobnymi składnikami.
Sumy i indeksy mogą być nieskończone.
\end{stwierdzenie}
\begin{proof}
W konstrukcji drzewa Bassa--Serre'a wierzchołki nad $v$
są klasami bocznymi grup $G_v$ w grupie fundamentalnej
grafu grup. Po ustaleniu $\widetilde v$ półkrawędzie
przypisane określonemu $e$ są dokładnie orbitą działania
$G_v$ na wybranej półkrawędzi. Jej stabilizator jest
obrazem $\iota_{e,-}(G_e)$, a z twierdzenia o orbitach
i stabilizatorach liczba tych półkrawędzi to podany indeks.
Zsumowanie po różnych orbitach krawędzi daje wzór.
\end{proof}
\begin{przyklad}
Dla iloczynu wolnego $C_2*C_3$ graf grup ma dwa wierzchołki
z grupami $C_2$ i $C_3$ oraz jedną krawędź z grupą trywialną.
Drzewo Bassa--Serre'a jest dwudzielne: stopnie wierzchołków
obu typów wynoszą odpowiednio $2$ i $3$. Natomiast
$\mathbb Z*\mathbb Z$ otrzymane z dwóch grup $\mathbb Z$
i trywialnej grupy krawędzi daje drzewo o nieskończonych
stopniach wierzchołków. Nie można więc zakładać lokalnej
skończoności drzew Bassa--Serre'a bez dodatkowych hipotez.
\end{przyklad}

\begin{uwaga}[Literatura]
Pełną konstrukcję drzew i grafów grup opisuje J.-P. Serre,
\emph{Trees} oraz wykład J. F. Manninga
\emph{Bass--Serre Theory and Complexes of Groups}.
\end{uwaga}


% !TeX encoding = UTF-8
% Teoria wezlow -- tresc bez preambuly.
% W ksiazce: \input{teoria_wezlow_tresc.tex}
% Wersja samodzielna: kompilowac teoria_wezlow.tex.
% Wymagane srodowiska: definicja, twierdzenie, lemat, stwierdzenie,
% wniosek, przyklad, uwaga, proof; pakiety w preambule pliku glownego.

% Rysunki TikZ: xcolor, tikz oraz biblioteki arrows.meta,calc,positioning.
\definecolor{ktBlue}{HTML}{2369A8}
\definecolor{ktTeal}{HTML}{178F87}
\definecolor{ktCoral}{HTML}{D96152}
\definecolor{ktPurple}{HTML}{7755A7}
\definecolor{ktGold}{HTML}{C9932D}
\definecolor{ktInk}{HTML}{283747}


\chapter{Teoria węzłów: od diagramów do rozmaitości czterowymiarowych}
\label{ch:teoria-wezlow}

\noindent\textbf{Związek z wcześniejszym materiałem.}
Przy wprowadzaniu teorii węzłów wykorzystujemy poznaną teorię
grafów planarnych i wielomian Tutte'a. Po podstawach teorii
węzłów następuje nowy rozdział pokazujący bezpośrednio, jak
z grafów Taita wyprowadzić niezmienniki Jonesa i Goeritza.
W dalszych rozdziałach skorzystamy z teorii Bassa--Serre'a,
chirurgii Dehna i rozkładów JSJ.


Teoria węzłów łączy elementarne pytania o zanurzenia okręgu w przestrzeni
trójwymiarowej z fundamentalnymi problemami geometrii i topologii.
Jeden i ten sam węzeł można opisać diagramem na płaszczyźnie, grupą
fundamentalną jego uzupełnienia, wielomianami i grupami homologii,
strukturą geometryczną uzupełnienia lub sposobem dołączenia uchwytu
wymiaru cztery. Celem rozdziału jest przeprowadzenie tych przejść,
z precyzyjnym rozróżnieniem argumentów przeprowadzonych w całości
i twierdzeń wymagających odrębnej, zaawansowanej teorii. W drugim
przypadku formułujemy dokładną tezę, określamy używane narzędzie
i wyjaśniamy zależność od wcześniejszego materiału; szkic nie jest
przedstawiany jako kompletny dowód.

Przyjmujemy kategorię gładką, o ile nie zaznaczono inaczej.
Rozmaitości i sploty są orientowane, jeśli orientacja jest potrzebna.
Sfera $S^3$ jest orientowanym brzegiem kuli $B^4$; utożsamiamy ją także
z jednopunktowym uzwarceniem $\mathbb R^3$. W szczególności
nie odróżniamy węzłów w $\mathbb R^3$ i w $S^3$, o ile ich nośniki
nie przechodzą przez punkt w nieskończoności.

\section{Węzły, sploty i izotopie}
\label{sec:kt-podstawy}

\begin{uwaga}[Konwencja terminologiczna]
W całym skrypcie słowa \emph{zanurzenie} używamy w znaczeniu
angielskiego \emph{embedding}: gładkie odwzorowanie
$f\colon N\to M$ jest zanurzeniem, gdy jest injektywną immersją
i homeomorfizmem $N$ na $f(N)$ z topologią podprzestrzeni.
W szczególności różniczka $\mathrm df_x$ jest injektywna w każdym
$x\in N$. Samo założenie injektywności różniczki definiuje
\emph{immersję} i nie wystarcza do zanurzenia w powyższym
znaczeniu. Jeżeli $N$ jest zwarte, injektywna immersja
$N\to M$ do rozmaitości Hausdorffa jest zanurzeniem.
Odwołując się do kategorii topologicznej, przez zanurzenie
lokalnie płaskie rozumiemy \emph{lokally flat embedding}.
\end{uwaga}


\begin{definicja}[Węzeł i splot]
Niech $M$ będzie gładką rozmaitością trójwymiarową.
\emph{Węzłem} w $M$ nazywamy obraz gładkiego zanurzenia
$k\colon S^1\hookrightarrow M$.
\emph{Splotem o $r$ składowych} nazywamy obraz gładkiego zanurzenia
$k\colon\bigsqcup_{i=1}^{r}S^1\hookrightarrow M$, gdzie $r\geq1$.
Splot $L$ jest \emph{zorientowany}, jeżeli na każdej jego składowej
wybrano orientację. Węzeł \emph{trywialny} $U\subset S^3$ to okrąg
ograniczający gładko zanurzony dysk w $S^3$. Splot trywialny o $r$
składowych to splot będący brzegiem $r$ parami rozłącznych,
gładko zanurzonych dysków w $S^3$.
\end{definicja}

Zastrzeżenie dotyczące dysków jest istotne: każda pojedyncza składowa
splotu może być węzłem trywialnym, mimo że cały splot jest
nietrywialny. Pierwszym przykładem jest splot Hopfa.

\begin{definicja}[Izotopia otoczenia]
Niech $K_0,K_1\subset M$ będą węzłami. \emph{Izotopią otoczenia}
przenoszącą $K_0$ na $K_1$ nazywamy gładkie odwzorowanie
$H\colon M\times[0,1]\to M$ takie, że $H_t:=H(\cdot,t)$
jest dyfeomorfizmem dla każdego $t$, $H_0=\mathrm{id}_M$ i
$H_1(K_0)=K_1$. Węzły nazywamy \emph{równoważnymi}, jeżeli
istnieje taka izotopia. Dla splotów definicja jest analogiczna,
z zachowaniem oznaczeń i orientacji składowych, jeśli są ustalone.
\end{definicja}

Twierdzenie o rozszerzaniu izotopii zanurzeń stwierdza precyzyjniej,
że dla gładkiej rodziny zanurzeń $k_t\colon S^1\hookrightarrow M$,
$t\in[0,1]$, istnieje izotopia otoczenia $H_t$ taka, że
$k_t=H_t\circ k_0$ dla wszystkich $t$. Twierdzenie stosuje się,
gdy źródło zanurzenia jest zwarte, a rodzina jest gładka względem
obu zmiennych. Jest to wynik niebanalny w ogólnej teorii podrozmaitości,
ale bardzo wygodny w praktyce. Wszystkie węzły rozpatrujemy zatem
\emph{z dokładnością do izotopii otoczenia}, nie tylko do homeomorfizmu
ich abstrakcyjnych nośników, które zawsze są okręgami.

\begin{przyklad}[Węzły torusowe]
Niech $p,q\in\mathbb Z$ będą względnie pierwsze. Na standardowym
torusie $T^2\subset S^3$, parametryzowanym przez
$S^1\times S^1$, rozważmy krzywą
$t\mapsto(e^{ipt},e^{iqt})$, gdzie $t\in\mathbb R/2\pi\mathbb Z$.
Jej obraz, po zanurzeniu standardowego torusa w $S^3$, oznaczamy
$T(p,q)$ i nazywamy \emph{węzłem torusowym}. Jeśli $|p|=1$
lub $|q|=1$, jest on trywialny. Węzeł $T(2,3)$ jest trójlistnikiem;
węzeł $T(2,-3)$ jest jego lustrzanym odbiciem, z odpowiednio
dobraną konwencją zanurzenia torusa.
\end{przyklad}

\begin{definicja}[Lustrzane odbicie i odwracalność]
Niech $K\subset S^3$ będzie zorientowanym węzłem.
Przez $\overline K$ oznaczamy jego obraz przy dyfeomorfizmie
$S^3\to S^3$ odwracającym orientację, z orientacją na obrazie
przeniesioną z $K$. Przez $K^{\mathrm{rev}}$ oznaczamy ten sam
węzeł z przeciwną orientacją. Węzeł jest \emph{amfichiralny},
gdy jest równoważny swemu odbiciu lustrzanemu po pominięciu
orientacji; jest \emph{odwracalny}, gdy $K$ i $K^{\mathrm{rev}}$
są równoważne jako węzły zorientowane.
\end{definicja}

\begin{uwaga}
Nie należy utożsamiać lustrzanego odbicia z odwróceniem orientacji.
Przykładowo trójlistnik jest odwracalny, ale nie jest amfichiralny.
Wykazanie drugiej własności wymaga niezmiennika rozróżniającego
jego dwa odbicia, np. sygnatury lub wielomianu Jonesa.
\end{uwaga}


\begin{center}
\begin{tikzpicture}[scale=.9,>=Stealth,line cap=round]
  \begin{scope}[shift={(-2.65,0)}]
    \fill[ktBlue!7] (0,0) circle (1.22);
    \draw[ktBlue,line width=3.5pt,-{Stealth[length=3mm]}] (0,0) circle (1.0);
    \node[text=ktInk,font=\small\bfseries] at (0,-1.55) {węzeł trywialny};
    \node[text=ktBlue,font=\small] at (0,-1.95) {$U$};
  \end{scope}
  \begin{scope}[shift={(2.65,0)}]
    \fill[ktTeal!7] (-.9,0) circle (.77);
    \fill[ktCoral!7] (.9,0) circle (.77);
    \draw[ktTeal,line width=3pt] (-.9,0) circle (.74);
    \draw[ktCoral,line width=3pt] (.9,0) circle (.74);
    \node[text=ktInk,font=\small\bfseries] at (0,-1.55) {splot rozłączny};
    \node[text=ktCoral,font=\small] at (0,-1.95) {$U\sqcup U$};
  \end{scope}
\end{tikzpicture}
\end{center}
\noindent Po prawej przedstawiono dwie rozłączne, niesplecione składowe.
Kolor oznacza przynależność do jednej składowej splotu.

\section{Diagramy i ruchy Reidemeistera}
\label{sec:kt-diagramy}

Wybierzmy standardowe rzutowanie $\pi\colon\mathbb R^3\to\mathbb R^2$.
Po dowolnie małym zaburzeniu zanurzenia możemy założyć, że $\pi|_K$
jest immersją poza skończoną liczbą punktów podwójnych, że wszystkie
przecięcia są poprzeczne oraz że nie występują punkty potrójne.
Zaznaczając przy każdym punkcie podwójnym, które z dwóch włókien
leży wyżej, otrzymujemy \emph{diagram węzła}. Odcinek przebiegający
pod drugim rysujemy z krótką przerwą. Rozważania są identyczne dla
splotów.

\begin{center}
\begin{tikzpicture}[scale=.88,line cap=round,line join=round,>=Stealth]
\begin{scope}[shift={(-3.5,0)}]
  \node[font=\small\bfseries,text=ktInk] at (0,1.55) {pasmo niebieskie nad czerwonym};
  \draw[ktCoral,line width=2.8pt] (-1.0,1.0)--(-.14,.14);
  \draw[ktCoral,line width=2.8pt] (.14,-.14)--(1,-1);
  \draw[ktBlue,line width=3.2pt,-{Stealth[length=3mm]}] (-1,-1)--(1,1);
  \draw[ktCoral,line width=2.8pt,-{Stealth[length=3mm]}] (-1,1)--(-.44,.44);
\end{scope}
\begin{scope}[shift={(3.5,0)}]
  \node[font=\small\bfseries,text=ktInk] at (0,1.55) {pasmo czerwone nad niebieskim};
  \draw[ktBlue,line width=2.8pt] (-1,-1)--(-.14,-.14);
  \draw[ktBlue,line width=2.8pt] (.14,.14)--(1,1);
  \draw[ktCoral,line width=3.2pt,-{Stealth[length=3mm]}] (-1,1)--(1,-1);
  \draw[ktBlue,line width=2.8pt,-{Stealth[length=3mm]}] (-1,-1)--(-.43,-.43);
\end{scope}
\node[font=\footnotesize,text=ktInk] at (0,-1.55)
  {Przerwa na diagramie oznacza pasmo przebiegające pod spodem.};
\end{tikzpicture}
\end{center}

\begin{definicja}[Liczba skrzyżowań]
Niech $D$ będzie diagramem splotu $L$. Oznaczamy przez $c(D)$ liczbę
punktów podwójnych diagramu. \emph{Liczbą skrzyżowań} splotu jest
$c(L):=\min\{c(D):D\text{ jest diagramem }L\}$. Diagram realizujący
minimum nazywamy \emph{minimalnym}.
\end{definicja}

\begin{twierdzenie}[Reidemeister]
Niech $D_0,D_1$ będą dwoma regularnymi diagramami splotów
na zorientowanej płaszczyźnie $\mathbb R^2$, wraz z informacją,
która gałąź przechodzi górą w każdym skrzyżowaniu. Diagramy reprezentują równoważne sploty wtedy
i tylko wtedy, gdy jeden z nich można przekształcić w drugi
skończonym ciągiem izotopii płaszczyzny oraz trzech lokalnych
ruchów Reidemeistera: $\mathrm{R I}$ (dodanie lub usunięcie
pętli z jednym skrzyżowaniem), $\mathrm{R II}$ (dodanie lub
usunięcie dwóch skrzyżowań) i $\mathrm{R III}$ (przesunięcie
jednego pasma nad skrzyżowaniem dwóch pozostałych).
\end{twierdzenie}

\begin{proof}[Uzasadnienie geometryczne; użycie transwersalności]
Każdy z ruchów realizuje się w małej kuli zawierającej wyłącznie
odpowiednie fragmenty zanurzenia. Na przykład w ruchu $\mathrm{R II}$
jeden łuk przesuwa się nad drugim, a w $\mathrm{R III}$ pasmo
przechodzi nad skrzyżowaniem pozostałych dwóch. Przedłużamy
te lokalne izotopie przez identyczność poza kulą. Otrzymujemy
izotopie otoczenia, więc operacje zachowują typ splotu.

Dla kierunku odwrotnego nie wystarcza arbitralne oglądanie
rysunków w kolejnych chwilach. Ustalmy gładką izotopię
$H_t$ przenoszącą $L_0$ na $L_1$. Złożenie
$\pi\circ H_t|_{L_0}$ jest jednoparametrową rodziną
odwzorowań krzywych do płaszczyzny. Stosujemy parametryczną
transwersalność: po zaburzeniu rodziny, nieruszając jej
końców, wszystkie rzuty poza skończoną liczbą chwil mają
wyłącznie poprzeczne punkty podwójne o różnych wysokościach.
W chwili wyjątkowej występuje dokładnie jedna osobliwość
najniższej kodymensji. Są trzy możliwości: styczność kierunku
rzutowania do łuku (powstanie lub zanik pętli), styczność
dwóch gałęzi obrazu (powstanie lub zanik pary punktów
podwójnych) albo przecięcie trzech gałęzi obrazu. Odpowiadają
one ruchom $\mathrm{R I}$, $\mathrm{R II}$ i $\mathrm{R III}$.

Między chwilami wyjątkowymi diagramy zmieniają się wyłącznie
izotopią płaszczyzny. Zwartość przedziału parametru i ogólne
położenie pozwalają wybrać skończoną liczbę chwil wyjątkowych,
a więc skończony ciąg ruchów. Przywołana klasyfikacja
osobliwości rodziny rzutowań jest jedynym krokiem dowodu,
w którym używamy parametrycznej transwersalności.
\end{proof}


\begin{center}
\begin{tikzpicture}[scale=.86,line cap=round,line join=round,>=Stealth]
  \node[text=ktInk,font=\footnotesize\bfseries] at (-3.32,1.62) {R I: pętla};
  \draw[ktBlue,line width=2.8pt] (-5.3,-1)--(-5.3,1);
  \draw[->,ktTeal,line width=2pt] (-4.75,0)--(-3.95,0);
  \draw[ktBlue,line width=2.8pt] (-2.4,-1) .. controls (-2.7,-.2)
    and (-1.75,-.4) .. (-1.55,.18)
    .. controls (-1.4,.83) and (-2.65,1.4) .. (-2.8,.45)
    .. controls (-2.9,-.2) and (-2.4,-.65) .. (-2.4,1);
  \node[text=ktInk,font=\footnotesize\bfseries] at (2.64,1.62) {R II: dwa skrzyżowania};
  \draw[ktCoral,line width=2.8pt] (.3,-1)--(.3,1);
  \draw[ktBlue,line width=2.8pt] (1.0,-1)--(1.0,1);
  \draw[->,ktTeal,line width=2pt] (1.5,0)--(2.32,0);
  \draw[ktCoral,line width=2.8pt] (2.9,-1)
    .. controls (4.6,-.8) and (4.7,-.25) .. (3.9,.05)
    .. controls (3.0,.32) and (2.8,.8) .. (2.9,1);
  \draw[ktBlue,line width=2.8pt] (3.7,-1)--(3.7,1);
  \draw[white,line width=6pt] (3.95,.035)
    .. controls (3.78,.12) and (3.57,.22) .. (3.46,.26);
  \draw[ktCoral,line width=2.8pt] (3.95,.035)
    .. controls (3.78,.12) and (3.57,.22) .. (3.46,.26);
  \node[text=ktInk,font=\footnotesize] at (0,-1.52)
    {Schematy lokalne: izotopia w przestrzeni nie wymaga przecinania pasm.};
\end{tikzpicture}
\end{center}

Twierdzenie Reidemeistera jest fundamentem konstrukcji niezmienników
diagramowych. Wystarczy zdefiniować wielkość dla diagramów, a następnie
sprawdzić jej niezmienniczość przy trzech rodzajach ruchów.

\begin{definicja}[Znak skrzyżowania i skręt]
Niech $D$ będzie zorientowanym diagramem. W punkcie skrzyżowania
$c$ wybierzmy dodatnio zorientowane współrzędne na płaszczyźnie.
Niech $v_{\mathrm{nad}},v_{\mathrm{pod}}$ oznaczają wektory styczne
do zorientowanych pasm, odpowiednio przechodzącego górą i dołem.
Definiujemy $\varepsilon(c):=\operatorname{sgn}\det(v_{\mathrm{pod}},v_{\mathrm{nad}})
\in\{-1,1\}$ oraz \emph{skręt diagramu}
$w(D):=\sum_{c}\varepsilon(c)$.
\end{definicja}

Wielkość $w(D)$ pozostaje niezmienna przy ruchach $\mathrm{R II}$
i $\mathrm{R III}$, lecz przy $\mathrm{R I}$ zmienia się o $\pm1$.
Nie jest więc niezmiennikiem węzła. Będzie natomiast koniecznym
czynnikiem normalizującym wielomian Jonesa.

\begin{przyklad}[Dlaczego sama liczba skrzyżowań nie wystarcza]
Na diagramie okręgu można wykonać dowolnie wiele ruchów $\mathrm{R I}$.
Każdy nowy diagram przedstawia węzeł trywialny, chociaż liczba
skrzyżowań jest coraz większa. Nie oznacza to, że $c(K)$ nie jest
niezmiennikiem: minimum po wszystkich diagramach nie zależy od
wybranego przedstawienia.
\end{przyklad}

\section{Liczba splecenia i dualność Alexandera}
\label{sec:kt-linking}

\begin{definicja}[Liczba splecenia]
Niech $K_1,K_2\subset S^3$ będą rozłącznymi, zorientowanymi węzłami.
Wybieramy zorientowaną, zwartą, gładko zanurzoną powierzchnię
$F\subset S^3$ o brzegu $\partial F=K_1$, poprzeczną do $K_2$.
\emph{Liczbą splecenia} nazywamy algebraiczną liczbę punktów
przecięcia $\operatorname{lk}(K_1,K_2):=F\cdot K_2$.
Znak punktu przecięcia wyznaczamy przez orientacje $F$, $K_2$ i $S^3$.
\end{definicja}

Istnienie powierzchni $F$ zostanie udowodnione konstruktywnie
w części poświęconej algorytmowi Seiferta. Następująca własność
pokazuje, że definicja nie zależy od wyboru $F$.

\begin{stwierdzenie}
Dla dowolnych rozłącznych, zorientowanych węzłów
$K_1,K_2\subset S^3$ liczba $\operatorname{lk}(K_1,K_2)$ jest
dobrze określona, symetryczna względem składników i niezmiennicza
przy izotopii zorientowanego splotu.
\end{stwierdzenie}

\begin{proof}
Jeżeli $F$ i $F'$ są dwoma łańcuchami powierzchniowymi z tym samym
brzegiem $K_1$, to $F-F'$ jest cyklem dwuwymiarowym w $S^3$.
Ponieważ $H_2(S^3;\mathbb Z)=0$, jest brzegiem łańcucha
trójwymiarowego. Algebraiczne przecięcie brzegu łańcucha z zamkniętą
krzywą $K_2$ wynosi zero, a więc $F\cdot K_2=F'\cdot K_2$.
Dla symetrii wybieramy powierzchnie $F_1,F_2$ z brzegami
$K_1,K_2$ i ustawiamy je poprzecznie. Ich przecięcie jest
zorientowaną jednowymiarową rozmaitością z brzegiem, a rachunek
orientacji na końcach tego przecięcia daje
$F_1\cdot K_2=F_2\cdot K_1$. Niezmienniczość przy izotopii wynika z
zachowania algebraicznej liczby przecięć w rodzinie poprzecznej.
\end{proof}

\begin{twierdzenie}[Wzór diagramowy]
Niech $D$ będzie zorientowanym diagramem dwuskładowego splotu
$K_1\cup K_2\subset S^3$. Niech $C_{12}(D)$ oznacza zbiór
skrzyżowań, w których uczestniczą różne składowe. Wówczas
\[
\operatorname{lk}(K_1,K_2)=\frac12\sum_{c\in C_{12}(D)}\varepsilon(c).
\]
W szczególności suma znaków skrzyżowań między dwiema składowymi
jest parzysta.
\end{twierdzenie}

\begin{proof}
Oznaczmy prawą stronę wzoru przez $I(D)$. W ruchu $\mathrm{R I}$
na jednej składowej nie powstaje skrzyżowanie międzyskładowe.
W ruchu $\mathrm{R II}$ dwa ewentualne skrzyżowania różnych
składowych mają przeciwne znaki. Ruch $\mathrm{R III}$ nie
zmienia znaków ani par składowych uczestniczących w
skrzyżowaniach. Zatem $I(D)$ zależy tylko od izotopii splotu.

Zmienienie jednego skrzyżowania między $K_1$ i $K_2$
zmienia $I(D)$ o $1$ albo $-1$. Taka sama zmiana zachodzi
w $F\cdot K_2$: w małej kuli realizującej zmianę
skrzyżowania krzywa $K_2$ przechodzi na przeciwną stronę
powierzchni rozpiętej na $K_1$, przekraczając ją dokładnie
raz z odpowiednim znakiem. Zmiana skrzyżowania wewnątrz
jednej składowej nie wpływa na żadną z obu wielkości.

Dowolny splot dwóch składowych można uczynić rozdzielnym
przez skończony ciąg zmian skrzyżowań: wybieramy najpierw
kolejność składowych, przeprowadzamy jedną z nich zawsze
ponad drugą, a następnie umieszczamy je w rozłącznych kulach.
Dla splotu rozdzielnego mamy $I(D)=0$ i
$\operatorname{lk}(K_1,K_2)=0$. Cofając zmiany skrzyżowań,
wnioskujemy, że obie wielkości są równe dla wyjściowego
splotu. Ponieważ liczba splecenia jest całkowita, suma znaków
międzyskładowych jest parzysta.
\end{proof}

\begin{przyklad}
Dla splotu Hopfa o zgodnych orientacjach dwóch skrzyżowań
otrzymujemy $\operatorname{lk}=1$. Dla splotu rozdzielnego
liczba splecenia wynosi zero. Odwrotna implikacja jest fałszywa:
trójskładowe pierścienie Borromeusza mają zerowe liczby splecenia
każdej pary składowych, ale cały splot jest nierozdzielny.
\end{przyklad}

\begin{twierdzenie}[Homologia uzupełnienia splotu]
Niech $L=K_1\cup\dots\cup K_r\subset S^3$ będzie splotem o $r\geq1$
składowych i niech $X_L:=S^3\setminus\operatorname{int}N(L)$,
gdzie $N(L)$ oznacza sumę rozłącznych, domkniętych otoczeń
tubularnych składowych. Wtedy
$H_1(X_L;\mathbb Z)\cong\mathbb Z^r$.
Izomorfizm wyznaczają klasy dodatnio zorientowanych meridianów
poszczególnych składowych.
\end{twierdzenie}

\begin{proof}
Przestrzeń $X_L$ jest homotopijnie równoważna $S^3\setminus L$.
Z dualności Alexandera otrzymujemy
$\widetilde H_i(S^3\setminus L;\mathbb Z)
\cong\widetilde H^{2-i}(L;\mathbb Z)$.
Dla $i=1$ prawa strona jest $H^1(L;\mathbb Z)\cong\mathbb Z^r$.
Meridian otaczający $K_i$ ma liczbę splecenia jeden z $K_i$
i zero z pozostałymi składowymi. Homomorfizm przypisujący
klasie pętli jej $r$ liczb splecenia identyfikuje więc
meridiany z bazą standardową, kończąc dowód.
\end{proof}

\section{Grupa węzła i prezentacja Wirtingera}
\label{sec:kt-wirtinger}

\begin{definicja}
Dla węzła $K\subset S^3$ \emph{grupą węzła} nazywamy
$G_K:=\pi_1(S^3\setminus K,x_0)$, gdzie $x_0\in S^3\setminus K$.
Zmiana punktu bazowego daje grupę izomorficzną, lecz niekanonicznie.
Dla splotu $L$ analogicznie oznaczamy $G_L:=\pi_1(S^3\setminus L,x_0)$.
\end{definicja}

Rysunek diagramu koduje prezentację tej grupy. Dzielimy go na
\emph{łuki Wirtingera}: maksymalne fragmenty przebiegające
bez przerwania pod innym pasmem. Każdemu takiemu łukowi
przypisujemy generator, reprezentowany przez meridian.

\begin{twierdzenie}[Prezentacja Wirtingera]
Niech $D$ będzie spójnym diagramem węzła $K\subset S^3$,
z co najmniej jednym skrzyżowaniem, a $x_1,\dots,x_n$ niech będą
meridianami odpowiadającymi łukom Wirtingera. W każdym skrzyżowaniu
trzy uczestniczące generatory spełniają relację postaci
$x_j=x_i x_k x_i^{-1}$ albo $x_j=x_i^{-1}x_kx_i$,
w zależności od orientacji i położenia pasm; $x_i$ odpowiada
pasmowi górnemu. Grupa $G_K$ ma prezentację o tych generatorach
i relacjach, przy czym jedna relacja jest konsekwencją pozostałych.
\end{twierdzenie}

\begin{proof}[Konstrukcja van Kampena]
Wybieramy regularny diagram i podnosimy łuki górne na
niewielką dodatnią wysokość względem płaszczyzny rzutowania.
Przestrzeń nad diagramem po usunięciu niewielkich rurek
wokół pasm deformacyjnie retraktuje się na układ, którego
grupa fundamentalna jest wolna na meridianach $x_1,\dots,x_n$.
Dołączając części przestrzeni pod diagramem, w każdym
skrzyżowaniu narzucamy jedną relację. Jeśli $x_i$ jest
meridianem łuku górnego, a $x_k,x_j$ oznaczają meridiany
łuków dolnych, przeciągnięcie małego meridianu spod pasma
górnego jest sprzężeniem, zatem $x_j=x_i^{\pm1}x_kx_i^{\mp1}$.
Twierdzenie van Kampena pokazuje, że nie ma innych relacji.

Żeby wyjaśnić redundancję, wybierzmy sferę obejmującą
cały diagram oraz jeden wspólny punkt bazowy dla meridianów.
Odpowiadające skrzyżowaniom relacje są brzegami
komórek dwuwymiarowych standardowego kompleksu Wirtingera.
Suma tych komórek tworzy zamkniętą powierzchnię po uwzględnieniu
orientacji łuków; stąd jedna relacja jest konsekwencją
pozostałych. Formalnie można usunąć jedną komórkę, nie
zmieniając $\pi_1$; pozostaje prezentacja o $n$ generatorach
i $n-1$ relacjach. Dla diagramu węzła bez skrzyżowań
grupa jest bezpośrednio równa $\mathbb Z$.
\end{proof}

\begin{uwaga}[Rola jednej zbędnej relacji Wirtingera]
Redundancja jednej relacji dotyczy spójnego diagramu węzła
i standardowej prezentacji meridionalnej. Nie oznacza,
że wszystkie pozostałe relacje są algebraicznie niezależne
w dowolnej prezentacji grupy: po wykonaniu transformacji
Tietzego można uzyskać inną liczbę generatorów i relacji.
Przy obliczaniu wielomianu Alexandera należy konsekwentnie
wybierać tę samą abelianizację wszystkich meridianów.
\end{uwaga}

\begin{przyklad}[Trójlistnik]
Dla standardowego trójskrzyżowaniowego diagramu trójlistnika
prezentacja Wirtingera upraszcza się do
$G_T\cong\langle a,b\mid aba=bab\rangle$.
Warto rozwinąć algebraicznie przejście do bardziej symetrycznej
prezentacji
$\langle x,y\mid x^2=y^3\rangle$.
Kładziemy $x:=aba$ oraz $y:=ab$. Z relacji $aba=bab$ wynika
\[
x^2=(aba)(bab)=(ab)^3=y^3.
\]
W drugą stronę, zakładając prezentację $\langle x,y\mid x^2=y^3\rangle$,
definiujemy
$a:=y^{-1}x$ oraz $b:=x^{-1}y^2$. Wtedy $ab=y$, zatem
$aba=(ab)a=ya=x$, a także
$bab=(x^{-1}y^2)(y^{-1}x)(x^{-1}y^2)=x^{-1}y^3=x$.
Stąd $aba=bab$. Otrzymaliśmy wzajemnie odwrotne podstawienia,
więc obie prezentacje opisują tę samą grupę.

Dodatkowo element $z:=x^2=y^3$ jest centralny, gdyż komutuje
zarówno z $x$, jak i z $y$. Po przejściu do ilorazu przez normalną
podgrupę generowaną przez $z$ dostajemy
$\langle x,y\mid x^2=y^3=1\rangle\cong C_2*C_3$.
Ta grupa jest nieprzemienna, a więc i sama grupa trójlistnika
jest nieprzemienna. W szczególności nie jest izomorficzna z
$G_U\cong\mathbb Z$, co ponownie dowodzi nietrywialności
trójlistnika.
\end{przyklad}

\begin{stwierdzenie}
Dla każdego węzła $K\subset S^3$ abelianizacja
$G_K^{\mathrm{ab}}:=G_K/[G_K,G_K]$ jest izomorficzna z $\mathbb Z$.
W szczególności wszystkie meridiany w grupie węzła mają
jednakowy obraz w abelianizacji (po ustaleniu orientacji).
\end{stwierdzenie}

\begin{proof}
Twierdzenie Hurewicza w stopniu pierwszym daje
$G_K^{\mathrm{ab}}\cong H_1(S^3\setminus K;\mathbb Z)$.
Poprzedni rezultat o homologii uzupełnienia splotu zastosowany
do jednej składowej daje $H_1\cong\mathbb Z$, generowaną
przez meridian.
\end{proof}

\begin{uwaga}[Dlaczego sama grupa nie wystarcza]
Grupa $G_K$ jest niezwykle silnym niezmiennikiem, jednak do
odtworzenia pełnej struktury uzupełnienia naturalnie dodaje się
dane \emph{peryferyjne}: wyróżnioną klasę sprzężoności meridianu
i długości (longitude). Przykładowo różne węzły torusowe mogą
mieć grupy abstrakcyjnie powiązane, a badanie włożenia podgrupy
peryferyjnej oddziela informacje gubione przy samej abelianizacji.
\end{uwaga}

\section{Kolorowania Foxa i wyznacznik węzła}
\label{sec:kt-fox}

Kolorowania Foxa stanowią pierwszy przykład niezmiennika, który
można obliczyć metodami elementarnej algebry liniowej.

\begin{definicja}[Kolorowanie $n$-kolorowe]
Niech $n\geq2$ będzie liczbą całkowitą, a $D$ diagramem węzła
$K\subset S^3$. \emph{$n$-kolorowaniem Foxa} nazywamy przypisanie
każdemu łukowi Wirtingera elementu $\mathbb Z/n\mathbb Z$ takiego,
że przy każdym skrzyżowaniu zachodzi
$2b\equiv a+c\pmod n$, gdzie $b$ jest kolorem pasma górnego,
a $a,c$ kolorami dwóch fragmentów pasma dolnego.
Kolorowanie jest \emph{trywialne}, gdy wszystkie łuki mają ten sam kolor.
\end{definicja}

Każdy diagram dopuszcza dokładnie $n$ kolorowań stałych.
Kolorowań nietrywialnych może nie być. Relacje są liniowe,
zatem ich zbiór jest jądrem macierzy nad $\mathbb Z/n\mathbb Z$.

\begin{stwierdzenie}
Dla każdego całkowitego $n\geq2$ liczba $n$-kolorowań Foxa diagramu
węzła nie zależy od diagramu i jest niezmiennikiem węzła.
\end{stwierdzenie}


\begin{proof}
Rozważmy lokalnie trzy ruchy Reidemeistera. W ruchu $\mathrm{R I}$
dodany łuk $b$ ma kolor $a$, gdyż $2a=a+b$ implikuje $b=a$.
W ruchu $\mathrm{R II}$ kolor nowego łuku jest jednoznacznie
wyznaczony przez dwa stare kolory; druga relacja odzyskuje
kolor wyjściowy.
W ruchu $\mathrm{R III}$ dwie kolejności propagowania kolorów
prowadzą do tej samej wartości: definiując $R_b(a)=2b-a$,
obliczamy $R_c(R_b(a))=a+2(c-b)$ oraz
$R_{R_c(b)}(R_c(a))=a+2(c-b)$. Są to dokładnie dwie
możliwe kolejności przejścia przez trzy pasma. Każde $R_b$
jest inwolucją, ponieważ $R_b(R_b(a))=a$; nigdzie nie
trzeba dzielić przez dwa. W ten sposób ruchy indukują
wzajemnie odwrotne bijekcje zbiorów kolorowań.
\end{proof}

\begin{przyklad}[Trójlistnik jest nietrywialny]
Niech $a,b,c\in\mathbb Z/3\mathbb Z$ będą kolorami trzech łuków
standardowego diagramu trójlistnika. Relacje przy skrzyżowaniach
mają postać $2a=b+c$, $2b=a+c$, $2c=a+b$.
W $\mathbb Z/3\mathbb Z$ każda z nich jest równoważna warunkowi
$a+b+c=0$. Istnieje zatem $3^2=9$ kolorowań, z których
trzy są stałe. Węzeł trywialny ma tylko trzy kolorowania.
Zatem trójlistnik nie jest równoważny węzłowi trywialnemu.
\end{przyklad}

\begin{definicja}[Wyznacznik węzła]
Dla węzła $K\subset S^3$ wybieramy diagram $D$ o $n$ łukach
i $n$ skrzyżowaniach. Niech $A_D$ będzie całkowitą macierzą
relacji $2b-a-c=0$, gdzie kolumny odpowiadają łukom, a wiersze
skrzyżowaniom. Po skreśleniu dowolnego jednego wiersza i jednej
kolumny otrzymujemy macierz $(n-1)\times(n-1)$, której wyznacznik
co do wartości bezwzględnej oznaczamy $\det(K)$.
Dla węzła trywialnego przyjmujemy $\det(U)=1$.
\end{definicja}

Niezależność od skreślonego wiersza i kolumny oraz od diagramu
można wykazać, wykorzystując związki macierzy Foxa z wielomianem
Alexandera. Otrzymamy poniżej wzór $\det(K)=|\Delta_K(-1)|$.

\begin{przyklad}
Dla trójlistnika macierz relacji może mieć postać
\[
A_D=\begin{pmatrix}2&-1&-1\\-1&2&-1\\-1&-1&2\end{pmatrix}.
\]
Wyznacznik minora $\begin{psmallmatrix}2&-1\\-1&2\end{psmallmatrix}$
wynosi $3$, zatem $\det(T)=3$. Dla węzła ósemkowego
analogiczny rachunek daje $\det(4_1)=5$.
\end{przyklad}

\begin{stwierdzenie}[Kolorowania a wyznacznik]
Niech $K\subset S^3$ będzie węzłem, $n\ge2$, a
$\det(K)$ jego wyznacznikiem. Jeśli
$\gcd(n,\det(K))=1$, to każde $n$-kolorowanie Foxa
jest stałe. Gdy liczba pierwsza $p$ dzieli $\det(K)$,
istnieje nietrywialne $p$-kolorowanie.
\end{stwierdzenie}

\begin{proof}
Niech $A_D$ będzie macierzą relacji kolorowania, a
$B$ otrzymamy, usuwając jeden wiersz i jedną kolumnę.
Usunięta relacja jest zależna od pozostałych, natomiast
usuniętą współrzędną można ustawić na zero, odejmując
od wszystkich kolorów wspólną stałą. Zatem kolorowania
modulo kolorowania stałe są identyfikowane z
$\ker(B\bmod n)$. Macierz $B$ jest kwadratowa i
$|\det B|=\det(K)$. Jeśli $\gcd(n,\det B)=1$,
macierz $B$ jest odwracalna nad $\mathbb Z/n\mathbb Z$,
więc ma zerowe jądro. Jeśli $p$ dzieli wyznacznik,
$B\bmod p$ jest macierzą osobliwą nad ciałem
$\mathbb F_p$ i posiada niezerowy wektor jądra.
Jest on nietrywialnym kolorowaniem po odtworzeniu
odjętej stałej.
\end{proof}

\section{Powierzchnie Seiferta i rodzaj węzła}
\label{sec:kt-seifert}

\begin{definicja}[Powierzchnia Seiferta]
Niech $K\subset S^3$ będzie zorientowanym węzłem.
\emph{Powierzchnią Seiferta} węzła $K$ nazywamy zwartą,
spójną, zorientowaną, gładko zanurzoną powierzchnię
$F\subset S^3$ taką, że $\partial F=K$ oraz orientacja
brzegu indukowana z $F$ zgadza się z orientacją $K$.
\emph{Rodzajem węzła} jest $g(K):=\min\{g(F):F\text{ jest
powierzchnią Seiferta }K\}$.
\end{definicja}

W pierwszej chwili nie jest oczywiste, że powierzchnia Seiferta
zawsze istnieje. Następne twierdzenie dostarcza konstrukcji.

\begin{twierdzenie}[Seifert; algorytm wygładzania]
Każdy zorientowany węzeł $K\subset S^3$ ogranicza powierzchnię
Seiferta. Jeżeli $D$ jest spójnym diagramem z $c(D)$ skrzyżowaniami,
a po zorientowanym wygładzeniu wszystkich skrzyżowań otrzymujemy
$s(D)$ okręgów Seiferta, to algorytm Seiferta konstruuje spójną
powierzchnię $F_D$ o brzegu $K$ i
\[
\chi(F_D)=s(D)-c(D),\qquad
 g(F_D)=\frac{1-s(D)+c(D)}{2}.
\]
W szczególności $g(K)\leq(1-s(D)+c(D))/2$.
\end{twierdzenie}

\begin{proof}
Przy każdym skrzyżowaniu wygładzamy oba przecinające się łuki
w jedyny sposób zgodny z ich orientacjami. Powstaje skończona
rodzina rozłącznych, zorientowanych okręgów na płaszczyźnie.
Każdy okrąg wypełniamy dyskiem, umieszczając dyski na różnych
wysokościach wtedy, gdy są zagnieżdżone. W otoczeniu każdego
oryginalnego skrzyżowania przytwierdzamy do dwóch odpowiednich
dysków skręcone pasmo, które odtwarza lokalny przebieg nad i pod.
Powstaje gładko zanurzona, zorientowana powierzchnia $F_D$.
Zgodność orientacji na jej fragmentach wynika z zastosowanego
zorientowanego wygładzenia. Jej brzeg jest izotopijny z $K$.
Ponieważ diagram węzła jest spójny jako graf, układ dysków
i pasm również jest spójny. Zaczynając od $s(D)$ dysków
i dołączając $c(D)$ jednowymiarowych uchwytów, otrzymujemy
$\chi(F_D)=s(D)-c(D)$. Dla spójnej orientowalnej powierzchni
z jednym brzegiem $\chi(F_D)=1-2g(F_D)$, co daje podany wzór.
\end{proof}

\begin{center}
\begin{tikzpicture}[scale=.95,>=Stealth,line cap=round,line join=round]
  \node[font=\small\bfseries,text=ktInk] at (-3,1.27) {skrzyżowanie};
  \node[font=\small\bfseries,text=ktInk] at (2.4,1.27) {wygładzenie orientowane};
  \begin{scope}[shift={(-3,0)}]
    \draw[ktBlue,line width=2.9pt,-{Stealth[length=2.7mm]}] (-.9,-.9)--(.9,.9);
    \draw[ktCoral,line width=2.9pt] (-.9,.9)--(-.14,.14);
    \draw[ktCoral,line width=2.9pt,-{Stealth[length=2.7mm]}] (.14,-.14)--(.9,-.9);
  \end{scope}
  \draw[->,very thick,ktTeal] (-1.3,0)--(.25,0);
  \node[text=ktTeal,font=\footnotesize] at (-.53,.3) {wygładź};
  \begin{scope}[shift={(2.4,0)}]
    \draw[ktBlue,line width=2.9pt,-{Stealth[length=2.7mm]}]
       (-.85,-.9) .. controls (-.85,.05) and (-.3,.22) .. (-.3,.9);
    \draw[ktCoral,line width=2.9pt,-{Stealth[length=2.7mm]}]
       (.85,.9) .. controls (.85,-.05) and (.3,-.22) .. (.3,-.9);
  \end{scope}
  \node[font=\footnotesize,text=ktInk] at (0,-1.5)
    {Po wygładzeniu nie ma skrzyżowania; pasma służą później do odbudowania powierzchni.};
\end{tikzpicture}
\end{center}


\begin{center}
\begin{tikzpicture}[scale=.9,line join=round,line cap=round]
  \fill[ktBlue!16,draw=ktBlue,line width=1.4pt] (-2.15,0) ellipse (1.00 and .72);
  \fill[ktCoral!16,draw=ktCoral,line width=1.4pt] (2.15,0) ellipse (1.00 and .72);
  \fill[ktPurple!24,draw=ktPurple,line width=1.25pt]
    (-1.18,.30)
      .. controls (-.72,.34) and (-.28,.48) .. (-.04,.28)
      .. controls (.23,.05) and (.23,-.05) .. (-.04,-.28)
      .. controls (-.28,-.48) and (-.72,-.34) .. (-1.18,-.30)
      .. controls (-.25,-.44) and (.25,.44) .. (1.18,.30)
      .. controls (.72,.34) and (.28,.48) .. (.04,.28)
      .. controls (-.23,.05) and (-.23,-.05) .. (.04,-.28)
      .. controls (.28,-.48) and (.72,-.34) .. (1.18,-.30)
      .. controls (.25,-.44) and (-.25,.44) .. cycle;
  \draw[ktPurple,line width=1.85pt]
    (-1.18,.30) .. controls (-.25,.44) and (.25,-.44) .. (1.18,-.30);
  \draw[ktPurple,line width=1.85pt]
    (-1.18,-.30) .. controls (-.25,-.44) and (.25,.44) .. (1.18,.30);
  \node[text=ktBlue,font=\small] at (-2.15,-1.08) {dysk};
  \node[text=ktCoral,font=\small] at (2.15,-1.08) {dysk};
  \draw[ktPurple,->,thick] (0,-1.08)--(.06,-.42);
  \node[text=ktPurple,font=\small] at (0,-1.32) {pasmo};
  \node[text=ktInk,font=\footnotesize] at (0,-1.80)
    {Poglądowy model dysków Seiferta połączonych pasmem o gładnym styku na brzegach.};
\end{tikzpicture}
\end{center}

\begin{przyklad}
Standardowy diagram trójlistnika ma $c(D)=3$ i $s(D)=2$,
więc algorytm daje $g(F_D)=1$. Standardowy diagram węzła
ósemkowego ma $c(D)=4$, $s(D)=3$, a więc również $g(F_D)=1$.
Dopiero obliczenie niezmienników algebraicznych wykaże, że
żaden z tych węzłów nie ogranicza dysku; ich rodzaj jest
rzeczywiście równy jeden.
\end{przyklad}

\begin{stwierdzenie}
Dla węzła $K\subset S^3$ następujące warunki są równoważne:
(a) $K$ jest trywialny; (b) $g(K)=0$; (c) $K$ ogranicza
gładko zanurzony dysk w $S^3$.
\end{stwierdzenie}

\begin{proof}
Równoważność (b) i (c) wynika z klasyfikacji zwartych,
spójnych, orientowalnych powierzchni z jednym brzegiem:
powierzchnia rodzaju zero jest dyskiem. Jeżeli $K$ ogranicza
dysk, jego otoczenie regularne w $S^3$ jest kulą, a brzeg
dysku można wyprostować izotopią otoczenia do okręgu standardowego.
W drugą stronę standardowy okrąg ogranicza dysk, którego obraz
przy izotopii ogranicza $K$.
\end{proof}

\section{Macierz Seiferta i wielomian Alexandera}
\label{sec:kt-alexander}

Niech $F\subset S^3$ będzie powierzchnią Seiferta węzła $K$,
o rodzaju $g$. Grupa $H_1(F;\mathbb Z)$ jest wolna rzędu $2g$.
Niech $[\alpha_1],\dots,[\alpha_{2g}]$ będzie jej bazą,
reprezentowaną przez zamknięte krzywe na $F$. Ustalmy dodatnią
normalną do $F$; przez $\alpha_i^+$ oznaczamy nieznaczne
przesunięcie $\alpha_i$ w tym kierunku, rozłączne z $F$.

\begin{definicja}[Macierz Seiferta]
Macierzą Seiferta powierzchni $F$ w wybranej bazie nazywamy
macierz całkowitą $V_F=(v_{ij})_{i,j=1}^{2g}$ daną wzorem
$v_{ij}:=\operatorname{lk}(\alpha_i^+,\alpha_j)$.
Jej symetryzacją jest $V_F+V_F^T$.
\end{definicja}

\begin{stwierdzenie}[Postać przecięcia]
Niech $F$ będzie zorientowaną powierzchnią Seiferta węzła
$K\subset S^3$, a $V=V_F$ jej macierzą Seiferta. Wtedy
$V-V^T$ reprezentuje, z dokładnością do znaku wynikającego
z konwencji normalnej, algebraiczną postać przecięcia
na $H_1(F;\mathbb Z)$. W szczególności $\det(V-V^T)=1$.
\end{stwierdzenie}

\begin{proof}
Różnica $\operatorname{lk}(\alpha_i^+,\alpha_j)
-\operatorname{lk}(\alpha_j^+,\alpha_i)$ równa się algebraicznej
liczbie przecięć krzywych $\alpha_i$ i $\alpha_j$ na powierzchni
(z ewentualną wspólną zmianą znaku zależnie od orientacji).
Jest to lokalny rachunek: gdy dwie krzywe przecinają się
poprzecznie, odpowiednie dwa przesunięcia tworzą jedno
skrzyżowanie liczone z właściwym znakiem. Dla spójnej
orientowanej powierzchni rodzaju $g$ z jednym brzegiem
postać przecięcia ma w bazie symplektycznej macierz
$\begin{psmallmatrix}0&I_g\\-I_g&0\end{psmallmatrix}$,
której wyznacznik wynosi jeden. Zmiana bazy w $GL_{2g}(\mathbb Z)$
nie zmienia wyznacznika.
\end{proof}

\begin{definicja}[Wielomian Alexandera]
Dla węzła $K\subset S^3$ o powierzchni Seiferta $F$
i macierzy $V$ definiujemy
$\Delta_K(t):=\det(tV-V^T)\in\mathbb Z[t,t^{-1}]$,
z dokładnością do mnożenia przez jednostkę $\pm t^k$.
Dla powierzchni będącej dyskiem przyjmujemy wyznacznik pustej
macierzy równy jeden. Jeśli zachodzi taka potrzeba, wybieramy
symetryczną normalizację spełniającą
$\Delta_K(t)=\Delta_K(t^{-1})$ oraz $\Delta_K(1)=1$.
\end{definicja}

\begin{uwaga}
W definicji należy dopiero wykazać, że wynik nie zależy
od powierzchni Seiferta ani od jej bazy. Wyznacznik $\det(tV-V^T)$
zmienia się pod zmianą bazy przez mnożenie po lewej i prawej
przez odpowiednie macierze unimodularne, więc nie zmienia
się wcale. Niezależność od powierzchni wymaga dodatkowej idei:
równoważności $S$ macierzy Seiferta lub modułu Alexandera.
\end{uwaga}

\begin{twierdzenie}[Opis przez nakrycie cykliczne]
Niech $K\subset S^3$ będzie węzłem, $X_K$ jego uzupełnieniem,
a $\widetilde X_K\to X_K$ nieskończonym nakryciem cyklicznym
wyznaczonym przez epimorfizm
$\pi_1(X_K)\to H_1(X_K;\mathbb Z)\cong\mathbb Z$.
Grupa $H_1(\widetilde X_K;\mathbb Z)$ jest modułem nad
$\Lambda:=\mathbb Z[t,t^{-1}]$, przy czym $t$ działa przez
wybrany generator grupy przekształceń pokładowych. Dla
macierzy Seiferta $V$ moduł ten posiada kwadratową macierz
prezentacji $tV-V^T$. Stąd $\Delta_K(t)$ jest niezmiennikiem
węzła z dokładnością do $\pm t^k$.
\end{twierdzenie}

\begin{proof}[Rozcięcie, nakrycie i ciąg Mayera--Vietorisa]
Niech $F$ będzie powierzchnią Seiferta rodzaju $g$, a
$Y$ domknięciem $X_K$ rozciętego wzdłuż właściwie
zanurzonej części $F\cap X_K$. W brzegu $Y$ znajdują się
dwie kopie tej powierzchni, oznaczone $F_+$ i $F_-$.
Nieskończone nakrycie cykliczne $\widetilde X_K$ otrzymujemy,
klejąc kopie $Y_j$, $j\in\mathbb Z$, tak, by strona
$F_+$ kopii $Y_j$ była utożsamiona ze stroną $F_-$
kopii $Y_{j+1}$. Działanie $t$ przesuwa numer $j$ o jeden.

Zapiszmy ciąg Mayera--Vietorisa nad
$\Lambda=\mathbb Z[t,t^{-1}]$:
\[
H_1(F)\otimes\Lambda\xrightarrow{\;t(i_+)_*-(i_-)_*\;}
H_1(Y)\otimes\Lambda\longrightarrow
H_1(\widetilde X_K)\longrightarrow
H_0(F)\otimes\Lambda\xrightarrow{\;t-1\;}
H_0(Y)\otimes\Lambda.
\]
Powierzchnia $F$ i rozcięta przestrzeń $Y$ są spójne.
Ponieważ $\Lambda$ jest dziedziną całkowitości, mnożenie
przez $t-1$ jest injekcją. Z dokładności otrzymujemy
$H_1(\widetilde X_K)\cong\operatorname{coker}
(t(i_+)_*-(i_-)_*)$.

Aby obliczyć macierz tej mapy, zauważmy, że $Y$ jest
homologicznie równoważne dopełnieniu powierzchni Seiferta.
Dualność Alexandera daje
$H_1(Y;\mathbb Z)\cong H^1(F;\mathbb Z)\cong\mathbb Z^{2g}$.
Wybierzmy bazę $[\alpha_1],\dots,[\alpha_{2g}]$ w $H_1(F)$,
a w $H_1(Y)$ bazę dualną, mierzoną za pomocą liczby splecenia
z tymi krzywymi. Macierze dwóch włożeń są wtedy macierzami
formy Seiferta $V$ i $V^T$ (zależnie od ustalenia numeracji
stron). Stąd macierz homomorfizmu wynosi $tV-V^T$.

Zmiana bazy nie zmienia modułu zdefiniowanego topologicznie.
Ponieważ $\det(V-V^T)=1$, wyznacznik $\det(tV-V^T)$
nie jest wielomianem zerowym. Zerowy ideał Fittinga
tego torsyjnego modułu jest zatem generowany przez
$\det(tV-V^T)$; inne bazy i powierzchnie dają ten sam
ideał. Jednostkami $\Lambda$ są dokładnie $\pm t^k$,
co kończy dowód niezależności wielomianu.
\end{proof}

\begin{stwierdzenie}[Własności wielomianu Alexandera]
Dla każdego węzła $K\subset S^3$ wielomian $\Delta_K$ spełnia:
\begin{enumerate}[label=(\roman*)]
\item $\Delta_K(1)=\pm1$ dla dowolnego reprezentanta, a w normalizacji
symetrycznej $\Delta_K(1)=1$;
\item $\Delta_K(t^{-1})=\pm t^m\Delta_K(t)$ dla pewnego $m\in\mathbb Z$;
\item $\det(K)=|\Delta_K(-1)|$, stąd $\det(K)$ jest dodatnią liczbą nieparzystą;
\item szerokość wykładników $\Delta_K$, czyli
$\deg_{\max}\Delta_K-\deg_{\min}\Delta_K$, nie przekracza $2g(K)$.
\end{enumerate}
\end{stwierdzenie}

\begin{proof}
Z twierdzenia o postaci przecięcia
$\Delta_K(1)=\det(V-V^T)=1$ dla przyjętej reprezentacji.
Przy $2g\times2g$ otrzymujemy
$\det(t^{-1}V-V^T)=t^{-2g}\det(V-tV^T)
=t^{-2g}\det(tV-V^T)$, gdyż wyznacznik jest niezmienniczy
przy transpozycji. Daje to symetrię. Wzór na wyznacznik
węzła wynika z identyfikacji macierzy Foxa przy $t=-1$
z odpowiednią prezentacją modułu Alexandera; jest to standardowa
redukcja relacji Wirtingera. Nieparzystość wynika z kongruencji
$\Delta_K(-1)\equiv\Delta_K(1)\equiv1\pmod2$.
Wreszcie macierz $tV-V^T$ ma rozmiar $2g$, więc jej
wyznacznik jest wielomianem stopnia co najwyżej $2g$
z nieujemnymi potęgami $t$. Minimalizując $g$ po wszystkich
powierzchniach, uzyskujemy ostatnią nierówność.
\end{proof}

\begin{przyklad}[Dwa najprostsze nietrywialne węzły]
Dla jednego z odbić trójlistnika można wybrać macierz
$V_T=\begin{psmallmatrix}1&0\\-1&1\end{psmallmatrix}$. Wtedy
\[
tV_T-V_T^T=
\begin{pmatrix}t-1&1\\-t&t-1\end{pmatrix},
\qquad
\Delta_T(t)=\det(tV_T-V_T^T)=(t-1)^2+t=t^2-t+1.
\]
Po symetrycznej normalizacji otrzymujemy
$\Delta_T^{\mathrm{sym}}(t)=t^{-1}-1+t$.
Ponieważ szerokość wykładników wynosi $2$, dostajemy
$g(T)\ge 1$. Z algorytmu Seiferta wiemy już, że $g(T)\le 1$,
a więc $g(T)=1$. Wyznacznik węzła jest równy
$\det(T)=|\Delta_T(-1)|=3$.

Ta sama macierz pozwala natychmiast odczytać sygnaturę.
Mamy bowiem
\[
V_T+V_T^T=
\begin{pmatrix}2&-1\\-1&2\end{pmatrix},
\]
której wartości własne wynoszą $1$ i $3$. Są dodatnie,
zatem $\sigma(T)=2$. W szczególności trójlistnik nie jest
amfichiralny, a z nierówności sygnaturowej otrzymujemy od razu
$g_4^{\mathrm{sm}}(T)\ge 1$.

Dla węzła ósemkowego można wziąć macierz
$V_8=\begin{psmallmatrix}1&0\\-1&-1\end{psmallmatrix}$. Otrzymujemy
\[
tV_8-V_8^T=
\begin{pmatrix}t-1&1\\-t&1-t\end{pmatrix},
\qquad
\Delta_{4_1}(t)=\det(tV_8-V_8^T)=-t^2+3t-1.
\]
Po normalizacji
$\Delta_{4_1}^{\mathrm{sym}}(t)=3-t-t^{-1}$,
stąd ponownie $g(4_1)=1$ oraz
$\det(4_1)=|\Delta_{4_1}(-1)|=5$.
Z kolei
\[
V_8+V_8^T=
\begin{pmatrix}2&-1\\-1&-2\end{pmatrix}
\]
ma wyznacznik ujemny, a więc jest nieokreślona; stąd
$\sigma(4_1)=0$. Oba węzły mają więc ten sam rodzaj,
ale różnią się wielomianem Alexandera, wyznacznikiem i sygnaturą.
\end{przyklad}

\begin{uwaga}[Ograniczenia wielomianu]
Równość $\Delta_K(t)=1$ nie pociąga za sobą trywialności węzła.
Istnieją nietrywialne węzły o trywialnym wielomianie Alexandera;
przykładowo odpowiednie satelity typu Whiteheada. Wielomian
Alexandera nie rozróżnia też węzła od jego odbicia lustrzanego.
Potrzebujemy zatem dalszych niezmienników.
\end{uwaga}

\begin{stwierdzenie}[Zmiana bazy i stabilizacja macierzy Seiferta]
Niech $K\subset S^3$ będzie węzłem i niech $V$ będzie
jego macierzą Seiferta. Przy zastąpieniu bazy w
$H_1(F;\mathbb Z)$ macierz przechodzi w $P^TVP$,
gdzie $P\in GL_{2g}(\mathbb Z)$. Przy elementarnym
powiększeniu powierzchni Seiferta macierz można, po
odpowiedniej zmianie bazy, zapisać jako
\[
V'=\begin{pmatrix}V&0&0\\ u^T&0&0\\v^T&1&0\end{pmatrix}
\quad\text{lub}\quad
V'=\begin{pmatrix}V&u&v\\0&0&1\\0&0&0\end{pmatrix},
\]
gdzie $u,v\in\mathbb Z^{2g}$. W obu przypadkach
$\det(tV'-{V'}^T)$ różni się od $\det(tV-V^T)$
co najwyżej przez czynnik będący jednostką w
$\mathbb Z[t,t^{-1}]$.
\end{stwierdzenie}

\begin{proof}
Przy zmianie bazy biliniowa forma Seiferta jest mnożona
przez $P^T$ i $P$, zatem
$\det(tP^TVP-P^TV^TP)=(\det P)^2\det(tV-V^T)
=\det(tV-V^T)$. Dla pierwszej macierzy $V'$ rozwijamy
$\det(tV'-{V'}^T)$ kolejno względem ostatniej kolumny
i ostatniego wiersza; otrzymujemy czynnik $t$ (z
możliwym znakiem zależnym od kolejności bazy) razy
$\det(tV-V^T)$. Dla drugiej postaci otrzymujemy
analogicznie czynnik $t$. Są one jednostkami w pierścieniu
wielomianów Laurenta.
Twierdzenie o $S$-równoważności stwierdza, że wszystkie macierze
Seiferta tego samego węzła można połączyć skończonym
ciągiem takich operacji i ich odwrotności. Ten ostatni
fakt jest niezależnym twierdzeniem topologicznym; powyższy
rachunek pokazuje, dlaczego wystarcza on do wykazania
niezmienniczości $\Delta_K$.
\end{proof}

\section{Rachunek Foxa i algorytm wyznaczania wielomianu Alexandera}
\label{sec:kt-fox-calculus}

Macierz Seiferta jest geometrycznym sposobem obliczania
wielomianu Alexandera. Jeżeli znamy jednak prezentację
Wirtingera, nie musimy konstruować powierzchni: wystarcza
specjalny rachunek różniczkowy w pierścieniu grupowym.
Metoda ta działa bezpośrednio na słowach w generatorach.

\begin{definicja}[Pochodne Foxa]
Niech $F_n=\langle x_1,\dots,x_n\rangle$ będzie wolną grupą,
a $\mathbb Z[F_n]$ jej całkowitym pierścieniem grupowym.
Dla $i\in\{1,\dots,n\}$ \emph{pochodną Foxa względem $x_i$}
nazywamy jedyne $\mathbb Z$-liniowe odwzorowanie
$\partial_i\colon\mathbb Z[F_n]\to\mathbb Z[F_n]$
spełniające $\partial_i(x_j)=\delta_{ij}$ i
$\partial_i(uv)=\partial_i(u)+u\partial_i(v)$
dla wszystkich $u,v\in F_n$. W szczególności
$\partial_i(1)=0$ oraz
$\partial_i(x_j^{-1})=-x_j^{-1}\delta_{ij}$.
\end{definicja}

Różnicą w stosunku do zwykłej pochodnej jest reguła
Leibniza: drugi składnik jest mnożony przez $u$, a nie
przez pochodną $u$ razy $v$. Operacja ta jest naturalną
pochodną w wolnym rachunku różniczkowym, dopasowaną
do algebry grupy.

\begin{lemat}[Fundamentalna tożsamość Foxa]
Dla każdego słowa $w\in F_n$ zachodzi równość
$w-1=\sum_{i=1}^n\partial_i(w)(x_i-1)$
w pierścieniu $\mathbb Z[F_n]$.
\end{lemat}

\begin{proof}
Dla generatora $w=x_j$ prawa strona wynosi $x_j-1$.
Dla odwrotności $w=x_j^{-1}$ mamy
$-x_j^{-1}(x_j-1)=x_j^{-1}-1$. Jeśli tożsamość
zachodzi dla $u$ i $v$, to dla $uv$ prawa strona wynosi
$\sum_i(\partial_i(u)+u\partial_i(v))(x_i-1)
=(u-1)+u(v-1)=uv-1$.
Indukcja względem długości słowa kończy dowód.
\end{proof}

\begin{definicja}[Macierz Alexandera z prezentacji grupy]
Niech $G=\langle x_1,\dots,x_n\mid r_1,\dots,r_m\rangle$
będzie skończoną prezentacją grupy, a
$\phi\colon G\to\mathbb Z=\langle t\rangle$ epimorfizmem.
Przez ten sam symbol oznaczamy indukowany homomorfizm
pierścieni $\mathbb Z[F_n]\to\mathbb Z[t,t^{-1}]$.
\emph{Macierzą Foxa--Alexandera} tej prezentacji i
homomorfizmu $\phi$ nazywamy macierz
$A=(\phi(\partial_jr_i))_{i=1,\dots,m;\ j=1,\dots,n}$
nad pierścieniem $\Lambda=\mathbb Z[t,t^{-1}]$.
\end{definicja}

\begin{twierdzenie}[Algorytm dla diagramu Wirtingera]
Niech $D$ będzie diagramem zorientowanego węzła
$K\subset S^3$ o $n\geq1$ skrzyżowaniach.
Z prezentacji Wirtingera wybierzmy $n$ generatorów
$x_1,\dots,x_n$ i $n-1$ niezależnych relacji
$r_1,\dots,r_{n-1}$. Niech
$\phi(x_j)=t$ dla wszystkich $j$. Zbudujmy macierz
$A=(\phi(\partial_jr_i))$ rozmiaru $(n-1)\times n$.
Po skreśleniu dowolnej jednej kolumny otrzymujemy
macierz kwadratową $A^{(j)}$, dla której
$\Delta_K(t)\doteq\det A^{(j)}$.
Jeżeli wybrana prezentacja ma dodatkowe generatory
lub relacje, stosujemy ogólną teorię ideałów
wyznacznikowych macierzy Alexandera.
\end{twierdzenie}

\begin{proof}
Niech $\Lambda=\mathbb Z[t,t^{-1}]$.
Standardowy kompleks dwuwymiarowy prezentacji Wirtingera
ma po podniesieniu do nakrycia cyklicznego łańcuchy
$C_2=\Lambda^{n-1}$, $C_1=\Lambda^n$ i $C_0=\Lambda$.
Operator $\partial_1$ jest dany wzorem
$\partial_1(u_1,\dots,u_n)=(t-1)\sum_{j=1}^n u_j$,
gdyż wszystkie meridiany przechodzą na $t$.
Ponieważ $t-1$ nie jest dzielnikiem zera w $\Lambda$,
$\ker\partial_1$ składa się z wektorów o sumie współrzędnych
równej zero. Usunięcie dowolnej współrzędnej daje więc
izomorfizm $\Lambda$-modułów
$\ker\partial_1\cong\Lambda^{n-1}$; brakującą współrzędną
odzyskuje się jako minus sumę pozostałych.

Macierz $\partial_2$, w konwencji kolumnowej, jest
transpozycją macierzy pochodnych Foxa. Istotnie, śledząc
brzeg podniesionej komórki odpowiadającej słowu
$r_i$, każdą literę $x_j$ odczytujemy z wagą równą obrazowi
prefiksu tego słowa w $\Lambda$. Jest to dokładnie reguła
Leibniza pochodnej Foxa; litera odwrotna wnosi wagę ujemną.
Fundamentalna tożsamość Foxa implikuje
$\partial_1\partial_2=0$. Dlatego po utożsamieniu
$\ker\partial_1$ z $\Lambda^{n-1}$ macierz
$\partial_2$ jest transpozycją $A^{(j)}$.

W pierwszej homologii mamy
$H_1(\widetilde X_K;\mathbb Z)=\ker\partial_1/
\operatorname{im}\partial_2$. Standardowy kompleks
Wirtingera odtwarza pierwszy typ homotopii dopełnienia,
więc oblicza tę samą pierwszą homologię co nakrycie.
Zatem moduł Alexandera ma prezentację $A^{(j)}$
(z dokładnością do transpozycji), a jego wielomianem
rzędu zerowego jest $\det A^{(j)}$. W połączeniu z
poprzednim twierdzeniem daje to
$\det A^{(j)}\doteq\Delta_K(t)$.
\end{proof}

\begin{przyklad}[Pełne obliczenie dla trójlistnika]
Weźmy prezentację
$G_T=\langle a,b\mid aba=bab\rangle$.
Relację zapisujemy jako słowo
$r=aba\,b^{-1}a^{-1}b^{-1}$.
Reguła Foxa daje
$\partial_a r=1+ab-aba\,b^{-1}a^{-1}$,
a po abelianizacji $\phi(a)=\phi(b)=t$ otrzymujemy
$\phi(\partial_a r)=1+t^2-t=t^2-t+1$.
Podobnie
$\partial_b r=a-aba\,b^{-1}-aba\,b^{-1}a^{-1}b^{-1}$,
więc $\phi(\partial_b r)=t-t^2-1$.
Macierz Alexandera jest zatem
$A=\begin{pmatrix}t^2-t+1&-(t^2-t+1)\end{pmatrix}$.
Usuwając dowolną kolumnę, otrzymujemy
$\Delta_T(t)\doteq t^2-t+1$, zgodnie z rachunkiem
z macierzy Seiferta. Widzimy, że metoda wymaga tylko
operacji na słowach grupowych i wielomianach Laurenta.
\end{przyklad}

\begin{uwaga}[Procedura obliczeniowa]
Dla praktycznych obliczeń można kolejno:
(1) zakodować orientowany diagram i jego łuki Wirtingera;
(2) zapisać relacje jako słowa w wolnej grupie;
(3) policzyć pochodne Foxa z reguł Leibniza;
(4) zastąpić każdy meridian zmienną $t$;
(5) skreślić kolumnę macierzy i policzyć wyznacznik;
(6) usunąć czynnik $\pm t^k$ przez symetryczną normalizację.
Kroki (3)--(5) są czysto algebraiczne i mogą być
zaimplementowane w systemie algebry komputerowej.
\end{uwaga}

\section{Wielomian Conwaya i niezmiennik Arfa}
\label{sec:kt-conway-arf}

\begin{definicja}[Wielomian Conwaya]
Dla każdego zorientowanego splotu $L\subset S^3$
przez $\nabla_L(z)\in\mathbb Z[z]$ oznaczamy
wielomian Conwaya, jednoznacznie scharakteryzowany
przez warunki $\nabla_U(z)=1$ oraz
$\nabla_{L_+}(z)-\nabla_{L_-}(z)=z\nabla_{L_0}(z)$
dla każdej zorientowanej trójki skein. Dla splotu
rozdzielnego o co najmniej dwóch składowych
$\nabla_L(z)=0$.
\end{definicja}

\begin{twierdzenie}[Conway--Alexander]
Dla dowolnego zorientowanego węzła $K\subset S^3$
wielomian Conwaya spełnia
$\Delta_K^{\mathrm{sym}}(t)=
\nabla_K(t^{1/2}-t^{-1/2})$,
gdzie $\Delta_K^{\mathrm{sym}}$ jest symetrycznie
znormalizowanym wielomianem Alexandera.
Ponadto $\nabla_K(z)$ zawiera wyłącznie parzyste
potęgi $z$, a jego wyraz stały wynosi jeden.
\end{twierdzenie}

\begin{proof}[Zarys indukcji skein]
Dla splotów wprowadzamy symetrycznie unormowany wielomian
Alexandera z uwzględnieniem liczby składowych; jego wartości
mogą należeć do $\mathbb Z[t^{1/2},t^{-1/2}]$. Relacja
Alexandera--Conwaya ma postać
$\Delta_{L_+}(t)-\Delta_{L_-}(t)=
(t^{1/2}-t^{-1/2})\Delta_{L_0}(t)$.
Można ją uzyskać bezpośrednio z dwóch macierzy Seiferta
różniących się lokalnym uchwytem w otoczeniu skrzyżowania
lub z odpowiednich macierzy prezentacji modułu Alexandera.
Po podstawieniu $z=t^{1/2}-t^{-1/2}$ przyjmuje ona dokładnie
postać relacji Conwaya.

Wyjaśnijmy jednoznaczność rekurencji. Dla diagramu splotu
wybieramy punkty bazowe na składowych i kolejność tych
składowych. Zmieniając skrzyżowania, w których po raz
pierwszy przechodzimy dołem podczas śledzenia diagramu,
uzyskujemy diagram splotu rozdzielnego. Każdą taką zmianę
wyrażamy relacją skein, przy czym wygładzenie zmniejsza
liczbę skrzyżowań. Indukcja po liczbie skrzyżowań, wraz
z wartością $1$ na węźle trywialnym i $0$ na splotach
rozdzielnych o co najmniej dwóch składowych, wyznacza
wielomian jednoznacznie. Dla węzła relacja symetrii
wielomianu Alexandera odpowiada niezmienniczości po
zamianie $z$ na $-z$, skąd znikają potęgi nieparzyste.
\end{proof}

\begin{definicja}[Forma kwadratowa i niezmiennik Arfa]
Niech $K\subset S^3$ będzie węzłem i niech
$F$ będzie jego powierzchnią Seiferta rodzaju $g$.
Na $H_1(F;\mathbb Z/2\mathbb Z)$ definiujemy
formę kwadratową $q_F([\gamma]):=
\operatorname{lk}(\gamma,\gamma^+)\pmod2$.
Jeżeli $(a_1,b_1,\dots,a_g,b_g)$ jest bazą symplektyczną
przestrzeni $H_1(F;\mathbb Z/2\mathbb Z)$ względem
modulo-dwójkowej formy przecięcia, kładziemy
$\operatorname{Arf}(K):=\sum_{i=1}^{g}q_F(a_i)q_F(b_i)
\pmod2\in\mathbb Z/2\mathbb Z$.
Niezależność od wyborów wykazujemy w następnym stwierdzeniu.
\end{definicja}

\begin{lemat}[Suma Gaussa dla formy kwadratowej nad $\mathbb F_2$]
\label{lem:tg-arf-gauss}
Niech $V$ będzie przestrzenią liniową wymiaru $2g$ nad
$\mathbb F_2$, a $\omega\colon V\times V\to\mathbb F_2$
niezdegenerowaną formą naprzemienną. Niech funkcja
$q\colon V\to\mathbb F_2$ spełnia
$q(x+y)=q(x)+q(y)+\omega(x,y)$ dla wszystkich $x,y\in V$.
Wówczas dla każdej bazy symplektycznej
$(a_1,b_1,\ldots,a_g,b_g)$ zachodzi
\[
\sum_{x\in V}(-1)^{q(x)}
=2^g(-1)^{\sum_{i=1}^g q(a_i)q(b_i)}.
\]
W szczególności suma $\sum_iq(a_i)q(b_i)\in\mathbb F_2$
nie zależy od wyboru bazy symplektycznej.
\end{lemat}

\begin{proof}
Każdy $x\in V$ ma jednoznaczny zapis
$x=\sum_i(\alpha_i a_i+\beta_i b_i)$, gdzie
$\alpha_i,\beta_i\in\mathbb F_2$. Ponieważ różne pary
$(a_i,b_i)$ są względem $\omega$ ortogonalne, a
$\omega(a_i,b_i)=1$, to z tożsamości kwadratowej wynika
\[
q(x)=\sum_{i=1}^g
\bigl(\alpha_i q(a_i)+\beta_i q(b_i)+\alpha_i\beta_i\bigr)
\pmod 2.
\]
Suma Gaussa jest zatem iloczynem $g$ sum dwuwymiarowych.
Dla $A,B\in\mathbb F_2$ mamy
\[
\sum_{\alpha,\beta\in\mathbb F_2}
(-1)^{\alpha A+\beta B+\alpha\beta}
=1+(-1)^A+(-1)^B-(-1)^{A+B}
=2(-1)^{AB}.
\]
Mnożąc te równości, dostajemy wzór. Lewa strona zależy
wyłącznie od $q$, a zatem wyznacza niezależnie od bazy
wykładnik po prawej stronie.
\end{proof}

\begin{stwierdzenie}[Niezmienniczość formy Arfa]
Dla każdego węzła $K\subset S^3$ i powierzchni Seiferta
$F$ funkcja $q_F$ spełnia
$q_F(x+y)=q_F(x)+q_F(y)+x\cdot y$ w $\mathbb F_2$,
gdzie $x\cdot y$ oznacza formę przecięcia na $H_1(F;\mathbb F_2)$.
Liczba $\sum_iq_F(a_i)q_F(b_i)$ jest niezależna od bazy
symplektycznej i powierzchni $F$.
\end{stwierdzenie}

\begin{proof}
Dla dwóch poprzecznych krzywych $\alpha,\beta\subset F$
reprezentant klasy $[\alpha]+[\beta]$ otrzymujemy przez
lokalne wygładzenie punktów przecięcia. Przy obliczaniu
samosplecenia pojawiają się składniki diagonalne
$q_F([\alpha])$, $q_F([\beta])$ oraz suma dwóch składników
mieszanych. Modulo dwa te ostatnie są równe algebraicznej
liczbie przecięć krzywych, stąd wzór kwadratowy.

Niezależność sumy $\sum_iq_F(a_i)q_F(b_i)$ od bazy
symplektycznej wynika dokładnie z lematu
\ref{lem:tg-arf-gauss}: jego lewa strona jest określona
bez użycia jakiejkolwiek bazy.
Przy zmianie powierzchni Seiferta macierze są związane
$S$-równoważnością. Elementarne powiększenie dodaje parę
hiperboliczną, dla której składnik formy Arfa wynosi zero.
Zmiany bazy i takie powiększenia nie zmieniają więc wyniku.
\end{proof}

\begin{stwierdzenie}
Niech $K\subset S^3$ będzie węzłem, a
$\nabla_K(z)=1+a_2z^2+a_4z^4+\cdots$.
Wtedy $\operatorname{Arf}(K)\equiv a_2\pmod2$.
Ponadto jeśli $K$ jest gładko slice, to
$\operatorname{Arf}(K)=0$.
\end{stwierdzenie}

\begin{uwaga}
Zerowanie niezmiennika Arfa jest warunkiem koniecznym
slicowości, ale nie wystarczającym. Niezmiennik ten
jest znacznie słabszy niż $\tau$ czy $s$, lecz
ma ważne znaczenie w teorii form kwadratowych
nad ciałem dwuelementowym i w czterowymiarowej
teorii chirurgii.
\end{uwaga}

\begin{przyklad}
Dla trójlistnika
$\Delta_T^{\mathrm{sym}}(t)=t+t^{-1}-1$.
Ponieważ
$(t^{1/2}-t^{-1/2})^2=t+t^{-1}-2$,
mamy $\nabla_T(z)=1+z^2$ i
$\operatorname{Arf}(T)=1$.
Dla węzła ósemkowego
$\Delta_{4_1}^{\mathrm{sym}}(t)=3-t-t^{-1}$,
więc $\nabla_{4_1}(z)=1-z^2$;
także $\operatorname{Arf}(4_1)=1$.
\end{przyklad}


\section{Sygnatura węzła i ograniczenia rodzaju}
\label{sec:kt-sygnatura}

\begin{definicja}[Sygnatura klasyczna]
Dla węzła $K\subset S^3$ i macierzy Seiferta $V$ definiujemy
$\sigma(K):=\operatorname{sign}(V+V^T)$,
gdzie $\operatorname{sign}$ oznacza liczbę dodatnich minus liczbę
ujemnych wartości własnych rzeczywistej macierzy symetrycznej,
liczonych z krotnościami. Wartość nie zależy od wyboru
powierzchni i bazy.
\end{definicja}

Niezależność od powierzchni wynika z $S$-równoważności macierzy
Seiferta: zmiana bazy zachowuje sygnaturę przez prawo
bezwładności Sylvestera, a elementarne powiększenia odpowiadają
dodaniu składnika hiperbolicznego o sygnaturze zero.

\begin{stwierdzenie}
Niech $K\subset S^3$ będzie węzłem. Wówczas
$|\sigma(K)|\leq2g(K)$.
Odbicie lustrzane zmienia znak sygnatury:
$\sigma(\overline K)=-\sigma(K)$.
\end{stwierdzenie}

\begin{proof}
Macierz $V+V^T$ ma rozmiar $2g(F)$, a sygnatura symetrycznej
macierzy wymiaru $2g(F)$ ma wartość bezwzględną nie większą
od $2g(F)$. Minimalizując po powierzchniach, dostajemy pierwszą
nierówność. Przy odbiciu lustrzanym macierz Seiferta zmienia się,
z dokładnością do odpowiedniej zmiany bazy, na $-V^T$;
symetryzacja zmienia znak, co daje drugą równość.
\end{proof}

\begin{przyklad}
Dla $V_T=\begin{psmallmatrix}1&0\\-1&1\end{psmallmatrix}$
macierz $V_T+V_T^T=\begin{psmallmatrix}2&-1\\-1&2\end{psmallmatrix}$
ma dwie dodatnie wartości własne, więc $\sigma(T)=2$
w tej konwencji wyboru odbicia trójlistnika. Jego odbicie
ma sygnaturę $-2$. W szczególności trójlistnik nie jest
amfichiralny. Dla $V_8$ macierz symetryczna ma jedną wartość
dodatnią i jedną ujemną, a więc $\sigma(4_1)=0$.
\end{przyklad}

\begin{uwaga}[Sygnatury Levine'a--Tristrama]
Pełniejszy układ niezmienników uzyskuje się z hermitowskich
macierzy $H_\omega:=(1-\omega)V+(1-\overline\omega)V^T$
dla $\omega\in S^1\setminus\{1\}$. Ich sygnatury, gdy macierz
jest niezdegenerowana, definiują funkcję sygnatur węzła.
Dla $\omega=-1$ otrzymujemy
$H_{-1}=2(V+V^T)$, a więc sygnaturę klasyczną.
Przy pierwiastkach wielomianu Alexandera na okręgu jednostkowym
mogą pojawiać się skoki funkcji sygnatur.
\end{uwaga}

\section{Wielomian Jonesa i nawias Kauffmana}
\label{sec:kt-jones}

Wielomian Jonesa jest pierwszym przykładem niezmiennika zbudowanego
wprost z lokalnych reguł rozwiązywania skrzyżowań, który wykracza
poza klasyczną teorię modułu Alexandera. Jego odkrycie wyrosło
z teorii algebr operatorowych, w szczególności z badań podczynników
algebr von Neumanna.

\begin{definicja}[Nawias Kauffmana]
Każdemu niezorientowanemu diagramowi splotu $D$ przyporządkowujemy
wielomian Laurenta $\langle D\rangle\in\mathbb Z[A,A^{-1}]$
rekurencyjnie przez trzy reguły:
\begin{enumerate}[label=(\roman*)]
\item $\langle\bigcirc\rangle=1$ dla diagramu pojedynczego okręgu;
\item $\langle D\sqcup\bigcirc\rangle=
(-A^2-A^{-2})\langle D\rangle$ dla diagramu zawierającego
co najmniej jedną składową;
\item dla lokalnego skrzyżowania
$\langle D_{\times}\rangle
=A\langle D_{\mathrm A}\rangle+A^{-1}\langle D_{\mathrm B}\rangle$,
gdzie $D_{\mathrm A},D_{\mathrm B}$ oznaczają dwa lokalne
wygładzenia w przedstawionej poniżej konwencji. W pokazanym
skrzyżowaniu pasmo biegnące z lewego górnego końca do prawego
dolnego przechodzi górą; wygładzenie A łączy lewy górny
koniec z lewym dolnym oraz prawy górny z prawym dolnym,
a wygładzenie B łączy oba końce górne i oba dolne.
Jeżeli zmienimy, które pasmo przechodzi górą, oznaczenia
A i B zamieniamy miejscami. Reguła nie zależy od orientacji
pasm i jest zgodna z obrotem lokalnego diagramu.
\end{enumerate}
Dla diagramu bez skrzyżowań o $m\geq1$ składowych otrzymujemy
$\langle D\rangle=(-A^2-A^{-2})^{m-1}$.
\end{definicja}

\begin{center}
\begin{tikzpicture}[scale=.86,line cap=round,line join=round,>=Stealth]
\begin{scope}[shift={(-4.0,0)}]
  \draw[ktBlue,line width=2.8pt] (-.73,1)--(.73,-1);
  \draw[ktCoral,line width=2.8pt] (.73,1)--(.13,.18);
  \draw[ktCoral,line width=2.8pt] (-.13,-.18)--(-.73,-1);
  \node[text=ktInk,font=\small] at (0,-1.4) {skrzyżowanie};
\end{scope}
\node[text=ktGold,font=\large\bfseries] at (-2.15,0) {$A$};
\begin{scope}
  \draw[ktTeal,line width=2.8pt] (-.73,-1)--(-.73,1);
  \draw[ktTeal,line width=2.8pt] (.73,-1)--(.73,1);
  \node[text=ktInk,font=\small] at (0,-1.4) {wygładzenie A};
\end{scope}
\node[text=ktGold,font=\large\bfseries] at (2.14,0) {$A^{-1}$};
\begin{scope}[shift={(4.0,0)}]
  \draw[ktPurple,line width=2.8pt] (-.73,1).. controls (-.73,.34) and (.73,.34) .. (.73,1);
  \draw[ktPurple,line width=2.8pt] (-.73,-1).. controls (-.73,-.34) and (.73,-.34) .. (.73,-1);
  \node[text=ktInk,font=\small] at (0,-1.4) {wygładzenie B};
\end{scope}
\end{tikzpicture}
\end{center}

Definicja jest wykonalna: przy $n$ skrzyżowaniach rozwijamy
rekurencję do najwyżej $2^n$ stanów wygładzonych. Dla każdego
stanu $s$ niech $a(s)$ oznacza liczbę wygładzeń typu A,
$b(s)$ liczbę wygładzeń typu B, a $|s|$ liczbę okręgów po
wygładzeniu. Wówczas
\[
\langle D\rangle=\sum_s A^{a(s)-b(s)}
(-A^2-A^{-2})^{|s|-1}.
\]
To pierwszy całkowicie jawny wzór obliczeniowy.

\begin{stwierdzenie}
Nawias Kauffmana jest niezmienny przy ruchach Reidemeistera
$\mathrm{R II}$ i $\mathrm{R III}$. Przy dodatnim lub ujemnym
ruchu $\mathrm{R I}$ zostaje pomnożony odpowiednio przez
$-A^3$ lub $-A^{-3}$, przy zgodnej konwencji znaku skrzyżowania.
\end{stwierdzenie}

\begin{proof}
Sprawdzenie pierwszego ruchu polega na jednorazowym rozwinięciu
skrzyżowania i użyciu relacji dla dodatkowego okręgu:
$A(-A^2-A^{-2})+A^{-1}=-A^3$.
Dla przeciwnego skrętu analogiczny rachunek daje $-A^{-3}$.
Dla ruchu $\mathrm{R II}$ rozwijamy dwa skrzyżowania. Cztery
stany wygładzenia grupują się w dwa typy połączenia końców;
współczynniki upraszczają się przy
$\delta=-A^2-A^{-2}$, gdyż $A^2+\delta+A^{-2}=0$.
W ruchu $\mathrm{R III}$ rozwinięcie jednego z trzech
skrzyżowań po obu stronach redukuje sprawdzenie do równości
dwóch diagramów związanych ruchami płaskimi i $\mathrm{R II}$.
Są to skończone rachunki lokalne, niezależne od reszty diagramu;
przy szczegółowym rozwinięciu trzeba zachować tę samą konwencję
wygładzeń A i B na wszystkich obrazach lokalnych.
\end{proof}

\begin{definicja}[Wielomian Jonesa]
Niech $D$ będzie zorientowanym diagramem węzła lub splotu
$L\subset S^3$, a $w(D)$ jego skrętem. Kładziemy
$f_D(A):=(-A^3)^{-w(D)}\langle D\rangle$.
Następnie definiujemy $V_L(t):=f_D(t^{-1/4})$,
rozumiejąc wynik jako wielomian Laurenta w $t^{1/2}$.
Dla węzła przy typowej normalizacji wykładniki są całkowite;
dla splotów mogą występować połówkowe.
\end{definicja}

\begin{twierdzenie}
Wielomian $V_L(t)$ nie zależy od diagramu zorientowanego
splotu $L\subset S^3$. W szczególności $V_U(t)=1$.
Przy konsekwentnie ustalonych konwencjach skrzyżowań spełnia
relację skein
$t^{-1}V_{L_+}(t)-tV_{L_-}(t)=
(t^{1/2}-t^{-1/2})V_{L_0}(t)$,
gdzie $L_+,L_-,L_0$ są jednakowymi poza dyskiem diagramami:
$L_+$ i $L_-$ różnią się znakiem skrzyżowania, zaś $L_0$
powstaje przez zorientowane wygładzenie.
\end{twierdzenie}

\begin{proof}
Wszystkie trzy ruchy Reidemeistera zachowują $f_D$.
Dla $\mathrm{R II}$ i $\mathrm{R III}$ nie zmienia się ani
nawias, ani skręt. Dla $\mathrm{R I}$ mnożnik $-A^{\pm3}$
w nawiasie skraca się z przyrostem $(-A^3)^{-w(D)}$.
Twierdzenie Reidemeistera zapewnia niezależność od diagramu.
Relację skein otrzymujemy, rozwijając lokalnie skrzyżowanie
dla $L_+$ i $L_-$ oraz eliminując oba wygładzenia z
uwzględnieniem różnicy ich skrętów. Po podstawieniu
$t=A^{-4}$ otrzymujemy wskazaną tożsamość.
\end{proof}

\begin{uwaga}[Orientacje i konwencje]
W literaturze spotyka się także zmienną $t^{-1}$ zamiast $t$,
co zamienia wartości dla odbić lustrzanych. Zawsze należy
sprawdzać przyjętą relację skein i normalizację węzła
trywialnego przed porównaniem tablic wielomianów.
\end{uwaga}

\begin{stwierdzenie}
Niech $K\subset S^3$ będzie węzłem, a $\overline K$ jego
odbiciem lustrzanym. Wówczas
$V_{\overline K}(t)=V_K(t^{-1})$.
\end{stwierdzenie}

\begin{proof}
Odbicie lustrzane wymienia wygładzenia A i B, więc
$\langle\overline D\rangle(A)=\langle D\rangle(A^{-1})$.
Jednocześnie $w(\overline D)=-w(D)$.
Po normalizacji i zamianie $A$ na $t^{-1/4}$ dostajemy wynik.
\end{proof}

\begin{przyklad}[Obliczenie wielomianu Jonesa trójlistnika]
Trójlistnik można przedstawić jako domknięcie dodatniego warkocza
$\sigma_1^3$. W dwupasmowej algebrze Temperleya--Lieba oznaczmy przez
$1$ diagram identyczności, a przez $e$ element z jednym połączeniem
między pasmami. Dla generatora warkoczowego mamy wzór
$b=A1+A^{-1}e$, relację $e^2=\delta e$ oraz
$\delta=-A^2-A^{-2}$.
Rozwijając $b^3$, dostajemy
\[
b^3=A^3 1+(A-A^{-3}+A^{-7})e.
\]
Po domknięciu warkocza diagram $1$ daje dwa okręgi, natomiast
$e$ daje jeden okrąg. Z definicji nawiasu Kauffmana wynika zatem
\[
\langle\widehat{\sigma_1^3}\rangle
=A^3\delta+(A-A^{-3}+A^{-7})
=-A^5-A^{-3}+A^{-7}.
\]
Skręt standardowego dodatniego diagramu trójlistnika wynosi $w=3$,
więc po normalizacji otrzymujemy
\[
f(A)=(-A^3)^{-3}\langle\widehat{\sigma_1^3}\rangle
=A^{-4}+A^{-12}-A^{-16}.
\]
Po podstawieniu $t=A^{-4}$ dostajemy
\[
V_T(t)=t+t^3-t^4.
\]
Dla odbicia lustrzanego dostajemy z poprzedniego stwierdzenia
$V_{\overline T}(t)=t^{-1}+t^{-3}-t^{-4}$. Wielomian Jonesa
rozróżnia więc oba odbicia trójlistnika.
\end{przyklad}

\section{Warkocze i grupy warkoczowe}
\label{sec:kt-warkocze}

\begin{definicja}[Warkocz geometryczny]
Niech $n\geq1$. \emph{Warkoczem o $n$ pasmach} nazywamy
zbiór $n$ parami rozłącznych, gładkich łuków w
$\mathbb R^2\times[0,1]$, których początki są ustalonymi
$n$ punktami w $\mathbb R^2\times\{0\}$, końce stanowią
permutację tych punktów w $\mathbb R^2\times\{1\}$,
a rzutowanie każdego łuku na współrzędną $[0,1]$
jest dyfeomorfizmem. Warkocze utożsamiamy z dokładnością
do izotopii zachowującej końce i warunek monotoniczności.
Składanie przez sklejanie końców i przeskalowanie wysokości
tworzy \emph{grupę warkoczową} $B_n$.
\end{definicja}

\begin{twierdzenie}[Prezentacja Artina]
Dla $n\geq2$ grupa warkoczowa $B_n$ ma prezentację
\[
B_n=\Big\langle\sigma_1,\ldots,\sigma_{n-1}\ \Big|\ 
\sigma_i\sigma_j=\sigma_j\sigma_i\ (|i-j|\geq2),\quad
\sigma_i\sigma_{i+1}\sigma_i=\sigma_{i+1}\sigma_i\sigma_{i+1}
\ (1\leq i\leq n-2)\Big\rangle.
\]
Generator $\sigma_i$ oznacza dodatnie skrzyżowanie pasm
o numerach $i$ oraz $i+1$; $\sigma_i^{-1}$ oznacza
skrzyżowanie przeciwne.
\end{twierdzenie}

\begin{proof}[Idea dowodu i wykorzystywane narzędzie]
Każdy warkocz można zaburzyć tak, aby jego rzut na płaszczyznę
poziomej i pionowej współrzędnej rzutowania miał skończenie wiele
skrzyżowań rozdzielonych wysokością. Odczytanie ich w kolejności
daje słowo w generatorach $\sigma_i^{\pm1}$.
W rodzinie izotopii ogólnego położenia zmiany słowa pochodzą
z odwrócenia sąsiednich skrzyżowań, jeśli dotyczą rozłącznych
par pasm, oraz z lokalnego ruchu trzech pasm. Dają one
dokładnie relacje komutacji i relację Artina. Szczegółowe
uzasadnienie zgodności z izotopią względną względem końców
wymaga argumentu o transwersalności.
\end{proof}

\begin{definicja}[Domknięcie warkocza]
Dla $\beta\in B_n$ \emph{domknięciem} $\widehat\beta\subset S^3$
nazywamy splot otrzymany przez połączenie każdego końca
górnego z odpowiadającym mu początkiem dolnym rozłącznymi
łukami biegnącymi poza cylindrem zawierającym warkocz.
\end{definicja}

\begin{twierdzenie}[Alexander]
Każdy zorientowany splot $L\subset S^3$ jest równoważny
domknięciu pewnego warkocza $\beta\in B_n$ dla pewnego $n\geq1$.
\end{twierdzenie}

\begin{proof}[Idea dowodu]
Ustalamy oś w $\mathbb R^3$ rozłączną z $L$ i rozważamy
współrzędną kątową wokół osi. Jeżeli podczas śledzenia splotu
kąt czasem maleje, zastępujemy odpowiednie fragmenty łukami
przebiegającymi po dodatniej stronie osi w taki sposób,
aby stał się lokalnie rosnący. Ta operacja jest wykonywana
skończenie wiele razy i nie zmienia izotopijnego typu splotu.
Po jej zakończeniu współrzędna kątowa jest regularna
na każdej składowej i definiuje zamknięty warkocz.
\end{proof}

\begin{twierdzenie}[Markow]
Niech $\alpha\in B_m$ i $\beta\in B_n$ będą warkoczami.
Ich domknięcia $\widehat\alpha$ i $\widehat\beta$ są
równoważnymi zorientowanymi splotami wtedy i tylko wtedy,
gdy można przejść od jednego warkocza do drugiego skończonym
ciągiem następujących operacji i ich odwrotności:
\begin{enumerate}[label=(\roman*)]
\item sprzęganie $\gamma\mapsto\eta\gamma\eta^{-1}$ wewnątrz $B_k$;
\item stabilizacja $\gamma\in B_k\mapsto\gamma\sigma_k^{\pm1}
\in B_{k+1}$, przy naturalnym włożeniu $B_k\hookrightarrow B_{k+1}$.
\end{enumerate}
\end{twierdzenie}

\begin{uwaga}
Twierdzenia Alexandera i Markowa są podstawą algebraicznych
konstrukcji niezmienników splotów: szukamy reprezentacji
$\rho_n\colon B_n\to\mathrm{GL}(V_n)$ i funkcjonałów typu
śladu, które są zgodne ze sprzęganiem i stabilizacją.
Tak właśnie powstają wielomian Jonesa i pokrewne niezmienniki
w teorii reprezentacji algebr Temperleya--Lieba i Hecke'go.
\end{uwaga}

\begin{center}
\begin{tikzpicture}[line cap=round,scale=.95]
  \begin{scope}[shift={(-2.5,0)}]
    \fill[ktBlue!6,rounded corners=5pt] (-1.35,-.7) rectangle (1.35,2.85);
    \draw[draw=ktBlue,line width=2.4pt] (-0.45,2.45) .. controls (-1.45,2.65) and (-1.45,0.0) .. (-0.45,0.0);
    \draw[draw=ktCoral,line width=2.4pt] ( 0.45,2.45) .. controls ( 1.45,2.65) and ( 1.45,0.0) .. ( 0.45,0.0);
    \foreach \yy in {2.4,1.6,0.8}{
      \draw[draw=ktBlue,line width=2.4pt] (-0.45,\yy)--(0.45,\yy-0.65);
      \draw[draw=ktCoral,line width=2.4pt] ( 0.45,\yy)--(0.12,\yy-0.24);
      \draw[draw=ktCoral,line width=2.4pt] (-0.12,\yy-0.41)--(-0.45,\yy-0.65);
    }
    \foreach \yt/\yb in {2.45/2.4,1.75/1.6,0.95/0.8}{
      \draw[ktTeal,line width=2.4pt] (-0.45,\yb)--(-0.45,\yt);
      \draw[ktGold,line width=2.4pt] ( 0.45,\yb)--( 0.45,\yt);
    }
    \node at (0,-.47) {$\widehat{\sigma_1^3}$: trójlistnik};
  \end{scope}
  \begin{scope}[shift={(2.4,0)}]
    \fill[ktCoral!6,rounded corners=5pt] (-1.35,-.7) rectangle (1.35,2.85);
    \draw[ktPurple,line width=2.4pt] (-0.45,1.65) .. controls (-1.3,1.95) and (-1.3,0.0) .. (-0.45,0.0);
    \draw[ktCoral,line width=2.4pt] ( 0.45,1.65) .. controls ( 1.3,1.95) and ( 1.3,0.0) .. ( 0.45,0.0);
    \foreach \yy in {1.6,0.8}{
      \draw[ktPurple,line width=2.4pt] (-0.45,\yy)--(0.45,\yy-0.65);
      \draw[ktCoral,line width=2.4pt] (0.45,\yy)--(0.12,\yy-0.24);
      \draw[ktCoral,line width=2.4pt] (-0.12,\yy-0.41)--(-0.45,\yy-0.65);
    }
    \foreach \yt/\yb in {1.65/1.6,0.95/0.8}{
      \draw[ktTeal,line width=2.4pt] (-0.45,\yb)--(-0.45,\yt);
      \draw[ktGold,line width=2.4pt] (0.45,\yb)--(0.45,\yt);
    }
    \node at (0,-.47) {$\widehat{\sigma_1^2}$: splot Hopfa};
  \end{scope}
\end{tikzpicture}
\end{center}

\begin{przyklad}[Trójlistnik jako zamknięty warkocz]
Domknięcie warkocza $\sigma_1^3\in B_2\cong\mathbb Z$
jest jednym z dwóch trójlistników. Domknięcie
$\sigma_1^{-3}$ jest jego odbiciem lustrzanym.
Dla $\sigma_1^{2k+1}$ otrzymujemy węzeł torusowy
$T(2,2k+1)$, z orientacją zależną od konwencji znaków.
\end{przyklad}

\section{Suma spójna, węzły pierwsze i rozkład Schuberta}
\label{sec:kt-suma}

\begin{definicja}[Suma spójna]
Niech $K_1,K_2\subset S^3$ będą zorientowanymi węzłami.
Wybieramy dwie rozłączne kule $B_1,B_2$, których przecięcia
z odpowiednimi węzłami są trywialnymi, właściwie zanurzonymi
łukami. Usuwamy wnętrza tych łuków i łączymy ich cztery końce
dwoma rozłącznymi łukami zgodnie z orientacjami. Otrzymany
węzeł oznaczamy $K_1\#K_2$ i nazywamy \emph{sumą spójną}.
Węzeł jest \emph{pierwszy}, jeśli jest nietrywialny i z równości
$K=K_1\#K_2$ wynika, że jeden ze składników jest trywialny.
\end{definicja}

\begin{stwierdzenie}
Suma spójna jest dobrze określona z dokładnością do izotopii,
łączna, przemienna i ma element neutralny $U$.
Ponadto $\Delta_{K_1\#K_2}(t)=\Delta_{K_1}(t)\Delta_{K_2}(t)$
z dokładnością do jednostki $\pm t^k$, a
$\sigma(K_1\#K_2)=\sigma(K_1)+\sigma(K_2)$.
\end{stwierdzenie}

\begin{proof}
Niezależność od punktów połączenia wynika z możliwości
przesuwania krótkich łuków wzdłuż zorientowanych składników.
Łączność i przemienność uzyskujemy przez zmianę położenia
kul łączących. Z powierzchni Seiferta $F_1,F_2$ budujemy
powierzchnię $F_1\mathbin{\natural}F_2$ przez sumę brzegową.
W odpowiedniej bazie jej macierz Seiferta jest blokowo diagonalna:
$V=V_1\oplus V_2$. Otrzymujemy zatem
$\det(tV-V^T)=\det(tV_1-V_1^T)\det(tV_2-V_2^T)$
i analogiczną addytywność sygnatury symetryzacji.
\end{proof}

\begin{twierdzenie}[Schubert; rozkład na czynniki pierwsze]
Każdy nietrywialny węzeł $K\subset S^3$ można przedstawić
jako skończoną sumę spójną węzłów pierwszych. Dwa takie
rozkłady mają tę samą liczbę czynników, a po ich odpowiednim
przestawieniu czynniki odpowiadające sobie są równoważne.
\end{twierdzenie}

\begin{uwaga}
Dowód twierdzenia Schuberta wykorzystuje sfery dekomponujące,
które przecinają węzeł w dwóch punktach, i analizę ich
wzajemnego położenia. Istnienie rozkładu i jego skończoność można uzasadniać
za pomocą teorii sfer dekomponujących i twierdzeń o skończoności
dla powierzchni istotnych w trójrozmaitościach; jednoznaczność
wymaga analizy przecięć takich sfer. Nie jest to
bezpośredni wniosek z jedyności faktoryzacji wielomianów
Alexandera, ponieważ różne węzły mogą mieć ten sam wielomian.
\end{uwaga}

\begin{twierdzenie}[Addytywność rodzaju; Schubert]
Dla dowolnych węzłów $K_1,K_2\subset S^3$ zachodzi
$g(K_1\#K_2)=g(K_1)+g(K_2)$.
\end{twierdzenie}

\begin{proof}[Redukcja do rozcięcia minimalnej powierzchni]
Biorąc powierzchnie Seiferta $F_1,F_2$ minimalnego rodzaju
oraz ich sumę brzegową $F_1\mathbin{\natural}F_2$,
uzyskujemy powierzchnię o jednym brzegu i
$\chi(F_1\mathbin{\natural}F_2)=\chi(F_1)+\chi(F_2)-1$.
Ponieważ $\chi(F_i)=1-2g(F_i)$, otrzymujemy
$g(K_1\#K_2)\le g(K_1)+g(K_2)$.

Niech teraz $F$ będzie powierzchnią Seiferta minimalnego rodzaju
dla $K_1\#K_2$. Wybierzmy sferę dekomponującą
$S\subset S^3$, która spotyka węzeł dokładnie w dwóch
punktach i rozdziela odpowiednie składniki. Ustawiamy $F$
i $S$ poprzecznie. Powierzchnia $F\cap S$ jest sumą
okręgów oraz właściwie zanurzonych łuków z końcami w dwóch
punktach $S\cap\partial F$; jedyny taki łuk musi łączyć
te punkty.

Eliminujemy składowe okręgowe metodą najbardziej wewnętrznego
dysku. Jeśli okrąg $C\subset F\cap S$ ogranicza dysk na
$F$, wybieramy dysk najbardziej wewnętrzny i usuwamy
$C$ przez izotopię. Gdy pozostałby okrąg istotny na $F$,
to rozcinając sferę $S$ wzdłuż najbardziej wewnętrznego
okręgu na $S$ i stosując nieściśliwość powierzchni
minimalnego rodzaju w zewnętrzu węzła, uzyskalibyśmy
sprzeczność. Wykorzystaliśmy tu standardowy lemat:
\emph{powierzchnia Seiferta minimalnego rodzaju jest
nieściśliwa w zewnętrzu węzła}. Jego dowód polega na
kompresji: istotny dysk kompresujący zmniejsza rodzaj,
nie zmieniając zadanej składowej brzegu.

Po usunięciu okręgów $F\cap S$ pozostaje jeden łuk.
Rozcinamy $F$ wzdłuż tego łuku, a brzegi powstałe na
sferze zaślepiamy odpowiednimi dyskami po obu stronach
$S$. Dostajemy dwie powierzchnie Seiferta $F_1',F_2'$
dla $K_1,K_2$. Rozcięcie wzdłuż łuku zwiększa
charakterystykę Eulera o jeden:
$\chi(F_1')+\chi(F_2')=\chi(F)+1$.
Ponieważ wszystkie trzy powierzchnie są spójne i mają
po jednej składowej brzegu, daje to
$g(F)=g(F_1')+g(F_2')\ge g(K_1)+g(K_2)$.
W połączeniu z poprzednią nierównością kończy to dowód.
\end{proof}

\section{Węzły torusowe, włóknistość i rodzaj}
\label{sec:kt-torusowe}

\begin{twierdzenie}[Grupa węzła torusowego]
Niech $p,q\geq2$ będą względnie pierwsze. Dla węzła torusowego
$T(p,q)\subset S^3$ zachodzi
$\pi_1(S^3\setminus T(p,q))\cong
\langle a,b\mid a^p=b^q\rangle$.
\end{twierdzenie}

\begin{proof}
Standardowy torus rozdziela $S^3$ na dwa torusy pełne
$V_1,V_2$. Węzeł $T(p,q)$ jest istotną krzywą na ich wspólnym
brzegu. Po wycięciu wąskiego otoczenia tej krzywej
uzupełnienie węzła jest sumą dwóch torusów pełnych sklejonych
wzdłuż pierścienia. Rdzenie torusów pełnych generują
$\pi_1(V_1)\cong\mathbb Z$ i $\pi_1(V_2)\cong\mathbb Z$.
Pierścień sklejający biegnie odpowiednio $p$ i $q$ razy
wokół tych rdzeni, zatem twierdzenie van Kampena daje
prezentację z relacją $a^p=b^q$.
\end{proof}

\begin{twierdzenie}[Rodzaj węzła torusowego]
Jeżeli $p,q\geq2$ są względnie pierwsze, to
$g(T(p,q))=(p-1)(q-1)/2$.
\end{twierdzenie}

\begin{proof}
Dodatni warkocz $\beta=(\sigma_1\dots\sigma_{p-1})^q$
ma domknięcie równe $T(p,q)$. Ma $(p-1)q$ skrzyżowań
i $p$ okręgów Seiferta, więc algorytm Seiferta daje
powierzchnię rodzaju $(1-p+(p-1)q)/2=(p-1)(q-1)/2$.
Dla oszacowania z dołu obliczamy z prezentacji grupy
węzła (lub z odpowiedniej macierzy Seiferta) wielomian
\[
\Delta_{T(p,q)}(t)=
\frac{(t^{pq}-1)(t-1)}{(t^p-1)(t^q-1)}
\]
z dokładnością do jednostki $\pm t^k$.
Licznik i mianownik skracają się do wielomianu o szerokości
wykładników $(p-1)(q-1)$, ponieważ po redukcji stopień
wynosi $pq+1-p-q$. Z nierówności
$\operatorname{span}\Delta_K\leq2g(K)$ otrzymujemy
$g(T(p,q))\geq(p-1)(q-1)/2$.
\end{proof}

\begin{definicja}[Węzeł włóknisty]
Węzeł $K\subset S^3$ nazywamy \emph{włóknistym}, jeżeli jego
uzupełnienie $X_K=S^3\setminus\operatorname{int}N(K)$
jest przestrzenią wiązki nad $S^1$ o włóknie będącym zwartą,
spójną powierzchnią $F$ z jednym brzegiem, przy czym
$\partial F$ odpowiada preferowanej długości na
$\partial X_K$. Dla pewnego dyfeomorfizmu
$\varphi\colon F\to F$ identycznego w pobliżu brzegu mamy
$X_K\cong (F\times[0,1])/(x,1)\sim(\varphi(x),0)$.
Dyfeomorfizm $\varphi$ nazywamy \emph{monodromią}.
\end{definicja}

\begin{twierdzenie}
Jeżeli $K\subset S^3$ jest węzłem włóknistym, to jego
wielomian Alexandera jest jednomianowy na obu krańcach,
tzn. współczynniki przy najwyższej i najniższej potędze
$t$ są równe $\pm1$. Każdy węzeł torusowy jest włóknisty.
\end{twierdzenie}

\begin{proof}[Uzasadnienie pierwszej tezy]
Dla węzła włóknistego nakrycie nieskończenie cykliczne
uzupełnienia jest homotopijnie równoważne włóknu $F$.
Na $H_1(F;\mathbb Z)$ działanie generatora przekształceń
pokładowych jest dane przez automorfizm całkowity
$\varphi_*\in GL_{2g}(\mathbb Z)$. Wielomian Alexandera
jest, z dokładnością do jednostki, wielomianem
charakterystycznym tego automorfizmu:
$\Delta_K(t)\doteq\det(tI-\varphi_*)$.
Jego najwyższy współczynnik wynosi jeden, a wyraz wolny
$\det(-\varphi_*)=\pm1$. Włóknistość węzłów torusowych
wynika z opisu ich uzupełnień jako rozwarstwienia z monodromią;
pełny dowód można otrzymać przez analizę osobliwości
$f(z,w)=z^p+w^q$ i fibration Milnora na małej sferze
w $\mathbb C^2$.
\end{proof}

\begin{uwaga}
Jednomianowość wielomianu Alexandera jest warunkiem koniecznym,
ale nie wystarczającym włóknistości. Później zobaczymy, że
homologia Floera węzła pozwala rozpoznawać włóknistość,
a nie tylko wykrywać tę przeszkodę.
\end{uwaga}



\chapter{Grafy Taita, wielomian Tutte'a i macierze Goeritza}
\label{ch:tg-tait-goeritz}

W rozdziale~\ref{ch:tg-tutte} i w sekcji~\ref{sec:kt-jones}
poznaliśmy dwa opisy diagramu węzła: rozkład płaszczyzny
na obszary oraz wybory wygładzeń w nawiasie Kauffmana.
Zaskakująco wiele informacji o diagramie można przechować
w grafie planarnym: obszary stają się wierzchołkami,
a skrzyżowania krawędziami. Ten rozdział pokazuje
jednocześnie, kiedy opis grafowy jest pełny, a kiedy
trzeba zachować także znaki lub dane zanurzenia.

\section{Kolorowanie szachownicowe i grafy Taita}
\label{sec:tg-tait-def}

\begin{definicja}[Graf Taita]
Niech $D$ będzie spójnym regularnym diagramem splotu
na orientowanej sferze $S^2$, z co najmniej jednym
skrzyżowaniem. Kolorujemy składowe $S^2\setminus D$
na czarno i biało tak, aby obszary sąsiadujące przez
odcinek diagramu miały różne kolory. \emph{Czarny graf
Taita} $G_B(D)$ ma jeden wierzchołek w każdym czarnym
obszarze oraz jedną krawędź za każde skrzyżowanie,
łączącą wierzchołki dwóch czarnych sektorów przy tym
skrzyżowaniu. Krawędzie prowadzimy w małych otoczeniach
skrzyżowań, tak aby otrzymać graf zanurzony w $S^2$.
Analogicznie definiujemy biały graf Taita $G_W(D)$.
Pętle i wielokrotne krawędzie są dozwolone.
\end{definicja}

\begin{stwierdzenie}[Dualność grafów Taita]
Dla każdego spójnego regularnego diagramu splotu $D$
z co najmniej jednym skrzyżowaniem istnieje kolorowanie
szachownicowe jego obszarów. Po ustaleniu kolorowania
grafy $G_B(D)$ i $G_W(D)$ są planarnymi grafami dualnymi.
\end{stwierdzenie}
\begin{proof}
Na sferze diagram jest zanurzonym grafem czteroregularnym.
Jest eulerowski, ponieważ każdy wierzchołek ma stopień
cztery. Z twierdzenia o parzystości obwodów ścianowych
równoważnego dwudzielności grafu dualnego wynika
istnienie kolorowania szachownicowego. Można też
skonstruować je bezpośrednio: ustalamy kolor obszaru
zewnętrznego i zmieniamy kolor za każdym przejściem
przez łuk diagramu; zamknięta krzywa poprzeczna przecina
diagram parzyście wiele razy modulo dwa, ponieważ diagram
jest cyklem nad $\mathbb F_2$ na sferze i
$H_1(S^2;\mathbb F_2)=0$. Zatem wynik nie zależy od
wybranej drogi między obszarami.

W każdym skrzyżowaniu czarne sektory stanowią przeciwległą
parę, a białe drugą. Umieszczamy krawędź czarnego grafu
przez środek skrzyżowania, a krawędź białego grafu
poprzecznie do niej. Każde skrzyżowanie diagramu odpowiada
dokładnie jednemu przecięciu krawędzi obu grafów.
Białe obszary diagramu są ścianami czarnego grafu,
więc otrzymujemy definicję planarnego grafu dualnego.
\end{proof}

\begin{przyklad}[Dwa grafy Taita trójlistnika]
Dla standardowego trójskrzyżowaniowego diagramu trójlistnika
jeden wybór kolorowania daje graf trójkątny $C_3$.
Drugi daje graf o dwóch wierzchołkach połączonych trzema
równoległymi krawędziami. Są to wzajemnie dualne grafy
planarne; w obu przypadkach liczba krawędzi wynosi trzy.
Sama liczba krawędzi nie wystarcza do odtworzenia splotu:
w ogólności niezbędne są znaki skrzyżowań i sposób
zanurzenia grafu.
\end{przyklad}

\begin{center}
\begin{tikzpicture}[scale=.9,line cap=round,>=Stealth]
 \begin{scope}[shift={(-2.65,0)}]
  \node[tgBlueVertex] (a) at (0,1.1) {};
  \node[tgBlueVertex] (b) at (-1,-.65) {};
  \node[tgBlueVertex] (c) at (1,-.65) {};
  \draw[tgAccent] (a)--(b)--(c)--cycle;
  \node[font=\small,text=tgNavy] at (0,-1.25) {czarny graf Taita: $C_3$};
 \end{scope}
 \draw[tgTeal,-{Stealth[length=3mm]},line width=1.4pt] (-.85,.2)--(.85,.2);
 \begin{scope}[shift={(2.65,0)}]
  \node[tgRedVertex] (u) at (-.7,.05) {};
  \node[tgRedVertex] (v) at (.7,.05) {};
  \draw[tgTree] (u)--(v);
  \draw[tgTree] (u) to[bend left=55] (v);
  \draw[tgTree] (u) to[bend right=55] (v);
  \node[font=\small,text=tgNavy] at (0,-1.25) {biały graf Taita: trzy krawędzie};
 \end{scope}
\end{tikzpicture}
\end{center}

\section{Nawias Kauffmana ze zbiorów krawędzi}
\label{sec:tg-tait-tutte}

\begin{twierdzenie}[Wielomian Tutte'a a nawias diagramu alternującego]
\label{thm:tg-tait-tutte}
Niech $D$ będzie spójnym, niezorientowanym, alternującym
diagramem splotu, a $G$ jego grafem Taita dla kolorowania
szachownicowego wybranego tak, by w każdym skrzyżowaniu
wygładzenie A łączyło czarne sektory. Niech
$v=|V(G)|$ oraz $e=|E(G)|$. W przyjętej wcześniej
normalizacji nawiasu Kauffmana, w której
$\langle\bigcirc\rangle=1$ i
$\delta=-A^2-A^{-2}$, zachodzi
\[
\langle D\rangle
=A^{2v-e-2}\,T_G(-A^{-4},-A^4).
\]
Dla diagramu z orientacją wielomian Jonesa otrzymujemy
przez dodatkową normalizację
$V_L(t)=(-A^3)^{-w(D)}\langle D\rangle\big|_{A=t^{-1/4}}$.
\end{twierdzenie}

\begin{proof}
Niech $F\subseteq E(G)$ oznacza krawędzie tych skrzyżowań,
w których wybieramy wygładzenie A; w pozostałych
wybieramy B. Tak określony stan ma współczynnik
$A^{|F|}A^{-(e-|F|)}=A^{2|F|-e}$.
Czarne obszary diagramu można potraktować jako dyski,
a krawędzie $F$ jako wąskie pasma łączące te dyski.
Otrzymana płaska powierzchnia regularnego otoczenia
podgrafu rozpinającego $(V(G),F)$ ma $k(F)$ składowych,
charakterystykę Eulera $v-|F|$ i, ponieważ leży na
sferze, rodzaj zero na każdej składowej. Jeżeli
$b(F)$ oznacza liczbę jej okręgów brzegowych, to
$2k(F)-b(F)=v-|F|$, czyli
$b(F)=2k(F)-v+|F|$.
Te właśnie okręgi są okręgami stanu Kauffmana.
Zatem
\[
\langle D\rangle=
\sum_{F\subseteq E(G)}
 A^{2|F|-e}\delta^{2k(F)-v+|F|-1}.
\]
Dla spójnego grafu $G$ wielomian Tutte'a można
zapisać jako sumę po podzbiorach krawędzi:
\[
T_G(x,y)=\sum_{F\subseteq E(G)}
(x-1)^{k(F)-1}(y-1)^{|F|-v+k(F)}.
\]
Ponieważ przy $x=-A^{-4}$ i $y=-A^4$ zachodzi
$x-1=A^{-2}\delta$ i $y-1=A^2\delta$,
otrzymujemy po podstawieniu
\[
A^{2v-e-2}T_G(-A^{-4},-A^4)
=\sum_{F\subseteq E(G)}
 A^{2|F|-e}\delta^{2k(F)-v+|F|-1},
\]
co jest dokładnie poprzednią sumą stanów.
Dla diagramu alternującego można wybrać kolorowanie,
przy którym wszystkie wygładzenia A łączą czarne
sektory; przy zmianie konwencji trzeba odpowiednio
zamienić $A$ z $A^{-1}$ i graf z jego dualnym.
Formuła na $V_L$ jest definicją normalizacji
podaną w części poświęconej wielomianowi Jonesa.
\end{proof}

\begin{przyklad}[Pełny rachunek dla trójlistnika]
Dla grafu trójkątnego $G=C_3$ wyznaczyliśmy wcześniej
$T_G(x,y)=x^2+x+y$. Dlatego
\[
T_G(-A^{-4},-A^4)=A^{-8}-A^{-4}-A^4,
\quad
\langle D\rangle=A^{2\cdot3-3-2}T_G(-A^{-4},-A^4)
=A^{-7}-A^{-3}-A^5.
\]
Dla dodatnio zorientowanego diagramu trójlistnika
$w(D)=3$, zatem
$(-A^3)^{-3}\langle D\rangle
=A^{-4}+A^{-12}-A^{-16}$.
Po przyjęciu $t=A^{-4}$ dostajemy
$V_T(t)=t+t^3-t^4$, zgodnie z obliczeniem
z algebry Temperleya--Lieba z sekcji~\ref{sec:kt-jones}.
Dwa zupełnie różne
rachunki prowadzą więc do tego samego niezmiennika.
\end{przyklad}

\section{Macierz Goeritza i wyznacznik węzła}
\label{sec:tg-goeritz}

\begin{definicja}[Podpisany laplasjan grafu Taita]
Dla wybranego kolorowania szachownicowego diagramu $D$
przypisujemy każdemu skrzyżowaniu wagę
$\eta(c)\in\{-1,1\}$: znak dodatni oznacza,
że lokalne wygładzenie A łączy czarne sektory,
a ujemny, że łączy białe. Niech $G=G_B(D)$, a
$B$ będzie dowolną macierzą incydencji grafu $G$,
z pominięciem krawędzi pętlowych (ich kolumny są zerowe).
Definiujemy macierz
\[
\mathcal G_D=B\,\operatorname{diag}(\eta(c))\,B^T.
\]
Jej elementy pozadiagonalne są ujemnymi sumami wag
krawędzi łączących odpowiednie wierzchołki; przekątna
jest sumą wag krawędzi incydentnych. \emph{Zredukowaną
macierzą Goeritza} nazywamy macierz otrzymaną przez
usunięcie jednego wiersza i odpowiadającej kolumny
z $\mathcal G_D$.
\end{definicja}

\begin{stwierdzenie}[Suma drzew z wagami]
Niech $G$ będzie spójnym grafem o wierzchołkach
$1,\ldots,n$ i wagach $w_e\in\mathbb R$ na krawędziach.
Dla macierzy $L=B\operatorname{diag}(w_e)B^T$
i dowolnego jej minora głównego $L^{(r)}$
uzyskanego przez skreślenie wiersza i kolumny $r$
zachodzi
\[
\det L^{(r)}=\sum_{T\text{ drzewo rozpinające }G}
\ \prod_{e\in E(T)}w_e.
\]
Tożsamość obowiązuje także dla wag ujemnych.
\end{stwierdzenie}
\begin{proof}
Niech $B_r$ oznacza macierz powstałą z $B$ przez
usunięcie wiersza $r$. Z formuły Cauchy'ego--Bineta
mamy
\[
\det(B_r\operatorname{diag}(w_e)B_r^T)
=\sum_{F\subseteq E,\ |F|=n-1}
\det(B_{r,F})^2\prod_{e\in F}w_e.
\]
Minor $B_{r,F}$ jest zerowy, jeżeli $F$ nie jest
drzewem rozpinającym; gdy $F$ jest drzewem,
jego wyznacznik wynosi $\pm1$, czego dowiedliśmy
przy twierdzeniu Kirchhoffa przez rozwijanie
względem liści. Stąd dokładnie wskazany wzór.
Nie korzystaliśmy z dodatniości wag; ujemne wartości
są dozwolone, a ewentualne znaki pojawiają się jedynie
w iloczynach $\prod w_e$.
\end{proof}

\begin{twierdzenie}[Goeritz; identyfikacja wyznacznika]
Niech $K\subset S^3$ będzie węzłem, a $D$ jego
spójnym regularnym diagramem. Niech
$\mathcal G_D^{(r)}$ będzie dowolną zredukowaną
macierzą Goeritza zdefiniowaną powyżej. Wtedy
\[
\det(K)=|\Delta_K(-1)|=|\det\mathcal G_D^{(r)}|.
\]
W szczególności, dla alternującego diagramu węzła,
po wyborze kolorowania, przy którym wszystkie wagi
są jednakowe, otrzymujemy
$\det(K)=\tau(G_B(D))$, gdzie $\tau(G)$ oznacza
liczbę drzew rozpinających grafu $G$.
\end{twierdzenie}
\begin{uwaga}[Co zostało udowodnione]
Przed chwilą wykazaliśmy w pełni formułę wyznacznika
podpisanego laplasjanu jako sumy po drzewach.
Utożsamienie tego wyznacznika z $|\Delta_K(-1)|$
jest osobnym, klasycznym twierdzeniem Goeritza,
przy standardowej konwencji znaków skrzyżowań
oraz po odpowiedniej redukcji macierzy,
które można udowodnić przez prezentację homologii
pierwszej dwukrotnego nakrycia $S^3$ rozgałęzionego
wzdłuż $K$. Jego pełny dowód wymaga przygotowania
teorii nakryć rozgałęzionych; przyjmujemy go tutaj
z materiału o węzłach, nie utożsamiając go z
kombinatoryczną częścią rachunku.
\end{uwaga}

\begin{przyklad}[Wyznacznik trójlistnika bez macierzy Seiferta]
Dla $G=C_3$ wszystkie wagi można przyjąć za $+1$.
Laplasjan grafu ma macierz
\[
\mathcal G_D=\begin{pmatrix}
 2&-1&-1\\ -1&2&-1\\ -1&-1&2
\end{pmatrix}.
\]
Po skreśleniu trzeciego wiersza i kolumny dostajemy
$\mathcal G_D^{(3)}=\begin{psmallmatrix}2&-1\\-1&2\end{psmallmatrix}$,
więc $\det(K)=|4-1|=3$.
Ta sama liczba jest liczbą trzech drzew rozpinających
trójkątnego grafu Taita i zgadza się z wyznacznikiem
wielomianu Alexandera $t^2-t+1$ przy $t=-1$.
\end{przyklad}

\begin{uwaga}[Literatura]
Związek wielomianu Jonesa i wielomianu Tutte'a dla diagramów
alternujących jest klasycznym wynikiem Thistlethwaite'a.
Metodę wyznaczników macierzy Goeritza i jej związek z
nakryciami rozgałęzionymi opisują podręczniki
W. B. R. Lickorisha i D. Rolfsena.
\end{uwaga}


\chapter{Uzupełnienia węzłów, geometria i chirurgia Dehna}
\label{ch:kt-chirurgia-3d}

W tym rozdziale przechodzimy od krzywych zanurzonych w $S^3$ do
trójrozmaitości z brzegiem. Kluczowym obiektem będzie uzupełnienie
węzła, a kluczową operacją chirurgia Dehna, która pozwala budować
zamknięte trójrozmaitości przez ponowne doklejenie torusów pełnych.
Dzięki temu teoria węzłów staje się narzędziem konstrukcyjnym
w topologii trójwymiarowej.

\section{Otoczenie tubularne, meridian i długość}
\label{sec:kt-peryferia}

\begin{definicja}[Zewnętrze węzła]
Niech $K\subset S^3$ będzie węzłem. Ustalmy domknięte otoczenie
tubularne $N(K)\cong S^1\times D^2$, którego rdzeniem jest $K$.
\emph{Zewnętrzem węzła} nazywamy zwartą trójrozmaitość
$X_K:=S^3\setminus\operatorname{int}N(K)$. Jej brzeg
$\partial X_K=\partial N(K)\cong T^2$ jest torusem.
\end{definicja}

\begin{definicja}[Meridian i długość preferowana]
Niech $K\subset S^3$ będzie zorientowanym węzłem.
\emph{Meridianem} $\mu\subset\partial X_K$ nazywamy
zorientowaną krzywą prostą, która ogranicza dysk
w $N(K)$ i spełnia $\operatorname{lk}(\mu,K)=1$.
\emph{Długością preferowaną} (\emph{preferred longitude})
$\lambda\subset\partial X_K$ nazywamy zorientowaną,
prostą krzywą o liczbie splecenia zero z $K$, reprezentującą
w $H_1(N(K);\mathbb Z)$ tę samą klasę co zorientowany rdzeń.
Klasę $[\lambda]$ można równoważnie określić jako generator
jądra $H_1(\partial X_K;\mathbb Z)\to H_1(X_K;\mathbb Z)$,
z orientacją zgodną z $K$.
\end{definicja}

\begin{stwierdzenie}
Dla każdego zorientowanego węzła $K\subset S^3$ klasy
$[\mu],[\lambda]$ tworzą bazę $H_1(\partial X_K;\mathbb Z)
\cong\mathbb Z^2$, przy czym odwzorowanie
$i_*\colon H_1(\partial X_K;\mathbb Z)\to H_1(X_K;\mathbb Z)$
spełnia $i_*[\mu]=1$ i $i_*[\lambda]=0$, po identyfikacji
$H_1(X_K;\mathbb Z)\cong\mathbb Z$ przez liczbę splecenia.
\end{stwierdzenie}

\begin{proof}
Meridian przecina dysk normalny w torusie pełnym dokładnie
jeden raz i dlatego jest krzywą pierwotną na torusie brzegowym.
Wybrana długość ma z meridianem algebraiczną liczbę przecięć
równą jeden, więc ich klasy stanowią bazę $H_1(T^2;\mathbb Z)$.
Wnętrze zewnętrza ma pierwszą homologię generowaną przez
meridian, a odwzorowanie $H_1(X_K)\to\mathbb Z$ dane przez
splecenie z $K$ jest izomorfizmem. Długość preferowana
ma liczbę splecenia zero, stąd znika w $H_1(X_K)$.
\end{proof}

\begin{stwierdzenie}[Jednoznaczność długości preferowanej]
Niech $K\subset S^3$ będzie zorientowanym węzłem, $N(K)$ jego
otoczeniem tubularnym, a $X_K=S^3\setminus\operatorname{int}N(K)$.
Istnieje dokładnie jedna klasa $\lambda\in H_1(\partial X_K;\mathbb Z)$,
która przechodzi na generator $[K]\in H_1(N(K);\mathbb Z)$
i jednocześnie znika w $H_1(X_K;\mathbb Z)$.
W szczególności długość preferowana nie zależy od ramowania
wybranego przy początkowej identyfikacji otoczenia tubularnego.
\end{stwierdzenie}

\begin{proof}
Wybierzmy dowolną krzywą podłużną $\lambda_0$ przechodzącą
na $[K]$ w $H_1(N(K))$. Każda inna klasa o tej własności
ma postać $\lambda_0+k\mu$ dla pewnej liczby całkowitej $k$,
bo jądro homomorfizmu $H_1(\partial N(K))\to H_1(N(K))$
jest generowane przez meridian. Odwzorowanie
$H_1(X_K)\to\mathbb Z$, $[\gamma]\mapsto
\operatorname{lk}(\gamma,K)$, jest izomorfizmem.
Istnieje zatem dokładnie jedna wartość $k$, dla której
$\operatorname{lk}(\lambda_0+k\mu,K)=0$: mianowicie
$k=-\operatorname{lk}(\lambda_0,K)$.
Ponieważ $\operatorname{lk}(\mu,K)=1$, otrzymana
klasa jest jedyną długością preferowaną.
\end{proof}

\begin{center}
\begin{tikzpicture}[scale=.91,line cap=round,>=Stealth]
  \fill[ktBlue!14] (0,0) ellipse (3 and 1.32);
  \draw[ktBlue!70!black,line width=1.25pt] (0,0) ellipse (3 and 1.32);
  \fill[white] (0,0) ellipse (1.15 and .48);
  \draw[ktInk,line width=1.15pt] (0,0) ellipse (1.15 and .48);
  \draw[ktTeal,line width=3.1pt,densely dashed,opacity=.7] (-2.09,0)
        arc[start angle=180,end angle=360,x radius=2.09,y radius=.76];
  \draw[ktTeal,line width=3.2pt,-{Stealth[length=3mm]}] (2.09,0)
        arc[start angle=0,end angle=170,x radius=2.09,y radius=.76];
  \draw[ktCoral,line width=3.1pt,-{Stealth[length=2.5mm]}]
        (2.08,0) ellipse (.43 and .62);
  \node[font=\small\bfseries,text=ktTeal] at (-.15,1.75) {$\lambda$: kierunek podłużny};
  \draw[ktTeal,->,thick] (1.15,1.5)--(1.6,.97);
  \node[font=\small\bfseries,text=ktCoral] at (3.68,.68) {$\mu$};
  \draw[ktCoral,->,thick] (3.42,.55)--(2.45,.37);
  \node[text=ktInk,font=\footnotesize] at (0,-1.75)
        {Schemat torusa $\partial N(K)$: dwa kierunki homologii $H_1(T^2;\mathbb Z)$.};
\end{tikzpicture}
\end{center}

\begin{uwaga}
Na rysunku długość i meridian pokazano schematycznie.
Dla dokładnej definicji kierunków istotne są ich klasy w
$H_1(\partial X_K;\mathbb Z)$, a nie wybór konkretnego
położenia na powierzchni torusa.
\end{uwaga}

\section{Nieściśliwe torusy i topologia uzupełnienia}
\label{sec:kt-torus-jsj}

\begin{definicja}[Powierzchnia nieściśliwa]
Niech $M$ będzie zwartą, orientowalną trójrozmaitością,
a $F\subset M$ dwustronną, właściwie zanurzoną powierzchnią.
Mówimy, że $F$ jest \emph{nieściśliwa}, jeżeli żadna istotna
krzywa prosta na $F$ nie ogranicza dysku zanurzonego w $M$,
którego wnętrze jest rozłączne z $F$; równoważnie, przy zwykłych
założeniach o dwustronności i wyłączeniu składników sferycznych,
$\pi_1(F)\to\pi_1(M)$ jest injekcją.
\end{definicja}

\begin{stwierdzenie}
Jeśli $K\subset S^3$ jest nietrywialnym węzłem, to
$\partial X_K$ jest nieściśliwy w $X_K$.
Dla węzła trywialnego jego zewnętrze jest torusem pełnym,
a torus brzegowy jest ściśliwy.
\end{stwierdzenie}

\begin{proof}
Przypuśćmy, że torus $\partial X_K$ jest ściśliwy i niech
$D\subset X_K$ będzie dyskiem kompresującym, to znaczy
$\partial D$ jest istotną krzywą na $\partial X_K$.
Jej klasa ma postać $p\mu+q\lambda$ z
$\gcd(p,q)=1$. Ponieważ $[\partial D]=0$ w $H_1(X_K)$,
a $[\mu]$ generuje tę grupę, koniecznie $p=0$.
Prymitywność klasy wymusza $q=\pm1$.
Tak więc $\partial D$ jest izotopijny na torusie do
$\lambda$ lub $-\lambda$. W otoczeniu $N(K)$ istnieje
pierścień o jednej składowej brzegu równiej $\lambda$
i drugiej równej rdzeniowi $K$: wynika to z faktu, że
$\lambda$ jest krzywą podłużną i reprezentuje generator
grupy $H_1(N(K))$ z odpowiednią orientacją.
Sklejając ten pierścień z $D$ i wygładzając narożniki,
otrzymujemy gładko zanurzony dysk o brzegu $K$.
Oznacza to, że $K$ jest trywialny.

Odwrotnie, dla węzła trywialnego jego zewnętrze jest
torusem pełnym. Dysk południkowy tego torusa stanowi
dysk kompresujący $\partial X_U$.
\end{proof}

\begin{definicja}[Węzeł satelitarny]
Węzeł $K\subset S^3$ nazywamy \emph{satelitarnym}, jeżeli jego
zewnętrze $X_K$ zawiera torus $T$ nieściśliwy w $X_K$,
który nie jest izotopijny do $\partial X_K$.
Węzeł satelitarny można geometrycznie rozumieć jako
węzeł zanurzony w nietrywialnym torusie pełnym, którego rdzeń
jest innym węzłem.
\end{definicja}

\begin{uwaga}
Istnienie torusa nieściśliwego w uzupełnieniu nie musi
wynikać z tego, że sam węzeł jest torusowy. Dla węzłów
torusowych $T(p,q)$, $p,q\geq2$, uzupełnienie jest
przestrzenią Seiferta i nie zawiera istotnego torusa
satelitarnego w powyższym sensie. Torusy otaczające
węzły towarzyszące odgrywają rolę dla węzłów satelitarnych.
\end{uwaga}

\section{Twierdzenie Gordona--Lueckego}
\label{sec:kt-gl}

\begin{twierdzenie}[Gordon--Luecke; węzły są wyznaczone przez uzupełnienia]
Niech $K_1,K_2\subset S^3$ będą węzłami. Jeżeli istnieje
homeomorfizm zachowujący orientację
$h\colon X_{K_1}\to X_{K_2}$, to istnieje dyfeomorfizm
$F\colon S^3\to S^3$ zachowujący orientację i przenoszący
$K_1$ na $K_2$. W konsekwencji węzły są równoważne
jako węzły niezorientowane. Stwierdzenie nie rozstrzyga
samo przez się równoważności ustalonych orientacji $K_i$.
\end{twierdzenie}

\begin{uwaga}
To głęboki wynik, którego dowód korzysta z teorii chirurgii
na węzłach i własności powierzchni nieściśliwych.
Nie wystarcza sam fakt, że $X_K$ ma brzeg torusowy:
przy dowolnej trójrozmaitości z takim brzegiem różne
wypełnienia mogą prowadzić do tej samej przestrzeni.
Specjalna rola $S^3$ jest tu zasadnicza. Twierdzenie
pokazuje, że problem klasyfikacji węzłów można przenieść
do topologii ich zewnętrzy.
\end{uwaga}

\section{Geometria hiperboliczna uzupełnień}
\label{sec:kt-hyperbolic}

\begin{definicja}[Węzeł hiperboliczny]
Węzeł $K\subset S^3$ jest \emph{hiperboliczny}, jeżeli
$S^3\setminus K$ dopuszcza zupełną metrykę riemannowską
o stałej krzywiźnie sekcyjnej $-1$ i skończonej objętości.
Objętość tę oznaczamy $\operatorname{Vol}(K)$.
\end{definicja}

\begin{twierdzenie}[Hiperbolizacja uzupełnień węzłów; Thurston]
Każdy nietrywialny węzeł $K\subset S^3$ należy do jednej
z trzech rozłącznych klas: torusowych, satelitarnych albo
hiperbolicznych. W ostatnim przypadku uzupełnienie posiada
zupełną, skończenie objętościową strukturę hiperboliczną.
\end{twierdzenie}

\begin{uwaga}
Twierdzenie jest konsekwencją teorii hiperbolizacji
Hakenowskich trójrozmaitości i rozkładu JSJ, a w szerszym
kontekście wpisuje się w geometryzację Thurstona.
Dla węzła trywialnego zewnętrze jest torusem pełnym,
a $S^3\setminus U$ nie posiada skończenie objętościowej
zupełnej metryki hiperbolicznej.
\end{uwaga}

\begin{twierdzenie}[Sztywność Mostowa--Prasada]
Niech $M_1,M_2$ będą spójnymi, zupełnymi, orientowalnymi
trójrozmaitościami hiperbolicznymi o skończonej objętości.
Jeżeli $\pi_1(M_1)\cong\pi_1(M_2)$, to istnieje izometria
$M_1\to M_2$ realizująca ten izomorfizm z dokładnością
do sprzężenia wewnętrznego. W szczególności objętość
hiperboliczna jest niezmiennikiem topologicznym.
\end{twierdzenie}

\begin{przyklad}[Węzeł ósemkowy]
Węzeł ósemkowy $4_1$ jest hiperboliczny. Jego uzupełnienie
można przedstawić jako sklejanie dwóch regularnych idealnych
czworościanów hiperbolicznych. Stąd
$\operatorname{Vol}(4_1)=2\,v_3\approx2.0298832128$,
gdzie $v_3$ jest objętością regularnego idealnego
czworościanu w $\mathbb H^3$. Trójlistnik nie jest
hiperboliczny, ponieważ jest węzłem torusowym.
\end{przyklad}

\begin{uwaga}[Koniec paraboliczny]
Zewnętrze hiperbolicznego węzła jest zwarte i ma brzeg
będący torusem, natomiast jego wnętrze $S^3\setminus K$
jest niezwartą rozmaitością hiperboliczną. Geometryczny
koniec ma postać $T^2\times[0,\infty)$ z metryką
$\mathrm ds^2=\mathrm dt^2+e^{-2t}g_0$,
gdzie $g_0$ jest metryką płaską na torusie.
Całkowita objętość końca jest skończona, ponieważ
$\int_0^\infty e^{-2t}\,\mathrm dt<\infty$.
\end{uwaga}

\section{Chirurgia Dehna: definicja i wybór nachylenia}
\label{sec:kt-dehn}

\begin{definicja}[Nachylenie na torusie]
Niech $K\subset S^3$ będzie zorientowanym węzłem, a
$([\mu],[\lambda])$ bazą $H_1(\partial X_K;\mathbb Z)$.
\emph{Nachyleniem} na $\partial X_K$ nazywamy niezorientowaną
klasę izotopii istotnej krzywej prostej, reprezentowanej przez
pierwotny element $p[\mu]+q[\lambda]$, gdzie
$p,q\in\mathbb Z$, $\gcd(p,q)=1$, z utożsamieniem
$(p,q)\sim(-p,-q)$. Oznaczamy je symbolem $p/q$;
przez $1/0=\infty$ rozumiemy nachylenie meridianu.
\end{definicja}

\begin{definicja}[Wypełnienie i chirurgia Dehna]
Dla węzła $K\subset S^3$ i nachylenia $p/q$ wybieramy
homeomorfizm $\varphi\colon\partial(D^2\times S^1)
\to\partial X_K$ wysyłający brzeg dysku południkowego
$\partial D^2\times\{*\}$ na krzywą o klasie
$p[\mu]+q[\lambda]$. Rozmaitość
\[
S^3_{p/q}(K):=X_K\cup_\varphi(D^2\times S^1)
\]
nazywamy \emph{chirurgią Dehna o współczynniku $p/q$}.
Wypełnienie dla $p/q=\infty$ odtwarza $S^3$.
\end{definicja}

\begin{stwierdzenie}
Rozmaitość $S^3_{p/q}(K)$ jest dobrze określona z dokładnością
do homeomorfizmu zachowującego orientację przez nieorientowaną
klasę nachylenia $p/q$; inne wybory homeomorfizmu
$\varphi$ o zadanym obrazie meridianu torusa dołączanego
prowadzą do homeomorficznych wypełnień.
\end{stwierdzenie}

\begin{proof}
Niech $\varphi_0,\varphi_1:\partial V\to\partial X_K$
będą dwoma sklejeniami, gdzie $V=D^2\times S^1$.
Zakładamy, że obrazy klasy meridianu $m=[\partial D^2\times\{*\}]$
wyznaczają to samo niezorientowane nachylenie. Wtedy
automorfizm $h:=\varphi_0^{-1}\circ\varphi_1$ torusa
$\partial V$ zachowuje nieorientowaną klasę $m$.
Przypomnijmy elementarny opis takich homeomorfizmów.
W bazie $(m,\ell)$ macierz $h_*$ ma postać
$\begin{psmallmatrix}\varepsilon & k\\ 0 & \delta\end{psmallmatrix}$,
gdzie $\varepsilon,\delta\in\{\pm1\}$ i $k\in\mathbb Z$;
zachowanie orientacji właściwego sklejenia wymaga
$\varepsilon\delta=1$. Macierz ta może być zrealizowana
przez skręt Dehna wzdłuż meridianu, ewentualnie z obrotem
obu kierunków. Każdy z tych homeomorfizmów przedłuża się
na torus pełny: skręt wzdłuż meridianu przedłużamy
przez obrót dysków $D^2\times\{z\}$ zależny od $z\in S^1$.
Zatem $h$ (po izotopii) przedłuża się do homeomorfizmu
$H:V\to V$. Odwzorowanie będące identycznością na $X_K$
i równe $H$ na dołączanym $V$ wyznacza homeomorfizm
wypełnień z odpowiednio dobraną orientacją.
Dowód pokazuje dokładnie, dlaczego wybór całego sklejenia
nie wnosi nowej informacji poza nachyleniem meridianu.
\end{proof}


\medskip
\begin{center}
\begin{tikzpicture}[scale=.92,>=Stealth,line cap=round,line join=round]
  \fill[ktBlue!6] (0,0) rectangle (4.2,2.5);
  \draw[ktBlue,line width=1.5pt] (0,0) rectangle (4.2,2.5);
  \draw[ktInk!25,densely dashed] (1.4,0)--(1.4,2.5) (2.8,0)--(2.8,2.5)
    (0,1.25)--(4.2,1.25);
  \draw[ktCoral,-{Stealth[length=3mm]},line width=2.8pt] (.45,.38)--(3.25,2.02);
  \node[text=ktCoral,anchor=west,font=\small] at (3.34,1.95) {$\mu+\lambda$};
  \draw[ktTeal,-{Stealth[length=2.8mm]},line width=1.6pt] (.35,-.42)--(1.55,-.42);
  \draw[ktPurple,-{Stealth[length=2.8mm]},line width=1.6pt] (-.35,.25)--(-.35,1.55);
  \node[text=ktTeal,font=\small] at (.95,-.67) {$\mu$};
  \node[text=ktPurple,font=\small] at (-.62,.92) {$\lambda$};
\end{tikzpicture}

{\footnotesize Model torusa $T^2=\mathbb R^2/\mathbb Z^2$ z zaznaczonym nachyleniem $1/1$. \\
W chirurgii Dehna doklejamy torus pełny tak, aby meridian torusa dołączanego \\
przeszedł na wybraną klasę nachylenia.}
\end{center}
\medskip

\begin{twierdzenie}[Homologia pierwsza po chirurgii]
Niech $K\subset S^3$ będzie węzłem, a $p,q\in\mathbb Z$
liczbami względnie pierwszymi. Wtedy
$H_1(S^3_{p/q}(K);\mathbb Z)\cong\mathbb Z/|p|\mathbb Z$
dla $p\ne0$, a dla $p=0$ jest izomorficzna z $\mathbb Z$.
Przyjmujemy $\mathbb Z/1\mathbb Z=0$.
\end{twierdzenie}

\begin{proof}
Niech $V=D^2\times S^1$ będzie dołączanym torusem pełnym,
a $\nu\in H_1(V;\mathbb Z)\cong\mathbb Z$ jego klasą rdzenia.
W ciągu Mayera--Vietorisa rozkładu
$M=X_K\cup_\varphi V$ rozważamy fragment
\[
H_1(\partial X_K)\longrightarrow H_1(X_K)\oplus H_1(V)
\longrightarrow H_1(M)\longrightarrow
H_0(\partial X_K).
\]
Ostatnia mapa do $H_0(X_K)\oplus H_0(V)$ jest injekcją,
gdyż wszystkie przestrzenie są spójne. Wobec tego $H_1(M)$
jest kokernem pierwszej strzałki.

W bazie $(\mu,\lambda)$ obraz elementu $p\mu+q\lambda$
w $H_1(V)$ jest zerowy, bo jest meridianem doklejonego
torusa. Jego obraz w $H_1(X_K)$ wynosi $p[\mu]$,
ponieważ $[\lambda]=0$. Otrzymujemy więc relację
$p[\mu]=0$. Druga, niezależna klasa na torusie brzegowym
przechodzi w generator $\nu$ z dokładnością do znaku,
dzięki czemu eliminuje się generator $H_1(V)$ i nie
powstaje nowa relacja na $[\mu]$. W rezultacie
$H_1(M)\cong\mathbb Z/(p)$, przy czym dla $p=0$
ilorazem jest $\mathbb Z$, a dla $p\ne0$ grupa cykliczna
rzędu $|p|$.
\end{proof}

\begin{przyklad}[Chirurgia na węźle trywialnym]
Dla węzła trywialnego $U$ jego zewnętrze $X_U$
jest torusem pełnym. Wypełnienie $p/q$ otrzymujemy
przez sklejenie dwóch torusów pełnych, a więc jest
przestrzenią soczewkową. Przy konwencji definiującej
$L(p,q)$ właśnie przez tę operację, mamy
$S^3_{p/q}(U)\cong L(p,q)$ dla $p>0$, z odpowiednimi
konwencjami orientacji przy $p<0$. Dla $p=1$ otrzymujemy
$S^3$, a dla nachylenia $0=0/1$ dostajemy $S^1\times S^2$.
W szczególności sama grupa $H_1\cong\mathbb Z/p\mathbb Z$
nie wystarcza do rozróżniania wszystkich przestrzeni
soczewkowych: parametr $q$ wpływa na ich topologię.
\end{przyklad}

\section{Chirurgia na splotach i macierz splecenia}
\label{sec:kt-chirurgia-link}

\begin{definicja}[Splot z ramowaniem]
Niech $L=K_1\cup\dots\cup K_r\subset S^3$ będzie
zorientowanym splotem. \emph{Ramowaniem} składowej $K_i$
nazywamy klasę homotopii nigdzie niezerowego pola normalnego
wzdłuż $K_i$, równoważnie klasę równoległej krzywej
$K_i'\subset\partial N(K_i)$. Względem długości preferowanej
$\lambda_i$ ramowanie jest całkowitą liczbą $n_i$, taką że
$[K_i']=n_i[\mu_i]+[\lambda_i]$.
Splot z ramowaniem oznaczamy $(L;n_1,\dots,n_r)$.
\end{definicja}

Ramowanie jest więc szczególnym, całkowitym nachyleniem
chirurgii. Wprowadza strukturę konieczną dla uchwytów
wymiaru cztery: bez ramowania nie można jednoznacznie
zdefiniować sposobu ich dołączenia.

\begin{definicja}[Macierz splecenia]
Dla splotu $L=K_1\cup\dots\cup K_r$ z całkowitymi
ramowaniami $n_1,\dots,n_r$ definiujemy symetryczną
macierz całkowitą $Q_L=(q_{ij})_{i,j=1}^r$ przez
$q_{ii}=n_i$ oraz $q_{ij}=\operatorname{lk}(K_i,K_j)$
dla $i\ne j$.
\end{definicja}

\begin{twierdzenie}[Homologia po integralnej chirurgii]
Niech $M_L$ będzie zamkniętą, orientowalną trójrozmaitością
otrzymaną przez integralną chirurgię na splocie
$(L;n_1,\ldots,n_r)\subset S^3$. Wtedy
\[
H_1(M_L;\mathbb Z)\cong\operatorname{coker}
(Q_L:\mathbb Z^r\longrightarrow\mathbb Z^r).
\]
W szczególności, jeśli $\det Q_L\ne0$, to
$|H_1(M_L;\mathbb Z)|=|\det Q_L|$.
\end{twierdzenie}

\begin{proof}
Zewnętrze $X_L$ ma pierwszą homologię swobodną z bazą
meridianów $[\mu_1],\ldots,[\mu_r]$. Każda długość
preferowana $\lambda_i$, po włączeniu do $X_L$,
ma klasę $\sum_{j\ne i}\operatorname{lk}(K_i,K_j)[\mu_j]$:
wynika to z obliczenia liczb splecenia z pozostałymi
składowymi. Przy integralnej chirurgii doklejony meridian
spełnia relację $n_i[\mu_i]+[\lambda_i]=0$.
Zapisanie wszystkich $r$ relacji daje właśnie macierz $Q_L$.
Jeżeli $Q_L$ jest nieosobliwa nad $\mathbb Q$, jej
kokernelem jest skończona grupa abelowa rzędu
$|\det Q_L|$, co wynika ze sprowadzenia do postaci
normalnej Smitha.
\end{proof}

\begin{uwaga}[Całkowite ramowanie i wzór na diagonalę]
Współczynnik całkowitej chirurgii na składowej $K_i$
oznacza krzywą $n_i\mu_i+\lambda_i$, gdzie $\lambda_i$
jest długością preferowaną w uzupełnieniu pojedynczego węzła
$K_i$. W zewnętrzu całego splotu jej klasa nie musi już
znikać: wynosi $\sum_{j\ne i}\operatorname{lk}(K_i,K_j)\mu_j$.
Dlatego przekątna macierzy $Q_L$ jest zadana ramowaniami,
a elementy pozadiagonalne liczbami splecenia. Nie należy
utożsamiać wartości $n_i$ z liczbą skrzyżowań w diagramie.
\end{uwaga}

\begin{przyklad}[Integralna chirurgia na splocie Hopfa]
Niech $L=K_1\cup K_2$ będzie dodatnim splotem Hopfa,
a $n_1,n_2\in\mathbb Z$ ustalonymi ramowaniami.
Wtedy $Q_L=\begin{psmallmatrix}n_1&1\\1&n_2\end{psmallmatrix}$,
$\det Q_L=n_1n_2-1$. Jeżeli $n_1=n_2=0$, to
$Q_L=\begin{psmallmatrix}0&1\\1&0\end{psmallmatrix}$,
więc $H_1(M_L;\mathbb Z)=0$. Nie oznacza to automatycznie,
że $M_L\cong S^3$ dla dowolnego splotu z taką macierzą:
homologia nie rozstrzyga typu topologicznego. Dla standardowego
splotu Hopfa z oboma ramowaniami zero rachunek Kirby'ego
rzeczywiście daje $S^3$.
\end{przyklad}

\section{Nakrycia rozgałęzione nad węzłami}
\label{sec:kt-branched}

Oprócz usuwania węzła i wykonywania chirurgii można
zbudować nową rozmaitość przez wzięcie nakrycia
rozgałęzionego. Ta operacja łączy liczbę skrzyżowań,
wyznacznik węzła i teorię form przecięcia.

\begin{definicja}[Dwukrotne nakrycie rozgałęzione]
Niech $K\subset S^3$ będzie węzłem. \emph{Dwukrotnym
nakryciem rozgałęzionym $S^3$ wzdłuż $K$} nazywamy
trójrozmaitość $\Sigma_2(K)$ z odwzorowaniem
$p\colon\Sigma_2(K)\to S^3$, które poza przeciwobrazem
$K$ jest dwukrotnym nakryciem, a w lokalnych zespolonych
współrzędnych normalnych do $K$ ma postać
$(\theta,z)\mapsto(\theta,z^2)$,
gdzie $\theta$ jest współrzędną wzdłuż węzła.
\end{definicja}

\begin{twierdzenie}[Wyznacznik z nakrycia]
Dla dowolnego węzła $K\subset S^3$ grupa
$H_1(\Sigma_2(K);\mathbb Z)$ jest skończona i
jej rząd wynosi
$|H_1(\Sigma_2(K);\mathbb Z)|
=\det(K)=|\Delta_K(-1)|$.
Jeżeli $V$ jest macierzą Seiferta węzła $K$,
to grupa ta jest izomorficzna z
$\operatorname{coker}(V+V^T:\mathbb Z^{2g}\to\mathbb Z^{2g})$.
\end{twierdzenie}

\begin{proof}[Macierz przecięć i ciąg dokładny pary]
Wybierzmy powierzchnię Seiferta $F\subset S^3$ rodzaju $g$
i przesuńmy jej wnętrze w głąb $B^4$, otrzymując właściwie
zanurzoną powierzchnię $F'\subset B^4$ o brzegu $K$.
Tworzymy dwukrotne nakrycie rozgałęzione
$W_F\to B^4$ wzdłuż $F'$. W lokalnym modelu rozgałęzienia
$z\mapsto z^2$ na normalnym dysku sprawdzamy, że brzeg
$W_F$ jest kanonicznie równy $\Sigma_2(K)$.

Powierzchnię $F$ można przedstawić jako dysk z dołączonymi
$2g$ pasmami. Opis jej nakrycia rozgałęzionego prowadzi do
modelu uchwytowego $W_F$ z jednym uchwytem indeksu zero
i $2g$ uchwytami indeksu dwa. Wynika stąd
$H_1(W_F;\mathbb Z)=0$ i $H_2(W_F;\mathbb Z)\cong\mathbb Z^{2g}$.
Przecięcia rdzeni odpowiednich uchwytów, po zamknięciu ich
powierzchniami, oblicza się przez wzajemne przesunięcia
krzywych z bazy $H_1(F)$. Diagonalne wkłady pochodzą z
ramowań pasm, a pozadiagonalne z obu kierunków przesunięcia.
Ostatecznie forma przecięcia jest reprezentowana przez
symetryzację macierzy Seiferta $S=V+V^T$.
To właśnie jest podstawowy lemat rachunku nakryć
rozgałęzionych; jego konstrukcja uchwytowa jest kluczowym
zewnętrznym elementem dowodu.

Ponieważ $\partial W_F=\Sigma_2(K)$, ciąg dokładny pary
i dualność Poincarégo--Lefschetza dają
\[
H_2(W_F)\xrightarrow{\;S\;}H_2(W_F,\partial W_F)
\longrightarrow H_1(\partial W_F)\longrightarrow H_1(W_F)=0.
\]
Obie grupy środkowe są wolne rzędu $2g$, a pierwsza mapa
ma w wybranych bazach macierz $S$. Zatem
$H_1(\Sigma_2(K);\mathbb Z)\cong\operatorname{coker}S$.
Ponadto $\det S=\det(V+V^T)=\Delta_K(-1)$ dla macierzy
rozmiaru parzystego, przy odpowiedniej normalizacji.
Jest to liczba nieparzysta, a więc niezerowa; stąd
$|H_1(\Sigma_2(K);\mathbb Z)|=|\Delta_K(-1)|$.
\end{proof}

\begin{przyklad}
Dla trójlistnika
$V+V^T=\begin{psmallmatrix}2&-1\\-1&2\end{psmallmatrix}$.
Postać normalna Smitha ma przekątną $(1,3)$,
zatem $H_1(\Sigma_2(T);\mathbb Z)\cong\mathbb Z/3\mathbb Z$.
Dla węzła ósemkowego
$V+V^T=\begin{psmallmatrix}2&-1\\-1&-2\end{psmallmatrix}$,
więc $H_1(\Sigma_2(4_1);\mathbb Z)\cong\mathbb Z/5\mathbb Z$.
Same rzędy tych grup odróżniają oba węzły.
\end{przyklad}

\begin{twierdzenie}[Homologia cyklicznych nakryć rozgałęzionych]
Niech $K\subset S^3$ będzie węzłem, a $d\ge2$ liczbą
całkowitą. Przez $\Sigma_d(K)$ oznaczmy $d$-krotne
cykliczne nakrycie $S^3$ rozgałęzione wzdłuż $K$.
Jeżeli $\zeta_d=e^{2\pi i/d}$ i
$\Delta_K(\zeta_d^j)\ne0$ dla $j=1,\dots,d-1$, to
$H_1(\Sigma_d(K);\mathbb Z)$ jest skończona i
\[
|H_1(\Sigma_d(K);\mathbb Z)|=
\left|\prod_{j=1}^{d-1}\Delta_K(\zeta_d^j)\right|.
\]
Jeżeli przynajmniej jeden z tych czynników jest zerem,
to $H_1(\Sigma_d(K);\mathbb Z)$ ma dodatnią rangę.
\end{twierdzenie}

\begin{proof}[Redukcja do modułu Alexandera]
Niech $A_K=H_1(\widetilde X_K;\mathbb Z)$ będzie modułem
Alexandera nad $\Lambda=\mathbb Z[t,t^{-1}]$.
Przejście do $d$-krotnego nakrycia nierozgałęzionego
odpowiada narzuceniu relacji $t^d=1$. Następnie
doklejamy torusy pełne tak, aby ich meridiany zabijały
odpowiednie podniesienia meridianu węzła. Z ciągu
Mayera--Vietorisa tej operacji uzyskujemy
$H_1(\Sigma_d(K);\mathbb Z)\cong A_K/(t^d-1)A_K$.

Z prezentacji modułu przez macierz $tV-V^T$
wynika, że ostatni iloraz jest kokernem całkowitej
macierzy otrzymanej przez podstawienie w $tV-V^T$
macierzy cyklicznego przesunięcia $d\times d$.
Po zespoleniu przesunięcie ma wartości własne
$1,\zeta_d,\dots,\zeta_d^{d-1}$. Wyznacznik
macierzy prezentacji jest więc równy
$\prod_{j=0}^{d-1}\det(\zeta_d^jV-V^T)$.
Czynnik dla $j=0$ to $\Delta_K(1)=1$.
Jeżeli iloczyn jest niezerowy, koker ma rząd równy
jego wartości bezwzględnej; jeśli jest zerowy, koker
posiada niezerową część wolną. Otrzymujemy tezę.
\end{proof}

\begin{przyklad}
Dla trójlistnika $\Delta_T(t)=t^2-t+1$ i
$\Delta_T(-1)=3$, co daje
$|H_1(\Sigma_2(T);\mathbb Z)|=3$.
Dla $d=3$ pierwiastek $\zeta_3$ spełnia
$\zeta_3^2+\zeta_3+1=0$, więc
$\Delta_T(\zeta_3)=-2\zeta_3$ i
$\Delta_T(\zeta_3^2)=-2\zeta_3^2$.
Stąd $|H_1(\Sigma_3(T);\mathbb Z)|=4$.
\end{przyklad}

\section{Forma splecenia na wymiernej sferze homologicznej}
\label{sec:kt-linking-pairing}

\begin{definicja}[Wymierna sfera homologiczna]
Zamkniętą, spójną, orientowalną trójrozmaitość $M$
nazywamy \emph{wymierną sferą homologiczną}, jeżeli
$H_i(M;\mathbb Q)\cong H_i(S^3;\mathbb Q)$
dla wszystkich $i\geq0$. Równoważnie $H_1(M;\mathbb Z)$
jest grupą skończoną.
\end{definicja}

\begin{definicja}[Forma splecenia torsyjnego]
Niech $M$ będzie orientowaną wymierną sferą homologiczną.
Dla $x,y\in H_1(M;\mathbb Z)$ wybieramy reprezentujące
je rozłączne, zorientowane 1-cykle $a,b\subset M$.
Ponieważ $x$ ma skończony rząd, istnieje $m\geq1$
i całkowity 2-łańcuch $C$ taki, że $\partial C=ma$.
Definiujemy
$\ell_M(x,y):=\frac{C\cdot b}{m}\pmod{\mathbb Z}
\in\mathbb Q/\mathbb Z$.
Tak otrzymana forma jest dobrze określona, symetryczna,
dwuliniowa i niezdegenerowana.
\end{definicja}

\begin{stwierdzenie}[Dobra określoność formy splecenia]
Dla każdej zamkniętej, spójnej, orientowanej wymiernej
sfery homologicznej $M$ liczba $\ell_M(x,y)$ z poprzedniej
definicji nie zależy od wyboru reprezentantów, liczby $m$
ani łańcucha $C$; określa niezdegenerowaną, symetryczną
formę dwuliniową na $H_1(M;\mathbb Z)$.
\end{stwierdzenie}

\begin{proof}
Niech $C,C'$ będą dwoma 2-łańcuchami z brzegiem
$m a$ (po sprowadzeniu do wspólnej wielokrotności,
co nie zmienia ilorazu). Ich różnica jest 2-cyklem.
Z dualności Poincarégo i faktu, że $H_1(M;\mathbb Q)=0$,
wynika $H_2(M;\mathbb Z)=0$, a więc różnica jest
brzegiem 3-łańcucha. Jego przecięcie z zamkniętym
1-cyklem wynosi zero. Zmiana reprezentanta $a$ lub $b$
o brzeg 2-łańcucha zmienia wartość ilorazu o liczbę
całkowitą. Stąd otrzymujemy wartość w $\mathbb Q/\mathbb Z$.

Dwuliniowość wynika z addytywności przecięć łańcuchów.
Symetrię uzyskujemy, biorąc $\partial C=ma$,
$\partial D=nb$ i porównując granicę przecięcia
$C\cap D$; daje to $n(C\cdot b)=m(a\cdot D)$,
a zatem $\ell_M(x,y)=\ell_M(y,x)$ w $\mathbb Q/\mathbb Z$.
Niezdegenerowanie jest konsekwencją dualności
Poincarégo z wartościami w $\mathbb Q/\mathbb Z$:
mapa $H_1(M)\to\operatorname{Hom}(H_1(M),\mathbb Q/\mathbb Z)$
indukowana przez $\ell_M$ jest izomorfizmem.
\end{proof}

\begin{twierdzenie}[Forma splecenia z macierzy chirurgii]
Niech $L\subset S^3$ będzie $r$-składowym splotem z
całkowitymi ramowaniami, a jego macierz splecenia
$Q\in M_r(\mathbb Z)$ niech spełnia $\det Q\ne0$.
Niech $M_L$ będzie wynikiem integralnej chirurgii na $L$.
Wówczas, przy identyfikacji
$H_1(M_L;\mathbb Z)\cong\mathbb Z^r/Q\mathbb Z^r$,
forma splecenia torsyjnego jest reprezentowana przez
$\ell_{M_L}([x],[y])\equiv -x^TQ^{-1}y\pmod{\mathbb Z}$,
przy przyjętej konwencji orientacji brzegu czterorozmaitości
uzyskanej przez dołączenie uchwytów.
\end{twierdzenie}

\begin{proof}
Dobrą określoność prawej strony sprawdzamy bezpośrednio:
po zastąpieniu $x$ przez $x+Qz$ wartość zmienia się
o $-z^Ty\in\mathbb Z$, a podobnie dla zmiany $y$.
Symetria wynika z $Q^T=Q$, zaś niezdegenerowanie
z nieosobliwości macierzy nad $\mathbb Q$.
Identyfikację z geometrycznie zdefiniowaną formą splecenia
otrzymujemy przez rozważenie powierzchni wypełniających
w czterorozmaitości $W_L$, wykorzystanie
$H_2(W_L,\partial W_L;\mathbb Z)$ i mapy
$H_2(W_L)\to H_2(W_L,\partial W_L)$ reprezentowanej
przez macierz $Q$. Znak minus wynika z orientacji
brzegu w odpowiedniej konwencji przecięć.
\end{proof}

\begin{przyklad}[Chirurgia na pojedynczej składowej]
Dla ramowania $p\ne0$ macierz chirurgii jest $Q=[p]$.
Stąd $H_1(M;\mathbb Z)\cong\mathbb Z/|p|\mathbb Z$,
a na generatorze $x$ zachodzi
$\ell_M(x,x)\equiv-1/p\pmod{\mathbb Z}$.
W szczególności dla przestrzeni soczewkowych sama
postać grupy homologii nie wystarcza do rozpoznania
orientowanej rozmaitości; forma splecenia przechowuje
więcej danych o topologii.
\end{przyklad}


\section{Twierdzenie Lickorisha--Wallace'a i rachunek Kirby'ego}
\label{sec:kt-wallace}

\begin{twierdzenie}[Lickorish--Wallace]
Niech $M$ będzie dowolną zamkniętą, spójną, orientowalną
trójrozmaitością. Istnieją liczba $r\geq0$, splot
$L\subset S^3$ o $r$ składowych i całkowite ramowania
$n_1,\dots,n_r$ takie, że $M$ jest homeomorficzna z zachowaniem orientacji
z trójrozmaitością otrzymaną przez integralną chirurgię
na $(L;n_1,\dots,n_r)$.
\end{twierdzenie}

\begin{proof}[Idea dowodu]
Z rozkładu Heegaarda $M=H_g\cup_\varphi H_g$ otrzymujemy
$M$ jako sklejenie dwóch brył uchwytowych wzdłuż
powierzchni rodzaju $g$. Grupę klas odwzorowań tej
powierzchni generują skręty Dehna wzdłuż odpowiednich
krzywych prostych. Każdy taki skręt można zrealizować
przez integralną chirurgię na krzywej zanurzonej
w powierzchni Heegaarda, z ramowaniem względem powierzchni.
Składając chirurgiczne realizacje generatorów i przenosząc
krzywe do $S^3$, dostajemy splot z ramowaniem, którego
chirurgia odtwarza $M$. Pełne rozwinięcie wymaga
zademonstrowania chirurgicznej realizacji skrętu Dehna
i zestawienia jej z prezentacją grupy klas odwzorowań.
\end{proof}

\begin{definicja}[Ruchy Kirby'ego dla splotów z ramowaniem]
\emph{Przesunięcie uchwytu} (handle slide) polega na zastąpieniu
jednej składowej splotu sumą pasmową z równoległą kopią
innej składowej, z odpowiednio przeniesionym ramowaniem.
\emph{Dodanie lub usunięcie wybuchu} (blow-up/down)
polega na dodaniu lub usunięciu odseparowanej od pozostałych
składowych trywialnej składowej o ramowaniu $+1$ albo $-1$,
w wersji lokalnej również z odpowiednimi zmianami sąsiednich
ramowań i skrzyżowań przy tzw. wkręceniu.
\end{definicja}

\begin{twierdzenie}[Kirby; wersja dla chirurgii integralnej]
Dwa sploty z całkowitymi ramowaniami w $S^3$ opisują
orientowanie homeomorficzne zamknięte, orientowalne
trójrozmaitości wtedy i tylko wtedy, gdy można przejść
od jednego do drugiego skończonym ciągiem izotopii,
przesunięć uchwytów oraz dodań i usunięć rozdzielnych
trywialnych składowych o ramowaniu $\pm1$.
\end{twierdzenie}

\begin{uwaga}
Chirurgia na odseparowanej od reszty składowej trywialnej
o ramowaniu $\pm1$ nie zmienia trójrozmaitości brzegowej,
ponieważ $\pm1$-chirurgia na węźle trywialnym daje $S^3$.
Nie należy jednak mylić tego ruchu z niezmienniczością
\emph{czterorozmaitości} wypełniającej brzeg: dodanie
$\pm1$-ramowanej składowej odpowiada sumie spójnej
z $\pm\mathbb{CP}^2$ i zmienia formę przecięcia.
Właśnie to rozróżnienie będzie kluczowe w następnym rozdziale.
\end{uwaga}

\section{Hiperboliczna chirurgia Dehna}
\label{sec:kt-hyper-dehn}

\begin{twierdzenie}[Thurston; hiperboliczna chirurgia Dehna]
Niech $K\subset S^3$ będzie węzłem hiperbolicznym.
Wówczas wszystkie z wyjątkiem skończenie wielu nachyleń
$p/q$ prowadzą do zamkniętych rozmaitości
$S^3_{p/q}(K)$ dopuszczających metrykę hiperboliczną.
Ponadto dla dowolnego ciągu różnych nachyleń, których
długości w ustalonej płaskiej metryce torusa horosferycznego
dążą do nieskończoności, objętości hiperbolicznych wypełnień
dążą do $\operatorname{Vol}(K)$ od dołu.
\end{twierdzenie}

\begin{uwaga}
Twierdzenie mówi, że jedna niezwarta rozmaitość hiperboliczna
z końcem torusowym generuje nieskończenie wiele zamkniętych
rozmaitości hiperbolicznych przez zmianę nachylenia.
W wielu przykładach otrzymane rozmaitości mają różne
objętości, a więc są topologicznie różne na mocy sztywności
Mostowa--Prasada. Wynik stanowi pomost między geometrią
hiperboliczną a konstrukcyjną topologią trójrozmaitości.
\end{uwaga}



\chapter{Grafy rozkładów JSJ i grupy trójrozmaitości}
\label{ch:tg-jsj-graphs}

Graf, którego wierzchołki opisują części rozmaitości,
a krawędzie powierzchnie ich sklejania, jest jednym z
najważniejszych modeli współczesnej topologii
trójwymiarowej. Po sekcji~\ref{sec:kt-dehn}, poświęconej
chirurgii Dehna, możemy go zastosować do zewnętrzy węzłów.
Odróżniamy \emph{graf dualny rozkładu}
od całej \emph{struktury grafu przestrzeni lub grup}:
sam graf zapomina, jakie rozmaitości sklejamy i jak
identyfikujemy ich brzegi. Różnica jest analogiczna
do tej, którą poznaliśmy przy triangulacjach powierzchni.

\section{Torusy nieściśliwe i definicja grafu JSJ}
\label{sec:tg-jsj-graph-def}

\begin{definicja}[Nieściśliwy torus]
Niech $M$ będzie zwartą, spójną, orientowalną
trójrozmaitością, a $T\subset\operatorname{int}M$
dwustronnym, topologicznie zanurzonym torusem.  Nazywamy $T$ \emph{nieściśliwym},
gdy indukowany inkluzją homomorfizm
$\pi_1(T)\to\pi_1(M)$ jest injekcją.
Torus $T$ nazywamy \emph{równoległym do brzegu},
gdy ogranicza wraz z pewną składową $\partial M$
podrozmaitość homeomorficzną z $T^2\times[0,1]$.
\end{definicja}

\begin{twierdzenie}[Jaco--Shalen--Johannson; rozkład JSJ]
\label{thm:tg-jsj-theorem}
Niech $M$ będzie zwartą, spójną, orientowalną,
nieprzywiedlną trójrozmaitością, której brzeg jest
pusty lub jest sumą torusów. Istnieje skończona
rodzina parami rozłącznych, nieściśliwych,
nierównoległych do brzegu torusów, parami
nieizotopowych, taka że każda składowa rozmaitości
rozciętej wzdłuż tych torusów jest rozmaitością
Seiferta lub atoroidalną. Wśród rodzin o tej
własności rodzina minimalna jest wyznaczona
jednoznacznie z dokładnością do izotopii.
\end{twierdzenie}

\begin{uwaga}[Status twierdzenia]
Nieprzywiedlność oznacza, że każda topologicznie zanurzona
sfera $S^2\subset M$ ogranicza w $M$ kulę trójwymiarową. Rozmaitość atoroidalna
to taka, w której każdy nieściśliwy torus jest
równoległy do brzegu. Rozmaitość Seiferta posiada
rozwarstwienie na okręgi z dopuszczalnymi włóknami
wyjątkowymi. Samo istnienie i jednoznaczność
minimalnego rozkładu JSJ przyjmujemy jako głęboki
wynik topologii trójwymiarowej. Nie należy go
wywodzić jedynie z klasycznego twierdzenia
van Kampena.
\end{uwaga}

\begin{definicja}[Graf dualny rozkładu JSJ]
Niech $\mathcal T=\{T_1,\ldots,T_m\}$ będzie
minimalną rodziną torusów z twierdzenia JSJ,
a $M\mid\mathcal T$ rozmaitością z brzegiem,
otrzymaną przez rozcięcie $M$ wzdłuż tych torusów.
\emph{Graf JSJ} $\Gamma_{\mathrm{JSJ}}(M)$ ma po jednym
wierzchołku dla każdej składowej $M\mid\mathcal T$
oraz po jednej krawędzi dla każdego $T_i$;
końce krawędzi odpowiadają składowym przylegającym
do dwóch stron $T_i$. Jeżeli obie strony leżą
w jednej składowej, otrzymujemy pętlę.
\end{definicja}

\begin{stwierdzenie}[Spójność grafu rozkładu]
Dla każdej rozmaitości $M$ spełniającej założenia
poprzedniej definicji graf $\Gamma_{\mathrm{JSJ}}(M)$
jest spójny. Gdy $m=0$, ma dokładnie jeden
wierzchołek i żadnej krawędzi.
\end{stwierdzenie}
\begin{proof}
Niech $M_i,M_j$ będą dwiema składowymi rozciętej
rozmaitości. Wybieramy punkty $p_i\in\operatorname{int}M_i$
i $p_j\in\operatorname{int}M_j$. Ponieważ $M$ jest
spójna i łukowo spójna, istnieje ścieżka $\gamma$
łącząca ich obrazy w $M$. Po małej perturbacji
możemy założyć, że $\gamma$ jest poprzeczna do
każdego torusa i przecina rodzinę $\mathcal T$
w skończenie wielu punktach. Kolejne fragmenty
ścieżki leżą w składowych rozciętej rozmaitości,
a każde przecięcie torusa odpowiada krawędzi grafu.
Otrzymujemy spacer z $M_i$ do $M_j$, co dowodzi
spójności. Jeżeli torusów nie ma, rozcięta
rozmaitość pozostaje równa spójnej $M$.
\end{proof}

\begin{center}
\begin{tikzpicture}[scale=.95,line cap=round]
 \node[tgBlueVertex] (A) at (-2,0) {};
 \node[tgTealVertex] (B) at (0,.85) {};
 \node[tgRedVertex] (C) at (2,0) {};
 \draw[tgEdge] (A)--node[above left,font=\scriptsize] {$T_1$}(B)
                   (B)--node[above right,font=\scriptsize] {$T_2$}(C);
 \node[font=\scriptsize,text=tgNavy] at (-2,-.62) {$M_1$};
 \node[font=\scriptsize,text=tgNavy] at (0,1.43) {$M_2$};
 \node[font=\scriptsize,text=tgNavy] at (2,-.62) {$M_3$};
 \node[font=\footnotesize,text=tgNavy] at (0,-1.12)
   {Wierzchołki oznaczają części; krawędzie -- torusy sklejania.};
\end{tikzpicture}
\end{center}

\section{Grupa fundamentalna z grafu JSJ}
\label{sec:tg-jsj-groups}

\begin{stwierdzenie}[Rozdzielający torus i amalgamacja]
Niech $M=M_1\cup_T M_2$, gdzie $M_1,M_2$ są
spójnymi trójrozmaitościami, a $T$ jest wspólnym,
spójnym torusem brzegowym. Załóżmy, że inkluzje
$T\hookrightarrow M_i$ indukują monomorfizmy na
$\pi_1$. Wtedy, po zgodnym wyborze punktów i ścieżek
bazowych,
\[
\pi_1(M)\cong\pi_1(M_1)
 *_{\pi_1(T)}\pi_1(M_2),
\qquad \pi_1(T)\cong\mathbb Z^2.
\]
\end{stwierdzenie}
\begin{proof}
Wybieramy kołnierz $T\times(-\varepsilon,\varepsilon)$
i nieco powiększamy $M_1,M_2$ o odpowiednie części
kołnierza, aby powstało pokrycie $M$ dwoma
otwartymi, łukowo spójnymi zbiorami. Ich część
wspólna deformacyjnie retraktuje się na $T$, a
każdy zbiór na odpowiednie $M_i$. W twierdzeniu
van Kampena dostajemy prezentację z generatorami
z obu grup $\pi_1(M_i)$ i relacjami utożsamiającymi
obrazy elementów $\pi_1(T)$. Ponieważ założyliśmy
injektywność obu odwzorowań, otrzymany iloraz
iloczynu wolnego jest dokładnie iloczynem wolnym
z amalgamacją nad $\mathbb Z^2$.
\end{proof}

\begin{uwaga}[Torus nierozdzielający]
Jeśli po rozcięciu wzdłuż nieściśliwego torusa
$T$ otrzymujemy spójną rozmaitość $N$ o dwóch
wyróżnionych składowych toroidalnego brzegu,
to sklejenie ich z powrotem jest konstrukcją
analogiczną do przypadku HNN z rozdziału o
Bassie--Serre'ze. Jego stabilna litera odpowiada
pętli przechodzącej przez sklejany torus.
Cały graf JSJ, wyposażony w grupy
$\pi_1(M_v)$ i $\pi_1(T_e)$ oraz odpowiednie
monomorfizmy, jest grafem grup. Nie wystarcza
sam nieoznakowany graf $\Gamma_{\mathrm{JSJ}}(M)$.
\end{uwaga}

\section{Przykłady: węzły hiperboliczne, torusowe i satelitarne}
\label{sec:tg-jsj-knots}

\begin{przyklad}[Trzy różne geometrie, różne dane grafowe]
Niech $K\subset S^3$ będzie węzłem, a
$X_K=S^3\setminus\operatorname{int}N(K)$ jego
zewnętrzem. Dla węzła ósemkowego $X_K$ jest
hiperboliczne o skończonej objętości i atoroidalne;
jego graf JSJ ma jeden wierzchołek i żadnej krawędzi.
Dla nietrywialnego węzła torusowego $T(p,q)$
zewnętrze jest rozmaitością Seiferta i również
ma graf JSJ z jednym wierzchołkiem. Sam graf nie
rozróżnia więc obu przypadków: potrzebna jest
etykieta rodzaju części oraz jej topologia.

Węzeł satelitarny ma w swoim zewnętrzu torus nieściśliwy,
który nie jest równoległy do toroidalnego brzegu.
W wielu przykładach, np. dla odpowiednio
nietrywialnego kabla węzła hiperbolicznego,
torus otaczający węzeł towarzyszący należy do
kanonicznego rozkładu JSJ i graf ma co najmniej jedną
krawędź: łączy część kablową z zewnętrzem
towarzysza. Po oznaczeniu części kablowej przez $P$
i zewnętrza towarzysza przez $X_J$ dostajemy przy
odpowiednich punktach bazowych
$\pi_1(X_K)\cong\pi_1(P)*_{\mathbb Z^2}\pi_1(X_J)$.
Torus krawędziowy wnosi dwa generatory, meridian i długość,
których obrazy po obu stronach sklejania zależą od typu kabla. Wyróżnienie grup obwodowych
$\pi_1(T^2)\cong\mathbb Z^2$ jest konieczne,
aby odzyskać sposób sklejenia tych części.
\end{przyklad}

\begin{uwaga}[Graf a geometria trójrozmaitości]
Rozkład JSJ ma szczególne znaczenie po
geometryzacji Thurstona. Części atoroidalne,
które nie są rozmaitościami Seiferta, w
odpowiednich założeniach hiperbolizacji
posiadają wnętrza z zupełną metryką
hiperboliczną o skończonej objętości.
Graf dualny porządkuje części, zaś grupy,
parametry torusów brzegowych i homeomorfizmy
sklejania zapisują dodatkowe dane topologiczne.
To właśnie te dane łączą teorię grafów z
chirurgią Dehna i rachunkiem Kirby'ego. Sama
izomorficzna struktura abstrakcyjnych grup na wierzchołkach
i krawędziach nie odzyskuje automatycznie pełnych map
sklejania torusów; do rekonstrukcji geometrycznej
należy zachować dane peryferyjne i ich utożsamienia.
\end{uwaga}

\begin{uwaga}[Literatura]
Klasyczne twierdzenie o rozkładzie pochodzi z prac
W. Jaco--P. Shalena oraz K. Johannsona.
Dalsze przykłady i strukturę grup fundamentalnych
można znaleźć w A. Aschenbrenner,
S. Friedl, H. Wilton, \emph{3-Manifold Groups}.
\end{uwaga}


\chapter{Węzły i rozmaitości czterowymiarowe}
\label{ch:kt-cztery}

Chirurgia Dehna opisuje zamknięte trójrozmaitości jako brzegi
czterorozmaitości zbudowanych z uchwytów. Ta obserwacja jest
punktem wyjścia rachunku Kirby'ego, teorii form przecięcia,
problemu rozpoznawania egzotycznych struktur gładkich oraz
teorii kobordyzmu węzłów. Poniżej precyzujemy konstrukcje
od podstaw.

\section{Uchwyt dwuwymiarowy w wymiarze cztery}
\label{sec:kt-uchwyty}

\begin{definicja}[Uchwyt indeksu dwa]
\emph{Uchwytem indeksu dwa w wymiarze cztery} nazywamy
$h^2=D^2\times D^2$. Jego \emph{obszar przyczepienia}
to $\partial D^2\times D^2\cong S^1\times D^2$,
a \emph{rdzeń} to $D^2\times\{0\}$.
Jeżeli $W$ jest zwartą, orientowaną, gładką czterorozmaitością
z brzegiem, to dołączenie uchwytu $h^2$ do $W$ wymaga
wyboru zanurzenia z ramowaniem
$\phi\colon S^1\times D^2\hookrightarrow\partial W$.
Wynikiem jest $W\cup_\phi(D^2\times D^2)$, po wygładzeniu
narożników.
\end{definicja}

Gdy $W=B^4$, krzywa $\phi(S^1\times\{0\})\subset\partial B^4=S^3$
jest węzłem $K$. Ramowanie decyduje o sposobie utożsamienia
otoczenia $K$ z $S^1\times D^2$. Tak więc \emph{węzeł wraz
z liczbą całkowitą} koduje konkretny sposób dołączenia
uchwytu indeksu dwa.

\begin{twierdzenie}[Brzeg po dołączeniu uchwytu]
Niech $K\subset S^3=\partial B^4$ będzie zorientowanym węzłem,
a $n\in\mathbb Z$ jego całkowitym ramowaniem względem długości
preferowanej. Niech
$W_n(K):=B^4\cup_{(K,n)}(D^2\times D^2)$.
Wtedy $\partial W_n(K)$, z orientacją brzegową, jest
orientowanie dyfeomorficzny z rozmaitością $S^3_n(K)$
otrzymaną przez integralną chirurgię Dehna o współczynniku $n$.
\end{twierdzenie}

\begin{proof}
Przy dołączeniu uchwytu usuwamy z $\partial B^4$
wnętrze obszaru przyczepienia $S^1\times D^2$.
Pozostałą część brzegu nowego uchwytu stanowi
$D^2\times S^1$, czyli torus pełny. Jest on przyklejany
do torusa $\partial N(K)$ w taki sposób, że jego
meridian $\partial D^2\times\{*\}$ odpowiada krzywej
o klasie $n[\mu]+[\lambda]$, wynikającej z ramowania
$K$. Jest to dokładnie definicja integralnej chirurgii
Dehna.
\end{proof}

\begin{center}
\begin{tikzpicture}[>=Stealth,scale=.90,line cap=round]
  \node[draw=ktBlue,line width=1.5pt,rounded corners=8pt,
        minimum width=2.4cm,minimum height=1.32cm,
        fill=ktBlue!10,text=ktInk] (b) at (-3.0,0) {$B^4$};
  \node[draw=ktPurple,line width=1.5pt,rounded corners=8pt,
        minimum width=3.2cm,minimum height=1.32cm,
        fill=ktPurple!10,text=ktInk] (h) at (3.0,0) {$D^2\times D^2$};
  \draw[->,ktTeal,line width=2.4pt] (b)--(h);
  \node[text=ktInk,font=\footnotesize,fill=white,inner sep=2pt] at (0,1.00)
      {$\phi\colon S^1\times D^2\hookrightarrow S^3$};
  \node[align=center,text=ktInk,font=\small] at (0,-1.18)
    {węzeł przyczepienia $K\subset S^3$ i ramowanie $n\in\mathbb Z$};
  \node[rounded corners=4pt,fill=ktGold!16,inner sep=6pt,
        text=ktInk,font=\small] at (0,-2.0)
    {$\partial W_n(K)\cong S^3_n(K)$};
\end{tikzpicture}
\end{center}

\section{Forma przecięcia a macierz splecenia}
\label{sec:kt-intersection}

\begin{twierdzenie}[Forma przecięcia z diagramu uchwytowego]
Niech $L=K_1\cup\dots\cup K_r\subset S^3$ będzie
zorientowanym splotem z całkowitymi ramowaniami
$n_1,\dots,n_r$. Niech $W_L$ będzie gładką, zwartą,
orientowaną czterorozmaitością otrzymaną z $B^4$
przez dołączenie do niego $r$ uchwytów indeksu dwa
wzdłuż $(L;n_1,\dots,n_r)$. Wtedy:
\begin{enumerate}[label=(\roman*)]
\item $\pi_1(W_L)=0$;
\item $H_2(W_L;\mathbb Z)\cong\mathbb Z^r$;
\item w bazie wyznaczonej przez rdzenie uchwytów po
uzupełnieniu ich powierzchniami Seiferta forma przecięcia
$Q_{W_L}$ jest reprezentowana przez macierz
$Q_L$ o przekątnej $(n_1,\dots,n_r)$ i elementach
pozadiagonalnych $\operatorname{lk}(K_i,K_j)$.
\end{enumerate}
\end{twierdzenie}

\begin{proof}
Zacznijmy od grupy fundamentalnej. Dla dołączenia
pojedynczego uchwytu indeksu dwa twierdzenie van Kampena
identyfikuje grupę nowej przestrzeni z ilorazem grupy
poprzedniej przez normalną podgrupę generowaną przez
klasę krzywej przyczepienia. Ponieważ $\pi_1(B^4)=1$,
po wszystkich dołączeniach nadal mamy $\pi_1(W_L)=1$.

Następnie skorzystajmy z względnego kompleksu komórkowego
rozkładu na uchwyty pary $(W_L,B^4)$.
W tej parze występuje dokładnie $r$ komórek wymiaru dwa,
nie ma natomiast względnych komórek wymiarów jeden i trzy.
Stąd
$C_2(W_L,B^4;\mathbb Z)\cong\mathbb Z^r$,
$C_1(W_L,B^4;\mathbb Z)=C_3(W_L,B^4;\mathbb Z)=0$,
a wszystkie różniczki są zerowe. Zatem
$H_2(W_L,B^4;\mathbb Z)\cong\mathbb Z^r$.
W odpowiednim fragmencie ciągu długiego pary mamy
\[
H_2(B^4)\longrightarrow H_2(W_L)
\longrightarrow H_2(W_L,B^4)
\longrightarrow H_1(B^4).
\]
Skrajne grupy są zerowe, więc środkowe odwzorowanie
jest izomorfizmem. Otrzymujemy
$H_2(W_L;\mathbb Z)\cong\mathbb Z^r$.

Aby obliczyć przecięcia, dla każdego $K_i$ wybieramy
powierzchnię Seiferta $F_i\subset S^3$. Jej wnętrze
przesuwamy nieznacznie do wnętrza $B^4$, a następnie
dołączamy rdzeń $D^2\times\{0\}$ odpowiedniego uchwytu.
Otrzymujemy zamkniętą powierzchnię reprezentującą generator
$H_2(W_L)$. Dla $i\ne j$ wybieramy reprezentantów w ogólnym położeniu.
Przesuwamy wnętrze $F_i$ do kołnierza $S^3$ na innej
głębokości niż wnętrze $F_j$; przecinania z rdzeniem
$j$-tego uchwytu zostają policzone przez przecięcia
$F_i$ z $K_j$ w $S^3$. Ich algebraiczna suma jest
$\operatorname{lk}(K_i,K_j)$. Dla $i=j$ bierzemy
równoległą kopię domkniętej powierzchni $F_i$; liczba
jej przecięć z powierzchnią pierwotną jest przesunięciem
normalnym zapisanym przez ramowanie $n_i$. Zatem macierz przecięć jest równa $Q_L$.
\end{proof}

\begin{wniosek}
Jeżeli $W_L$ jest jak wyżej, to $\partial W_L$ jest
homologiczną sferą trójwymiarową, tzn.
$H_*(\partial W_L;\mathbb Z)\cong H_*(S^3;\mathbb Z)$,
wtedy i tylko wtedy, gdy $\det Q_L=\pm1$.
\end{wniosek}

\begin{proof}
W szczególności rozpiszmy zależność, której używamy.
Dla pary $(W_L,Y)$, gdzie $Y=\partial W_L$, mamy
fragment ciągu dokładnego
\[
H_2(W_L;\mathbb Z)\xrightarrow{j}
H_2(W_L,Y;\mathbb Z)\longrightarrow
H_1(Y;\mathbb Z)\longrightarrow H_1(W_L;\mathbb Z)=0.
\]
Dualność Poincarégo--Lefschetza i uniwersalne twierdzenie
o współczynnikach identyfikują obie pierwsze grupy z
$\mathbb Z^r$. W bazach uchwytowych homomorfizm $j$
jest reprezentowany przez macierz przecięcia $Q_L$.
Zatem $H_1(Y;\mathbb Z)\cong\operatorname{coker}Q_L$.
Grupa ta znika dokładnie wtedy, gdy $Q_L$ jest unimodularna. Dla zamkniętej spójnej orientowalnej
trójrozmaitości dualność Poincarégo, uniwersalne twierdzenie
o współczynnikach i orientowalność pokazują, że z
$H_1=0$ wynika $H_2=0$, a $H_0=H_3=\mathbb Z$.
\end{proof}

\begin{przyklad}[Trzy klasyczne formy]
Dla pojedynczej składowej trywialnej o ramowaniu $+1$
mamy macierz przecięcia $Q=[1]$. Jej wyznacznik jest równy $1$,
więc brzeg jest homologiczną sferą trójwymiarową; w istocie
chodzi o $S^3$, a po dołączeniu uchwytu indeksu cztery
otrzymujemy $\mathbb{CP}^2$.
Dla ramowania $-1$ dostajemy analogicznie
$\overline{\mathbb{CP}}^{\,2}$ z formą przecięcia $[-1]$.

Dla standardowego splotu Hopfa z dwoma ramowaniami zero mamy
\[
Q=\begin{pmatrix}0&1\\1&0\end{pmatrix}.
\]
Macierz ta ma wyznacznik $-1$, a więc ponownie brzeg jest
homologiczną sferą trójwymiarową. Jej wartości własne wynoszą
$1$ i $-1$, zatem forma jest hiperboliczna i ma sygnaturę zero.
Po odpowiednim dołączeniu uchwytów wyższych indeksów dostajemy
rozkład na uchwyty rozmaitości $S^2\times S^2$.

Ten przykład pokazuje, jak bezpośrednio z diagramu odczytać zarówno
homologię brzegu, bo $H_1(\partial W;\mathbb Z)\cong\operatorname{coker}Q$,
jak i informację o formie przecięcia we wnętrzu. Dla splotu Hopfa
z ramowaniami $n_1,n_2$ otrzymalibyśmy ogólniej macierz
$\begin{psmallmatrix}n_1&1\\1&n_2\end{psmallmatrix}$.
Jej wyznacznik równy $n_1n_2-1$ kontroluje pierwszą homologię brzegu,
a znak własnych wartości decyduje o sygnaturze formy.
\end{przyklad}

\section{Przesuwanie uchwytów i operacja blow-up}
\label{sec:kt-slides}

\begin{stwierdzenie}
Niech $W_L$ będzie czterorozmaitością zbudowaną z kuli
$B^4$ przez dołączenie uchwytów indeksu dwa opisanych
przez ramowany splot $L$ z macierzą przecięcia $Q_L$.
Przesunięcie jednego uchwytu po innym realizuje zmianę
bazy w $H_2(W_L;\mathbb Z)$ postaci
$e_i\mapsto e_i\pm e_j$, pozostawiając pozostałe
wektory bez zmian. Odpowiada to kongruencji całkowitej
$Q_L\mapsto P^TQ_LP$ dla pewnej macierzy
$P\in GL_r(\mathbb Z)$.
\end{stwierdzenie}

\begin{proof}
Przesunięcie uchwytu indeksu dwa zastępuje jego rdzeń
sumą pasmową z równoległą kopią rdzenia drugiego uchwytu.
Po domknięciu powierzchni w $B^4$ odpowiada to sumie
klas homologii $e_i\pm e_j$. W nowej bazie biliniowa
forma przecięcia przyjmuje macierz $P^TQ_LP$.
Sama operacja przesunięcia jest realizowana dyfeomorfizmem
czterorozmaitości zachowującym zewnętrzny brzeg.
\end{proof}

\begin{przyklad}[Zmiana wszystkich współczynników przy przesunięciu]
Niech dwie składowe ramowanego splotu mają ramowania
$n_i,n_j$ i liczbę splecenia $\ell$. W bazie klas
rdzeni uchwytów macierz przecięcia wynosi
$Q=\begin{psmallmatrix}n_i&\ell\\\ell&n_j\end{psmallmatrix}$.
Przesunięcie $i$-tego uchwytu po $j$-tym odpowiada
zamianie $e_i'=e_i+e_j$, $e_j'=e_j$, a więc
\[
Q'=\begin{pmatrix}n_i+2\ell+n_j&\ell+n_j\\
\ell+n_j&n_j\end{pmatrix}.
\]
Nowe ramowanie $i$-tej składowej nie jest zatem zwykłą
sumą $n_i+n_j$: pojawia się dodatkowy składnik
$2\operatorname{lk}(K_i,K_j)$. Rachunek wynika z
równości $(e_i+e_j)^2=e_i^2+2(e_i\cdot e_j)+e_j^2$.
\end{przyklad}

\begin{stwierdzenie}[Blow-up]
Niech $W$ będzie gładką, orientowaną czterorozmaitością.
Po wykonaniu we wnętrzu $W$ operacji blow-up otrzymujemy
$W\#\overline{\mathbb{CP}}^{\,2}$; po zmianie znaku
operacji odpowiednio $W\#\mathbb{CP}^2$.
Forma przecięcia zmienia się przez sumę prostą
$Q_W\oplus[-1]$ albo $Q_W\oplus[1]$,
przy zachowaniu interpretacji formy na odpowiednich
klasach względnych, jeśli $W$ ma brzeg.
\end{stwierdzenie}

\begin{proof}
Usuwamy z wnętrza $W$ standardową kulę $B^4$ i wklejamy
$\mathbb{CP}^2\setminus\operatorname{int}B^4$ lub
$\overline{\mathbb{CP}}^{\,2}\setminus\operatorname{int}B^4$.
Definicja sumy spójnej daje opis topologiczny.
Ponieważ $H_2(S^3)=H_1(S^3)=0$, ciąg Mayera--Vietorisa
rozkłada drugą homologię na sumę prostą, a reprezentanci
pochodzący z obu składników mogą być wybrani rozłącznie.
Wynikowa forma jest zatem sumą prostą form składników.
\end{proof}

\begin{uwaga}[Granica rachunku macierzowego]
Dwie macierze przecięcia całkowicie izometryczne nie muszą
pochodzić od dyfeomorficznych czterorozmaitości.
Rachunek na macierzach zapomina dodatkowe dane zanurzeń,
gładkości i uchwytów wyższych indeksów. To jeden z powodów,
dla których wymiar cztery jest szczególny: klasyfikacja
samych form przecięcia jest znacznie łatwiejsza niż
klasyfikacja gładkich czterorozmaitości.
\end{uwaga}

\section{Węzły jako brzegi powierzchni w \texorpdfstring{$B^4$}{B4}}
\label{sec:kt-slice}

\begin{definicja}[Rodzaj czterowymiarowy]
Niech $K\subset S^3=\partial B^4$ będzie zorientowanym węzłem.
\emph{Gładkim rodzajem czterowymiarowym} węzła nazywamy
\[
\begin{aligned}
g_4^{\mathrm{sm}}(K):=\min\big\{g(F):\;&F\subset B^4
\text{ jest spójną, zwartą,}\\
&\text{orientowaną, gładko i właściwie zanurzoną powierzchnią},\\
&\partial F=K\big\}.
\end{aligned}
\]
Analogicznie \emph{topologiczny rodzaj czterowymiarowy}
$g_4^{\mathrm{top}}(K)$ definiujemy, dopuszczając powierzchnie
topologicznie lokalnie płasko i właściwie zanurzone.
Węzeł jest \emph{gładko slice}, gdy $g_4^{\mathrm{sm}}(K)=0$,
a \emph{topologicznie slice}, gdy $g_4^{\mathrm{top}}(K)=0$.
\end{definicja}

\begin{stwierdzenie}
Dla każdego węzła $K\subset S^3$ zachodzi
$0\leq g_4^{\mathrm{top}}(K)\leq g_4^{\mathrm{sm}}(K)\leq g(K)$.
\end{stwierdzenie}

\begin{proof}
Każda powierzchnia gładko zanurzona jest lokalnie płaska,
więc minimum topologiczne nie przekracza gładkiego.
Każdą powierzchnię Seiferta w $S^3$ można przesunąć
nieznacznie w głąb kołnierza $S^3\times[0,\varepsilon)$
we wnętrzu $B^4$, pozostawiając jej brzeg na $S^3$.
Dostajemy właściwie zanurzoną powierzchnię tego samego
rodzaju, stąd $g_4^{\mathrm{sm}}(K)\leq g(K)$.
Nieujemność jest oczywista.
\end{proof}

\begin{definicja}[Kobordyzm węzłów, czyli konkordancja]
Niech $K_0,K_1\subset S^3$ będą zorientowanymi węzłami.
\emph{Gładką konkordancją} między $K_0$ i $K_1$ nazywamy
gładkie, właściwe zanurzenie
$A\colon S^1\times[0,1]\hookrightarrow S^3\times[0,1]$
takie, że $A(S^1\times\{i\})=K_i\times\{i\}$ dla $i=0,1$
oraz orientacja $A$ indukuje na brzegu $K_1\sqcup(-K_0)$.
Analogicznie definiujemy \emph{konkordancję topologiczną}
przez lokalnie płaskie zanurzenie pierścienia.
\end{definicja}

\begin{stwierdzenie}
Dwa zorientowane węzły $K_0,K_1\subset S^3$ są gładko
konkordantne wtedy i tylko wtedy, gdy
$K_0\#(-K_1)$ jest gładko slice, gdzie $-K_1$
oznacza odbicie lustrzane $K_1$ z odwróconą orientacją.
Analogiczna równoważność zachodzi w kategorii topologicznej
lokalnie płaskiej.
\end{stwierdzenie}

\begin{proof}[Idea argumentu]
Przecinamy konkordancję $A\subset S^3\times[0,1]$
wzdłuż łuku poprzecznego i przenosimy otrzymaną powierzchnię
do $B^4$ przez standardową identyfikację odpowiedniego
uzupełnienia kołnierza. Brzegiem powstałego dysku jest
$K_0\#(-K_1)$. Odwrotnie, z dysku ograniczonego przez
$K_0\#(-K_1)$ usuwamy otoczenie pasma realizującego sumę
spójną; odpowiednie rozcięcie daje pierścień między
$K_0$ i $K_1$ w $S^3\times[0,1]$.
Pełna wersja wymaga wygładzenia narożników i kontroli
orientacji brzegowych.
\end{proof}

\begin{definicja}[Grupa konkordancji]
Niech $\mathcal C^{\mathrm{sm}}$ oznacza zbiór klas gładkiej
konkordancji zorientowanych węzłów w $S^3$.
Operacją jest suma spójna $[K_1]+[K_2]:=[K_1\#K_2]$.
Elementem neutralnym jest klasa węzła trywialnego,
a elementem przeciwnym do $[K]$ jest $[-K]$.
Tak zdefiniowana grupa abelowa nazywa się
\emph{gładką grupą konkordancji węzłów}.
\end{definicja}

\begin{proof}[Uzasadnienie struktury grupy]
Suma spójna jest łączna i przemienna z dokładnością
do izotopii, a więc także z dokładnością do konkordancji.
Węzeł $K\#(-K)$ jest ribbon, a zatem gładko slice:
w standardowym diagramie parę przeciwnych kopii można
połączyć pasmami i uprościć do dysku z osobliwościami
wstęgowymi, które wypychamy do $B^4$. To pokazuje, że
$[-K]$ jest elementem przeciwnym. Klasa węzła trywialnego
jest neutralna.
\end{proof}

\section{Węzły ribbon i nierówności rodzaju}
\label{sec:kt-ribbon}

\begin{definicja}[Węzeł ribbon]
Węzeł $K\subset S^3$ jest \emph{ribbon}, jeżeli ogranicza
immersję dysku w $S^3$, której wszystkie osobliwości
są poprzecznymi łukami podwójnymi typu wstęgowego:
przeciwobraz każdego takiego łuku składa się z jednego
łuku właściwie zanurzonego w dysku i jednego łuku leżącego
we wnętrzu dysku. Równoważnie $K$ ogranicza w $B^4$
gładki dysk, dla którego funkcja radialna jest funkcją
Morse'a bez lokalnych maksimów we wnętrzu dysku.
\end{definicja}

\begin{stwierdzenie}
Każdy węzeł ribbon jest gładko slice.
\end{stwierdzenie}

\begin{proof}
W zanurzonym dysku typu ribbon przesuwamy lokalnie jedną z gałęzi
przecinających się w stronę wnętrza $B^4$ zgodnie
z wysokością radialną. Ponieważ osobliwości mają typ wstęgowy,
przesunięcie eliminuje wszystkie łuki podwójne bez tworzenia
nowych, a otrzymany dysk jest gładko zanurzony w $B^4$.
Brzeg pozostaje równy $K$.
\end{proof}

\begin{uwaga}[Hipoteza slice--ribbon]
Odwrotność poprzedniego stwierdzenia, tzn. pytanie,
czy każdy gładko slice węzeł jest ribbon, to klasyczna
hipoteza Foxa. Jest to problem otwarty w klasycznej
teorii konkordancji w $S^3$.
Nie należy go mylić z implikacją topologicznie slice
$\Rightarrow$ gładko slice, która jest fałszywa.
\end{uwaga}

\begin{twierdzenie}[Nierówność Murasugiego dla rodzaju czterowymiarowego]
Dla każdego zorientowanego węzła $K\subset S^3$
zachodzi $|\sigma(K)|\leq2g_4^{\mathrm{sm}}(K)$.
W szczególności, jeżeli $K$ jest gładko slice, to
$\sigma(K)=0$.
\end{twierdzenie}

\begin{proof}[Redukcja do podprzestrzeni izotropowej]
Niech $F$ będzie powierzchnią Seiferta węzła $K$ o
rodzaju $g$, a $G\subset B^4$ właściwie zanurzoną
zorientowaną powierzchnią o brzegu $K$ i rodzaju $h$.
Rozważmy symetryczną formę $S=V+V^T$ na
$H_1(F;\mathbb R)$, gdzie $V$ jest macierzą Seiferta.

Potrzebujemy klasycznego lematu Murasugiego:
\emph{przy powyższych założeniach istnieje rzeczywista
podprzestrzeń $L\subset H_1(F;\mathbb R)$
o wymiarze co najmniej $g-h$, na której forma $V$
znika identycznie}. Dowodzi się go przez połączenie
$F$ i $G$ w zamkniętą powierzchnię w $S^4$, wybranie
ograniczającej ją trójrozmaitości oraz zastosowanie
ciągu dokładnego pary i dualności Poincarégo--Lefschetza.
W przypadku $h=0$ daje on metabolizer opisany wcześniej.

Na $L$ znika również symetryzacja $S$. Ponieważ
$\det S=\pm\det(K)\ne0$, forma $S$ jest niezdegenerowana. Niech więc
$s_+$ i $s_-$ będą liczbami dodatnich i ujemnych wartości
własnych $S$, dla $S$. Dla niezdegenerowanej rzeczywistej formy symetrycznej
maksymalny wymiar całkowicie izotropowej
podprzestrzeni wynosi $\min(s_+,s_-)$. Ponieważ
wymiar całej przestrzeni to $2g=s_++s_-$, otrzymujemy
\[
|\operatorname{sign}S|=|s_+-s_-|
\le 2g-2\dim L\le2h.
\]
Minimalizując $h$ po wszystkich gładkich powierzchniach
w $B^4$, dostajemy $|\sigma(K)|\le2g_4^{\mathrm{sm}}(K)$.
\end{proof}

\begin{uwaga}
Ważne jest, że w twierdzeniu występuje rodzaj w $B^4$,
a nie tylko rodzaj powierzchni Seiferta w $S^3$.
Dowód wykorzystuje formy Seiferta i ograniczenie
na wymiar podprzestrzeni izotropowej pochodzącej
z powierzchni zanurzonej w $B^4$; jest istotnie głębszy
od prostego oszacowania $|\sigma(K)|\leq2g(K)$.
\end{uwaga}

\begin{twierdzenie}[Fox--Milnor]
Jeśli $K\subset S^3$ jest gładko slice, to istnieje
wielomian Laurenta $f(t)\in\mathbb Z[t,t^{-1}]$ taki, że
$\Delta_K(t)\doteq f(t)f(t^{-1})$,
gdzie $\doteq$ oznacza równość z dokładnością do
mnożenia przez $\pm t^k$.
\end{twierdzenie}

\begin{uwaga}
Warunek Foxa--Milnora nie jest wystarczający do gładkiej
slicowości. Ograniczenie to jest algebraiczne i nie wykrywa
wielu obstrukcji pochodzących z gładkiej topologii
czterowymiarowej. Zobaczymy niżej przykład o
$\Delta_K(t)=1$, który jest topologicznie slice,
ale nie jest gładko slice.
\end{uwaga}

\begin{definicja}[Metabolizer formy Seiferta]
Niech $F$ będzie powierzchnią Seiferta węzła $K\subset S^3$
o rodzaju $g$ i niech $V$ będzie macierzą jej formy Seiferta.
\emph{Metabolizerem} nazywamy podgrupę $L\subset H_1(F;\mathbb Z)$
która jest składnikiem prostym rzędu $g$ oraz na której
forma Seiferta znika identycznie, tj.
$\operatorname{lk}(x^+,y)=0$ dla wszystkich $x,y\in L$.
Węzeł nazywamy \emph{algebraicznie slice}, jeżeli jego
forma Seiferta dopuszcza metabolizer.
\end{definicja}

\begin{twierdzenie}[Slicowość gładka implikuje algebraiczną]
Niech $K\subset S^3$ będzie węzłem, który ogranicza gładko
zanurzony dysk $D\subset B^4$. Wtedy dla każdej powierzchni
Seiferta $F$ węzła $K$ forma Seiferta na $H_1(F;\mathbb Z)$
posiada metabolizer.
\end{twierdzenie}

\begin{proof}[Argument przez rozmaitość trójwymiarową]
Niech $F$ ma rodzaj $g$. Łącząc $F$ z dyskiem $D$
wzdłuż wspólnego brzegu, otrzymujemy zamkniętą,
zorientowaną powierzchnię $\Sigma$ rodzaju $g$ w
czterowymiarowej sferze powstałej po doklejeniu kuli
do $B^4$. Powierzchnia $\Sigma$ ogranicza
odpowiednią trójwymiarową rozmaitość Seiferta $M$;
w wersji dostosowanej do rozcięcia na dwie kule jej
brzeg jest sumą $F$ i dysku $D$ z przeciwnymi orientacjami.
Istnienie takiej rozmaitości $M$ jest klasycznym
lematem o hiperpłaskiej powierzchni Seiferta dla
zanurzenia kodymensji dwa.

Po zastosowaniu ciągu dokładnego pary $(M,\partial M)$
i dualności Poincarégo--Lefschetza otrzymujemy zasadę
\emph{połowa żyje, połowa ginie}: jądro
$H_1(\partial M;\mathbb Q)\to H_1(M;\mathbb Q)$
ma wymiar $g$ i jest lagranżowskie względem
formy przecięcia. Jej część całkowita, po wysyceniu
w $H_1(F;\mathbb Z)$, jest składnikiem prostym rzędu $g$.
Jeśli $x,y$ należą do tego składnika, to odpowiednie
wielokrotności krzywych ograniczają 2-łańcuchy w $M$.
Przesuwając jedną z nich w dodatnim kierunku normalnym
do $F$, można wybrać łańcuchy bez przecięć z drugą krzywą.
Stąd $\operatorname{lk}(x^+,y)=0$.
Otrzymany lagranżowski składnik jest metabolizerem.
\end{proof}

\begin{stwierdzenie}[Algebraiczna przyczyna faktoryzacji Foxa--Milnora]
Niech $K\subset S^3$ będzie algebraicznie slice,
a $V$ macierzą Seiferta rzędu $2g$.
Wtedy istnieje $f(t)\in\mathbb Z[t,t^{-1}]$ takie, że
$\det(tV-V^T)\doteq f(t)f(t^{-1})$.
\end{stwierdzenie}

\begin{proof}
Wybieramy całkowitą bazę $H_1(F;\mathbb Z)$,
której pierwsze $g$ elementów tworzy bazę metabolizera.
W tej bazie macierz Seiferta ma postać blokową
$V=\begin{psmallmatrix}0&A\\B&C\end{psmallmatrix}$,
gdzie bloki są kwadratowe rozmiaru $g$.
Stąd
\[
tV-V^T=\begin{pmatrix}
0&tA-B^T\\tB-A^T&tC-C^T
\end{pmatrix}.
\]
Rozwinięcie wyznacznika względem pierwszych $g$ wierszy
daje
$\det(tV-V^T)=(-1)^g\det(tA-B^T)\det(tB-A^T)$.
Niech $f(t):=\det(tA-B^T)$. Z transpozycji wynika
$\det(tB-A^T)=(-t)^g f(t^{-1})$.
Ostatecznie
$\Delta_K(t)=t^g f(t)f(t^{-1})$,
co dowodzi tezy z dokładnością do jednostki $t^g$.
\end{proof}

\begin{uwaga}
Klasyczna teoria algebraicznej konkordancji pokazuje,
że gładko slice węzeł jest algebraicznie slice.
Samo istnienie metabolizera nie gwarantuje jednak
zanurzonego dysku w $B^4$. Wyprowadzone twierdzenie
odtwarza zatem algebraiczną część przeszkody Foxa--Milnora,
ale nie jej odwrotność.
\end{uwaga}


\section{Twierdzenie Freedmana i rozdzielenie kategorii TOP/DIFF}
\label{sec:kt-freedman}

\begin{twierdzenie}[Freedman; trywialny wielomian Alexandera]
Niech $K\subset S^3$ będzie węzłem spełniającym
$\Delta_K(t)=1$ (w symetrycznej normalizacji).
Wówczas istnieje lokalnie płasko zanurzony dysk
$D^2\hookrightarrow B^4$ o brzegu $K$. Innymi słowy,
$g_4^{\mathrm{top}}(K)=0$.
\end{twierdzenie}

\begin{uwaga}
Dowód wykorzystuje topologiczną teorię zanurzania dysków,
uchwyty Cassona i czterowymiarową topologiczną chirurgię.
Nie otrzymujemy stąd gładko zanurzonego dysku. Wymiar
cztery jest szczególny właśnie dlatego, że topologicznie
dopuszczalne operacje na dyskach nie muszą mieć odpowiedników
gładkich.
\end{uwaga}

\begin{przyklad}[Podwojenie Whiteheada]
Dla dowolnego węzła $J\subset S^3$ jego nieskręcony
(zerowo ramowany) podwójny węzeł Whiteheada $\operatorname{Wh}(J)$
ma wielomian Alexandera $\Delta_{\operatorname{Wh}(J)}(t)=1$.
Z twierdzenia Freedmana jest więc topologicznie slice.
Dla dodatniego nieskręconego podwojenia prawoskrętnego
trójlistnika $T_{2,3}$ zachodzi jednak
$\tau(\operatorname{Wh}^{+}(T_{2,3}))=1$,
gdzie $\tau$ jest niezmiennikiem Floera opisanym w następnym
rozdziale. Z nierówności $|\tau(K)|\leq g_4^{\mathrm{sm}}(K)$
wynika, że ten węzeł nie jest gładko slice, choć jest
topologicznie slice.
\end{przyklad}

\begin{uwaga}[Znaczenie dla egzotycznych $\mathbb R^4$]
Rozbieżność między topologicznymi i gładkimi dyskami
w $B^4$ wiąże się z istnieniem niezgodnych struktur
gładkich na rozmaitościach homeomorficznych.
Szczegółowe konstrukcje egzotycznych $\mathbb R^4$
wymagają dodatkowo teorii uchwytów Cassona, twierdzeń
Donaldsona i innych narzędzi czterowymiarowych.
Sam fakt $g_4^{\mathrm{top}}\ne g_4^{\mathrm{sm}}$
jest ważnym, ale niepełnym opisem tej problematyki.
\end{uwaga}

\section{Węzły torusowe w \texorpdfstring{$B^4$}{B4} i problem Milnora}
\label{sec:kt-milnor}

\begin{twierdzenie}[Rozstrzygnięcie hipotezy Milnora]
Niech $p,q\geq2$ będą względnie pierwsze i niech
$T(p,q)\subset S^3$ będzie dodatnim węzłem torusowym.
Wtedy
$g_4^{\mathrm{sm}}(T(p,q))=(p-1)(q-1)/2$.
\end{twierdzenie}

\begin{uwaga}
Równość $g(T(p,q))=(p-1)(q-1)/2$ została udowodniona
wcześniej metodami klasycznymi. Znacznie trudniejsza
jest równość z rodzajem czterowymiarowym: a priori
powierzchnia w $B^4$ mogłaby mieć mniejszy rodzaj niż
powierzchnia w $S^3$. Pierwszy dowód dla całej rodziny
pochodzi z teorii cechowania Kronheimera--Mrowki.
Inne dowody wykorzystują teorię Seiberga--Wittena,
homologię Floera lub niezmiennik Rasmussena.
\end{uwaga}

\begin{przyklad}
Dla $T(2,3)$ mamy $g_4^{\mathrm{sm}}=1$;
dla $T(3,4)$ mamy $g_4^{\mathrm{sm}}=3$;
dla $T(3,5)$ mamy $g_4^{\mathrm{sm}}=4$.
Jest to rodzina, w której klasyczny rodzaj powierzchni
Seiferta i gładki rodzaj czterowymiarowy są sobie równe.
\end{przyklad}



% Iteracja 7: rozdzial autorski o plumbing, formach przeciecia i E8.
\chapter{Grafy plumbingowe, formy przecięcia i krata \texorpdfstring{$E_8$}{E8}}
\label{ch:tg-plumbing}

W poprzednich rozdziałach rozpatrywaliśmy dwie konstrukcje: z jednej
strony grafy, których wierzchołki kodują części rozmaitości, z drugiej
strony ramowane sploty opisujące dołączanie uchwytów indeksu dwa.
\emph{Plumbing} łączy te punkty widzenia. Z prostego grafu z wagami
całkowitymi budujemy gładką czterorozmaitość $P_\Gamma$, a następnie
odczytujemy z grafu jej formę przecięcia i homologię brzegu.
To jeden z najbardziej użytecznych sposobów przechodzenia od
kombinatoryki do topologii czterowymiarowej.

Zaczniemy od drzew o wierzchołkach reprezentujących sfery. W tym
przypadku konstrukcja jest elementarna i pozwala przeprowadzić
rachunki bez rozbudowanego aparatu rozmaitości Seiferta. Na końcu
omówimy ograniczenia metody oraz jej związek z teorią Donaldsona
i klasyfikacją Freedmana. Wszystkie rozmaitości plumbingowe
rozważamy jako \emph{zwarte rozmaitości gładkie z brzegiem};
nie są one automatycznie rozmaitościami zamkniętymi.

\section{Wiązki dyskowe i lokalna operacja plumbing}
\label{sec:tg-plumbing-definition}

\begin{definicja}[Zorientowana wiązka dysków nad sferą]
Niech $e\in\mathbb Z$. Przez $E(e)\to S^2$ oznaczamy zorientowaną
wiązkę $D^2$ o grupie strukturalnej $SO(2)$, której skojarzona
zorientowana wiązka wektorowa rzędu dwa ma klasę Eulera
$e\,u\in H^2(S^2;\mathbb Z)$, gdzie $u$ jest generatorem
wyznaczonym przez orientację sfery. \emph{Sekcją zerową}
nazywamy obraz odwzorowania $s\colon S^2\to E(e)$ przypisującego
punktowi bazowemu wektor zerowy jego włókna.
\end{definicja}

Każdą zorientowaną wiązkę płaszczyzn nad $S^2$ klasyfikuje
liczba całkowita: stopień funkcji przejścia
$S^1\to SO(2)\cong S^1$ między dwiema trywializacjami
na półkulach. Jest to właśnie liczba Eulera $e$.
Jeżeli przez $S$ oznaczymy obraz sekcji zerowej,
to jej wiązka normalna w $E(e)$ jest pierwotną wiązką
płaszczyzn, zatem $S\cdot S=e$. Równość tę należy
rozumieć jako liczbę przecięcia $S$ z niewielkim,
generycznym przesunięciem w kierunku normalnym.
Jej znak zależy od ustalonych orientacji.

\begin{definicja}[Dodatni plumbing dwóch wiązek]
Niech $E_1\to S^2$ i $E_2\to S^2$ będą zorientowanymi
wiązkami dyskowymi. W każdej wybieramy dysk $D_i\subset S^2$
nieprzecinający dysków wybranych dla innych operacji.
Nad $D_i$ ustalamy trywializację, która daje blok
$D^2_{\text{baza}}\times D^2_{\text{włókno}}$.
W lokalnym modelu plumbingowania identyfikujemy wybrane
bloki obu wiązek (po odpowiednim ograniczeniu ich
trywializacji), używając odwzorowania
\[
(z,w)\longmapsto(w,z), \qquad (z,w)\in D^2\times D^2,
\]
z dobranymi orientacjami, tak aby sekcje zerowe przecinały się
w powstałym bloku dodatnio i poprzecznie w jednym punkcie.
Po wygładzeniu narożników otrzymujemy \emph{dodatni plumbing}.
\end{definicja}

\begin{uwaga}[Lokalna geometria i znaki]
W bloku $D^2\times D^2$ sekcje zerowe dwóch wiązek
reprezentują odpowiednio dyski $D^2\times\{0\}$ i
$\{0\}\times D^2$. Mają jeden punkt przecięcia $(0,0)$.
Zamiana czynników zachowuje orientację czterowymiarową,
ponieważ permutuje dwa czynniki rzeczywiste wymiaru dwa.
Orientacje samych sfer i wiązek trzeba dobrać konsekwentnie:
wówczas znak lokalnego przecięcia wynosi $+1$. Zmiana
orientacji jednej sfery odwraca znaki wszystkich jej przecięć
z sąsiednimi sferami, ale nie zmienia jej samoprzecięcia.
\end{uwaga}

\begin{center}
\begin{tikzpicture}[scale=.98,line cap=round,>=Stealth]
  \begin{scope}[shift={(-2.65,0)}]
    \fill[tgBlue!15] (-1,-.65) rectangle (1,.65);
    \draw[tgBlue,very thick] (-1,-.65) rectangle (1,.65);
    \fill[tgCoral!22] (-.27,-.65) rectangle (.27,.65);
    \draw[tgCoral,line width=2.6pt] (0,-.65)--(0,.65);
    \draw[tgBlue,line width=2.4pt] (-1,0)--(1,0);
    \node[text=tgNavy,font=\footnotesize] at (0,-1.05) {$D^2_{\text{baza}}\times D^2_{\text{włókno}}$};
  \end{scope}
  \draw[tgTeal,-{Stealth[length=3mm]},line width=1.6pt] (-.75,.05)--(.75,.05);
  \node[text=tgTeal,font=\footnotesize] at (0,.42) {zamiana ról};
  \begin{scope}[shift={(2.65,0)}]
    \fill[tgCoral!15] (-1,-.65) rectangle (1,.65);
    \draw[tgCoral,very thick] (-1,-.65) rectangle (1,.65);
    \draw[tgBlue,line width=2.4pt] (0,-.65)--(0,.65);
    \draw[tgCoral,line width=2.6pt] (-1,0)--(1,0);
    \fill[tgPurple] (0,0) circle (.09);
    \node[text=tgNavy,font=\footnotesize] at (0,-1.05) {przecięcie dwóch sekcji};
  \end{scope}
  \node[text=tgNavy,font=\scriptsize] at (0,-1.52)
    {Model lokalny: jeden punkt przecięcia, bez utożsamiania całych sfer.};
\end{tikzpicture}
\end{center}

\begin{definicja}[Ważone drzewo plumbingowe]
Niech $\Gamma=(V,E)$ będzie skończonym, spójnym drzewem
prostym, a $e\colon V\to\mathbb Z$ funkcją wag.
Dla $v\in V$ bierzemy wiązkę $E(e_v)\to S^2$.
Dla każdej krawędzi $uv\in E$ wykonujemy dodatni plumbing
wiązek przypisanych do $u$ i $v$, używając rozłącznych
dysków bazowych przy różnych krawędziach. Po wygładzeniu
narożników otrzymujemy zwartą, zorientowaną, gładką
czterorozmaitość z brzegiem $P_\Gamma$.
\end{definicja}

W ogólnym grafie można także stosować wagi ujemne na
krawędziach i dopuszczać pętle; trzeba wtedy ustalić dodatkowe
dane orientacyjne i sposób sklejenia. W tej części, jeżeli nie
zaznaczono inaczej, graf oznacza \emph{ważone drzewo} z dodatnim
lokalnym znakiem przecinania, tak jak w powyższej definicji.

\section{Forma przecięcia z grafu}
\label{sec:tg-plumbing-Q}

\begin{twierdzenie}[Forma przecięcia drzewa plumbingowego]
\label{thm:tg-plumbing-intersection}
Niech $\Gamma=(V,E)$ będzie skończonym, spójnym, ważonym
drzewem plumbingowym, a $P_\Gamma$ czterorozmaitością
z poprzedniej definicji. Wtedy $\pi_1(P_\Gamma)=0$ oraz
$H_2(P_\Gamma;\mathbb Z)\cong\mathbb Z^{|V|}$.
W bazie $[S_v]$ wyznaczonej przez obrazy sekcji zerowych
macierz formy przecięcia wynosi
\[
Q_\Gamma=(q_{uv})_{u,v\in V},\qquad
q_{uv}=\begin{cases}
e_v,&u=v,\\
1,&uv\in E,\\
0,&u\ne v\text{ i }uv\notin E.
\end{cases}
\]
\end{twierdzenie}

\begin{proof}
Najpierw rozpatrujemy pojedynczy wierzchołek $v$.
Wiązka $E(e_v)$ retraktuje się na sekcję zerową $S^2$,
a więc jest po prostu spójna i ma drugą grupę homologii
izomorficzną z $\mathbb Z$. Dodajmy teraz do zbudowanego
już plumbingowania nową wiązkę $E(e_w)$ przyłączaną
przez krawędź drzewa. Wybierając kolejne liście drzewa,
możemy zapewnić, że nowa część styka się z poprzednią
w jednym bloku plumbingowym homeomorficznym z $D^4$.
W szczególności to przecięcie jest ściągalne.
Twierdzenie van Kampena daje więc nadal trywialną grupę
fundamentalną, natomiast ciąg Mayera--Vietorisa w stopniu
drugim ma postać
\[
0\longrightarrow H_2(P_{\rm stary};\mathbb Z)\oplus
H_2(E(e_w);\mathbb Z)\longrightarrow
H_2(P_{\rm nowy};\mathbb Z)\longrightarrow0.
\]
Indukcja względem liczby wierzchołków daje opis grupy
$H_2$ oraz bazę sekcji zerowych.

Dla $v=u$ liczbę $S_v\cdot S_v$ wyznacza klasa Eulera
wiązki normalnej, zatem wynosi ona $e_v$. Gdy $uv\in E$,
sekcje $S_u$ i $S_v$ przecinają się dokładnie w jednym
bloku plumbingowym i znak ich przecięcia, zgodnie z
przyjętą orientacją, jest dodatni. Gdy brak krawędzi,
sekcje można reprezentować rozłącznie. Otrzymujemy
wskazane elementy macierzy. Symetria macierzy wynika
z symetrii formy przecięcia w wymiarze cztery.
\end{proof}

\begin{przyklad}[Dwa wierzchołki]
Dla drzewa z jednym bokiem i wagami $-2,-3$ dostajemy
\[
Q_\Gamma=\begin{pmatrix}-2&1\\1&-3\end{pmatrix},
\qquad \det Q_\Gamma=5.
\]
Dwie sfery bazowe mają samoprzecięcia $-2$ i $-3$ oraz
przecinają się dodatnio w jednym punkcie. Jak zobaczymy,
ze względu na niezerowy wyznacznik brzeg tej czterorozmaitości
jest wymierną sferą homologiczną i ma pierwszą grupę homologii
rzędu pięć.
\end{przyklad}

\begin{center}
\begin{tikzpicture}[scale=1,>=Stealth]
  \node[tgBlueVertex,minimum size=12mm] (a) at (-1.5,0) {$-2$};
  \node[tgRedVertex,minimum size=12mm] (b) at (1.5,0) {$-3$};
  \draw[tgTree] (a)--(b);
  \node[font=\footnotesize,text=tgNavy] at (0,-.95)
    {Jedna krawędź: liczba przecięcia $[S_a]\cdot[S_b]=1$.};
\end{tikzpicture}
\end{center}

\begin{uwaga}[Konstrukcja przez uchwyty]
Ten sam $P_\Gamma$ można opisać za pomocą ramowanego splotu
w $S^3$. Każdemu wierzchołkowi odpowiada niezawiązany okrąg
z ramowaniem $e_v$, a każdej krawędzi -- dodatnie splecenie
tych dwóch okręgów równe jeden. Dla drzewa można wybrać
standardowy splot plumbingowy, w którym pozostałe liczby
splecenia wynoszą zero. Dołączenie do $B^4$ uchwytów indeksu
dwa wzdłuż tego splotu daje dyfeomorficzną czterorozmaitość.
Identyfikacja ta wynika z lokalnego modelu plumbing i
standardowego rachunku uchwytów; jej szczegółowa konstrukcja
za pomocą diagramów Kirby'ego została omówiona wcześniej.
W szczególności macierz splecenia tego splotu jest macierzą
$Q_\Gamma$.
\end{uwaga}

\section{Homologia i forma splecenia brzegu}
\label{sec:tg-plumbing-boundary}

\begin{twierdzenie}[Homologia brzegu plumbingowania]
\label{thm:tg-plumbing-boundary}
Niech $\Gamma$ będzie skończonym, spójnym, ważonym drzewem,
a $Y_\Gamma:=\partial P_\Gamma$ brzegiem odpowiadającej
mu czterorozmaitości. Wtedy
\[
H_1(Y_\Gamma;\mathbb Z)\cong
\operatorname{coker}(Q_\Gamma:\mathbb Z^{|V|}\to\mathbb Z^{|V|}),
\qquad
H_2(Y_\Gamma;\mathbb Z)\cong\ker Q_\Gamma.
\]
Jeżeli $\det Q_\Gamma\ne0$, to $Y_\Gamma$ jest wymierną
sferą homologiczną oraz
$|H_1(Y_\Gamma;\mathbb Z)|=|\det Q_\Gamma|$.
Jeżeli $|\det Q_\Gamma|=1$, to $Y_\Gamma$ jest całkowitą
sferą homologiczną.
\end{twierdzenie}

\begin{proof}
Oznaczmy $W=P_\Gamma$, $Y=\partial W$ i $r=|V|$.
W poprzednim twierdzeniu wykazaliśmy, że $H_1(W;\mathbb Z)=0$
i $H_2(W;\mathbb Z)\cong\mathbb Z^r$.
W długim ciągu dokładnym pary $(W,Y)$ występuje fragment
\[
H_3(W,Y)\longrightarrow H_2(Y)\longrightarrow H_2(W)
\xrightarrow{\ j\ }H_2(W,Y)\longrightarrow H_1(Y)
\longrightarrow H_1(W).
\]
Z dualności Poincarégo--Lefschetza i uniwersalnego twierdzenia
o współczynnikach wynika
$H_3(W,Y)\cong H^1(W)=0$ oraz
$H_2(W,Y)\cong H^2(W)\cong\operatorname{Hom}(H_2(W),\mathbb Z)$.
W wybranych bazach odwzorowanie $j$ zapisuje się jako macierz
formy przecięcia $Q_\Gamma$: obraz klasy $x$ jest funkcjonałem
$y\mapsto x\cdot y$. Ciąg dokładny upraszcza się więc do
\[
0\longrightarrow H_2(Y)\longrightarrow\mathbb Z^r
\xrightarrow{Q_\Gamma}\mathbb Z^r
\longrightarrow H_1(Y)\longrightarrow0,
\]
co dowodzi dwóch pierwszych izomorfizmów. Jeśli $Q_\Gamma$
ma niezerowy wyznacznik, jego jądro nad $\mathbb Z$ jest zerowe,
a iloraz $\mathbb Z^r/Q_\Gamma\mathbb Z^r$ jest skończony
o liczbie elementów $|\det Q_\Gamma|$ (np. z postaci normalnej
Smitha). Dualność Poincarégo na zamkniętej, spójnej,
zorientowanej trójrozmaitości daje wtedy pozostałe grupy
homologii takie jak dla sfery nad $\mathbb Q$.
Przy wyznaczniku $\pm1$ również $H_1(Y)=0$, co daje
analogiczne stwierdzenie nad $\mathbb Z$.
\end{proof}

\begin{przyklad}[Macierz i relacje w pierwszej homologii]
Dla drzewa o wagach $-2,-3$ macierz z poprzedniego
przykładu daje prezentację
$H_1(Y;\mathbb Z)=\langle x,y\mid -2x+y=0,\ x-3y=0\rangle$.
Pierwsza relacja daje $y=2x$, a druga $-5x=0$.
Zatem $H_1(Y;\mathbb Z)\cong\mathbb Z/5\mathbb Z$.
Zauważmy, że wyznacznik macierzy wystarcza do obliczenia
\emph{rzędu} tej grupy, ale nie zawsze do wyznaczenia
jej struktury: przykładem będzie graf typu $D_4$.
\end{przyklad}

\begin{stwierdzenie}[Forma splecenia brzegu]
Niech $\Gamma$ będzie jak wyżej i niech $Q_\Gamma$ będzie
nieosobliwa nad $\mathbb Q$. Dla $Y=\partial P_\Gamma$
forma splecenia na skończonej grupie
$H_1(Y;\mathbb Z)\cong\mathbb Z^r/Q_\Gamma\mathbb Z^r$
ma, przy konwencji orientacyjnej brzegu jako $\partial P_\Gamma$,
postać
\[
\lambda_Y([x],[y])\equiv-\,x^TQ_\Gamma^{-1}y
\pmod{\mathbb Z}, \qquad x,y\in\mathbb Z^r.
\]
\end{stwierdzenie}

\begin{proof}
Macierz $Q_\Gamma$ jest całkowita, symetryczna i nieosobliwa.
Jeżeli zastąpimy $x$ przez $x+Q_\Gamma z$, wyrażenie po prawej
zmienia się o $-z^Ty\in\mathbb Z$; analogicznie dla $y$.
Forma jest zatem dobrze określona modulo $\mathbb Z$.
Jej niezdegenerowanie wynika z tego, że jeśli
$x^TQ_\Gamma^{-1}y\in\mathbb Z$ dla każdego $y\in\mathbb Z^r$,
to $Q_\Gamma^{-1}x\in\mathbb Z^r$, a więc
$[x]=0$ w grupie ilorazowej. Identyfikację tej formy
z geometryczną liczbą splecenia otrzymujemy następująco.
Klasę skrętnej krzywej na brzegu reprezentuje cykl, którego
wielokrotność ogranicza łańcuch dwuwymiarowy.
Po umieszczeniu takiego łańcucha w $W$ jego klasy
względne mierzy odwrotność macierzy $j=Q_\Gamma$ z
dowodu poprzedniego twierdzenia. Liczba przecięcia tego
łańcucha z drugim cyklem, podzielona przez odpowiednią
wielokrotność, daje właśnie współczynnik $Q_\Gamma^{-1}$;
znak minus pochodzi z konwencji orientacji brzegu.
Jest to szczególny przypadek rachunku formy splecenia
przez formę przecięcia czterorozmaitości.
\end{proof}

\begin{uwaga}
Przy odwróceniu orientacji $Y$ forma splecenia zmienia znak.
Nie należy więc porównywać formuł $Q^{-1}$ i $-Q^{-1}$
bez sprawdzenia orientacji i definicji dodatniego meridianu.
Sama grupa $H_1(Y;\mathbb Z)$ na ogół nie rozróżnia
różnych trójrozmaitości.
\end{uwaga}

\section{Łańcuchy ważone, ułamki łańcuchowe i przestrzenie soczewkowe}
\label{sec:tg-plumbing-chains}

Najłatwiejsze rachunki wykonujemy dla grafów liniowych.
Złożenie kolejnych bloków można wtedy opisać zarówno
macierzą trójdiagonalną, jak i ułamkiem łańcuchowym.

\begin{definicja}[Łańcuch plumbingowy]
Niech $r\ge1$ i niech $a_1,\dots,a_r\in\mathbb Z$.
Przez $C(a_1,\dots,a_r)$ oznaczamy ważone drzewo o
wierzchołkach $v_1,\dots,v_r$, krawędziach
$v_iv_{i+1}$ dla $1\le i<r$ i wagach $e(v_i)=a_i$.
Jego macierz przecięcia oznaczamy przez $Q(a_1,\dots,a_r)$.
\end{definicja}

\begin{lemat}[Rekurencja wyznaczników]
Dla $r\ge2$ wyznaczniki
$d_r:=\det Q(a_1,\dots,a_r)$ spełniają
\[
d_r=a_r d_{r-1}-d_{r-2},\qquad
 d_0=1,\quad d_1=a_1.
\]
\end{lemat}
\begin{proof}
Ostatni wiersz macierzy $Q(a_1,\dots,a_r)$ ma
jedynie dwa potencjalnie niezerowe wyrazy: $1$ na pozycji
$r-1$ oraz $a_r$ na pozycji $r$. Rozwijamy wyznacznik
Laplace'a wzdłuż ostatniego wiersza. Składnik z $a_r$
równa się $a_r d_{r-1}$. Minor odpowiadający wyrazowi $1$
na pozycji $r-1$ ma w ostatniej kolumnie tylko jedną
jedynkę; po drugim rozwinięciu dostajemy $-d_{r-2}$.
Dla $r=2$ otrzymujemy bezpośrednio
$d_2=a_1a_2-1$, co zgadza się z rekurencją.
\end{proof}

\begin{stwierdzenie}[Ułamek łańcuchowy dla chirurgii]
Niech $a_1,\dots,a_r\in\mathbb Z$ będą liczbami,
dla których wszystkie mianowniki w poniższym wyrażeniu
są niezerowe. Niech $L$ będzie standardowym ramowanym
splotem będącym łańcuchem niezawiązanych okręgów o
ramowaniach $a_1,\dots,a_r$. Operacje \emph{slam-dunk}
(racjonalnej eliminacji kolejnego torusa pełnego)
prowadzą do wymiernej chirurgii na jednym węźle trywialnym
o współczynniku
\[
[a_1,\dots,a_r]^-:=a_1-\cfrac1{a_2-
\cfrac1{\ddots-\cfrac1{a_r}}}.
\]
W szczególności brzeg odpowiadającego plumbingowania
jest przestrzenią soczewkową, gdy licznik tego ułamka
jest niezerowy (z uwzględnieniem orientacji i konwencji
parametrów przestrzeni soczewkowych).
\end{stwierdzenie}

\begin{proof}[Rachunek współczynnika]
Ostatni element łańcucha jest meridianem poprzedniego
węzła i ma ramowanie $a_r$. Z lokalnego rachunku chirurgii
Dehna wynika, że po wypełnieniu tego elementu
współczynnik chirurgii na poprzednim zmienia się z
$a_{r-1}$ na $a_{r-1}-1/a_r$. Powtarzając operację,
otrzymujemy wskazany ułamek łańcuchowy. Operacja lokalna
polega na zmianie bazy meridian--długość na torusie
brzegowym i jest równoważna sklejaniu dwóch torusów
pełnych; ma więc również interpretację geometryczną.
Szczegółowa klasyfikacja przestrzeni soczewkowych
zależy od konwencji $L(p,q)$ i nie wynika z samej
pierwszej homologii. Ostatnie wypełnienie jest chirurgią
na węźle trywialnym, a zatem sklejeniem dwóch torusów
pełnych, czyli przestrzenią soczewkową.
\end{proof}

\begin{uwaga}[Orientacja współczynnika]
Dla łańcucha ujemnych wag $-b_i$, gdzie $b_i\ge2$,
wartość $[-b_1,\dots,-b_r]^-$ jest ujemna.
Przy przejściu od tej liczby do oznaczenia $L(p,q)$
trzeba ustalić znak licznika oraz orientację brzegu;
różne podręczniki stosują konwencje $L(p,q)$ i
$L(p,q^{-1})$ lub ich odbicia lustrzane.
Poniższe obliczenia homologii nie zależą od tych wyborów.
\end{uwaga}

\begin{twierdzenie}[Łańcuch typu $A_n$]
\label{thm:tg-plumbing-An}
Niech $n\ge1$, a $\Gamma=A_n$ będzie łańcuchem
$n$ wierzchołków o wagach $-2$. Niech $P_{A_n}$
będzie jego czterorozmaitością plumbingową. Wtedy
\[
\det Q_{A_n}=(-1)^n(n+1),\qquad
H_1(\partial P_{A_n};\mathbb Z)\cong\mathbb Z/(n+1)\mathbb Z.
\]
Forma przecięcia $Q_{A_n}$ jest ujemnie określona.
Jej brzeg jest przestrzenią soczewkową o pierwszej
grupie homologii rzędu $n+1$.
\end{twierdzenie}
\begin{proof}
Niech $d_n$ będzie wyznacznikiem odpowiedniej macierzy.
Z lematu o wyznaczniku łańcucha mamy
$d_n=-2d_{n-1}-d_{n-2}$, gdzie $d_0=1$, $d_1=-2$.
Podstawienie $d_n=(-1)^n(n+1)$ spełnia rekurencję
i warunki początkowe, zatem wzór jest prawdziwy dla
wszystkich $n$.

Aby wyznaczyć grupę, oznaczmy jej generatory odpowiadające
kolejnym wierzchołkom przez $x_1,\dots,x_n$. Dla $n=1$ grupa jest bezpośrednio
$\mathbb Z/2\mathbb Z$, więc dalej przyjmijmy $n\ge2$.
Relacje
macierzowe mają postać
$-2x_1+x_2=0$ na początku,
$x_{i-1}-2x_i+x_{i+1}=0$ dla $1<i<n$ oraz
$x_{n-1}-2x_n=0$ na końcu.
Przyjmując formalnie $x_0=x_{n+1}=0$, zapisujemy je
jednolicie jako $x_{i-1}-2x_i+x_{i+1}=0$.
Indukcyjnie otrzymujemy $x_i=i x_1$, a ostatnia relacja
daje $(n+1)x_1=0$. Grupa jest więc cykliczna i ma
rząd $n+1$ na mocy poprzedniego twierdzenia.

Wszystkie wiodące minory główne macierzy $Q_{A_n}$
mają wyznaczniki $(-1)^k(k+1)$. Z kryterium Sylvestera
wynika, że $Q_{A_n}$ jest ujemnie określona.
Wreszcie współczynnik chirurgii wyrażony przez ułamek
łańcuchowy wynosi
$[-2,\dots,-2]^-=-(n+1)/n$ dla $n\ge1$.
Dlatego brzeg otrzymuje się przez wymierną chirurgię
na węźle trywialnym, co kończy dowód.
\end{proof}

\begin{center}
\begin{tikzpicture}[scale=.95,line cap=round]
  \foreach \x in {0,1,2,3,4}{
    \node[tgBlueVertex,minimum size=9mm] (v\x) at ({1.45*\x},0) {$-2$};
  }
  \foreach \x in {0,1,2,3}{
    \pgfmathtruncatemacro{\y}{\x+1}
    \draw[tgTree] (v\x)--(v\y);
  }
  \node[text=tgNavy,font=\small] at (2.9,-.85) {Przykład: $A_5$, wyznacznik $-6$};
\end{tikzpicture}
\end{center}

\section{Rozgałęzienia: typ \texorpdfstring{$D_n$}{Dn} i różne grupy brzegowe}
\label{sec:tg-plumbing-Dn}

\begin{przyklad}[Graf $D_4$: wyznacznik nie wystarcza]
Rozważmy gwiazdę z czterema wierzchołkami, w której
wierzchołek centralny łączy się z trzema liśćmi,
a wszystkie wagi wynoszą $-2$. W kolejności
(centrum, trzy liście) forma przecięcia ma macierz
\[
Q_{D_4}=\begin{pmatrix}
-2&1&1&1\\1&-2&0&0\\1&0&-2&0\\1&0&0&-2
\end{pmatrix}.
\]
Blok odpowiadający trzem liściom jest równy $-2I_3$.
Jego dopełnienie Schura to
$-2-(1,1,1)(-\tfrac12I_3)(1,1,1)^T=-\tfrac12$.
Dlatego $\det Q_{D_4}=(-2)^3(-\tfrac12)=4$.

Rachunek modulo dwa pokazuje jednak, że
$H_1(\partial P_{D_4};\mathbb Z)$ nie jest grupą
$\mathbb Z/4\mathbb Z$. Rzeczywiście, macierz
$Q_{D_4}$ po redukcji modulo dwa ma rząd $2$, więc
\[
\dim_{\mathbb F_2}
(H_1(\partial P_{D_4};\mathbb Z)\otimes\mathbb F_2)=2.
\]
Grupa rzędu cztery z dwoma niezależnymi elementami
modulo dwa musi być izomorficzna z
$\mathbb Z/2\mathbb Z\oplus\mathbb Z/2\mathbb Z$.
Pokazuje to, że rząd $|\det Q|$ nie określa całej
grupy; potrzebna jest np. postać normalna Smitha.
\end{przyklad}

\begin{center}
\begin{tikzpicture}[scale=.9]
  \node[draw=tgPurple,fill=tgPurple!14,circle,
    minimum size=10mm] (c) at (0,0) {$-2$};
  \node[tgBlueVertex,minimum size=10mm] (a) at (-1.6,0) {$-2$};
  \node[tgRedVertex,minimum size=10mm] (b) at (1.6,0) {$-2$};
  \node[tgTealVertex,minimum size=10mm] (d) at (0,1.35) {$-2$};
  \draw[tgEdge] (a)--(c)--(b) (c)--(d);
  \node[font=\footnotesize,text=tgNavy] at (0,-.9)
    {$D_4$: centralny wierzchołek o stopniu trzy};
\end{tikzpicture}
\end{center}

\begin{uwaga}[Cała rodzina $D_n$]
Dla standardowego grafu Dynkina $D_n$, $n\ge4$, z wagami
$-2$ i dodatnimi przecięciami na krawędziach, zachodzi
$|\det Q_{D_n}|=4$. Dokładniejszy rachunek postaci normalnej
Smitha daje grupę brzegu $\mathbb Z/4\mathbb Z$ dla $n$
nieparzystego oraz
$\mathbb Z/2\mathbb Z\oplus\mathbb Z/2\mathbb Z$ dla $n$
parzystego. Oto szczegóły. Usunięcie centralnego wierzchołka rozdziela
graf na ramiona długości $n-3,1,1$. Z wyznaczników
łańcuchów $A_m$ oraz dopełnienia Schura otrzymujemy
$|\det Q_{D_n}|=4$. Aby rozpoznać strukturę grupy,
przechodzimy do prezentacji jej generatorów: wzdłuż długiego
ramienia relacje $-2x_i+x_{i-1}+x_{i+1}=0$ wyrażają
kolejne $x_i$ przez generator $x$ przy wolnym końcu.
Jeśli $u,v$ oznaczają dwa krótkie liście, relacje przy
nich dają $2u=2v=(n-2)x$, zaś relacja centralna daje
$u+v=(n-1)x$. Eliminacja $v$ prowadzi zatem do
\[
H_1(\partial P_{D_n};\mathbb Z)
\cong \langle x,u\mid 2u=(n-2)x,\ 2x=0\rangle.
\]
Dla $n$ parzystego pierwsza relacja sprowadza się do
$2u=0$, więc otrzymujemy $(\mathbb Z/2)^2$.
Dla $n$ nieparzystego jest równoważna $2u=x$,
a wtedy $4u=0$ i otrzymujemy $\mathbb Z/4$.
Rachunek uzasadnia podaną wcześniej klasyfikację.
\end{uwaga}

\section{Graf \texorpdfstring{$E_8$}{E8} i jego macierz przecięcia}
\label{sec:tg-plumbing-E8}

Najważniejszym przykładem jest rozgałęzione drzewo z ośmioma
wierzchołkami, występujące również jako diagram Dynkina
układu pierwiastków $E_8$. Wagi wszystkich wierzchołków
przyjmujemy równe $-2$.

\begin{definicja}[Ujemna forma $E_8$]
Niech $\Gamma_{E_8}$ będzie drzewem o wierzchołkach
$v_1,\ldots,v_8$, którego krawędzie to
$v_iv_{i+1}$ dla $1\le i\le6$ oraz $v_3v_8$.
Niech $Q_{E_8}$ ma elementy diagonalne równe $-2$,
pozadiagonalne równe $1$ dla sąsiednich wierzchołków
i zero w pozostałych miejscach. \emph{Ujemną formą
$E_8$} nazywamy formę całkowitą o macierzy $Q_{E_8}$.
Plumbing odpowiadający temu grafowi oznaczamy
przez $P_{E_8}$.
\end{definicja}

\begin{center}
\begin{tikzpicture}[scale=1.0,line cap=round]
  \foreach \i in {1,...,7}{
    \node[tgBlueVertex,minimum size=10mm] (v\i) at ({1.15*(\i-1)},0) {$-2$};
  }
  \foreach \i in {1,...,6}{
    \pgfmathtruncatemacro{\j}{\i+1}
    \draw[tgTree] (v\i)--(v\j);
  }
  \node[tgRedVertex,minimum size=10mm] (v8) at (2.30,1.40) {$-2$};
  \draw[tgAccent] (v3)--(v8);
  \node[font=\footnotesize,text=tgNavy] at (3.45,-.85)
    {Drzewo $E_8$: długości ramion $2$, $4$ i $1$.};
\end{tikzpicture}
\end{center}

\begin{twierdzenie}[Rachunek macierzy $E_8$]
\label{thm:tg-E8-unimodular}
Forma $Q_{E_8}$ jest parzysta, unimodularna i ujemnie
określona. W szczególności
\[
\det Q_{E_8}=1,\qquad \operatorname{sign}(Q_{E_8})=-8,
\qquad H_1(\partial P_{E_8};\mathbb Z)=0.
\]
\end{twierdzenie}

\begin{proof}
\emph{Krok pierwszy: parzystość.}
Dla $x=(x_v)_{v\in V}\in\mathbb Z^8$ mamy
\[
x^TQ_{E_8}x=-2\sum_vx_v^2+2\sum_{uv\in E}x_ux_v.
\]
Liczba ta jest parzysta, ponieważ oba składniki są
podzielne przez dwa. Jest to definicja parzystości
całkowitej formy biliniowej symetrycznej.

\emph{Krok drugi: wyznacznik i określoność.}
Usuńmy z drzewa wierzchołek trójwalentny $v_3$.
Powstają trzy rozłączne łańcuchy o długościach odpowiednio
$2$, $4$ i $1$, z macierzami $A_2,A_4,A_1$ z poprzedniej
sekcji. Każda z tych macierzy jest ujemnie określona.
Ich wyznaczniki wynoszą odpowiednio $3$, $5$ i $-2$.
Wszystkie trzy bloki razem mają więc wyznacznik $-30$.

Obliczmy dopełnienie Schura odpowiadające środkowemu
wierzchołkowi. Dla ujemnego łańcucha $A_m$ z
wpisami $-2$ na przekątnej i $1$ na sąsiednich
pozycjach wartość $(A_m^{-1})_{11}$ przy jego końcu
wynosi $-m/(m+1)$. Aby to wykazać, wystarczy użyć
wzoru na element macierzy odwrotnej jako iloraz
minorów: po usunięciu pierwszego wiersza i kolumny
pozostaje $A_{m-1}$, zatem
$(A_m^{-1})_{11}=\det A_{m-1}/\det A_m=-m/(m+1)$.
Wynik pozostaje taki sam, gdy ramie zapisujemy w
odwrotnej kolejności.

Dopełnienie Schura macierzy z trzema ramionami wynosi
\[
S=-2-\left(-\frac23-\frac45-\frac12\right)
=-2+\frac23+\frac45+\frac12=-\frac1{30}.
\]
Stąd ze wzoru wyznacznika blokowego
$\det Q_{E_8}=(-30)(-1/30)=1$. Ponieważ wszystkie
bloki ramion są ujemnie określone i ich dopełnienie
Schura także jest ujemne, kryterium określoności
dla macierzy blokowych daje ujemną określoność
całej macierzy $Q_{E_8}$. Ma ona zatem osiem
ujemnych wartości własnych i żadnej dodatniej,
więc jej sygnatura jest równa $-8$.

\emph{Krok trzeci: brzeg.}
Wyznacznik jest równy $1$, więc macierz $Q_{E_8}$
należy do $GL_8(\mathbb Z)$. Twierdzenie
\ref{thm:tg-plumbing-boundary} daje
$H_1(\partial P_{E_8};\mathbb Z)=0$.
Razem z dualnością Poincarégo w wymiarze trzy
pokazuje to, że brzeg jest całkowitą sferą homologiczną.
\end{proof}

\begin{uwaga}[Dlaczego to wyjątkowy przykład]
Wśród dodatnio określonych parzystych krat unimodularnych
najmniejsza możliwa dodatnia ranga wynosi osiem; jej
pierwszym przykładem jest dodatnia krata $E_8$.
Tutaj otrzymaliśmy jej wersję ujemnie określoną.
Parzystość oznacza, że wszystkie samoprzecięcia klas
całkowitych są parzyste, natomiast unimodularność
wiąże formę z homologiczną trywialnością brzegu.
Nie jest jednak prawdą, że każda całkowita sfera
homologiczna jest dyfeomorficzna z $S^3$.
\end{uwaga}

\begin{stwierdzenie}[Brzeg $E_8$ i sfera homologiczna Poincarégo]
Brzeg ujemnie określonego plumbingowania $P_{E_8}$
jest, z dokładnością do konwencji orientacyjnej,
\emph{sferą homologiczną Poincarégo}
$\Sigma(2,3,5)$. Jej grupa fundamentalna jest
nieprzemienna i nietrywialna; w szczególności
$\partial P_{E_8}$ nie jest homeomorficzny z $S^3$.
\end{stwierdzenie}

\begin{proof}[Identyfikacja przez rozwarstwienie Seiferta]
Graf $E_8$ jest gwiazdą z centralną wagą $-2$ i trzema
ramionami złożonymi odpowiednio z $1$, $2$ i $4$
wierzchołków o wadze $-2$. Ułamki łańcuchowe tych ramion
wynoszą $-2/1$, $-3/2$ i $-5/4$, co sprawdza się
indukcyjnie ze wzoru $[-2,\ldots,-2]^-=-(m+1)/m$.
Twierdzenie~\ref{thm:tg-star-seifert}, udowodnione w
następnym rozdziale przez redukcję ramion do wypełnień Dehna,
identyfikuje brzeg z rozmaitością o znormalizowanych danych
Seiferta $M(0;-2;(2,1),(3,2),(5,4))$.
Jej identyfikacja ze sferą Brieskorna $\Sigma(2,3,5)$
wykorzystuje klasyfikację rozwarstwień Seiferta; ten ostatni,
nieelementarny krok przyjmujemy z literatury.

Dla niebanalności grupy fundamentalnej nie trzeba
znać jej pełnej prezentacji. Prezentacja grupy rozmaitości
Seiferta posiada centralny generator $h$ regularnego
włókna i generatory $x_1,x_2,x_3$ wyjątkowych włókien.
Po dodaniu relacji $h=1$ otrzymujemy grupę trójkątną
\[
\Delta(2,3,5)=\langle x_1,x_2,x_3\mid
x_1^2=x_2^3=x_3^5=x_1x_2x_3=1\rangle.
\]
Jest ona izomorficzna z grupą obrotów dwudziestościanu,
czyli z grupą $A_5$ rzędu $60$. Istnieje więc
surjekcja $\pi_1(\partial P_{E_8})\twoheadrightarrow A_5$.
Grupa fundamentalna brzegu jest zatem nietrywialna.
Z drugiej strony, jak wykazaliśmy, jej abelianizacja
jest zerowa. Otrzymujemy przykład przestrzeni będącej
sferą homologiczną, ale nie homotopijną.
\end{proof}

\begin{uwaga}[Pełna grupa fundamentalna]
W istocie $\pi_1(\Sigma(2,3,5))$ jest binarną grupą
dwudziestościenną rzędu $120$, czyli pewnym centralnym
rozszerzeniem $A_5$ przez $\mathbb Z/2\mathbb Z$.
Pełne obliczenie tej grupy nie jest potrzebne do
wykazania, że brzeg jest różny od $S^3$.
\end{uwaga}

\section{Rachunek Kirby'ego, przesuwanie uchwytów i wydmuchiwanie}
\label{sec:tg-plumbing-kirby}

\begin{stwierdzenie}[Przesunięcie uchwytu a kongruencja macierzy]
Niech $W$ będzie czterorozmaitością zbudowaną z $B^4$
przez dołączenie $r$ uchwytów indeksu dwa, o macierzy
przecięcia $Q$ w bazie klas powierzchni domkniętych.
Przesunięcie $i$-tego uchwytu po $j$-tym odpowiada
zastąpieniu $e_i$ przez $e_i\pm e_j$. Macierz formy
w nowej bazie ma postać $P^TQP$, gdzie
$P\in GL_r(\mathbb Z)$ realizuje taką zmianę bazy.
W szczególności wyznacznik oraz sygnatura formy
pozostają niezmienione.
\end{stwierdzenie}
\begin{proof}
Forma przecięcia jest formą dwuliniową na
$H_2(W;\mathbb Z)$. Jeśli kolumny macierzy $P$
reprezentują nowe wektory bazowe w starej bazie,
to ich wzajemne iloczyny skalarne tworzą macierz
$P^TQP$. W rozkładzie na uchwyty przesunięcie
przeprowadza na poziomie klas homologii dokładnie dodawanie lub
odejmowanie jednej klasy rdzenia od drugiej.
Macierz $P$ ma wyznacznik $\pm1$, a zatem
$\det(P^TQP)=(\det P)^2\det Q=\det Q$.
Z prawa bezwładności Sylvestera wynika niezmienniczość
liczby dodatnich i ujemnych wartości własnych.
\end{proof}

\begin{przyklad}[Jawny rachunek przesunięcia]
Niech $Q=\begin{psmallmatrix}a&\ell\\\ell&b\end{psmallmatrix}$.
Po zastąpieniu pierwszego wektora $e_1$ przez
$e_1+e_2$ otrzymujemy
\[
Q'=\begin{pmatrix}a+2\ell+b&\ell+b\\\ell+b&b\end{pmatrix},
\qquad \det Q'=ab-\ell^2=\det Q.
\]
Rachunek pokazuje dokładnie, skąd bierze się wyraz
$2\ell$ w nowym ramowaniu: pochodzi z obu iloczynów
mieszanych formy symetrycznej.
\end{przyklad}

\begin{stwierdzenie}[Algebraiczne usunięcie liścia o wadze $-1$]
Niech $\Gamma$ będzie ważonym drzewem i niech $v$
będzie liściem o wadze $-1$, połączonym tylko z
wierzchołkiem $u$ o wadze $b$. Wtedy macierz
$Q_\Gamma$ jest kongruentna nad $\mathbb Z$
z sumą ortogonalną $[-1]\oplus Q_{\Gamma'}$, gdzie
$\Gamma'$ otrzymujemy przez usunięcie liścia $v$
i zwiększenie wagi $u$ z $b$ do $b+1$.
W szczególności $\det Q_\Gamma=-\det Q_{\Gamma'}$.
\end{stwierdzenie}
\begin{proof}
Niech $e_v,e_u$ oznaczają odpowiednie wektory bazowe.
Dane przecięcia to $e_v^2=-1$, $e_v\cdot e_u=1$
i $e_v\cdot e_w=0$ dla $w\notin\{u,v\}$.
Wprowadźmy nową bazę: zostawiamy $e_v$, a $e_u$
zastępujemy przez $f_u=e_u+e_v$. Wtedy
$e_v\cdot f_u=1-1=0$ oraz
$f_u^2=b+2-1=b+1$. Dla pozostałych wierzchołków
$w$ mamy $f_u\cdot e_w=e_u\cdot e_w$, ponieważ
$e_v\cdot e_w=0$. Otrzymaliśmy więc dokładnie
wymaganą sumę ortogonalną. Zmiana bazy jest
unimodularna, a równość wyznaczników wynika ze wzoru
na wyznacznik sumy blokowej.
\end{proof}

\begin{uwaga}[Geometria operacji blow-down]
Jeżeli klasę $e_v$ reprezentuje gładko zanurzona sfera
o samoprzecięciu $-1$, to jej otoczenie tubularne jest
ujemną wiązką dyskową $E(-1)$, której brzeg jest
$S^3$. Można zastąpić to otoczenie przez kulę
$B^4$; nazywa się to \emph{wydmuchiwaniem w dół}
(\emph{blow-down}). Dla liścia plumbingowego operacja
prowadzi do grafu $\Gamma'$ ze zmienioną wagą.
Przeciwnie, \emph{wydmuchanie} (\emph{blow-up}) dodaje
lokalną sferę wyjątkową o samoprzecięciu $-1$.
Nie jest to dowolne skreślanie wiersza macierzy:
konieczna jest geometryczna obecność odpowiedniej sfery.
\end{uwaga}

\begin{center}
\begin{tikzpicture}[scale=.95,>=Stealth]
  \node[tgRedVertex,minimum size=11mm] (a) at (-2.5,0) {$-1$};
  \node[tgBlueVertex,minimum size=11mm] (b) at (-.8,0) {$b$};
  \draw[tgAccent] (a)--(b);
  \draw[tgTeal,-{Stealth[length=3.3mm]},line width=1.6pt] (.25,0)--(1.40,0);
  \node[tgTealVertex,minimum size=11mm] (c) at (2.65,0) {$b+1$};
  \node[text=tgNavy,font=\footnotesize] at (.05,-1.00)
    {Algebraiczne usunięcie liścia $-1$; w geometrii -- blow-down.};
\end{tikzpicture}
\end{center}

\section{Plumbing \texorpdfstring{$E_8$}{E8}, twierdzenia Freedmana i Donaldsona}
\label{sec:tg-plumbing-freedman}

\begin{twierdzenie}[Donaldson; diagonalizacja form określonych]
Niech $X$ będzie gładką, zamkniętą, spójną i po prostu
spójną zorientowaną czterorozmaitością. Jeśli całkowita
forma przecięcia $Q_X$ jest dodatnio lub ujemnie określona,
to jest izomorficzna odpowiednio z formą diagonalną
$\langle1\rangle^{\oplus b_2(X)}$ albo
$\langle-1\rangle^{\oplus b_2(X)}$.
\end{twierdzenie}

Wynik Simona Donaldsona jest głębokim twierdzeniem
teorii połączeń antysamodualnych. Przyjmujemy go bez dowodu;
w naszej konstrukcji potrzebujemy jedynie jego konsekwencji:
\emph{żadna gładka zamknięta po prostu spójna czterorozmaitość
nie ma określonej formy przecięcia $\pm E_8$}.
Rzeczywiście, forma $\pm E_8$ jest parzysta, natomiast każda
niezerowa forma diagonalna $\pm I_8$ ma wektor bazowy
z nieparzystym samoprzecięciem $\pm1$.

\begin{twierdzenie}[Niedopuszczalność gładkiego wypełnienia kontraktywnego]
\label{thm:tg-E8-no-smooth-fill}
Niech $Y=\partial P_{E_8}$ będzie zorientowanym brzegiem
ujemnego plumbingowania $E_8$. Nie istnieje zwarta,
gładka, kontraktywna czterorozmaitość $C$ taka, że
$\partial C$ jest orientacyjnie przeciwna do $Y$.
W szczególności sfera homologiczna Poincarégo nie może
być brzegiem gładkiej rozmaitości kontraktywnej.
\end{twierdzenie}
\begin{proof}
Przypuśćmy, że taka $C$ istnieje. Sklejmy ją z
$P_{E_8}$ wzdłuż przeciwstawnie zorientowanych brzegów:
$X=P_{E_8}\cup_Y C$. Ponieważ obie części są gładkie,
po użyciu kołnierza i wygładzeniu sklejenia otrzymujemy
gładką, zamkniętą, orientowaną czterorozmaitość.
Twierdzenie van Kampena, wykorzystujące
$\pi_1(P_{E_8})=0$ i $\pi_1(C)=0$, daje
$\pi_1(X)=0$. Z ciągu Mayera--Vietorisa,
$H_1(Y)=H_2(Y)=0$ i kontraktywności $C$ wynika
$H_2(X;\mathbb Z)\cong H_2(P_{E_8};\mathbb Z)$.
Ponadto przecięcia klas reprezentowanych przez sfery
we wnętrzu $P_{E_8}$ nie ulegają zmianie przy sklejeniu.
Stąd forma przecięcia zamkniętej $X$ jest dokładnie
$Q_{E_8}$: parzysta, unimodularna i ujemnie określona.
Przeczy to przytoczonemu twierdzeniu Donaldsona,
gdyż $Q_{E_8}$ nie jest diagonalizowalna nad $\mathbb Z$.
\end{proof}

\begin{twierdzenie}[Freedman; topologiczne wypełnianie sfer homologicznych]
Każda zamknięta, zorientowana całkowita sfera homologiczna
trójwymiarowa $Y$ jest brzegiem zwartej,
kontraktywnej \emph{topologicznej} czterorozmaitości $C$.
Nie zakłada się, że $C$ ma strukturę gładką.
\end{twierdzenie}

Przyjmujemy to twierdzenie z literatury. Jego dowód
korzysta z czterowymiarowej topologicznej teorii chirurgii
i twierdzenia o zanurzaniu dysków. Właśnie różnica między
słowami \emph{topologiczna} i \emph{gładka} jest tu istotą
zagadnienia, a nie jedynie techniczną uwagą.

\begin{wniosek}[Zamknięta topologiczna rozmaitość $E_8$]
Istnieje zamknięta, spójna, zorientowana, po prostu
spójna topologiczna czterorozmaitość $X_{E_8}$,
której całkowita forma przecięcia jest równa
ujemnej formie $E_8$. Nie dopuszcza ona struktury gładkiej.
\end{wniosek}
\begin{proof}
Niech $Y=\partial P_{E_8}$. Z udowodnionego rachunku
homologii wynika, że $Y$ jest całkowitą sferą homologiczną.
Z twierdzenia Freedmana istnieje kontraktywna topologiczna
czterorozmaitość $C$ o brzegu $-Y$. Sklejamy
$X_{E_8}=P_{E_8}\cup_Y C$. Argument z twierdzenia
\ref{thm:tg-E8-no-smooth-fill}, oparty na van Kampenie
i Mayerze--Vietorisie, nie używa gładkości przy
obliczaniu grupy fundamentalnej ani formy przecięcia.
Stąd $X_{E_8}$ jest po prostu spójna, zamknięta
i ma formę $Q_{E_8}$. Gdyby dopuszczała strukturę gładką,
twierdzenie Donaldsona wymuszałoby diagonalizację
jej ujemnie określonej formy przecięcia. Jest to
niemożliwe ze względu na parzystość $Q_{E_8}$.
\end{proof}

\begin{uwaga}[Nie należy utożsamiać trzech rozmaitości]
Występują tu trzy różne obiekty:
\begin{enumerate}[label=(\roman*)]
\item $P_{E_8}$ jest \emph{gładką zwartą} czterorozmaitością
z niepustym brzegiem $\Sigma(2,3,5)$;
\item $C$ jest \emph{topologiczną kontraktywną}
czterorozmaitością z przeciwnie zorientowanym brzegiem;
\item $X_{E_8}=P_{E_8}\cup C$ jest \emph{zamkniętą
czterorozmaitością topologiczną}, która nie może być wygładzona.
\end{enumerate}
Sama obecność formy $E_8$ na rozmaitości \emph{z brzegiem}
nie przeczy twierdzeniu Donaldsona. Sprzeczność powstaje
dopiero po gładkim zamknięciu rozmaitości, którego
w tym przykładzie nie da się wykonać.
\end{uwaga}

\section{Odczytywanie topologii z grafu: zestawienie i dalsze kierunki}
\label{sec:tg-plumbing-conclusions}

\begin{przyklad}[Trzy grafy i trzy różne informacje]
\begin{enumerate}[label=(\roman*)]
\item Łańcuch $A_3$ z wagami $-2$ ma macierz
$Q=\begin{psmallmatrix}-2&1&0\\1&-2&1\\0&1&-2\end{psmallmatrix}$,
wyznacznik $-4$ i brzeg o grupie $H_1\cong\mathbb Z/4\mathbb Z$.
\item Gwiazda $D_4$ ma wyznacznik $4$, ale
$H_1\cong\mathbb Z/2\oplus\mathbb Z/2$.
\item Graf $E_8$ ma wyznacznik $1$, więc jego brzeg jest
sferą homologiczną, lecz nie jest sferą $S^3$.
\end{enumerate}
Wyznacznik dostarcza pierwszej informacji o brzegu,
postać Smitha rozróżnia możliwe struktury grupy,
a grupa fundamentalna i geometryczne dane sklejania
potrafią rozróżnić także rozmaitości o identycznej homologii.
\end{przyklad}

\begin{uwaga}[Zakres stosowalności metody]
Wielkości $\det Q$, $\operatorname{sign}Q$ i
$\operatorname{coker}Q$ są dostępne przez elementarne rachunki
na macierzach. Nie wyznaczają jednak samodzielnie
typu homeomorfizmu ani dyfeomorfizmu brzegu.
Różne ważone grafy mogą opisywać tę samą rozmaitość,
a ich porównanie wymaga operacji Kirby'ego i odpowiednich
ruchów na grafach. Z kolei przy konstrukcji z pętlami,
powierzchniami bazowymi o dodatnim rodzaju czy grafami
o znakowanych krawędziach trzeba zachować dodatkowe
dane geometryczne: sam wektor wag nie wystarcza.
\end{uwaga}

\begin{uwaga}[Dalsze zastosowania]
Grafy plumbingowe pojawiają się w rozdzielczościach
osobliwości powierzchni zespolonych, teorii splotów
algebraicznych, rozmaitościach Seiferta, rachunku Kirby'ego,
a także w badaniach niezmienników Floera i nierówności
dla gładkich wypełnień. Bardziej zaawansowana konstrukcja
\emph{rational blow-down} zastępuje odpowiednie ujemnie
określone plumbingowanie przez racjonalną kulę homologiczną;
wymaga ona dodatkowej wiedzy o gładkich czterorozmaitościach.
Tutaj ograniczamy się do fundamentów potrzebnych do
rozumienia form przecięcia i przypadku $E_8$.
\end{uwaga}

\medskip
\noindent\textbf{Literatura do rozdziału.}
R. E. Gompf, A. I. Stipsicz, \emph{4-Manifolds and Kirby Calculus};
M. H. Freedman, F. Quinn, \emph{Topology of 4-Manifolds};
W. D. Neumann, \emph{A calculus for plumbing applied to the
topology of complex surface singularities and degenerating complex
curves}; R. Kirby, \emph{The topology of $4$-manifolds}.


% Iteracja 8: uzupelnienie teorii Seiferta i rachunku Kirby'ego.
\chapter{Rozmaitości Seiferta i grafy gwiaździste}
\label{ch:tg-seifert}

Z grafów plumbingowych odczytujemy formy przecięcia.
Przy odpowiednich założeniach dotyczących grafu możemy
ponadto rozpoznać rozwarstwienie jego brzegu na okręgi.
To prowadzi do \emph{rozmaitości Seiferta}. Zaczniemy od lokalnego
modelu, a następnie wyprowadzimy prezentację grupy fundamentalnej,
formułę na pierwszą homologię i związek z ułamkami łańcuchowymi.
Wszystkie trójrozmaitości w tym rozdziale są zwarte, spójne i zorientowane;
podstawowa orbipowierzchnia jest orientowalna. Nie omawiamy tutaj
odmian z nieorientowalną bazą ani z brzegiem lustrzanym.

\section{Włókna zwykłe i wyjątkowe: model lokalny}
\label{sec:tg-seifert-model}

\begin{definicja}[Rozwarstwienie Seiferta]
Niech $M$ będzie spójną, zorientowaną trójrozmaitością.
\emph{Rozwarstwieniem Seiferta} nazywamy rozkład $M$ na rozłączne
okręgi, dla którego każdy okrąg ma nasycone otoczenie dyfeomorficzne
z torusem pełnym $V_{\alpha,q}=(D^2\times S^1)/\mathbb Z_\alpha$,
z rozkładem na obrazy zbiorów $\{z\}\times S^1$.
Przy $\alpha\geq1$ grupa cykliczna działa według wzoru
\[
(z,u)\longmapsto(\zeta z,\zeta^q u),\qquad
\zeta=e^{2\pi i/\alpha},\quad \gcd(\alpha,q)=1.
\]
Dla $\alpha=1$ przyjmujemy działanie trywialne. Obraz
$\{0\}\times S^1$ jest \emph{włóknem wyjątkowym krotności $\alpha$},
gdy $\alpha>1$; włókna krotności jeden nazywamy \emph{zwykłymi}.
\end{definicja}

\begin{uwaga}[Dlaczego iloraz jest rozmaitością]
Jeśli niezerowa potęga generatora ustala $(z,u)$, to
$\zeta^{qk}u=u$, czyli $\alpha\mid qk$. Ponieważ
$\gcd(\alpha,q)=1$, wynika stąd $\alpha\mid k$.
Działanie jest więc wolne. Lokalny model nie zawiera
osobliwości rozmaitości; osobliwe są jedynie punkty
\emph{przestrzeni włókien}, gdzie pojawia się kąt stożkowy
$2\pi/\alpha$. Parametr $q$ wpływa na sposób sklejenia
włókna zwykłego z meridianem torusa pełnego.
\end{uwaga}

Przestrzeń włókien jest dwuwymiarową \emph{orbipowierzchnią}
$B$. W rozpatrywanym przypadku ma orientowalną powierzchnię
podstawową rodzaju $g$ i skończenie wiele punktów stożkowych
rzędów $\alpha_1,\ldots,\alpha_r$. Oznaczamy ją
$B(g;\alpha_1,\ldots,\alpha_r)$. Topologiczna
przestrzeń podkładowa orbipowierzchni nie pamięta rzędów
punktów stożkowych; dlatego samej powierzchni $B$ nie wolno
utożsamiać z pełnymi danymi rozwarstwienia.

\begin{center}
\begin{tikzpicture}[scale=.91,>=Stealth,line cap=round]
 \fill[tgBlue!9] (-3.6,-.1) ellipse (1.65 and .72);
 \draw[tgBlue,line width=1.4pt] (-3.6,-.1) ellipse (1.65 and .72);
 \draw[tgNavy!50,dashed] (-3.6,-.1) ellipse (.66 and .72);
 \foreach \x in {-4.1,-3.1}{\draw[tgTeal,line width=1.1pt] (\x,-.65) arc[start angle=-90,end angle=90,x radius=.19,y radius=.53];}
 \draw[tgCoral,line width=2.6pt] (-3.6,-.08) ellipse (.46 and .30);
 \node[font=\scriptsize,text=tgCoral] at (-3.6,-1.10) {włókno centralne};
 \draw[tgTeal,-{Stealth[length=2.8mm]},line width=1.6pt] (-1.3,-.1)--(1.05,-.1);
 \node[font=\small,text=tgTeal] at (-.1,.25) {iloraz};
 \fill[tgBlue!8] (3,-.1) ellipse (1.2 and .75);
 \draw[tgBlue,line width=1.4pt] (3,-.1) ellipse (1.2 and .75);
 \fill[tgCoral] (3,-.1) circle (.095);
 \node[font=\scriptsize,text=tgNavy] at (3,-1.11) {punkt stożkowy rzędu $\alpha$};
 \node[font=\scriptsize,text=tgNavy] at (0,-1.66)
 {Model lokalny: okręgi w torusie pełnym i odpowiadający im dysk bazowy.};
\end{tikzpicture}
\end{center}

\section{Niezmienniki Seiferta jako dane sklejenia}
\label{sec:tg-seifert-gluing}

Aby uniknąć niejednoznaczności znaków, zdefiniujemy
niezmienniki Seiferta przez konkretne nachylenia chirurgii.
Niech $B_{g,r+1}$ będzie zwartą, spójną, zorientowaną
powierzchnią rodzaju $g$ z $r+1$ brzegami, oznaczonymi
$\gamma_0,\gamma_1,\ldots,\gamma_r$.
Dla krzywych na powierzchni wybieramy punkt bazowy i drogi
do brzegów tak, aby w $\pi_1(B_{g,r+1})$ zachodziła relacja
\[
\gamma_0\gamma_1\cdots\gamma_r
\prod_{j=1}^g[a_j,b_j]=1.
\]
Na każdym torusie brzegowym $B_{g,r+1}\times S^1$ mamy
klasę krzywej $\gamma_i$ oraz klasę włókna $h$.

\begin{definicja}[Zorientowane dane Seiferta]
Niech $g\geq0$, $b\in\mathbb Z$, $r\geq0$,
$\alpha_i\geq2$ i $0<\beta_i<\alpha_i$ oraz
$\gcd(\alpha_i,\beta_i)=1$ dla $1\leq i\leq r$.
Przez
$M(g;b;(\alpha_1,\beta_1),\ldots,(\alpha_r,\beta_r))$
oznaczamy orientowaną rozmaitość powstającą z
$B_{g,r+1}\times S^1$ przez dołączenie torusów pełnych
w taki sposób, aby klasy meridianów dołączanych torusów
były odpowiednio
$[\gamma_0]+b[h]$ oraz
$\alpha_i[\gamma_i]+\beta_i[h]$ dla $1\leq i\leq r$.
W tej konwencji \emph{wymierną liczbą Eulera} nazywamy
$e_{\mathbb Q}(M):=-(b+\sum_{i=1}^r\beta_i/\alpha_i)$.
\end{definicja}

Wypełnienie torusa odpowiadającego $\gamma_0$ jest zwykłe
(tj. nie tworzy włókna wyjątkowego) i koduje liczbę
całkowitą $b$. Pozostałe tworzą włókna wyjątkowe
krotności $\alpha_i$. Są to \emph{dane rozwarstwienia},
a nie tylko niezmienniki samej rozmaitości: niektóre
trójrozmaitości dopuszczają więcej niż jedno rozwarstwienie
Seiferta. Zmiana konwencji orientacji i krzywych brzegowych
może prowadzić do innych zapisów tych samych danych.

\begin{center}
\begin{tikzpicture}[scale=.92,>=Stealth,line cap=round]
 \fill[tgBlue!9] (0,0) ellipse (2.1 and 1.15);
 \draw[tgBlue,line width=1.4pt] (0,0) ellipse (2.1 and 1.15);
 \foreach \x/\col/\lab in {-1.1/tgTeal/$\alpha_1$,0/tgCoral/$\alpha_2$,1.1/tgPurple/$\alpha_3$}{
  \draw[\col,line width=1.2pt] (\x,.1) circle (.33);
  \fill[\col!20] (\x,.1) circle (.29);
  \node[text=\col,font=\small] at (\x,.1) {\lab};
 }
 \node[font=\small,text=tgNavy] at (0,-1.54) {baza $S^2$ z trzema punktami stożkowymi};
 \draw[tgGold,line width=1.4pt,->] (2.45,.35)--(3.4,.35);
 \node[font=\footnotesize,text=tgGold,anchor=west] at (3.46,.35)
 {$S^1$ nad punktami zwykłymi};
\end{tikzpicture}
\end{center}

\section{Grupa fundamentalna: dowód przez van Kampena}
\label{sec:tg-seifert-pi1}

\begin{twierdzenie}[Prezentacja grupy rozmaitości Seiferta]
\label{thm:tg-seifert-presentation}
Niech $M=M(g;b;(\alpha_i,\beta_i)_{i=1}^r)$ będzie
rozmaitością z poprzedniej definicji. Wówczas jej grupa
fundamentalna ma prezentację
\[
\pi_1(M)\cong\Big\langle a_j,b_j,c_i,h\ \Big|\
\begin{array}{l}
[h,a_j]=[h,b_j]=[h,c_i]=1,\\
c_i^{\alpha_i}h^{\beta_i}=1\quad (1\leq i\leq r),\\
c_1\cdots c_r\prod_{j=1}^g[a_j,b_j]=h^b
\end{array}\Big\rangle.
\]
Indeksy $j$ przebiegają $1,\ldots,g$, a $i$ przebiegają
$1,\ldots,r$; dla pustych zbiorów indeksów pomijamy
dane generatory i relacje. Element $h$ jest reprezentowany
przez regularne włókno i należy do centrum grupy.
\end{twierdzenie}

\begin{proof}
Przed wypełnieniem torusów rozważamy produkt
$M_0=B_{g,r+1}\times S^1$. Ponieważ jest produktem,
$\pi_1(M_0)=\pi_1(B_{g,r+1})\times\langle h\rangle$,
gdzie $h$ komutuje z wszystkimi generatorami
powierzchni. Grupa $\pi_1(B_{g,r+1})$ ma
prezentację na generatorach $a_j,b_j,c_0,\ldots,c_r$
z relacją $c_0c_1\cdots c_r\prod_j[a_j,b_j]=1$;
wynika to np. z deformacyjnej retrakcji powierzchni
z brzegiem na odpowiedni graf i wycięcia dysków.

Dołączenie torusa pełnego wzdłuż wybranego
\emph{nachylenia} na torusie brzegowym narzuca relację,
że klasa tego nachylenia jest trywialna w grupie
fundamentalnej wypełnionej rozmaitości. Rzeczywiście
klasa ta jest obrazem meridianu torusa dołączanego,
a meridian ogranicza w nim dysk. Na mocy twierdzenia
van Kampena nie pojawia się inna niezależna relacja:
nowy generator, pochodzący z rdzenia torusa, można
wyeliminować, wykorzystując pierwotność nachylenia.
Nie chodzi zatem o zabicie samej klasy $\gamma_i$,
lecz o zabicie konkretnych klas
$\gamma_0h^b$ oraz $\gamma_i^{\alpha_i}h^{\beta_i}$.
Otrzymujemy dokładnie relacje $c_0h^b=1$ oraz
$c_i^{\alpha_i}h^{\beta_i}=1$.

Z pierwszej relacji mamy $c_0=h^{-b}$. Po podstawieniu
do relacji powierzchniowej uzyskujemy
$c_1\cdots c_r\prod_j[a_j,b_j]=h^b$.
Usunięcie generatora $c_0$ przez przekształcenie
Tietzego daje dokładnie wskazaną prezentację.
Centralność $h$ została odziedziczona z iloczynu
$B_{g,r+1}\times S^1$.
\end{proof}

\begin{wniosek}[Grupa orbipowierzchni jako iloraz]
Dla rozmaitości $M$ z twierdzenia
\ref{thm:tg-seifert-presentation} iloraz
$\pi_1(M)/\langle h\rangle$ ma prezentację
\[
\big\langle a_j,b_j,c_i\mid c_i^{\alpha_i}=1,
\ c_1\cdots c_r\prod_j[a_j,b_j]=1\big\rangle.
\]
Jest to grupa fundamentalna orbipowierzchni
$B(g;\alpha_1,\ldots,\alpha_r)$.
\end{wniosek}
\begin{proof}
Ponieważ $h$ jest centralny, podgrupa cykliczna
$\langle h\rangle$ jest normalna. W prezentacji
ilorazu dodajemy relację $h=1$, która usuwa $h$
ze wszystkich istniejących relacji. Pozostają dokładnie
zapisane generatory i relacje.
\end{proof}

\begin{uwaga}[Uwaga na rząd włókna]
Nie zawsze $\langle h\rangle\cong\mathbb Z$:
w rozmaitościach Seiferta o skończonej grupie fundamentalnej
włókno zwykłe może mieć skończony rząd. Dlatego poprawny
zapis sekwencji opisującej grupę jest
$1\to\langle h\rangle\to\pi_1(M)\to
\pi_1^{\mathrm{orb}}(B)\to1$,
a nie bezwarunkowo $1\to\mathbb Z\to\pi_1(M)\to\cdots$.
\end{uwaga}

\section{Pierwsza homologia i wymierna liczba Eulera}
\label{sec:tg-seifert-homology}

\begin{twierdzenie}[Homologia Seiferta nad sferą]
\label{thm:tg-seifert-h1}
Niech $M=M(0;b;(\alpha_i,\beta_i)_{i=1}^r)$,
gdzie $b\in\mathbb Z$, $\alpha_i\ge2$,
$0<\beta_i<\alpha_i$ oraz $\gcd(\alpha_i,\beta_i)=1$.
Grupa $H_1(M;\mathbb Z)$ ma prezentację abelową
z generatorami $c_1,\ldots,c_r,h$ i relacjami
$\alpha_ic_i+\beta_ih=0$ oraz
$c_1+\cdots+c_r-bh=0$.
Jeśli $e_{\mathbb Q}(M)\ne0$, grupa ta jest skończona
rzędu
\[
|H_1(M;\mathbb Z)|=
\left|\alpha_1\cdots\alpha_r
\left(b+\sum_{i=1}^r\frac{\beta_i}{\alpha_i}\right)\right|.
\]
Jeśli $e_{\mathbb Q}(M)=0$, to $H_1(M;\mathbb Z)$ ma
rangę jeden; jej ewentualną część torsyjną wyznacza
postać Smitha macierzy relacji.
\end{twierdzenie}

\begin{proof}
Abelianizacja prezentacji z twierdzenia
\ref{thm:tg-seifert-presentation} zabija wszystkie
komutatory. Otrzymujemy więc kwadratową macierz
relacji $A$ wymiaru $r+1$: elementy diagonalne
pierwszego bloku to $\alpha_i$, ostatnia kolumna
ma wartości $\beta_i$, a ostatni wiersz to
$(1,\ldots,1,-b)$.
Nad $\mathbb Q$ eliminujemy pierwsze $r$ generatorów.
W blokowej postaci
\[
A=\begin{pmatrix}
\operatorname{diag}(\alpha_1,\ldots,\alpha_r)&\beta\\
\mathbf 1^T&-b
\end{pmatrix}
\]
wyznacznik obliczamy przez dopełnienie Schura:
$\det A=\alpha_1\cdots\alpha_r
(-b-\sum_i\beta_i/\alpha_i)$.
Gdy ta liczba jest niezerowa, homologia jest skończonym
ilorazem $\mathbb Z^{r+1}/A^T\mathbb Z^{r+1}$
(o tej samej wartości wyznacznika), a jej rząd
wynosi $|\det A|$. Jeśli wyznacznik jest zerowy,
pierwszy blok ma nadal pełny rząd $r$, zatem
$\operatorname{rank}_{\mathbb Q}A=r$ i grupa
ilorazowa ma rangę jeden. To dowodzi wszystkich tez.
\end{proof}

\begin{przyklad}[Trzy wypełnienia $S^2(2,3,5)$]
Rozważmy dane $b=-2$ oraz
$(\alpha_1,\beta_1)=(2,1)$,
$(\alpha_2,\beta_2)=(3,2)$,
$(\alpha_3,\beta_3)=(5,4)$.
Wtedy
\[
b+\frac{\beta_1}{\alpha_1}
+\frac{\beta_2}{\alpha_2}
+\frac{\beta_3}{\alpha_3}
=-2+\frac12+\frac23+\frac45=-\frac1{30}.
\]
Z twierdzenia otrzymujemy $|H_1(M;\mathbb Z)|=30/30=1$,
czyli $H_1(M;\mathbb Z)=0$. Nie oznacza to
jeszcze, że $M\cong S^3$: grupa fundamentalna może
być nieprzemienna i doskonała. Pokażemy to poniżej.
\end{przyklad}

\begin{wniosek}[Ranga pierwszej homologii dla bazy rodzaju $g$]
\label{cor:tg-seifert-general-h1}
Niech $M=M(g;b;(\alpha_i,\beta_i)_{i=1}^r)$, przy czym
$g\geq0$, $r\geq0$ oraz pozostałe parametry są jak
w definicji danych Seiferta. Wówczas
\[
H_1(M;\mathbb Z)\cong\mathbb Z^{2g}\oplus A,
\]
gdzie $A$ jest grupą abelową o generatorach
$c_1,\ldots,c_r,h$ i relacjach
$\alpha_i c_i+\beta_i h=0$ oraz
$\sum_i c_i-bh=0$.
W szczególności, jeśli $b+\sum_i\beta_i/\alpha_i\ne0$,
to $\operatorname{rank}H_1(M;\mathbb Z)=2g$,
a gdy suma ta jest zerowa, ranga wynosi $2g+1$.
\end{wniosek}
\begin{proof}
Abelianizujemy prezentację z twierdzenia
\ref{thm:tg-seifert-presentation}. Wszystkie komutatory
$[a_j,b_j]$ stają się trywialne. Generatory
$a_1,b_1,\ldots,a_g,b_g$ nie występują w żadnej
pozostałej relacji, wyznaczają więc wolny składnik
$\mathbb Z^{2g}$. Pozostała grupa ma tę samą macierz
relacji, którą obliczyliśmy dla $g=0$.
Jej ranga nad $\mathbb Q$ wynosi zero albo jeden,
zależnie od tego, czy dopełnienie Schura
$-b-\sum_i\beta_i/\alpha_i$ jest niezerowe,
czy równe zeru.
\end{proof}

\begin{przyklad}[Wiązka trywialna nad torusem]
Dla $g=1$, $r=0$ oraz $b=0$ otrzymujemy
$M(1;0;)=T^2\times S^1=T^3$. Z prezentacji grupy
mamy $\pi_1(M)=\langle a,b,h\mid [a,b]=[a,h]=[b,h]=1\rangle
\cong\mathbb Z^3$, w zgodzie z równością
$H_1(T^3;\mathbb Z)\cong\mathbb Z^3$.
\end{przyklad}

\section{Od ramion plumbingowych do rozwarstwienia Seiferta}
\label{sec:tg-seifert-stars}

Niech $\Gamma$ będzie drzewem gwiaździstym z centralną
wagą $b$ oraz $r$ ramionami. W $i$-tym ramieniu wagi
wynoszą $-a_{i1},\ldots,-a_{ik_i}$, licząc od centrum,
gdzie $a_{ij}\ge2$. Załóżmy, że
\[
\frac{\alpha_i}{\beta_i}=
[a_{i1},\ldots,a_{ik_i}]^-
:=a_{i1}-\cfrac1{a_{i2}-\cfrac1{\ddots-\cfrac1{a_{ik_i}}}},
\quad \alpha_i>\beta_i>0,\quad\gcd(\alpha_i,\beta_i)=1.
\]
Dane $(\alpha_i,\beta_i)$ wyznaczają zatem ułamki
pochodzące z poszczególnych ramion.

\begin{stwierdzenie}[Wyznacznik grafu gwiaździstego]
\label{prop:tg-star-determinant}
Przy powyższych założeniach macierz przecięcia
$Q_\Gamma$ spełnia
\[
\big|\det Q_\Gamma\big|=
\left|\prod_{i=1}^r\alpha_i
\left(b+\sum_{i=1}^r\frac{\beta_i}{\alpha_i}\right)\right|.
\]
\end{stwierdzenie}
\begin{proof}
Niech $A_i$ będzie macierzą przecięcia $i$-tego ramienia.
Jest to ujemnie określona macierz trójdiagonalna o
przekątnej $(-a_{i1},\ldots,-a_{ik_i})$ i jedynkach
przy przekątnej. Rekurencja wyznaczników prowadzi do
$|\det A_i|=\alpha_i$ oraz
$(A_i^{-1})_{11}=-\beta_i/\alpha_i$.
Drugą tożsamość można także uzyskać wzorem na iloraz
minorów: usunięcie pierwszego wiersza i kolumny
zostawia wyznacznik odpowiadający ogonowi ułamka
łańcuchowego; po podzieleniu przez $\det A_i$
otrzymujemy $-1/[a_{i1},\ldots,a_{ik_i}]^-$.

Po usunięciu centralnego wierzchołka pozostały blok
$A=A_1\oplus\cdots\oplus A_r$ jest odwracalny.
Wektor sprzęgający centrum z ramionami ma jedynkę
przy pierwszym wierzchołku każdego ramienia.
Dopełnienie Schura wynosi zatem
$S=b-\sum_i(A_i^{-1})_{11}=b+\sum_i\beta_i/\alpha_i$.
Ponieważ $\det Q_\Gamma=(\prod_i\det A_i)S$,
po wzięciu modułu otrzymujemy tezę.
\end{proof}

\begin{twierdzenie}[Brzeg plumbingowania gwiaździstego]
\label{thm:tg-star-seifert}
Niech $\Gamma$ będzie ujemnie ważonym drzewem
gwiaździstym z centralną wagą $b$ i ramionami
$(-a_{i1},\ldots,-a_{ik_i})$, gdzie $a_{ij}\ge2$.
Niech $(\alpha_i,\beta_i)$ będą określone powyższymi
ułamkami łańcuchowymi. Wówczas, po ustaleniu opisanych
wyżej orientacji i baz meridian--włókno, brzeg
$\partial P_\Gamma$ jest rozmaitością Seiferta
$M(0;b;(\alpha_i,\beta_i)_{i=1}^r)$.
\end{twierdzenie}
\begin{proof}[Geometryczna redukcja do wypełnień]
Brzeg wiązki dyskowej $E(b)\to S^2$ jest okręgową
wiązką nad $S^2$. Usuwamy z bazy $r$ małych dysków,
w których przyłączane są ramiona. Nad pozostałą
powierzchnią otrzymujemy produkt z $S^1$, gdyż
każda zorientowana wiązka okręgowa nad powierzchnią
z niepustym brzegiem jest trywialna. Wybór
trywializacji koduje centralną liczbę Eulera $b$:
po dołączeniu brakującego zwykłego torusa pełnego
meridian ma w bazie $([\gamma_0],[h])$ nachylenie
$[\gamma_0]+b[h]$.

Każde ramię składa się z kolejnych wiązek dyskowych
nad sferami. Na poziomie brzegu jego diagram chirurgii
tworzy liniowy łańcuch składowych o ramowaniach
$-a_{i1},\ldots,-a_{ik_i}$. Reguła \emph{slam-dunk}
rachunku Kirby'ego pozwala kolejno usunąć końcowe
składowe tego łańcucha. Jeśli aktualny współczynnik
końcowej składowej jest $-r$, a składowa poprzednia
ma ramowanie $-a$, po eliminacji otrzymuje ona
współczynnik $-a+1/r=-(a-1/r)$. Rekurencja prowadzi
właśnie do ułamka
$-[a_{i1},\ldots,a_{ik_i}]^-=-\alpha_i/\beta_i$.
W wybranych w tym rozdziale orientacjach meridianu
i włókna oznacza to doklejenie torusa pełnego
wzdłuż klasy $\alpha_i[\gamma_i]+\beta_i[h]$.
Równość ułamków jest rachunkiem algebraicznym;
topologiczna poprawność samego ruchu \emph{slam-dunk}
stanowi lokalny lemat rachunku chirurgii
\ref{lem:tg-kirby-slam-dunk}. Korzystamy z jego opisu
przez sklejanie torusów,
a nie wywodzimy homeomorfizmu z równości wyznaczników.
Operacja ta jest lokalna w otoczeniu odpowiedniego
torusa brzegowego: nie zmienia pozostałych części
rozwarstwienia. Po zredukowaniu wszystkich ramion
otrzymujemy dokładnie konstrukcję z definicji danych
Seiferta. Dla pełnej zgodności orientacji trzeba używać
tej samej dodatniej normalnej do wszystkich torusów
sklejania; wprowadzone wcześniej znaki $+1$ przy
plumbingu gwarantują zgodność z tą konwencją.
\end{proof}

\begin{center}
\begin{tikzpicture}[scale=.95,>=Stealth]
  \node[tgRedVertex,minimum size=11mm] (c) at (0,0) {$b$};
  \node[tgBlueVertex] (a1) at (-1.35,.95) {$-a_{11}$};
  \node[tgBlueVertex] (a2) at (-2.8,.95) {$-a_{12}$};
  \node[tgTealVertex] (b1) at (1.35,.95) {$-a_{21}$};
  \node[tgPurple!80!black,draw=tgPurple,fill=tgPurple!12,circle,
       minimum size=9mm] (d1) at (0,-1.35) {$-a_{31}$};
  \node[tgPurple!80!black,draw=tgPurple,fill=tgPurple!12,circle,
       minimum size=9mm] (d2) at (0,-2.55) {$-a_{32}$};
  \draw[tgAccent] (c)--(a1);
  \draw[tgEdge] (a1)--(a2);
  \draw[tgTree] (c)--(b1);
  \draw[tgEdge] (c)--(d1)--(d2);
  \node[text=tgNavy,font=\footnotesize] at (0,-3.20)
  {Każde ramię zastępujemy jedną parą $(\alpha_i,\beta_i)$.};
\end{tikzpicture}
\end{center}

\begin{przyklad}[Brzeg plumbingowania $E_8$: grupa Poincarégo]
\label{ex:tg-seifert-E8}
Dla $E_8$ centralna waga wynosi $b=-2$,
a trzy ramiona mają odpowiednio długości $1,2,4$
z wagami $-2$. Dlatego
$[2]^-=2/1$, $[2,2]^-=3/2$ i
$[2,2,2,2]^-=5/4$. Na mocy
\ref{thm:tg-star-seifert} brzeg jest rozmaitością
\[
M(0;-2;(2,1),(3,2),(5,4)).
\]
Jej grupa fundamentalna ma prezentację
\[
\left\langle c_1,c_2,c_3,h\ \middle|\
h\text{ centralny},\ c_1^2h=1,\ c_2^3h^2=1,
\ c_3^5h^4=1,\ c_1c_2c_3=h^{-2}\right\rangle.
\]
Po przejściu do ilorazu przez centralną podgrupę
$\langle h\rangle$ otrzymujemy grupę
$\Delta(2,3,5)=\langle x,y,z\mid x^2=y^3=z^5=xyz=1\rangle$.
Aby zobaczyć, że iloraz ten jest nietrywialny,
można wziąć w grupie alternującej $A_5$ permutacje
$s=(12)(34)$ i $t=(235)$. Wtedy $s^2=t^3=(st)^5=1$,
a $s,t$ generują $A_5$. Przypisania
$x\mapsto s$, $y\mapsto t$, $z\mapsto(st)^{-1}$
określają surjekcję $\Delta(2,3,5)\twoheadrightarrow A_5$.
W szczególności $\pi_1(\partial P_{E_8})\ne1$.
Z poprzedniego przykładu mamy jednocześnie
$H_1(\partial P_{E_8};\mathbb Z)=0$.
To dowodzi, że brzeg jest nietrywialną sferą homologiczną.

Standardowa identyfikacja rozwarstwienia o tych danych
ze sferą Brieskorna $\Sigma(2,3,5)$ noszącą nazwę
\emph{sfery homologicznej Poincarégo} jest osobnym
rezultatem geometrii rozwarstwień Seiferta. Przyjmujemy
tę identyfikację z literatury; obliczenia homologii i
dowód nietrywialności grupy fundamentalnej jej nie
potrzebują. W istocie grupa fundamentalna ma rząd $120$
i jest binarną grupą dwudziestościenną.
\end{przyklad}

\begin{uwaga}[Przypadki z co najwyżej dwoma włóknami wyjątkowymi]
Gdy podstawą jest $S^2$ i liczba włókien wyjątkowych
jest nie większa niż dwa, rozmaitość Seiferta jest
uzyskiwana przez sklejenie dwóch torusów pełnych.
Otrzymujemy przestrzeń soczewkową (w przypadkach
granicznych także $S^1\times S^2$). Wniosek wynika
z rozcięcia $S^2$ wzdłuż okręgu oddzielającego
punkty stożkowe: przedobrazy obu dysków są torusami
pełnymi. Należy odróżnić ten opis od klasyfikacji
orientowanych przestrzeni soczewkowych według
parametrów $(p,q)$, której nie wyprowadzamy tutaj.
\end{uwaga}

\section{Orbicharakterystyka Eulera i geometria}
\label{sec:tg-seifert-orbifold}

\begin{definicja}[Orbicharakterystyka]
Dla orientowalnej orbipowierzchni
$B(g;\alpha_1,\ldots,\alpha_r)$ definiujemy
\[
\chi_{\mathrm{orb}}(B)=2-2g-
\sum_{i=1}^r\left(1-\frac1{\alpha_i}\right).
\]
\end{definicja}

\begin{przyklad}
Dla bazy $S^2(2,3,5)$ mamy
$\chi_{\mathrm{orb}}=2-[(1-1/2)+(1-1/3)+(1-1/5)]
=1/30>0$. Dla $S^2(2,3,6)$ otrzymujemy
$\chi_{\mathrm{orb}}=0$, a dla $S^2(2,3,7)$ mamy
$\chi_{\mathrm{orb}}=-1/42<0$.
Znaki te odpowiadają odpowiednio trójkątnym
geometriom sferycznej, euklidesowej i hiperbolicznej
orbipowierzchni. Nie określają one samodzielnie geometrii
trójrozmaitości: znaczenie ma również wymierna
liczba Eulera danego rozwarstwienia.
\end{przyklad}

\begin{uwaga}[Miejsce w geometryzacji]
Wśród ośmiu geometrii Thurstona rozmaitości Seiferta
realizują, zależnie od danych i przypadków wyjątkowych,
geometrie $S^3$, $S^2\times\mathbb R$, $\mathbb E^3$,
$\mathrm{Nil}$, $\mathbb H^2\times\mathbb R$
i $\widetilde{\mathrm{SL}}_2(\mathbb R)$.
Odpowiednie pełne twierdzenie klasyfikacyjne wymaga
teorii geometryzacji i jest tutaj przyjmowane bez dowodu.
\end{uwaga}

\chapter{Rachunek Kirby'ego: operacje na ramowanych splotach i grafach}
\label{ch:tg-kirby-full}

Rachunek Kirby'ego pozwala porównywać różne opisy
tej samej trójrozmaitości jako brzegu czterorozmaitości
uchwytowej. Ważne jest rozróżnienie dwóch pytań:
czy dwie operacje dają tę samą rozmaitość \emph{brzegową}
i czy zachowują całą \emph{czterorozmaitość}.
Przesunięcie uchwytu zachowuje obie;
rozdmuchanie i kontrakcja składnika o ramowaniu $\pm1$
mogą natomiast zmienić czterorozmaitość,
choć nie zmienia zorientowanej trójrozmaitości brzegowej.

\section{Ramowanie, równoległa kopia i przesunięcie uchwytu}
\label{sec:tg-kirby-slides}

\begin{definicja}[Przesunięcie uchwytu indeksu dwa]
Niech $L=K_1\sqcup\cdots\sqcup K_r\subset S^3$
będzie zorientowanym splotem o ustalonych całkowitych
ramowaniach. \emph{Przesunięciem uchwytu $K_i$ po $K_j$}
($i\ne j$) nazywamy zastąpienie $K_i$ sumą pasmową
z równoległą kopią $K_j'$ wyznaczoną ramowaniem $K_j$.
Pasmo wybieramy rozłącznie od pozostałych składowych.
Orientację kopii możemy wybrać zgodną lub przeciwną;
w obliczeniach algebraicznych odpowiada to znakom $+$ i $-$.
\end{definicja}

\begin{stwierdzenie}[Zmiana macierzy przecięcia]
\label{prop:tg-kirby-congruence}
Niech $W_L$ będzie czterorozmaitością złożoną z $B^4$
i $r$ uchwytów indeksu dwa, a $Q_L$ jej macierzą przecięcia.
Po przesunięciu uchwytu $i$ po uchwycie $j$ ze znakiem
$\varepsilon\in\{\pm1\}$ otrzymujemy macierz
$Q_{L'}=P^TQ_LP$, gdzie $P\in GL_r(\mathbb Z)$
zastępuje wektor bazowy $e_i$ przez
$e_i+\varepsilon e_j$ i pozostawia pozostałe bez zmian.
W szczególności $W_L$ i $W_{L'}$ są dyfeomorficzne
(z odpowiednimi identyfikacjami granic), a ich brzegi
są orientacyjnie dyfeomorficzne.
\end{stwierdzenie}

\begin{proof}
W rachunku uchwytów rdzenie dołączanych uchwytów indeksu dwa
po uzupełnieniu powierzchniami rozpiętymi na krzywych
przyczepienia wyznaczają bazę $H_2(W_L;\mathbb Z)$.
Przesunięcie uchwytu odpowiada połączeniu tych powierzchni
pasmem, więc realizuje w homologii operację
$e_i\mapsto e_i+\varepsilon e_j$.
Przy zmianie bazy formy dwuliniowej jej macierz
przechodzi w $P^TQ_LP$. Elementarna zmiana bazy
jest unimodularna; stąd nie zmienia wyznacznika ani
sygnatury. Fakt, że ruch pasmowy realizuje dyfeomorfizm
czterorozmaitości, jest standardowym lokalnym lematem
rachunku uchwytów: przesuwamy obszar przyczepienia
jednego uchwytu przez kołnierz drugiego, nie naruszając
pozostałych części. Do pełnej konstrukcji dyfeomorfizmu
potrzebne jest izotopijne twierdzenie o przyczepianiu
uchwytów, wcześniej wprowadzone w teorii Morse'a.
\end{proof}

\begin{przyklad}[Dwie krzywe z liczbą splecenia $\ell$]
Niech $Q=\begin{psmallmatrix}a&\ell\\\ell&b\end{psmallmatrix}$.
Po dodaniu drugiego wektora do pierwszego otrzymujemy
\[
Q'=\begin{pmatrix}a+2\ell+b&\ell+b\\\ell+b&b\end{pmatrix}.
\]
Wpis diagonalny zwiększył się o $2\ell+b$,
a nowa liczba splecenia wynosi $\ell+b$.
Bezpośrednio
$\det Q'=(a+2\ell+b)b-(\ell+b)^2=ab-\ell^2$.
Rachunek ten jednocześnie kontroluje ramowanie,
liczby splecenia i niezmienniczość wyznacznika.
\end{przyklad}

\begin{lemat}[Reguła \emph{slam-dunk}]
\label{lem:tg-kirby-slam-dunk}
Niech $U\cup V\subset S^3$ będzie standardowym zorientowanym
splotem Hopfa, którego składowe mają liczbę splecenia $+1$.
Niech $r,s\in\mathbb Q$, przy czym $s\ne0$.
Wykonanie chirurgii o współczynnikach $r$ na $U$
i $s$ na jego geometrycznym meridianie $V$ jest
równoważne wykonaniu na pojedynczym węźle trywialnym
chirurgii o współczynniku $r-1/s$. Wszystkie
współczynniki odczytujemy względem preferowanych
długości i z ustaloną orientacją $S^3$.
\end{lemat}
\begin{proof}[Opis sklejenia torusów]
Zewnętrze splotu Hopfa jest dyfeomorficzne
z $T^2\times[0,1]$. Na jego dwóch torusach brzegowych
możemy wybrać meridiany $\mu_U,\mu_V$ i długości
$\lambda_U,\lambda_V$ tak, że
$[\lambda_U]=[\mu_V]$ i $[\lambda_V]=[\mu_U]$
po identyfikacji w przestrzeni produktu.
Po dołączeniu torusa pełnego wzdłuż nachylenia
$s[\mu_V]+[\lambda_V]$ pozostaje rozmaitość
z jednym torusowym brzegiem. Złożenie odpowiadających
mu identyfikacji południka i długości z mapą sklejającą
$T^2\times[0,1]$ przekształca pozostałe nachylenie
$r[\mu_U]+[\lambda_U]$ w nachylenie opisane
ułamkiem $r-1/s$. Współczynnik $-1/s$ pochodzi
z odwrócenia bazy przy przejściu przez zewnętrze
splotu Hopfa. Dokładnie tę samą regułę otrzymujemy
z macierzy sklejenia torusów w $SL_2(\mathbb Z)$.
Gdy $r,s$ są ułamkami, wszystkie wyrażenia należy
interpretować przez pierwotne całkowite pary
współczynników nachylenia; wynik nie zależy od ich
pomnożenia przez wspólną niezerową liczbę całkowitą.
\end{proof}

\begin{przyklad}[Początek redukcji łańcucha]
Dwa kolejne wierzchołki plumbingowe o wagach
$-5,-2$ odpowiadają splotowi Hopfa o tych samych
ramowaniach. Po usunięciu składowej o ramowaniu
$-2$ otrzymujemy współczynnik
$-5-1/(-2)=-9/2$. Jest to właśnie ułamek
$-[5,2]^-$, który pojawi się przy konfiguracji
$C_3$ w rozdziale o wymiernej kontrakcji.
\end{przyklad}

\section{Rozdmuchanie i kontrakcja składnika o ramowaniu \texorpdfstring{$\pm1$}{+/-1}}
\label{sec:tg-kirby-blowdown}

\begin{twierdzenie}[Usuwanie składowej o ramowaniu $\pm1$]
\label{thm:tg-kirby-blowdown}
Niech $L=U\sqcup K_1\sqcup\cdots\sqcup K_r$
będzie całkowicie ramowanym splotem w $S^3$,
przy czym $U$ jest węzłem trywialnym o ramowaniu
$\varepsilon\in\{+1,-1\}$.
Niech $\ell_i=\operatorname{lk}(U,K_i)$ i niech
$B$ będzie macierzą splecenia składowych $K_i$.
Wtedy macierz całego splotu ma postać
$Q=\begin{psmallmatrix}\varepsilon&\ell^T\\\ell&B\end{psmallmatrix}$.
Po kontrakcji $U$ i odpowiednim pełnym skręceniu
pozostałych pasm otrzymujemy splot $L'$ o macierzy
\[
B'=B-\varepsilon\,\ell\ell^T.
\]
Chirurgie na $L$ i $L'$ dają orientacyjnie dyfeomorficzne
trójrozmaitości brzegowe. Na poziomie form przecięcia
zachodzi całkowita kongruencja
$Q\cong[\varepsilon]\oplus B'$.
Nie twierdzimy, że odpowiadające czterorozmaitości
są dyfeomorficzne: usunięcie składnika o ramowaniu
$\pm1$ może zmienić formę przecięcia we wnętrzu.
\end{twierdzenie}

\begin{proof}
Dla generatorów homologii $u,e_1,\ldots,e_r$
formy przecięcia mamy $u^2=\varepsilon$,
$u\cdot e_i=\ell_i$, $e_i\cdot e_j=B_{ij}$.
Przyjmijmy
$f_i=e_i-\varepsilon\ell_i u$. Ponieważ
$\varepsilon^2=1$, otrzymujemy
$u\cdot f_i=\ell_i-\varepsilon\ell_i\varepsilon=0$.
Ponadto
\[
f_i\cdot f_j=B_{ij}
-\varepsilon\ell_i\ell_j
-\varepsilon\ell_i\ell_j
+\varepsilon^2\ell_i\ell_j\varepsilon
=B_{ij}-\varepsilon\ell_i\ell_j.
\]
Zmiana bazy jest całkowita i unimodularna,
a zatem $Q\cong[\varepsilon]\oplus B'$.

Aby uzyskać twierdzenie o brzegach, trzeba ponadto
wykonać geometryczny ruch Kirby'ego: operacja
$\varepsilon$-chirurgii na węźle trywialnym $U$
przywraca $S^3$, a podczas tego wypełnienia łuki
pozostałych składowych przecinające dysk rozpięty na $U$
otrzymują pełny skręt o znaku przeciwnym do $\varepsilon$.
Współczynniki skręcenia wyznaczają nowe ramowania
i liczby splecenia opisane przez macierz $B'$.
Ta część argumentu jest geometrycznym lematem Kirby'ego:
obejmuje zmianę nachyleń na torusie brzegowym
realizowaną przez odpowiedni skręt Dehna.
Nie wynika ona z samej kongruencji macierzy.
Łącząc ten ruch z rachunkiem formy przecięcia,
otrzymujemy tezę.
\end{proof}

\begin{przyklad}[Liść o wadze $-1$ i centralny wierzchołek]
Załóżmy, że $U$ jest liściem grafu plumbingowego,
ma wagę $-1$ i łączy się tylko z wierzchołkiem
o wadze $b$, z liczbą splecenia $1$.
Wtedy $\varepsilon=-1$ i $\ell=(1,0,\ldots,0)$.
Z twierdzenia wynika $b'=b+1$, a pozostałe wagi
nie zmieniają się. Jest to klasyczna kontrakcja
liścia o wadze $-1$. Dla liścia z wagą $+1$
otrzymujemy odpowiednio $b'=b-1$.
\end{przyklad}

\begin{przyklad}[Kontrakcja węzła o dwóch sąsiadach]
Rozpatrzmy łańcuch z trzema wierzchołkami i wagami
$a,-1,b$. Macierz przecięcia wynosi
\[
Q=\begin{pmatrix}a&1&0\\1&-1&1\\0&1&b\end{pmatrix}.
\]
Usunięcie środkowego wierzchołka daje macierz
$B'=\begin{psmallmatrix}a+1&1\\1&b+1\end{psmallmatrix}$.
Nowe składowe są splecione raz, choć pierwotnie
końcowe wierzchołki nie były sąsiadami.
W szczególności nie wystarcza usunięcie wierzchołka
i zmiana samych wag: trzeba również dodać nową
krawędź grafu. Wynika to z pozadiagonalnego składnika
$-(-1)\ell_1\ell_2=1$.
\end{przyklad}

\begin{center}
\begin{tikzpicture}[scale=.92,>=Stealth]
 \node[tgBlueVertex] (a) at (-3,0) {$a$};
 \node[tgRedVertex] (x) at (-1.55,0) {$-1$};
 \node[tgTealVertex] (b) at (-.10,0) {$b$};
 \draw[tgEdge] (a)--(x)--(b);
 \draw[tgGold,-{Stealth[length=3mm]},line width=1.6pt] (.6,0)--(2.0,0);
 \node[tgBlueVertex] (c) at (2.8,0) {$a+1$};
 \node[tgTealVertex] (d) at (4.45,0) {$b+1$};
 \draw[tgAccent] (c)--(d);
 \node[font=\footnotesize,text=tgNavy] at (.55,-.97)
 {Usunięcie wewnętrznego wierzchołka $-1$ tworzy nową krawędź.};
\end{tikzpicture}
\end{center}

\begin{uwaga}[Gdy wynik nie jest zwykłym grafem]
Jeżeli usuwany wierzchołek o wadze $\pm1$ ma
kilku sąsiadów, eliminacja zmienia nie tylko ich
wagi, ale i wszystkie wzajemne elementy pozadiagonalne.
Może to prowadzić do krawędzi wielokrotnych,
ujemnych liczb splecenia lub niezerowych elementów
między wierzchołkami już sąsiadującymi. Dlatego
ogólny rachunek Kirby'ego wygodniej prowadzić
na \emph{ramowanych splotach} lub na macierzach
ze znakowanymi krawędziami, a nie na samych
nieskierowanych drzewach prostych.
\end{uwaga}

\section{Kontrola brzegu: wyznacznik i forma splecenia}
\label{sec:tg-kirby-linking}

\begin{stwierdzenie}[Postać formy splecenia po ruchu Kirby'ego]
Niech $Q$ będzie nieosobliwą macierzą splecenia
całkowicie ramowanego splotu i niech $Y$ będzie
brzegiem odpowiadającej mu czterorozmaitości uchwytowej.
Wtedy forma splecenia $\lambda_Y$ na
$H_1(Y;\mathbb Z)\cong\operatorname{coker}Q$ jest
opisywana macierzą $-Q^{-1}$ modulo $\mathbb Z$.
Przesunięcie uchwytu i kontrakcja składowej
$\varepsilon=\pm1$ zachowują klasę izometrii
tej formy, mimo że przy kontrakcji może zmienić
się ranga macierzy przecięcia.
\end{stwierdzenie}
\begin{proof}
Formułę na formę splecenia wyprowadziliśmy wcześniej
z dualności Poincarégo--Lefschetza. Przesunięcie
uchwytu daje $Q'=P^TQP$. Wówczas
${Q'}^{-1}=P^{-1}Q^{-1}P^{-T}$, a zmiana
prezentacji grupy ilorazowej przez unimodularną
macierz przenosi formę $-Q^{-1}$ na $-{Q'}^{-1}$.

Przy kontrakcji otrzymujemy
$Q\cong[\varepsilon]\oplus B'$.
Składnik $[\varepsilon]$ jest unimodularny i ma
trywialną grupę ilorazową; nie wnosi zatem nic do
formy splecenia na brzegu. Pozostaje forma
odpowiadająca $B'$. Stąd w obu operacjach
forma brzegowa jest zachowana.
\end{proof}

\begin{przyklad}[Dwa równoważne diagramy i grupa $\mathbb Z/5$]
Niech $Q=\begin{psmallmatrix}-1&1\\1&-6\end{psmallmatrix}$.
Po usunięciu wierzchołka $-1$ dostajemy $B'=[-5]$.
Istotnie, $-6-(-1)^{-1}\cdot1^2=-5$.
Wyznaczniki to odpowiednio $5$ i $-5$,
więc ich moduły są jednakowe. Obie chirurgie
mają pierwszą grupę homologii $\mathbb Z/5$,
a forma splecenia w każdym przypadku jest
przedstawiona, po odpowiedniej identyfikacji
generatora, przez $-[-5]^{-1}=1/5\pmod{\mathbb Z}$.
Ten przykład pokazuje, dlaczego wyznacznik
przy usuwaniu składnika $-1$ zmienia znak, podczas
gdy grupa brzegowa i jej forma pozostają niezmienione.
\end{przyklad}

\section{Granice rachunku macierzowego i twierdzenie Kirby'ego}
\label{sec:tg-kirby-scope}

\begin{twierdzenie}[Kirby; postać klasyczna, bez dowodu]
Na mocy twierdzenia Lickorisha--Wallace'a każda zamknięta,
spójna, zorientowana trójrozmaitość powstaje przez
całkowitą chirurgię Dehna wzdłuż odpowiednio
ramowanego splotu w $S^3$.
Dwa takie przedstawienia dają dyfeomorficzne
trójrozmaitości z zachowaniem orientacji wtedy i tylko wtedy,
gdy można je połączyć skończoną liczbą izotopii
ramowanych splotów, przesunięć uchwytów oraz
dodawania lub usuwania \emph{rozłącznej} składowej
trywialnej z ramowaniem $+1$ albo $-1$.
\end{twierdzenie}

Przyjmujemy pełną implikację odwrotną twierdzenia
Kirby'ego bez dowodu; wymaga ona niezależnego
rachunku uchwytów. Ważne jest, że \emph{kongruencja
macierzy} nie jest kryterium równoważności chirurgii:
różne ramowane sploty mogą mieć tę samą macierz
splecenia, lecz dawać niehomeomorficzne brzegi.
Macierz rejestruje homologię i formę splecenia,
a nie cały typ topologiczny trójrozmaitości.

\begin{uwaga}[Granica między 3D i 4D]
Dodanie oddzielnego węzła trywialnego o ramowaniu
$+1$ lub $-1$ nie zmienia brzegu chirurgii,
ale na poziomie czterowymiarowym jest dołączeniem
składnika $\mathbb{CP}^2$ lub
$\overline{\mathbb{CP}}^{\,2}$ (z usuniętą kulą).
Przesunięcia uchwytów nie zmieniają samej
czterorozmaitości. Przy zastosowaniach do twierdzeń
Freedmana i Donaldsona trzeba więc zawsze pamiętać,
który rodzaj równoważności jest aktualnie rozważany.
\end{uwaga}

\begin{uwaga}[Połączenie z rozdziałem o $E_8$]
Grafy $A_n$, $D_n$ i $E_8$ są przykładami, w których
konstrukcje brzegu, jego homologii i formy przecięcia
można obliczyć bezpośrednio z macierzy grafu. Natomiast
identyfikacja konkretnej trójrozmaitości brzegowej
wymaga informacji geometrycznej: rozwarstwienia
Seiferta, grupy fundamentalnej, diagramu Kirby'ego
lub innych niezmienników. Dlatego oba nowe rozdziały
uzupełniają, ale nie zastępują, wcześniejszą teorię
węzłów i rozmaitości czterowymiarowych.
\end{uwaga}


% Iteracja 9. Nowy rozdział po rachunku Kirby'ego.
\begin{lemat}[Algebraiczna postać kontrakcji składowej o ramowaniu $\pm1$]
\label{lem:tg-audit-schur}
Niech $\varepsilon\in\{-1,1\}$ i
$Q=\begin{psmallmatrix}B&\ell\\\ell^T&\varepsilon\end{psmallmatrix}$
będzie symetryczną macierzą całkowitą, gdzie $B$ jest
kwadratowa. Wówczas jest ona kongruentna nad
$\mathbb Z$ do
$\operatorname{diag}(B-\varepsilon\ell\ell^T,\varepsilon)$.
W szczególności
$\det Q=\varepsilon\det(B-\varepsilon\ell\ell^T)$.
\end{lemat}

\begin{proof}
Wystarczy wykonać całkowitą zmianę bazy, która zastępuje
$e_i$ przez $e_i-\varepsilon\ell_i e_{r+1}$ i zachowuje
$e_{r+1}$. W nowej bazie iloczyn
$(e_i-\varepsilon\ell_i e_{r+1})\cdot e_{r+1}$
wynosi $\ell_i-\varepsilon\ell_i\varepsilon=0$,
a iloczyn dwóch nowych wektorów jest równy
$B_{ij}-\varepsilon\ell_i\ell_j$.
Macierz przejścia jest trójkątna z wyznacznikiem jeden,
więc jest unimodularna. Daje to żądaną kongruencję
i, po wzięciu wyznaczników, wzór determinantowy.
Dowodzi to wyłącznie części macierzowej operacji;
geometryczna realizacja kontrakcji jest osobnym krokiem
rachunku Kirby'ego.
\end{proof}

\chapter{Wymierna kontrakcja: od grafów plumbingowych do egzotycznych czterorozmaitości}
\label{ch:tg-rational-blowdown}

Przy zwykłym \emph{rozdmuchaniu} (\emph{blow-up}) gładkiej
czterorozmaitości wprowadzamy sferę o samoprzecięciu $-1$,
a przy kontrakcji (\emph{blow-down}) taką sferę usuwamy.
Istnieje istotnie subtelniejsza operacja: można usunąć całe
\emph{ujemnie określone plumbingowanie sfer} i zastąpić je
czterorozmaitością, która ma wymierną homologię punktu.
Nazywa się to \emph{wymierną kontrakcją} (\emph{rational
blow-down}). Wprowadzona przez Fintushela i Sterna operacja
dostarcza metod konstruowania egzotycznych struktur gładkich.

Kluczowe jest rozróżnienie trzech kroków: (i) obliczamy brzeg
plumbingowania, (ii) korzystamy z istnienia \emph{innego}
wypełnienia tego samego brzegu, (iii) dopiero wtedy
przeprowadzamy operację wycięcia i sklejenia. Rachunek
macierzy przecięcia uzasadnia pierwszy krok homologicznie,
ale sam nie dowodzi istnienia drugiego wypełnienia.

\section{Łańcuch \texorpdfstring{$C_p$}{Cp} i jego ujemnie określona forma przecięcia}
\label{sec:tg-rb-config}

\begin{definicja}[Konfiguracja Fintushela--Sterna]
Niech $p\geq2$ będzie liczbą całkowitą.
Przez $C_p$ oznaczamy zwartą, spójną, zorientowaną,
gładką czterorozmaitość plumbingową związaną z liniowym
grafem o $p-1$ wierzchołkach, z wagami
$-(p+2),-2,\ldots,-2$; waga $-2$ występuje
$p-2$ razy, a kolejne krawędzie mają dodatni znak
przecięcia. W bazie sfer zerowych forma przecięcia jest opisana
macierzą symetryczną $Q_p$ stopnia $p-1$, której elementy
określamy wzorem
\[
(Q_p)_{ij}=\begin{cases}
-(p+2),&i=j=1,\\
-2,&i=j\geq2,\\
1,&|i-j|=1,\\
0,&\text{w pozostałych przypadkach}.
\end{cases}
\]
W szczególności $Q_2=[-4]$,
$Q_3=\begin{psmallmatrix}-5&1\\1&-2\end{psmallmatrix}$.
Dla $p\geq4$ jest to macierz trójdiagonalna,
z początkowym wpisem $-(p+2)$, dalszymi wpisami
$-2$ i jedynkami bezpośrednio przy przekątnej.
\end{definicja}

\begin{center}
\begin{tikzpicture}[>=Stealth,scale=.95]
  \node[tgRedVertex,minimum size=10mm] (a) at (-3.35,0) {$-(p+2)$};
  \node[tgBlueVertex] (b) at (-1.35,0) {$-2$};
  \node[tgBlueVertex] (c) at (.20,0) {$-2$};
  \node[font=\large,text=tgNavy] (d) at (1.43,0) {$\cdots$};
  \node[tgBlueVertex] (e) at (2.80,0) {$-2$};
  \draw[tgTree] (a)--(b)--(c)--(d)--(e);
  \node[font=\footnotesize,text=tgNavy] at (0,-.85)
    {$p-1$ wierzchołków; rysunek schematyczny dla $p\geq5$.};
\end{tikzpicture}
\end{center}

\begin{lemat}[Minory główne i określoność]
\label{lem:tg-rb-minors}
Dla każdej liczby całkowitej $p\geq2$ i każdego
$1\leq k\leq p-1$ wyznacznik górnego lewego bloku
$Q_p$ rozmiaru $k$ jest równy
$(-1)^k(1+k(p+1))$. W szczególności
\[
\det Q_p=(-1)^{p-1}p^2,
\qquad \operatorname{sign}(Q_p)=-(p-1).
\]
\end{lemat}

\begin{proof}
Niech $d_k$ będzie modułem wyznacznika $k\times k$
początkowego bloku macierzy $-Q_p$.
Przyjmujemy $d_0=1$. Mamy $d_1=p+2$.
Dla $2\leq k\leq p-1$ rozwinięcie wyznacznika
trójdiagonalnego względem ostatniego wiersza daje
$d_k=2d_{k-1}-d_{k-2}$. Dowodzimy indukcyjnie,
że $d_k=1+k(p+1)$: sprawdza się to dla $k=0,1$,
a jeśli wzór zachodzi dla dwóch poprzednich indeksów,
to $2[1+(k-1)(p+1)]-[1+(k-2)(p+1)]
=1+k(p+1)$.
Każdy początkowy minor główny macierzy $-Q_p$
jest dodatni, więc kryterium Sylvestera dowodzi,
że $-Q_p$ jest dodatnio określona. Wstawienie
$k=p-1$ daje $d_{p-1}=1+(p-1)(p+1)=p^2$.
Dla $p=2$ rozumowanie obejmuje pojedynczy minor $4$.
\end{proof}

\begin{twierdzenie}[Homologia brzegu $C_p$]
\label{thm:tg-rb-h1boundary}
Dla każdego $p\geq2$ brzeg $Y_p:=\partial C_p$
jest wymierną sferą homologiczną i zachodzi
$H_1(Y_p;\mathbb Z)\cong\mathbb Z/p^2\mathbb Z$.
\end{twierdzenie}

\begin{proof}
Twierdzenie o brzegu drzewa plumbingowego daje
$H_1(Y_p;\mathbb Z)\cong\operatorname{coker}Q_p$.
Z poprzedniego lematu znamy już rząd tej grupy,
równej $p^2$. Pokażemy, że jest cykliczna;
nie wynika to z samego wyznacznika.

Niech $e_1,\ldots,e_{p-1}$ oznaczają generatory
prezentacji ilorazowej, w kolejności od wierzchołka
o wadze $-(p+2)$. Jeśli $p\geq3$, pierwsza relacja
ma postać $-(p+2)e_1+e_2=0$, a kolejne relacje
wewnętrzne dają $e_{j+1}=2e_j-e_{j-1}$.
Stąd przez indukcję
$e_j=[1+(j-1)(p+1)]e_1$ dla $1\leq j\leq p-1$.
Ostatnia relacja $e_{p-2}-2e_{p-1}=0$
(dla $p=3$ interpretowana jako pierwsza i ostatnia
relacja w macierzy $2\times2$) staje się
$p^2e_1=0$, z pominięciem nieistotnego znaku.
Tak więc cała grupa jest generowana przez $e_1$,
a jej rząd wynosi $p^2$. Dla $p=2$ jedyna relacja
$-4e_1=0$ daje ten sam wynik.
\end{proof}

\section{Brzeg jako przestrzeń soczewkowa}
\label{sec:tg-rb-lens}

\begin{lemat}[Ułamek łańcuchowy konfiguracji]
\label{lem:tg-rb-cfrac}
Dla każdej liczby całkowitej $p\geq2$ zachodzi
\[
[p+2,\underbrace{2,\ldots,2}_{p-2}]^-
=\frac{p^2}{p-1}.
\]
\end{lemat}

\begin{proof}
Przez $[a_1,\ldots,a_k]^-$ oznaczamy
$a_1-1/(a_2-1/(\cdots-1/a_k))$.
Z indukcji otrzymujemy $[\underbrace{2,\ldots,2}_{m}]^-
=(m+1)/m$ dla $m\geq1$: dla $m=1$ jest to $2$,
a dodanie kolejnego wyrazu $2$ z lewej daje
$2-m/(m+1)=(m+2)/(m+1)$.
Gdy $p\geq3$, ogon ułamka ma $p-2$ dwójki,
a więc wartość $(p-1)/(p-2)$. Otrzymujemy
\[
[p+2,2,\ldots,2]^-
=p+2-\frac{p-2}{p-1}
=\frac{(p+2)(p-1)-(p-2)}{p-1}
=\frac{p^2}{p-1}.
\]
Dla $p=2$ lewa strona to $[4]^-=4$.
\end{proof}

\begin{twierdzenie}[Identyfikacja brzegu]
\label{thm:tg-rb-lens}
Przy konwencji, w której $L(n,q)$ oznacza wynik
$n/q$-chirurgii na dodatnio zorientowanym węźle
trywialnym w $S^3$, brzeg zorientowanej
czterorozmaitości $C_p$ jest orientacyjnie
dyfeomorficzny z $L(p^2,1-p)$.
\end{twierdzenie}

\begin{proof}
Przedstawienie uchwytowe plumbingowania $C_p$
jest liniowym ramowanym splotem Hopfa o ramowaniach
$-(p+2),-2,\ldots,-2$.
Kolejne usuwanie końcowych składowych przez
regułę \emph{slam-dunk} w rachunku chirurgii
zamienia ten splot na pojedynczy węzeł trywialny
o \emph{wymiernym} współczynniku chirurgii
\[
-[p+2,\underbrace{2,\ldots,2}_{p-2}]^-
=-\frac{p^2}{p-1}=\frac{p^2}{1-p}.
\]
Z poprzedniego lematu wyprowadziliśmy wartość
ułamka; przejście od łańcucha do współczynnika
jest lokalnym twierdzeniem rachunku chirurgii
(\emph{slam-dunk}, lemat \ref{lem:tg-kirby-slam-dunk}),
udowadnianym przez śledzenie
południków i długości na kolejnych torusach brzegowych.
Nie należy zastępować tego kroku samym rachunkiem
$\det Q_p$: wyznacznik nie wyznacza typu soczewkowej
trójrozmaitości. Otrzymany współczynnik daje dokładnie
$L(p^2,1-p)$ przy konwencji ustalonej w tezie.
\end{proof}

\begin{przyklad}[Przypadki $p=2,3,4$]
Dla $p=2$ graf składa się z jednego wierzchołka
$-4$, zatem $\partial C_2=L(4,-1)$.
Dla $p=3$ mamy wagi $-5,-2$ i ułamek
$[5,2]^-=5-1/2=9/2$, więc
$\partial C_3=L(9,-2)$.
Dla $p=4$ z wag $-6,-2,-2$ otrzymujemy
$[6,2,2]^-=6-2/3=16/3$ oraz
$\partial C_4=L(16,-3)$.
W każdym przypadku rząd $H_1$ jest kwadratem $p^2$,
ale parametr drugiego argumentu $L(p^2,1-p)$
wymaga pełnej informacji o sklejeniu.
\end{przyklad}

\begin{center}
\begin{tikzpicture}[>=Stealth,scale=.94]
\begin{scope}[shift={(-3.25,0)}]
  \node[font=\small\bfseries,text=tgNavy] at (0,1.10) {$C_2$};
  \node[tgRedVertex] at (0,0) {$-4$};
  \node[font=\scriptsize,text=tgNavy] at (0,-.75) {$L(4,-1)$};
\end{scope}
\begin{scope}[shift={(0,0)}]
  \node[font=\small\bfseries,text=tgNavy] at (0,1.10) {$C_3$};
  \node[tgRedVertex] (c) at (-.8,0) {$-5$};
  \node[tgBlueVertex] (d) at (.8,0) {$-2$};
  \draw[tgTree] (c)--(d);
  \node[font=\scriptsize,text=tgNavy] at (0,-.75) {$L(9,-2)$};
\end{scope}
\begin{scope}[shift={(3.6,0)}]
  \node[font=\small\bfseries,text=tgNavy] at (0,1.10) {$C_4$};
  \node[tgRedVertex] (a) at (-1.25,0) {$-6$};
  \node[tgBlueVertex] (b) at (0,0) {$-2$};
  \node[tgBlueVertex] (c) at (1.25,0) {$-2$};
  \draw[tgTree] (a)--(b)--(c);
  \node[font=\scriptsize,text=tgNavy] at (0,-.75) {$L(16,-3)$};
\end{scope}
\end{tikzpicture}
\end{center}

\section{Wymierna kula homologii Cassona--Harera}
\label{sec:tg-rb-rational-ball}

\begin{definicja}[Wymierna kula homologii]
Zwartą, spójną, zorientowaną czterorozmaitość
$B$ z niepustym brzegiem nazywamy
\emph{wymierną kulą homologii}, jeżeli
$H_*(B;\mathbb Q)\cong H_*(B^4;\mathbb Q)$,
tzn. $H_0(B;\mathbb Q)\cong\mathbb Q$ oraz
$H_k(B;\mathbb Q)=0$ dla każdego $k>0$.
Nie wymagamy, aby $B$ była ściągalna ani
aby $H_1(B;\mathbb Z)$ było zerowe.
\end{definicja}

\begin{twierdzenie}[Casson--Harer; istnienie wypełnienia]
\label{thm:tg-rb-casson-harer}
Dla każdej liczby całkowitej $p\geq2$ istnieje
zwarta, gładka, zorientowana czterorozmaitość $B_p$
taka, że
\[
\partial B_p\cong L(p^2,1-p),\qquad
H_1(B_p;\mathbb Z)\cong\mathbb Z/p\mathbb Z,
\quad H_k(B_p;\mathbb Q)=0\ (k>0).
\]
Można ją wybrać z rozkładem na jeden uchwyt
indeksu zero, jeden uchwyt indeksu jeden i jeden
uchwyt indeksu dwa; dla tej konstrukcji
$\pi_1(B_p)\cong\mathbb Z/p\mathbb Z$, a
$\pi_1(\partial B_p)\to\pi_1(B_p)$ jest surjekcją.
\end{twierdzenie}

\begin{proof}[Zakres dowodu i rachunek homologii]
Istnienie takiego \emph{konkretnego rozkładu na uchwyty},
a przede wszystkim identyfikacja jego brzegu jako
$L(p^2,1-p)$, stanowi nietrywialną treść konstrukcji
Cassona--Harera. Przyjmujemy ten krok z literatury.
Wyjaśnijmy natomiast szczegółowo konsekwencje
jego opisu uchwytowego.

W odpowiadającym rozkładowi kompleksie komórkowym
mamy jedną komórkę stopnia zero, jedną stopnia jeden
i jedną stopnia dwa. Krzywa przyczepienia dwukomórki
przebiega algebraicznie $p$ razy przez jedyny
uchwyt indeksu jeden, czyli w homotopii reprezentuje $x^p$
w grupie wolnej na jednym generatorze $x$.
Komórkowy kompleks łańcuchowy w dodatnich stopniach
ma zatem postać
\[
0\longrightarrow\mathbb Z
\xrightarrow{\ \cdot p\ }\mathbb Z
\longrightarrow\mathbb Z\longrightarrow0.
\]
Wynika stąd $H_2(B_p;\mathbb Z)=0$,
$H_1(B_p;\mathbb Z)\cong\mathbb Z/p\mathbb Z$,
a po tensorowaniu z $\mathbb Q$ również
$H_k(B_p;\mathbb Q)=0$ dla $k>0$.
Ponadto van Kampen daje
$\pi_1(B_p)\cong\langle x\mid x^p=1\rangle$
i stąd grupę cykliczną rzędu $p$.
Surjektywność homomorfizmu brzegowego jest
oddzielną własnością tego modelu uchwytowego;
nie wynika wyłącznie z obliczenia homologii.
\end{proof}

\begin{uwaga}[Dlaczego $B_p$ nie jest zwykłą kulą]
Dla $p\geq2$ grupa $H_1(B_p;\mathbb Z)$
nie jest zerowa. Zatem $B_p$ ma homologię kuli
wyłącznie ze współczynnikami wymiernymi;
nie jest całkowitą kulą homologii ani
nie może być homeomorficzna z $B^4$.
\end{uwaga}

\section{Forma splecenia i podgrupa metabolizująca}
\label{sec:tg-rb-linkingform}

\begin{stwierdzenie}[Forma splecenia brzegu]
\label{prop:tg-rb-linkingform}
Niech $p\geq2$ i niech $e_1$ będzie generatorem
$H_1(Y_p;\mathbb Z)\cong\mathbb Z/p^2\mathbb Z$
wskazanym w dowodzie twierdzenia
\ref{thm:tg-rb-h1boundary}. Przy konwencji
$\lambda_{Y_p}$ reprezentowanej przez $-Q_p^{-1}$
zachodzi
\[
\lambda_{Y_p}(e_1,e_1)=\frac{p-1}{p^2}
\quad\text{w }\mathbb Q/\mathbb Z.
\]
Podgrupa $p\langle e_1\rangle$ rzędu $p$ jest
izotropowa, a nawet jest równa swemu ortogonalnemu
do niej dopełnieniu.
\end{stwierdzenie}

\begin{proof}
Wzór na element odwrotności macierzy, wyrażony
przez minor, daje
$(Q_p^{-1})_{11}=\det Q_p^{(1)}/\det Q_p$,
gdzie $Q_p^{(1)}$ powstaje przez skreślenie pierwszego
wiersza i kolumny. Dla $p\geq3$ jest to
macierz łańcucha $p-2$ wierzchołków o wadze $-2$,
więc $\det Q_p^{(1)}=(-1)^{p-2}(p-1)$.
Dla $p=2$ używamy konwencji, że wyznacznik
macierzy pustej wynosi jeden. W obu przypadkach
\[
-(Q_p^{-1})_{11}=-\frac{(-1)^{p-2}(p-1)}{(-1)^{p-1}p^2}
=\frac{p-1}{p^2}.
\]
Następnie
$\lambda_{Y_p}(pe_1,pe_1)=p^2(p-1)/p^2=p-1=0$
w $\mathbb Q/\mathbb Z$. Cała podgrupa
$p\langle e_1\rangle$ jest więc izotropowa.
Forma splecenia jest niezdegenerowana na grupie
rzędu $p^2$, dlatego dopełnienie ortogonalne
podgrupy rzędu $p$ również ma rząd $p$.
Skoro podgrupa jest zawarta w swoim dopełnieniu,
to oba zbiory są równe.
\end{proof}

\begin{uwaga}[Związek z wypełnieniem]
Surjekcja $H_1(Y_p;\mathbb Z)\to H_1(B_p;\mathbb Z)$
ma jądro rzędu $p$; jest to właśnie
$p\langle e_1\rangle$.
Taki układ jest zgodny z ogólną zasadą,
że jądro mapy brzegowej dla wymiernej kuli homologii
jest metabolizerem formy splecenia.
W szczególności istnienie $B_p$ nakłada ograniczenie
na brzeg znacznie silniejsze niż samo $|H_1(Y_p)|=p^2$.
\end{uwaga}

\section{Definicja operacji i zmiana niezmienników}
\label{sec:tg-rb-definition}

\begin{definicja}[Wymierna kontrakcja]
\label{def:tg-rb-operation}
Niech $p\geq2$, a $X$ niech będzie zamkniętą,
spójną, zorientowaną, gładką czterorozmaitością.
Załóżmy, że $X$ zawiera w swoim wnętrzu
gładko zanurzoną kopię $C_p$ jako podrozmaitość
z brzegiem i z ustaloną identyfikacją
regularnego otoczenia. Oznaczmy
$W:=X\setminus\operatorname{int}C_p$ i
$Y_p:=\partial C_p$.
\emph{Wymierną kontrakcją $X$ wzdłuż $C_p$}
nazywamy rozmaitość
\[
X_{(p)}:=W\cup_{\varphi}B_p,
\]
gdzie $\varphi:\partial B_p\to\partial C_p$ jest
dyfeomorfizmem \emph{zachowującym} orientację
brzegową $\partial C_p$ i $\partial B_p$.
Ponieważ $\partial W=-\partial C_p$, to po
utożsamieniu $\partial W$ z $\partial C_p$
to samo sklejenie jest odwracające orientację
między brzegami $W$ i $B_p$.
W klasycznym modelu $B_p$ typ otrzymanej
gładkiej rozmaitości jest niezależny od wyboru
takiej identyfikacji; korzystamy tu z twierdzenia
o przedłużaniu samodyfeomorfizmów $\partial B_p$
do $B_p$, przyjmowanego z literatury.
\end{definicja}

\begin{center}
\begin{tikzpicture}[>=Stealth,scale=.92]
  \node[draw=tgNavy,fill=tgNavy!6,rounded corners=3pt,
    minimum width=23mm,minimum height=13mm,align=center,
    font=\footnotesize] (W) at (0,1.7) {$W$\\{\scriptsize wspólne dopełnienie}};
  \node[draw=tgCoral,fill=tgCoral!14,rounded corners=3pt,
    minimum width=27mm,minimum height=13mm,align=center,
    font=\footnotesize] (C) at (-3.0,0) {$C_p$\\{\scriptsize plumbing sfer}};
  \node[draw=tgTeal,fill=tgTeal!14,rounded corners=3pt,
    minimum width=27mm,minimum height=13mm,align=center,
    font=\footnotesize] (B) at (3.0,0) {$B_p$\\{\scriptsize wymierna kula}};
  \draw[tgAccent,-{Stealth[length=3mm]}] (W.south west)--(C.north east);
  \draw[tgTree,-{Stealth[length=3mm]}] (W.south east)--(B.north west);
  \node[font=\footnotesize,text=tgCoral] at (-3.0,-1.05) {$X=W\cup_{Y_p}C_p$};
  \node[font=\footnotesize,text=tgTeal] at (3.0,-1.05) {$X_{(p)}=W\cup_{Y_p}B_p$};
  \node[font=\scriptsize,text=tgNavy] at (0,-1.65)
    {Ten sam brzeg $Y_p\cong L(p^2,1-p)$, różne wypełnienia.};
\end{tikzpicture}
\end{center}

\begin{twierdzenie}[Skutek liczbowy wymiernej kontrakcji]
\label{thm:tg-rb-invariants}
Niech $p\geq2$ i niech $X$, $X_{(p)}$ oraz
$C_p\subset X$ będą jak w definicji
\ref{def:tg-rb-operation}.
Niech $e$ oznacza charakterystykę Eulera,
$\sigma$ sygnaturę formy przecięcia, a
$b_2^\pm$ wymiary maksymalnych dodatnio i
ujemnie określonych podprzestrzeni
$H_2(-;\mathbb R)$.
Wtedy
\[
\begin{aligned}
e(X_{(p)})&=e(X)-(p-1), &
\sigma(X_{(p)})&=\sigma(X)+(p-1),\\
b_2^+(X_{(p)})&=b_2^+(X), &
b_2^-(X_{(p)})&=b_2^-(X)-(p-1).
\end{aligned}
\]
W szczególności, jeśli wprowadzimy oznaczenia
$\chi_h(Z):=(e(Z)+\sigma(Z))/4$ oraz
$c_{\mathrm{num}}(Z):=2e(Z)+3\sigma(Z)$,
to $\chi_h(X_{(p)})=\chi_h(X)$ i
$c_{\mathrm{num}}(X_{(p)})=c_{\mathrm{num}}(X)+(p-1)$.
Dla zamkniętych czterorozmaitości prawie zespolonych,
orientowanych przez strukturę prawie zespoloną,
ostatnia wielkość jest liczbą Chernowską $c_1^2$.
Nie zakładamy tu, że każda rozmaitość występująca
w twierdzeniu jest prawie zespolona. W takim przypadku
$\chi_h$ jest tylko oznaczeniem liczbowym, a nie holomorficzną
charakterystyką Eulera.
\end{twierdzenie}

\begin{proof}
Brzeg $Y_p$ jest wymierną sferą homologiczną,
więc $H_1(Y_p;\mathbb Q)=H_2(Y_p;\mathbb Q)=0$.
Rozpatrzmy rozkłady $X=W\cup C_p$ oraz
$X_{(p)}=W\cup B_p$. Ciągi Mayera--Vietorisa
nad $\mathbb Q$ dają izomorfizmy
\[
H_2(X;\mathbb Q)\cong H_2(W;\mathbb Q)
\oplus H_2(C_p;\mathbb Q),\qquad
H_2(X_{(p)};\mathbb Q)\cong H_2(W;\mathbb Q).
\]
Reprezentantów klas z wnętrz obu części można
wybrać rozłącznie, toteż pierwsza suma jest
ortogonalna względem formy przecięcia.
Na $H_2(C_p;\mathbb Q)$ forma jest ujemnie określona
rangi $p-1$, jak wykazaliśmy w lemacie
\ref{lem:tg-rb-minors}. Otrzymujemy stąd równości
dla $b_2^\pm$ oraz dla sygnatury.

Wzór na charakterystykę Eulera korzysta z addytywności
$e(A\cup_YB)=e(A)+e(B)-e(Y)$ dla zwartych
kompleksów triangulowalnych. Ponieważ
$e(Y_p)=0$, $e(C_p)=1+(p-1)=p$ oraz
$e(B_p)=1$, mamy
$e(X_{(p)})-e(X)=e(B_p)-e(C_p)=1-p$.
Równania dla $\chi_h$ i $c_{\mathrm{num}}$ są
następnie bezpośrednim podstawieniem.
\end{proof}

\begin{przyklad}[Kontrola dla $p=2$ i $p=3$]
Dla $p=2$ usuwamy jedną sferę o samoprzecięciu
$-4$. Ponieważ $e(C_2)=2$, $\sigma(C_2)=-1$,
a $e(B_2)=1$, $\sigma(B_2)=0$, otrzymujemy
$\Delta e=-1$ i $\Delta\sigma=1$.
Dla $p=3$ usuwamy łańcuch sfer o wagach $-5,-2$:
$e(C_3)=3$, $\sigma(C_3)=-2$, zatem
$\Delta e=-2$ i $\Delta\sigma=2$.
W obu przypadkach wzrost $2e+3\sigma$ wynosi
odpowiednio $1$ i $2$.
\end{przyklad}

\section{Grupa fundamentalna po wymiernej kontrakcji}
\label{sec:tg-rb-pi1}

\begin{stwierdzenie}[Kryterium prostej spójności]
\label{prop:tg-rb-pi1}
Dla $p\geq2$, przy założeniach definicji
\ref{def:tg-rb-operation},
jeżeli $W:=X\setminus\operatorname{int}C_p$ jest
po prostu spójna, to $X_{(p)}$ jest po prostu spójna.
\end{stwierdzenie}

\begin{proof}
Wybierzmy punkt bazowy na wspólnym brzegu $Y_p$.
Twierdzenie van Kampena daje prezentację
\[
\pi_1(X_{(p)})\cong\pi_1(W)*_{\pi_1(Y_p)}\pi_1(B_p).
\]
Jeśli $\pi_1(W)=1$, pozostaje iloraz
$\pi_1(B_p)$ przez normalne domknięcie obrazu
$\pi_1(Y_p)\to\pi_1(B_p)$. Na mocy twierdzenia
Cassona--Harera obraz ten jest całą grupą
$\pi_1(B_p)\cong\mathbb Z/p\mathbb Z$,
więc iloraz jest trywialny.
\end{proof}

\begin{uwaga}[Dlaczego $\pi_1(X)=1$ nie wystarcza w dowodzie]
Z prostej spójności $X$ nie wynika automatycznie,
że $W=X\setminus\operatorname{int}C_p$ jest po prostu
spójna. Grupy przy sklejaniu należy obliczyć oddzielnie.
W konstrukcjach egzotycznych czterorozmaitości
jest to jeden z istotnych etapów dowodu, którego
nie wolno zastąpić rachunkiem samej homologii.
\end{uwaga}

\section{Zastosowanie do egzotycznych struktur gładkich}
\label{sec:tg-rb-exotic}

\begin{definicja}[Egzotyczna struktura gładka]
Dwie zamknięte, spójne, zorientowane, gładkie
czterorozmaitości $X_1,X_2$ nazywamy
\emph{egzotyczną parą}, jeżeli są homeomorficzne
z zachowaniem orientacji, ale nie są dyfeomorficzne.
\end{definicja}

Sama wymierna kontrakcja nie dowodzi egzotyczności.
Trzeba wykonać trzy niezależne czynności:
obliczyć $\pi_1$ nowej rozmaitości, ustalić
jej formę przecięcia oraz znaleźć niezmiennik
\emph{gładki}, który odróżnia wynik od znanego modelu.
Twierdzenie Freedmana służy do klasyfikacji
topologicznej po spełnieniu właściwych hipotez;
niezmienniki Seiberga--Wittena lub Donaldsona
służą do rozróżniania struktur gładkich.

\begin{przyklad}[Dane modelu $\mathbb{CP}^2\#7\overline{\mathbb{CP}}^{\,2}$]
Niech $Z=\mathbb{CP}^2\#7\overline{\mathbb{CP}}^{\,2}$
z orientacją standardową na pierwszym składniku.
Z rozkładu sumy spójnej mamy
$H_2(Z;\mathbb Z)\cong\mathbb Z^8$ oraz
\[
Q_Z=\langle1\rangle\oplus7\langle-1\rangle,
\qquad b_2^+(Z)=1,\quad b_2^-(Z)=7.
\]
Ponieważ $Z$ jest po prostu spójna,
$e(Z)=2+b_2(Z)=10$, a
$\sigma(Z)=1-7=-6$. W szczególności
$2e(Z)+3\sigma(Z)=20-18=2$ i
$\chi_h(Z)=(10-6)/4=1$.

Te liczby są \emph{konieczne}, lecz nie wystarczają,
aby dowolna rozmaitość była homeomorficzna z $Z$.
Wymagana jest m.in. prosta spójność i ustalenie
całkowitej formy przecięcia wraz z jej parzystością.
Przy tych założeniach twierdzenie Freedmana
może dać klasyfikację topologiczną, lecz nie
rozstrzyga problemu dyfeomorfizmu.
\end{przyklad}

\begin{twierdzenie}[Park; zastosowanie wymiernej kontrakcji,
przyjmowane z literatury]
Istnieje zamknięta, po prostu spójna, gładka,
zorientowana czterorozmaitość $X_{\mathrm{Park}}$
homeomorficzna z
$\mathbb{CP}^2\#7\overline{\mathbb{CP}}^{\,2}$,
lecz z nią niedyfeomorficzna. W konstrukcji
wykorzystuje się m.in. wymierną kontrakcję
odpowiedniego plumbingowania oraz niezmienniki
Seiberga--Wittena.
\end{twierdzenie}

\begin{uwaga}[Co w tym przykładzie pozostaje bez dowodu]
Powyższe twierdzenie jest głębokim wynikiem J. Parka.
Nasz rachunek $e$, $\sigma$ i $Q$ \emph{nie} jest
jego dowodem. Pełna konstrukcja wymaga znalezienia
odpowiedniego zanurzenia konfiguracji w wyjściowej
rozmaitości, sprawdzenia grupy fundamentalnej
dopełnienia oraz obliczenia gładkich niezmienników.
Ponadto dla $b_2^+=1$ niezmienniki
Seiberga--Wittena zależą od wyboru komory;
porównanie z rozmaitościami wymiernymi wymaga
konsekwentnego ustalenia tej komory.
\end{uwaga}

\begin{twierdzenie}[Symington; wersja symplektyczna, bez dowodu]
Jeżeli gładka rozmaitość symplektyczna zawiera
konfigurację $C_p$ zanurzoną symplektycznie,
ze standardowym modelem otoczenia i warunkami
zgodności struktury kontaktowej na brzegu,
to można wykonać odpowiednią wymierną kontrakcję
w kategorii symplektycznej: otrzymana czterorozmaitość
posiada formę
symplektyczną zgodną z formą na dopełnieniu.
\end{twierdzenie}

\section{Przejście do uogólnień}
\label{sec:tg-rb-general}
Gdy $p>q\geq1$ są względnie pierwsze, analogiczne
plumbingowanie $C_{p,q}$ wyznacza ujemne rozwinięcie
$\frac{p^2}{pq-1}$. Rachunek grupy brzegowej
$H_1(\partial C_{p,q};\mathbb Z)\cong\mathbb Z/p^2\mathbb Z$
nie wystarcza do identyfikacji jej orientowanego typu.
W następnym rozdziale wyprowadzimy ujemne ułamki łańcuchowe,
pełny rachunek formy przecięcia i rzeczywisty współczynnik
chirurgii, a następnie zastosujemy je do wymiernych kontrakcji.
Aby sprawdzić, kiedy otrzymujemy nową strukturę gładką,
potrzebujemy jeszcze niezmienników Seiberga--Wittena.

% Iteracja 10: uogolniona wymierna kontrakcja i podstawy niezmiennikow SW.
\begin{uwaga}[Granica rachunku wyznacznika]
Sama tożsamość $|\det Q_{C_p}|=p^2$ wyznacza rząd
grupy $H_1(\partial C_p;\mathbb Z)$, ale nie rozstrzyga
ani izomorfizmu tej grupy z $\mathbb Z/p^2\mathbb Z$,
ani typu dyfeomorfizmu brzegu. Cykliczność trzeba
wyprowadzić z prezentacji relacji lub jej postaci Smitha,
a identyfikację przestrzeni soczewkowej z odpowiedniego
ułamka łańcuchowego. Istnienie zastępującej ją wymiernej
kuli homologii jest kolejnym, niezależnym wynikiem
geometrycznym (Casson--Harer).
\end{uwaga}

\chapter{Uogólniona wymierna kontrakcja i łańcuchy \texorpdfstring{$C_{p,q}$}{Cpq}}
\label{ch:tg-rb-general}

W poprzednim rozdziale rozpatrzyliśmy konfigurację $C_p$ opisaną przez
łańcuch sfer o wagach $-(p+2),-2,\ldots,-2$. W niniejszym
rozdziale wyprowadzimy uogólnienie $C_{p,q}$, gdzie liczby $p>q\geq1$
są względnie pierwsze. Podstawowym narzędziem będzie \emph{ujemne
rozwinięcie ułamka łańcuchowego}. Odróżniamy przy tym rachunki
algebraiczne, które wykonamy w całości, od geometrycznego twierdzenia
o istnieniu wymiernej kuli homologii.

\section{Ujemne ułamki łańcuchowe: definicja i jednoznaczność}
\label{sec:tg-general-cf}

\begin{definicja}
Dla liczb całkowitych $a_1,\ldots,a_r$, gdzie $r\geq1$ oraz
$a_i\geq2$, definiujemy \emph{ujemny ułamek łańcuchowy}
rekurencyjnie przez $[a_r]^-:=a_r$ oraz
$[a_i,\ldots,a_r]^-:=a_i-1/[a_{i+1},\ldots,a_r]^-$.
\end{definicja}

Nie jest to zwykłe dodatnie rozwinięcie Euklidesa. W szczególności
w każdym kroku odejmujemy \emph{odwrotność} kolejnego ogona.

\begin{twierdzenie}[Algorytm ujemnego rozwinięcia]
\label{thm:tg-cf-algorithm}
Niech $A,B\in\mathbb N$ spełniają $A>B\geq1$ i
$\gcd(A,B)=1$. Istnieje dokładnie jeden skończony ciąg
$a_1,\ldots,a_r\geq2$ taki, że $A/B=[a_1,\ldots,a_r]^-$. 
Otrzymujemy go, biorąc $a_1=\lceil A/B\rceil$, a jeśli
$B>1$, zastępując $A/B$ ułamkiem $B/(a_1B-A)$ i powtarzając
postępowanie. Proces kończy się, gdy mianownik wynosi jeden.
\end{twierdzenie}

\begin{proof}
Jeżeli $B=1$, jedynym rozwinięciem jest $[A]^-$: każde
rozwinięcie o co najmniej dwóch wyrazach ma wartość niecałkowitą,
ponieważ jego pierwszy ogon jest większy od jedności. Niech
$B>1$. Ponieważ $A/B$ jest nieskracalny, nie jest całkowity.
Definiujemy $a_1:=\lceil A/B\rceil\geq2$ oraz
$B':=a_1B-A$. Wówczas $0<B'<B$ i $\gcd(B,B')=1$.
Zatem
$A/B=a_1-B'/B=a_1-1/(B/B')$.
Dalszy ułamek ma mianownik $B'$ mniejszy od $B$, toteż
indukcja kończy konstrukcję.

Dla dowodu jednoznaczności zauważmy, że każdy ogon
$[a_i,\ldots,a_r]^->$ jest większy od jedności. Stąd
$a_1-1<[a_1,\ldots,a_r]^-\leq a_1$, przy czym równość
po prawej zachodzi wyłącznie dla $r=1$.
Pierwszy wyraz jest więc wymuszony jako sufit danego ułamka;
po jego usunięciu wymuszony jest także następny ogon.
Indukcja daje jednoznaczność wszystkich wyrazów.
\end{proof}

\begin{przyklad}[Trzy kolejne obliczenia]
Dla $(p,q)=(5,1)$ zachodzi $25/4=[7,2,2,2]^-$.
Dla $(p,q)=(5,2)$ obliczamy
$25/9=3-2/9=3-1/(9/2)=3-1/(5-1/2)$, czyli
$25/9=[3,5,2]^-$. Dla $(p,q)=(5,3)$ mamy
$25/14=2-3/14=2-1/(14/3)=2-1/(5-1/3)$,
zatem $25/14=[2,5,3]^-$. Kolejność wag ma znaczenie dla
konwencji orientacji i oznaczania generatorów brzegu.
\end{przyklad}

\section{Definicja \texorpdfstring{$C_{p,q}$}{Cpq} i forma przecięcia}
\label{sec:tg-general-config}

\begin{definicja}[Konfiguracja $C_{p,q}$]
Niech $p,q\in\mathbb N$, $p>q\geq1$, $\gcd(p,q)=1$, i niech
$p^2/(pq-1)=[a_1,\ldots,a_r]^-$, $a_i\geq2$.
Przez $C_{p,q}$ oznaczamy zwartą, spójną, zorientowaną,
gładką czterorozmaitość otrzymaną przez liniowe plumbingowanie
wiązek dyskowych nad $S^2$ o liczbach Eulera
$-a_1,\ldots,-a_r$, przy czym każda para sąsiednich sfer
przecina się poprzecznie i dodatnio w jednym punkcie.
Forma przecięcia $Q_{p,q}$ w bazie klas tych sfer ma macierz
\[
Q_{p,q}=\begin{pmatrix}
-a_1&1&0&\cdots&0\\
1&-a_2&1&\ddots&\vdots\\
0&1&\ddots&\ddots&0\\
\vdots&\ddots&\ddots&-a_{r-1}&1\\
0&\cdots&0&1&-a_r
\end{pmatrix}
\]
(dla $r=1$ jest to macierz $[-a_1]$).
\end{definicja}

\begin{center}
\begin{tikzpicture}[>=Stealth,scale=.9]
  \node[font=\small\bfseries,text=tgNavy] at (-3.7,.92) {$C_{5,1}$};
  \node[tgRedVertex] (a) at (-5.2,0) {$-7$};
  \node[tgBlueVertex] (b) at (-4.2,0) {$-2$};
  \node[tgBlueVertex] (c) at (-3.2,0) {$-2$};
  \node[tgBlueVertex] (d) at (-2.2,0) {$-2$};
  \draw[tgTree] (a)--(b)--(c)--(d);
  \node[font=\small\bfseries,text=tgNavy] at (2.4,.92) {$C_{5,2}$};
  \node[tgRedVertex] (u) at (1,0) {$-3$};
  \node[tgBlueVertex] (v) at (2.3,0) {$-5$};
  \node[tgBlueVertex] (w) at (3.6,0) {$-2$};
  \draw[tgTree] (u)--(v)--(w);
  \node[font=\scriptsize,text=tgNavy] at (-3.7,-.68) {$25/4=[7,2,2,2]^-$};
  \node[font=\scriptsize,text=tgNavy] at (2.35,-.68) {$25/9=[3,5,2]^-$};
\end{tikzpicture}
\end{center}

\begin{lemat}[Wyznaczniki kolejnych ogonów]
\label{lem:tg-general-minors}
Niech $a_1,\ldots,a_r\geq2$. Oznaczmy przez $D_i$ wyznacznik
macierzy trójdiagonalnej o przekątnej
$a_i,\ldots,a_r$ i elementach $-1$ bezpośrednio przy przekątnej,
a ponadto przyjmijmy $D_{r+1}=1$, $D_{r+2}=0$. Wtedy
$D_i=a_iD_{i+1}-D_{i+2}$ i
$[a_i,\ldots,a_r]^-=D_i/D_{i+1}$. Wszystkie $D_i$ dla
$1\leq i\leq r+1$ są dodatnie, zaś $D_i>D_{i+1}$ dla
$1\leq i\leq r$.
\end{lemat}

\begin{proof}
Rozwinięcie wyznacznika względem pierwszego wiersza daje
$D_i=a_iD_{i+1}-D_{i+2}$. Dla $i=r$ mamy $D_r=a_r$ i
$D_{r+1}=1$. Indukcja wsteczna prowadzi do
$D_i\geq2D_{i+1}-D_{i+2}>D_{i+1}>0$.
Dzieląc rekurencję przez $D_{i+1}$ otrzymujemy
$D_i/D_{i+1}=a_i-D_{i+2}/D_{i+1}$, co jest dokładnie
rekurencją definiującą ujemny ułamek łańcuchowy.
\end{proof}

\begin{twierdzenie}[Ujemna określoność, wyznacznik i homologia]
\label{thm:tg-general-algebra}
Dla dowolnych względnie pierwszych $p>q\geq1$ macierz
$Q_{p,q}$ jest ujemnie określona i
$\det Q_{p,q}=(-1)^r p^2$.
Ponadto jej cokernel jest cykliczny i
\[
H_1(\partial C_{p,q};\mathbb Z)
\cong\operatorname{coker}Q_{p,q}
\cong\mathbb Z/p^2\mathbb Z.
\]
W szczególności $\partial C_{p,q}$ jest wymierną sferą
homologiczną.
\end{twierdzenie}

\begin{proof}
Niech $A:=-Q_{p,q}$. Z poprzedniego lematu mamy
$\det A=D_1=p^2$, gdyż
$D_1/D_2=p^2/(pq-1)$ jest ułamkiem nieskracalnym,
a $D_1,D_2$ są całkowite i względnie pierwsze
(ten ostatni fakt wynika rekurencyjnie z
$\gcd(D_i,D_{i+1})=\gcd(D_{i+1},D_{i+2})=1$).
W szczególności $D_2=pq-1$.

Dowód określoności nie wymaga wyliczania wartości własnych.
Dla $r\geq2$ i $x=(x_1,\ldots,x_r)\in\mathbb R^r$
zachodzi tożsamość
\[
x^TAx=\sum_{i=1}^{r-1}(x_i-x_{i+1})^2
+(a_1-1)x_1^2+(a_r-1)x_r^2
+\sum_{i=2}^{r-1}(a_i-2)x_i^2.
\]
Każdy składnik po prawej jest nieujemny, a pierwszy i ostatni
składnik brzegowy wymuszają $x_1=x_r=0$ przy zerowej sumie;
różnice sąsiednich współrzędnych wymuszają wtedy $x=0$.
Zatem $A$ jest dodatnio określona. Dla $r=1$ mamy $A=[a_1]$.
Otrzymujemy ujemną określoność $Q_{p,q}$ i wzór na wyznacznik.

Cokernel $Q_{p,q}$ jest generowany przez klasy
$e_1,\ldots,e_r$ z relacjami
$-a_1e_1+e_2=0$, $e_{i-1}-a_ie_i+e_{i+1}=0$
oraz $e_{r-1}-a_re_r=0$ (dla $r=1$ tylko
$-a_1e_1=0$). Pierwsza relacja wyraża $e_2$ przez $e_1$,
a każda następna wyraża $e_{i+1}$ przez poprzednie.
W konsekwencji grupa jest generowana przez $e_1$.
Ponieważ macierz jest nieosobliwa, grupa jest skończona rzędu
$|\det Q_{p,q}|=p^2$, więc musi być cykliczna rzędu $p^2$.
Identyfikację z $H_1$ brzegu uzyskaliśmy wcześniej z
ciągu dokładnego pary i dualności Poincarégo--Lefschetza.
\end{proof}

\begin{przyklad}[Rachunek dla $C_{5,2}$]
Dla $25/9=[3,5,2]^-$ otrzymujemy
\[
Q_{5,2}=\begin{pmatrix}-3&1&0\\1&-5&1\\0&1&-2\end{pmatrix},
\qquad\det Q_{5,2}=-25.
\]
Relacje $-3e_1+e_2=0$ oraz
$e_1-5e_2+e_3=0$ dają $e_2=3e_1$ i $e_3=14e_1$.
Ostatnia relacja $e_2-2e_3=0$ staje się
$-25e_1=0$. W ten sposób bez użycia postaci Smitha
obliczamy $H_1(\partial C_{5,2};\mathbb Z)\cong\mathbb Z/25\mathbb Z$.
\end{przyklad}

\section{Identyfikacja brzegu i forma splecenia}
\label{sec:tg-general-boundary}

\begin{twierdzenie}[Brzeg konfiguracji $C_{p,q}$]
\label{thm:tg-general-lens}
Przy konwencji $L(n,s)=S^3_{n/s}(U)$ dla
niezorientowanego węzła trywialnego $U\subset S^3$,
z orientacją wyznaczoną przez chirurgię, istnieje
orientacyjny dyfeomorfizm
$\partial C_{p,q}\cong L(p^2,1-pq)$.
\end{twierdzenie}

\begin{proof}
Diagram Kirby'ego $C_{p,q}$ jest liniowym łańcuchem
ramowanych okręgów trywialnych o współczynnikach
$-a_1,\ldots,-a_r$ i liczbach splecenia sąsiednich
składowych równych $+1$.
Stosując kolejno geometryczną regułę \emph{slam-dunk}
(udowodnioną wcześniej dla usuwania sąsiedniej składowej),
redukujemy diagram do pojedynczej składowej z wymiernym
współczynnikiem
$-[a_1,\ldots,a_r]^-= -p^2/(pq-1)=p^2/(1-pq)$.
Jest to chirurgiczne przedstawienie $L(p^2,1-pq)$.
Argument nie korzysta wyłącznie z wyznacznika macierzy:
niezależnie ustala klasę krzywej południkowej podczas sklejenia.
\end{proof}

\begin{stwierdzenie}[Generator, forma splecenia i metabolizer]
\label{prop:tg-general-linking}
Niech $e_1$ oznacza klasę pierwszego wierzchołka w
$H_1(\partial C_{p,q};\mathbb Z)=\operatorname{coker}Q_{p,q}$.
Przy konwencji, w której formę splecenia brzegu reprezentuje
$-Q_{p,q}^{-1}$, zachodzi
\[
\lambda(e_1,e_1)=\frac{pq-1}{p^2}\pmod{\mathbb Z}.
\]
Jedyna podgrupa rzędu $p$, czyli $p\langle e_1\rangle$,
jest równa swojemu dopełnieniu ortogonalnemu względem
$\lambda$. Jest więc \emph{metabolizerem} tej formy.
\end{stwierdzenie}

\begin{proof}
Wyznacznik macierzy powstałej po usunięciu pierwszego
wiersza i kolumny $Q_{p,q}$ jest równy
$(-1)^{r-1}D_2$, więc z reguły dopełnień algebraicznych
\[
-(Q_{p,q}^{-1})_{11}=
-\frac{(-1)^{r-1}D_2}{(-1)^rD_1}
=\frac{D_2}{D_1}=
\frac{pq-1}{p^2}.
\]
Dla $r=1$ przyjmujemy wyznacznik macierzy pustej równy jeden.
Ponieważ $e_1$ ma rząd $p^2$ i $\gcd(p,pq-1)=1$,
forma jest niezdegenerowana. Dla generatora $pe_1$ podgrupy
rzędu $p$ mamy
$\lambda(pe_1,pe_1)=pq-1\equiv0\pmod{\mathbb Z}$.
Cała podgrupa jest zatem izotropowa.
Jej dopełnienie ortogonalne ma rząd $p^2/p=p$,
więc z zawierania i zgodności rzędów otrzymujemy równość.
\end{proof}

\begin{uwaga}[Konwencje znaku]
W literaturze często definiuje się $L(n,s)$ przez
\emph{dodatni} ułamek chirurgii $n/s$, a niektóre teksty
przyjmują przeciwną orientację brzegu plumbingowania.
Wówczas można spotkać $L(p^2,pq-1)$ zamiast
$L(p^2,1-pq)$. Przy korzystaniu z takich wyników
należy jawnie odwrócić orientację odpowiedniego brzegu,
a nie tylko zmienić znak w liczniku lub mianowniku.
\end{uwaga}

\section{Wymierna kula homologii i zmiana niezmienników}
\label{sec:tg-general-operation}

\begin{twierdzenie}[Istnienie $B_{p,q}$; wynik przyjmowany]
\label{thm:tg-general-ball}
Dla każdego $p>q\geq1$ ze $\gcd(p,q)=1$ istnieje
zwarta, spójna, zorientowana, gładka czterorozmaitość
$B_{p,q}$ o brzegu orientacyjnie dyfeomorficznym z
$\partial C_{p,q}$, dla której
\[
H_*(B_{p,q};\mathbb Q)\cong H_*(B^4;\mathbb Q),
\qquad H_1(B_{p,q};\mathbb Z)\cong\mathbb Z/p\mathbb Z.
\]
Można wybrać standardowy model, w którym
$\pi_1(B_{p,q})\cong\mathbb Z/p\mathbb Z$, a mapa
$\pi_1(\partial B_{p,q})\to\pi_1(B_{p,q})$ jest surjektywna.
\end{twierdzenie}

\begin{uwaga}
Istnienie $B_{p,q}$ z tym brzegiem jest nietrywialnym
twierdzeniem geometrycznym, wynikającym z konstrukcji
wymiernych kul homologii (Casson--Harer i ich uogólnienia).
Powyższych faktów nie da się wyprowadzić z samej równości
$|H_1(\partial C_{p,q})|=p^2$ ani z istnienia metabolizera.
Traktujemy je tu jako wynik zewnętrzny, wskazując jego dokładną
rolę w dalszym rozumowaniu.
\end{uwaga}

\begin{definicja}[Uogólniona wymierna kontrakcja]
Niech $X$ będzie zamkniętą, spójną, zorientowaną,
gładką czterorozmaitością zawierającą gładko zanurzoną
kopię $C_{p,q}$ jako podrozmaitość z brzegiem.
Kładziemy $W:=X\setminus\operatorname{int}C_{p,q}$.
Dla ustalonego dyfeomorfizmu
$\varphi:\partial B_{p,q}\to\partial C_{p,q}$ zachowującego
orientację tych dwóch brzegów definiujemy
$X_{p,q}:=W\cup_{\varphi}B_{p,q}$.
Ponieważ $\partial W=-\partial C_{p,q}$, jest to
sklejenie odwracające orientację między $\partial W$
a $\partial B_{p,q}$. Wynik nazywamy
\emph{uogólnioną wymierną kontrakcją} $X$ wzdłuż $C_{p,q}$.
\end{definicja}

\begin{twierdzenie}[Obliczenie niezmienników]
\label{thm:tg-general-numeric}
Niech $r$ będzie długością ujemnego rozwinięcia
$p^2/(pq-1)$, a $X_{p,q}$ powstanie z zamkniętej,
zorientowanej, gładkiej czterorozmaitości $X$
przez uogólnioną wymierną kontrakcję wzdłuż $C_{p,q}$.
Wówczas
\[
\begin{aligned}
e(X_{p,q})&=e(X)-r,&
\sigma(X_{p,q})&=\sigma(X)+r,\\
b_2^+(X_{p,q})&=b_2^+(X),&
b_2^-(X_{p,q})&=b_2^-(X)-r.
\end{aligned}
\]
Stąd $\chi_h:=(e+\sigma)/4$ pozostaje bez zmian,
natomiast $c_{\mathrm{num}}:=2e+3\sigma$ wzrasta o $r$.
Przez $\chi_h$ rozumiemy tu jedynie liczbowe wyrażenie
$(e+\sigma)/4$; nie zakładamy istnienia struktury
zespolonej, która uzasadniałaby interpretację holomorficzną.
\end{twierdzenie}

\begin{proof}
Brzeg $Y=\partial C_{p,q}$ jest wymierną sferą homologiczną.
Ciągi Mayera--Vietorisa nad $\mathbb Q$ dają rozkłady
$H_2(X;\mathbb Q)\cong H_2(W;\mathbb Q)\oplus
H_2(C_{p,q};\mathbb Q)$ oraz
$H_2(X_{p,q};\mathbb Q)\cong H_2(W;\mathbb Q)$.
Sumę pierwszą można przedstawić przez powierzchnie o nośnikach
w rozłącznych częściach, zatem jest ona ortogonalna względem
formy przecięcia. Ponieważ $Q_{p,q}$ jest ujemnie określona
rangi $r$, usunięcie konfiguracji zmniejsza $b_2^-$ o $r$
i nie zmienia $b_2^+$. Wynika stąd zmiana sygnatury o $r$.

Z rozkładu uchwytowego plumbingowania, który ma jedną
komórkę wymiaru zero i $r$ dwukomórek, otrzymujemy
$e(C_{p,q})=1+r$. Natomiast $e(B_{p,q})=1$
i $e(Y)=0$. Addytywność charakterystyki Eulera prowadzi do
$e(X_{p,q})-e(X)=1-(1+r)=-r$.
Pozostałe wzory są prostym podstawieniem.
\end{proof}

\begin{wniosek}[Prosta spójność po sklejeniu]
Jeżeli w sytuacji poprzedniego twierdzenia dopełnienie
$W=X\setminus\operatorname{int}C_{p,q}$ jest po prostu spójne,
a model $B_{p,q}$ wybieramy jak w twierdzeniu
\ref{thm:tg-general-ball}, to $X_{p,q}$ również jest po prostu
spójna.
\end{wniosek}

\begin{proof}
Z twierdzenia van Kampena otrzymujemy
$\pi_1(X_{p,q})\cong\pi_1(W)*_{\pi_1(Y)}\pi_1(B_{p,q})$.
Gdy $\pi_1(W)=1$, pierwsza grupa znika, a iloraz
$\pi_1(B_{p,q})$ przez normalne domknięcie obrazu
$\pi_1(Y)$ jest trywialny, ponieważ obraz jest całą grupą.
\end{proof}

\begin{przyklad}[Co wnosi parametr $q$?]
Dla $(p,q)=(5,1)$ mamy łańcuch długości $r=4$ o wagach
$-7,-2,-2,-2$. Dla $(p,q)=(5,2)$ łańcuch ma długość
$r=3$ i wagi $-3,-5,-2$. Oba brzegi mają tę samą grupę
$H_1\cong\mathbb Z/25\mathbb Z$, lecz odpowiednie formy
splecenia w kanonicznych generatorach mają wartości
odpowiednio $4/25$ i $9/25$; również współczynniki chirurgii
są różne. Jeśli konfiguracje występują jako podrozmaitości
w różnych czterorozmaitościach, operacja zmienia ich sygnatury
o odpowiednio $4$ i $3$. Z samych danych homologicznych
nie wolno wnioskować, że otrzymane rozmaitości są dyfeomorficzne.
\end{przyklad}

\begin{center}
\begin{tikzpicture}[>=Stealth]
  \node[draw=tgNavy,fill=tgBlue!9,rounded corners=2pt,
    minimum width=27mm,minimum height=11mm,align=center,font=\small] (w) at (0,1.6) {$W$};
  \node[draw=tgCoral,fill=tgCoral!13,rounded corners=3pt,
    minimum width=29mm,minimum height=13mm,align=center,font=\small] (c) at (-3,0) {$C_{p,q}$};
  \node[draw=tgTeal,fill=tgTeal!15,rounded corners=3pt,
    minimum width=29mm,minimum height=13mm,align=center,font=\small] (b) at (3,0) {$B_{p,q}$};
  \draw[tgAccent,-{Stealth[length=2.9mm]}] (w.south west)--(c.north east);
  \draw[tgTree,-{Stealth[length=2.9mm]}] (w.south east)--(b.north west);
  \node[font=\footnotesize,text=tgNavy] at (-3,-1.1) {$X=W\cup C_{p,q}$};
  \node[font=\footnotesize,text=tgNavy] at (3,-1.1) {$X_{p,q}=W\cup B_{p,q}$};
  \node[font=\scriptsize,text=tgNavy] at (0,-1.8)
    {$\partial C_{p,q}\cong\partial B_{p,q}\cong L(p^2,1-pq)$};
\end{tikzpicture}
\end{center}

\section{Granica metody i jej zastosowania}
\label{sec:tg-general-limit}

Wymierna kontrakcja zmienia formę przecięcia i liczby Betti'ego,
ale te dane nie wystarczają do dowodu egzotyczności.
W konstrukcjach gładkich trzeba jeszcze wykazać odpowiednią
prostą spójność i skorzystać z niezmienników zależnych od
struktury różniczkowej. W szczególności wykrywanie
nierównoważnych struktur gładkich wymaga teorii cechowania.
Jednym z jej podstawowych narzędzi są niezmienniki
Seiberga--Wittena, które wprowadzimy w kolejnym rozdziale.

\begin{uwaga}[Uogólniona kontrakcja a typ gładki]
Ciąg Mayera--Vietorisa i algebra macierzy przecięcia
kontrolują homologię, sygnaturę i część informacji
o formie splecenia. Nie dowodzą jednak istnienia
zanurzonej konfiguracji $C_{p,q}$ w zadanej
czterorozmaitości, ani dyfeomorficznego typu jej
dopełnienia. W zastosowaniach do egzotycznych struktur
gładkich każdy z tych punktów wymaga osobnego
argumentu geometrycznego.
\end{uwaga}

\chapter{Wprowadzenie do teorii Seiberga--Wittena i zastosowań w wymiarze cztery}
\label{ch:tg-sw}

Rachunek grafów plumbingowych daje nam formy przecięcia,
grupy homologii i sygnatury. Nie odpowiada jednak na pytanie,
czy dwie homeomorficzne czterorozmaitości są dyfeomorficzne.
Do rozstrzygania tego typu problemów wykorzystuje się
niezmienniki teorii cechowania. Naszym celem jest tu
precyzyjne wprowadzenie niezbędnych obiektów, wyprowadzenie
rachunków topologicznych i wyjaśnienie zastosowania tych
niezmienników do wymiernej kontrakcji. Zaawansowane
fakty analityczne jawnie przyjmujemy z literatury.

\section{Formy samodualne i klasy charakterystyczne}
\label{sec:tg-sw-selfdual}

Niech $X$ będzie zamkniętą, spójną, zorientowaną,
gładką czterorozmaitością z metryką riemannowską $g$.
Operator Hodge'a $*_g:\Omega^2(X)\to\Omega^2(X)$ spełnia
$*_g^2=\mathrm{id}$. Rozkłada zatem wiązkę form na podwiązki
$\Lambda^2=\Lambda^2_+\oplus\Lambda^2_-$, będące wiązkami
własnymi dla wartości własnych odpowiednio $+1$ i $-1$.
Dla formy dwuwymiarowej $\eta$ oznaczamy
$\eta^+:=(\eta+*_g\eta)/2$.

\begin{stwierdzenie}[Formy samodualne i liczba $b_2^+$]
Niech $X$ będzie zamkniętą, spójną, zorientowaną
czterorozmaitością riemannowską. Przestrzeń harmonicznych
form samodualnych ma wymiar $b_2^+(X)$, równy dodatniemu
indeksowi bezwładności formy przecięcia na $H^2(X;\mathbb R)$.
Analogicznie wymiar harmonicznych form antysamodualnych
jest równy $b_2^-(X)$.
\end{stwierdzenie}

\begin{proof}
Z twierdzenia Hodge'a każda klasa de Rhama stopnia dwa
ma dokładnie jednego reprezentanta harmonicznego.
Na przestrzeni harmonicznych form operator $*_g$ zachowuje
harmoniczność, bo komutuje z laplasjanem Hodge'a.
Dla harmonicznej formy samodualnej $\omega$ zachodzi
$\int_X\omega\wedge\omega=\int_X\omega\wedge*_g\omega
=\int_X|\omega|^2\,\mathrm d\mu_g>0$, jeśli $\omega\ne0$.
Dla formy antysamodualnej analogiczna całka jest ujemna.
Obie przestrzenie są ortogonalne względem formy przecięcia,
a ich suma daje całe $H^2(X;\mathbb R)$.
Tak więc ich wymiary są właśnie dodatnim i ujemnym
indeksem tej formy.
\end{proof}

\begin{definicja}[Klasa charakterystyczna]
Niech $X$ będzie jak wyżej. Klasę $c\in H^2(X;\mathbb Z)$
nazywamy \emph{charakterystyczną}, jeżeli
$\langle c\smile x,[X]\rangle\equiv
\langle x\smile x,[X]\rangle\pmod2$
dla każdego $x\in H^2(X;\mathbb Z)$.
Jeżeli nie występuje torsja w $H^2(X;\mathbb Z)$,
warunek ten jest równoważny równości
$c\bmod2=w_2(TX)\in H^2(X;\mathbb Z/2)$,
na mocy wzoru Wu. W obecności torsji sama zgodność
parowań z klasami całkowitymi nie musi wykrywać całej
informacji modulo dwa. Dlatego dla wiązki wyznacznikowej
struktury $\mathrm{Spin}^c$ wymagamy jawnie
$c_1(L)\bmod2=w_2(TX)$.
\end{definicja}

\begin{przyklad}[Klasy charakterystyczne na sumie spójnej]
Na $X=\mathbb{CP}^2\#k\overline{\mathbb{CP}}^{\,2}$
przyjmijmy bazę $H,E_1,\ldots,E_k$ o kwadratach
$H^2=1$, $E_i^2=-1$ i zerowych przecięciach mieszanych.
Ponieważ macierz przecięcia jest diagonalna z nieparzystymi
wpisami, klasa $c=aH-\sum b_iE_i$ jest charakterystyczna
wtedy i tylko wtedy, gdy wszystkie liczby $a,b_i$ są
nieparzyste. W szczególności warunek charakterystyczności
ogranicza zbiór klas, które mogą być podstawą
niezmienników Seiberga--Wittena, ale nie orzeka,
że każda taka klasa daje niezerowy niezmiennik.
\end{przyklad}

\section{Struktura \texorpdfstring{$\mathrm{Spin}^c$}{Spinc} i równania}
\label{sec:tg-sw-equations}

\begin{definicja}[Dane struktury $\mathrm{Spin}^c$]
Niech $X$ będzie zamkniętą zorientowaną gładką
czterorozmaitością. Struktura $\mathrm{Spin}^c$ na $X$
składa się z odpowiedniego podniesienia wiązki
ortogonalnych układów odniesienia do grupy
$\mathrm{Spin}^c(4)$; po wyborze metryki indukuje
zespolone wiązki spinorowe $S^+$ i $S^-$ rangi dwa
oraz zespoloną wiązkę liniową $L=\det S^+$.
Klasa $c_1(L)$ spełnia $c_1(L)\bmod2=w_2(TX)$.
Przy ustalonej metryce i strukturze $\mathrm{Spin}^c$
połączenie unitarne $A$ na $L$ wyznacza operator Diraca
$D_A:\Gamma(S^+)\to\Gamma(S^-)$.
\end{definicja}

\begin{uwaga}
Konstrukcja wiązek spinorowych wymaga algebry Clifforda
i pojęcia struktur spinorowych. Jej pełne wyprowadzenie
odsyłamy do rozdziałów książki o spinorach.
W dalszym ciągu korzystamy z podstawowych własności
operatora Diraca oraz z twierdzenia indeksowego
Atiyaha--Singera; nie odtwarzamy tu wszystkich narzędzi
analitycznych teorii cechowania.
\end{uwaga}

\begin{definicja}[Równania Seiberga--Wittena]
\label{def:tg-sw-eqn}
Niech $g$ będzie metryką riemannowską na zamkniętej,
spójnej, zorientowanej czterorozmaitości $X$, a
$\mathfrak s$ strukturą $\mathrm{Spin}^c$ o wiązce
wyznacznikowej $L$. Niech $\eta\in\Omega^2_+(X)$
będzie ustaloną rzeczywistą formą samodualną.
Szukamy pary $(A,\psi)$, gdzie $A$ jest połączeniem
unitarnym na $L$, a $\psi\in\Gamma(S^+)$ spinorem,
spełniającej
\[
D_A\psi=0,\qquad F_A^+=q(\psi)+i\eta.
\]
Tutaj $F_A\in i\Omega^2(X)$ jest krzywizną połączenia,
$F_A^+$ jej częścią samodualną, a
$q:S^+\to i\Lambda^2_+$ standardowym kwadratowym
odwzorowaniem związanym z mnożeniem Clifforda;
wybór jego stałej normalizacyjnej ustalamy zgodnie
z zapisem drugiego równania.
\end{definicja}

\begin{definicja}[Przestrzeń moduli]
\emph{Grupą cechowania} nazywamy
$\mathcal G=C^\infty(X,S^1)$.
Działa ona na parach $(A,\psi)$ przez
$u\cdot(A,\psi)=(A-2u^{-1}\mathrm du,u\psi)$,
przy standardowej konwencji połączenia wyznacznikowego.
\emph{Przestrzenią moduli} $\mathcal M_{\mathfrak s}(g,\eta)$
nazywamy zbiór rozwiązań równań z definicji
\ref{def:tg-sw-eqn}, podzielony przez działanie $\mathcal G$.
\end{definicja}

\begin{uwaga}[Rozwiązania redukowalne]
Pary z $\psi=0$ nazywamy \emph{redukowalnymi}.
Spełniają one równanie $F_A^+=i\eta$.
Dla $b_2^+>0$ można generycznie dobrać perturbację
tak, aby uniknąć rozwiązań redukowalnych. Przy odpowiedniej
transwersalności przestrzeń moduli jest wtedy gładka,
a twierdzenie o zwartości zapewnia jej zwartość.
Przy $b_2^+>1$ generyczna jednoparametrowa rodzina danych
może dodatkowo ominąć ścianę redukowalnych. Nie jest to jednak wyłącznie argument
kombinatoryczny: wymaga analizy eliptycznej, transwersalności
i twierdzenia o zwartości.
\end{uwaga}

\section{Wymiar oczekiwany i definicja niezmiennika}
\label{sec:tg-sw-dimension}

\begin{twierdzenie}[Oczekiwany wymiar przestrzeni moduli]
\label{thm:tg-sw-dimension}
Dla zamkniętej, spójnej, zorientowanej,
gładkiej czterorozmaitości $X$ i struktury
$\mathrm{Spin}^c$ z klasą $c=c_1(L)$ rzeczywisty
wymiar oczekiwany przestrzeni moduli wynosi
\[
d(\mathfrak s)=\frac{\langle c^2,[X]\rangle-2e(X)-3\sigma(X)}4.
\]
Jest on liczbą całkowitą. Gdy $b_2^+(X)>1$ i dla danych
generycznych nie występują rozwiązania redukowalne,
jest to wymiar gładkiej części przestrzeni moduli.
\end{twierdzenie}

\begin{proof}[Rachunek indeksu; użycie twierdzenia Atiyaha--Singera]
Kompleks deformacyjny teorii cechowania zawiera część
Diraca oraz część rzeczywistą z operatorem
$\mathrm d^*\oplus\mathrm d^+:\Omega^1(X)\to
\Omega^0(X)\oplus\Omega^2_+(X)$.
Twierdzenie Atiyaha--Singera daje zespolony indeks operatora
Diraca $\operatorname{ind}_{\mathbb C}D_A
=(\langle c^2,[X]\rangle-\sigma(X))/8$.
Część rzeczywista wnosi do indeksu deformacji
$b_1(X)-b_0(X)-b_2^+(X)=b_1(X)-1-b_2^+(X)$.
Łączny indeks rzeczywisty wynosi wobec tego
\[
2\operatorname{ind}_{\mathbb C}D_A+b_1-1-b_2^+
=\frac{\langle c^2,[X]\rangle-\sigma(X)}4+b_1-1-b_2^+.
\]
Ponieważ $e(X)=2-2b_1+b_2^++b_2^-$ i
$\sigma(X)=b_2^+-b_2^-$, prawa strona po uproszczeniu
jest równa wyrażeniu w tezie.
Całkowitość wynika z indeksowej interpretacji.
Nie dowodzimy tu samego twierdzenia Atiyaha--Singera ani
analitycznej regularności przestrzeni moduli.
\end{proof}

\begin{definicja}[Niezmiennik Seiberga--Wittena w wymiarze zero]
Niech $X$ będzie zamkniętą, spójną, zorientowaną,
gładką czterorozmaitością z $b_1(X)=0$ i $b_2^+(X)>1$.
Ustalmy orientację $H^2_+(X;\mathbb R)$.
Dla struktury $\mathrm{Spin}^c$ o $d(\mathfrak s)=0$
i generycznych danych $(g,\eta)$ przestrzeń moduli
składa się ze skończonej liczby punktów, liczonych
z odpowiednimi znakami orientacji. Ich algebraiczną
liczbę oznaczamy przez $SW_X(\mathfrak s)\in\mathbb Z$.
Gdy $d(\mathfrak s)<0$, przyjmujemy $SW_X(\mathfrak s)=0$.
W przypadku $d(\mathfrak s)>0$ do zdefiniowania
liczbowych niezmienników używa się dodatkowych klas
kohomologii na przestrzeni moduli; nie są one częścią
powyższej definicji.
\end{definicja}

\begin{twierdzenie}[Niezmienniczość; wynik analityczny]
Dla $X$ jak w poprzedniej definicji, po ustaleniu
orientacji homologicznej, $SW_X(\mathfrak s)$ nie zależy
od wyboru danych generycznych $(g,\eta)$ i jest
niezmiennikiem gładkiej struktury na $X$, z dokładnością
do naturalnego działania dyfeomorfizmów na strukturach
$\mathrm{Spin}^c$ i na orientacji homologicznej.
\end{twierdzenie}

\begin{uwaga}[Przypadek $b_2^+=1$]
Dla $b_2^+(X)=1$ niezmiennik może zależeć od wyboru
\emph{komory}: podczas zmiany metryki i perturbacji
można przejść przez ścianę, na której występują
rozwiązania redukowalne. W związku z tym porównanie
z rozmaitościami wymiernymi, takimi jak
$\mathbb{CP}^2\#7\overline{\mathbb{CP}}^{\,2}$,
wymaga dodatkowej formuły zmiany komory i świadomego
wyboru tej komory. Nie wystarcza samo wyliczenie
$d(\mathfrak s)$.
\end{uwaga}

\section{Klasy podstawowe i ograniczenia geometryczne}
\label{sec:tg-sw-basic}

\begin{definicja}[Klasa podstawowa]
Dla $X$ z $b_1=0$, $b_2^+>1$ i dla wersji
niezmienników o wymiarze oczekiwanym zero klasę
charakterystyczną $c\in H^2(X;\mathbb Z)$
nazywamy \emph{klasą podstawową Seiberga--Wittena},
jeżeli istnieje struktura $\mathrm{Spin}^c$ o
$c_1(\det S^+)=c$ i $SW_X(\mathfrak s)\ne0$.
Przy obecności torsji odróżniamy klasę $c$ od samej
struktury $\mathrm{Spin}^c$.
\end{definicja}

\begin{twierdzenie}[Zerowanie przy dodatniej krzywiźnie skalarnej]
\label{thm:tg-sw-psc}
Jeżeli zamknięta, spójna, zorientowana,
gładka czterorozmaitość $X$ ma $b_2^+(X)>1$
i dopuszcza metrykę riemannowską o wszędzie dodatniej
krzywiźnie skalarnej, to wszystkie jej niezmienniki
Seiberga--Wittena z powyższej definicji są zerowe.
\end{twierdzenie}

\begin{proof}[Argument Weitzenböcka; krok analityczny wskazany jawnie]
Dla rozwiązania nieperturbowanych równań tożsamość
Weitzenböcka, po połączeniu z równaniem krzywizny,
daje całkę zawierającą nieujemne składniki
$|\nabla_A\psi|^2$, stałą dodatnią razy $|\psi|^4$
i wyraz $s_g|\psi|^2/4$, gdzie $s_g$ jest krzywizną
skalarną. Gdy $s_g>0$, suma może być zerowa tylko
przy $\psi=0$. Redukowalne rozwiązania można wykluczyć
małą generyczną perturbacją w sytuacji $b_2^+>1$,
zachowując oszacowanie Weitzenböcka i dodatniość krzywizny.
Dla tych danych przestrzeń moduli jest pusta, więc
jej liczba algebraiczna wynosi zero.
Ścisłe uzasadnienie oszacowania przy perturbacji
korzysta ze standardowej analizy równań Seiberga--Wittena.
\end{proof}

\begin{twierdzenie}[Taubes; niezerowanie dla rozmaitości symplektycznych]
\label{thm:tg-sw-taubes}
Jeżeli $(X,\omega)$ jest zamkniętą, spójną,
symplektyczną czterorozmaitością z $b_2^+(X)>1$,
to struktury $\mathrm{Spin}^c$ związane z klasami
$\pm K_\omega$ (gdzie $K_\omega=-c_1(TX,J)$ dla
zgodnej z $\omega$ struktury prawie zespolonej $J$)
mają niezerowe niezmienniki Seiberga--Wittena.
W odpowiedniej orientacji homologicznej ich wartości
są równe $\pm1$. Wynik ten przyjmujemy bez dowodu.
\end{twierdzenie}

\begin{wniosek}
Zamknięta symplektyczna czterorozmaitość
$X$ z $b_2^+(X)>1$ nie dopuszcza metryki o wszędzie
dodatniej krzywiźnie skalarnej.
\end{wniosek}

\begin{proof}
Twierdzenie Taubesa zapewnia niezerowy niezmiennik,
a poprzednie twierdzenie o dodatniej krzywiźnie
wymuszałoby zerowanie wszystkich niezmienników.
Otrzymujemy sprzeczność.
\end{proof}

\section{Wymierna kontrakcja a niezmienniki gładkie}
\label{sec:tg-sw-rb}

\begin{twierdzenie}[Fintushel--Stern; formuła podniesienia, bez dowodu]
\label{thm:tg-sw-rbprinciple}
Niech $X$ będzie zamkniętą, po prostu spójną,
gładką, zorientowaną czterorozmaitością z
$b_2^+(X)>1$, zawierającą standardową konfigurację
$C_p$, $p\geq2$. Niech $W=X\setminus\operatorname{int}C_p$
i $X_{(p)}=W\cup B_p$ będzie wymierną kontrakcją.
Ustalmy zgodne orientacje homologiczne.

Niech $\mathfrak s'$ będzie strukturą $\mathrm{Spin}^c$
na $X_{(p)}$ z $d_{X_{(p)}}(\mathfrak s')=0$.
Przy ustalonym sklejeniu brzegu nazywamy
\emph{dopuszczalnym podniesieniem} $\mathfrak s'$
strukturę $\mathfrak s$ na $X$, dla której
$\mathfrak s|_W\cong\mathfrak s'|_W$, restrykcje na
wspólnym brzegu zgadzają się jako struktury
$\mathrm{Spin}^c$, a ich przedłużenie przez $C_p$
spełnia $d_X(\mathfrak s)=d_{X_{(p)}}(\mathfrak s')=0$.
Ostatni warunek można równoważnie zapisać jako
$c_1(\mathfrak s)^2=c_1(\mathfrak s')^2-(p-1)$.
Dla takiego dopuszczalnego podniesienia zachodzi
\[
SW_{X_{(p)}}(\mathfrak s')=SW_X(\mathfrak s)
\]
przy odpowiednim zgodnym wyborze orientacji homologicznych.
Istnienie i własności dopuszczalnych podniesień są częścią
teorii Fintushela--Sterna: nie można zastąpić ich
warunkiem samej zgodności klas na $W$.
Wynik przyjmujemy z literatury; jego dowód wymaga
formuły sklejenia dla równań Seiberga--Wittena.
\end{twierdzenie}

\begin{uwaga}[Dlaczego sformułowanie jest ostrożne]
Pełna formuła sklejenia wymaga zidentyfikowania struktur
$\mathrm{Spin}^c$ na brzegu oraz dopasowania ich przedłużeń
przez $B_{p,q}$, a także uwzględnienia wymiarów
oczekiwanych i orientacji przestrzeni moduli. Samo
porównanie $c^2$ lub wyznacznika formy przecięcia
nie wyznacza liczby rozwiązań równań cechowania.
Dla $b_2^+=1$ trzeba dodatkowo kontrolować komory,
których definicję i formułę zmiany przedstawimy w rozdziale
\ref{ch:tg-sw-analiza}.
\end{uwaga}

\begin{przyklad}[Kontrola równości wymiarów oczekiwanych]
Niech $X\to X_{p,q}$ będzie uogólnioną wymierną kontrakcją
konfiguracji długości $r$. Oznaczmy
$\kappa(Z):=2e(Z)+3\sigma(Z)$, a przez
$c\in H^2(X;\mathbb Z)$ i $c'\in H^2(X_{p,q};\mathbb Z)$
klasy charakterystyczne. Jeżeli
$\langle c'^2,[X_{p,q}]\rangle
=\langle c^2,[X]\rangle+r$, to ze wzoru indeksowego i
$\kappa(X_{p,q})=\kappa(X)+r$ wynika
$d_{X_{p,q}}(c')=d_X(c)$.
Warunek o kwadratach klas nie jest automatyczny;
wybór dopuszczalnego podniesienia musi go zapewnić.
Ten rachunek wyjaśnia, jakie dane poza samą formą
przecięcia są potrzebne do porównania niezmienników.
\end{przyklad}

\begin{przyklad}[Wymierna kontrakcja sekcji w $E(4)$]
Niech $E(4)$ oznacza standardową względnie minimalną,
po prostu spójną powierzchnię eliptyczną o
$e(E(4))=48$ i $\sigma(E(4))=-32$.
Ma ona $b_2^+=7$ i $b_2^-=39$.
Jej sekcja jest sferą o samoprzecięciu $-4$,
a więc stanowi konfigurację $C_2$. Po wymiernej
kontrakcji otrzymujemy gładką czterorozmaitość $Y$
o liczbach $e(Y)=47$, $\sigma(Y)=-31$,
$b_2^+(Y)=7$ i $b_2^-(Y)=38$.

Dla $E(4)$ znane obliczenie daje wielomian
Seiberga--Wittena $(t-t^{-1})^2$, w którym
$t$ symbolizuje klasę włókna eliptycznego.
Klasy podstawowe mają więc postać $-2F,0,2F$.
Gdy $S$ jest sekcją, zachodzi $F\cdot S=1$,
a zatem $(\pm2F)\cdot S=\pm2$.
Twierdzenie o wymiernej kontrakcji, zastosowane
z odpowiednimi strukturami $\mathrm{Spin}^c$,
prowadzi do dwóch ocalałych klas podstawowych
nowej rozmaitości $Y$ i wielomianu postaci
$t_Y+t_Y^{-1}$ (z dokładnością do ogólnego znaku).
Jest to wynik głębokiej formuły sklejenia, a nie
bezpośredni wniosek z liczb Betti'ego.
\end{przyklad}

\section{Jak wykorzystać teorię do rozróżniania struktur gładkich?}
\label{sec:tg-sw-exotic}

\begin{stwierdzenie}[Zasada rozróżniania]
Niech $X_1,X_2$ będą zamkniętymi, spójnymi,
zorientowanymi, gładkimi czterorozmaitościami.
Załóżmy, że są homeomorficzne z zachowaniem orientacji,
lecz ich zbiory niezerowych niezmienników
Seiberga--Wittena nie odpowiadają sobie przy żadnej
izometrii ich całkowitych form przecięcia indukowanej
przez możliwy dyfeomorfizm (z uwzględnieniem orientacji
homologicznej i, gdy trzeba, komory). Wówczas $X_1$
i $X_2$ nie są dyfeomorficzne z zachowaniem orientacji.
\end{stwierdzenie}

\begin{proof}
Gdyby istniał dyfeomorfizm zachowujący orientację,
indukowałby izometrię form przecięcia i odpowiadające
sobie struktury $\mathrm{Spin}^c$. Niezmienniczość
Seiberga--Wittena wymuszałaby wówczas zgodność
wartości niezmienników po przeniesieniu tych struktur,
wbrew założeniu.
\end{proof}

\begin{center}
\begin{tikzpicture}[>=Stealth,node distance=8mm]
  \node[draw=tgBlue,fill=tgBlue!9,rounded corners=3pt,
    minimum width=29mm,minimum height=11mm,align=center,font=\small] (a) at (-3,0)
    {graf plumbingowy\\$Q_{p,q}$};
  \node[draw=tgTeal,fill=tgTeal!12,rounded corners=3pt,
    minimum width=31mm,minimum height=11mm,align=center,font=\small] (b) at (.4,0)
    {wymierna kontrakcja\\$X\longmapsto X_{p,q}$};
  \node[draw=tgCoral,fill=tgCoral!12,rounded corners=3pt,
    minimum width=28mm,minimum height=11mm,align=center,font=\small] (c) at (3.9,0)
    {równania SW\\klasy podstawowe};
  \draw[tgTree,-{Stealth[length=2.7mm]}] (a)--(b);
  \draw[tgAccent,-{Stealth[length=2.7mm]}] (b)--(c);
  \node[font=\footnotesize,text=tgNavy,align=center] at (.4,-1.08)
    {Topologia: $\pi_1,Q,e,\sigma$\quad i osobno\quad gładkość: $SW$};
\end{tikzpicture}
\end{center}

\begin{uwaga}[Granice obecnego rozdziału]
Nie przeprowadziliśmy pełnego dowodu twierdzeń o zwartości
i transwersalności przestrzeni moduli, indeksu Atiyaha--Singera,
twierdzenia Taubesa ani formuły Fintushela--Sterna.
Przedstawiliśmy natomiast precyzyjnie niezbędne obiekty,
pełne rachunki topologiczne i logiczne miejsce, w którym
niezmienniki analityczne rozstrzygają problem gładkości.
W szczególności w przypadkach $b_2^+=1$, takich jak egzotyczne
$\mathbb{CP}^2\#7\overline{\mathbb{CP}}^{\,2}$, potrzebna
jest jeszcze osobna teoria komór i bardziej szczegółowa
formuła sklejenia. Tych kroków nie zastępujemy
rachunkiem charakterystyki Eulera.
\end{uwaga}


% Iteracja 11. Bez preambuły; wykorzystuje istniejące środowiska i barwy.
\begin{uwaga}[Wymiar oczekiwany a niezerowość niezmiennika]
Wartość $d(\mathfrak s)=0$ oznacza jedynie zerowy indeks
formalnego operatora deformacyjnego. Nie zapewnia istnienia
żadnego rozwiązania równań Seiberga--Wittena, regularności
wszystkich rozwiązań ani niezerowego zliczenia.
Do obliczenia niezmiennika potrzebne są niezależnie
zwartość, transwersalność i orientacje; przy $b_2^+=1$
również wybór komory. Dlatego sama zgodność klas
charakterystycznych i kwadratów ich pierwszych klas
Cherna nie stanowi dowodu równości niezmienników.
\end{uwaga}

\chapter{Analityczne podstawy teorii Seiberga--Wittena: przestrzenie moduli i komory}
\label{ch:tg-sw-analiza}

Poprzedni rozdział zdefiniował równania Seiberga--Wittena, wskazał ich
wymiar oczekiwany i zastosowania do wymiernej kontrakcji. Aby zrozumieć,
dlaczego liczenie rozwiązań daje niezmiennik różniczkowy, potrzebujemy
jeszcze aparatu analitycznego. Omówimy kolejno dopełnienia Sobolewa,
iloraz przez grupę cechowania, kompleks deformacyjny, zwartość,
transwersalność i orientacje. Następnie wyjaśnimy, dlaczego przypadek
$b_2^+=1$ wymaga teorii komór. Podamy pełne argumenty algebraiczne
i zasadnicze oszacowania; dwa głębokie kroki analityczne, tj. globalną
teorię przekrojów Fredholma i lokalną teorię Kuranishiego w pobliżu
rozwiązań redukowalnych, sformułujemy jawnie jako przyjmowane twierdzenia.

\section{Dopełnienia Sobolewa i grupa cechowania}
\label{sec:tg-sw-sobolev}

Przez cały rozdział $X$ oznacza zamkniętą, spójną, zorientowaną,
gładką czterorozmaitość z ustaloną metryką riemannowską $g$,
a $\mathfrak s$ --- strukturę $\mathrm{Spin}^c$ z wiązkami $S^\pm$
i wiązką wyznacznikową $L$. Oznaczamy przez
$\Omega^j(X;i\mathbb R)$ przestrzeń gładkich $j$-form o wartościach
w $i\mathbb R$ oraz przez $\Gamma(S^\pm)$ przestrzenie spinorów.
Wybieramy raz na zawsze gładkie połączenie unitarne $A_0$ na $L$.

\begin{definicja}[Przestrzenie Sobolewa i konfiguracje]
Dla całkowitego $k\geq3$ niech $W^{k,2}(E)$ będzie dopełnieniem
przestrzeni gładkich przekrojów wiązki euklidesowej lub hermitowskiej
$E\to X$ w normie Sobolewa rzędu $k$, zdefiniowanej przez
wybraną koneksję i metrykę. \emph{Przestrzenią konfiguracji} nazywamy
przestrzeń afiniczną
\[
\mathcal C_k=(A_0+W^{k,2}(iT^*X))\times W^{k,2}(S^+).
\]
\emph{Grupą cechowania} jest $\mathcal G_{k+1}:=W^{k+1,2}(X,S^1)$,
z działaniem $u\cdot(A,\psi)=(A-2u^{-1}\mathrm du,u\psi)$.
Konfigurację nazywamy \emph{nieredukowalną}, jeżeli
$\psi\not\equiv0$; zbiór takich konfiguracji oznaczamy
$\mathcal C_k^*$.
\end{definicja}

Wymóg $k\geq3$ zapewnia przez twierdzenie Sobolewa
$W^{k,2}(X)\hookrightarrow C^0(X)$; w szczególności
iloczyny i nieliniowe wyrażenia w równaniach mają właściwy sens.
Do operatora równań używamy przestrzeni docelowej
$W^{k-1,2}(S^-\oplus i\Lambda^2_+)$. W rachunkach
linearyzacyjnych należy odróżnić parametr rzeczywisty $f$,
odpowiadający jednoparametrowej grupie $u_t=e^{itf}$, od form
i-funkcji o wartościach w $i\mathbb R$.

\begin{stwierdzenie}[Stabilizator konfiguracji]
\label{prop:tg-sw-stabilizer}
Niech $(A,\psi)\in\mathcal C_k$. Jeżeli $u\in\mathcal G_{k+1}$
zachowuje tę konfigurację, to $u$ jest stałą funkcją o wartości
w $S^1$. Jeżeli ponadto $\psi\not\equiv0$, to $u=1$.
Wobec tego działanie $\mathcal G_{k+1}$ na $\mathcal C_k^*$
jest wolne, a stabilizatorem konfiguracji redukowalnej jest $S^1$.
\end{stwierdzenie}

\begin{proof}
Z równości $A-2u^{-1}\mathrm du=A$ wynika
$u^{-1}\mathrm du=0$. Funkcja $u$ jest więc lokalnie stała,
a ponieważ $X$ jest spójna, jest stała globalnie.
Równość $u\psi=\psi$ przy $\psi\not\equiv0$ wymusza
$u=1$ w punkcie, w którym $\psi$ jest niezerowy,
a więc na całym $X$. Przy $\psi=0$ każda stała faza
zachowuje konfigurację.
\end{proof}

\begin{uwaga}[Dlaczego potrzebny jest wybór przekroju?]
Iloraz $\mathcal C_k/\mathcal G_{k+1}$ nie jest na ogół gładką
rozmaitością w pobliżu konfiguracji redukowalnych, ponieważ
stabilizator nie jest tam trywialny. Na części nieredukowalnej
twierdzenie o przekroju lokalnym (ang. \emph{slice theorem})
daje lokalny model ilorazu. Warunek $\mathrm d^*a=0$
na różnicy połączeń $A-A_0=ia$ opisuje cechowanie Coulomba
samej części połączeniowej, ale nie usuwa resztkowego działania
stałych faz na spinor. Pełny liniowy warunek przekroju
w punkcie $(A,\psi)$ ma postać $\mathrm d_0^*(ia,\phi)=0$,
a zatem uwzględnia również kierunek spinorowy.
Nie należy utożsamiać takiego lokalnego przekroju
z globalnym wyborem cechowania na całej przestrzeni konfiguracji.
\end{uwaga}

\section{Przekrój Coulomba i kompleks deformacyjny}
\label{sec:tg-sw-deform}

Wybierzmy rozwiązanie $(A,\psi)$ równań z poprzedniego
rozdziału. Wygodnie zapisać ich lewą stronę jako
$\mathfrak F_\eta(A,\psi):=(D_A\psi,F_A^+-q(\psi)-i\eta)$.
Niech $\rho$ oznacza mnożenie Clifforda, a
$\mathrm dq_\psi$ różniczkę odwzorowania kwadratowego $q$
w punkcie $\psi$.

\begin{stwierdzenie}[Linearyzacja i różniczka działania]
\label{prop:tg-sw-linear}
W punkcie będącym rozwiązaniem równania różniczka działania
cechowania w kierunku $f\in\Omega^0(X;\mathbb R)$ oraz
różniczka operatora $\mathfrak F_\eta$ mają postać
\[
\begin{aligned}
\mathrm d_0f&=(-2i\,\mathrm df,\,if\psi),\\
\mathrm d_1(a,\phi)&=\bigl(D_A\phi+\tfrac12\rho(a)\psi,\,
\mathrm d^+a-(\mathrm dq)_\psi\phi\bigr),
\end{aligned}
\]
gdzie $a\in\Omega^1(X;i\mathbb R)$ i $\phi\in\Gamma(S^+)$.
Zachodzi $\mathrm d_1\circ\mathrm d_0=0$.
\end{stwierdzenie}

\begin{proof}
Dla $u_t=e^{itf}$ mamy
$u_t^{-1}\mathrm du_t=it\,\mathrm df$; różniczkując działanie
w $t=0$, otrzymujemy wzór na $\mathrm d_0$.
Krzywizna połączenia $A+ta$ jest równa
$F_A+t\,\mathrm da$, gdyż grupa strukturalna $U(1)$
jest przemienna. Stąd różniczka równania krzywizny
to $\mathrm d^+a-(\mathrm dq)_\psi\phi$.
Dla przyjętej konwencji połączenia na wiązce wyznacznikowej
różniczka $D_{A+ta}$ wynosi $\rho(a)/2$, co daje
pierwszą składową $\mathrm d_1$.
Wreszcie $\mathfrak F_\eta$ jest równoważny względem
cechowania, a zatem przy rozwiązaniu
$\mathfrak F_\eta(A,\psi)=0$ różniczka w kierunku orbity
cechowania znika. Otrzymujemy $\mathrm d_1\mathrm d_0=0$.
\end{proof}

\begin{lemat}[Warunek przekroju i stałe transformacje cechowania]
\label{lem:tg-sw-full-coulomb}
Niech $(A,\psi)$ będzie gładką konfiguracją nieredukowalną.
Na przestrzeni stycznej wyposażonej w rzeczywisty iloczyn
skalarny $L^2$ formalne sprzężenie operatora
$\mathrm d_0f=(-2i\,\mathrm df,if\psi)$ jest dane wzorem
\[
\mathrm d_0^*(i\alpha,\phi)
=-2\,\mathrm d^*\alpha
+\operatorname{Re}\langle i\psi,\phi\rangle,
\]
gdzie $\alpha\in\Omega^1(X;\mathbb R)$, a iloczyn
spinorów jest hermitowski. Ponadto
$\mathrm d_0^*\mathrm d_0=4\,\mathrm d^*\mathrm d+|\psi|^2$
na funkcjach rzeczywistych i operator ten ma zerowe jądro.
W konsekwencji w otoczeniu konfiguracji nieredukowalnej
równanie $\mathrm d_0^*(i\alpha,\phi)=0$ jest liniowym
warunkiem lokalnego przekroju dla działania cechowania.
\end{lemat}

\begin{proof}
Z definicji sprzężenia i dla dowolnej rzeczywistej funkcji
$f$ obliczamy
\[
\begin{aligned}
\langle \mathrm d_0f,(i\alpha,\phi)\rangle_{L^2}
&=-2\langle\mathrm df,\alpha\rangle_{L^2}
+\int_X f\,\operatorname{Re}\langle i\psi,\phi\rangle
\,\mathrm d\mu_g\\
&=\int_X f\bigl(-2\,\mathrm d^*\alpha
+\operatorname{Re}\langle i\psi,\phi\rangle\bigr)
\,\mathrm d\mu_g.
\end{aligned}
\]
Stąd otrzymujemy wzór na sprzężenie oraz
$\mathrm d_0^*\mathrm d_0f=4\,\mathrm d^*\mathrm df+|\psi|^2f$.
Jeśli ten operator zeruje $f$, to przez całkowanie
$4\|\mathrm df\|_{L^2}^2+\|f\psi\|_{L^2}^2=0$.
Zatem $f$ jest stała i $f\psi=0$. Ponieważ
$\psi\not\equiv0$, musi być $f=0$.
Eliptyczny, samosprzężony operator $\mathrm d_0^*\mathrm d_0$
ma więc trywialne jądro i, na odpowiednich dopełnieniach
Sobolewa, jest odwracalny. Lokalny przekrój otrzymuje się
z twierdzenia o funkcji uwikłanej dla działania grupy
cechowania. Zastosowanie tego ostatniego twierdzenia
jest standardowym krokiem analizy banachowej.
\end{proof}

\begin{definicja}[Kompleks deformacyjny]
\emph{Kompleksem deformacyjnym Seiberga--Wittena} w punkcie
$(A,\psi)$ będącym rozwiązaniem nazywamy kompleks rzeczywisty
\[
0\longrightarrow\Omega^0(X;\mathbb R)
\xrightarrow{\mathrm d_0}\Omega^1(X;i\mathbb R)\oplus\Gamma(S^+)
\xrightarrow{\mathrm d_1}\Omega^2_+(X;i\mathbb R)\oplus\Gamma(S^-)
\longrightarrow0.
\]
Pierwszy wyraz zapisujemy jako funkcje rzeczywiste,
ponieważ przyjęliśmy $u_t=e^{itf}$ oraz
$\mathrm d_0f=(-2i\,\mathrm df,if\psi)$.
Równoważna konwencja z $if\in\Omega^0(X;i\mathbb R)$
prowadzi do tej samej teorii po odpowiedniej zmianie wzoru
na $\mathrm d_0$. Jego kohomologie oznaczamy przez
$H^0_{A,\psi}$, $H^1_{A,\psi}$ i $H^2_{A,\psi}$.
\end{definicja}

\begin{uwaga}[Rola poszczególnych kohomologii]
$H^0_{A,\psi}$ jest algebrą Liego stabilizatora,
$H^1_{A,\psi}$ jest przestrzenią styczną do przestrzeni
moduli w przypadku regularnym, a $H^2_{A,\psi}$ jest
przestrzenią przeszkód w zastosowaniu twierdzenia o funkcji
uwikłanej. Dla rozwiązania nieredukowalnego
$H^0_{A,\psi}=0$. Gdy ponadto $H^2_{A,\psi}=0$, rozwiązanie
jest \emph{regularne} i przestrzeń moduli jest lokalnie gładka.
W przypadku nieregularnym $H^1_{A,\psi}$ opisuje
przestrzeń deformacji pierwszego rzędu, lecz nie musi być
przestrzenią styczną do gładkiej rozmaitości.
Są to stwierdzenia lokalne, po wykonaniu dopełnień Sobolewa
i użyciu twierdzenia o przekroju.
\end{uwaga}

\begin{twierdzenie}[Eliptyczność i indeks]
\label{thm:tg-sw-elliptic}
Kompleks deformacyjny jest eliptyczny, a dla $\mathfrak s$
o klasie $c=c_1(L)$ jego indeks rzeczywisty
$\dim H^1-\dim H^0-\dim H^2$ wynosi
\[
d(\mathfrak s)=\frac{\langle c^2,[X]\rangle-2e(X)-3\sigma(X)}4.
\]
\end{twierdzenie}

\begin{proof}[Wskazanie części elementarnej i twierdzenia indeksowego]
Symbole członów liniowych nie zależą od $\psi$:
człony z mnożeniem przez $\psi$ i różniczką $q$ mają rząd zero.
Główna część rozpada się więc na symbol kompleksu
$\Omega^0\xrightarrow{\mathrm d}\Omega^1
\xrightarrow{\mathrm d^+}\Omega^2_+$ oraz symbol operatora Diraca.
Dla niezerowego kowektora $\xi$ symbol pierwszego
kompleksu jest dokładny w środkowym wyrazie, a symbol
Diraca jest izomorfizmem zespolonych przestrzeni spinorowych.
To wyjaśnia eliptyczność.
Indeks części de Rhama wynosi $b_1-1-b_2^+$,
a indeks rzeczywisty Diraca to
$2\operatorname{ind}_{\mathbb C}D_A=(c^2-\sigma)/4$.
Po ich dodaniu i podstawieniu
$e=2-2b_1+b_2^++b_2^-$ oraz
$\sigma=b_2^+-b_2^-$ dostajemy wzór z tezy.
Wartość indeksu Diraca jest tutaj zastosowaniem
twierdzenia Atiyaha--Singera, którego nie dowodzimy.
\end{proof}

\section{Oszacowanie spinora i zwartość}
\label{sec:tg-sw-compact}

Kluczowa różnica między teorią Seiberga--Wittena a niektórymi
innymi teoriami cechowania polega na tym, że na zamkniętej
czterorozmaitości przestrzeń rozwiązań modulo cechowanie
jest zwarta bez dodawania punktów odpowiadających
koncentracji energii. Wyprowadzimy najważniejsze oszacowanie.

\begin{lemat}[Oszacowanie z formuły Weitzenböcka]
\label{lem:tg-sw-bound}
Przy ustalonej gładkiej metryce $g$, gładkiej perturbacji
$\eta\in\Omega^2_+(X)$ i normalizacji odwzorowania $q$
jak w rozdziale poprzednim istnieją stałe $c_0,C_0>0$,
zależne tylko od przyjętych konwencji algebraicznych, takie że
każde gładkie rozwiązanie równań Seiberga--Wittena spełnia
\[
\|\psi\|_{L^\infty}^2\le
c_0^{-1}\bigl(\tfrac14\|s_g^-\|_{L^\infty}
+C_0\|\eta\|_{L^\infty}\bigr),
\qquad s_g^-:=\max\{0,-s_g\}.
\]
\end{lemat}

\begin{proof}
Formuła Weitzenböcka dla operatora $D_A$ wyraża
$D_A^*D_A$ jako sumę laplasjanu koneksyjnego,
członu $s_g/4$ i mnożenia Clifforda przez krzywiznę
samodualną. Po podstawieniu $D_A\psi=0$ i
$F_A^+=q(\psi)+i\eta$ otrzymujemy punktową równość,
w której wyraz kwadratowy w spinorze jest
nieujemny i ma postać $c_0|\psi|^4$, natomiast
wyraz perturbacyjny jest ograniczony z góry przez
$C_0|\eta||\psi|^2$.

Wybierzmy punkt $x_0$, w którym funkcja $|\psi|^2$
osiąga maksimum. Dla ustalonej konwencji znaku
laplasjanu zasada maksimum i tożsamość Weitzenböcka
dają w tym punkcie
\[
c_0|\psi(x_0)|^4\le
\bigl(\tfrac14s_g^-(x_0)+C_0|\eta(x_0)|\bigr)
|\psi(x_0)|^2.
\]
Jeżeli $\psi(x_0)=0$, spinor jest identycznie zerowy.
W przeciwnym razie dzielimy przez $|\psi(x_0)|^2$
i uzyskujemy żądane oszacowanie globalne.
\end{proof}

\begin{twierdzenie}[Zwartość przestrzeni moduli; twierdzenie analityczne]
\label{thm:tg-sw-compact}
Dla zamkniętej czterorozmaitości $X$, ustalonej struktury
$\mathfrak s$, metryki $g$ i gładkiej perturbacji $\eta$
przestrzeń moduli $\mathcal M_{\mathfrak s}(g,\eta)$
jest zwarta. Każda sekwencja gładkich rozwiązań ma po
zastosowaniu odpowiednich przekształceń cechowania
podciąg zbieżny w topologii $C^\infty$.
\end{twierdzenie}

\begin{proof}[Redukcja do standardowych oszacowań eliptycznych]
Poprzedni lemat daje jednolite ograniczenie
$\|\psi\|_{L^\infty}$. Z drugiego równania wynika
ograniczenie $\|F_A^+\|_{L^2}$.
Ponieważ
$[F_A/(2\pi i)]=c_1(L)$, liczba
$\int_X(F_A/i)\wedge(F_A/i)=(2\pi)^2\langle c_1(L)^2,[X]\rangle$
nie zależy od połączenia. Rozkład na części samodualną
i antysamodualną daje zatem także ograniczenie
$\|F_A^-\|_{L^2}$.

Po zapisaniu $A=A_0+a$ wybieramy cechowanie Coulomba
dla części połączeniowej, $\mathrm d^*a=0$, a część harmoniczną $a/i$ przenosimy do
ustalonego zwartego zbioru reprezentantów torusa
$H^1(X;\mathbb R)/(4\pi H^1(X;\mathbb Z))$,
gdzie $H^1(X;\mathbb Z)$ rozumiemy jako obraz w kohomologii
rzeczywistej. Czynnik $4\pi$ wynika z przyjętego działania
$A\mapsto A-2u^{-1}\mathrm du$: forma $u^{-1}\mathrm du/i$
ma okresy należące do $2\pi\mathbb Z$.
Jest to konieczne, gdyż sam warunek Coulomba
nie kontroluje składowej harmonicznej.
Oszacowanie eliptyczne dla $\mathrm d\oplus\mathrm d^*$
daje jednolitą kontrolę $W^{1,2}$ połączeń w takiej postaci.
Równanie Diraca i równanie krzywizny pozwalają następnie
na standardowe podnoszenie regularności Sobolewa,
połączone z osadzeniami Sobolewa i lokalną poprawą cechowania.
Uzyskujemy ograniczenia wszystkich norm $W^{m,2}$.
Z twierdzenia Rellicha oraz wyboru podciągu diagonalnego
dostajemy zbieżność $C^\infty$ po cechowaniu.

Ostatnie przejście --- pełna indukcja oszacowań
nieliniowych i regularność cechowań --- jest standardowym
wynikiem analizy eliptycznej i tutaj zostaje przyjęte.
Nie utożsamiamy przedstawionej redukcji z niezależnym
dowodem wszystkich tych twierdzeń analitycznych.
\end{proof}

\begin{uwaga}[Dlaczego nie ma zjawiska bubblingu?]
Jednolita kontrola spinora oraz abelowy charakter
połączenia pozwalają kontrolować pełną krzywiznę
i jej pochodne po ustaleniu cechowania. Nie pojawiają
się osobne koncentracje krzywizny odpowiadające
nieabelowym instantonom, dla których zwykle trzeba
uzupełniać przestrzeń moduli obiektami osobliwymi.
\end{uwaga}

\section{Transwersalność, orientacje i niezmienniczość}
\label{sec:tg-sw-transversality}

\begin{definicja}[Rozwiązanie regularne]
Nieredukowalne rozwiązanie $(A,\psi)$ nazywamy
\emph{regularnym}, jeżeli $H^2_{A,\psi}=0$.
W takim przypadku $H^1_{A,\psi}$ jest przestrzenią
styczną do lokalnej przestrzeni moduli i ma wymiar
$d(\mathfrak s)$.
\end{definicja}

\begin{twierdzenie}[Regularność generyczna; twierdzenie analityczne]
\label{thm:tg-sw-generic}
Niech $X$ będzie zamkniętą, spójną, zorientowaną
gładką czterorozmaitością z $b_2^+(X)>0$,
a $\mathfrak s$ ustaloną strukturą $\mathrm{Spin}^c$.
Dla odpowiedniego rezydualnego zbioru gładkich perturbacji
samodualnych $\eta$ nie występują rozwiązania redukowalne,
a wszystkie rozwiązania nieredukowalne są regularne.
Wówczas $\mathcal M_{\mathfrak s}(g,\eta)$ jest zwartą,
gładką rozmaitością wymiaru $d(\mathfrak s)$
(lub jest pusta, gdy $d(\mathfrak s)<0$).
\end{twierdzenie}

\begin{proof}[Mechanizm dowodu]
Rozważamy przestrzeń rozwiązań równań ze zmienną
perturbacją $\eta$ i odwzorowanie rzutowania na przestrzeń
perturbacji. Na nieredukowalnej części przestrzeni konfiguracji
działanie grupy cechowania jest wolne. Różniczka równań
łącznie z kierunkiem perturbacji ma składnik
$\delta\eta\mapsto-i\delta\eta$ w równaniu krzywizny.
Standardowy argument z jedyności kontynuacji dla operatora
Diraca i odpowiedniego lematu o surjektywności pokazuje,
że otrzymujemy gładki przekrój uniwersalny o właściwościach
Fredholma. Twierdzenie Sarda--Smale'a daje wtedy regularność
rozwiązań dla rezydualnego zbioru perturbacji.

Rozwiązanie redukowalne wymaga spełnienia równania w
$\mathcal H_g^+$, które narzuca $b_2^+(X)$ niezależnych
warunków rzeczywistych. Ponieważ $b_2^+>0$, także ten
zbiór perturbacji można generycznie ominąć.
Zwartość wynika z twierdzenia \ref{thm:tg-sw-compact}.
Kroki o strukturze Fredholma i transwersalności uniwersalnego
przekroju są tutaj przyjętymi wynikami analitycznymi.
\end{proof}

Orientację przestrzeni moduli określa orientacja
\emph{homologiczna}, tj. orientacja rzeczywistej
przestrzeni $H^1(X;\mathbb R)\oplus H^2_+(X;\mathbb R)$.
Jeżeli $b_1=0$, jest to po prostu wybór orientacji
$H^2_+(X;\mathbb R)$. Część kompleksowa operatora Diraca
ma orientację kanoniczną, a pozostała część jest orientowana
właśnie przez tę orientację homologiczną.

\begin{stwierdzenie}[Niezmienniczość przy braku redukowalnych]
\label{prop:tg-sw-cobordism}
Niech $b_1(X)=0$, $b_2^+(X)>1$ oraz
$d(\mathfrak s)=0$. Ustalmy orientację homologiczną.
Dla dowolnych dwóch generycznych perturbacji i metryk
$(g_0,\eta_0)$, $(g_1,\eta_1)$ otrzymujemy
$SW_{X,g_0,\eta_0}(\mathfrak s)
=SW_{X,g_1,\eta_1}(\mathfrak s)$.
\end{stwierdzenie}

\begin{proof}[Jednowymiarowy kobordyzm]
Łączymy dane generyczne gładką ścieżką
$t\mapsto(g_t,\eta_t)$. Warunek istnienia rozwiązania
redukowalnego ma kowymiar $b_2^+(X)>1$,
więc generyczna ścieżka jednoparametrowa go omija.
Parametryczna wersja twierdzenia Sarda--Smale'a
daje gładką, zorientowaną, zwartą przestrzeń moduli
wymiaru $d(\mathfrak s)+1=1$.
Jej brzeg stanowi rozłączną sumę
$\mathcal M_{\mathfrak s}(g_1,\eta_1)$ i przeciwnego
zbioru $\mathcal M_{\mathfrak s}(g_0,\eta_0)$.
Algebraiczna liczba punktów brzegowych zwartej,
zorientowanej jednowymiarowej rozmaitości wynosi zero;
stąd równość liczebności ze znakami.
Wykorzystane istnienie regularnej parametrycznej
przestrzeni moduli jest częścią teorii Fredholma.
\end{proof}

\section{Rozwiązania redukowalne, ściany i komory}
\label{sec:tg-sw-chambers}

Niech $\mathcal H_g^+\subset\Omega^2_+(X)$ oznacza
przestrzeń samodualnych form harmonicznych, a
$\Pi_g^+$ --- rzut $L^2$ na tę przestrzeń.
Dla każdej struktury $\mathfrak s$ z wiązką $L$
ustalmy gładkie połączenie $A_0$ na $L$.

\begin{stwierdzenie}[Równanie ściany]
\label{prop:tg-sw-wall}
Dla ustalonej metryki $g$ i samodualnej perturbacji
$\eta$ istnieje rozwiązanie redukowalne równań
Seiberga--Wittena wtedy i tylko wtedy, gdy
\[
\Pi_g^+\bigl(F_{A_0}/i-\eta\bigr)=0.
\]
Warunek ten nie zależy od wyboru $A_0$. Zbiór perturbacji
spełniających to równanie jest afiniczną podprzestrzenią
o kowymiarze $b_2^+(X)$ w przestrzeni perturbacji.
\end{stwierdzenie}

\begin{proof}
Dla konfiguracji redukowalnej mamy $\psi=0$, więc
$F_A^+=i\eta$. Każde inne połączenie na $L$ ma postać
$A=A_0+ia$ z rzeczywistą formą $a$, a jego krzywizna
wynosi $F_A=F_{A_0}+i\,\mathrm da$.
Szukamy zatem rozwiązania równania
$\mathrm d^+a=\eta-(F_{A_0}/i)^+$.
Na zamkniętej czterorozmaitości operator
$\mathrm d^+:\Omega^1\to\Omega^2_+$ ma domknięty obraz,
a jego dopełnieniem ortogonalnym jest dokładnie
$\mathcal H_g^+$. Jest to następstwo teorii Hodge'a
i eliptyczności operatora $\mathrm d^*\oplus\mathrm d^+$.
Równanie ma więc rozwiązanie dokładnie wtedy, gdy
prawa strona jest ortogonalna do $\mathcal H_g^+$,
czyli gdy zachodzi warunek z tezy.
Zmiana $A_0$ dodaje do krzywizny formę dokładną,
której rzut na $\mathcal H_g^+$ jest zerowy.
Ponieważ rzut $\eta\mapsto\Pi_g^+\eta$ jest surjektywny,
ściana ma kowymiar $\dim\mathcal H_g^+=b_2^+(X)$.
\end{proof}

Dla $b_2^+=1$ ustalmy jedną składową dodatniego stożka
$\{h\in H^2(X;\mathbb R):h^2>0\}$ i niech
$\omega_g$ będzie unormowaną harmoniczną formą
samodualną reprezentującą wybraną orientację $H_g^+$.
Wprowadzamy rzeczywistą funkcję
\[
\mathfrak w_{\mathfrak s}(g,\eta)
:=2\pi\langle c_1(L),[\omega_g]\rangle
-\int_X\eta\wedge\omega_g.
\]
Warunek $\mathfrak w_{\mathfrak s}(g,\eta)=0$ jest dokładnie
równaniem ściany. Komory odpowiadają znakom
$\mathfrak w_{\mathfrak s}>0$ i $\mathfrak w_{\mathfrak s}<0$
po ustaleniu wybranej składowej dodatniego stożka.

\begin{center}
\begin{tikzpicture}[scale=.93,>=Stealth]
  \fill[tgBlue!9] (-4,-1.25) rectangle (0,1.5);
  \fill[tgCoral!10] (0,-1.25) rectangle (4,1.5);
  \draw[tgNavy!55,->] (-4.2,0)--(4.25,0) node[right,font=\scriptsize] {parametr};
  \draw[tgPurple,line width=2.2pt] (0,-1.32)--(0,1.52);
  \node[font=\small,text=tgBlue] at (-2.4,1.02) {$\mathfrak w<0$};
  \node[font=\small,text=tgCoral] at (2.35,1.02) {$\mathfrak w>0$};
  \node[font=\small,text=tgPurple,rotate=90] at (.27,-.63) {ściana redukowalnych};
  \draw[tgTeal,-{Stealth[length=2.9mm]},line width=1.8pt]
     (-2.7,-.48)--(2.6,-.48);
  \node[font=\scriptsize,text=tgNavy] at (-2.7,-.75) {$(g_0,\eta_0)$};
  \node[font=\scriptsize,text=tgNavy] at (2.6,-.75) {$(g_1,\eta_1)$};
  \fill[tgTeal] (-2.7,-.48) circle (.065) (2.6,-.48) circle (.065);
\end{tikzpicture}

{\footnotesize Schemat dwóch komór przy $b_2^+=1$; ścieżka poprzeczna
przecina ścianę w jednym punkcie.}
\end{center}

\begin{twierdzenie}[Zmiana komory w wymiarze zero; wynik analityczny]
\label{thm:tg-sw-wallcross}
Niech $X$ będzie zamkniętą, spójną, zorientowaną,
gładką czterorozmaitością z $b_1(X)=0$ i $b_2^+(X)=1$.
Niech $\mathfrak s$ będzie strukturą $\mathrm{Spin}^c$
o $d(\mathfrak s)=0$. Po wyborze orientacji homologicznej
i spójnej konwencji koorientacji ściany istnieją dwie
liczby $SW_X^+(\mathfrak s)$ i $SW_X^-(\mathfrak s)$,
odpowiadające dwóm komorom, i spełniają one
\[
\bigl|SW_X^+(\mathfrak s)-SW_X^-(\mathfrak s)\bigr|=1.
\]
Znak różnicy zależy od konwencji orientacji oraz tego,
którą komorę oznaczono znakiem $+$.
\end{twierdzenie}

\begin{proof}[Mechanizm lokalny; użycie twierdzenia Kuranishiego]
Wybieramy generyczną jednoparametrową rodzinę danych
przecinającą ścianę poprzecznie dokładnie raz.
Poza ścianą przestrzeń parametryczna ma wymiar jeden.
W punkcie redukowalnym, przy $b_1=0$, jedynym ciągłym
stabilizatorem jest $S^1$. Z rachunku indeksu mamy
\[
0=d(\mathfrak s)=2\operatorname{ind}_{\mathbb C}D_A-2,
\qquad\operatorname{ind}_{\mathbb C}D_A=1.
\]
Dla generycznego przejścia i przy użyciu analitycznej
teorii transwersalności uzyskuje się w punkcie redukowalnym
jednowymiarowe zespolone jądro operatora Diraca oraz
zerowy cokernel. Nie jest to formalny wniosek z samej
równości indeksowej. Lokalna redukcja Lyapunova--Schmidta
i iloraz przez $S^1$ prowadzą do modelu o współrzędnych
$(\tau,z)\in\mathbb R\times\mathbb C$, w którym równanie
ma, z dokładnością do składników wyższego rzędu, postać
$\tau+\varepsilon|z|^2=0$, $\varepsilon\ne0$.
Po podzieleniu przez działanie $z\mapsto e^{i\theta}z$
zostaje półprosta $|z|^2\geq0$, przyłączona do punktu
redukowalnego. Podczas przejścia przez ścianę
pojawia się zatem lub znika jeden punkt nieredukowalny.
Liczenie z orientacjami daje zmianę równą $\pm1$.

Wykorzystany tutaj fakt, że rzeczywisty lokalny model
ma wskazaną postać z niezerowym współczynnikiem,
wymaga analitycznej teorii Kuranishiego i generyczności.
Nie wynika on z samego rachunku indeksu; jest głęboką
częścią twierdzenia o zmianie komory.
\end{proof}

\begin{przyklad}[\texorpdfstring{$\mathbb{CP}^2\#k\overline{\mathbb{CP}}^{\,2}$}{CP2}]
Dla $0\leq k\leq8$ rozważmy
$X_k=\mathbb{CP}^2\#k\overline{\mathbb{CP}}^{\,2}$
z bazą $H,E_1,\ldots,E_k$ taką, że $H^2=1$ i $E_i^2=-1$.
Mamy $e(X_k)=3+k$, $\sigma(X_k)=1-k$, $b_1=0$ i $b_2^+=1$.
Klasa charakterystyczna $c=3H-\sum_{i=1}^kE_i$ spełnia
\[
c^2=9-k=2e(X_k)+3\sigma(X_k),\qquad d(c)=0.
\]
Rozmaitość $X_k$ dopuszcza metrykę o dodatniej krzywiźnie
skalarnej (dla $k>0$ korzystamy z zachowywania tej własności
przez sumę spójną w wymiarze co najmniej trzy).
Przyjmijmy taką metrykę i perturbację $\eta=0$.
Z oszacowania Weitzenböcka wynika, że spinor każdego
rozwiązania musiałby być zerowy. Jednocześnie $c^2=9-k>0$,
a więc samodualny rzut harmoniczny klasy $c$ nie może
być zerowy: gdyby $c^+=0$, mielibyśmy $c^2\leq0$.
Równanie ściany wyklucza zatem także rozwiązania
redukowalne. W tej komorze przestrzeń moduli jest pusta
i niezmiennik wynosi zero. Po przejściu przez ścianę
niezmiennik dla tej samej struktury $\mathrm{Spin}^c$
ma wartość $\pm1$. Nie wolno więc mówić o jednym
niezmienniku Seiberga--Wittena bez określenia komory.
\end{przyklad}

\begin{uwaga}[Dlaczego ma to znaczenie przy wymiernej kontrakcji?]
Wymierna kontrakcja zachowuje $b_2^+$, więc gdy rozmaitość
początkowa ma $b_2^+=1$, rozmaitość wynikowa również ma
$b_2^+=1$. Formuły sklejenia dotyczą wówczas konkretnych
komór zgodnych z klasami użytymi w konstrukcji.
Porównanie z rozdmuchaniami $\mathbb{CP}^2$ musi być
przeprowadzone w tej samej konwencji, a sam brak niezerowych
niezmienników w jednej komorze nie dowodzi niczego
o wszystkich innych komorach.
\end{uwaga}

\section{Zastosowanie: granica dotychczasowych rachunków}
\label{sec:tg-sw-what-next}

Rozpatrzmy raz jeszcze przykład sekcji o samoprzecięciu
$-4$ na powierzchni eliptycznej $E(4)$. Wyjściowa rozmaitość
ma $b_2^+=7$, więc nie występuje zależność niezmiennika od komory.
Przy wymiernej kontrakcji otrzymaliśmy niezmienniki
$e=47$, $\sigma=-31$, a więc $2e+3\sigma=1$.
Dla struktur $\mathrm{Spin}^c$ o oczekiwanym wymiarze zero
koniecznie $c_1(L)^2=1$. Przejście od dwóch ocalałych
klas podstawowych do ich niezerowych wartości wymaga
formuły Fintushela--Sterna oraz danych o przedłużeniach
przez wymierną kulę homologii, nie samego rachunku indeksu.

W przykładach z $b_2^+=1$ dochodzi dodatkowo
wybór komory i weryfikacja poprzeczności rodzin danych.
Dopiero po skompletowaniu: (i) obliczenia grupy fundamentalnej,
(ii) klasyfikacji topologicznej, (iii) zgodnych przedłużeń
$\mathrm{Spin}^c$ i (iv) obliczenia niezmienników w komorze
można poprawnie dowodzić, że skonstruowana rozmaitość
jest homeomorficzna, lecz nie dyfeomorficzna ze standardowym
modelem.

\begin{uwaga}[Zakres dowodów w tej iteracji]
Zostały tutaj bezpośrednio udowodnione: opis stabilizatorów,
linearyzacja równań, równanie ściany, rachunek indeksu
pod warunkiem twierdzenia Atiyaha--Singera, zasada maksimum
dająca oszacowanie spinora oraz konsekwencja kobordyzmu
jednowymiarowego. Przyjęliśmy natomiast analityczne
twierdzenia o przekrojach Coulomba, właściwych własnościach
Fredholma, transwersalności Sarda--Smale'a, pełnej
regularności i zwartości eliptycznej oraz lokalnym modelu
Kuranishiego przy zmianie komory. Oddzielenie tych dwóch
poziomów stanowi część świadomie przyjętego rygoru wykładu.
\end{uwaga}


% ========================================================================
% ITERACJA 12. Umiescic po rozdziale „Analityczne podstawy teorii
% Seiberga--Wittena: przestrzenie moduli i komory”, przed rozdzialem
% o niezmiennikach wezlow. Nie wymaga nowej preambuly.
% ========================================================================
\chapter{Analiza Seiberga--Wittena: od lokalnego przekroju do niezmiennika}
\label{ch:tg-sw-analiza-uzup}

W poprzednim rozdziale opisaliśmy mechanizm konstrukcji przestrzeni moduli.
Pozostały jednak dwa rodzaje przejść, których nie należy utożsamiać:
rachunki liniowe na przestrzeniach spinorów oraz twierdzenia analityczne
zapewniające regularność, zwartość i istnienie ilorazu. W tym rozdziale
rozpiszemy przejścia między tymi poziomami. Wykorzystamy twierdzenia
o oszacowaniach eliptycznych, twierdzenie Sarda--Smale'a i zasadę
jednoznacznego przedłużenia, podając wyraźnie moment ich użycia.
Nie będziemy dowodzić od podstaw teorii eliptycznych operatorów Fredholma.

Przez cały rozdział $X$ oznacza \emph{zamkniętą, spójną, gładką,
zorientowaną czterorozmaitość}, $g$ --- ustaloną metrykę riemannowską,
a $\mathfrak s$ --- ustaloną strukturę $\mathrm{Spin}^c$.
Piszemy $L=\det S^+$, gdzie $S^\pm\to X$ są jej wiązkami spinorowymi.
Formy o wartościach w $i\mathbb R$ traktujemy jako rzeczywiste
przestrzenie wektorowe. Oznaczenia $\mathcal C_k$, $\mathcal G_{k+1}$
oraz $\mathfrak F_\eta$ zachowują znaczenie z poprzedniego rozdziału.

\section{Normalizacja mnożenia Clifforda i oszacowanie spinora}
\label{sec:tg-sw12-clifford}

Iloczyn hermitowski na $S^+$ przyjmujemy liniowy w pierwszym
argumencie. Mnożenie Clifforda na formach imaginarnych daje, po
wyborze standardowego przeskalowania, izomorfizm rzeczywistych
wiązek $\rho\colon i\Lambda^2_+\longrightarrow
\operatorname{Herm}_0(S^+)$, gdzie po prawej stronie stoją
bezśladowe operatory hermitowskie. Aby uniknąć ukrytych stałych,
ustalmy następującą konwencję.

\begin{definicja}[Kwadratowe odwzorowanie spinorowe]
\label{def:tg-sw12-q}
Dla $x\in X$ i $\psi\in S_x^+$ niech
$P_\psi(v):=\langle v,\psi\rangle\psi$ oraz
$P_\psi^0:=P_\psi-\tfrac12|\psi|^2\operatorname{id}_{S_x^+}$.
Definiujemy $q(\psi)\in i\Lambda^2_{+,x}$ równością
$\rho(q(\psi))=P_\psi^0$.
Wybieramy konwencję operatora Diraca tak, aby formuła
Weitzenböcka miała postać
\[
D_A^*D_A=\nabla_A^*\nabla_A+\frac{s_g}{4}
 +\frac12\rho(F_A^+),
\]
gdzie $s_g$ jest krzywizną skalarną. Równanie krzywizny
zapisujemy jako $F_A^+=q(\psi)+i\eta$, dla rzeczywistej
samodualnej formy $\eta$.
\end{definicja}

W innych podręcznikach $q$ może różnić się czynnikiem $2$
lub znakiem, wraz z odpowiadającą mu konwencją Clifforda.
Przy porównywaniu formuł z różnych źródeł należy porównywać
\emph{całe równanie}, a nie sam symbol $q$.

\begin{lemat}[Dwie tożsamości algebraiczne]
\label{lem:tg-sw12-q}
Dla każdego $x\in X$ i $\psi,\phi\in S_x^+$ zachodzą
\[
\langle\rho(q(\psi))\psi,\psi\rangle=\frac12|\psi|^4,
\qquad
\rho((\mathrm dq)_\psi\phi)
=(\phi\otimes\psi^*+\psi\otimes\phi^*)_0.
\]
W szczególności $q(\psi)=0$ wtedy i tylko wtedy, gdy $\psi=0$.
\end{lemat}

\begin{proof}
Macierz operatora $P_\psi$ w bazie ortonormalnej, której
pierwszy wektor jest równoległy do niezerowego $\psi$,
ma wartości własne $|\psi|^2$ i $0$. Po odjęciu połowy
śladu otrzymujemy wartości $|\psi|^2/2$ i $-|\psi|^2/2$.
Stąd $P_\psi^0\psi=(|\psi|^2/2)\psi$ i pierwsza
równość. W przypadku $\psi=0$ obie strony znikają.
Dla drugiej równości różniczkujemy
$P_{\psi+t\phi}$ względem $t\in\mathbb R$ w zerze:
\[
\frac{\mathrm d}{\mathrm dt}\bigg|_{t=0}P_{\psi+t\phi}
=\phi\otimes\psi^*+\psi\otimes\phi^*.
\]
Przejście do części bezśladowej komutuje z różniczkowaniem.
Ponieważ $\rho$ jest izomorfizmem, otrzymujemy tezę.
\end{proof}

\begin{twierdzenie}[Jawne oszacowanie typu Weitzenböcka]
\label{thm:tg-sw12-spinor-bound}
Niech $\eta\in\Omega^2_+(X)$ będzie gładka i niech $(A,\psi)$
będzie gładkim rozwiązaniem równań Seiberga--Wittena.
Niech $C_\rho>0$ spełnia
$\|\rho(i\omega)\|_{\mathrm{op}}\le C_\rho|\omega|$
dla każdej rzeczywistej formy samodualnej $\omega$.
Wtedy
\[
\|\psi\|_{L^\infty}^2\le
\|s_g^-\|_{L^\infty}+2C_\rho\|\eta\|_{L^\infty},
\qquad s_g^-:=\max\{0,-s_g\}.
\]
Jeżeli $s_g>0$ na całym $X$ i $\eta=0$, każde rozwiązanie
jest redukowalne.
\end{twierdzenie}

\begin{proof}
Z równania $D_A\psi=0$, tożsamości Weitzenböcka i
lematu \ref{lem:tg-sw12-q} otrzymujemy punktowo
\[
0=\frac12\Delta|\psi|^2+|\nabla_A\psi|^2
 +\frac{s_g}{4}|\psi|^2+\frac14|\psi|^4
 +\frac12\langle\rho(i\eta)\psi,\psi\rangle,
\]
gdzie $\Delta=\mathrm d^*\mathrm d$ na funkcjach.
W punkcie maksimum $|\psi|^2$ mamy $\Delta|\psi|^2\ge0$.
Odrzucając nieujemny wyraz $|\nabla_A\psi|^2$ oraz
szacując wyraz perturbacyjny z dołu przez
$-C_\rho|\eta||\psi|^2/2$, dostajemy
\[
\frac14|\psi|^4\le
\left(\frac{s_g^-}{4}+\frac{C_\rho}{2}|\eta|\right)
|\psi|^2.
\]
Jeżeli maksimum jest dodatnie, możemy podzielić przez
$|\psi|^2$; dla maksimum równego zero nierówność jest
oczywista. Gdy $s_g>0$ i $\eta=0$, w punkcie dodatniego
maksimum prawa strona pierwszej tożsamości byłaby dodatnia,
co daje sprzeczność. Zatem $\psi\equiv0$.
\end{proof}

\begin{przyklad}[Krzywizna dodatnia a niezmiennik]
Dla $X_k=\mathbb{CP}^2\#k\overline{\mathbb{CP}}^{\,2}$,
$0\le k\le8$, i $c=3H-\sum_{j=1}^kE_j$ mamy
$e(X_k)=3+k$, $\sigma(X_k)=1-k$ i $c^2=9-k>0$.
Metryka o dodatniej krzywiźnie skalarnej oraz $\eta=0$
wykluczają rozwiązania nieredukowalne. Rozwiązanie redukowalne
wymagałoby zarazem $c^+=0$, a wtedy $c^2\le0$.
Wobec tego w komorze zawierającej tę parę danych przestrzeń
moduli jest pusta. Nie jest to dowód zanikania niezmiennika
we wszystkich komorach, ponieważ $b_2^+(X_k)=1$.
\end{przyklad}

\section{Cechowanie: lokalny przekrój i globalna część harmoniczna}
\label{sec:tg-sw12-slice}

Lokalne ustalenie cechowania oraz globalne usunięcie części
dokładnej połączenia są różnymi operacjami. Pierwsza służy
budowaniu gładkiej struktury ilorazu w pobliżu rozwiązania
nieredukowalnego, druga --- dowodowi zwartości rodzin rozwiązań.

\begin{uwaga}[Zakres twierdzenia o przekroju]
Twierdzenie dotyczy konfiguracji nieredukowalnej oraz
małego otoczenia tożsamości w grupie cechowania.
W globalnym ilorazie występują tzw. duże transformacje
cechowania, klasyfikowane przez $H^1(X;\mathbb Z)$,
i nie można globalizować otrzymanej funkcji $f(\xi)$
bez dodatkowych warunków topologicznych.
Przy rozwiązaniach redukowalnych operator
$\mathrm d_0^*\mathrm d_0$ ma jądro funkcji stałych,
co uniemożliwia bezpośrednie zastosowanie tej samej
wersji twierdzenia o funkcji uwikłanej.
\end{uwaga}

\begin{twierdzenie}[Lokalny przekrój w punkcie nieredukowalnym]
\label{thm:tg-sw12-slice}
Ustalmy $k\ge3$ i gładką konfigurację
$\mathfrak c=(A,\psi)\in\mathcal C_k^*$.
Niech $\mathrm d_0$ oznacza różniczkę działania grupy
cechowania, zapisaną w poprzednim rozdziale, a
$\mathrm d_0^*$ --- jej formalne sprzężenie w rzeczywistym
iloczynie $L^2$. Istnieją otoczenie $U$ konfiguracji $\mathfrak c$ oraz
otoczenie $\mathcal U$ identyczności w grupie cechowania
takie, że dla każdej $\mathfrak c'\in U$ istnieje
dokładnie jedna transformacja $u\in\mathcal U$,
dla której $u\cdot\mathfrak c'\in
\mathfrak c+\ker\mathrm d_0^*$.
Jedyność dotyczy zatem małych transformacji
w ustalonym lokalnym modelu ilorazu.
\end{twierdzenie}

\begin{proof}
Dla $\mathfrak c=(A,\psi)$ ustalonej i małych
$\xi=(a,\phi)$ rozważamy funkcję
\[
\mathcal F(f,\xi):=\mathrm d_0^*
\bigl(e^{if}\cdot(\mathfrak c+\xi)-\mathfrak c\bigr),
\]
gdzie $f\in W^{k+1,2}(X;\mathbb R)$ jest mała, a $\xi$
ma klasę $W^{k,2}$. Działanie cechowania jest gładkie
na odpowiednich przestrzeniach Sobolewa, ponieważ $k\ge3$
i zachodzi właściwe twierdzenie o mnożeniu funkcji Sobolewa.
Mamy $\mathcal F(0,0)=0$, a różniczka względem $f$
w zerze jest równa $\mathrm d_0^*\mathrm d_0$.
Z lematu \ref{lem:tg-sw-full-coulomb} wiadomo, że
$\mathrm d_0^*\mathrm d_0=4\mathrm d^*\mathrm d+|\psi|^2$
ma zerowe jądro. Jest eliptyczny i samosprzężony o indeksie
zero, więc jest izomorfizmem
$W^{k+1,2}\to W^{k-1,2}$.
Twierdzenie o funkcji uwikłanej daje jedyną małą funkcję
$f=f(\xi)$ spełniającą $\mathcal F(f,\xi)=0$.
Jest to lokalna jednoznaczność, nie globalne twierdzenie
o wyborze cechowania; nie wyklucza transformacji dużych,
niehomotopijnych do jedynki.
\end{proof}

\begin{stwierdzenie}[Krata okresów grupy cechowania]
\label{prop:tg-sw12-lattice}
Niech $u\in C^\infty(X,S^1)$. Forma
$\alpha_u:=\frac{1}{2\pi i}u^{-1}\mathrm du$ jest zamknięta
i reprezentuje klasę w obrazie
$H^1(X;\mathbb Z)\to H^1(X;\mathbb R)$.
Dla $A=A_0+i a$, $a\in\Omega^1(X)$, transformacja $u$
zmienia część rzeczywistą $a$ o $-4\pi\alpha_u$.
W konsekwencji harmoniczna część $a$ jest określona
modulo kratę $4\pi H^1(X;\mathbb Z)$, rozumianą
w przestrzeni harmonicznych form stopnia jeden.
\end{stwierdzenie}

\begin{proof}
Ograniczenie funkcji $u$ do dowolnej gładkiej pętli
$\gamma$ ma całkowity stopień, więc
$\int_\gamma u^{-1}\mathrm du/(2\pi i)\in\mathbb Z$.
Klasa formy jest tym samym całkowita. Ze wzoru
$u\cdot A=A-2u^{-1}\mathrm du$ otrzymujemy
$u\cdot A=A_0+i(a-4\pi\alpha_u)$.
Każda całkowita klasa pierwszej kohomologii ma reprezentanta
postaci $\alpha_u$ dla pewnego odwzorowania $u:X\to S^1$,
bo $[X,S^1]\cong H^1(X;\mathbb Z)$. Po wybraniu
reprezentantów harmonicznych klasy daje to dokładnie
wskazaną kratę.
\end{proof}

\begin{przyklad}[Cechowanie harmoniczne na $S^1\times S^3$]
Niech $t\in\mathbb R/\mathbb Z$ będzie współrzędną na
$S^1$ i niech $X=S^1\times S^3$ z metryką produktową.
Dla $u_n(t,x)=e^{2\pi int}$, $n\in\mathbb Z$,
otrzymujemy $u_n^{-1}\mathrm du_n=2\pi in\,\mathrm dt$.
Połączenia $A_0+i\theta\,\mathrm dt$ i
$A_0+i(\theta-4\pi n)\,\mathrm dt$ są cechowaniowo
równoważne. Parametr harmoniczny należy więc do okręgu
$\mathbb R/(4\pi\mathbb Z)$, a nie do całej prostej.
Ten przykład dotyczy mechanizmu cechowania; sam w sobie
nie wyznacza żadnego niezmiennika Seiberga--Wittena.
\end{przyklad}

\begin{center}
\begin{tikzpicture}[>=Stealth,scale=.9]
  \fill[tgBlue!10] (-4.2,-1.3) rectangle (4.2,1.5);
  \draw[tgNavy!50,->] (-4.45,-.05)--(4.5,-.05);
  \draw[tgPurple!75,dashed,line width=1.3pt]
    (-3.55,-1.28)--(-3.55,1.45)
    (3.55,-1.28)--(3.55,1.45);
  \draw[tgTeal,line width=2.6pt] (-3.55,.48)--(3.55,.48);
  \draw[tgCoral,-{Stealth[length=2.8mm]},line width=2pt]
    (3.05,.72) .. controls (4.1,1.18) and (4.1,-.65) .. (3.18,-.84);
  \node[font=\small,text=tgNavy] at (0,1.06)
    {parametr harmoniczny $h$};
  \node[font=\scriptsize,text=tgPurple] at (-3.55,-1.52) {$h$};
  \node[font=\scriptsize,text=tgPurple] at (3.55,-1.52) {$h+4\pi\alpha$};
  \node[font=\scriptsize,text=tgCoral] at (2.0,-1.0)
    {utożsamienie przez cechowanie};
\end{tikzpicture}

{\footnotesize Schemat kraty okresów w przestrzeni części harmonicznych
połączenia. Pokazano jeden kierunek kraty; przy $b_1(X)>1$
otrzymujemy torus wyższego wymiaru.}
\end{center}

\section{Symbol kompleksu i indeks bez pomijania przejść}
\label{sec:tg-sw12-symbol}

\begin{lemat}[Dokładność symbolu kompleksu samodualnego]
\label{lem:tg-sw12-symbol}
Niech $V$ będzie zorientowaną euklidesową przestrzenią
rzeczywistą wymiaru cztery, a $\xi\in V^*\setminus\{0\}$.
Wówczas ciąg
\[
0\longrightarrow\mathbb R
\xrightarrow{f\mapsto f\xi}V^*
\xrightarrow{a\mapsto(\xi\wedge a)^+}\Lambda^2_+V^*
\longrightarrow0
\]
jest dokładny.
\end{lemat}

\begin{proof}
Normalizując $\xi$, możemy przyjąć $\xi=e^1$ w dodatnio
zorientowanej bazie ortonormalnej $e^1,\dots,e^4$.
Dla $a=\sum_{j=1}^4a_je^j$ dostajemy
\[
(e^1\wedge a)^+=\frac12\bigl[
 a_2(e^{12}+e^{34})+a_3(e^{13}-e^{24})
 +a_4(e^{14}+e^{23})\bigr].
\]
Trzy formy w nawiasach stanowią bazę $\Lambda^2_+V^*$.
Odwzorowanie drugie jest więc surjektywne i ma jądro
$\mathbb Re^1$, dokładnie równe obrazowi pierwszego.
\end{proof}

\begin{twierdzenie}[Eliptyczność operatora cechowania i równań]
\label{thm:tg-sw12-fredholm}
Dla rozwiązania $(A,\psi)$ równań Seiberga--Wittena operator
\[
\mathcal D_{A,\psi}:=\mathrm d_0^*\oplus\mathrm d_1:
\Omega^1(X;i\mathbb R)\oplus\Gamma(S^+)
\longrightarrow
\Omega^0(X;\mathbb R)\oplus\Gamma(S^-)
\oplus\Omega^2_+(X;i\mathbb R)
\]
jest eliptyczny. Po dopełnieniu w przestrzeniach Sobolewa
jest operatorem Fredholma o indeksie
$\operatorname{ind}_\mathbb R\mathcal D_{A,\psi}
=\bigl(c_1(L)^2-2e(X)-3\sigma(X)\bigr)/4$.
\end{twierdzenie}

\begin{proof}
Człony zależne od $\psi$ mają rząd zero, więc symbol
operatora jest sumą prostą symbolu
$\mathrm d^*\oplus\mathrm d^+$ oraz symbolu Diraca.
Lemat \ref{lem:tg-sw12-symbol} pokazuje, że pierwszy jest
izomorfizmem dla każdego niezerowego kowektora: jego jądro
ma wymiar zero, a dziedzina i przeciwdziedzina mają
wymiar cztery. Drugi symbol jest ograniczeniem mnożenia Clifforda
$\rho(\xi)\colon S^+\to S^-$. Jest odwracalny,
ponieważ $\rho(\xi)^*\rho(\xi)=
|\xi|^2\operatorname{id}_{S^+}$ dla $\xi\ne0$.
Wynika to z relacji Clifforda na całej wiązce spinorowej;
obie półwiązki mają ten sam wymiar zespolony.
Dostajemy eliptyczność.

Po odkształceniu członów rzędu zero indeks nie zmienia
się. Indeks rzeczywisty części $\mathrm d^*\oplus\mathrm d^+$
wynosi $b_1(X)-b_0(X)-b_2^+(X)=b_1-1-b_2^+$.
Indeks rzeczywisty Diraca jest dwukrotnością indeksu
zespolonego, a twierdzenie Atiyaha--Singera daje
$2\operatorname{ind}_\mathbb C D_A=(c_1(L)^2-\sigma)/4$.
Po dodaniu i użyciu $e=2-2b_1+b_2^++b_2^-$,
$\sigma=b_2^+-b_2^-$, otrzymujemy wzór z tezy.
Twierdzenie indeksowe jest tu wyraźnie przyjęte z literatury.
\end{proof}

\begin{przyklad}[Rachunek wymiaru na rozdmuchaniach $\mathbb{CP}^2$]
Niech $X_k=\mathbb{CP}^2\#k\overline{\mathbb{CP}}^{\,2}$
i $c=aH-\sum_{j=1}^k b_jE_j$, gdzie wszystkie liczby
$a,b_j$ są nieparzyste. Wtedy
\[
d(c)=\frac{a^2-\sum_{j=1}^k b_j^2-9+k}{4}.
\]
Dla $a=3$ i $b_j=1$ otrzymujemy $d(c)=0$ dla każdego $k$.
Równość wymiaru oczekiwanego zeru nie mówi jednak,
że istnieje rozwiązanie: przy metryce o dodatniej
krzywiźnie skalarnej i $\eta=0$ przestrzeń może być pusta.
\end{przyklad}

\section{Zwartość: od ograniczenia spinora do zbieżności konfiguracji}
\label{sec:tg-sw12-compact}

Największy skok w standardowych krótkich dowodach
zwartości polega na przejściu od ograniczenia $|\psi|$
do zwartości \emph{połączeń}. Rozpiszemy go przy użyciu
klasycznej teorii Hodge'a i oszacowań eliptycznych.

\begin{lemat}[Kontrola pełnej krzywizny]
\label{lem:tg-sw12-curvature}
Dla połączenia unitarnego $A$ na $L$ zachodzi
\[
\|F_A\|_{L^2}^2=2\|F_A^+\|_{L^2}^2
-4\pi^2\langle c_1(L)^2,[X]\rangle.
\]
W szczególności przy ustalonej klasie $c_1(L)$
i jednostajnie ograniczonej normie $\|F_A^+\|_{L^2}$
pełna krzywizna jest jednostajnie ograniczona w $L^2$.
\end{lemat}

\begin{proof}
Piszemy $F_A=i\Omega$, gdzie $\Omega$ jest rzeczywistą
zamkniętą formą. Z teorii Czern--Weila
$[\Omega/(2\pi)]=c_1(L)\otimes1$.
Rozkład Hodge'a daje
\[
\int_X\Omega\wedge\Omega
=\|\Omega^+\|_{L^2}^2-\|\Omega^-\|_{L^2}^2
=4\pi^2\langle c_1(L)^2,[X]\rangle.
\]
Dodając i odejmując kwadraty obu norm, otrzymujemy wzór.
\end{proof}

\begin{twierdzenie}[Zwartość rozwiązań przy ustalonych danych]
\label{thm:tg-sw12-compact}
Niech $g$, $\mathfrak s$ i gładka forma $\eta$ będą ustalone.
Dla dowolnego ciągu gładkich rozwiązań
$(A_j,\psi_j)$ równań Seiberga--Wittena istnieją podciąg
i gładkie transformacje cechowania $u_j$ takie, że
$u_j\cdot(A_j,\psi_j)$ zbiega w topologii $C^\infty$
do gładkiego rozwiązania. W szczególności przestrzeń
moduli $\mathcal M_{\mathfrak s}(g,\eta)$ jest zwarta.
\end{twierdzenie}

\begin{proof}[Pełny schemat z jawnym użyciem oszacowań eliptycznych]
\emph{Krok 1.} Z twierdzenia \ref{thm:tg-sw12-spinor-bound}
spinory $\psi_j$ są jednostajnie ograniczone w $L^\infty$.
Równanie $F_{A_j}^+=q(\psi_j)+i\eta$ daje zatem
jednostajne ograniczenie $F_{A_j}^+$ w $L^\infty$.
Lemat \ref{lem:tg-sw12-curvature} zapewnia dodatkowo
ograniczenie pełnej krzywizny w $L^2$.

\emph{Krok 2.} Wybieramy gładkie połączenie $A_0$ na $L$
i zapisujemy $A_j=A_0+i a_j$.
Rozkład Hodge'a formy rzeczywistej $a_j$ ma postać
$a_j=\mathrm df_j+\mathrm d^*\beta_j+h_j$.
Transformacja $e^{if_j/2}$ usuwa część dokładną.
Następnie transformacje odpowiadające klasom całkowitym
pozwalają wybrać części harmoniczne $h_j$ w ustalonej
zwartej dziedzinie fundamentalnej kraty z
wniosku \ref{prop:tg-sw12-lattice}. Otrzymujemy
$\mathrm d^*a_j=0$ i jednostajnie ograniczone $h_j$.

\emph{Krok 3.} Z równania krzywizny mamy
$\mathrm d^+a_j=(F_{A_j}^+-F_{A_0}^+)/i$.
Jego prawa strona jest ograniczona w każdym $L^p$,
$1<p<\infty$. Eliptyczne oszacowanie dla operatora
$\mathrm d^*\oplus\mathrm d^+$ daje
\[
\|a_j\|_{W^{1,p}}\le C_p
\bigl(\|\mathrm d^+a_j\|_{L^p}
+\|\mathrm d^*a_j\|_{L^p}+\|h_j\|_{L^p}\bigr)\le C'_p.
\]
Wybieramy $p>4$. Z zanurzenia Sobolewa
$W^{1,p}\hookrightarrow C^{0,\alpha}$ otrzymujemy
jednostajny moduł ciągłości form $a_j$.

\emph{Krok 4.} Względem połączenia odniesienia
$D_{A_0}\psi_j=-\tfrac12\rho(i a_j)\psi_j$.
Prawa strona jest ograniczona w $L^p$ dzięki
$\|a_j\|_{L^p}$ i $\|\psi_j\|_{L^\infty}$.
Z oszacowania eliptycznego dla operatora Diraca
wynika jednostajne ograniczenie $\psi_j$ w $W^{1,p}$.
Ponieważ $q$ jest wielomianem kwadratowym, z tego
ograniczenia wynika $q(\psi_j)$ ograniczone w $W^{1,p}$.
Stąd prawa strona równania dla $\mathrm d^+a_j$ jest
ograniczona w $W^{1,p}$ i eliptyczność
$\mathrm d^*\oplus\mathrm d^+$ daje $a_j$ ograniczone w
$W^{2,p}$. Powtarzając ten argument, uzyskujemy
ograniczenia wszystkich norm $W^{m,p}$ spinorów i połączeń.

\emph{Krok 5.} Z twierdzenia Rellicha i osadzeń Sobolewa
wydzielamy metodą przekątniową podciąg zbieżny w $C^m$
dla każdego $m$. Granica spełnia równania, ponieważ
operator Diraca, krzywizna i odwzorowanie $q$ są ciągłe
w tych topologiach. To dowodzi zwartości modulo cechowanie.

W krokach 3--5 korzystaliśmy z \emph{standardowych globalnych
oszacowań eliptycznych} i twierdzeń Sobolewa--Rellicha.
Ich pełne dowody nie należą do teorii grafów ani do
algebry spinorowej; argument wskazuje natomiast dokładnie,
które szacunki są potrzebne i dlaczego nie wystarczy samo
ograniczenie spinora.
\end{proof}

\begin{uwaga}[Dlaczego nie pojawia się koncentracja energii?]
W teorii połączeń o nieabelowej grupie strukturalnej
zbieżność może wymagać usunięcia punktów koncentracji.
Tutaj mamy grupę $U(1)$, równanie samodualnej krzywizny
bez wyrazu $A\wedge A$ oraz bezpośredni szacunek
$F_A^+=q(\psi)+i\eta$. Po ustaleniu cechowania jest to
równanie eliptyczne dla formy $a$, a jego oszacowania
dają zbieżność bez dodawania idealnych połączeń.
\end{uwaga}

\section{Transwersalność dla nieredukowalnych rozwiązań}
\label{sec:tg-sw12-transversality}

Zwartość nie przesądza, że przestrzeń moduli jest gładka.
Aby liczyć rozwiązania ze znakami, trzeba jeszcze wykazać
surjektywność linearyzacji, zwykle po generycznym wyborze
perturbacji.

\begin{lemat}[Cokernel uniwersalnej linearyzacji]
\label{lem:tg-sw12-surj}
Niech $(A,\psi)$ będzie gładkim rozwiązaniem nieredukowalnym.
Rozważmy linearyzację operatora równań, w której oprócz
zmiany $(A,\psi)$ dopuszczamy dowolną zmianę perturbacji
$\eta\in\Omega^2_+(X)$. Wtedy operator
\[
(a,\phi,\zeta)\longmapsto
\bigl(D_A\phi+\tfrac12\rho(a)\psi,\,
\mathrm d^+a-(\mathrm dq)_\psi\phi-i\zeta\bigr)
\]
ma gęsty obraz w przestrzeni gładkich wariacji równań
w sensie topologii $L^2$. Dla odpowiednio dobranych
uzupełnień Sobolewa operator uniwersalny (z nieskończenie
wymiarowym parametrem perturbacji) jest surjektywny,
jeśli wykorzystamy jego eliptyczną część i fakt, że
jej obraz jest domknięty modulo skończeniewymiarowy cokernel.
\end{lemat}

\begin{proof}
Przypuśćmy, że para $(\chi,\beta)$ jest ortogonalna
do obrazu. Testując dowolne $\zeta\in\Omega^2_+(X)$,
otrzymujemy $\beta=0$. Testując dowolne $\phi$,
otrzymujemy $D_A^*\chi=0$ w sensie dystrybucyjnym,
a z regularności eliptycznej $\chi$ jest gładka.
Testując wszystkie $a\in\Omega^1(X;i\mathbb R)$,
dostajemy punktowo
$\operatorname{Re}\langle\chi,\rho(a)\psi\rangle=0$
dla wszystkich $a$. W każdym punkcie, w którym $\psi\ne0$,
odwzorowanie rzeczywiste
$T_x^*X\to S_x^-$ dane przez $v\mapsto\rho(iv)\psi$
jest izomorfizmem (wynika to z relacji Clifforda i równości
wymiarów obu przestrzeni rzeczywistych, równych cztery).
Stąd $\chi=0$ na niepustym zbiorze otwartym
$\{\psi\ne0\}$. Dla operatora Diraca obowiązuje
zasada jednoznacznego przedłużenia: gładki spinor w jądrze
$D_A^*$, znikający na niepustym zbiorze otwartym,
znika na spójnej rozmaitości. Wynika zatem $\chi=0$.

Otrzymaliśmy brak niezerowej pary ortogonalnej do obrazu.
Aby uzyskać surjektywność, korzystamy z faktu, że
linearyzacja przy ustalonej perturbacji, po ustaleniu
cechowania, jest Fredholmowska: jej cokernel jest
skończeniewymiarowy. Wariacje perturbacji eliminują
wszystkie elementy tego cokernelu, co daje surjektywność
pełnej linearyzacji w wybranych przestrzeniach Sobolewa.
\end{proof}

\begin{twierdzenie}[Generyczna regularność; użycie Sarda--Smale'a]
\label{thm:tg-sw12-regular}
Załóżmy $b_2^+(X)>0$ i ustalmy metrykę $g$ oraz
strukturę $\mathfrak s$. Dla generycznej gładkiej
samodualnej perturbacji $\eta$ przestrzeń moduli
$\mathcal M_{\mathfrak s}(g,\eta)$ nie ma rozwiązań
redukowalnych, a wszystkie jej rozwiązania są regularne.
Jeśli jest niepusta, jest więc zwartą, gładką
zorientowalną rozmaitością rzeczywistą wymiaru
$d(\mathfrak s)$.
\end{twierdzenie}

\begin{proof}[Redukcja do twierdzeń analitycznych]
Z równania ściany (stwierdzenie \ref{prop:tg-sw-wall})
wynika, że perturbacje dopuszczające rozwiązanie
redukowalne tworzą afiniczną podprzestrzeń kowymiaru
$b_2^+(X)>0$. Poza nią wszystkie rozwiązania są
nieredukowalne. Lemat \ref{lem:tg-sw12-surj} zapewnia
surjektywność uniwersalnej linearyzacji przy takich
rozwiązaniach. Twierdzenie Sarda--Smale'a zastosowane
do projekcji uniwersalnej przestrzeni rozwiązań
na przestrzeń perturbacji daje generyczną regularność.
Twierdzenie \ref{thm:tg-sw12-compact} zapewnia zwartość.
Orientowalność wynika z opisu linii wyznacznikowej
w kolejnej sekcji. Pełnym narzędziem analitycznym jest
w tym miejscu parametryczne twierdzenie Sarda--Smale'a.
\end{proof}

\section{Orientacja przestrzeni moduli i zliczanie rozwiązań}
\label{sec:tg-sw12-orientation}

\begin{definicja}[Linia wyznacznikowa operatora Fredholma]
Niech $D:E\to F$ będzie rzeczywistym operatorem Fredholma
między przestrzeniami Banacha. Jego \emph{linią wyznacznikową}
nazywamy jednowymiarową rzeczywistą przestrzeń
$\det D:=\Lambda^{\max}\ker D\otimes
(\Lambda^{\max}\operatorname{coker}D)^*$.
Przez $\Lambda^{\max}0$ rozumiemy $\mathbb R$.
\end{definicja}

\begin{stwierdzenie}[Wybór orientacji homologicznej]
\label{prop:tg-sw12-hom-orient}
Po wybraniu orientacji rzeczywistej przestrzeni
$H^1(X;\mathbb R)\oplus H^2_+(X;\mathbb R)$
wszystkie regularne nieredukowalne przestrzenie moduli
Seiberga--Wittena na $X$ otrzymują zgodne orientacje.
\end{stwierdzenie}

\begin{proof}[Opis linii orientacji]
Operator deformacyjny $\mathcal D_{A,\psi}$ można odkształcić
przez człony rzędu zero do sumy
$(\mathrm d^*\oplus\mathrm d^+)\oplus D_A$.
Wyznacznik zespolonego operatora $D_A$ ma orientację
kanoniczną, ponieważ jądro i cokernel są zespolonymi
przestrzeniami wektorowymi. Linia wyznacznikowa
$\mathrm d^*\oplus\mathrm d^+$ odpowiada, przez teorię
Hodge'a, przestrzeniom
$H^1(X;\mathbb R)$, $H^0(X;\mathbb R)$ i $H^2_+(X;\mathbb R)$.
Czynnik $H^0$ ma orientację kanoniczną wyznaczoną przez
funkcję stałą $1$, a orientacja pozostałych dwóch
przestrzeni wyznacza orientację całej linii.
Ciągłość tego wyboru w rodzinach operatorów wynika
ze standardowej teorii wiązek wyznacznikowych.
\end{proof}

\begin{stwierdzenie}[Przeliczanie punktów moduli]
\label{prop:tg-sw12-count}
Załóżmy $b_1(X)=0$, $b_2^+(X)>0$ i $d(\mathfrak s)=0$.
Ustalmy orientację homologiczną oraz generyczne dane
$(g,\eta)$ bez rozwiązań redukowalnych. Wówczas
$\mathcal M_{\mathfrak s}(g,\eta)$ jest zbiorem skończonym,
a suma znaków jego punktów
$\sum_{[A,\psi]}\operatorname{sgn}([A,\psi])\in\mathbb Z$
jest dobrze określona. Przy $b_2^+>1$ nie zależy
od wyboru generycznych danych; przy $b_2^+=1$
należy dodatkowo ustalić komorę.
\end{stwierdzenie}

\begin{proof}
Z poprzedniego twierdzenia przestrzeń moduli jest
zwartą gładką rozmaitością wymiaru zero, a więc
skończonym zbiorem punktów. Ich znaki wyznacza
orientacja linii wyznacznikowej. Przy zmianie danych
zachodzą dwie możliwości: gdy jednoparametrowa rodzina
unika redukowalnych, liczbę punktów zachowuje
jednowymiarowy zorientowany kobordyzm, dokładnie jak
w poprzednim rozdziale. Dla $b_2^+>1$ generyczna
ścieżka omija ścianę kowymiaru co najmniej dwa.
Dla $b_2^+=1$ ściana ma kowymiar jeden i
liczba może się zmienić przy jej przekroczeniu.
\end{proof}

\section{Lokalny model ściany i rachunek zmiany komory}
\label{sec:tg-sw12-wall}

\begin{stwierdzenie}[Indeks w punkcie redukowalnym]
\label{prop:tg-sw12-wall-index}
Niech $b_1(X)=0$, $b_2^+(X)=1$, a
$d(\mathfrak s)=0$. Wtedy
$\operatorname{ind}_\mathbb C D_A=1$
dla każdego połączenia $A$ zgodnego z $\mathfrak s$.
\end{stwierdzenie}

\begin{proof}
Wzór na indeks operatora deformacyjnego ma w tym
przypadku postać
$0=d(\mathfrak s)=2\operatorname{ind}_\mathbb C D_A
+b_1-1-b_2^+=2\operatorname{ind}_\mathbb C D_A-2$.
Stąd żądana równość. Sam indeks nie mówi jeszcze,
że jądro Diraca ma wymiar zespolony jeden: wymagałoby
to także zanikania jego cokernelu.
\end{proof}

\begin{twierdzenie}[Lokalny mechanizm zmiany komory]
\label{thm:tg-sw12-wall-model}
Przy założeniach stwierdzenia
\ref{prop:tg-sw12-wall-index} wybierzmy generyczną
jednoparametrową rodzinę danych przecinającą ścianę
redukowalnych poprzecznie w jednym punkcie. Jeżeli
w punkcie tym operator Diraca jest surjektywny, to
jego jądro ma wymiar zespolony jeden. Sam warunek
przekroczenia ściany bez tego dodatkowego założenia
nie pozwala utożsamić jądra z $\mathbb C$. Przy dodatkowym
standardowym warunku niedegeneracji lokalnego modelu
Kuranishiego, po ilorazowaniu przez stabilizator $S^1$
lokalna przestrzeń moduli jest homeomorficzna z półprostą,
a na przejściu pojawia się lub znika jeden punkt
liczony ze znakiem. W konsekwencji
$|SW_X^+(\mathfrak s)-SW_X^-(\mathfrak s)|=1$.
\end{twierdzenie}

\begin{proof}[Redukcja lokalna z oznaczeniem użytego twierdzenia]
Operator Diraca ma indeks zespolony jeden; jego
surjektywność daje $\ker D_A\cong\mathbb C$.
Przy konfiguracji redukowalnej grupa stałych
transformacji cechowania $S^1$ działa na tym jądrze
przez mnożenie przez fazę. Parametr $\tau\in\mathbb R$
wybierzmy poprzecznie do ściany. Redukcja
Lyapunova--Schmidta, z uwzględnieniem
$S^1$-ekwiwariantności, prowadzi do równania
\[
\tau+\varepsilon|z|^2+R(\tau,z)=0,
\qquad z\in\mathbb C,\quad\varepsilon\ne0,
\]
gdzie $R$ ma wyższy rząd i nie zmienia lokalnego typu
zerozbioru. Po podstawieniu $r=|z|^2\ge0$ i ilorazowaniu
przez działanie faz otrzymujemy lokalną półprostą.
Jej punkt brzegowy odpowiada rozwiązaniu redukowalnemu.
W parametrycznej przestrzeni moduli wkład tego punktu
do brzegu wynosi $+1$ albo $-1$ w zależności od
wyboru koorientacji ściany i orientacji homologicznej.
Zliczenie brzegu daje żądaną różnicę.

Współczynnik $\varepsilon$ zależy od normalizacji
i wyboru lokalnych współrzędnych; jego niezerowość,
istnienie takiego równania oraz zachowanie orientacji
są głęboką częścią teorii
Kuranishiego oraz transwersalności rodzin; tutaj
przyjmujemy je wprost jako analityczne twierdzenie,
a nie jako wniosek z samej wartości indeksu.
\end{proof}

\begin{center}
\begin{tikzpicture}[>=Stealth,scale=.9]
  \draw[tgNavy!60,->] (-4.2,0)--(4.25,0)
       node[right,font=\scriptsize] {$\tau$};
  \draw[tgNavy!60,->] (0,-.15)--(0,2.5)
       node[above,font=\scriptsize] {$r=|z|^2$};
  \fill[tgBlue!12] (-4,-.05) rectangle (0,2.4);
  \draw[tgTeal,line width=2.6pt] (0,0)--(-2.9,2.2);
  \fill[tgCoral] (0,0) circle (.09);
  \node[text=tgCoral,font=\small,anchor=west] at (.23,.36)
       {punkt redukowalny};
  \node[text=tgTeal,font=\small] at (-2.55,2.55)
       {$\tau+\varepsilon r=0$};
  \node[text=tgNavy,font=\scriptsize] at (2.1,1.85)
       {brak pobliskich rozwiązań};
  \node[text=tgNavy,font=\scriptsize] at (2.1,1.58)
       {nieredukowalnych};
\end{tikzpicture}

{\footnotesize Lokalny model ilorazu przez $S^1$ przy przejściu przez
ścianę redukowalnych ($\varepsilon>0$). Po jednej stronie ściany
pojawia się rodzina nieredukowalnych rozwiązań w przekroju.}
\end{center}

\section{Co zostało dowiedzione, a co wymaga dalszej teorii?}
\label{sec:tg-sw12-scope}

W tej części dowiedliśmy jawnie tożsamości algebraicznych
dla $q$, dokładności symbolu kompleksu samodualnego,
kontroli pełnej krzywizny i wzorów na indeks
\emph{przy założeniu} twierdzenia Atiyaha--Singera.
Pokazaliśmy także krok po kroku, jak oszacowanie spinora,
Hodge'owskie ustalenie cechowania i bootstrap eliptyczny
prowadzą do zwartości. Dowód generycznej regularności
zredukowaliśmy do zasady jednoznacznego przedłużenia
i twierdzenia Sarda--Smale'a.

Nie udowodniliśmy w tym skrypcie od podstaw:
(i) pełnej teorii Fredholma w skali Sobolewa,
(ii) globalnych oszacowań eliptycznych i twierdzeń o
zwartości zanurzeń Sobolewa,
(iii) twierdzenia indeksowego Atiyaha--Singera,
(iv) parametrycznej wersji Sarda--Smale'a,
(v) konstrukcji i orientacji mapy Kuranishiego
w pobliżu punktu redukowalnego. Ich precyzyjne
sformułowania i dowody należą do analizy globalnej
oraz teorii cechowania.

\begin{uwaga}[Most do następnej iteracji]
Dla wykazania egzotyczności konkretnej czterorozmaitości
nie wystarczy skonstruować diagram Kirby'ego i obliczyć
$e$, $\sigma$ oraz formę przecięcia. Potrzebujemy dodatkowo
policzyć $\pi_1$ po zabiegach chirurgicznych,
ustalić topologiczny typ rozmaitości i wykazać, że
jej niezmiennik Seiberga--Wittena w odpowiedniej
komorze różni się od standardowego modelu.
Następna iteracja będzie poświęcona takiej konkretnej
konstrukcji, a nie kolejnym abstrakcyjnym definicjom.
\end{uwaga}


% ========================================================================
% Iteracja 13: konstrukcja Parka. Umiescic po analizie SW (iteracja 12),
% przed rozdzialem o nowoczesnych niezmiennikach wezlow.
% ========================================================================
\chapter{Egzotyczna czterorozmaitość Parka: konstrukcja i topologia}
\label{ch:tg-park}

Dotychczasowy aparat pozwala połączyć trzy pozornie odległe
konstrukcje: grafy plumbingowe, wymierną kontrakcję i teorię
czterorozmaitości symplektycznych. Celem rozdziału jest
\emph{skonstruowanie i zbadanie topologicznie} rozmaitości
$X_7$, której topologicznym modelem jest
$\mathbb{CP}^{2}\#7\overline{\mathbb{CP}}^{\,2}$.
Rozróżnienie jej struktury gładkiej od standardowej wymaga
niezmienników Seiberga--Wittena oraz wyboru komory; ten
ostatni argument przedstawimy w następnym rozdziale.

Wszystkie rozmaitości w tym rozdziale są zwarte i gładkie,
chyba że wyraźnie zaznaczono kategorię topologiczną.
Zakładamy orientacje zgodne z zespoloną orientacją
$\mathbb{CP}^{2}$ i standardową orientacją plumbingu.
Piszemy $\overline{\mathbb{CP}}^{\,2}$ dla przestrzeni
rzutowej z orientacją przeciwną.

\section{Racjonalna powierzchnia eliptyczna jako punkt wyjścia}
\label{sec:tg-park-e1}

\begin{definicja}[Powierzchnia eliptyczna $E(1)$]
Przez $E(1)$ oznaczamy rozdmuchanie dziewięciu punktów
bazowych odpowiednio wybranego pęku krzywych sześciennych
na $\mathbb{CP}^{2}$. W rezultacie dostajemy powierzchnię
zespoloną z odwzorowaniem na $\mathbb{CP}^{1}$, którego
włókno ogólne jest gładką krzywą eliptyczną. Topologicznie
$E(1)\cong\mathbb{CP}^{2}\#9\overline{\mathbb{CP}}^{\,2}$.
\end{definicja}

Wybieramy bazę $h,e_1,\ldots,e_9$ w drugiej homologii,
przy czym $h$ jest klasą prostej w $\mathbb{CP}^{2}$,
a $e_j$ klasą sfery wyjątkowej. Forma przecięcia jest
diagonalna: $h^2=1$, $e_j^2=-1$, $h\cdot e_j=0$
i $e_i\cdot e_j=0$ dla $i\ne j$.
Klasa ogólnego włókna eliptycznego wynosi
$f=3h-e_1-\cdots-e_9$.

\begin{lemat}
\label{lem:tg-park-fiber}
W opisanej bazie zachodzą $f^2=0$ oraz $f\cdot e_j=1$
dla $1\le j\le9$. Po czterech dalszych rozdmuchaniach
otrzymujemy $R:=E(1)\#4\overline{\mathbb{CP}}^{\,2}$,
o niezmiennikach
\[
e(R)=16,\qquad \sigma(R)=-12,\qquad
b_2^+(R)=1,\quad b_2^-(R)=13.
\]
\end{lemat}

\begin{proof}
Z postaci diagonalnej otrzymujemy
$f^2=9-\sum_{j=1}^{9}1=0$,
a $f\cdot e_j=-e_j\cdot e_j=1$.
Każde rozdmuchanie zwiększa charakterystykę Eulera
o jeden, zmniejsza sygnaturę o jeden i dodaje jedną
ujemną składową do formy przecięcia. Ponieważ
$e(\mathbb{CP}^{2})=3$ i
$\sigma(\mathbb{CP}^{2})=1$, dla $R$ mamy
$e(R)=3+13=16$ oraz $\sigma(R)=1-13=-12$.
Równości dla $b_2^\pm$ wynikają z diagonalizacji formy.
\end{proof}

\begin{przyklad}[Klasa sfery po czterech rozdmuchaniach]
\label{ex:tg-park-S}
Niech $e_{10},\ldots,e_{13}$ oznaczają klasy nowych
sfer wyjątkowych. Rozważmy klasę
\[
s:=4f+e_9-2(e_{10}+e_{11}+e_{12}+e_{13})\in H_2(R;\mathbb Z).
\]
Korzystając z $f^2=0$, $f\cdot e_9=1$,
$e_9^2=-1$ oraz z ortogonalności $f$ i $e_9$ do nowych klas
wyjątkowych, obliczamy
\[
s^2=8(f\cdot e_9)+e_9^2
   +4\sum_{j=10}^{13}e_j^2=8-1-16=-9.
\]
Jest to rachunek homologiczny, nie dowód istnienia
gładko zanurzonej sfery reprezentującej $s$.
Konstrukcja geometryczna takiej sfery korzysta z włókien
osobliwych pęku eliptycznego i rozwiązywania dodatnich
przecięć, jak wyjaśniamy niżej.
\end{przyklad}

\begin{center}
\begin{tikzpicture}[>=Stealth,line join=round,line cap=round]
  \node[draw=tgBlue,rounded corners=5pt,fill=tgBlue!9,
        minimum width=2.65cm,minimum height=1cm,align=center,text width=2.3cm,font=\small] (a)
        at (-4.6,0) {$\mathbb{CP}^{2}$\\pęk sześciennych};
  \node[draw=tgTeal,rounded corners=5pt,fill=tgTeal!10,
        minimum width=2.65cm,minimum height=1cm,align=center,text width=2.3cm,font=\small] (b)
        at (0,0) {$E(1)$\\włókna eliptyczne};
  \node[draw=tgPurple,rounded corners=5pt,fill=tgPurple!11,
        minimum width=2.65cm,minimum height=1cm,align=center,text width=2.3cm,font=\small] (c)
        at (4.6,0) {$R$\\$4$ nowe rozdmuchania};
  \draw[-{Stealth[length=2.8mm]},line width=1.65pt,tgTeal]
        (a)--(b);
  \draw[-{Stealth[length=2.8mm]},line width=1.65pt,tgPurple]
        (b)--(c);
\end{tikzpicture}

{\footnotesize Pierwsza strzałka: $9$ rozdmuchań punktów bazowych;
druga: $4$ rozdmuchania punktów podwójnych włókien osobliwych.
$E(1)=\mathbb{CP}^{2}\#9\overline{\mathbb{CP}}^{\,2}$ oraz
$R=E(1)\#4\overline{\mathbb{CP}}^{\,2}$.}
\end{center}

\section{Sześć sfer i graf plumbingowy \texorpdfstring{$C_7$}{C7}}
\label{sec:tg-park-c7}

\begin{definicja}[Konfiguracja $C_7$]
Niech $\Gamma_7$ będzie grafem łańcuchowym o sześciu
wierzchołkach, z wagami $(-9,-2,-2,-2,-2,-2)$.
Konfiguracją $C_7$ nazywamy zwartą czterorozmaitość
uzyskaną przez plumbing wiązek dyskowych zgodny z tym grafem.
Przez $S_0,\ldots,S_5$ oznaczamy sfery zerowe wiązek
dyskowych. Ich klasy tworzą bazę $H_2(C_7;\mathbb Z)$
i spełniają
\[
S_0^2=-9,\quad S_j^2=-2\ (1\le j\le5),\quad
S_j\cdot S_{j+1}=1\ (0\le j\le4),
\]
a pozostałe przecięcia różnych sfer są zerowe.
\end{definicja}

\begin{center}
\begin{tikzpicture}[x=1.23cm,y=1cm,line cap=round]
 \draw[tgTeal,line width=2.1pt] (0,0)--(5,0);
 \foreach \i in {0,...,5}{
   \ifnum\i=0
     \node[circle,draw=tgCoral,line width=1.7pt,fill=tgCoral!15,
       minimum size=10mm] (v\i) at (\i,0) {$-9$};
   \else
     \node[circle,draw=tgBlue,line width=1.4pt,fill=tgBlue!12,
       minimum size=9mm] (v\i) at (\i,0) {$-2$};
   \fi
   \node[font=\scriptsize,text=tgNavy] at (\i,-.74) {$S_{\i}$};
 }
 \node[font=\footnotesize,text=tgNavy] at (2.5,1.00)
   {graf ważony konfiguracji Parka $C_7$};
\end{tikzpicture}
\end{center}

\begin{twierdzenie}[Istnienie konfiguracji Parka; wynik zewnętrzny]
\label{thm:tg-park-embedding}
Istnieje wybór włókien osobliwych eliptycznego rozwarstwienia
$E(1)\to\mathbb{CP}^{1}$ oraz symplektyczne zanurzenie
$\iota\colon C_7\hookrightarrow R$ takie, że pięć sfer
$(-2)$ pochodzi z liniowego podgrafu osobliwego włókna
typu $\widetilde E_6$, a sfera $(-9)$ reprezentuje klasę
$s$ z przykładu \ref{ex:tg-park-S}. Ponadto dla
$V:=R\setminus\operatorname{int}\iota(C_7)$ zachodzi
$\pi_1(V)=1$.
\end{twierdzenie}

\begin{proof}[Konstrukcja i wskazanie nietrywialnego kroku]
W pracy J.~Parka, \emph{Simply connected symplectic
$4$-manifolds with $b_2^+=1$ and $c_1^2=2$},
Invent. Math. 159 (2005), 657--667, twierdzenie to
wynika z konstrukcji konfiguracji w Proposition~3.1 i
obliczenia grupy fundamentalnej w dowodzie Theorem~3.1.
Wybiera się cztery osobliwe włókna typu \emph{fishtail},
z których każde jest zanurzoną sferą z jednym dodatnim
punktem podwójnym, i rozdmuchuje ten punkt. Przekształcenie
właściwe takiego włókna ma klasę $f-2e_j$ i
samoprzecięcie $-4$; przecina sekcję $e_9$ dodatnio
w jednym punkcie. Po rozwiązaniu czterech przecięć z sekcją
powstaje jedna symplektyczna sfera w klasie
$s=4f+e_9-2\sum_{j=10}^{13}e_j$, której kwadrat
wynosi $-9$.

Dodatkowo włókno typu $\widetilde E_6$ zawiera łańcuch
pięciu sfer o samoprzecięciach $-2$. Ich przecięcia
z nową sferą dają dokładnie pokazany graf. Istnienie
odpowiednio dobranych włókien i ich symplektyczne
zanurzenie wymagają geometrii osobliwych rozwarstwień
eliptycznych, dlatego tego fragmentu nie zastępujemy
samym rachunkiem macierzy.

W dowodzie prostej spójności dopełnienia Park wskazuje
meridian konfiguracji reprezentujący generator
$\pi_1(Y_7)\cong\mathbb Z/49\mathbb Z$, który ogranicza
dysk w pozostałej części osobliwego włókna.
Wyjaśnijmy logiczny krok od tego faktu do tezy.
Ponieważ $R=V\cup_{Y_7} C_7$, obie części są spójne,
$\pi_1(R)=1$ i $\pi_1(C_7)=1$, twierdzenie van Kampena
mówi, że $\pi_1(V)$ jest normalnie generowana przez
obraz $\pi_1(Y_7)$. Skoro obraz generatora znika w
$\pi_1(V)$ dzięki wskazanemu dyskowi, to cały ten obraz
jest trywialny. Wynika stąd $\pi_1(V)=1$.
Istnienie wspomnianego dysku wymaga konkretnego układu
włókien z pracy Parka; ten geometryczny krok przyjmujemy
tutaj jako wynik zewnętrzny. Źródło: \url{https://doi.org/10.1007/s00222-004-0404-1}.
\end{proof}


\begin{lemat}[Jawne klasy pięciu sfer w $C_7$]
\label{lem:tg-park-S-classes}
W bazie $h,e_1,\ldots,e_{13}$ grupy $H_2(R;\mathbb Z)$
pięć sfer o kwadracie $-2$ można reprezentować klasami
\[
\begin{aligned}
 v_1&=e_4-e_7,& v_2&=e_1-e_4,&
 v_3&=h-e_1-e_2-e_3,\\
 v_4&=e_2-e_5,& v_5&=e_5-e_9.
\end{aligned}
\]
Ich kolejne przecięcia są równe $1$, a pozostałe przecięcia
różnych klas są zerowe. Klasa $s$ z przykładu
\ref{ex:tg-park-S} spełnia $s\cdot v_5=1$ i
$s\cdot v_j=0$ dla $1\le j\le4$.
\end{lemat}

\begin{proof}
Przez bezpośrednie podstawienie do diagonalnej formy przecięcia
obliczamy $v_j^2=-2$ oraz
$v_i\cdot v_{i+1}=1$ dla $1\le i\le4$.
Jeżeli $|i-j|>1$, odpowiednie klasy są ortogonalne.
Ponieważ $f\cdot v_j=0$, nowe klasy wyjątkowe
$e_{10},\ldots,e_{13}$ są ortogonalne do $v_j$, a
$e_9\cdot v_j=\delta_{5j}$, mamy
$s\cdot v_j=\delta_{5j}$.
Wynik nie zastępuje geometrycznego twierdzenia o realizacji
klas przez symplektyczne sfery; to pozostaje zewnętrznym
rezultatem Parka.
\end{proof}

\begin{uwaga}[Co rzeczywiście przedstawia diagram?]
Graf plumbingowy koduje samoprzecięcia i wzajemne
przecięcia sześciu sfer. \emph{Nie jest} pełnym diagramem
Kirby'ego rozmaitości $R$. Aby uzyskać pełny diagram,
trzeba również przedstawić uchwyty opisujące dopełnienie
$V$ oraz ich mapy przyczepienia. W tej iteracji rysujemy
poprawny lokalny diagram uchwytowy $C_7$, nie twierdząc,
że sam stanowi diagram całej konstrukcji Parka.
\end{uwaga}

\section{Macierz przecięcia i brzeg \texorpdfstring{$C_7$}{C7}}
\label{sec:tg-park-boundary}

\begin{stwierdzenie}[Macierz, określoność i wyznacznik]
\label{prop:tg-park-Q}
Macierz przecięcia $C_7$ w bazie $[S_0],\ldots,[S_5]$
ma postać
\[
Q_7=\begin{pmatrix}
-9&1&0&0&0&0\\
1&-2&1&0&0&0\\
0&1&-2&1&0&0\\
0&0&1&-2&1&0\\
0&0&0&1&-2&1\\
0&0&0&0&1&-2
\end{pmatrix}.
\]
Jest ujemnie określona, a jej wyznacznik wynosi $49$.
\end{stwierdzenie}

\begin{proof}
Wyraz diagonalny macierzy jest samoprzecięciem
odpowiedniej sfery; wyraz pozadiagonalny jest liczbą
przecięć dwóch sfer. Stąd od razu postać $Q_7$.

Niech $d_k$ będzie wyznacznikiem górnego lewego bloku
$k\times k$, z konwencją $d_0=1$. Wówczas
$d_1=-9$ i, przez rozwinięcie względem ostatniego wiersza,
$d_k=-2d_{k-1}-d_{k-2}$ dla $2\le k\le6$.
Indukcja daje $d_k=(-1)^k(8k+1)$: dla $k=0,1$
jest to prawda, a krok wynika z równości
$2(8k-7)-(8k-15)=8k+1$.
Dla $k=6$ otrzymujemy $\det Q_7=49$.

Ponadto $(-1)^k d_k=8k+1>0$ dla wszystkich $k$.
Z kryterium Sylvestera zastosowanego do $-Q_7$
wynika dodatnia określoność $-Q_7$, czyli ujemna
określoność $Q_7$.
\end{proof}

\begin{stwierdzenie}[Homologia brzegu]
\label{prop:tg-park-h1-boundary}
Dla $Y_7:=\partial C_7$ zachodzi
$H_1(Y_7;\mathbb Z)\cong\mathbb Z/49\mathbb Z$.
\end{stwierdzenie}

\begin{proof}
Ponieważ $C_7$ powstała z jednej $0$-komórki i sześciu
uchwytów indeksu dwa, jest po prostu spójna, a
$H_1(C_7)=0$, $H_2(C_7)\cong\mathbb Z^6$. Z ciągu
dokładnego pary $(C_7,Y_7)$ i dualności Poincarégo--Lefschetza
otrzymujemy
\[
H_2(C_7;\mathbb Z)\xrightarrow{Q_7}
H_2(C_7,Y_7;\mathbb Z)\longrightarrow
H_1(Y_7;\mathbb Z)\longrightarrow0.
\]
Zatem $H_1(Y_7)\cong\operatorname{coker}Q_7$.
Aby ustalić nie tylko rząd, lecz także strukturę grupy,
niech $x_0,\ldots,x_5$ odpowiadają generatorom prezentacji
przez kolumny $Q_7$. Ostatnie pięć relacji daje kolejno
$x_4=2x_5$, $x_3=3x_5$, $x_2=4x_5$,
$x_1=5x_5$ i $x_0=6x_5$.
Pierwsza relacja $-9x_0+x_1=0$ daje
$(-54+5)x_5=-49x_5=0$.
Wszystkie generatory są więc wielokrotnościami $x_5$,
a rząd grupy jest $|\det Q_7|=49$. Stąd grupa jest
cykliczna i ma podaną postać.
\end{proof}

\begin{stwierdzenie}[Wypełnienie soczewkowe]
\label{prop:tg-park-lens}
W konwencji, w której ujemny łańcuch $(-a_1,\ldots,-a_r)$
przedstawia chirurgię o współczynniku
$- [a_1,\ldots,a_r]^-$ na węźle trywialnym,
brzeg $C_7$ jest zorientowanie dyfeomorficzny z
$L(49,-6)$, gdzie
$[a_1,\ldots,a_r]^-=a_1-1/(a_2-\cdots-1/a_r)$.
\end{stwierdzenie}

\begin{proof}
Pięć końcowych dwójek redukuje się do ułamka $6/5$.
Stąd $[9,2,2,2,2,2]^-=9-5/6=49/6$.
Kolejne zastosowanie operacji \emph{slam-dunk}
(rachunku chirurgicznego torusów) zastępuje sześć
sprzężonych składowych pojedynczą wymierną chirurgią
o współczynniku $-49/6$. W konwencji chirurgii Dehna
przyjętej w skrypcie jest to przestrzeń soczewkowa
$L(49,-6)$. Samo obliczenie wyznacznika $49$ nie
wyznaczałoby parametru $-6$ ani orientowanego typu
rozmaitości, dlatego potrzebny jest rachunek ułamka.
\end{proof}

\begin{przyklad}[Forma splecenia: generator i orientacja]
Klasa meridianu przy składowej o ramowaniu $-9$
jest generatorem $H_1(Y_7)$. Rzeczywiście, lewy górny
minor macierzy $Q_7$ ma wyznacznik $-6$ (jest to macierz
łańcucha pięciu sfer o wagach $-2$). Wzór na odwrotność
macierzy daje więc $(Q_7^{-1})_{00}=-6/49$.
Ponieważ $6$ i $49$ są względnie pierwsze, element
takiej formy ma rząd $49$ w $\mathbb Q/\mathbb Z$;
stąd sama klasa meridianu ma rząd $49$.
Forma splecenia brzegu, z konwencją orientacji zgodną
z równaniem $\lambda_{Y_7}=Q_7^{-1}\pmod{\mathbb Z}$,
spełnia $\lambda_{Y_7}([\mu_0],[\mu_0])=-6/49$
w $\mathbb Q/\mathbb Z$. Przy odwróceniu orientacji
brzegu znak formy ulega zmianie. W literaturze spotyka
się też konwencję $-Q_7^{-1}$; nie należy mieszać
jej ze znakiem współczynnika chirurgii.
\end{przyklad}

\begin{center}
\begin{tikzpicture}[x=1.08cm,y=.95cm,>=Stealth,line cap=round]
  \foreach \i in {0,...,5}{
    \ifnum\i=0
      \draw[tgCoral,line width=1.9pt] (\i,0) circle (.32);
      \node[tgCoral,font=\scriptsize] at (\i,-.68) {$-9$};
    \else
      \draw[tgBlue,line width=1.9pt] (\i,0) circle (.32);
      \node[tgBlue,font=\scriptsize] at (\i,-.68) {$-2$};
    \fi
  }
  \foreach \i [evaluate=\i as \j using int(\i+1)] in {0,...,4}{
    \draw[tgTeal,line width=1.65pt] (\i+.32,.23)
          .. controls (\i+.52,.74) and (\j-.52,.74) .. (\j-.32,.23);
    \draw[tgTeal,line width=1.65pt] (\i+.32,-.23)
          .. controls (\i+.52,-.74) and (\j-.52,-.74) .. (\j-.32,-.23);
  }
  \node[font=\footnotesize,text=tgNavy] at (2.5,1.28)
       {Schemat przyległości składowych ramowania};
  \node[font=\footnotesize,text=tgNavy] at (2.5,-1.14)
       {Współczynnik ramowania jest podany pod składową};
\end{tikzpicture}

{\footnotesize Poglądowy schemat ramowań i sprzężeń: liczby pod okręgami
mają oznaczać ramowania, a połączenie sąsiadów ma
reprezentować liczbę splecenia $1$. Nie jest to regularny
diagram węzłowy z określonymi przejściami nad i pod;
ścisły diagram Kirby'ego wymaga narysowania tych skrzyżowań.}
\end{center}

\section{Zastąpienie \texorpdfstring{$C_7$}{C7} wymierną kulą homologii}
\label{sec:tg-park-blowdown}

\begin{twierdzenie}[Casson--Harer; wynik zewnętrzny]
\label{thm:tg-park-B7}
Istnieje gładka, zwarta, orientowana czterorozmaitość
$B_7$ o brzegu $\partial B_7\cong L(49,-6)$, dla której
$H_*(B_7;\mathbb Q)\cong H_*(D^4;\mathbb Q)$ oraz
$\pi_1(B_7)\cong\mathbb Z/7\mathbb Z$.
Naturalne odwzorowanie
$\pi_1(\partial B_7)\to\pi_1(B_7)$ jest surjektywne.
\end{twierdzenie}

\begin{definicja}[Rozmaitość $X_7$]
Niech $V:=R\setminus\operatorname{int}\iota(C_7)$,
gdzie $\iota$ pochodzi z twierdzenia
\ref{thm:tg-park-embedding}. Wybierzmy
dyfeomorfizm $\partial B_7\to\partial V$ odwracający
orientacje indukowane na obu brzegach.
Definiujemy
\[
X_7:=V\cup_{\partial}B_7.
\]
Jest to \emph{wymierna kontrakcja konfiguracji $C_7$}
w powierzchni $R$.
\end{definicja}

\begin{center}
\begin{tikzpicture}[>=Stealth,line join=round,line cap=round]
  \node[draw=tgBlue,rounded corners=5pt,fill=tgBlue!9,
    minimum width=3.15cm,minimum height=1cm,align=center] (r)
     at (-3.3,1.08) {$R=V\cup C_7$\\$e=16,\ \sigma=-12$};
  \node[draw=tgCoral,rounded corners=5pt,fill=tgCoral!10,
    minimum width=2.8cm,minimum height=1cm,align=center] (c)
     at (1.85,1.08) {$C_7$\\$e=7,\ \sigma=-6$};
  \node[draw=tgTeal,rounded corners=5pt,fill=tgTeal!10,
    minimum width=2.8cm,minimum height=1cm,align=center] (b)
     at (1.85,-1.18) {$B_7$\\$e=1,\ \sigma=0$};
  \node[draw=tgPurple,rounded corners=5pt,fill=tgPurple!10,
    minimum width=3.0cm,minimum height=1cm,align=center] (x)
     at (-3.3,-1.18) {$X_7=V\cup B_7$\\$e=10,\ \sigma=-6$};
  \draw[-{Stealth[length=3mm]},line width=1.6pt,tgCoral]
        (r)-- node[above,font=\scriptsize] {usuń $C_7$} (c);
  \draw[-{Stealth[length=3mm]},line width=1.6pt,tgTeal]
        (c)-- node[right,font=\scriptsize,align=left]
        {zastąp przez\\$B_7$} (b);
  \draw[-{Stealth[length=3mm]},line width=1.6pt,tgPurple]
        (b)-- node[below,font=\scriptsize] {sklej z $V$} (x);
\end{tikzpicture}

{\footnotesize Schemat wymiernej kontrakcji. Równość brzegów umożliwia
operację wycięcia i wklejenia; sam schemat nie dowodzi
zachowania prostej spójności.}
\end{center}

\begin{twierdzenie}[Grupa fundamentalna wyniku]
\label{thm:tg-park-pi1}
Rozmaitość $X_7$ jest po prostu spójna.
\end{twierdzenie}

\begin{proof}
Rozważmy rozkład $X_7=V\cup B_7$,
którego część wspólna jest kołnierzem spójnej
przestrzeni soczewkowej $Y_7$.
Twierdzenie van Kampena daje
\[
\pi_1(X_7)\cong
\pi_1(V)*_{\pi_1(Y_7)}\pi_1(B_7).
\]
Z twierdzenia \ref{thm:tg-park-embedding} mamy
$\pi_1(V)=1$. Zatem prawa strona jest ilorazem
$\pi_1(B_7)$ przez domknięcie normalne obrazu
$\pi_1(Y_7)$. Z twierdzenia \ref{thm:tg-park-B7}
odwzorowanie $\pi_1(Y_7)\to\pi_1(B_7)$ jest
surjektywne, a więc iloraz jest trywialny.
\end{proof}

\begin{twierdzenie}[Charakterystyki liczbowe]
\label{thm:tg-park-numbers}
Zachodzą
\[
e(X_7)=10,\quad \sigma(X_7)=-6,\quad
b_1(X_7)=b_3(X_7)=0,\quad
b_2^+(X_7)=1,\quad b_2^-(X_7)=7.
\]
Ponadto $H_2(X_7;\mathbb Z)\cong\mathbb Z^8$.
\end{twierdzenie}

\begin{proof}
$C_7$ powstaje z jednego uchwytu indeksu zero
i sześciu uchwytów indeksu dwa. Wobec tego
$e(C_7)=1+6=7$; nie jest to $8$, ponieważ
liczba wierzchołków grafu liczy uchwyty indeksu
dwa, a nie wszystkie komórki. Jej forma jest
ujemnie określona rzędu sześć, więc
$\sigma(C_7)=-6$.
Wymierna kula homologii ma $e(B_7)=1$ i
$\sigma(B_7)=0$.
Ponieważ brzeg $Y_7$ ma zerową charakterystykę
Eulera, z addytywności otrzymujemy
\[
e(X_7)=e(R)-e(C_7)+e(B_7)=16-7+1=10.
\]
Z addytywności sygnatury (twierdzenie Novikova)
przy rozkładzie wzdłuż wymiernej sfery homologicznej
wynika $\sigma(X_7)=\sigma(R)-\sigma(C_7)=-12+6=-6$.
Prosta spójność daje $H_1(X_7;\mathbb Z)=0$, więc
$b_1=0$. Dualność Poincarégo daje $b_3=0$.
Zatem $e=2+b_2$, stąd $b_2(X_7)=8$.
Rozwiązanie równań $b_2^++b_2^-=8$ i
$b_2^+-b_2^-=-6$ prowadzi do $(b_2^+,b_2^-)=(1,7)$.
Wreszcie uniwersalne twierdzenie o współczynnikach
i dualność Poincarégo dają brak torsji w
$H_2(X_7;\mathbb Z)$ dla po prostu spójnej zamkniętej
orientowanej czterorozmaitości. Skoro rząd jest osiem,
$H_2(X_7;\mathbb Z)\cong\mathbb Z^8$.
\end{proof}

\begin{przyklad}[Kontrola wzoru $c_1^2$]
Twierdzenie Symington o symplektycznej wymiernej kontrakcji
zapewnia, że $X_7$ ma strukturę symplektyczną, ponieważ
konfiguracja Parka składa się z symplektycznych sfer
przecinających się dodatnio. Dla zgodnej struktury prawie
zespolonej na takiej zamkniętej czterorozmaitości zachodzi
$c_1^2=2e+3\sigma$. Dla $X_7$ otrzymujemy
$c_1^2(X_7)=20-18=2$.
Zgadza się to z twierdzeniem Parka. Sama równość
$c_1^2=2$ nie rozróżnia jednak struktur gładkich:
tę samą liczbę ma standardowa racjonalna powierzchnia
$\mathbb{CP}^{2}\#7\overline{\mathbb{CP}}^{\,2}$.
\end{przyklad}

\section{Forma przecięcia i typ homeomorfizmu}
\label{sec:tg-park-homeo}

\begin{lemat}[Nieparzystość formy przecięcia]
\label{lem:tg-park-odd}
Forma przecięcia $X_7$ jest nieparzysta.
\end{lemat}

\begin{proof}
Gdyby forma przecięcia $Q_{X_7}$ była parzysta,
rozmaitość $X_7$ byłaby spin: dla po prostu spójnej
orientowanej czterorozmaitości warunek
$w_2(X_7)=0$ jest równoważny parzystości formy.
Twierdzenie Rochlina mówi, że sygnatura zamkniętej
gładkiej czterorozmaitości spin jest podzielna przez $16$.
Tymczasem $\sigma(X_7)=-6$, sprzeczność.
Wynika stąd nieparzystość formy.
\end{proof}

\begin{twierdzenie}[Topologiczny typ $X_7$]
\label{thm:tg-park-homeo}
Istnieje zachowujący orientację homeomorfizm
\[
X_7\cong_{\mathrm{TOP}}
\mathbb{CP}^{2}\#7\overline{\mathbb{CP}}^{\,2}.
\]
\end{twierdzenie}

\begin{proof}
Dla każdej zamkniętej, spójnej, orientowanej, po prostu
spójnej czterorozmaitości forma przecięcia na
$H_2/\operatorname{Tor}$ jest unimodularna.
Twierdzenie \ref{thm:tg-park-numbers} i lemat
\ref{lem:tg-park-odd} dają więc nieparzystą,
nieokreśloną formę unimodularną o sygnaturze $-6$
i rzędzie osiem. Klasyfikacja nieparzystych
nieokreślonych form unimodularnych nad $\mathbb Z$
implikuje izometrię z formą
$\langle1\rangle\oplus7\langle-1\rangle$.

Aby przejść od formy do rozmaitości, stosujemy twierdzenie
klasyfikacyjne Freedmana dla zamkniętych, po prostu
spójnych topologicznych czterorozmaitości.
Obie rozmaitości są wygładzalne, zatem ich niezmiennik
Kirby'ego--Siebenmanna jest zerowy. Dla ustalonej
nieparzystej formy i tego niezmiennika dostajemy
homeomorfizm zachowujący orientację.
Twierdzenie Freedmana przyjmujemy tu jako wynik
zewnętrzny, natomiast identyfikację samej formy
wyprowadziliśmy z poprzednich rachunków.
\end{proof}

\begin{uwaga}[Czego jeszcze nie dowiedliśmy]
W dotychczasowych rachunkach uzyskaliśmy homeomorfizm,
nie zaś dyfeomorfizm. Istnienie formy przecięcia
$\langle1\rangle\oplus7\langle-1\rangle$ nie pozwala
rozpoznać gładkiej struktury. Należy porównać
niezmienniki Seiberga--Wittena w odpowiedniej komorze
przy $b_2^+=1$, wykorzystując symplektyczność $X_7$,
odpowiednie twierdzenie Taubesa--Li--Liu i nierówność
dodatniego parowania klasy kanonicznej z klasą formy
symplektycznej. W następnym rozdziale przeprowadzimy ten argument.
\end{uwaga}

\section{Kontrolny rachunek macierzy i granice metody}
\label{sec:tg-park-summary}

\begin{przyklad}[Rekurencja ułamka łańcuchowego]
Oznaczmy przez $r_j$ ujemny ułamek łańcuchowy
złożony z $j$ dwójek: $r_1=2$ oraz
$r_{j+1}=2-1/r_j$. Indukcja daje
$r_j=(j+1)/j$. Stąd dla pięciu dwójek
otrzymujemy $r_5=6/5$, a więc
$[9,2,2,2,2,2]^-=9-1/r_5=49/6$.
To drugi, niezależny od rachunku wyznacznika sposób
ustalenia liczby $49$ i parametru chirurgii $-6$.
\end{przyklad}

\begin{przyklad}[Bilans liczb Betti'ego]
Rozmaitość $R$ ma $b_2^-=13$. Wycięta konfiguracja
$C_7$ wnosi sześć ujemnych kierunków, a $B_7$ ma
zerową drugą homologię wymierną. Zatem $b_2^-$
spada do $7$, podczas gdy $b_2^+=1$ pozostaje
bez zmian. Zgadza się to z obliczeniem przez
$(e,\sigma)$: w obu metodach otrzymujemy $b_2=8$.
\end{przyklad}

\begin{uwaga}[Zakres dowodu a konstrukcja globalna]
Z rozdziału można samodzielnie odtworzyć wszystkie
rachunki na grafie $C_7$, obliczenie brzegu, zastosowanie
van Kampena \emph{przy znanej prostej spójności dopełnienia},
liczby Betti'ego oraz formę przecięcia i typ TOP.
Zaawansowanymi danymi wejściowymi pozostają
(1) istnienie konkretnego symplektycznego zanurzenia
$C_7\hookrightarrow R$ i dowód $\pi_1(V)=1$ z pracy Parka,
(2) istnienie wymiernej kuli homologii $B_7$ i jej
własności z konstrukcji Cassona--Harera,
(3) twierdzenie Freedmana.

W kolejnym rozdziale do tej listy dojdzie obliczenie
niezmienników Seiberga--Wittena i porównanie komór.
Żaden z tych wyników nie wynika z samej macierzy $Q_7$.
\end{uwaga}


% ========================================================================
% Iteracja 14. Umiescic bezposrednio po rozdziale o konstrukcji Parka.
% Konwencja: u_0=S_0 (sfera -9), u_1=S_5 z pracy Parka,
% ..., u_5=S_1 z pracy Parka. Tylko ponizszy dokument zawiera pelne
% doprowadzenie do nierownosci K_{X_7}.[omega_7]>0.
% ========================================================================
\begin{uwaga}[Geometria konfiguracji Parka]
Rachunek $Q_{C_7}$ oraz równości
$e(X_7)=10$, $\sigma(X_7)=-6$ nie konstruują same
sześciu symplektycznych sfer w rozdmuchanej powierzchni
eliptycznej. Ich konkretne klasy i istnienie zgodnej
formy symplektycznej są danymi geometrycznymi konstrukcji
Parka. Podobnie prosta spójność dopełnienia musi zostać
wykazana przez odpowiedni argument o meridianie,
a nie wywnioskowana wyłącznie z wyznacznika formy.
\end{uwaga}

\chapter{Niezmienniki Seiberga--Wittena i egzotyczność rozmaitości Parka}
\label{ch:tg-park-sw}

W poprzednim rozdziale opisaliśmy konstrukcję gładkiej,
zamkniętej, orientowanej i po prostu spójnej
czterorozmaitości $X_7$, korzystając z geometrycznych
wyników Parka i Cassona--Harera.
Wiemy, że ma ona formę przecięcia izometryczną z formą
$\langle1\rangle\oplus7\langle-1\rangle$, a zatem jest
homeomorficzna z $\mathbb{CP}^2\#7\overline{\mathbb{CP}}^{\,2}$.
Nie wynika stąd jednak, że obie rozmaitości są dyfeomorficzne.
Pokażemy, dlaczego nie są, korzystając z teorii Seiberga--Wittena.

Wymiar dodatniej części formy przecięcia wynosi tu $b_2^+=1$.
Przy tym założeniu liczba otrzymana przez zliczenie rozwiązań
równań Seiberga--Wittena zależy na ogół od \emph{komory}.
Z tego powodu musimy szczególnie starannie oddzielić niezmiennik
z małą perturbacją, który dla rozważanej klasy charakterystycznej
jest niezależny od metryki, od niezmiennika w \emph{komorze
symplektycznej}. Wybór właściwej komory jest zasadniczym krokiem
rozróżnienia struktur gładkich.

\section{Klasy sfer i ich przecięcia: uzupełnienie konstrukcji}
\label{sec:tg-park14-spheres}

Przypomnijmy bazę $h,e_1,\ldots,e_{13}$ grupy
$H_2(R;\mathbb Z)$, gdzie $R=E(1)\#4\overline{\mathbb{CP}}^{\,2}$,
$h^2=1$, $e_i^2=-1$ i wszystkie pozostałe przecięcia między
różnymi elementami bazy są zerowe. Oznaczmy
$f=3h-e_1-\cdots-e_9$. Wybierzmy pięć klas
\[
\begin{aligned}
 v_1&=e_4-e_7, & v_2&=e_1-e_4, \\
 v_3&=h-e_1-e_2-e_3, & v_4&=e_2-e_5, \\
 v_5&=e_5-e_9,
\end{aligned}
\]
oraz $s=4f+e_9-2(e_{10}+e_{11}+e_{12}+e_{13})$.
Park wykorzystuje realizację klas $v_i$ przez pięć sfer
jednego włókna osobliwego typu $\widetilde E_6$, a klasy $s$
przez szóstą sferę uzyskaną po rozwiązaniu przecięć.
Geometryczne istnienie takiej realizacji przyjmujemy z jego pracy.

\begin{stwierdzenie}[Ustalenie numeracji]
\label{prop:tg-park14-classes}
Połóżmy $u_0:=s$ oraz $u_j:=v_{6-j}$ dla $1\le j\le5$.
Wówczas $u_0^2=-9$, $u_j^2=-2$ dla $j\ge1$,
$u_j\cdot u_{j+1}=1$ dla $0\le j\le4$, a pozostałe
przecięcia różnych sfer są zerowe.
\end{stwierdzenie}

\begin{proof}
Jest to dokładnie wynik lematu
\ref{lem:tg-park-S-classes} z poprzedniego rozdziału,
po odwróceniu numeracji pięciu sfer o samoprzecięciu $-2$.
Podane tam obliczenie wykorzystuje jedynie diagonalną
formę przecięcia na $H_2(R;\mathbb Z)$.
\end{proof}

\begin{uwaga}[Zakres dowodu]
Lemat jest \emph{rachunkiem homologicznym}.
Nie dowodzi, że dowolne sześć klas o takiej macierzy
ma reprezentantów będących zanurzonymi sferami
symplektycznymi o zadanym układzie przecięć.
To właśnie niezależna geometryczna treść konstrukcji Parka.
Jest ona potrzebna do zastosowania twierdzenia Symington.
\end{uwaga}

\section{Klasa kanoniczna i odwrotność macierzy plumbingowej}
\label{sec:tg-park14-canonical}

Niech $(R,\omega)$ będzie symplektyczną powierzchnią
racjonalną z konfiguracją Parka, a $K_R=-c_1(TR,J)$
jej klasą kanoniczną dla zgodnej struktury prawie
zespolonej $J$. Po czterech rozdmuchaniach mamy
\[
K_R=-3h+e_1+\cdots+e_{13}
    =-f+e_{10}+e_{11}+e_{12}+e_{13}.
\]
W szczególności $K_R\cdot u_i=0$ dla $1\le i\le5$.
Aby policzyć $K_R\cdot u_0$, można wykorzystać
formułę adjunkcji dla symplektycznej sfery $u_0$:
$-2=u_0^2+K_R\cdot u_0=-9+K_R\cdot u_0$.
Dostajemy $K_R\cdot u_0=7$. Ten sam wynik otrzymujemy
bezpośrednio z wyrażenia klasy $s$ w bazie $h,e_i$.

\begin{lemat}[Macierz odwrotna konfiguracji $C_7$]
\label{lem:tg-park14-inverse}
Niech $Q_7$ będzie macierzą przecięcia konfiguracji $C_7$.
Jej odwrotność jest określona nad $\mathbb Q$ i ma pierwszy
wiersz postaci
\[
(Q_7^{-1})_{0*}=-\frac1{49}(6,5,4,3,2,1).
\]
\end{lemat}

\begin{proof}
Przypomnijmy, że $\det Q_7=49$, więc macierz jest
odwracalna nad $\mathbb Q$. Aby znaleźć pierwszy wiersz,
szukamy wektora $z=(z_0,\ldots,z_5)^T$ spełniającego
$Q_7z=(1,0,\ldots,0)^T$; symetria $Q_7$ kończy dowód.
Pięć ostatnich równań to kolejno
$z_{j-1}-2z_j+z_{j+1}=0$ dla $1\le j\le4$
oraz $z_4-2z_5=0$. Wynika z nich
$z_j=(6-j)z_5$ dla $0\le j\le5$.
Pierwsze równanie $-9z_0+z_1=1$ daje
$-49z_5=1$. Otrzymujemy
$z_j=-(6-j)/49$, jak w tezie.
\end{proof}

\begin{stwierdzenie}[Rachunek parowania po wycięciu konfiguracji]
\label{prop:tg-park14-splitting}
Niech $V=R\setminus\operatorname{int}C_7$,
a $X_7=V\cup B_7$ będzie symplektyczną wymierną
kontrakcją, przy której formy symplektyczne na $V$
są identyfikowane przez symplektomorfizm.
Niech $\alpha_i:=[\omega]\cdot u_i$ dla $0\le i\le5$.
Wtedy
\[
K_{X_7}\cdot[\omega_7]
 =K_R\cdot[\omega]+
   \frac17(6\alpha_0+5\alpha_1+4\alpha_2
                    +3\alpha_3+2\alpha_4+\alpha_5).
\]
\end{stwierdzenie}

\begin{proof}
Ponieważ $H_1(\partial C_7;\mathbb Q)
=H_2(\partial C_7;\mathbb Q)=0$,
ciąg Mayera--Vietorisa pozwala rozłożyć, nad $\mathbb Q$,
klasy kohomologii z $R$ na części pochodzące z $V$ i $C_7$.
Oznaczmy przez $\gamma_i$ bazę dualną do $u_i$ w
$H^2(C_7;\mathbb Q)$. Ograniczenie klasy kanonicznej ma
wektor współrzędnych $(7,0,\ldots,0)$, a $[\omega]$
ma współrzędne $(\alpha_0,\ldots,\alpha_5)$.
Formę przecięcia w bazie dualnej reprezentuje $Q_7^{-1}$.
Lemat \ref{lem:tg-park14-inverse} daje
\[
K_R|_{C_7}\cdot[\omega]|_{C_7}
  =7(1,0,\ldots,0)Q_7^{-1}\alpha
  =-\frac17\sum_{j=0}^{5}(6-j)\alpha_j.
\]
Przy wymiernej kontrakcji $H^2(B_7;\mathbb Q)=0$.
Klasy kanoniczne i symplektyczne na $V$ są utożsamione,
a wymierna część formy przecięcia pochodząca z $B_7$
nie wnosi wkładu. Stąd, w ujęciu twierdzenia
Symington i lematu porównującego parowania,
$K_{X_7}\cdot[\omega_7]
=K_R\cdot[\omega]-K_R|_{C_7}\cdot[\omega]|_{C_7}$.
Podstawiając obliczoną wartość, otrzymujemy tezę.
\end{proof}

\section{Ścisła dodatniość parowania \texorpdfstring{$K_{X_7}\cdot[\omega_7]$}{K omega}}
\label{sec:tg-park14-positivity}

Kluczowa pozostaje możliwość wyboru struktury symplektycznej,
zgodnej z opisaną konfiguracją, dla której współrzędne klasy
symplektycznej spełniają nierówności potrzebne w rachunku.
Wykorzystamy tu następującą postać twierdzeń o stożku
symplektycznym powierzchni racjonalnych, używaną przez Parka.

\begin{twierdzenie}[Geometryczne założenie o klasie symplektycznej; wynik zewnętrzny]
\label{thm:tg-park14-cone}
W konstrukcji Parka istnieje zgodna struktura symplektyczna
na $R$, dla której zanurzone sfery konfiguracji $C_7$
są symplektyczne, a klasa $[\omega]\in H^2(R;\mathbb R)$
ma reprezentację
\[
[\omega]=a h-\sum_{j=1}^{13}b_j e_j,
\qquad
 a\ge b_1\ge\cdots\ge b_{13}>0,
\qquad 3a>\sum_{j=1}^{13}b_j.
\]
Istnienie jednocześnie takiej formy i takiej symplektycznej
realizacji konfiguracji jest założeniem geometrycznym
importowanym z konstrukcji Parka i teorii stożka
symplektycznego Li--Liu; nie jest konsekwencją
samych nierówności między liczbami $a,b_j$.
\end{twierdzenie}

\begin{uwaga}[Zakres założenia geometrycznego]
Powyższa teza dotyczy konkretnej realizacji sfer w $R$.
Nie wolno wybrać dowolnych dodatnich współczynników
$a,b_j$ i na podstawie samych nierówności twierdzić,
że reprezentują one formę symplektyczną, względem której
wszystkie sfery konfiguracji są symplektyczne.
Istnienie formy, jej zgodność z ustalonymi klasami sfer
i znaki ich pól symplektycznych muszą zostać zweryfikowane
w rzeczywistej konstrukcji geometrycznej.
\end{uwaga}

Dla takiego wyboru wyprowadzimy znak parowania
$K_{X_7}\cdot[\omega_7]$ przez całkowicie jawny rachunek.

\begin{twierdzenie}[Nierówność kluczowa w dowodzie Parka]
\label{thm:tg-park14-positive}
Niech $X_7$ pochodzi z symplektycznej wymiernej kontrakcji
w poprzednim rozdziale i niech $\omega_7$ odpowiada klasie
$[\omega]$ z twierdzenia \ref{thm:tg-park14-cone}.
Wtedy $K_{X_7}\cdot[\omega_7]>0$.
\end{twierdzenie}

\begin{proof}
Wprowadźmy skrót $B:=b_{10}+b_{11}+b_{12}+b_{13}$.
Z klas $u_0=s$, $u_1=v_5$, $u_2=v_4$, $u_3=v_3$,
$u_4=v_2$ i $u_5=v_1$ otrzymujemy bezpośrednio
\[
\begin{aligned}
\alpha_0&=12a-4\sum_{j=1}^{9}b_j+b_9-2B,
 &\alpha_1&=b_5-b_9,\\
\alpha_2&=b_2-b_5,
 &\alpha_3&=a-b_1-b_2-b_3,\\
\alpha_4&=b_1-b_4,
 &\alpha_5&=b_4-b_7.
\end{aligned}
\]
Na przykład w wyrażeniu dla $\alpha_0$ współczynnik
przy $b_9$ to $-3$, a nie $-4$, ponieważ klasa
$s=4f+e_9-2\sum_{j=10}^{13}e_j$ zawiera dodatkowy
składnik $+e_9$. Ponadto
$K_R\cdot[\omega]=-3a+\sum_{j=1}^{13}b_j$.

Korzystając ze stwierdzenia \ref{prop:tg-park14-splitting}
i grupując wyrazy, otrzymujemy
\[
\begin{aligned}
7K_{X_7}\cdot[\omega_7]
={}&54a-18b_1-16b_2-20b_3-18b_4-16b_5\\
 &-17b_6-18b_7-17b_8-16b_9-5B.
\end{aligned}
\]
Nie jest jeszcze oczywiste, że prawa strona jest dodatnia.
Rozbijmy ją zatem na dwie części:
\[
\begin{aligned}
7K_{X_7}\cdot[\omega_7]
={}&18\left(3a-\sum_{j=1}^{13}b_j\right)\\
 &+2(b_2-b_3)+2b_5+b_6+b_8+2b_9+13B.
\end{aligned}
\]
Pierwszy nawias jest ściśle dodatni z twierdzenia
o wyborze klasy symplektycznej. Pozostałe wyrazy
są nieujemne, ponieważ $b_2\ge b_3$ oraz $b_j>0$.
Wynika stąd $K_{X_7}\cdot[\omega_7]>0$.
\end{proof}

\begin{przyklad}[Kontrola algebraiczna wzoru]
Przyjmijmy wyłącznie do sprawdzenia rachunku
$a=40$ oraz $(b_1,\ldots,b_{13})=(13,12,11,10,9,8,7,6,5,1,1,1,1)$.
Wtedy $3a-\sum b_j=35>0$, a wektor przecięć
$(\alpha_0,\ldots,\alpha_5)=(153,4,3,4,3,3)$.
Zatem $K_R\cdot[\omega]=-35$ oraz
\[
K_{X_7}\cdot[\omega_7]
 =-35+\frac{6\cdot153+5\cdot4+4\cdot3
                    +3\cdot4+2\cdot3+3}{7}
 =\frac{726}{7}>0.
\]
Liczby te sprawdzają tożsamość algebraiczną;
nie twierdzimy, że dowolna zadana lista współczynników
opisuje formę symplektyczną zachowującą cały geometryczny
układ sfer Parka.
\end{przyklad}

\begin{center}
\begin{tikzpicture}[>=Stealth,scale=.95,line cap=round,line join=round]
  \node[draw=tgBlue,fill=tgBlue!8,rounded corners=4pt,
     minimum width=3.1cm,minimum height=1.1cm,align=center] (r)
     at (-3.3,0) {\small $R=E(1)\#4\overline{\mathbb{CP}}^{\,2}$\\[-1pt]
             \scriptsize $K_R\cdot[\omega]<0$};
  \node[draw=tgCoral,fill=tgCoral!10,rounded corners=4pt,
     minimum width=3.0cm,minimum height=1.1cm,align=center] (c)
     at (0,0) {\small konfiguracja $C_7$\\[-1pt]
             \scriptsize ujemne parowanie};
  \node[draw=tgTeal,fill=tgTeal!10,rounded corners=4pt,
     minimum width=3.1cm,minimum height=1.1cm,align=center] (x)
     at (3.3,0) {\small $X_7$ po kontrakcji\\[-1pt]
             \scriptsize $K_{X_7}\cdot[\omega_7]>0$};
  \draw[-{Stealth[length=2.6mm]},tgPurple,line width=1.7pt] (r)--(c);
  \draw[-{Stealth[length=2.6mm]},tgPurple,line width=1.7pt] (c)--(x);
  \node[text=tgNavy,font=\footnotesize] at (0,-1.05)
      {Wymierna kontrakcja zmienia znak istotnego parowania.};
\end{tikzpicture}
\end{center}

\section{Komory przy \texorpdfstring{$b_2^+=1$}{b2+=1} i małe perturbacje}
\label{sec:tg-park14-chamber}

Przypomnijmy, że dla struktury $\mathrm{Spin}^c$ o klasie
$c:=c_1(L)$ i wymiarze oczekiwanym $d=0$ niezmiennik
Seiberga--Wittena jest liczbą otrzymaną ze zorientowanego
zliczenia rozwiązań dla danych generycznych. Jeśli $b_2^+=1$,
przestrzeń metryk i perturbacji ma ścianę z rozwiązaniami
redukowalnymi. Rozdziela ona dwa rodzaje komór.
Do porównania z metryką dodatniej krzywizny skalarnej
wybieramy perturbacje dostatecznie małe, osobno wzdłuż
każdej zwartej ścieżki metryk; uzyskaną wartość oznaczamy
$SW^{\circ}_X(\mathfrak s)$.

\begin{lemat}[Brak ściany dla małych perturbacji]
\label{lem:tg-park14-no-wall}
Niech $M$ będzie zamkniętą, spójną, orientowaną
czterorozmaitością z $b_2^+(M)=1$ i ustaloną składową
stożka dodatniego. Jeżeli $c\in H^2(M;\mathbb R)$
spełnia $c^2>0$, to jego rzut harmoniczny samodualny
$c_g^+$ jest niezerowy dla każdej metryki $g$.
Co więcej, znak $c\cdot H_g$ dla generatora $H_g$
ustalonej dodatniej półprostej $H_g^2>0$ jest niezależny
od metryki. W konsekwencji, po wyborze wystarczająco małych
perturbacji wzdłuż każdej rozpatrywanej ścieżki metryk,
nie przecina się ściany redukowalnej.
\end{lemat}

\begin{proof}
Dekompozycja Hodge'a dla metryki $g$ daje ortogonalny
rozkład $c=c_g^++c_g^-$ na dodatni i ujemny podprzestrzenie
formy przecięcia. Gdyby $c_g^+=0$, mielibyśmy
$c^2=(c_g^-)^2\le0$, sprzecznie z założeniem.
W wymiarze Lorentzowskim dwa wektory o dodatnim kwadracie,
należące do ustalonych składowych stożka, nie są ortogonalne.
Znak iloczynu $c\cdot H_g$ jest więc stały, gdy $H_g$
przebiega ustaloną składową stożka dodatniego.
Dla ustalonej zwartej ścieżki metryk wartość bezwzględna
tego parowania jest jednostajnie oddalona od zera.
Wobec tego perturbacja dostatecznie mała na całej ścieżce
nie może doprowadzić do równości definiującej ścianę.
\end{proof}

\begin{wniosek}[Niezmiennik z małą perturbacją]
\label{cor:tg-park14-small-invariant}
Jeśli dodatkowo $M$ jest po prostu spójna, $d(\mathfrak s)=0$
i $c_1(\mathfrak s)^2>0$, to przy ustalonej orientacji
homologicznej niezmiennik $SW_X^{\circ}(\mathfrak s)$
można określić niezależnie od wyboru metryki, używając
małej generycznej perturbacji. Jeżeli na $M$ istnieje
metryka dodatniej krzywizny skalarnej, to
$SW_M^{\circ}(\mathfrak s)=0$.
\end{wniosek}

\begin{proof}
Niezmienniczość liczby rozwiązań w rodzinie generycznych
danych bez redukowalnych wynika z parametrycznej
transwersalności, zwartości przestrzeni moduli i orientacji;
są to twierdzenia analityczne wyprowadzone zewnętrznymi
narzędziami w poprzednich rozdziałach. Lemat
\ref{lem:tg-park14-no-wall} wyklucza ścianę w dowolnej
zwartej ścieżce metryk przy dostatecznie małej perturbacji.

Jeżeli krzywizna skalarna jest wszędzie dodatnia, formuła
Weitzenböcka przy zerowej perturbacji wyklucza rozwiązania
z niezerowym spinorem $\psi$: całkując równanie Diraca
otrzymujemy sumę nieujemnych wyrazów, wśród których
występują dodatnie $s_g|\psi|^2$ i $|\psi|^4$.
Rozwiązanie redukowalne wymagałoby, aby samodualna
część krzywizny była zerowa. Wtedy harmoniczna klasa
$c_g^+$ musiałaby być zerowa, wbrew poprzedniemu lematowi.
Równania nie mają więc rozwiązań dla perturbacji zerowej,
a przez zwartość moduli pozostają bez rozwiązań dla
wszystkich dostatecznie małych perturbacji.
Stąd zliczenie rozwiązań wynosi zero.
\end{proof}

\begin{uwaga}[Dlaczego nie wystarczy po prostu porównać komór?]
Dla $b_2^+=1$ symplektyczna komora Taubesa nie musi
pokrywać się z komorą, w której pewna metryka dodatniej
krzywizny skalarnej daje pustą przestrzeń moduli.
Bez lematu \ref{lem:tg-park14-no-wall} porównanie
niezerowego zliczenia w jednej komorze z zerowym zliczeniem
w drugiej nie prowadziłoby do sprzeczności.
W zastosowaniu do Parka decydujące jest to, że
$c^2=K_{X_7}^2=2>0$.
\end{uwaga}

\section{Twierdzenie Taubesa--Li--Liu i brak dyfeomorfizmu}
\label{sec:tg-park14-exotic}

\begin{twierdzenie}[Taubes--Li--Liu; wynik zewnętrzny]
\label{thm:tg-park14-taubes}
Niech $(M,\omega)$ będzie zamkniętą, spójną,
symplektyczną czterorozmaitością z $b_2^+(M)=1$,
a $K_M=-c_1(TM,J)$ jej klasą kanoniczną.
Dla struktury $\mathrm{Spin}^c$ o klasie wyznacznikowej
$c_1(L)=-K_M$, przy orientacji homologicznej
zgodnej z $[\omega]$, niezmiennik w odpowiedniej
komorze symplektycznej spełnia
\[
SW_M^{\mathrm{symp}}(-K_M)=\pm1.
\]
W konwencji, w której komora oznaczona znakiem
$-$ jest komorą zawierającą silną perturbację
symplektyczną, zapisujemy $SW_M^-(-K_M)=\pm1$.
Twierdzenie to przyjmujemy z teorii cechowania;
jego dowód korzysta z odpowiednio przeskalowanych
równań Seiberga--Wittena i analizy granicznej Taubesa.
\end{twierdzenie}

\begin{twierdzenie}[Egzotyczność rozmaitości Parka]
\label{thm:tg-park14-exotic}
Rozmaitość $X_7$ z poprzedniego rozdziału jest
homeomorficzna, lecz nie jest dyfeomorficzna z
$\mathbb{CP}^{2}\#7\overline{\mathbb{CP}}^{\,2}$.
Nie istnieje też na $X_7$ żadna metryka riemannowska
o wszędzie dodatniej krzywiźnie skalarnej.
\end{twierdzenie}

\begin{proof}
Homeomorfizm został już wykazany w twierdzeniu
\ref{thm:tg-park-homeo}. Skupmy się na strukturze gładkiej.
Symplektyczność $X_7$ wynika z geometrycznej konstrukcji
Parka i twierdzenia Symington. Dla klasy kanonicznej $K_7$
wykazaliśmy $K_7\cdot[\omega_7]>0$ w twierdzeniu
\ref{thm:tg-park14-positive}. Zatem
$(-K_7)\cdot[\omega_7]<0$.
Ponadto $K_7^2=2e(X_7)+3\sigma(X_7)=2$,
a wymiar oczekiwany moduli dla struktury o klasie
$c_1(L)=-K_7$ wynosi
\[
d(\mathfrak s)=\frac{(-K_7)^2-2e(X_7)-3\sigma(X_7)}4=0.
\]
Ponieważ $(-K_7)^2>0$, lemat
\ref{lem:tg-park14-no-wall} pokazuje, że znak
parowania $(-K_7)\cdot H_g$ pozostaje ujemny dla
każdej metryki i ustalonej składowej stożka dodatniego.
Wobec tego komora małych perturbacji dla tej klasy
jest właśnie komorą symplektyczną, w której twierdzenie
Taubesa--Li--Liu daje
\[
SW_{X_7}^{\circ}(-K_7)
 =SW_{X_7}^{\mathrm{symp}}(-K_7)=\pm1.
\]

Gdyby $X_7$ miała metrykę dodatniej krzywizny skalarnej,
wniosek \ref{cor:tg-park14-small-invariant} dałby
$SW_{X_7}^{\circ}(-K_7)=0$, sprzecznie z powyższą równością.
Zatem taka metryka nie istnieje.

Z kolei $\mathbb{CP}^{2}$ i $\overline{\mathbb{CP}}^{\,2}$
dopuszczają metryki dodatniej krzywizny skalarnej,
a twierdzenie o zachowywaniu dodatniej krzywizny
skalarnej przy sumie spójnej w wymiarze co najmniej trzy
zapewnia taką metrykę na
$\mathbb{CP}^{2}\#7\overline{\mathbb{CP}}^{\,2}$.
Wspomniane twierdzenie o sumie spójnej przyjmujemy z
geometrii riemannowskiej. Istnienie metryki dodatniej
krzywizny skalarnej jest niezmiennikiem typu dyfeomorfizmu.
Stąd standardowa rozmaitość nie może być dyfeomorficzna
z $X_7$.
\end{proof}

\begin{center}
\begin{tikzpicture}[>=Stealth,line cap=round,line join=round]
  \node[draw=tgTeal,fill=tgTeal!10,rounded corners=5pt,
       text width=4.15cm,align=center,inner sep=6pt] (a)
       at (-3.45,0) {\small $X_7$ — konstrukcja Parka\\
          \footnotesize $K_7\cdot[\omega_7]>0$, $(-K_7)^2=2$\\
          \footnotesize $SW^{\circ}(-K_7)=\pm1$};
  \node[draw=tgBlue,fill=tgBlue!10,rounded corners=5pt,
       text width=4.15cm,align=center,inner sep=6pt] (b)
       at (3.45,0) {\small $\mathbb{CP}^{2}\#7\overline{\mathbb{CP}}^{\,2}$\\
          \footnotesize dodatnia krzywizna skalarna\\
          \footnotesize $SW^{\circ}(c)=0$ dla $c^2>0$};
  \node[text=tgCoral,font=\scriptsize,align=center] at (0,.35)
        {nie są\\dyfeomorficzne};
  \node[text=tgPurple,font=\scriptsize,align=center] at (0,-.42)
        {są\\homeomorficzne};
\end{tikzpicture}
\end{center}

\begin{uwaga}[Znaczenie wyniku]
Przykład Parka pokazuje, że forma przecięcia i grupa
fundamentalna, nawet w połączeniu z charakterystyką Eulera
i sygnaturą, nie klasyfikują gładkich czterorozmaitości.
Wymierna kontrakcja zmienia gładką strukturę w sposób,
który nie jest wykrywany przez klasyfikację Freedmana.
Rozstrzygające są dwa dodatkowe fakty: dodatniość
$K_7\cdot[\omega_7]$ oraz niezerowość odpowiedniego
niezmiennika Seiberga--Wittena.
\end{uwaga}

\section{Dokładny status użytych twierdzeń i dalszy krok}
\label{sec:tg-park14-status}

Dowód egzotyczności jest \emph{kompletny warunkowo względem
jawnie wymienionych wyników zewnętrznych}:
(1) istnienia symplektycznej konfiguracji Parka i prostej
spójności jej dopełnienia, (2) twierdzenia Symington,
(3) wyboru klasy symplektycznej z teorii Li--Liu,
(4) twierdzenia Taubesa--Li--Liu i zwartości moduli,
(5) twierdzenia o dodatniej krzywiźnie skalarnej
na sumach spójnych.
Elementarne rachunki form przecięcia, rozkładu klas,
wyznacznika, nierówności i porównania komór zostały
przedstawione szczegółowo. Osobnego wykładu wymagałby
pełny dowód analityczny twierdzenia Taubesa.

W następnej części skryptu będziemy mogli wykorzystać
podobną strategię do operacji chirurgii wzdłuż węzła
(\emph{knot surgery}), gdzie wielomian Alexandera
wchodzi bezpośrednio do formuł na niezmienniki
Seiberga--Wittena. W ten sposób teoria grafów,
teoria węzłów i topologia czterowymiarowa tworzą
jeden ciąg argumentów, a nie trzy odrębne tematy.


% ================================================================
% ITERACJA 15 — rozdzial wstawiany po egzotycznosci Parka
% Konwencje: długosc preferowana lambda_K; zmienna u = exp(2[T]);
% twierdzenia Fintushela--Sterna cytowane bez pelnego dowodu.
% ================================================================
\chapter{Chirurgia wzdłuż węzła i niezmienniki czterorozmaitości}
\label{ch:tg-knot-surgery}

W rozdziałach o macierzach Seiferta pokazaliśmy, jak wielomian
Alexandera koduje własności trójwymiarowego uzupełnienia węzła.
Teraz wyjaśnimy, jak ten sam wielomian pozwala rozróżniać
\emph{struktury gładkie} na czterorozmaitościach. Podstawowym
narzędziem jest konstrukcja Fintushela--Sterna, zwana
\emph{chirurgią wzdłuż węzła} (\emph{knot surgery}).

Istotne jest odróżnienie tej operacji od chirurgii Dehna.
W chirurgii Dehna zastępujemy torus pełny w rozmaitości
trójwymiarowej. Tutaj usuwamy otoczenie torusa $T^2$ w
czterorozmaitości, a w jego miejsce wstawiamy iloczyn okręgu
z zewnętrzem węzła. Operacja zachowuje wiele klasycznych
niezmienników rozmaitości, może jednak zmienić jej typ
dyfeomorfizmu.

\section{Dane geometryczne i meridian torusa}
\label{sec:tg-ks-data}

\begin{definicja}[Torus o trywialnej wiązce normalnej]
Niech $X$ będzie gładką, zamkniętą, spójną i zorientowaną
czterorozmaitością, a $T\subset X$ gładko zanurzonym,
zorientowanym torusem. Zakładamy, że jego wiązka normalna
jest trywialna jako zorientowana wiązka rzeczywista
rzędu dwa. Wybieramy trywializację otoczenia
$\nu T\cong T^2\times D^2$. Na torusie brzegowym
$\partial\nu T\cong T^3$ oznaczamy przez $a,b$ dwie
klasy stanowiące bazę $H_1(T;\mathbb Z)$, a przez
$m_T$ klasę okręgu $\{*\}\times\partial D^2$.
\end{definicja}

Dla zorientowanego torusa w orientowanej czterorozmaitości
warunek trywialności wiązki normalnej jest równoważny
$T\cdot T=0$. Rzeczywiście zorientowane wiązki
rzeczywiste rzędu dwa nad $T^2$ są klasyfikowane
przez klasę Eulera w $H^2(T^2;\mathbb Z)\cong\mathbb Z$,
a jej wartość na klasie fundamentalnej jest
liczbą samoprzecięcia torusa. Ta obserwacja
wyjaśnia geometryczną treść założenia $T^2=0$;
nie mówi jednak nic o grupie fundamentalnej dopełnienia.

\begin{definicja}[Zewnętrze węzła i długość preferowana]
Niech $K\subset S^3$ będzie zorientowanym węzłem.
Wybieramy domknięte otoczenie tubularne $N(K)$ i kładziemy
$E_K:=S^3\setminus\operatorname{int}N(K)$.
Na $\partial E_K$ oznaczamy przez $\mu_K$ klasę
meridianu, a przez $\lambda_K$ klasę długości
preferowanej, to znaczy jedyną klasę prymitywną
mającą algebraiczną liczbę przecięcia z $\mu_K$
równą $\pm1$ i przechodzącą na zero w
$H_1(E_K;\mathbb Z)$. Znak $\lambda_K$ ustalamy
zgodnie z orientacją $K$.
\end{definicja}

\begin{lemat}[Homologia zewnętrza]
\label{lem:tg-ks-exterior}
Dla każdego węzła $K\subset S^3$ zachodzi
$H_0(E_K;\mathbb Z)\cong\mathbb Z$,
$H_1(E_K;\mathbb Z)\cong\mathbb Z$ i
$H_i(E_K;\mathbb Z)=0$ dla $i\ge2$.
Meridian $\mu_K$ reprezentuje generator $H_1(E_K;\mathbb Z)$,
natomiast $[\lambda_K]=0$ w tej grupie.
\end{lemat}

\begin{proof}
Zewnętrze jest homotopijnie równoważne $S^3\setminus K$.
Dualność Alexandera daje
$\widetilde H_i(S^3\setminus K;\mathbb Z)
\cong\widetilde H^{2-i}(S^1;\mathbb Z)$.
W szczególności jedyna niezerowa zredukowana grupa
homologii pojawia się w wymiarze jeden i ma rząd jeden.
Izomorfizm z $\mathbb Z$ jest zadany liczbą splecenia
z $K$, dlatego meridian jest generatorem. Z definicji
długość preferowana ma zerowy obraz w $H_1(E_K)$.
\end{proof}

\begin{wniosek}[Homologia części dołączanej]
\label{cor:tg-ks-piece}
Niech $P_K:=S^1\times E_K$, a $s$ niech oznacza
klasę pierwszego czynnika $S^1$. Wówczas
$H_1(P_K;\mathbb Z)\cong\mathbb Z[s]\oplus\mathbb Z[\mu_K]$,
$H_2(P_K;\mathbb Z)\cong\mathbb Z[s\times\mu_K]$
oraz $H_i(P_K;\mathbb Z)=0$ dla $i\ge3$.
Ponadto $\chi(P_K)=0$.
\end{wniosek}

\begin{proof}
Wszystkie grupy homologii z lematu
\ref{lem:tg-ks-exterior} są wolne, więc wzór Künnetha
nie ma składników $\operatorname{Tor}$. Otrzymujemy
$H_1(S^1\times E_K)\cong H_1(S^1)\oplus H_1(E_K)$
oraz $H_2(S^1\times E_K)\cong H_1(S^1)\otimes H_1(E_K)$.
Pozostałe grupy dodatniego wymiaru są zerowe.
Wreszcie $\chi(S^1\times E_K)=\chi(S^1)\chi(E_K)=0$.
\end{proof}

\section{Definicja chirurgii i sposób sklejenia}
\label{sec:tg-ks-glue}

\begin{definicja}[Chirurgia wzdłuż węzła Fintushela--Sterna]
\label{def:tg-ks-operation}
Niech $X,T,K$ spełniają warunki poprzedniej sekcji.
Niech $A:=X\setminus\operatorname{int}\nu T$.
Wybieramy dyfeomorfizm $\varphi:\partial P_K\to\partial A$
odwracający orientację brzegową i spełniający na
pierwszych grupach homologii równości
$\varphi_*([s])=a$, $\varphi_*([\mu_K])=b$
oraz $\varphi_*([\lambda_K])=m_T$.
Rozmaitość
\[
X_K:=A\cup_\varphi(S^1\times E_K)
\]
nazywamy \emph{wynikiem chirurgii wzdłuż węzła $K$ na torusie $T$}.
Identyfikację przyjmujemy jako część danych standardowej
operacji: jest ona zgodna z wybranym ramowaniem torusa
i z orientacjami na brzegach. W zastosowaniach
formuły Seiberga--Wittena żądamy dodatkowo, aby $T$
leżał w otoczeniu typu \emph{cusp}. Przez $[T_K]$
oznaczamy klasę torusa pochodzącego od
$S^1\times\mu_K$ po niewielkim przesunięciu do wnętrza $P_K$.
\end{definicja}

\begin{uwaga}[Orientacje i zakres definicji]
Równanie $\varphi_*([\lambda_K])=m_T$ to istota
operacji: meridian usuwanego torusa jest zastępowany
przez długość preferowaną węzła. Sam zapis na $H_1$
nie jest stwierdzeniem, że dowolne sklejenie trójtorusów
o takich własnościach daje zawsze tę samą rozmaitość
gładką. W twierdzeniach poniżej używamy standardowego
sklejenia Fintushela--Sterna przy torusie o ustalonym
ramowaniu. W szczególności nie utożsamiamy bez dowodu
wyników wszystkich możliwych sklejeń.
\end{uwaga}

\begin{center}
\begin{tikzpicture}[>=Stealth,line cap=round,line join=round]
  \node[draw=tgBlue,fill=tgBlue!10,rounded corners=4pt,
        text width=4.0cm,align=center,inner sep=7pt] (left) at (-3.25,0)
        {\small $A=X\setminus\operatorname{int}\nu T$\\
         \footnotesize brzeg $T^3$; klasy $a,b,m_T$};
  \node[draw=tgTeal,fill=tgTeal!11,rounded corners=4pt,
        text width=4.0cm,align=center,inner sep=7pt] (right) at (3.25,0)
        {\small $S^1\times E_K$\\
         \footnotesize brzeg $T^3$; klasy $s,\mu_K,\lambda_K$};
  \draw[-{Stealth[length=2.5mm]},tgPurple,line width=1.8pt] (left)--(right);
  \node[fill=white,text=tgNavy,font=\scriptsize,align=center,inner sep=2pt] at (0,.35)
        {$m_T\longleftrightarrow\lambda_K$};
  \node[text=tgCoral,font=\footnotesize] at (0,-1.12)
        {$X_K=A\cup_\varphi(S^1\times E_K)$};
\end{tikzpicture}
\end{center}

\begin{przyklad}[Węzeł trywialny]
Dla węzła trywialnego $U$ mamy $E_U\cong S^1\times D^2$,
przy czym meridian węzła odpowiada kierunkowi $S^1$,
a długość preferowana ogranicza dysk w $E_U$.
Stąd $P_U\cong T^2\times D^2$. Przy standardowym sklejeniu
$X_U\cong X$. Przykład ten ustala, która krzywa
pełni rolę meridianu otoczenia torusa po wymianie.
\end{przyklad}

\section{Homologia, charakterystyka Eulera i sygnatura}
\label{sec:tg-ks-homology}

\begin{stwierdzenie}[Homologia dwóch modeli otoczenia]
\label{prop:tg-ks-boundary-maps}
Przy identyfikacji brzegów z definicji
\ref{def:tg-ks-operation} odwzorowania indukowane
przez włączenia $T^3\hookrightarrow T^2\times D^2$
i $T^3\hookrightarrow S^1\times E_K$ mają izomorficzne
obrazy i jądra we wszystkich stopniach homologii
całkowitej, po utożsamieniu generatorów $a,b,m_T$
z odpowiednio $s,\mu_K,\lambda_K$.
\end{stwierdzenie}

\begin{proof}
W stopniu jeden w pierwszym modelu zachodzi
$a\mapsto a$, $b\mapsto b$, $m_T\mapsto0$.
W drugim modelu mamy $s\mapsto s$,
$\mu_K\mapsto\mu_K$, $\lambda_K\mapsto0$.
Zatem jądro jest generowane w obu przypadkach
przez trzeci element, a dwa pierwsze przechodzą
na bazę grupy $\mathbb Z^2$.

W stopniu dwa homologia torusa brzegowego jest
wolna rzędu trzy, z bazą reprezentowaną przez
$a\times b$, $a\times m_T$, $b\times m_T$.
W pierwszym modelu do generatora $H_2(T^2\times D^2)$
przechodzi wyłącznie $a\times b$.
W drugim modelu analogicznie do generatora $H_2(P_K)$
przechodzi $s\times\mu_K$; pozostałe dwie klasy
zawierają $\lambda_K$ i zerują się, ponieważ
$[\lambda_K]=0$ w $H_1(E_K)$.
W stopniu zero oba modele są spójne, a w stopniu trzy
grupy homologii części dołączanych znikają. To
wyczerpuje możliwości.
\end{proof}

\begin{wniosek}[Klasyczne niezmienniki operacji]
\label{cor:tg-ks-invariants}
W sytuacji z definicji \ref{def:tg-ks-operation}
rozmaitości $X$ i $X_K$ mają izomorficzne
całkowite grupy homologii, a ponadto
$\chi(X_K)=\chi(X)$ i $\sigma(X_K)=\sigma(X)$.
\end{wniosek}

\begin{proof}
Dwie części dołączane mają takie same grupy homologii,
a stwierdzenie \ref{prop:tg-ks-boundary-maps} identyfikuje
odwzorowania z homologii wspólnego brzegu.
Dla każdego $i$ ciąg Mayera--Vietorisa zawiera krótki ciąg
$0\to\operatorname{coker}j_i\to H_i(X_K;\mathbb Z)
\to\ker j_{i-1}\to0$, gdzie $j_i$ oznacza mapę z
$H_i(T^3;\mathbb Z)$ do $H_i(A;\mathbb Z)\oplus H_i(P_K;\mathbb Z)$.
Odpowiedni ciąg dla $X$ ma izomorficzne jądra i kokernel.
Grupa $\ker j_{i-1}$ jest podgrupą wolnej grupy
$H_{i-1}(T^3;\mathbb Z)$, a więc sama jest wolna;
oba krótkie ciągi rozszczepiają się jako ciągi grup abelowych.
Otrzymujemy stąd izomorfizmy całkowitych grup homologii. Addytywność charakterystyki Eulera
względem rozkładu daje
$\chi(X_K)=\chi(A)+\chi(P_K)-\chi(T^3)$,
$\chi(X)=\chi(A)+\chi(T^2\times D^2)-\chi(T^3)$.
Ponieważ wszystkie trzy składniki zastępowane w tych
równaniach mają charakterystykę Eulera równą zero,
wzór dla $\chi$ jest udowodniony.

Równość sygnatur wynika z addytywności Nowikowa:
wymieniane części mają zerową sygnaturę swoich form
przecięcia, ponieważ ich druga homologia jest
generowana przez torus o zerowym samoprzecięciu.
W tym kroku używamy twierdzenia o addytywności
sygnatur przy sklejaniu wzdłuż wspólnego brzegu;
nie jest ono elementarną konsekwencją samych
liczb Betti'ego.
\end{proof}

\begin{uwaga}[Forma przecięcia a homologia]
Izomorfizm grup $H_2(X_K;\mathbb Z)\cong H_2(X;\mathbb Z)$
nie implikuje automatycznie izometrii form przecięcia.
W standardowej konstrukcji Fintushela--Sterna
odpowiednia izometria wynika z bardziej precyzyjnego
porównania par homologicznych i ich przecięć.
Będziemy jawnie powoływać się na ten fakt przy
klasyfikacji homeomorficznej. Sam rachunek $\chi$ i
$\sigma$ nie wystarcza do ustalenia parzystości formy.
\end{uwaga}

\section{Grupa fundamentalna i zachowanie typu homeomorfizmu}
\label{sec:tg-ks-pi1}

\begin{twierdzenie}[Warunek prostej spójności]
\label{thm:tg-ks-pi1}
Niech $X$ będzie gładką, spójną, zamkniętą
czterorozmaitością, a $T\subset X$ torusem o
trywialnej wiązce normalnej i połączonym dopełnieniu.
Jeżeli $\pi_1(X\setminus\operatorname{int}\nu T)=1$,
to dla dowolnego węzła $K\subset S^3$
rozmaitość $X_K$ ma trywialną grupę fundamentalną.
\end{twierdzenie}

\begin{proof}
Oznaczmy $A:=X\setminus\operatorname{int}\nu T$.
Z twierdzenia van Kampena zastosowanego do
$X_K=A\cup_\varphi P_K$ (po poszerzeniu kawałków
przez otwarte kołnierze) dostajemy iloraz wolnego
iloczynu $\pi_1(A)*\pi_1(P_K)$ przez normalną
domkniętą podgrupę generowaną przez relacje
$i_A(\varphi_*(\gamma))=i_P(\gamma)$,
gdzie $\gamma\in\pi_1(\partial P_K)$, zaś $i_A,i_P$
są homomorfizmami indukowanymi przez włączenia.
Ponieważ $\pi_1(A)=1$, iloraz ten ma postać
\[
\pi_1(X_K)\cong
\pi_1(P_K)\big/\langle\!\langle
  i_P\bigl(\pi_1(\partial P_K)\bigr)
\rangle\!\rangle.
\]
Obraz grupy brzegu jest generowany przez $s$, $\mu_K$
i $\lambda_K$. Ponieważ
$\pi_1(P_K)\cong\mathbb Z[s]\times\pi_1(E_K)$,
a grupa węzła $\pi_1(E_K)$ jest normalnie generowana
przez meridian $\mu_K$, to relacje $s=1$ i
$\mu_K=1$ już zabijają całą grupę. Otrzymujemy
$\pi_1(X_K)=1$.

Przypomnijmy uzasadnienie normalnej generacji:
wypełnienie meridionalne zewnętrza $E_K$
odtwarza $S^3$. Twierdzenie van Kampena dla tego
wypełnienia daje
$\pi_1(E_K)/\langle\!\langle\mu_K\rangle\!\rangle
\cong\pi_1(S^3)=1$.
\end{proof}

\begin{twierdzenie}[Wniosek topologiczny; wynik Fintushela--Sterna]
\label{thm:tg-ks-homeo}
Niech $X$ będzie gładką, zamkniętą, spójną,
zorientowaną i po prostu spójną czterorozmaitością.
Niech $T\subset X$ będzie torusem o $T^2=0$,
o otoczeniu typu \emph{cusp} i po prostu spójnym
dopełnieniu. Dla standardowej chirurgii
Fintushela--Sterna $X_K$ jest homeomorficzna z $X$
dla każdego węzła $K\subset S^3$.
\end{twierdzenie}

\begin{proof}[Odwołanie do twierdzenia zewnętrznego]
Udowodniliśmy, że $X_K$ jest po prostu spójna i ma
tę samą homologię, charakterystykę Eulera i sygnaturę.
Do zastosowania klasyfikacji Freedmana nie wystarczają
jednak te dane: trzeba jeszcze wykazać izometrię
całkowitej formy przecięcia oraz zgodność niezmiennika
Kirby'ego--Siebenmanna. W standardowej konstrukcji
Fintushela--Sterna wynikają one z analizy sklejenia
i ramowania torusa. Samą tę identyfikację przyjmujemy
z twierdzenia Fintushela--Sterna, a następnie
stosujemy twierdzenie Freedmana.
\end{proof}

\begin{uwaga}[Różne hipotezy mają różne role]
Prosta spójność dopełnienia pozwala zastosować
van Kampena. Warunek $b_2^+(X)>1$ będzie potrzebny
dopiero do prostej, bezkomorowej formuły Seiberga--Wittena.
Natomiast otoczenie typu \emph{cusp} jest naturalnym
założeniem zapewniającym standardową wersję twierdzenia
Fintushela--Sterna o niezmiennikach gładkich.
\end{uwaga}

\section{Wielomian Alexandera w formule Seiberga--Wittena}
\label{sec:tg-ks-sw}

Niech $X$ będzie zamkniętą, po prostu spójną czterorozmaitością
z $b_2^+(X)>1$. Niezmiennik Seiberga--Wittena można zapisać
jako skończony element pierścienia grupowego
$\mathbb Z[H_2(X;\mathbb Z)/\operatorname{Tor}]$:
\[
\mathcal{SW}_X=\sum_{\beta}SW_X(\beta)\,\exp(\beta),
\]
gdzie suma biegnie po klasach podstawowych, a
$\exp(\beta)$ jest formalnym symbolem o własności
$\exp(\beta)\exp(\gamma)=\exp(\beta+\gamma)$.
Nie jest to funkcja wykładnicza w sensie analizy.
Identyfikację odpowiednich grup homologii $X$
i $X_K$ ustalamy zgodnie z konstrukcją chirurgiczną.
Wielomianową postać niezmiennika stosujemy przede wszystkim
do rozmaitości typu prostego w sensie Seiberga--Wittena;
dla rozmaitości wymagających niezerowymiarowych moduli
konieczne są dodatkowe dane i pełniejsze wersje niezmiennika.

\begin{twierdzenie}[Fintushel--Stern: formuła \emph{knot surgery}]
\label{thm:tg-ks-formula}
Niech $X$ będzie gładką, zamkniętą, spójną,
zorientowaną i po prostu spójną czterorozmaitością
z $b_2^+(X)>1$. Niech $T\subset X$ będzie homologicznie
nietrywialnym, gładko zanurzonym torusem o $T^2=0$,
leżącym w otoczeniu typu \emph{cusp} i mającym
po prostu spójne dopełnienie. Niech
$K\subset S^3$ będzie zorientowanym węzłem,
a $\Delta_K(z)$ niech będzie symetrycznie
znormalizowanym wielomianem Alexandera, tj.
$\Delta_K(z^{-1})=\Delta_K(z)$ i $\Delta_K(1)=1$.
Wówczas dla standardowej chirurgii z definicji
\ref{def:tg-ks-operation} zachodzi
\[
\mathcal{SW}_{X_K}=\mathcal{SW}_X\,\Delta_K(u),
\qquad u:=\exp(2[T_K]).
\]
Dopuszczamy identyfikację $[T_K]$ z odpowiadającą
klasą $[T]$ w przyjętym izomorfizmie homologii.
\end{twierdzenie}

\begin{uwaga}[Status dowodu i podwójna zmienna]
Twierdzenie Fintushela--Sterna jest \emph{głębokim
wynikiem zewnętrznym}. Jego dowód opiera się na teorii
Seiberga--Wittena, odpowiednich względnych niezmiennikach
i relacjach skein. Nie wynika z ciągu Mayera--Vietorisa.
Niektóre źródła zapisują formułę jako
$\mathcal{SW}_{X_K}=\mathcal{SW}_X\Delta_K(t^2)$,
używając $t=\exp([T])$. W naszym zapisie
$u=\exp(2[T])$, więc są to dokładnie te same równania.
\end{uwaga}

\begin{stwierdzenie}[Klasy podstawowe po chirurgii]
\label{prop:tg-ks-basic}
W hipotezach twierdzenia \ref{thm:tg-ks-formula}
załóżmy dodatkowo, że $\mathcal{SW}_X=1$.
Jeżeli $\Delta_K(z)=\sum_{j=-g}^{g}a_jz^j$,
to $\mathcal{SW}_{X_K}=\sum_{j=-g}^{g}a_j\exp(2j[T_K])$.
W szczególności niezerowe współczynniki $a_j$
określają klasy podstawowe $2j[T_K]$, gdy
$[T_K]$ ma nieskończony rząd; licząc różne wykładniki,
uzyskujemy dokładnie liczbę niezerowych współczynników.
\end{stwierdzenie}

\begin{proof}
Wstawiamy rozwinięcie Laurenta wielomianu Alexandera
do formuły Fintushela--Sterna. Formalne mnożenie
w pierścieniu grupowym jest rozdzielne względem dodawania.
Jeżeli $[T_K]$ ma nieskończony rząd, klasy $2j[T_K]$
dla różnych $j$ są różne, a zatem nie następuje
łączenie odpowiadających im współczynników.
\end{proof}

\section{Przykłady: trójlistnik, węzeł ósemkowy i K3}
\label{sec:tg-ks-examples}

Rozważmy powierzchnię eliptyczną $E(2)$,
czyli powierzchnię K3. Włókno regularne
eliptycznej fibracji jest torusem o samoprzecięciu
zero. Można je wybrać tak, by miało otoczenie typu
\emph{cusp} oraz po prostu spójne dopełnienie.
Włókno $F$ ma klasę pierwotną w
$H_2(E(2);\mathbb Z)$: przekrój fibracji przecina
włókno z liczbą przecięcia jeden, więc gdyby
$[F]=m\alpha$ z $|m|>1$, otrzymalibyśmy
$1=[F]\cdot[S]=m(\alpha\cdot[S])$, sprzeczność.
W szczególności wykładniki $2j[F]$ dla różnych $j$
są różne w wolnej drugiej homologii.
Dla $E(2)$ zachodzi
$\pi_1(E(2))=1$, $b_2^+(E(2))=3$,
$\chi(E(2))=24$, $\sigma(E(2))=-16$, a jego
wielomian Seiberga--Wittena jest równy $1$ przy
odpowiednim ustaleniu orientacji homologicznej.
Stąd
\[
\mathcal{SW}_{E(2)_K}=\Delta_K(u),
\qquad u=\exp(2[T_K]).
\]

\begin{przyklad}[Chirurgia trójlistnika]
Dla trójlistnika $T_{2,3}$ mamy
$\Delta_{T_{2,3}}(z)=z^{-1}-1+z$. Wobec tego
\[
\mathcal{SW}_{E(2)_{T_{2,3}}}=u^{-1}-1+u.
\]
Rozmaitość $E(2)_{T_{2,3}}$ ma trzy klasy podstawowe
$-2[T_K]$, $0$ i $2[T_K]$ o współczynnikach
odpowiednio $1,-1,1$. Natomiast K3 ma tylko
jedną klasę podstawową, mianowicie zero.
Liczba klas podstawowych jest niezmiennikiem
dyfeomorfizmu, dlatego rozmaitości te nie są
dyfeomorficzne. Na mocy twierdzenia
\ref{thm:tg-ks-homeo} są natomiast homeomorficzne.
\end{przyklad}

\begin{przyklad}[Węzeł ósemkowy]
Dla węzła ósemkowego $4_1$ mamy
$\Delta_{4_1}(z)=3-z-z^{-1}$.
Zatem
$\mathcal{SW}_{E(2)_{4_1}}=3-u-u^{-1}$.
Oba omawiane węzły mają rodzaj jeden, lecz
ich wielomiany Alexandera są różne.
W konsekwencji powstałe czterorozmaitości mają
inny multizbiór bezwzględnych wartości niezerowych
współczynników SW. To wystarcza do wykazania,
że nie są dyfeomorficzne, bez konieczności ustalania
konkretnego izomorfizmu ich drugich homologii.
\end{przyklad}

\begin{center}
\begin{tikzpicture}[>=Stealth,line cap=round,line join=round]
 \node[draw=tgBlue,fill=tgBlue!10,rounded corners=4pt,
       text width=3cm,align=center,inner sep=5pt] (k3) at (0,1.2)
       {\small $E(2)$ (K3)\\ \footnotesize $\mathcal{SW}=1$};
 \node[draw=tgTeal,fill=tgTeal!12,rounded corners=4pt,
       text width=3.4cm,align=center,inner sep=5pt] (tr) at (-3.4,-1.1)
       {\small chirurgia z trójlistnikiem\\ \footnotesize $u^{-1}-1+u$};
 \node[draw=tgCoral,fill=tgCoral!10,rounded corners=4pt,
       text width=3.4cm,align=center,inner sep=5pt] (f) at (3.4,-1.1)
       {\small chirurgia z węzłem ósemkowym\\ \footnotesize $3-u-u^{-1}$};
 \draw[-{Stealth[length=2.5mm]},tgPurple,line width=1.7pt] (k3.south west)--(tr.north);
 \draw[-{Stealth[length=2.5mm]},tgPurple,line width=1.7pt] (k3.south east)--(f.north);
 \node[font=\footnotesize,text=tgNavy] at (0,-2.12)
       {Wspólny punkt wyjścia; różne wielomiany SW.};
\end{tikzpicture}
\end{center}

\section{Nieskończenie wiele egzotycznych struktur na K3}
\label{sec:tg-ks-infinite}

\begin{twierdzenie}[Rodzina egzotycznych struktur]
\label{thm:tg-ks-infinite}
Dla każdego $n\ge1$ niech $K_n=T(2,2n+1)$ będzie
węzłem torusowym o parametrach $(2,2n+1)$,
a $X_n:=E(2)_{K_n}$ wynikiem standardowej chirurgii
Fintushela--Sterna na regularnym włóknie eliptycznym.
Wtedy wszystkie $X_n$ są homeomorficzne z powierzchnią
K3, lecz każda para $X_n,X_m$ dla $n\ne m$
nie jest dyfeomorficzna.
\end{twierdzenie}

\begin{proof}
Wielomian Alexandera węzła torusowego
$T(2,2n+1)$ wyznaczymy ze wzoru podanego wcześniej:
\[
\Delta_{K_n}(z)
 =\frac{z^{2n+1}+1}{z+1}\,z^{-n}
 =\sum_{j=-n}^{n}(-1)^{n-j}z^j.
\]
Pierwsza równość jest tożsamością w pierścieniu
$\mathbb Z[z,z^{-1}]$: dzielenie jest dokładne,
ponieważ $2n+1$ jest nieparzyste. Wzór po prawej
wynika z rozwinięcia
$(z^{2n+1}+1)/(z+1)
=z^{2n}-z^{2n-1}+\cdots-z+1$.
Każdy z $2n+1$ współczynników jest równy $\pm1$
i jest niezerowy.

Twierdzenie Fintushela--Sterna daje
$\mathcal{SW}_{X_n}=\Delta_{K_n}(u)$.
Wynika stąd, że $X_n$ ma dokładnie $2n+1$
klas podstawowych. Dla $n\ne m$ liczby te są różne,
więc rozmaitości nie są dyfeomorficzne.
Twierdzenie \ref{thm:tg-ks-homeo} daje jednocześnie
homeomorfizmy $X_n\cong_{\mathrm{TOP}}E(2)$.
\end{proof}

\begin{uwaga}[Symplektyczność i włóknistość]
Węzły torusowe $T(2,2n+1)$ są włókniste.
Wówczas standardowa operacja Fintushela--Sterna
na włóknie eliptycznym może być zrealizowana
jako suma włóknista z symplektycznym
$S^1\times M_{K_n}$, gdzie $M_K$ oznacza
trójrozmaitość po chirurgii współczynnika zero
na $K$. Z konstrukcji Thurstona struktury symplektycznej
na $S^1\times M_K$ dla $K$ włóknistego oraz
z twierdzenia Gompfa o symplektycznej sumie
włóknistej wynika, że wszystkie $X_n$ mają struktury
symplektyczne. Te dwa twierdzenia wykorzystujemy
jako wyniki zewnętrzne. W szczególności jest to
rodzina \emph{symplektycznych} struktur egzotycznych
na topologicznej powierzchni K3.
\end{uwaga}

\section{Czego wielomian Alexandera nie wykrywa?}
\label{sec:tg-ks-limits}

\begin{stwierdzenie}[Ograniczenie formuły]
\label{prop:tg-ks-limit}
Niech $K_1,K_2\subset S^3$ będą dwoma węzłami,
dla których symetrycznie znormalizowane wielomiany
Alexandera są równe:
$\Delta_{K_1}=\Delta_{K_2}$.
Jeśli $X$ i $T$ spełniają hipotezy twierdzenia
\ref{thm:tg-ks-formula}, to po odpowiednich
identyfikacjach homologii zachodzi
$\mathcal{SW}_{X_{K_1}}=\mathcal{SW}_{X_{K_2}}$.
Nie wynika stąd, że $X_{K_1}$ i $X_{K_2}$
są dyfeomorficzne.
\end{stwierdzenie}

\begin{proof}
Pierwsza teza jest natychmiastowym podstawieniem
do wzoru Fintushela--Sterna. Druga wypowiedź oznacza
jedynie ograniczenie logiczne metody: równość tych
niezmienników nie jest twierdzeniem klasyfikacyjnym
wymuszającym dyfeomorfizm.
\end{proof}

\begin{przyklad}[Wielomian trywialny nie oznacza węzła trywialnego]
Istnieją nietrywialne węzły $K$ z $\Delta_K(z)=1$,
na przykład węzeł Kinoshity--Terasaki (Kinoshita--Terasaka).
Dla takich węzłów w formule z K3 otrzymujemy
$\mathcal{SW}_{E(2)_K}=1=\mathcal{SW}_{E(2)}$.
Jest to równość obliczonych niezmienników, a nie
dowód dyfeomorficzności. Z drugiej strony pokazuje,
że sam wielomian Alexandera nie może rozstrzygnąć
wszystkich pytań o strukturę gładką.
\end{przyklad}

\section{Porównanie z wymierną kontrakcją Parka}
\label{sec:tg-ks-compare}

Operacje wymiernej kontrakcji i chirurgii wzdłuż węzła
prowadzą do egzotycznych struktur gładkich przy
odmiennych założeniach. Dla Parka punktem wyjścia
jest ujemnie określony łańcuch sfer o samoprzecięciach
różnych od zera, zastępowany wymierną kulą homologii.
W konstrukcji Fintushela--Sterna modyfikujemy otoczenie
torusa o samoprzecięciu zero.

\begin{center}
\begin{tikzpicture}[>=Stealth,line cap=round,line join=round]
  \node[draw=tgCoral,fill=tgCoral!11,rounded corners=4pt,
        text width=4.2cm,align=center,inner sep=7pt] (a) at (-3.1,0)
        {\small wymierna kontrakcja\\
         \footnotesize łańcuch sfer $C_p$\\
         \footnotesize $b_2^+$ może wynosić $1$};
  \node[draw=tgTeal,fill=tgTeal!11,rounded corners=4pt,
        text width=4.2cm,align=center,inner sep=7pt] (b) at (3.1,0)
        {\small chirurgia wzdłuż węzła\\
         \footnotesize torus $T^2=0$\\
         \footnotesize klasyczna formuła SW: $b_2^+>1$};
  \draw[<->,tgPurple,line width=1.7pt] (a)--(b);
  \node[font=\footnotesize,text=tgNavy] at (0,-1.22)
        {Różne operacje, wspólny cel: wykrywanie typu DIFF.};
\end{tikzpicture}
\end{center}

Nie wolno bezpośrednio zastosować twierdzenia
\ref{thm:tg-ks-formula} do rozmaitości Parka $X_7$,
ponieważ $b_2^+(X_7)=1$. W tym przypadku
niezbędne jest rozróżnienie komór i osobne
argumenty analityczne omówione wcześniej.
Ponadto sama obecność zanurzonego torusa o
$T^2=0$ nie gwarantuje istnienia otoczenia typu
\emph{cusp} ani prostej spójności dopełnienia.

\begin{uwaga}[Co jest dowiedzione, a co przyjęte]
W tym rozdziale samodzielnie przeprowadziliśmy
obliczenia homologii zewnętrza węzła i części
$S^1\times E_K$, dowód prostej spójności
$X_K$ przy trywialnej grupie dopełnienia torusa
oraz rachunki wielomianowe dla wielu przykładów.
Korzystamy zewnętrznie z addytywności sygnatur,
izometrii form przecięcia w standardowej konstrukcji,
twierdzenia Freedmana, formuły Fintushela--Sterna
i konstrukcji symplektycznej sumy włóknistej.
Dopiero ostatni z tych pakietów twierdzeń
umożliwia wnioski o egzotycznych strukturach gładkich.
\end{uwaga}


% ITERACJA 16: synteza rachunków grafowych, chirurgicznych i cechowania
\chapter{Od grafów plumbingowych do egzotycznych struktur gładkich: synteza}
\label{ch:tg-synteza16}

Poprzednie rozdziały wprowadzały kolejne narzędzia w różnym języku:
macierze incydencji grafów, formy przecięcia, chirurgie trójrozmaitości,
rozkłady na uchwyty i niezmienniki Seiberga--Wittena. Obecnie zestawimy
je w jeden ciąg rozumowania, pilnując, by nie wyprowadzać wniosków
gładkich z danych wyłącznie topologicznych. Rozważymy dwie konstrukcje:
wymierną kontrakcję konfiguracji $C_7$ prowadzącą do rozmaitości Parka
oraz chirurgię Fintushela--Sterna na torusie w powierzchni K3.

\section{Jaki rodzaj informacji koduje graf?}
\label{sec:tg-syn-graph}

Niech $\Gamma$ będzie skończonym, spójnym drzewem o zbiorze
wierzchołków $V(\Gamma)=\{v_1,\ldots,v_r\}$, którego wierzchołkom
przypisano całkowite wagi $e_1,\ldots,e_r$. Zakładamy, że każda
krawędź oznacza dodatnie, poprzeczne skrzyżowanie dwóch bazowych sfer
w standardowym, zorientowanym plumbingowaniu wiązek dyskowych.
Przez $P_\Gamma$ oznaczamy wynikowe plumbingowanie, a przez
$Q_\Gamma$ macierz określoną wzorem
\[
 (Q_\Gamma)_{ij}=\begin{cases}
 e_i,&i=j,\\
 1,&i\ne j\text{ i }\{v_i,v_j\}\in E(\Gamma),\\
 0,&\text{w pozostałych przypadkach}.
 \end{cases}
\]
Macierz reprezentuje formę przecięcia na
$H_2(P_\Gamma;\mathbb Z)\cong\mathbb Z^r$ w bazie klas sfer zerowych.
To stwierdzenie zależy od określonego sposobu plumbingowania;
nie należy utożsamiać dowolnego grafu ważonego z jednoznacznie
wyznaczoną czterorozmaitością bez podania danych konstrukcji.

\begin{twierdzenie}[Homologia brzegu ujemnie określonego plumbingowania]
\label{thm:tg-syn-boundary}
Niech $\Gamma$ będzie drzewem z całkowitymi wagami, a
$P_\Gamma$ standardowym plumbingowaniem zorientowanych wiązek
dyskowych nad $S^2$. Niech $Y=\partial P_\Gamma$ z orientacją
brzegową i załóżmy, że $\det Q_\Gamma\ne0$. Wówczas
$H_1(Y;\mathbb Z)\cong\operatorname{coker}Q_\Gamma$,
$|H_1(Y;\mathbb Z)|=|\det Q_\Gamma|$ i
$H_1(Y;\mathbb Q)=0$. Jeśli $Q_\Gamma$ jest ujemnie określona,
to $\sigma(P_\Gamma)=-r$.
\end{twierdzenie}

\begin{proof}
Drzewo plumbingowe daje rozkład na uchwyty z jednym uchwytem indeksu
zero i $r$ uchwytami indeksu dwa, bez uchwytów indeksu jeden.
Stąd $P_\Gamma$ jest po prostu spójna,
$H_2(P_\Gamma;\mathbb Z)\cong\mathbb Z^r$ i
$H_1(P_\Gamma;\mathbb Z)=0$. W długim ciągu dokładnym pary
$(P_\Gamma,Y)$ otrzymujemy
\[
 H_2(P_\Gamma;\mathbb Z)\xrightarrow{j}
 H_2(P_\Gamma,Y;\mathbb Z)\longrightarrow
 H_1(Y;\mathbb Z)\longrightarrow0.
\]
Z dualności Poincar\'ego--Lefschetza i uniwersalnego twierdzenia
o współczynnikach wynika
$H_2(P_\Gamma,Y;\mathbb Z)\cong H^2(P_\Gamma;\mathbb Z)
\cong\operatorname{Hom}(H_2(P_\Gamma;\mathbb Z),\mathbb Z)$.
W tych identyfikacjach $j$ jest homomorfizmem wyznaczonym przez
formę przecięcia, a więc ma macierz $Q_\Gamma$. Jego kokernel
jest zatem grupą $H_1(Y;\mathbb Z)$. Niezerowy wyznacznik
implikuje, że kokernel jest skończony, o rzędzie równym
$|\det Q_\Gamma|$. Po przejściu do $\mathbb Q$ mapa staje się
izomorfizmem, więc $H_1(Y;\mathbb Q)=0$.
Jeżeli forma jest ujemnie określona, wszystkie jej $r$
wartości własnych są ujemne, stąd $\sigma(P_\Gamma)=-r$.
\end{proof}

\begin{uwaga}[Forma splecenia na brzegu]
W hipotezach twierdzenia i przy orientacji $Y=\partial P_\Gamma$
forma splecenia na $H_1(Y;\mathbb Z)$ jest reprezentowana
(z użyciem konwencji odpowiadającej tej orientacji) przez
$-Q_\Gamma^{-1}$ modulo $\mathbb Z$. Zmiana orientacji $Y$
zmienia znak tej formy. Sama grupa $H_1(Y)$, a nawet
forma splecenia, nie wyznaczają ogólnie typu dyfeomorfizmu $Y$.
Na przykład rozpoznanie konkretnej przestrzeni soczewkowej
wymaga dodatkowo informacji o sposobie sklejenia torusów pełnych.
\end{uwaga}

\begin{przyklad}[Łańcuch $C_p$ i ułamek ciągły]
\label{ex:tg-syn-cp}
Niech $p\ge2$ i niech $C_p$ będzie łańcuchem długości $p-1$
o wagach $(-(p+2),-2,\ldots,-2)$. Dla wyznacznika
$D_k$ lewego głównego minora rozmiaru $k$ macierzy przecięcia
mamy $D_0=1$, $D_1=-(p+2)$ i
$D_k=-2D_{k-1}-D_{k-2}$ dla $2\le k\le p-1$.
Indukcyjnie otrzymujemy
$D_k=(-1)^k(1+k(p+1))$: wzór zachodzi dla $k=0,1$,
a podstawienie do rekurencji daje
\[
 -2(-1)^{k-1}(1+(k-1)(p+1))
 -(-1)^{k-2}(1+(k-2)(p+1))
 =(-1)^k(1+k(p+1)).
\]
Stąd $\det Q_{C_p}=D_{p-1}=(-1)^{p-1}p^2$.
Ponieważ kolejne ilorazy $D_k/D_{k-1}$ nie zerują się,
a znaki $(-1)^kD_k$ są dodatnie, kryterium Sylvestera
wykazuje ujemną określoność $Q_{C_p}$.
Z kolei ujemny ułamek łańcuchowy opisujący brzeg ma postać
\[
[p+2,\underbrace{2,\ldots,2}_{p-2\text{ razy}}]^-
=p+2-\cfrac{1}{2-\cfrac{1}{\ddots-\cfrac12}}
=\frac{p^2}{p-1}.
\]
Aby uzasadnić ostatnią równość, zauważamy, że łańcuch
$k$ dwójek ma ułamek $(k+1)/k$ dla $k\ge1$;
po podstawieniu $k=p-2$ dostajemy
$p+2-(p-2)/(p-1)=p^2/(p-1)$.
Dla $p=2$ lista dwójek jest pusta i ułamek wynosi $4$,
co również zgadza się z wzorem.
Ten drugi rachunek, a nie sama wartość wyznacznika,
pozwala zidentyfikować brzeg jako odpowiednio zorientowaną
przestrzeń soczewkową rzędu $p^2$.
\end{przyklad}

\section{Dwie operacje wymiany: co rzeczywiście zachowują?}
\label{sec:tg-syn-exchange}

W każdej z dwóch omawianych operacji czterorozmaitość
jest sumą dwóch zwartych części o wspólnym trójwymiarowym
brzegu. Zastępujemy jedną część inną, nie zmieniając dopełnienia.
Dane potrzebne do obliczeń można zorganizować następująco.

\begin{center}
\begin{tikzpicture}[>=Stealth,line cap=round,line join=round,
    every node/.style={align=center}]
\node[draw=tgCoral,fill=tgCoral!10,rounded corners=4pt,
    text width=4.3cm,inner sep=6pt] (p) at (-3.25,1.2)
    {\small \textbf{wymierna kontrakcja}\\
     \footnotesize $C_p\rightsquigarrow B_p$\\
     \footnotesize wspólny brzeg $L(p^2,1-p)$};
\node[draw=tgTeal,fill=tgTeal!10,rounded corners=4pt,
    text width=4.3cm,inner sep=6pt] (k) at (3.25,1.2)
    {\small \textbf{chirurgia wzdłuż węzła}\\
     \footnotesize $T^2\times D^2\rightsquigarrow S^1\times E_K$\\
     \footnotesize wspólny brzeg $T^3$};
\node[draw=tgPurple,fill=tgPurple!9,rounded corners=4pt,
    text width=9cm,inner sep=6pt] (d) at (0,-1.05)
    {\small zachowane dopełnienie $A$ i identyfikacja brzegu\\
     \footnotesize Mayer--Vietoris, van Kampen oraz formy przecięcia};
\draw[-{Stealth[length=2.7mm]},tgPurple,line width=1.5pt] (p.south)--(-2.2,-.64);
\draw[-{Stealth[length=2.7mm]},tgPurple,line width=1.5pt] (k.south)--(2.2,-.64);
\end{tikzpicture}
\end{center}

\begin{twierdzenie}[Bilans wymiernej kontrakcji]
\label{thm:tg-syn-rational}
Niech $X$ będzie gładką, zamkniętą, spójną i zorientowaną
czterorozmaitością zawierającą standardową konfigurację
plumbingową $C_p$ o $p-1$ wierzchołkach. Niech $B_p$
będzie wymierną kulą homologii o brzegu odwzorowanym na
$\partial C_p$ przez dyfeomorfizm odwracający orientację,
a $X_{(p)}$ wynikiem zastąpienia $C_p$ przez $B_p$.
Wówczas
\[
 \chi(X_{(p)})=\chi(X)-(p-1),\qquad
 \sigma(X_{(p)})=\sigma(X)+(p-1).
\]
Jeśli ponadto $b_1(X_{(p)})=b_1(X)=0$, to
$\,b_2^+(X_{(p)})=b_2^+(X)$ oraz
$b_2^-(X_{(p)})=b_2^-(X)-(p-1)$.
\end{twierdzenie}

\begin{proof}
Plumbingowanie $C_p$ ma rozkład na uchwyty z jednym uchwytem
indeksu zero i $p-1$ uchwytami indeksu dwa,
zatem $\chi(C_p)=p$. Wymierna kula homologii ma
homologię wymierną punktu, więc $\chi(B_p)=1$.
Przy sklejeniu wzdłuż zamkniętej trójrozmaitości $Y$
zachodzi $\chi(X)=\chi(A)+\chi(C_p)-\chi(Y)$ i analogicznie
dla $X_{(p)}$, przy czym $\chi(Y)=0$.
Po odjęciu obu wzorów dostajemy
$\chi(X_{(p)})-\chi(X)=1-p$.
Forma przecięcia $C_p$ jest ujemnie określona rzędu $p-1$,
więc $\sigma(C_p)=-(p-1)$. Forma przecięcia $B_p$
ma zerowy wymiar nad $\mathbb Q$, dlatego $\sigma(B_p)=0$.
Z addytywności sygnatury Nowikowa przy sklejaniu otrzymujemy
drugi wzór.

Jeżeli $b_1=0$, dualność Poincar\'ego daje
$\chi=2+b_2$. Zatem $b_2$ maleje o $p-1$.
Rozwiązując układ $b_2=b_2^++b_2^-$,
$\sigma=b_2^+-b_2^-$, otrzymujemy oba wzory dla liczb
Betti'ego. Warunek $b_1=0$ jest istotny dla tego ostatniego
rachunku i nie wynika automatycznie z samej wymiany części.
\end{proof}

\begin{stwierdzenie}[Bilans chirurgii wzdłuż węzła]
\label{prop:tg-syn-knot}
Niech $X$ będzie zamkniętą, gładką i zorientowaną
czterorozmaitością z gładko zanurzonym torusem $T$
o trywialnej zorientowanej wiązce normalnej.
Niech $X_K$ będzie wynikiem standardowego sklejenia
z rozdziału \ref{ch:tg-knot-surgery}. Wówczas
$\chi(X_K)=\chi(X)$ i $\sigma(X_K)=\sigma(X)$.
Przy tym sklejeniu całkowite grupy homologii $X_K$ i $X$
są izomorficzne, ale poprawny dowód zachowania formy
przecięcia wymaga dodatkowego argumentu geometrycznego.
\end{stwierdzenie}

\begin{proof}
Dla $P_K=S^1\times E_K$ z poprzedniego rozdziału mamy
$H_i(P_K;\mathbb Z)\cong H_i(T^2\times D^2;\mathbb Z)$
we wszystkich stopniach. Co więcej, ustalone sklejenie
identyfikuje odwzorowania z $H_i(T^3)$ do homologii tych dwóch
części. Dla dwóch rozkładów z dopełnieniem $A$ ciągi
Mayera--Vietorisa dają więc w każdym stopniu krótkie ciągi
\[
 0\longrightarrow\operatorname{coker} j_i
 \longrightarrow H_i(X_K;\mathbb Z)
 \longrightarrow\ker j_{i-1}\longrightarrow0,
\]
gdzie $j_i$ jest sumą map z $H_i(T^3;\mathbb Z)$ do
$H_i(A;\mathbb Z)\oplus H_i(P_K;\mathbb Z)$.
Analogiczne krótkie ciągi występują dla $X$ i $T^2\times D^2$.
Ponieważ $H_i(T^3;\mathbb Z)$ jest grupą wolną, także
$\ker j_{i-1}$ jest wolna. Każdy taki ciąg rozszczepia się
jako ciąg grup abelowych. Identyfikacja jąder i kokerneli
zapewnia zatem izomorfizmy całkowitych grup homologii.

Obie części dołączane mają zerową charakterystykę Eulera,
a ich formy przecięcia mają zerową sygnaturę; po
zastosowaniu addytywności Nowikowa otrzymujemy
$\chi(X_K)=\chi(X)$ i $\sigma(X_K)=\sigma(X)$.
Samo rozszczepienie krótkich ciągów nie ustanawia jednak
izometrii form przecięcia: do jej ustalenia należy zbadać
geometrycznie reprezentantów klas i wybrane ramowanie torusa.
\end{proof}

\section{Grupa fundamentalna jako kontrola po sklejeniu}
\label{sec:tg-syn-pi1}

Gdy $X=A\cup_Y W$, twierdzenie van Kampena daje prezentację
$\pi_1(X)$ jako sumy amalgamowanej grup $\pi_1(A)$ i
$\pi_1(W)$ nad obrazami $\pi_1(Y)$. Poniższe dwa zastosowania
wyjaśniają, dlaczego samo obliczenie $H_1(Y)$ nie wystarcza.

\begin{stwierdzenie}[Prosta spójność po chirurgii wzdłuż węzła]
\label{prop:tg-syn-pi1}
Niech $X$ będzie spójną, zamkniętą, gładką czterorozmaitością,
$T\subset X$ torusem o trywialnej wiązce normalnej i niech
$A=X\setminus\operatorname{int}\nu T$ będzie spójna i po prostu
spójna. Wówczas $\pi_1(X_K)=1$ dla dowolnego węzła
$K\subset S^3$ i standardowego sklejenia.
\end{stwierdzenie}

\begin{proof}
Mamy $\pi_1(S^1\times E_K)\cong\mathbb Z\times\pi_1(E_K)$.
Przy $\pi_1(A)=1$ twierdzenie van Kampena daje
\[
 \pi_1(X_K)\cong
 \frac{\mathbb Z\times\pi_1(E_K)}
 {\langle\!\langle s,\mu_K,\lambda_K\rangle\!\rangle}.
\]
Meridian $\mu_K$ normalnie generuje $\pi_1(E_K)$: istotnie
meridionalne wypełnienie Dehna odtwarza $S^3$, więc
$\pi_1(E_K)/\langle\!\langle\mu_K\rangle\!\rangle=1$.
Ponadto relacja $s=1$ usuwa pozostały czynnik $\mathbb Z$.
Zatem grupa wynikowa jest trywialna.
\end{proof}

Dla wymiernej kontrakcji argument jest inny:
$\pi_1(B_p)$ i $\pi_1(\partial C_p)$ mogą być nietrywialne.
Aby uzyskać prostą spójność, trzeba sprawdzić obrazy ich
generatorów w dopełnieniu konfiguracji. W konstrukcji Parka
niezbędny fakt geometryczny polega na wskazaniu krzywej
brzegowej, która ogranicza dysk w dopełnieniu. Bez tego
kroku sam rachunek $\det Q_{C_7}=49$ nie wyznacza
$\pi_1(X_7)$.

\section{Formy przecięcia i klasyfikacja topologiczna}
\label{sec:tg-syn-freedman}

\begin{twierdzenie}[Obliczenie typu topologicznego rozmaitości Parka]
\label{thm:tg-syn-park-homeo}
Niech $X_7$ będzie gładką, zamkniętą, spójną,
zorientowaną i po prostu spójną czterorozmaitością
uzyskaną z $R=E(1)\#4\overline{\mathbb{CP}}^{\,2}$
przez standardową wymierną kontrakcję konfiguracji $C_7$.
Wtedy jej forma przecięcia jest nieparzysta, unimodularna
i nieokreślona, o sygnaturze $-6$ i randze $8$.
W konsekwencji, przy wykorzystaniu twierdzenia Freedmana,
$X_7$ jest homeomorficzna z
$\mathbb{CP}^2\#7\overline{\mathbb{CP}}^{\,2}$.
\end{twierdzenie}

\begin{proof}
Ponieważ $E(1)\cong\mathbb{CP}^2\#9\overline{\mathbb{CP}}^{\,2}$,
mamy $\chi(E(1))=12$ oraz $\sigma(E(1))=-8$.
Rozdmuchanie czterech punktów zwiększa charakterystykę Eulera
o $4$, a zmniejsza sygnaturę o $4$; stąd
$\chi(R)=16$, $\sigma(R)=-12$.
Konfiguracja $C_7$ ma sześć sfer, więc
z twierdzenia \ref{thm:tg-syn-rational} dostajemy
$\chi(X_7)=16-6=10$, $\sigma(X_7)=-12+6=-6$.
Z prostej spójności wynika $b_1(X_7)=0$ i
$\chi(X_7)=2+b_2(X_7)$, zatem $b_2(X_7)=8$.
Stąd $b_2^+=1$ i $b_2^-=7$. Forma przecięcia
na $H_2(X_7;\mathbb Z)$ jest unimodularna
z dualności Poincar\'ego, gdyż $H_1(X_7;\mathbb Z)=0$
i nie występuje torsja w $H_2$.
Forma nie może być parzysta: każda parzysta unimodularna
forma całkowita ma sygnaturę podzielną przez $8$,
a tutaj sygnatura wynosi $-6$. Z klasyfikacji nieparzystych
nieokreślonych form unimodularnych wynika więc
$Q_{X_7}\cong\langle1\rangle\oplus7\langle-1\rangle$.
Ponieważ rozmaitość jest gładka, jej niezmiennik
Kirby'ego--Siebenmanna wynosi zero. W tym miejscu stosujemy
głębokie twierdzenie Freedmana o klasyfikacji po prostu spójnych
zamkniętych topologicznych czterorozmaitości i otrzymujemy
też zgodność typu homeomorfizmu z rozmaitością standardową.
\end{proof}

\begin{uwaga}[Dlaczego nie wystarcza sam wzór na $\chi$ i $\sigma$?]
W dowodzie wykorzystano dodatkowo prostą spójność, dualność
Poincar\'ego, parzystość formy i klasyfikację unimodularnych
form całkowitych. Z dwóch liczb $\chi$ i $\sigma$ bez tych
hipotez nie można wyprowadzić homeomorfizmu.
\end{uwaga}

\begin{przyklad}[Topologiczny model chirurgii na K3]
Powierzchnia K3 ma $\chi=24$, $\sigma=-16$,
$b_2^+=3$, $b_2^-=19$ oraz parzystą unimodularną
formę przecięcia $2(-E_8)\oplus3H$, gdzie
$H=\begin{psmallmatrix}0&1\\1&0\end{psmallmatrix}$.
Z twierdzenia Fintushela--Sterna o standardowej chirurgii
na włóknie eliptycznym wynika, że wszystkie $E(2)_K$
mają tę samą formę przecięcia i są po prostu spójne.
Twierdzenie Freedmana daje ich homeomorfizm z K3.
Sama równość $\chi=24$ i $\sigma=-16$ nie gwarantowałaby
parzystości formy ani takiego wniosku.
\end{przyklad}

\section{Dwa mechanizmy wykrywania różnic gładkich}
\label{sec:tg-syn-sw}

W tym miejscu kończą się metody klasycznej topologii.
Gładka klasyfikacja rozmaitości wymaga odmiennych informacji.
Poniżej porównamy dwa sposoby użycia niezmienników
Seiberga--Wittena.

\begin{stwierdzenie}[Dwa reżimy niezmiennika SW]
\label{prop:tg-syn-regimes}
Niech $X$ będzie gładką, zamkniętą, zorientowaną
czterorozmaitością z $b_1(X)=0$.
Jeżeli $b_2^+(X)>1$, zliczanie rozwiązań równań
Seiberga--Wittena dla generycznych danych wyznacza
niezmienniki niezależne od wyboru komory. Jeśli natomiast
$b_2^+(X)=1$, w ogólności trzeba wskazać komorę
danych metrycznych i perturbacji; wartości w dwóch
komorach mogą się różnić. W szczególności bezkomorowa
formuła Fintushela--Sterna nie może być automatycznie
zastosowana do rozmaitości Parka $X_7$.
\end{stwierdzenie}

\begin{proof}[Uzasadnienie zakresu]
Przy $b_2^+>1$ locus danych, dla których pojawiają się
rozwiązania redukowalne, ma kodimension co najmniej dwa,
a więc generyczną jednoparametrową rodzinę można wybrać
tak, by tego locus unikała. Nie ma wtedy zmiany komory
podczas porównania dwóch generycznych danych.
Przy $b_2^+=1$ locus jest hiperpłaszczyzną o kodimension jeden:
może rozdzielać przestrzeń danych, a przekraczanie ściany
prowadzi do poprawki w zliczaniu rozwiązań.
Dokładny dowód transwersalności oraz formuły przejścia przez
ścianę wymaga analizy przestrzeni moduli i został opisany
w rozdziałach o równaniach Seiberga--Wittena. Tutaj korzystamy
z tych wyników jako zewnętrznej teorii analitycznej.
\end{proof}

\begin{przyklad}[Park: jeden niezerowy niezmiennik wystarcza]
Dla $X_7$ mamy $b_2^+=1$. W skrypcie obliczono
$K_{X_7}^2=2$ i udowodniono, przy wykorzystaniu
geometrycznego wyboru formy symplektycznej, że
$K_{X_7}\cdot[\omega_7]>0$. Twierdzenia Taubesa--Li--Liu
pozwalają wybrać właściwą komorę, w której odpowiedni
niezmiennik SW ma wartość $\pm1$. Argument z metryką
o dodatniej krzywiźnie skalarnej wyklucza dyfeomorfizm
ze standardową $\mathbb{CP}^2\#7\overline{\mathbb{CP}}^{\,2}$.
Nie jest to konsekwencja samej formy przecięcia, lecz
głębokiej teorii gładkiej i poprawnego doboru komory.
\end{przyklad}

\begin{twierdzenie}[K3: nieograniczona rodzina klas podstawowych]
\label{thm:tg-syn-k3}
Dla każdego $n\ge1$ niech $K_n=T(2,2n+1)$,
a $X_n=E(2)_{K_n}$ niech będzie wynikiem standardowej
chirurgii Fintushela--Sterna na włóknie eliptycznym K3.
Wówczas $X_n$ są homeomorficzne z K3, a jeśli $m\ne n$,
to $X_m$ i $X_n$ nie są dyfeomorficzne.
\end{twierdzenie}

\begin{proof}
Powierzchnia K3 ma $b_2^+=3$ i $\mathcal{SW}_{E(2)}=1$.
Twierdzenie Fintushela--Sterna daje więc
$\mathcal{SW}_{X_n}=\Delta_{K_n}(u)$,
$u=\exp(2[T_n])$. Dla węzła torusowego $K_n$,
\[
 \Delta_{K_n}(u)=\sum_{j=-n}^n(-1)^{n-j}u^j.
\]
Klasa włókna ma nieskończony rząd, więc wykładniki
$2j[T_n]$ są parami różne. Każdy ze współczynników
jest niezerowy, zatem $X_n$ ma dokładnie $2n+1$
klas podstawowych Seiberga--Wittena.
Liczba klas podstawowych jest niezmiennikiem dyfeomorfizmu,
więc dla różnych $n$ rozmaitości są niedyfeomorficzne.
Homeomorfizm z K3 wynika z zachowania formy przecięcia,
prostej spójności i twierdzenia Freedmana, zgodnie
z poprzednią sekcją.
\end{proof}

\begin{uwaga}[Rachunek współczynników, nie tylko liczby klas]
Dla dwóch dowolnych węzłów liczba niezerowych współczynników
nie zawsze rozróżnia rozmaitości. Jeżeli
$\mathcal{SW}_X=\sum_\beta c_\beta\exp(\beta)$ oraz
$\Delta_K(u)=\sum_j a_ju^j$, to współczynnik przy
$\exp(\gamma)$ po chirurgii jest równy
\[
 \sum_{\beta+2j[T]=\gamma}c_\beta a_j.
\]
Różne pary $(\beta,j)$ mogą prowadzić do tej samej klasy
$\gamma$ i wtedy ich wkłady się sumują, niekiedy się znosząc.
Dlatego sama liczba wyrazów $\Delta_K$ w ogólnym przypadku
nie wyznacza liczby klas podstawowych $X_K$.
\end{uwaga}

\section{Schemat hipotez i zależności między twierdzeniami}
\label{sec:tg-syn-assumptions}

Podajemy zestawienie ułatwiające rozpoznanie, jakie dane
należy sprawdzić, zanim zastosuje się określone twierdzenie.
Rzędy tabeli są niezależnymi implikacjami, nie automatycznym
łańcuchem równoważności.

\begin{center}
\begingroup
\footnotesize
\renewcommand{\arraystretch}{1.4}
\begin{tabular}{@{}p{.255\linewidth}p{.365\linewidth}p{.30\linewidth}@{}}
\hline
\textbf{Narzędzie} & \textbf{Kluczowe założenia} & \textbf{Wniosek}\\
\hline
{\raggedright Plumbing i dualność Poincar\'ego--Lefschetza\par}
 & {\raggedright zwarty orientowany blok, znana całkowita macierz przecięcia\par}
 & {\raggedright $H_1(\partial P)\cong\operatorname{coker}Q$ przy $H_1(P)=0$\par}\\
{\raggedright Wymierna kontrakcja\par}
 & {\raggedright konfiguracja $C_p$, zgodne sklejenie $B_p$, kontrola $\pi_1$\par}
 & {\raggedright zmiana $\chi,\sigma$; przy dalszych warunkach prosta spójność\par}\\
{\raggedright Freedman\par}
 & {\raggedright zamknięta, po prostu spójna topologiczna $4$-rozmaitość, forma i dane KS\par}
 & {\raggedright klasyfikacja topologiczna, nie gładka\par}\\
{\raggedright Fintushel--Stern\par}
 & {\raggedright $b_2^+>1$, torus w otoczeniu typu \emph{cusp}, po prostu spójne dopełnienie\par}
 & {\raggedright formuła $\mathcal{SW}_{X_K}=\mathcal{SW}_X\Delta_K(u)$\par}\\
{\raggedright Taubes--Li--Liu\par}
 & {\raggedright zwarta rozmaitość symplektyczna; właściwa komora przy $b_2^+=1$\par}
 & {\raggedright niezerowość odpowiedniego SW i ograniczenia gładkie\par}\\
\hline
\end{tabular}
\endgroup
\end{center}

\begin{center}
\begin{tikzpicture}[>=Stealth,line cap=round,line join=round,
  every node/.style={align=center}]
  \node[draw=tgCoral,fill=tgCoral!10,rounded corners=4pt,
    text width=3.35cm,inner sep=6pt] (p0) at (-4.0,1.65)
    {\small $E(1)\#4\overline{\mathbb{CP}}^{\,2}$\\
     \footnotesize $C_7$, sześć sfer};
  \node[draw=tgCoral,fill=tgCoral!10,rounded corners=4pt,
    text width=3.35cm,inner sep=6pt] (p1) at (4.0,1.65)
    {\small $X_7$\\ \footnotesize $b_2^+=1$, $K\cdot\omega>0$};
  \node[draw=tgTeal,fill=tgTeal!10,rounded corners=4pt,
    text width=3.35cm,inner sep=6pt] (k0) at (-4.0,-1.65)
    {\small $E(2)$ (K3)\\ \footnotesize torus $T^2=0$};
  \node[draw=tgTeal,fill=tgTeal!10,rounded corners=4pt,
    text width=3.35cm,inner sep=6pt] (k1) at (4.0,-1.65)
    {\small $E(2)_K$\\ \footnotesize $b_2^+=3$, $\Delta_K(u)$};
  \draw[-{Stealth[length=2.7mm]},tgPurple,line width=1.7pt]
    (p0)--node[fill=white,font=\scriptsize,above=3pt]{wymierna kontrakcja}(p1);
  \draw[-{Stealth[length=2.7mm]},tgPurple,line width=1.7pt]
    (k0)--node[fill=white,font=\scriptsize,above=3pt]{chirurgia wzdłuż węzła}(k1);
  \node[text=tgNavy,font=\footnotesize] at (0,0)
    {ten sam język topologii, inne hipotezy teorii cechowania};
\end{tikzpicture}
\end{center}

\section{Granice przedstawionej syntezy}
\label{sec:tg-syn-limits}

Otrzymany obraz nie jest dowodem, że każdy węzeł albo każdy graf
plumbingowy daje egzotyczną czterorozmaitość. Potrzebne są
precyzyjne warunki istnienia zanurzeń i kontroli dopełnienia,
a także niezerowość właściwych niezmienników gładkich.

Z grafu plumbingowego można odczytać macierz przecięcia
\emph{określonego plumbingowania}, lecz nie wszystkie dane sklejenia
dowolnej rozmaitości. Z wielomianu Alexandera można obliczyć
pewne niezmienniki SW po chirurgii Fintushela--Sterna,
lecz dwa różne węzły mogą mieć ten sam wielomian.
Wymienione operacje są narzędziami konstrukcyjnymi,
a nie pełnymi klasyfikacjami w kategorii DIFF.

\begin{uwaga}[Jak korzystać z tego rozdziału]
Przy rozwiązywaniu konkretnego problemu należy przechodzić
w podanej kolejności: (1) dokładnie zdefiniować części i sklejenie;
(2) obliczyć homologię i grupę fundamentalną;
(3) wyznaczyć formę przecięcia i rozważyć typ homeomorfizmu;
(4) dopiero potem wybrać właściwe niezmienniki gładkie,
sprawdzając wszystkie ich hipotezy. Taka kolejność zapobiega
przeskokom od diagramu kombinatorycznego do wniosku o strukturze
gładkiej bez dowodu.
\end{uwaga}


\begin{uwaga}[Zakres wniosków klasyfikacyjnych]
Zgodność $e$, $\sigma$ i liczb Betti'ego nie klasyfikuje
czterorozmaitości. W przypadku rozmaitości po prostu
spójnych trzeba kontrolować całą całkowitą formę
przecięcia i odpowiednią hipotezę twierdzenia Freedmana.
Jeżeli rozmaitości mają niezerową grupę fundamentalną,
potrzebne są dalsze dane i twierdzenia klasyfikacyjne.
Niezmienniki Seiberga--Wittena badają strukturę gładką,
ale także nie stanowią pełnego układu niezmienników
klasyfikujących wszystkie czterorozmaitości.
\end{uwaga}

\chapter{Nowoczesne niezmienniki i zastosowania teorii węzłów}
\label{ch:kt-nowoczesne}

Klasyczne wielomiany i grupy wystarczają do rozwiązania wielu
problemów, ale nie rozróżniają wszystkich subtelności gładkiej
topologii czterowymiarowej. Współczesne niezmienniki przypisują
węzłowi nie tylko liczbę albo wielomian, lecz cały kompleks
łańcuchowy i jego homologie. Dają one nowe ograniczenia
na gładki rodzaj czterowymiarowy oraz rozpoznają własności
takie jak włóknistość.

\section{Homologia Khovanova: od wielomianu do kompleksu}
\label{sec:kt-khovanov}

\begin{definicja}[Algebra Frobeniusa]
Niech $A:=\mathbb Z[X]/(X^2)$ będzie wolną grupą abelową
z bazą $(1,X)$ i gradacją $\deg(1)=1$, $\deg(X)=-1$.
Mnożenie $m\colon A\otimes A\to A$ jest indukowane przez
mnożenie wielomianów, a komnożenie
$\Delta\colon A\to A\otimes A$ jest dane wzorami
$\Delta(1)=1\otimes X+X\otimes1$,
$\Delta(X)=X\otimes X$. Odwzorowania $m$ i $\Delta$
są jednorodne stopnia $-1$ względem powyższej gradacji.
\end{definicja}

\begin{definicja}[Stany kostki rozwiązań]
Niech $D$ będzie zorientowanym diagramem splotu
$L\subset S^3$ o $n$ uporządkowanych skrzyżowaniach.
Dla każdego $s=(s_1,\dots,s_n)\in\{0,1\}^n$
wygładzamy skrzyżowanie numer $j$ według konwencji
0 lub 1. Otrzymujemy rozłączną sumę $k(s)$ okręgów.
Niech $|s|:=s_1+\dots+s_n$.
Przypisujemy stanowi moduł $A^{\otimes k(s)}$, przyjmując
$A^{\otimes0}=\mathbb Z$. Oznaczmy przez $n_+$ oraz $n_-$
liczby dodatnich i ujemnych skrzyżowań diagramu.
\end{definicja}

Dla jasności podajemy jedną z najczęściej stosowanych
konwencji przesunięć gradacji. Jeżeli $M\{r\}$ oznacza
moduł, w którym stopnie elementów zostały zwiększone o $r$,
to definiujemy grupę łańcuchową w stopniu homologicznym $i$:
\[
C^{i}(D):=
\bigoplus_{\substack{s\in\{0,1\}^n\\|s|-n_-=i}}
A^{\otimes k(s)}\{|s|+n_+-2n_-\}.
\]
Krawędź kostki zmieniająca jedno 0 w 1 odpowiada
albo połączeniu dwóch okręgów w jeden, albo rozszczepieniu
jednego okręgu na dwa. Odpowiedni homomorfizm jest więc
$m$ albo $\Delta$. Sumując te homomorfizmy ze znakami
kostki, otrzymujemy różniczkę $d\colon C^i(D)\to C^{i+1}(D)$.

\begin{uwaga}[Znaki różniczki i jej obliczanie]
Uporządkujmy skrzyżowania $1,\dots,n$. Jeżeli stan
$s\in\{0,1\}^n$ ma $s_j=0$, to krawędzi zmieniającej
$s_j$ na $1$ przypisujemy znak
$(-1)^{\sum_{k<j}s_k}$. Pozostała część operatora tej
krawędzi jest odpowiednim mnożeniem $m$ albo
komnożeniem $\Delta$. Dla dowolnych $j<k$ dwie
drogi w kwadracie kostki otrzymują przeciwne znaki,
ponieważ w jednej kolejności zmiana $s_j$ zwiększa
o jeden wykładnik znaku przy zmianie $s_k$, a w drugiej
nie. Otrzymujemy w ten sposób jednoznaczną, skończoną
procedurę zapisu macierzy różniczki.
\end{uwaga}

\begin{stwierdzenie}[Różniczka kwadratuje się do zera]
Dla dowolnego zorientowanego diagramu splotu $D$ i ustalonego
porządku jego skrzyżowań operator kostkowy
$d:C^i(D)\to C^{i+1}(D)$ spełnia $d^2=0$.
\end{stwierdzenie}

\begin{proof}
Każda składowa $d^2$ prowadząca ze stanu $s$ do stanu,
w którym dwa różne skrzyżowania $j,k$ zmieniły wartość
z $0$ na $1$, jest sumą dwóch kompozycji: najpierw $j$,
potem $k$, oraz w odwrotnej kolejności. Operacje
łączenia i rozszczepiania okręgów opisane przez $m$
i $\Delta$ tworzą przemienny kwadrat algebraiczny.
Jego przemienność wynika odpowiednio z łączności
$m$, współłączności $\Delta$ albo tożsamości Frobeniusa
\[
\Delta m=(m\otimes\mathrm{id})(\mathrm{id}\otimes\Delta)
=(\mathrm{id}\otimes m)(\Delta\otimes\mathrm{id}).
\]
Przy $j<k$ wykładniki znaku przy drogach $j\to k$
i $k\to j$ różnią się dokładnie o jeden: zmiana $s_j$
powoduje dodatkowy składnik w sumie $\sum_{\ell<k}s_\ell$,
a zmiana $s_k$ nie wpływa na znak krawędzi $j$.
Dlatego obie kompozycje mają przeciwne znaki
i ich suma jest zerowa. Nie ma składnika $d^2$
zmieniającego tylko jeden stan, ponieważ każda krawędź
kostki zwiększa liczbę jedynek o jeden. Stąd $d^2=0$.
\end{proof}

\begin{twierdzenie}[Niezmienniczość homologii Khovanova; wynik zewnętrzny]
Po uwzględnieniu wskazanych przesunięć gradacji
homologia kompleksu $(C^*(D),d)$ nie zależy od wyboru
regularnego zorientowanego diagramu danego splotu,
z dokładnością do izomorfizmu zachowującego obie gradacje.
\end{twierdzenie}

\begin{uwaga}[Zakres dowodu]
Pełna niezmienniczość wymaga skonstruowania dla każdego
z trzech ruchów Reidemeistera równoważności homotopijnych
odpowiednich kompleksów, z rozdzieleniem składników
kontraktowalnych i kontrolą przesunięć gradacji.
Nie jest wnioskiem z samego równania $d^2=0$.
Przyjmujemy ją w tym miejscu jako twierdzenie Khovanova;
poprzedni dowód dotyczy tylko formalnej konstrukcji kompleksu.
\end{uwaga}

\begin{definicja}[Homologia Khovanova]
Dla zorientowanego splotu $L\subset S^3$ przez
$\operatorname{Kh}^{i,j}(L;\mathbb Z)$ oznaczamy dwugradowaną
homologię kompleksu zbudowanego z dowolnego diagramu $L$
według poprzedniej konstrukcji, w stopniach homologicznym $i$
i kwantowym $j$. \emph{Zredukowana homologia Khovanova}
$\widetilde{\operatorname{Kh}}(L;\mathbb Z)$ powstaje
z odpowiednio oznaczonego diagramu przez redukcję
jednego czynnika algebry $A$; dla węzła wynik jest
niezależny od wyboru składowej bazowej.
\end{definicja}

\begin{uwaga}[Kategoryfikacja]
Wielomian Jonesa otrzymujemy, po odpowiednim podstawieniu
zmiennej i przyjęciu normalizacji, jako gradującą
charakterystykę Eulera kompleksu Khovanova:
$\sum_{i,j}(-1)^i q^j\operatorname{rank}
\operatorname{Kh}^{i,j}(L;\mathbb Z)$.
Dwa sploty mogą mieć jednakowy wielomian Jonesa, lecz
różne homologie Khovanova. Homologia zachowuje informacje
pominięte przy obliczeniu samej charakterystyki Eulera.
\end{uwaga}

\begin{przyklad}[Węzeł trywialny]
Diagram węzła trywialnego nie ma skrzyżowań.
Jego kompleks Khovanova składa się wyłącznie z
$C^0(U)=A$ i zerowej różniczki. Stąd
$\operatorname{Kh}(U;\mathbb Z)\cong\mathbb Z\oplus\mathbb Z$
z generatorami w stopniach kwantowych $\pm1$.
Po redukcji $\widetilde{\operatorname{Kh}}(U;\mathbb Z)
\cong\mathbb Z$ (z odpowiednio wybraną normalizacją gradacji).
\end{przyklad}

\begin{twierdzenie}[Kronheimer--Mrowka; homologia Khovanova wykrywa węzeł trywialny]
Niech $K\subset S^3$ będzie węzłem. Jeżeli
$\operatorname{rank}_{\mathbb Q}
\widetilde{\operatorname{Kh}}(K;\mathbb Q)=1$,
to $K$ jest trywialny. Odwrotność wynika bezpośrednio
z obliczenia homologii Khovanova dla węzła trywialnego.
\end{twierdzenie}

\begin{uwaga}
To twierdzenie jest zdecydowanie silniejsze niż
stwierdzenie, że odpowiednio znormalizowany wielomian
Jonesa węzła trywialnego wynosi jeden. Nie wiadomo,
czy warunek $V_K(t)=1$ sam w sobie wystarcza
rozstrzygnąć trywialność $K$; homologia Khovanova
rozpoznaje ją na poziomie dodatkowej struktury homologicznej.
Dowód Kronheimera--Mrowki wykorzystuje teorię instantonów.
\end{uwaga}

\section{Niezmiennik Rasmussena i gładki rodzaj}
\label{sec:kt-rasmussen}

Homologia Khovanova posiada deformację Lee nad $\mathbb Q$,
w której dla węzła homologia ma wymiar dwa. Filtracja
tej homologii pozwala zdefiniować całkowitą liczbę parzystą
$s(K)$, zwaną \emph{niezmiennikiem Rasmussena}.

\begin{twierdzenie}[Rasmussen]
Istnieje niezmiennik $s$ zorientowanych węzłów w $S^3$
o wartościach w $2\mathbb Z$, który spełnia:
\begin{enumerate}[label=(\roman*)]
\item $s(U)=0$;
\item $s(K_1\#K_2)=s(K_1)+s(K_2)$;
\item $s(\overline K)=-s(K)$;
\item $|s(K)|\leq2g_4^{\mathrm{sm}}(K)$;
\item dla dodatnich względnie pierwszych $p,q\geq2$,
$s(T(p,q))=(p-1)(q-1)$.
\end{enumerate}
\end{twierdzenie}

\begin{wniosek}
Dla dodatnich węzłów torusowych $T(p,q)$,
$p,q\geq2$ względnie pierwszych, mamy
$g_4^{\mathrm{sm}}(T(p,q))=(p-1)(q-1)/2$.
\end{wniosek}

\begin{proof}
Z nierówności Rasmussena i jego obliczenia dla węzłów
torusowych otrzymujemy
$(p-1)(q-1)/2=|s(T(p,q))|/2\leq g_4^{\mathrm{sm}}(T(p,q))$.
Z drugiej strony $g_4^{\mathrm{sm}}(K)\leq g(K)$,
a klasyczny rachunek rodzaju węzła torusowego daje
$g(T(p,q))=(p-1)(q-1)/2$. Obie nierówności są więc
równościami.
\end{proof}

\begin{uwaga}
Niezmiennik $s$ jest przykładem rezultatu typowo
czterowymiarowego wydobytego z teorii diagramowej.
Jego obliczanie wymaga homologii i filtracji, ale nie
wymaga bezpośredniej analizy równań teorii cechowania.
To jest jedna z największych zalet podejścia Khovanova.
\end{uwaga}

\section{Homologia Floera węzła}
\label{sec:kt-hfk}

\begin{definicja}[Kompleks Floera węzła; zakres konstrukcji]
Niech $K\subset S^3$ będzie zorientowanym węzłem, a
$\mathbb F=\mathbb Z/2\mathbb Z$. Wybieramy dopuszczalny
dwupunktowy diagram Heegaarda $(\Sigma,\boldsymbol\alpha,
\boldsymbol\beta,w,z)$ opisujący parę $(S^3,K)$.
Jego generatory są punktami przecięcia torusów
$\mathbb T_{\alpha}$ i $\mathbb T_{\beta}$ w
$\operatorname{Sym}^{g}(\Sigma)$, a różniczka liczona jest
przez klasy pseudoholomorficznych dysków Whitney'a
o indeksie Maslova jeden. Zliczenie ich z wagą
$U^{n_w(\phi)}$ daje kompleks $CFK^-(K)$ nad
$\mathbb F[U]$ z filtracją Alexandera.

\emph{Wersja z daszkiem} jest związana z ustawieniem $U=0$
i filtrowanym kompleksem $\widehat{CFK}(K)$. Przez
$\widehat{HFK}(K,s)$ oznaczamy homologię odpowiedniego
składnika stopnia Alexandera $s$ w stowarzyszonym
kompleksie stopniowanym. W przeciwieństwie do samej
homologii $\widehat{HF}(S^3)$ nie zapominamy więc filtracji
Alexandera. Klasa izomorfizmu grup $\widehat{HFK}(K,s)$
nie zależy od wyboru dopuszczalnego diagramu Heegaarda.
\end{definicja}

\begin{uwaga}
Powyższa konstrukcja nie jest elementarna: do
pełnego dowodu, że $\partial^2=0$ i że homologia
nie zależy od diagramu Heegaarda, potrzebna jest
teoria przestrzeni moduli krzywych pseudoholomorficznych.
W praktyce obliczeniowej używa się również kombinatorycznych
modeli siatkowych (grid homology), które pozwalają liczyć
niezmienniki Floera dla konkretnych diagramów.
\end{uwaga}

\begin{twierdzenie}[Ozsváth--Szabó; wykrywanie rodzaju]
Niech $K\subset S^3$ będzie węzłem. Wtedy
\[
g(K)=\max\{s\in\mathbb Z:\widehat{HFK}(K,s)\ne0\}.
\]
W szczególności najwyższy niezerowy stopień Alexandera
w homologii Floera wyznacza minimalny rodzaj powierzchni
Seiferta.
\end{twierdzenie}

\begin{twierdzenie}[Ni; wykrywanie włóknistości]
Niech $K\subset S^3$ będzie węzłem rodzaju $g(K)$.
Węzeł $K$ jest włóknisty wtedy i tylko wtedy, gdy
$\operatorname{rank}_{\mathbb F}\widehat{HFK}(K,g(K))=1$
(dla standardowej wersji homologii Floera nad ciałem $\mathbb F$).
\end{twierdzenie}

\begin{definicja}[Niezmiennik $\tau$]
Dla węzła $K\subset S^3$ niezmiennik
$\tau(K)\in\mathbb Z$ definiujemy za pomocą filtracji
Alexandera na kompleksie Floera: jest to najmniejszy
poziom filtracji, którego włączenie indukuje
niezerowy homomorfizm do jednowymiarowej
$\widehat{HF}(S^3)$ w odpowiednim stopniu Maslova.
\end{definicja}

\begin{twierdzenie}[Ozsváth--Szabó; własności $\tau$]
Dla dowolnych zorientowanych węzłów $K,K_1,K_2\subset S^3$
zachodzą:
$\tau(U)=0$,
$\tau(K_1\#K_2)=\tau(K_1)+\tau(K_2)$,
$\tau(\overline K)=-\tau(K)$ oraz
$|\tau(K)|\leq g_4^{\mathrm{sm}}(K)$.
Dla dodatnich węzłów torusowych $p,q\geq2$,
$\tau(T(p,q))=(p-1)(q-1)/2$.
\end{twierdzenie}

\begin{przyklad}[Rachunek, który odróżnia TOP od DIFF]
Niech $T=T(2,3)$ będzie dodatnim trójlistnikiem,
a $D=\operatorname{Wh}^{+}(T)$ jego dodatnim,
nieskręconym podwojeniem Whiteheada. Z definicji
satelity otrzymujemy $\Delta_D(t)=1$.
Z twierdzenia Freedmana $g_4^{\mathrm{top}}(D)=0$.
Z obliczenia Heddena $\tau(D)=1$, a więc
$1\leq g_4^{\mathrm{sm}}(D)$. Diagram Whiteheada daje
powierzchnię Seiferta rodzaju jeden, zatem
$g_4^{\mathrm{sm}}(D)\leq1$.
Ostatecznie
$g_4^{\mathrm{top}}(D)=0$ oraz $g_4^{\mathrm{sm}}(D)=1$.
To przykład konkretnego węzła, dla którego kategorie
topologiczna i gładka dają różne odpowiedzi.
\end{przyklad}

\section{Związki z algebrami operatorów i fizyką kwantową}
\label{sec:kt-operator}

Wielomian Jonesa pojawił się w teorii podczynników algebr
von Neumanna, gdzie bada się inkluzje czynników
$N\subset M$ oraz indeks $[M:N]$. W odpowiednim
formalizmie projekcje Jonesa spełniają relacje
algebry Temperleya--Lieba. Z ich pomocą konstruuje się
reprezentacje grup warkoczowych, a funkcjonały śladowe
po odpowiedniej normalizacji dają niezmienniki domknięć
warkoczy.

\begin{definicja}[Algebra Temperleya--Lieba]
Niech $R$ będzie przemiennym pierścieniem z jedynką,
a $\delta\in R$. \emph{Algebrą Temperleya--Lieba}
$TL_n(\delta)$ nazywamy unitalną $R$-algebrę generowaną
przez $1,e_1,\dots,e_{n-1}$, z relacjami
$e_i^2=\delta e_i$,
$e_ie_{i\pm1}e_i=e_i$ dla dopuszczalnych indeksów
i $e_ie_j=e_je_i$ dla $|i-j|\geq2$.
\end{definicja}

\begin{stwierdzenie}[Reprezentacja warkoczy]
Niech $R=\mathbb Z[A,A^{-1}]$ i $\delta=-A^2-A^{-2}$.
Wtedy elementy $b_i:=A\,1+A^{-1}e_i\in TL_n(\delta)$
są odwracalne i określają homomorfizm
$\rho\colon B_n\to TL_n(\delta)^\times$ przez
$\rho(\sigma_i)=b_i$.
\end{stwierdzenie}

\begin{proof}
Ponieważ $e_i^2=\delta e_i$, mamy
$(A1+A^{-1}e_i)(A^{-1}1+Ae_i)
=1+(A^2+A^{-2}+\delta)e_i=1$.
Pozostałe relacje grupy warkoczowej wynikają po
rozwinięciu odpowiednich iloczynów i zastosowaniu
$e_ie_{i+1}e_i=e_i$ oraz komutacji generatorów
z rozłącznymi indeksami. W relacji z trzema pasmami
jedynymi nietrywialnymi składnikami są wyrazy
z $e_ie_{i+1}e_i$ i $e_{i+1}e_ie_{i+1}$, które
oba redukują się do odpowiednich generatorów.
\end{proof}

\begin{przyklad}[Weryfikacja relacji Artina w $TL_3$]
Niech $e=e_i$, $f=e_{i+1}$, $a=A$, $b=A^{-1}$ oraz
$\delta=-a^2-b^2$. Po rozwinięciu iloczynów
$(a+be)(a+bf)(a+be)$ i $(a+bf)(a+be)(a+bf)$
części asymetryczne różnią się o
\[
\bigl(a^2b+ab^2\delta+b^3\bigr)(e-f).
\]
Ponieważ $ab=1$, współczynnik ten wynosi
$a+b\delta+b^3=a+b(-a^2-b^2)+b^3=0$.
Wyrazy zawierające $efe$ i $fef$ redukują się za pomocą
relacji Temperleya--Lieba. Oba iloczyny są zatem równe.
Jest to bezpośrednie rachunkowe uzasadnienie homomorfizmu
$B_n\to TL_n(\delta)^\times$.
\end{przyklad}

\begin{uwaga}[Chern--Simons i anyony]
W kwantowej teorii pola Chern--Simonsa wartości średnie
operatorów Wilsona związanych z zamkniętymi krzywymi
prowadzą do kwantowych niezmienników splotów. W modelach
topologicznych obliczeń kwantowych warkocze trajektorii
nieabelowych anyonów działają poprzez reprezentacje
grup warkoczowych na przestrzeni stanów.
Nie należy utożsamiać matematycznego istnienia takiej
reprezentacji z fizyczną realizacją niezawodnego
komputera kwantowego: ta ostatnia wymaga dodatkowych
warunków eksperymentalnych i kontroli dekoherencji.
\end{uwaga}

\section{Biologia, chemia i dynamika płynów}
\label{sec:kt-zastosowania}

\begin{przyklad}[Koliste DNA]
Zamknięta cząsteczka dwuniciowego DNA może tworzyć
węzły i sploty. Enzymy topoizomerazy zmieniają jej
topologię, dokonując kontrolowanego przecięcia,
przełożenia i ponownego połączenia nici. W doświadczeniach
z rekombinacją DNA rozkład typów węzłów produktów
może stanowić informację o mechanizmie działania enzymu.
Ważne zastrzeżenie: rzeczywista cząsteczka ma grubość,
energię sprężystą i może ulegać lokalnym zmianom topologii;
model węzła dotyczy konfiguracji bez przecięć.
\end{przyklad}

\begin{przyklad}[Helikalność pola wirowości]
Niech $\Omega\subset\mathbb R^3$ będzie obszarem,
a $u\colon\Omega\to\mathbb R^3$ gładkim polem prędkości
nieściśliwego płynu, dla którego
$\nabla\cdot u=0$. Wirowością nazywamy
$\omega:=\nabla\times u$.
\emph{Helikalność} ma postać
$\mathcal H(u):=\int_\Omega u(x)\cdot\omega(x)\,
\mathrm d\lambda_3(x)$, jeśli całka istnieje.
Gdy wirowość koncentruje się w rozłącznych, cienkich
rurkach wirowych o strumieniach $\Phi_i$ i osiach
$K_i$, wkład wzajemnego splecenia w przybliżeniu
ma postać $2\sum_{i<j}\Phi_i\Phi_j\operatorname{lk}(K_i,K_j)$.
Występują także wkłady samosplecenia i skręcenia
poszczególnych rurek. Pokazuje to, że helikalność jest
geometrycznym odpowiednikiem informacji o spleceniu
linii pola, przy odpowiednich warunkach brzegowych.
\end{przyklad}

\begin{uwaga}[Orbity okresowe pól wektorowych]
Jeśli $X$ jest gładkim polem wektorowym na orientowanej
trójrozmaitości $M$, a $\gamma\colon S^1\hookrightarrow M$
jest prostą zamkniętą orbitą okresową, to jej obraz
jest węzłem w $M$. Typ węzła stanowi niezmiennik
geometryczny orbity względem izotopii układu.
W dynamice hiperbolicznej opisuje się niekiedy całe rodziny
orbit za pomocą rozgałęzionych powierzchni zwanych
\emph{szablonami} (templates), co prowadzi do badań
węzłów realizowanych przez przepływ.
\end{uwaga}

\section{Rozpoznawanie węzłów i algorytmy}
\label{sec:kt-algorytmy}

Teoria węzłów łączy głęboką topologię z obliczeniami:
konkretne pytania o równoważność węzłów można kodować
skończonymi diagramami. Niezmienniki, takie jak kolorowania
Foxa, wyznaczniki i wielomiany, są skutecznymi testami
\emph{nierównoważności}: jeśli dwa węzły mają różne wartości,
na pewno nie są równoważne. Zwykle nie są jednak pełnymi
testami równoważności.

\begin{stwierdzenie}
Jeżeli dla dwóch węzłów $K_1,K_2\subset S^3$ zachodzi
$\Delta_{K_1}(t)\not\doteq\Delta_{K_2}(t)$, to węzły
nie są równoważne. Analogiczny wniosek zachodzi,
gdy $V_{K_1}\ne V_{K_2}$ lub
$\operatorname{Kh}(K_1)\not\cong\operatorname{Kh}(K_2)$
jako odpowiednio dwugradowane grupy.
\end{stwierdzenie}

\begin{proof}
Wszystkie wymienione obiekty są niezmiennikami węzła.
Równoważność wymuszałaby równość wartości niezmienników,
więc ich różność wyklucza równoważność.
\end{proof}

\begin{uwaga}
Problem rozpoznawania węzła trywialnego jest rozstrzygalny
algorytmicznie metodami teorii powierzchni normalnych
w trójrozmaitościach. Istnieją także algorytmy do porównywania
topologii uzupełnień węzłów, a w przypadku hiperbolicznym
istotną rolę odgrywają triangulacje, struktury geometryczne
i obliczenia na grupach fundamentalnych.
Dla praktyki ważne jest rozróżnienie między istnieniem
algorytmu a jego efektywnością na dużych diagramach.
\end{uwaga}



\chapter{Literatura i podsumowanie}
\label{ch:kt-literatura-podsumowanie}

Wcześniejsze wydanie skryptu zawierało osobny rozdział ze szczegółowymi rachunkami. W bieżącej wersji obliczenia zostały wplecione bezpośrednio w odpowiednie części wykładu: grupa trójlistnika pojawia się przy prezentacji Wirtingera, rachunki z macierzy Seiferta przy wielomianie Alexandera i sygnaturze, obliczenie wielomianu Jonesa w sekcji o nawiasie Kauffmana, przykład przestrzeni soczewkowych przy chirurgii Dehna, a obliczenia formy przecięcia przy uchwytach w wymiarze cztery. Dzięki temu aparat obliczeniowy występuje dokładnie tam, gdzie jest używany.

\section{Literatura podstawowa i zaawansowana}
\label{sec:kt-literatura}

\begin{enumerate}[label={[\arabic*]},leftmargin=*,itemsep=3pt]
\item D. Rolfsen, \emph{Knots and Links}, AMS Chelsea Publishing.
Wprowadzenie klasyczne: diagramy, grupy, wielomiany i sploty.
\item C. Livingston, \emph{Knot Theory}, Carus Mathematical Monographs.
Wykład z naciskiem na rachunki i powierzchnie Seiferta.
\item W. B. R. Lickorish, \emph{An Introduction to Knot Theory}, Springer.
Wielomiany Jonesa, skein i chirurgia Dehna.
\item G. Burde, H. Zieschang, M. Heusener, \emph{Knots}, De Gruyter.
Grupy, rozwarstwienia i klasyczna topologia uzupełnień.
\item J. Hempel, \emph{3-Manifolds}, AMS Chelsea.
Tło dla powierzchni nieściśliwych i topologii trójrozmaitości.
\item W. P. Thurston, \emph{The Geometry and Topology of Three-Manifolds}.
Geometria hiperboliczna i wypełnienia Dehna.
\item R. E. Gompf, A. I. Stipsicz, \emph{4-Manifolds and Kirby Calculus}, AMS.
Podstawowa monografia o uchwytach, chirurgii i formach przecięcia.
\item M. H. Freedman, F. Quinn, \emph{Topology of 4-Manifolds}, Princeton.
Topologiczna kategoria wymiaru cztery i zanurzanie dysków.
\item P. Ozsváth, Z. Szabó, \emph{Holomorphic disks and knot invariants}.
Konstrukcja homologii Floera węzłów.
\item M. Khovanov, \emph{A categorification of the Jones polynomial}.
Pierwotna konstrukcja homologii Khovanova.
\item J. Rasmussen, \emph{Khovanov homology and the slice genus}.
Niezmiennik $s$ i rozwiązanie problemu Milnora metodami homologii.
\item P. Kronheimer, T. Mrowka, \emph{Khovanov homology is an unknot-detector}.
Detekcja węzła trywialnego przy użyciu teorii instantonów.
\item P. Orlik, \emph{Seifert Manifolds}, Lecture Notes in Mathematics 291.
Klasyczne rozwarstwienia okręgowe, grupy i ich klasyfikacja.
\item W. D. Neumann, \emph{A calculus for plumbing applied to the topology of complex surface singularities}.
Rachunek ważonych grafów, ich operacji i brzegów.
\item M. Hedden, \emph{Knot Floer homology of Whitehead doubles}.
Obliczenia $\tau$ i rozróżnianie slicowości gładkiej i topologicznej.
\item R. Fintushel, R. J. Stern, \emph{Rational Blowdowns of Smooth
$4$-Manifolds}, J. Differential Geom. 46 (1997), 181--235.
Wprowadzenie wymiernej kontrakcji i formuły dla niezmienników gładkich.
\item A. Casson, J. Harer, \emph{Some homology lens spaces which bound
rational homology balls}, Pacific J. Math. 96 (1981), 23--36.
Konstrukcje wymiernych kul homologii.
\item J. Park, \emph{Simply connected symplectic $4$-manifolds with $b_2^+=1$ and $c_1^2=2$}, Invent. Math. 159 (2005), 657--667. Konstrukcja $C_7$ w $E(1)\#4\overline{\mathbb{CP}}^{\,2}$ oraz argument prostej spójności dopełnienia. \url{https://doi.org/10.1007/s00222-004-0404-1}
\item J. Park, A. I. Stipsicz, Z. Szabó, \emph{Exotic smooth structures
on $\mathbb{CP}^2\#5\overline{\mathbb{CP}}^{\,2}$},
Math. Res. Lett. 12 (2005), 701--712.
Przykłady uogólnionych wymiernych kontrakcji.
\item J. Park, \emph{Seiberg--Witten invariants of generalised rational
blow-downs}, Bull. Austral. Math. Soc. 56 (1997), 363--384.
Zależności niezmienników od warunków przedłużenia klas.
\item R. Fintushel, R. J. Stern, \emph{Six Lectures on Four-Manifolds},
IAS/Park City Math. Ser. 15 (2009), 265--315.
Wprowadzenie do niezmienników Seiberga--Wittena i racjonalnych kontrakcji.
\item J. W. Morgan, \emph{The Seiberg--Witten Equations and Applications
to the Topology of Smooth Four-Manifolds}, Princeton University Press.
Analityczne podstawy teorii cechowania i konstrukcja niezmienników.
\item R. Fintushel, R. J. Stern, \emph{Knots, Links, and
$4$-Manifolds}, Invent. Math. 134 (1998), 363--400.
Oryginalne twierdzenie o chirurgii wzdłuż węzła, formuła
Seiberga--Wittena i konstrukcje egzotycznych rozmaitości K3.
\url{https://arxiv.org/abs/dg-ga/9612014}
\item R. Fintushel, R. J. Stern, \emph{Constructions of Smooth
$4$-Manifolds}, Documenta Math. Extra Volume ICM 1998.
Dodatkowe objaśnienia formuły chirurgii wzdłuż węzła.
\item T.-J. Li, A.-K. Liu, \emph{General wall crossing formula},
Math. Res. Lett. 2 (1995), 797--810.
Klasyczna formuła zmiany komory dla $b_2^+=1$.
\item L. I. Nicolaescu, \emph{Notes on Seiberg--Witten Theory},
Graduate Studies in Mathematics, AMS. Oszacowania,
przestrzenie moduli, twierdzenia o zwartości i zmiana komory.
\end{enumerate}

\section*{Podsumowanie zależności}
Najważniejsze przejścia, które warto zachować przy ponownym
wykorzystaniu skryptu w rozdziałach o wymiarze cztery, są następujące:
\begin{enumerate}[label=(\roman*)]
\item diagram węzła $\longrightarrow$ prezentacja Wirtingera,
powierzchnia Seiferta, macierz Seiferta i wielomian Alexandera;
\item diagram zorientowany $\longrightarrow$ nawias Kauffmana,
wielomian Jonesa oraz kompleks Khovanova;
\item uzupełnienie węzła $\longrightarrow$ torus peryferyjny,
nachylenia, chirurgia Dehna i nowe zamknięte trójrozmaitości;
\item ramowany splot $\longrightarrow$ uchwyty indeksu dwa,
forma przecięcia i rachunek Kirby'ego w wymiarze cztery;
\item dysk w $B^4$ $\longrightarrow$ gładka i topologiczna
konkordancja, twierdzenie Freedmana oraz niezmienniki
$\tau$ i $s$ rozróżniające kategorie TOP i DIFF;
\item włókno torusowe $T\subset E(2)$ oraz węzeł $K\subset S^3$
$\longrightarrow$ chirurgia Fintushela--Sterna i formuła
$\mathcal{SW}_{E(2)_K}=\Delta_K(\exp(2[T_K]))$,
prowadzące do egzotycznych struktur na K3.
\end{enumerate}

% Koniec tresci w wersji zintegrowanej. Przy wstawieniu do wiekszej ksiazki
% rozdzialy sa automatycznie numerowane zgodnie z jej struktura.


\end{document}
